\documentclass[11pt,letterpaper]{article}
\usepackage{balloons-style}
\hypersetup{pdftitle={The balloon process on the circle},
 pdfauthor={AI-generated working proofs prepared for Jacob Richey}}
\makeatletter
\renewcommand*\l@section{\@dottedtocline{1}{0em}{2.3em}}
\makeatother
\newcommand{\PP}{\mathbb P}\newcommand{\EE}{\mathbb E}
\newcommand{\T}{\mathbb R/\mathbb Z}\newcommand{\Leb}{\operatorname{Leb}}
\newcommand{\ind}{\mathbf 1}\newcommand{\Exp}{\operatorname{Exp}}
\newcommand{\Beta}{\operatorname{Beta}}\newcommand{\GammaD}{\operatorname{Gamma}}
\newcommand{\status}[1]{\textbf{Status: #1.}}
\newcommand{\rf}[2]{(#1)^{\overline{#2}}}
\newcommand{\ff}[2]{(#1)_{\underline{#2}}}
\title{The balloon process on the circle\\\large Technical companion: mathematical arguments}
\author{AI-generated working proofs prepared for Jacob Richey}
\date{20 October 2026}
\begin{document}
\maketitle
\vspace{-1.3em}
\begin{abstract}
Equally growing balloons launched from an odd iid uniform sample on the
unit circle annihilate in pairs. Resampling one point changes the final
survivor with probability $p_N$. We prove a uniform positive lower bound
and convergence $p_N\to p_\infty$. The limit lies between exact integral
constants with certified decimal bounds $0.2016982228$ and
$0.8112015673$. The proof combines
insertion maps, an explicit last-pair law, a positive renewal
representation, and a contraction for backward mass queries. We give
finite-tree approximation and finite-sample error bounds, exact
available-set moments, and reproducible checks. The proposed value near
$0.658$ remains uncertified. We also obtain an odd-sample random
counting-measure process, its reverse fragmentation rule, an explicit
last-collision law, and deterministic empirical limits at every vanishing
cutoff, including the full intermediate-time regime. At the small-cutoff
end of the entrance process, the centered counting measure has Gaussian
white-noise fluctuations and its normalized version has a Gaussian bridge
limit. These extend to temporal Gaussian fields and cumulative spatial
path limits throughout the finite mesoscopic regime, with centering
$Nq(Nrt)$. Their covariance is an explicit renewal tail, and the
log-time normalized Gaussian count limit is not Markov.
Replacing the finite-intensity centering by $c/(rt)$ requires a further condition.
A factorial renewal estimate gives certified early-label error bounds.
For perturbations of several starting points, we prove contraction of
the active-set discrepancy and an alternating-path description of the
changed matching. Exact spectral formulas organize the full resampling
profile. A separated-insertion argument proves unique joint limits with
shared identities, convergence of the full spectrum, and a stationary
refresh-time evolution. The limiting spectrum has no mass at zero:
any growing number of resamplings changes the survivor with probability
tending to one. Uniformly over sizes $1\le k\le N$, the expected
number of changed matching edges has order $k(1+\log(N/k))$.
For $k=o(N)$ its leading mean is
$(1+e^{-\gamma})k\log(N/k)$, and the history cost concentrates
relative to its mean uniformly over every nonzero resampling count.
Fixed edit families have joint Gaussian functional limits on logarithmic
separation scales, with a strictly positive variance coefficient.
For every growing sparse edit family, centered cumulative histories
have a Brownian limit up to the reciprocal edit cutoff and remain
constant at that fluctuation scale through completion.
For every growing $k=o(N)$, the completed history centered by its
deterministic mean profile has a Gaussian limit with variance asymptotic
$8kv_{\rm H}\log(N/k)$, for both spatial and label costs.
A uniform late-band locality bound controls the part after cutoff $1/k$.
At positive edit density $k/N\to\theta\in(0,1)$, the history centered
by its limiting mean profile has a $\sqrt N$ Gaussian limit with positive continuous
variance $v_{\rm D}(\theta)$. This variance is a conditioned
Poisson block limit and has small-density asymptotic
$8v_{\rm H}\theta\log(1/\theta)$.
Jointly over all edit fractions, the profile-centered history
converges to a continuous Gaussian process pinned at zero at
both endpoints, with covariance given by conditioned block limits.
Its small-edit tangent is Brownian motion on the scale
$\sqrt{8v_{\rm H}\delta\log(1/\delta)}$.
At the other endpoint, scaling by $\sqrt{2\pi_2}\delta$
gives the tangent $B(u^2)$ in finite-dimensional laws and $L^2$.
Near complete refresh, with $M=N-k=o(N)$, shared spatial edges
have a Poisson limit when $M^2/N$ has a finite limit and a Gaussian
limit when it diverges. The Poisson rate coefficient is the
certified constant $\kappa=\pi_2/2$; label histories have an
additional independent mean-$1/2$ Poisson overlap in the bounded
regime. The dense variance is asymptotic to $2\pi_2(1-\theta)^2$
near one. Profile centering holds throughout the Gaussian endpoint;
the quadratic center $N-1-\pi_2M^2/N$ gives a centered normal
limit precisely when $M=o(N^{3/4})$ in that regime.
Under rare refreshes the centered history has a compound Poisson limit.
At positive edit density it has a deterministic profile with an
integer-valued matching spectrum, distinct from the survivor spectrum.
Every coefficient of this matching spectrum is strictly positive,
with a squared Poisson-difference representation. Consequently every
fixed spectral truncation has a sharp Gaussian-centering window.
Finite mixtures of insertions and uniform deletions have the same
limiting response law when their total edit count agrees.
The limiting insertion basin measures are purely atomic, and their
ranked weights and normalized deletion indegrees converge in $\ell^1$
to the same random probability vector.
Their tree-depth approximation error decays exponentially, giving
power bounds for spatial noise and total early basin mass. Separate
optimization of the empirical interval discrepancy sharpens the
deletion-basin convergence rate.
The complete positive-cutoff history supports a reversible Markov
refresh rule built from commuting insertion involutions. Essential
continuity at protected endpoints and an adaptive invariant insertion
rule exclude atoms from the conditional survivor law. More generally,
growing resampling makes the active configurations disjoint at any fixed
positive cutoff. Independence of those configurations is not asserted.
The proofs retain their working status pending independent human review.
\end{abstract}

\textbf{Provenance and authority.} This is an AI-written synthesis of the
research discussion with Jacob Richey, prepared at his request. It is not
Jacob's verbatim writing or evidence of his independent verification.
Statements labelled \textsc{proved} have arguments included here; this is a
working research record, not a refereed manuscript. Numerical reports from
earlier sessions lack an archived reproducibility package and are labelled
accordingly. The new finite-recursion checks have archived source, seeds,
and output; they support the implementation, not an asymptotic theorem.
No claim of literature novelty is made.

{\small\setlength{\parskip}{0pt}\setcounter{tocdepth}{1}\tableofcontents}
\clearpage

\section{Problem, conventions, and current status}\label{sec:problem}

Work on $\T$ with circumference one and intrinsic distance
$d(x,y)=\min\{|x-y|,1-|x-y|\}$. A point $x$ has a balloon of radius $t$ at
time $t$, until it first meets another surviving balloon. Both then disappear.
Pair distances are distinct almost surely for independent uniform points;
all statements below disregard null sets of ties.

The process is equivalently the greedy matching: repeatedly remove the
closest pair of remaining points. A pair at separation $r$ dies at time
$r/2$. \emph{All cutoffs below use separation $r$, not physical time.}
With $N=2n+1$ points there is exactly one survivor. With $M=2n$ points the
matching is perfect.

Let $B=\{X_1,\ldots,X_M\}$ be the common background, and let $U,V$ be
independent uniform insertions, independent of $B$. Let $S_B(z)$ denote the
survivor's position in $B\cup\{z\}$. The inserted point has the same
\emph{label} $*$ in the two systems, even though its position changes.
Write
\[
 p_N=\PP(\text{survivor labels in }B+U\text{ and }B+V\text{ differ}),\qquad
 p_N^{\rm pos}=\PP(S_B(U)\ne S_B(V)).
\]
The user's question is whether $\liminf_{N\to\infty,\ N\ \mathrm{odd}}p_N>0$;
the answer below also implies the statement for positions. This is the
specified one-coordinate noise question; we do not identify it with every
standard Boolean-function definition of noise sensitivity.

For $v\in B$, define the basin mass and available mass
\[
 \theta_v(B)=\Leb\{z:S_B(z)=v\},\qquad
 E_B=\{z:S_B(z)=z\},\qquad \Theta_N=\Leb(E_B).
\]
These masses sum to one. Conditional independence of $U,V$ gives
\begin{equation}\label{eq:noise-identity}
 p_N=1-\EE\sum_{v\in B}\theta_v^2-\EE\Theta_N^2,
 \qquad p_N^{\rm pos}=p_N+\EE\Theta_N^2.
\end{equation}
Exchangeability gives $\EE\Theta_N=1/N$.

\begin{center}
\begin{tabular}{p{.54\textwidth}p{.36\textwidth}}
\toprule
Claim & Current status\\\midrule
Uniform positive lower bound for $p_N$ & Proved, Section~\ref{sec:lower}\\
Available-set correlations and all integer moments & Proved, Section~\ref{sec:available}\\
Conditional Beta law and component generating function & Proved, Section~\ref{sec:components}\\
Exact finite insertion-map reduction & Proved, Section~\ref{sec:maps}\\
Closed finite spatial moment recursion and agreement readout & Proved, Section~\ref{sec:maps}\\
Deletion-graph reformulation of agreement & Proved, Section~\ref{sec:matching}\\
Full last-pair limit and stronger noise lower bound & Proved, Section~\ref{sec:scalar-limit}\\
Active configuration law and early-label agreement bound & Proved, Section~\ref{sec:active-limit}\\
Convergence of $p_N$ and marked weights at each positive cutoff & Proved, Section~\ref{sec:marked-limit}\\
Active-location process, reverse rule, and all vanishing-cutoff empirical limits & Proved, Section~\ref{sec:measure-process}\\
Spatial Gaussian limits throughout the finite mesoscopic regime & Proved, Sections~\ref{sec:measure-process}, \ref{sec:full-finite-spatial-CLT}\\
Temporal Gaussian fields, cumulative path limits, and explicit covariance & Proved, Section~\ref{sec:temporal-spatial-CLT}\\
Improved noise lower bound from half-interval masses & Proved, Section~\ref{sec:sharper-budgets}\\
Several-point coupling and finite noise profiles & Proved, Section~\ref{sec:perturbations}\\
Unique coupled limits and the limiting noise spectrum & Proved, Section~\ref{sec:coupled-limit}\\
Expected history scale $k(1+\log(N/k))$, uniformly in $k$ & Proved, Section~\ref{sec:growing-perturbations}\\
Edit-type universality and atomic limiting basin weights & Proved, Section~\ref{sec:edit-universality}\\
Spatial noise profile and power bounds for small moves & Proved, Section~\ref{sec:spatial-noise}\\
Reversible Markov rule on complete entrance histories & Proved, Section~\ref{sec:history-markov}\\
No zero spectral atom; disjointness at fixed cutoffs under growing resampling & Proved, Section~\ref{sec:history-markov}\\
Polynomial finite-sample noise rate, with sharper exponent & Proved, Sections~\ref{sec:full-tree-rate}, \ref{sec:complex-coefficient-rate}\\
Logarithmic spectral tails and shrinking-cutoff entrance overlap & Proved, Section~\ref{sec:spectral-rate}\\
Uniform finite-sample logarithmic noise and overlap rates & Proved, Section~\ref{sec:finite-family-rate}\\
Sharp mean history coefficient for every $k=o(N)$ & Proved, Section~\ref{sec:sublinear-history}\\
Dense-edit profile and strictly positive integer matching spectrum & Proved, Sections~\ref{sec:dense-history}, \ref{sec:matching-Poisson-spectrum}\\
Relative history concentration uniformly over all nonzero counts & Proved, Section~\ref{sec:fixed-history-LLN}\\
Terminal laws, sparse sheets, Gaussian tangents, and explicit count correction & Proved, Sections~\ref{sec:late-history-variance}--\ref{sec:dense-count-correction}\\
Certified value near $0.658$ and useful sharp rate & Open\\
Earlier small-$N$ rational integrations and simulations & Reported evidence; see Section~\ref{sec:ledger}\\
\bottomrule
\end{tabular}
\end{center}

\section{Three deterministic facts}\label{sec:deterministic}

The closest pair on the circle is adjacent in cyclic order. If its current
gap is $m$ and the neighbouring gaps are $a,c$, deleting its endpoints
replaces the three gaps by $a+m+c$. The last two points require the smaller
of their two circular gaps. Normalized independent exponential spacings
have the same law as the circular gaps of independent uniform points,
rooted at a distinguished sampled point.

\begin{lemma}[Stability characterization]\label{lem:stable}
For an even finite set with distinct pair distances, the greedy matching is
the unique perfect matching with no blocking pair: there are no unmatched-to-
each-other vertices $u,v$ with
$d(u,v)<\min\{r_u,r_v\}$, where $r_u,r_v$ are their matching distances.
\end{lemma}
\begin{proof}
The globally closest pair must be matched in any stable perfect matching,
since otherwise it blocks. Remove it and iterate. The greedy construction
has no blocking pair: when the earlier of two proposed blocking vertices
is removed, the other is still present and cannot be strictly closer than
the partner chosen then.
\end{proof}

\begin{lemma}[A single moving discrepancy]\label{lem:discrepancy}
Couple the processes on $B$ and $B+z$ by scanning all pair events in order
of increasing distance. Their active sets have symmetric difference exactly
one until both have reached their terminal states. When an event involves
the unique extra point and a common active point, the discrepancy transfers
to that common point and changes which system contains the extra point.
\end{lemma}
\begin{proof}
A common pair disappears in both systems. A pair involving the extra point
can disappear only in its system; the other endpoint then becomes extra in
the other system. These are the only possibilities.
\end{proof}

\begin{lemma}[Shielding an interval]\label{lem:shield}
Put $k$ independent uniform interior points in an interval of length $x$.
Keep balloons at the two endpoints, and ignore their possible meeting with
each other. The probability that all interior points disappear without any
meeting involving an endpoint is zero for odd $k$, and $1/(k+1)$ for even
$k$. On this event all interior matching distances are at most $x/3$.
\end{lemma}
\begin{proof}
Identify the endpoints to obtain a circle of circumference $x$ with a
distinguished point. For even $k$, endpoint shielding is equivalent to this
point being the survivor. A proposed pair using the shortcut through the
distinguished point would have met that point earlier, so it cannot occur
while that point survives. This proves the equivalence with the interval
process. By rotating $k+1$ independent uniform circle points to put the
distinguished one at zero and using exchangeability, its survival probability
is $1/(k+1)$. While an odd circle process has at least three active points,
its shortest gap is at most $x/3$, proving the distance cutoff. Odd $k$
cannot clear by internal pairs.
\end{proof}

\section{A complete positive lower bound}\label{sec:lower}

Let $a,b$ be the last pair of the even background $B$, and put
$L=d(a,b)\in(0,1/2)$. Its two arcs have lengths $L$ and $1-L$.

\begin{lemma}[Last-pair comparison]\label{lem:last-comparison}
For $N=M+1$ odd,
\begin{equation}\label{eq:comparison}
 p_N\ge\EE[L(1-L)].
\end{equation}
More generally, for $0<h\le1/4$,
\begin{equation}\label{eq:spatial}
 \PP\bigl(d(S_B(U),S_B(V))\ge h\bigr)
 \ge\EE\bigl[L(1-L)\ind_{\{L\ge2h\}}\bigr].
\end{equation}
\end{lemma}
\begin{proof}
Until the discrepancy from Lemma~\ref{lem:discrepancy} reaches $a$ or $b$,
it stays in its insertion arc. It cannot jump across a surviving endpoint:
such a pair would have met that endpoint earlier. If the augmented system
removes $a$ first, the base-only discrepancy at $a$ persists until the last
base collision $a$--$b$; consequently $b$ is the augmented survivor. The
analogous statement holds with $a,b$ exchanged.

Any insertion in an arc of length less than $2L$ must therefore yield
$a$ or $b$. Otherwise, immediately before the last base collision, that
arc would contain an extra active point along with $a,b$. Its distance to
an endpoint is less than $L$, contradicting their survival to that time.
In particular, insertion in the short arc always yields a boundary survivor.

Condition on the last endpoints, their positions, and the allocation of
the remaining labelled points to the two arcs. The last-pair event is
exactly the intersection of two separate events: each arc clears internally
without touching its endpoints and with all internal distances less than
$L$. Necessity is immediate. For sufficiency, a cross-arc collision would
first meet an endpoint, so the two internal clearances run independently.
The conditional interior configurations in the two arcs are consequently
independent. Their individual laws are invariant under reflection of the
arc, including the distance cutoff.

A uniform insertion into the short arc has, after averaging its interior
configuration, equal probabilities of yielding $a$ and $b$. Its outcome
depends only on that arc. The long-arc insertion outcome depends only on
the long arc and is independent of this fair choice. It can equal at most
one of $a,b$, so conditional disagreement has probability at least $1/2$.
The probability that $U,V$ lie in opposite arcs is $2L(1-L)$, proving
\eqref{eq:comparison}. For the spatial statement, any point on the circle
is at distance at least $L/2$ from at least one of $a,b$, by the triangle
inequality. The same fair-choice argument proves \eqref{eq:spatial}.
\end{proof}

\begin{proposition}[An exact part of the last-pair law]\label{prop:last-density}
For even $M\ge2$, the last separation $L_M$ has density
\begin{equation}\label{eq:last-density}
 f_M(r)=\frac{1-(1-2r)^M}{2r(1-r)},\qquad \frac14\le r\le\frac12.
\end{equation}
Consequently, on the same range,
\begin{equation}\label{eq:last-tail}
 \PP(L_M\ge r)=\sum_{j=0}^{M/2-1}\frac{(1-2r)^{2j+1}}{2j+1}.
\end{equation}
\end{proposition}
\begin{proof}
Fix candidate endpoints at separation $r$. The numbers of other points in
their two arcs are $K$ and $M-2-K$, where $K$ is binomial with parameters
$M-2,r$. For $r\ge1/4$, both arc lengths are at most $3r$. Thus the cutoff
in the proof of Lemma~\ref{lem:last-comparison} is automatic on shielding,
by Lemma~\ref{lem:shield}. Conditional on even $K=k$, the candidate-last-pair
probability is $1/((k+1)(M-1-k))$. There are $\binom M2$ candidates and
each candidate separation has density two. The desired density is
\[
 M(M-1)\sum_{\substack{0\le k\le M-2\\ k\ {\rm even}}}
 \binom{M-2}{k}\frac{r^k(1-r)^{M-2-k}}{(k+1)(M-1-k)}.
\]
Use $M(M-1)\binom{M-2}{k}/((k+1)(M-1-k))=\binom M{k+1}$ and sum the odd
terms of the binomial expansion to obtain \eqref{eq:last-density}.
Integrating, or substituting $u=1-2r$, gives \eqref{eq:last-tail}.
\end{proof}

\begin{theorem}[One-point resampling changes the survivor]\label{thm:noise}
For every odd $N\ge3$,
\begin{equation}\label{eq:main-bound}
 \boxed{\displaystyle p_N\ge\frac18-\frac{1}{N2^{N+2}}.}
\end{equation}
The same bound holds for
$\PP(d(S_B(U),S_B(V))\ge1/8)$. In particular,
$\liminf p_N\ge1/8$.
\end{theorem}
\begin{proof}
Combine Lemma~\ref{lem:last-comparison} and Proposition~\ref{prop:last-density}:
\[
 \EE[L_M(1-L_M)\ind_{\{L_M\ge1/4\}}]
 =\frac12\int_{1/4}^{1/2}\bigl(1-(1-2r)^M\bigr)\,dr
 =\frac18-\frac{1}{(M+1)2^{M+3}}.
\]
For the spatial bound use \eqref{eq:spatial} with $h=1/8$.
\end{proof}

There is also a useful upper bound. If $L>1/3$, both arcs have length
less than $2L$, so every insertion yields $a$ or $b$. For a fixed background
the disagreement probability is then at most $1/2$. Therefore
\[
 p_N\le1-\tfrac12\PP(L_{N-1}>1/3),\qquad
 \limsup p_N\le1-\frac{\log2}{4}\approx0.826713.
\]

\section{The available set: exact correlations and moments}\label{sec:available}

For $v\in B$, let $r_v$ be its distance to its partner in the base matching.
The available set has the deterministic description
\begin{equation}\label{eq:available-balls}
 E_B=\{z:d(z,v)>r_v\text{ for every }v\in B\}
 =\T\setminus\bigcup_{v\in B}\overline B(v,r_v),
\end{equation}
up to endpoints. Indeed, if $z$ never meets a base point, the base process
is unchanged. Conversely, prior to the first such meeting the base process
is unchanged, and $d(z,v)>r_v$ precludes a meeting before $v$ dies.
The radius in \eqref{eq:available-balls} is the full matching distance,
twice the balloon's radius at its death.

\begin{lemma}[Several available probes]\label{lem:probes}
Fix distinct probe locations $z_1,\ldots,z_q$. All belong to $E_B$ if and
only if each of their intervening arcs contains an even number of base
points and those points clear without touching the arc endpoints.
\end{lemma}
\begin{proof}
An available probe cannot lie on a shortest matching arc: an endpoint of
that arc would be closer to the probe than to its partner. Thus every
matching pair lies within one probe arc, and the local configuration shields
its endpoints. Conversely, form the shielded matching in each arc. The
distance from each point to either adjacent probe exceeds its matching
distance. A path between different arcs crosses a probe and is longer
than at least one of the proposed pair's existing matching distances; in
fact its first and last probe crossings give the sum of those distances
as a lower bound. It cannot block the union matching. Alternative circular
paths for within-arc pairs also cross a probe and cannot block. Apply
Lemma~\ref{lem:stable}. The $q=1$ case is the endpoint identification in
Lemma~\ref{lem:shield}.
\end{proof}

\begin{proposition}[Two-point function and second moment]\label{prop:two-point}
For $0<x<1$ and odd $N$,
\begin{align}
 K_N(x):=\PP(0,x\in E_B)
 &=\sum_{\substack{0\le k\le N-1\\k\ {\rm even}}}
 \binom{N-1}{k}\frac{x^k(1-x)^{N-1-k}}{(k+1)(N-k)}\notag\\
 &=\frac{1-(1-2x)^{N+1}}{2N(N+1)x(1-x)},\label{eq:two-point}\\
 \EE\Theta_N^2
 &=\frac{2}{N(N+1)}\sum_{j=0}^{(N-1)/2}\frac1{2j+1}\notag\\
 &=\frac{\log N+\gamma+\log2+o(1)}{N^2}.\label{eq:theta-second}
\end{align}
\end{proposition}
\begin{proof}
Condition on the binomial number of base points in the first arc and apply
Lemmas~\ref{lem:probes} and \ref{lem:shield}. The same odd-binomial identity
as above gives the rational form. Translation invariance and Fubini give
$\EE\Theta_N^2=\int_0^1K_N(x)\,dx$. Expand the resulting geometric sum in
$1-2x$ to obtain the harmonic sum. Its asymptotic follows from
$\sum_{j=0}^{n}(2j+1)^{-1}=H_{2n+2}-\tfrac12H_{n+1}$.
\end{proof}

Put
\begin{equation}\label{eq:h}
 h(z)=\frac{\operatorname{arctanh}\sqrt z}{\sqrt z}
 =\sum_{j\ge0}\frac{z^j}{2j+1},\qquad h(0)=1.
\end{equation}
The expression with the power series is the definition near zero.

\begin{theorem}[All available-mass moments]\label{thm:moments}
For $N=2n+1$ and every positive integer $q$,
\begin{equation}\label{eq:moments}
 \EE\Theta_{2n+1}^q
 =\frac{(2n)!(q-1)!}{(2n+q-1)!}[z^n]h(z)^q.
\end{equation}
For fixed $q$ and $N\to\infty$ through odd integers,
\begin{equation}\label{eq:moment-asymptotic}
 \EE\Theta_N^q\sim\frac{q!}{2^{q-1}}
 \frac{(\log N)^{q-1}}{N^q}.
\end{equation}
\end{theorem}
\begin{proof}
Use $q$ independent uniform probes. Rooting cyclic order at the first probe,
their arc lengths have the Dirichlet$(1,\ldots,1)$ law. After integrating
these lengths, the counts $(k_1,\ldots,k_q)$ of $M=2n$ base points are
uniform over all weak compositions of $M$ into $q$ parts: each has
probability $M!(q-1)!/(M+q-1)!$. By Lemma~\ref{lem:probes}, only even counts
contribute, with probability $\prod_i(k_i+1)^{-1}$. Summing gives
\eqref{eq:moments} (for $q=1$ it is simply $1/(2n+1)$).

Near $z=1$, $h(z)=\tfrac12\log(4/(1-z))+
O((1-z)\log(1/(1-z)))$. For fixed $q$, logarithmic coefficient extraction
gives
\[
 [z^n]h(z)^q\sim\frac{q}{2^q n}(\log n)^{q-1}.
\]
One may obtain the leading coefficient by differentiating a fixed power
of $\log(1/(1-z))$ and summing its coefficients; the lower logarithmic
powers contribute one fewer power of $\log n$. Inserting this estimate in
\eqref{eq:moments} proves \eqref{eq:moment-asymptotic}.
\end{proof}

\section{Available components and their geometric representation}\label{sec:components}

Let $C_N$ be the number of connected components of the open set $E_B$.
For $M=2n\ge2$, $C_N\in\{0,\ldots,n\}$ and $C_N=0$ exactly when
$\Theta_N=0$, outside a null set. Every covered block between available
intervals contains an even positive number of points.

\begin{theorem}[Conditional Beta law]\label{thm:beta}
For every $c\ge1$ with $\PP(C_N=c)>0$,
\begin{equation}\label{eq:beta}
 \Theta_N\mid\{C_N=c\}\ \sim\ \Beta(c,M-c),\qquad M=N-1.
\end{equation}
\end{theorem}
\begin{proof}
Represent the unnormalized circular gaps by independent rate-one
exponentials. Fix the set $I$ of $c$ original gaps containing available
intervals. Cutting at these intervals partitions the points into $c$
consecutive blocks. Within each block the restriction of the global matching
is the unique stable matching for the line metric: an internal blocking
pair would also block on the circle. It therefore depends only on the
$M-c$ internal gaps.

For a block with first and last points $x_{\rm first},x_{\rm last}$, put
\[
 L=\min_v(x_v-r_v),\quad R=\max_v(x_v+r_v),\quad
 \beta=x_{\rm first}-L,\quad\alpha=R-x_{\rm last}.
\]
Its forbidden intervals have connected union $[L,R]$, since otherwise
there would be an additional available component inside the block.
Conversely, take any internal block configurations whose forbidden unions
are connected. Separate successive blocks by gaps
\[
 g_i=\alpha_{\text{left block}}+\beta_{\text{right block}}+s_i,
 \qquad i\in I,\quad s_i>0.
\]
Their internal matchings remain the global stable matching. Any competing
path that crosses an available interval has length at least $r_u+r_v$
for its endpoint vertices $u,v$, so cannot block. This also excludes a
shorter circular route for an internal matched pair. Thus opening or
compressing the positive $s_i$ preserves the block matchings and creates
exactly the chosen available intervals.

This is a bijective parametrization by the internal gaps and the positive
slacks. Its Jacobian is one: the unchanged internal coordinates are followed
by coordinates $g_i=s_i+\alpha+\beta$. The possible internal coordinates
form a scale-invariant set $D_I\subset(0,\infty)^{M-c}$; matching changes
partition it into finitely many regions, and all formulas are valid almost
everywhere. The remaining circumference
\[
 W=\sum_{\text{blocks}}(R-L)
\]
is positively homogeneous of degree one in these coordinates. The full
circumference is $W+\sum_i s_i$, so the exponential density factors as
$e^{-W}\prod_i e^{-s_i}$. Since
$\operatorname{vol}\{y\in D_I:W(y)\le t\}=t^{M-c}
\operatorname{vol}\{y\in D_I:W(y)\le1\}$,
the conditional law of $W$ is $\GammaD(M-c,1)$. Independently,
$\sum_i s_i\sim\GammaD(c,1)$. Normalizing the circumference proves
\eqref{eq:beta}. Mixing over the possible $I$ leaves this law unchanged.
\end{proof}

This proof supplies the geometric product decomposition that was omitted
in the earlier conversational argument. Homogeneity alone would not have
proved the assertion without the matching-preserving parametrization.

\begin{corollary}[Component generating function]\label{cor:pgf}
For $n\ge1$,
\begin{align}
 \EE u^{C_{2n+1}}
 &=1-2n[z^n]\log\bigl(u+(1-u)h(z)\bigr),\label{eq:pgf}\\
 \EE C_N&=\frac{N-1}{N},\label{eq:mean-components}\\
 \EE\ff{C_{2n+1}}q
 &=2n(q-1)![z^n](h(z)-1)^q.\label{eq:falling}
\end{align}
\end{corollary}
\begin{proof}
The Beta moment formula and Theorem~\ref{thm:moments} give
$\EE\rf{C}q=M(q-1)![z^n]h(z)^q$. Hence, initially for sufficiently small
$t$,
\[
 \EE(1-t)^{-C}=1-M[z^n]\log(1-th(z)).
\]
Set $u=(1-t)^{-1}$ and use $[z^n]\log u=0$ to get \eqref{eq:pgf}.
Both sides are polynomials in $u$, so the identity extends to all $u$.
Differentiation at $u=1$ proves the remaining statements. Independently,
$-K_N'(0+)=\EE C_N$ follows by differentiating the overlap length of a finite
union of intervals, and agrees with \eqref{eq:mean-components}.
\end{proof}

For example, the distributions $(\PP(C=0),\PP(C=1),\ldots)$ are
\[
 \begin{array}{c|l}
 M=2 &(1/3,\;2/3)\\
 M=4 &(19/45,\;16/45,\;2/9)\\
 M=6 &(443/945,\;88/315,\;8/45,\;2/27).
 \end{array}
\]
In particular, $\PP(C_{2n+1}=n)=2/3^n$.

\subsection{Rare availability and its conditional scale}

Set $a_n=\PP(\Theta_{2n+1}>0)$ for $n\ge1$, and $a_0=1$. Equation
\eqref{eq:pgf}, including its value at $u=0$, gives
\begin{equation}\label{eq:a-generating}
 a_n=2n[z^n]\log h(z),\qquad
 \sum_{n\ge0}a_nz^n=\frac{1}{(1-z)h(z)}.
\end{equation}
Here the second identity follows from
$2zh'(z)=1/(1-z)-h(z)$. This derivation does \emph{not} analytically
continue an integer-moment identity to $q=0$.
Equivalently,
\begin{equation}\label{eq:renewal}
 a_n=\sum_{j=1}^n\frac{2}{(2j-1)(2j+1)}a_{n-j}.
\end{equation}
Indeed $(1-z)h(z)=1-\sum_{j\ge1}2z^j/((2j-1)(2j+1))$.

An integral representation useful for asymptotics is
\begin{equation}\label{eq:a-integral}
 a_n=\int_0^\infty \tanh^{2n}(v/2)\frac{2\,dv}{v^2+\pi^2}.
\end{equation}
For completeness, apply the Cauchy integral representation to
$A(z)=1/((1-z)h(z))$ on the slit plane $\mathbb C\setminus[1,\infty)$.
The integral $h(z)=\int_0^1(1-zt^2)^{-1}dt$ has positive imaginary part in
the upper half-plane and is positive on $(-\infty,1)$, so has no zeros
there. For $x>1$, write $b=\log((\sqrt x+1)/(\sqrt x-1))$; the upper boundary
value of $h$ is $(b+i\pi)/(2\sqrt x)$. Taking the jump of $A$ yields
\[
 A(z)=\int_0^1\frac{1}{1-zt},
 \frac{2\,dt}{\sqrt t(1-t)
 [\log^2((1+\sqrt t)/(1-\sqrt t))+\pi^2]}.
\]
The contour at infinity vanishes since $A(z)=O(|z|^{-1/2})$, and the
small contour at $1$ contributes $O(1/|\log\varepsilon|)$.
Substitute $\sqrt t=\tanh(v/2)$ and extract coefficients to obtain
\eqref{eq:a-integral}. Its integrand changes from nearly zero to nearly one
around $v=\log n$. Bounds at $(1\pm\varepsilon)\log n$, followed by
$\varepsilon\downarrow0$, give
\begin{equation}\label{eq:positive-asymptotic}
 \PP(\Theta_N>0)\sim\frac{2}{\log N}.
\end{equation}
Together with \eqref{eq:moment-asymptotic}, this proves by moments
\begin{equation}\label{eq:conditional-exp}
 \left.\frac{2N\Theta_N}{\log N}\ \right|\ \Theta_N>0
 \ \Longrightarrow\ \Exp(1).
\end{equation}
The limiting conditional $q$th moments equal $q!$, which determine the
exponential law. The same reasoning using \eqref{eq:falling} gives the
joint limit
\begin{equation}\label{eq:joint-limit}
 \left.\left(\frac{2C_N}{\log N},\frac{2N\Theta_N}{\log N}\right)
 \ \right|\ C_N>0\ \Longrightarrow\ (Z,Z),\qquad Z\sim\Exp(1).
\end{equation}
To check the joint assertion, the Beta law gives
$\EE[M\Theta_N\mid C_N]=C_N$ and
$\operatorname{Var}(M\Theta_N\mid C_N)=C_N(M-C_N)/(M+1)\le C_N$.
Since $\EE[C_N\mid C_N>0]\sim\tfrac12\log N$, the difference divided by
$\log N$ vanishes in conditional $L^2$.

\section{Deletion statistics and an agreement graph}\label{sec:matching}

For an even sample $A$ of $2n$ points, let $D_A(u)$ be the survivor after
deleting $u$. Call an unordered pair $\{u,v\}$ \emph{mutually surviving
under deletion} if $D_A(u)=v$ and $D_A(v)=u$, and let $J_n$ count such pairs.
This terminology is specific to this note; it is not a literature claim.

\begin{proposition}\label{prop:deletion}
Every such pair is an edge of the original greedy matching. Moreover,
\begin{equation}\label{eq:deletion}
 \EE J_n=\sum_{j=0}^{n-1}\frac1{2j+1},
\end{equation}
and their expected edge-length counting measure has density
\begin{equation}\label{eq:deletion-density}
 \frac{1-(1-2r)^{2n}}{2r(1-r)},\qquad 0<r<1/2.
\end{equation}
\end{proposition}
\begin{proof}
For a fixed pair, put $B=A\setminus\{u,v\}$. The deletion condition says
exactly that $u,v\in E_B$. When both are added, they cannot affect the base
matching before their own meeting, since each lies outside every forbidden
ball \eqref{eq:available-balls}. Thus they pair with each other and the
base matching is unchanged. The two-point function \eqref{eq:two-point},
multiplied by $2\binom{2n}{2}$, gives \eqref{eq:deletion-density}; integration
gives \eqref{eq:deletion}.
\end{proof}

For $r\ge1/4$, such a pair is necessarily the last pair, by the shielding
cutoff. For smaller $r$ it need not be last. \emph{Do not extend the
last-pair density \eqref{eq:last-density} below $1/4$ using
\eqref{eq:deletion-density}.} Also $J_n$ and $C_N$ are different statistics.

\begin{proposition}[Deletion graph]\label{prop:deletion-graph}
For an even configuration $A$ of size $m=N+1$, draw the directed edge
$u\longrightarrow D_A(u)$ and let
$d_v=\#\{u:D_A(u)=v\}$ be the indegree of $v$. For an iid uniform sample,
\begin{align}
 A_N:=\EE\sum_{v\in B}\theta_v^2
 &=\frac{\EE\sum_{v\in A}d_v(d_v-1)}{N(N+1)},\label{eq:indegree}\\
 p_N&=1-\frac{\EE\sum_{v\in A}d_v(d_v-1)+2\EE J_{(N+1)/2}}
 {N(N+1)}.\label{eq:graph-noise}
\end{align}
Every directed cycle is a mutually surviving original matching pair.
Consequently the graph has exactly $J_{m/2}$ weak components, whose
expected number is $\sum_{j=0}^{m/2-1}(2j+1)^{-1}\sim\tfrac12\log m$.
\end{proposition}
\begin{proof}
Choose a uniformly random ordered pair of distinct labels $i,j$ in $A$.
Then $A\setminus\{i,j\}$ is the common background and the two omitted
points are independent uniform insertions. The survivors of $A-i$ and
$A-j$ agree at an existing point precisely when $D_A(i)=D_A(j)$.
There are $\sum_v d_v(d_v-1)$ such ordered pairs. The two resampled labels
also agree when $D_A(i)=j$ and $D_A(j)=i$, accounting for $2J_{m/2}$ more
ordered pairs. Divide by $m(m-1)$ and average.

For the cycle assertion, write $r_u$ for the separation at which $u$ is
matched in the even process on $A$. In the coupling with $A-u$, the first
discrepancy transfer occurs at separation $r_u$, from $u$ to its original
partner. Later transfers occur at strictly larger event separations. The
final survivor $D_A(u)$ must have become the extra point of the odd system
at its own removal event in the even system. Hence
$r_{D_A(u)}\ge r_u$, with equality only if $D_A(u)$ is the original
partner of $u$. Around a directed cycle all inequalities must be equalities;
each edge is therefore an original partner edge, forcing a two-cycle.
A finite directed graph with outdegree one has exactly one directed cycle
in each weak component. Apply \eqref{eq:deletion}.
\end{proof}

Thus the basin second moment also has a purely finite combinatorial target:
control the squared indegrees of this deletion graph. The logarithmic
expected component count alone gives no such control.

The matching on cyclically ordered indices is noncrossing: each removal
joins two currently adjacent vertices. Its distribution is invariant under
cyclic rotations and reversal of indices. Natural further questions are
the probability of a specified noncrossing matching, the distribution of
edge spans and nesting depths, and the relation of those statistics to
deletion sensitivity. Uniformity over noncrossing matchings has not been
established and must not be assumed.

\section{Exact insertion maps and a closed finite moment recursion}\label{sec:maps}

\subsection{The deterministic map}

The basin problem has a deterministic finite reduction that is more useful
than just knowing the unmarked gap lengths. Let $u,v$ be the closest pair,
with gap $m$ along its shorter arc. Let $a,c$ be the preceding and following
gaps. In local lifted coordinates with $u<v$, define
\[
 \ell=u-\min(m,a/2),\qquad h=(u+v)/2,\qquad
 r=v+\min(m,c/2),
\]
and the map
\begin{equation}\label{eq:T}
 T(z)=\begin{cases}
 v,&\ell<z<h,\\
 u,&h<z<r,\\
 z,&\text{otherwise}.
 \end{cases}
\end{equation}
Intervals are interpreted on the circle. With only two base points,
$a=c=1-m$. Values at the finitely many endpoints do not affect masses.

\begin{proposition}\label{prop:map}
The survivor position satisfies the exact recursion
\begin{equation}\label{eq:map-recursion}
 S_B(z)=S_{B\setminus\{u,v\}}(T(z)).
\end{equation}
\end{proposition}
\begin{proof}
The interval $(\ell,h)$ consists precisely of insertions that meet $u$
before the first base collision and before meeting another base point.
Such an insertion disappears with $u$, leaving $v$ as the discrepancy after
the base collision $u$--$v$. Starting the reduced background with insertion
$v$ gives the same subsequent active configuration. The other interval is
symmetric. For an insertion outside these two intervals, $u,v$ still remove
each other and can be omitted. If the insertion meets a different base
point before distance $m$, the resulting base-only discrepancy is frozen
until that point's base matching time, which is at least $m$; it cannot
interfere with the removal of $u,v$. This covers all cases.
\end{proof}

Compose these maps in base-removal order and push forward Lebesgue measure.
The final measure is
\begin{equation}\label{eq:final-measure}
 \nu_B=\ind_{E_B}(z)\,dz+\sum_{v\in B}\theta_v\delta_v.
\end{equation}
Indeed each map is the identity off two collapsed intervals; any point that
has once been collapsed stays in the discrete part, while an uncollapsed
point has retained its original position. A finite implementation can
therefore track continuous intervals and atoms exactly, without a probe grid.
\emph{The recursion preserves positions.} When a collapsed insertion is
represented by a deleted base endpoint, that endpoint keeps its original
base label for the final agreement calculation; it is not the label $*$.

For measures $\mu_i$ on independent daughter intervals, a merger sends their
translated sum through $T$. Its tensor square expands as
\begin{equation}\label{eq:tensor}
 (T\times T)_\#\left(
 \sum_i\mu_i\otimes\mu_i+
 \sum_{i\ne j}\mu_i\otimes\mu_j\right).
\end{equation}
The following construction makes this algebra into a closed, exact
recursion at fixed sample size. The asymptotic argument is supplied
separately in Section~\ref{sec:marked-limit}.

\subsection{Marked interval measures and the last internal pair}
\label{sec:finite-recursion}

Fix $x>0$ and $t\ge0$. For $2n$ ordered interior points
$0<y_1<\cdots<y_{2n}<x$, let $H_t$ be the event that all interior points
clear without touching the protected endpoints $0,x$, with every interior
pair separation at most $t$. As in Lemma~\ref{lem:shield}, ignore the
eventual endpoint--endpoint collision. On $H_t$, push $\Leb_{[0,x]}$
through all internal removal maps, using the active endpoints as the
boundary neighbours; call the resulting measure $\mu$. It has total mass
$x$. Define finite Borel measures by
\begin{align}
 W_n(t,x)&=\int_{0<y_1<\cdots<y_{2n}<x}\ind_{H_t}\,d\mathbf y,
 \label{eq:marked-definition}\\
 \mathcal M_n(t,x)&=\int\ind_{H_t}\mu\,d\mathbf y,\nonumber\\
 \mathcal Q_n(t,x)&=\int\ind_{H_t}(\mu\otimes\mu)\,d\mathbf y.
 \nonumber
\end{align}
Measure integrals here and below mean integration against nonnegative
Borel test functions. These are \emph{unnormalized ordered-simplex
integrals}, not conditional moments. For $2n$ iid uniform points in this
interval, the corresponding event-weighted expectations are multiplied
by $(2n)!/x^{2n}$.

For $n=0$, the initialization, for every $t\ge0$, is
\begin{equation}\label{eq:marked-initial}
 W_0(t,x)=1,\qquad \mathcal M_0(t,x)=\Leb_{[0,x]},\qquad
 \mathcal Q_0(t,x)=\Leb_{[0,x]}\otimes\Leb_{[0,x]}.
\end{equation}
For $n\ge1$, introduce the integration operator
\[
 \mathcal I_{n,t,x}[F]=
 \int_0^{\min(t,x/3)}\!\int_s^{x-2s}
 \sum_{i_1+i_2+i_3=n-1} F\,da\,ds.
\]
In each integrand let $c=x-a-s$, let $(x_1,x_2,x_3)=(a,s,c)$, and
evaluate all three daughter fields at cutoff $s$:
\[
 W_\ell=W_{i_\ell}(s,x_\ell),\qquad
 \mathcal M_\ell=\mathcal M_{i_\ell}(s,x_\ell),\qquad
 \mathcal Q_\ell=\mathcal Q_{i_\ell}(s,x_\ell).
\]
Translate their spatial coordinates by $0,a,a+s$, respectively, before
using them in a parent measure. In this integrand $T$ is the map
\eqref{eq:T} for the pair $u=a,v=a+s$ with neighbouring gaps $a,c$.

\begin{theorem}[Closed finite moment recursion]\label{thm:marked}
With these conventions, for every $n\ge1$,
\begin{align}
 W_n(t,x)&=\mathcal I_{n,t,x}[W_1W_2W_3],\label{eq:marked-W}\\
 \mathcal M_n(t,x)&=\mathcal I_{n,t,x}\left[
 T_\#\left(\sum_{\ell=1}^3\mathcal M_\ell\prod_{k\ne\ell}W_k\right)
 \right],\label{eq:marked-M}\\
 \mathcal Q_n(t,x)&=\mathcal I_{n,t,x}\left[
 (T\times T)_\#\left(
 \sum_{\ell=1}^3\mathcal Q_\ell\prod_{k\ne\ell}W_k
 +\sum_{\ell\ne j}\mathcal M_\ell\otimes\mathcal M_j
       \prod_{k\notin\{\ell,j\}}W_k\right)\right].
 \label{eq:marked-Q}
\end{align}
The second sum has six ordered terms. Together with
\eqref{eq:marked-initial}, these equations determine all the fields by
induction on $n$; there is no unspecified finite-size boundary term.
\end{theorem}
\begin{proof}
In a nonempty shielded interval, let the last internal pair be $a,a+s$.
At that instant the remaining gaps are $a,s,c$. The protected endpoints
survive this event precisely when $a,c>s$, up to null boundaries. Thus
$s\le\min(t,x/3)$ and $s\le a\le x-2s$.

Before this pair disappears, its two endpoints separate the three open
daughter intervals. No matching edge can pass across a still active
separating point. Each daughter consequently contains an even number
$2i_\ell$ of points, clears without touching its endpoints, and clears by
cutoff $s$. Conversely, three such daughter configurations with
$a,c>s$ combine into a shielded parent whose last internal pair is
$a,a+s$: no daughter endpoint--endpoint separation is smaller than $s$,
and there is no other possible cross-daughter first collision.

This is a bijection outside null equality sets. The change from the last
pair positions to $(a,s)$ has Jacobian one. The ordered-simplex volume
factors into the three daughter volumes; no binomial or factorial weight
is present. This proves \eqref{eq:marked-W}.

Every internal daughter collapse acts only inside that daughter's
interval. Hence, just before the last collision, the parent pushed measure
is the sum of the three translated daughter measures. Its final value is
$T_\#(\mu_1+\mu_2+\mu_3)$. Integrate this identity and its tensor square
against the product daughter volume measures. Fubini's theorem gives
\eqref{eq:marked-M}--\eqref{eq:marked-Q}. Finally each daughter index is
at most $n-1$, proving the inductive assertion.
\end{proof}

Useful normalizations and checks, valid for every Borel $F\subseteq[0,x]$,
are
\begin{align}
 \mathcal M_n(t,x)([0,x])&=xW_n(t,x),\nonumber\\
 \mathcal Q_n(t,x)(F\times[0,x])&=x\mathcal M_n(t,x)(F),
 &\mathcal Q_n(t,x)([0,x]^2)&=x^2W_n(t,x).\label{eq:marked-mass}
\end{align}
Lemma~\ref{lem:shield} gives
\[
 W_n(t,x)=\frac{x^{2n}}{(2n+1)!}\quad(t\ge x/3),\qquad
 W_1(t,x)=xb-\tfrac32b^2,\quad b=\min(t,x/3).
\]
Writing $D_x(z)=xz$, spatial scaling gives
\begin{align*}
 W_n(t,x)&=x^{2n}W_n(t/x,1),\\
 \mathcal M_n(t,x)&=x^{2n+1}(D_x)_\#\mathcal M_n(t/x,1),\\
 \mathcal Q_n(t,x)&=x^{2n+2}(D_x\times D_x)_\#\mathcal Q_n(t/x,1).
\end{align*}

\textbf{Relation to Poisson formulations.}
For intensity $\lambda$, the fields
$\lambda e^{-\lambda x}\sum_{n\ge0}\lambda^{2n}
(W_n,\mathcal M_n,\mathcal Q_n)$ are $\lambda$ times the
corresponding Poisson event-weighted moments. Their series converge
absolutely on every finite interval: the total of the order-$q$ measure
is bounded by $x^{2n+q}/(2n)!$, for $q=0,1,2$.
Resumming \eqref{eq:marked-W}--\eqref{eq:marked-Q} gives the same
three-daughter integral equations, with the zero-point term
$\lambda e^{-\lambda x}(1,\Leb,\Leb\otimes\Leb)$ and with the index
sum removed. The factor $\lambda$ makes the three daughter weights
multiply correctly. This justifies the finite-intensity equations;
by itself it does not establish their high-intensity limit. The
fixed-count formulation already applies directly to the original
problem; Section~\ref{sec:marked-limit} proves the marked convergence
needed for the noise limit without a Poisson conditioning argument.

\subsection{An exact readout of existing-point agreement}
\label{sec:readout}

For $x\ge r>0$, put $h=h_r(x)=\min(r,x/2)$ and define
\begin{align*}
 m_{j,r}(x)&=\mathcal M_j(r,x)([0,h]),\\
 q_{j,r}(x)&=\mathcal Q_j(r,x)([0,h]^2),\\
 d_{j,r}(x)&=\mathcal Q_j(r,x)
       (\{(z,z):h<z<x-h\}).
\end{align*}
The diagonal quantity retains squared atomic masses in the middle; a
diffuse measure contributes zero there. Null integration boundaries can
be assigned either way.

\begin{theorem}[Fixed-count agreement formula]\label{thm:readout}
For $n\ge1$, let $x=1-r$ inside the following integral. Then
\begin{align}
 A_{2n+1}=(2n)!\int_0^{1/2}\sum_{j+k=n-1}\big[&
 2q_{j,r}(r)W_k(r,x)+2W_j(r,r)q_{k,r}(x)\nonumber\\
 &+2rW_j(r,r)m_{k,r}(x)+W_j(r,r)d_{k,r}(x)
 \big]\,dr.\label{eq:readout}
\end{align}
Thus the finite recursion determines the exact basin second moment and,
using \eqref{eq:noise-identity}, the exact finite-$N$ noise probability.
No additional moments of the continuous part are needed for this readout.
\end{theorem}
\begin{proof}
Condition by integration on the final background pair. Its shorter arc
has length $r<1/2$ and its longer arc has length $x=1-r$. The two arc
interiors must be shielded and clear by separation $r$. They contain
$2j,2k$ points, where $j+k=n-1$. Let their pushed measures be $\mu_I$,
$\mu_J$, with local coordinates starting at the respective arc endpoints.
Before the final collapse define
\[
 P=\mu_I([0,r/2]),\qquad L=\mu_J([0,h]),\qquad
 R=\mu_J([x-h,x]),\quad h=\min(r,x/2).
\]
The two final-endpoint basin masses are $r-P+L$ and $P+R$.
Every other surviving atom lies in the middle of the long arc. Therefore
the existing-point agreement for this configuration is exactly
\begin{equation}\label{eq:readout-pathwise}
 (r-P+L)^2+(P+R)^2+
 (\mu_J\otimes\mu_J)(\{(z,z):h<z<x-h\}).
\end{equation}
Reflection of a daughter interval preserves its shielded volume measure.
The unnormalized means of $P,r-P$ both equal $rW_j(r,r)/2$, and their
second moments both equal $q_{j,r}(r)$. Similarly $L,R$ have first moment
$m_{k,r}(x)$ and second moment $q_{k,r}(x)$. The two daughter integration
measures factor, so expanding \eqref{eq:readout-pathwise} gives precisely
the integrand in \eqref{eq:readout}.

For completeness the prefactor is $(2n)!$: choose the unordered final
pair in $\binom{2n}{2}$ ways, use the shorter-separation density $2\,dr$,
and allocate the other labels to the two arcs. The allocation and ordering
factors are
$\binom{2n-2}{2j}(2j)!(2k)!=(2n-2)!$. The unique final pair ensures
that this sum counts each configuration once. Integration over its
overall rotation contributes one.
\end{proof}

For $n=1$ both daughter intervals are empty. Then $m_{0,r}(x)=h$,
$q_{0,r}(x)=h^2$, and $d_{0,r}(x)=0$, giving the explicit check
\[
 A_3=2!\int_0^{1/2}2\big(r/2+\min(r,(1-r)/2)\big)^2\,dr
     =\frac5{18},\qquad p_3=1-\frac5{18}-\frac29=\frac12.
\]

\subsection{The full scalar last-pair law and a quantitative limit}
\label{sec:scalar-limit}

The scalar clearance field from Section~\ref{sec:finite-recursion} is
explicitly solvable. The resulting fixed-count limit does not require
an unproved de-Poissonization step.

\begin{theorem}[Scalar transform]\label{thm:scalar-transform}
For $\lambda>0$ define
\[
 w_r^\lambda(x)=\lambda e^{-\lambda x}
       \sum_{n\ge0}\lambda^{2n}W_n(r,x),\qquad
 F_r^\lambda(s)=\int_r^\infty e^{-sx}w_r^\lambda(x)\,dx.
\]
For $r,s>0$,
\begin{equation}\label{eq:scalar-transform}
 F_r^\lambda(s)=\tanh\!\left(
 \frac{E_1(sr)-E_1((s+2\lambda)r)}2\right),\qquad
 E_1(u)=\int_u^\infty\frac{e^{-v}}v\,dv.
\end{equation}
\end{theorem}
\begin{proof}
Put $b_\lambda(x)=(1-e^{-2\lambda x})/(2x)$. The full-shielding
normalization gives $w_r^\lambda(r)=b_\lambda(r)$ and the useful bound
$0\le w_r^\lambda(x)\le b_\lambda(x)$ for every $x>0$.
Resum \eqref{eq:marked-W} and differentiate in the cutoff, almost everywhere:
\[
 \partial_r w_r^\lambda(x)=b_\lambda(r)\ind_{\{x\ge3r\}}
       \int_r^{x-2r}w_r^\lambda(a)w_r^\lambda(x-a-r)\,da.
\]
The moving lower limit in the Laplace transform is essential. It gives
\[
 \partial_r F_r^\lambda(s)
 =-e^{-sr}b_\lambda(r)\bigl(1-(F_r^\lambda(s))^2\bigr),
 \qquad F_0^\lambda(s)=\frac{\lambda}{s+\lambda}.
\]
All integrals and differentiations follow from the nonnegative integral
recursion and the displayed bounds; the identity is an absolutely
continuous differential equation in $r$. Its initial value is below one.
Integrating its $\operatorname{arctanh}$ form gives
\[
 \operatorname{arctanh}F_r^\lambda(s)
 =\int_r^\infty e^{-su}\frac{1-e^{-2\lambda u}}{2u}\,du,
\]
which proves the formula.
\end{proof}

Let $c_k=[z^k]\tanh z$ and put $g(y)=\ind_{\{y\ge1\}}/(2y)$.
Only odd $k$ occur, and $|c_k|\le1$. For the latter bound write
$\tan z=\sum_{j\ge0}a_jz^{2j+1}$; its differential equation gives
$a_0=1$ and $(2j+1)a_j=\sum_{h=0}^{j-1}a_ha_{j-1-h}\le j$
inductively. Define the locally finite convolution sum
\begin{equation}\label{eq:scalar-density-series}
 f(y)=\sum_{\substack{k\ge1\ \mathrm{odd}\\ k\le y}}c_k g^{*k}(y),
 \qquad y\ge1.
\end{equation}

\begin{theorem}[Fixed-count last-pair limit]\label{thm:full-last-law}
The function $f$ is a probability density on $[1,\infty)$ and
\begin{equation}\label{eq:limit-transform}
 \EE e^{-sY}=\int_1^\infty e^{-sy}f(y)\,dy
       =\tanh\!\left(\frac12E_1(s)\right),\qquad Y\sim f.
\end{equation}
For even $M\to\infty$, the last separation converges in total variation:
\[
 L_M\ \longrightarrow\ L_\infty:=\frac1{1+Y},\qquad
 f_\infty(r)=r^{-2}f((1-r)/r),\quad 0<r<1/2.
\]
For density versions defined by the formulas below,
\begin{equation}\label{eq:density-error}
 |f_M(r)-f_\infty(r)|
 \le r^{-2}e^{1/r}(1-2r)^M\qquad(0<r\le1/2).
\end{equation}
In particular, with total variation defined as $\sup_E|\PP_1(E)-\PP_2(E)|$,
for $M\ge4$ one has
\begin{align}
 \|\mathcal L(L_M)-\mathcal L(L_\infty)\|_{\rm TV}
 &\le B_M:=\min\left\{1,
 K e^{-(\sqrt M-1)/4}+\frac M2e^{-\sqrt M}\right\},
 \label{eq:last-TV}\\
 K&=\frac{1+q}{1-q}<7,\qquad
 q=\frac{e^\gamma}{4}\exp\!\left(
       \sum_{\ell\ge1}\frac{4^{-\ell}}{\ell\,\ell!}\right).
 \nonumber
\end{align}
Here $\gamma$ is Euler's constant. Thus the error is bounded by a constant
times $e^{-\sqrt M/4}$ for all sufficiently large $M$.
\end{theorem}
\begin{proof}
\emph{Local inversion and positivity.}
Set $b_{\lambda,r}(x)=b_\lambda(x)\ind_{\{x\ge r\}}$.
The transform on the right of \eqref{eq:scalar-transform} is
$\tanh(\mathcal L b_{\lambda,r})$. Expanding at sufficiently large $s$
and inverting gives
\begin{equation}\label{eq:finite-convolution}
 w_r^\lambda(x)=\sum_{\substack{k\ge1\ \mathrm{odd}\\ kr\le x}}
                 c_k b_{\lambda,r}^{*k}(x),\qquad x\ge r.
\end{equation}
This use of the power series is legitimate: on each bounded $x$ interval
the sum is finite, while the bound $|c_k|\le1$ and the simplex-volume
formula give an exponential bound in $x$ for the absolute sum. Its
Laplace transform therefore converges for sufficiently large $s$, where
the Taylor expansion is absolutely convergent, and uniqueness applies.
For fixed $r,x>0$, let $\lambda\to\infty$. The finite convolutions
converge to $r^{-1}f(x/r)$. Since $w_r^\lambda\ge0$, this limit is
nonnegative. The bound $w_r^\lambda(x)\le1/(2x)$ permits dominated
convergence of every Laplace transform with $s>0$. Taking $r=1$ gives
\eqref{eq:limit-transform}. Letting $s\downarrow0$ shows that $\int f=1$.

\emph{Direct coefficient extraction at fixed count.}
The same final-pair decomposition used in Theorem~\ref{thm:readout} gives
\begin{equation}\label{eq:last-generating}
 e^{-\lambda}\sum_{n\ge1}\frac{\lambda^{2n}}{(2n)!}f_{2n}(r)
 =w_r^\lambda(r)w_r^\lambda(1-r).
\end{equation}
In the $k$th term of \eqref{eq:finite-convolution}, put $y_0=r$ and
integrate over
\[
 \mathcal D_k(r)=\{y_1,\ldots,y_k\ge r:\ y_1+\cdots+y_k=1-r\},
\]
using $dy_1\cdots dy_{k-1}$ with $y_k$ determined by the sum.
For $k=1$ this is evaluation at $y_1=1-r$. Write
$y_S=\sum_{i\in S}y_i$. Extracting $\lambda^M$ in
\eqref{eq:last-generating} yields the exact finite sum
\begin{equation}\label{eq:last-inclusion}
 f_M(r)=\sum_{\substack{k\ge1\ \mathrm{odd}\\ (k+1)r\le1}}c_k
 \int_{\mathcal D_k(r)}
 \frac{\displaystyle\sum_{S\subseteq\{0,\ldots,k\}}
       (-1)^{|S|}(1-2y_S)^M}
      {2^{k+1}\prod_{i=0}^k y_i}\,d\mathbf y.
\end{equation}
Indeed $e^\lambda\prod_i(1-e^{-2\lambda y_i})$ is a finite
sum of exponentials with those bases. There is no asymptotic conditioning
argument in this step.

For even $M$, the empty and full subsets contribute two. Every proper
nonempty subset has $r\le y_S\le1-r$, so its contribution has absolute
value at most $(1-2r)^M$. Keeping just the two endpoint subsets gives
$2b_\infty(r)r^{-1}f((1-r)/r)=f_\infty(r)$.
The simplex volume is
$(1-(k+1)r)^{k-1}/(k-1)!\le1/(k-1)!$ and every $y_i\ge r$.
Consequently the sum of the remaining absolute contributions is bounded by
\[
 (1-2r)^M\sum_{k\ge1}\frac{r^{-(k+1)}}{(k-1)!}
 =r^{-2}e^{1/r}(1-2r)^M.
\]
This proves \eqref{eq:density-error}, in particular uniform convergence
on $[\varepsilon,1/2]$. Since $f_\infty$ integrates to one under the
change of variables $y=(1-r)/r$, no mass escapes to $r=0$.

\emph{Tail control and the quantitative bound.}
The identity
\[
 E_1(s)=-\gamma-\log s-
           \sum_{\ell\ge1}\frac{(-s)^\ell}{\ell\,\ell!}
\]
rewrites \eqref{eq:limit-transform} as $(1-Z(s))/(1+Z(s))$, where
$Z(s)=e^\gamma s\exp(\sum_{\ell\ge1}(-s)^\ell/(\ell\,\ell!))$ is
entire. On $|s|\le1/4$, $|Z(s)|\le q<3/4$: use
$e^\gamma<2$, $\sum4^{-\ell}/(\ell\,\ell!)<1/3$, and
$e^{1/3}<3/2$. The transform is therefore analytic on a larger disk.
Its derivatives at zero are the moments of $Y$, by monotone convergence
from positive $s$; the Taylor series and monotone convergence then give
\[
 \EE e^{Y/4}=\frac{1+q}{1-q}=K<7,
 \qquad \EE Y=2e^\gamma.
\]
Thus $\PP(L_\infty<\varepsilon)
\le K\exp(-(1/\varepsilon-1)/4)$.
On $[\varepsilon,1/2]$, the density error integral is at most
$\tfrac12\varepsilon^{-2}\exp(1/\varepsilon-2M\varepsilon)$.
Because both laws have mass one, their total variation is at most this
integral plus $\PP(L_\infty<\varepsilon)$: the unknown small-interval
mass of the finite law is controlled by the discrepancy on its complement.
Take $\varepsilon=M^{-1/2}$ and also use the trivial upper bound one.
\end{proof}

\begin{corollary}[Finite piecewise polynomial densities]\label{cor:polynomial}
For each even $M\ge2$, $f_M(r)$ is a polynomial with rational coefficients
of degree at most $M-2$ on each interval cut out by
$0,1/M,1/(M-2),\ldots,1/2$. In particular,
\[
 f_2(r)=2,\qquad
 f_4(r)=\begin{cases}
 24r-56r^2,&0<r<1/4,\\
 4-8r+8r^2,&1/4\le r<1/2.
 \end{cases}
\]
\end{corollary}
\begin{proof}
In \eqref{eq:last-inclusion}, before coefficient extraction, use
\[
 e^\lambda\prod_{i=0}^k b_\lambda(y_i)
 =\prod_{i=0}^k\frac{\sinh(\lambda y_i)}{y_i}.
\]
The coefficient of $\lambda^M$ is zero for $k+1>M$. Otherwise it is a
rational polynomial in the lengths, homogeneous of degree $M-k-1$.
Translate $y_i=r+z_i$ for $i\ge1$; the simplex has side-sum
$1-(k+1)r$. Integrating monomials uses
\[
 \int_{\substack{z_i\ge0\\ \sum_{i=1}^kz_i=L}}
       \prod_{i=1}^kz_i^{h_i}\,dz_1\cdots dz_{k-1}
 =\frac{\prod_i h_i!}{(\sum_i h_i+k-1)!}
       L^{\sum_i h_i+k-1}.
\]
The resulting degree is at most $M-2$, and the only new activation
threshold for the $k$th term is $r=1/(k+1)$. For $M=4$ the $k=3$ term
subtracts $4(1-4r)^2$ below $1/4$ from the $k=1$ polynomial, giving the
stated example.
\end{proof}

There is also a positive delay identity for the limiting density. Differentiate
\eqref{eq:limit-transform} and invert its Laplace transform to obtain
\begin{equation}\label{eq:density-delay}
 y f(y)=\frac12\left(1-\int_0^{y-1}(f*f)(u)\,du\right)
 =\frac12\PP(Y_1+Y_2>y-1),\qquad y\ge1,
\end{equation}
where $Y_1,Y_2$ are independent copies. This recovers $f(y)=1/(2y)$ on
$[1,3]$ and determines successive intervals using only earlier values.

\begin{corollary}[Stronger noise lower bound]\label{cor:strong-noise}
The explicit constant
\begin{align}
 c_*&=\EE[L_\infty(1-L_\infty)]
 =\frac12\int_0^\infty e^{-2s}
          \operatorname{sech}^2\!\left(\frac{E_1(s)}2\right)\,ds
 \label{eq:cstar}\\
 &\approx0.1868187712\nonumber
\end{align}
satisfies $\liminf_{N\to\infty,\ N\ \mathrm{odd}}p_N\ge c_*$.
More quantitatively, for $N=M+1\ge5$,
\begin{equation}\label{eq:strong-finite}
 p_N\ge c_*-\tfrac14B_M.
\end{equation}
The integral defines the rigorous constant; a certified decimal
enclosure is supplied in Section~\ref{sec:sharper-budgets}.
\end{corollary}
\begin{proof}
Apply Lemma~\ref{lem:last-comparison}. The function $r(1-r)$ takes values
in $[0,1/4]$, so \eqref{eq:last-TV} bounds the error in its expectation
by $B_M/4$. To evaluate the limit, write $F(s)=\EE e^{-sY}$ and use
\[
 \EE\frac{Y}{(1+Y)^2}
 =\int_0^\infty(1-s)e^{-s}F(s)\,ds
 =-\int_0^\infty se^{-s}F'(s)\,ds.
\]
Since $F'(s)=-e^{-s}(1-F(s)^2)/(2s)$, this is \eqref{eq:cstar}.
The integration by parts has zero boundary terms.
\end{proof}

\begin{corollary}[An improved noise upper bound]\label{cor:upper-noise}
With $F(s)=\tanh(E_1(s)/2)$, set
\begin{equation}\label{eq:ustar}
 u_*:=\frac{23}{24}-\frac14\log\frac32
        -\frac12\int_0^\infty se^{-s}F(s)\,ds
       \approx0.8112015673.
\end{equation}
Then $\limsup p_N\le u_*$. For $N=M+1\ge5$,
$p_N\le u_*+B_M/2$. Section~\ref{sec:sharper-budgets} supplies a
certified decimal enclosure of this integral expression.
\end{corollary}
\begin{proof}
All short-arc insertions end at one of the two last endpoints, so their
combined basin mass is at least $L_M$. Their contribution to
$\sum_v\theta_v^2$ is therefore at least $L_M^2/2$. When $L_M>1/3$,
both arcs have length less than $2L_M$, and the combined mass is one.
Consequently, with
$\psi(r)=\tfrac12r^2\ind_{\{r\le1/3\}}+\tfrac12\ind_{\{r>1/3\}}$,
\[
 p_{M+1}\le1-\EE\psi(L_M)
 \le1-\EE\psi(L_\infty)+\tfrac12B_M.
\]
Here $\PP(L_\infty>1/3)=\tfrac12\log2$ and
\[
 \EE[L_\infty^2\ind_{\{L_\infty>1/3\}}]
 =\frac12\left(\log\frac43-\frac16\right),\qquad
 \EE L_\infty^2=\int_0^\infty se^{-s}F(s)\,ds.
\]
Substitute these identities to obtain \eqref{eq:ustar}.
\end{proof}

The scalar law now controls the averaging weights in the marked recursion.
It does \emph{not} identify the limiting conditional basin measures, or
the limit of $A_N$ or $p_N$.

\begin{corollary}[Removing very small final separations]\label{cor:truncation}
For $0<\delta<1/2$ and $N=M+1\ge5$, define
\[
 A_N^{(\delta)}=\EE\left[\sum_v\theta_v^2
                            \ind_{\{L_M\ge\delta\}}\right].
\]
Then
\begin{equation}\label{eq:truncation}
 0\le A_N-A_N^{(\delta)}
 \le K e^{-(1/\delta-1)/4}+B_M.
\end{equation}
Consequently it suffices to prove convergence of $A_N^{(\delta)}$ along
any sequence of positive $\delta$ decreasing to zero in order to prove
convergence of $A_N$, and hence of $p_N$.
\end{corollary}
\begin{proof}
The omitted agreement lies in $[0,\ind_{\{L_M<\delta\}}]$.
Apply the limiting tail bound and total variation. If the truncated
quantities converge, the displayed estimate makes the upper and lower
limits of $A_N$ differ by at most $K e^{-(1/\delta-1)/4}$; let
$\delta\downarrow0$. Finally use $\EE\Theta_N^2\to0$.
\end{proof}

\subsection{The positive-cutoff configuration and an early-label bound}
\label{sec:active-limit}

The scalar field also determines the entire configuration present at a
positive cutoff. This gives a finite-dimensional reduction of the marked
problem, with a quantitative bound on the discarded second moment.

For even $M$, let $C_M(\varepsilon)$ be the background points still
present after all pairs of separation at most $\varepsilon$ have been
removed, and put $J_M(\varepsilon)=|C_M(\varepsilon)|$. Fix
$0<\varepsilon<1/2$. A nonempty configuration has an even number $k$ of
points, with cyclic gaps $x_i>\varepsilon$ summing to one; in particular
$k\le\lfloor1/\varepsilon\rfloor$. In this subsection use the truncated
fields
\[
 h_\varepsilon^\lambda(x)=\ind_{\{x\ge\varepsilon\}}
                  w_\varepsilon^\lambda(x),\qquad
 h_\varepsilon(x)=\varepsilon^{-1}f(x/\varepsilon)
                  \ind_{\{x\ge\varepsilon\}}.
\]
The latter is a probability density. Parametrize a cyclic configuration
by its root $a\in\T$ and gaps $\mathbf x=(x_1,\ldots,x_k)$; its points
are $a,a+x_1,\ldots,a+x_1+\cdots+x_{k-1}$ modulo one. Each unrooted
configuration is represented $k$ times, so its reference integration is
$k^{-1}\,da\,dx_1\cdots dx_{k-1}$, with $x_k$ determined by the sum.

\begin{theorem}[The full configuration at a positive cutoff]
\label{thm:active-limit}
On the nonempty $k$-point sector the exact density relative to this
reference integration is
\begin{equation}\label{eq:active-finite-density}
 G_{M,k}(\mathbf x)=M![\lambda^M]\,
             e^\lambda\prod_{i=1}^k h_\varepsilon^\lambda(x_i).
\end{equation}
As even $M\to\infty$, $C_M(\varepsilon)$ converges in total variation to
a finite random configuration $C_\infty(\varepsilon)$ whose density is
\begin{equation}\label{eq:active-limit-density}
 G_{\infty,k}(\mathbf x)=2\prod_{i=1}^k h_\varepsilon(x_i).
\end{equation}
Its count distribution is
\begin{align}
 \PP(|C_\infty(\varepsilon)|=k)
   &=\frac2k h_\varepsilon^{*k}(1),\qquad k\ge2\text{ even},
       \label{eq:active-count}\\
 \PP(C_\infty(\varepsilon)=\varnothing)
   &=\PP(L_\infty\le\varepsilon).\nonumber
\end{align}
One explicit error bound is
\begin{align}
 \|\mathcal L(C_M(\varepsilon))-
           \mathcal L(C_\infty(\varepsilon))\|_{\rm TV}
 &\le\Delta_M(\varepsilon)
   :=\min\{1,D_\varepsilon(1-2\varepsilon)^M\},
       \label{eq:active-TV}\\
 D_\varepsilon&=
 \sum_{\substack{2\le k\le\lfloor1/\varepsilon\rfloor\\ k\ {\rm even}}}
       \frac{(\varepsilon^{-1}e^{1/\varepsilon})^k}{k!}.
       \nonumber
\end{align}
Thus the complete future collision history above this cutoff also
converges in total variation. The limiting laws at different positive
cutoffs are consistent under greedy evolution.
\end{theorem}
\begin{proof}
Between consecutive active points, the original interior points are even
in number, shield their endpoints, and clear by $\varepsilon$. Conversely,
such gap interiors, together with gaps longer than $\varepsilon$, leave
precisely the specified active points. Apply the shielding decomposition
independently in all $k$ gaps. For $2n_i$ interior points in gap $i$,
$M=k+2\sum n_i$. Multiplying the series for $h_\varepsilon^\lambda$
and cancelling $e^{-\lambda\sum x_i}=e^{-\lambda}$ gives
\eqref{eq:active-finite-density}. Ordered-simplex integration of the
original labelled sample supplies $M!$, while the choice of cyclic root
supplies $1/k$. No further allocation factor is needed.

For fixed $\mathbf x$, expand each factor using
\eqref{eq:finite-convolution}. The expansion is finite: its $i$th
convolution uses an odd number $j_i$ of component lengths, each at least
$\varepsilon$, summing to $x_i$. The total number $q=\sum j_i$ is even,
and all component lengths together sum to one. Expanding
$e^\lambda\prod(1-e^{-2\lambda y})$ and extracting its $M$th coefficient
gives exactly the same subset sum as in \eqref{eq:last-inclusion}.
Its empty and full subsets contribute two; every other subset contributes
at most $(1-2\varepsilon)^M$ in absolute value. The endpoint terms sum
to \eqref{eq:active-limit-density}.

For an explicit uniform bound, each component denominator is at least
$\varepsilon$, the $j_i$-component simplex volume is at most
$1/(j_i-1)!$, and $|c_{j_i}|\le1$. The factor $2^{-q}$ cancels the bound
$2^q$ on the number of subsets. Extending the sum to all $j_i\ge1$ gives
\[
 |G_{M,k}(\mathbf x)-G_{\infty,k}(\mathbf x)|
 \le(\varepsilon^{-1}e^{1/\varepsilon})^k(1-2\varepsilon)^M.
\]
The gap simplex has volume at most $1/(k-1)!$; integrate with the root
factor $1/k$ and sum over the finitely many even $k$.

The probability of no active point is $\PP(L_M\le\varepsilon)$, whose
limit is known from Theorem~\ref{thm:full-last-law}. Hence the displayed
limiting sector densities and the stated empty atom have total mass one.
The discrepancy of that atom is the negative of the integrated density
discrepancy on the nonempty sectors. Total variation is consequently at
most their integrated absolute discrepancy, proving
\eqref{eq:active-TV}. Finally the future history is a measurable
deterministic function of the active configuration. Pushforward cannot
increase total variation. Applying the same fact between two cutoffs
proves consistency. This statement does not include the insertion basin
masses as marks.
\end{proof}

For example, when $1/3\le\varepsilon<1/2$, only the two-point sector
has positive mass and
\[
 \PP(|C_\infty(\varepsilon)|=2)
       =\frac12\log\frac{1-\varepsilon}{\varepsilon}.
\]
The formula follows by integrating
$h_\varepsilon(x)h_\varepsilon(1-x)=1/(4x(1-x))$ over
$\varepsilon<x<1-\varepsilon$.

\begin{lemma}[A bound on the limiting gap square sum]\label{lem:gap-squares}
Let $x_i$ be the gaps of $C_\infty(\varepsilon)$ when it is nonempty.
Put $R_\varepsilon=\sum_{m\ge1,\ m\ {\rm odd}}
h_\varepsilon^{*m}$, a locally finite sum. Then
\begin{align}
 \EE\left[\ind_{\{C_\infty(\varepsilon)\ne\varnothing\}}
                         \sum_i x_i^2\right]
 &=2\int_0^1 x^2 h_\varepsilon(x)R_\varepsilon(1-x)\,dx
      \label{eq:gap-square-identity}\\
 &\le 8e^\gamma(1+e^\gamma)\varepsilon.\nonumber
\end{align}
For every $a>0$ there is also the mesh estimate
\begin{equation}\label{eq:gap-mesh}
 \PP\left(C_\infty(\varepsilon)\ne\varnothing,
                     \max_i x_i\ge a\right)
 \le\frac{2K}{\varepsilon}e^{-a/(4\varepsilon)}.
\end{equation}
\end{lemma}
\begin{proof}
Formula \eqref{eq:active-limit-density} and symmetry in the $k$ gaps
give \eqref{eq:gap-square-identity}: the sum over $k$ coordinates cancels
the factor $1/k$, leaving the odd convolution sum for the other gaps.

We claim $R_\varepsilon(u)\le1/\varepsilon$ for every $u$. By scaling
it is enough to prove $R_1\le1$. The delay identity
\eqref{eq:density-delay} shows that $f$ is nonincreasing on $[1,\infty)$,
with $f(1)=1/2$ and integral one. Let $Y_i$ be iid with density $f$ and
$S_{2j}=Y_1+\cdots+Y_{2j}$, including $S_0=0$. Then
\[
 R_1(u)=\EE\sum_{j\ge0}f(u-S_{2j}),
\]
where $f$ is zero below one. The arguments contributing to this sum
are at least two apart, since $Y_i\ge1$. At most one lies in each
half-open interval $[1+2\ell,3+2\ell)$. Monotonicity therefore gives,
pathwise,
\[
 \sum_{j\ge0}f(u-S_{2j})
 \le\sum_{\ell\ge0}f(1+2\ell)
 \le\frac12+\frac12\int_1^\infty f(y)\,dy=1.
\]
The gap square expectation is thus at most
$2\varepsilon\EE Y^2$. Expanding \eqref{eq:limit-transform} at zero
gives $\EE Y^2=4e^\gamma(1+e^\gamma)$, proving the bound. The same
symmetry argument bounds the expected number of gaps of length at least
$a$ by
\[
 2\int_a^1 h_\varepsilon(x)R_\varepsilon(1-x)\,dx
 \le\frac2\varepsilon\PP(Y\ge a/\varepsilon)
 \le\frac{2K}{\varepsilon}e^{-a/(4\varepsilon)}.
\]
This dominates the probability in \eqref{eq:gap-mesh}.
\end{proof}

\begin{theorem}[Early labels have a small agreement contribution]
\label{thm:early-labels}
For $N=M+1$ define the agreement retained by the background labels still
present at cutoff $\varepsilon$:
\[
 A_N^{\rm late}(\varepsilon)
   =\EE\sum_{v\in C_M(\varepsilon)}\theta_v(B)^2.
\]
These are the original full-process basin masses; they are not recomputed
by starting with $C_M(\varepsilon)$. For even $M\ge2$,
\begin{align}
 0\le A_N-A_N^{\rm late}(\varepsilon)
 \le{}&8e^\gamma(1+e^\gamma)\varepsilon
        +K e^{-(1/\varepsilon-1)/4}
        +\Delta_M(\varepsilon).\label{eq:early-label-bound}
\end{align}
Consequently
\begin{equation}\label{eq:early-label-double-limit}
 \lim_{\varepsilon\downarrow0}\limsup_{N\to\infty,\ N\ {\rm odd}}
       \big(A_N-A_N^{\rm late}(\varepsilon)\big)=0.
\end{equation}
Convergence of $A_N^{\rm late}(\varepsilon)$ for any sequence of positive
$\varepsilon$ decreasing to zero would therefore imply convergence of
$A_N$ and $p_N$.
\end{theorem}
\begin{proof}
Suppose first that $C_M(\varepsilon)$ is nonempty, with cyclic gaps $I_i$
of lengths $x_i$. All removal maps before the cutoff act inside one of
these gaps and preserve that gap's total insertion mass $x_i$. Every map
after the cutoff has its two destinations in $C_M(\varepsilon)$.
Accordingly, if an insertion ends at an original label
$v\notin C_M(\varepsilon)$, its position at the cutoff was already $v$
and no subsequent map moved it. An insertion outside $v$'s gap cannot
contribute to this basin. Thus, in each gap,
\[
 \sum_{v\in B\cap I_i}\theta_v\le x_i,\qquad
 \sum_{v\in B\cap I_i}\theta_v^2\le x_i^2.
\]
Here the gap interior contains only labels removed before the cutoff.
If the active configuration is empty, use the bound one instead. This
proves the pathwise estimate for the discarded agreement by the bounded
functional
\[
 H(C)=\begin{cases}1,&C=\varnothing,\\
                    \sum_i x_i^2,&C\ne\varnothing.
       \end{cases}
\]
Apply \eqref{eq:active-TV} to $0\le H\le1$, then
Lemma~\ref{lem:gap-squares} and the exponential bound for
$\PP(L_\infty\le\varepsilon)$. This yields
\eqref{eq:early-label-bound}. First send $M$ to infinity at fixed
$\varepsilon$, then $\varepsilon$ to zero. The convergence criterion
follows by bounding the difference between upper and lower limits by
the right side without $\Delta_M$, and using $\EE\Theta_N^2\to0$.
\end{proof}

The reduction above does not by itself prove convergence of the basin
weights: small gaps alone do not control later discontinuous collapse
maps. The next subsection supplies the needed stability estimate by
following cumulative-mass queries backwards through the interval tree.

\subsection{A contraction for the marks and convergence of the noise probability}
\label{sec:marked-limit}

The remaining convergence step can be closed without assuming continuity
of squared atomic masses under arbitrary weak convergence. First obtain
uniform bounds on the scalar density, then prove that a mass query forgets
the unspecified lower levels of its collision tree.

\begin{lemma}[A positive representation and two-sided exponential bounds]
\label{lem:positive-renewal}
There are constants $\kappa>0$ and $0<b_-\le b_+<\infty$ such that
\begin{equation}\label{eq:two-sided-f}
 b_-e^{-\kappa y}\le f(y)\le b_+e^{-\kappa y}\qquad(y\ge1).
\end{equation}
More precisely, there is a probability density $a_0$ on $[1,\infty)$,
positive and continuous there, such that
\begin{equation}\label{eq:positive-geometric}
 f=\sum_{j\ge1}2^{-j}a_0^{*j}.
\end{equation}
All exponential moments of $a_0$ are finite. With
\[
 \Psi(s)=\exp\left(\int_0^1\frac{e^{-su}-1}{u}\,du\right),
\]
$\kappa$ is the unique positive solution of
$e^\gamma\kappa\Psi(-\kappa)=1$.
\end{lemma}
\begin{proof}
Let $\Pi$ be a Poisson process on $(0,1)$ with intensity $du/u$ and
put $V=\sum_{u\in\Pi}u$. The sum is finite almost surely since its mean
is one, and its Laplace transform is $\Psi(s)$ for every real $s$.
For example, derive this first on $(\delta,1)$ and then let
$\delta\downarrow0$; the exponential moment formula stays finite for
every real argument. The largest point of $\Pi$ is uniform on $(0,1)$,
and the process below it, rescaled by that point, is an independent copy.
Equivalently use the exponential arrival times after taking negative
logarithms. Thus $V$ satisfies $V\overset d=U(1+V')$ with independent
uniform $U$ and $V'\overset d=V$.

It follows that $V$ has the positive continuous density
\[
 \varphi(x)=\EE\left[\frac{\ind_{\{V>x-1\}}}{1+V}\right]\quad(x>0).
\]
It is constant, say $c$, on $(0,1)$, and
$\varphi'(x)=-\varphi(x-1)/x$ for $x>1$. Positivity for every finite $x$
follows because the number of points of $\Pi$ in $(1/2,1)$ has unbounded
support. Also $\varphi(x)\le\PP(V>x-1)$, so this density decays faster
than every fixed exponential. The identity
$\Psi(s)=e^{-\gamma}s^{-1}e^{-E_1(s)}$ for $s>0$ gives
$c=\lim_{s\to\infty}s\Psi(s)=e^{-\gamma}$.

Set $a_0(x)=e^\gamma\varphi(x-1)/x$ for $x\ge1$, and zero below one.
Since $a_0=-\varphi'/c$ above one, it is a probability density.
Integration by parts gives
\[
 \widehat a_0(s)=1-e^\gamma s\Psi(s),\qquad s\in\mathbb R.
\]
For positive $s$, the transform in \eqref{eq:limit-transform} is
$\widehat a_0(s)/(2-\widehat a_0(s))$. Expanding this positive geometric
series and using uniqueness of Laplace transforms proves
\eqref{eq:positive-geometric}. All its sums are locally finite because
$a_0$ is supported on $[1,\infty)$.

The positive function $\kappa\mapsto e^\gamma\kappa\Psi(-\kappa)$
increases strictly from zero to infinity. At its unique value one,
\[
 a_\kappa(x):=\tfrac12 e^{\kappa x}a_0(x)
\]
is a probability density. It is positive on $[1,\infty)$, has every
exponential moment, and has an exponentially decreasing upper envelope.
Equation \eqref{eq:positive-geometric} becomes
\begin{equation}\label{eq:tilted-renewal}
 e^{\kappa x}f(x)=u_\kappa(x)
       :=\sum_{j\ge1}a_\kappa^{*j}(x).
\end{equation}
We give elementary uniform upper and lower bounds for this renewal
density; no renewal asymptotic is being assumed.

Let $T_0=0,T_n=Z_1+\cdots+Z_n$, where the $Z_i$ have density
$a_\kappa$ and hence $Z_i\ge1$. Then
$u_\kappa(x)=\EE\sum_{n\ge0}a_\kappa(x-T_n)$.
The $T_n$ are at least one apart, so this is bounded above by
\[
 b_+:=\sum_{j\ge0}\sup_{1+j\le v<2+j}a_\kappa(v)<\infty.
\]
Choose $R>0$ large enough that
$\sum_{j\ge0}\PP(Z_1>R+j)\le1/2$; finite mean suffices. For any $y\ge0$,
the probability that there is no renewal in $(y,y+R]$ equals
\[
 \EE\sum_{n:T_n\le y}\PP(Z_1>y+R-T_n)
 \le\sum_{j\ge0}\PP(Z_1>R+j)\le\tfrac12.
\]
The equality conditions on the last renewal before $y$; the displayed
sum is valid by independence of the next increment. The inequality
again uses unit spacing of the $T_n$.
Put $m_R=\min_{1\le v\le R+1}a_\kappa(v)>0$.
If $x\ge R+1$, take $y=x-R-1$: with probability at least $1/2$,
some summand in $u_\kappa(x)$ is at least $m_R$.
For $1\le x\le R+1$, already $u_\kappa(x)\ge a_\kappa(x)\ge m_R$.
Thus $b_-=m_R/2$ and \eqref{eq:tilted-renewal} prove the lemma.
\end{proof}

The auxiliary law of $V$ is the standard Dickman law; its uniform
perpetuity characterization is discussed by
\href{https://arxiv.org/abs/1706.08192}{Bhattacharjee and Goldstein,
\emph{Dickman approximation in simulation, summations and perpetuities}
(2017; revised 2018)}. The construction and estimates needed here were
proved above. This reference is not a novelty claim for the balloon
argument.

\begin{lemma}[Limiting interval trees, directly at fixed count]
\label{lem:finite-tree-limit}
For $0<t\le x$ write $w_t(x)=t^{-1}f(x/t)$. A shielded interval of
length $x$ and cutoff $t$ has the following limiting last-pair kernel:
\begin{equation}\label{eq:split-kernel}
 \mathcal K_{t,x}(a,s)
   =\frac{w_s(a)w_s(s)w_s(x-a-s)}{w_t(x)},\quad
 0<s<\min(t,x/3),\quad s<a<x-2s.
\end{equation}
It is a probability density. Conditional on this split, the three
daughter trees are independent and use states $(s,a)$, $(s,s)$, and
$(s,x-a-s)$.

Fix $0<\varepsilon<1/2$ and a finite depth $d$. Jointly with the active
configuration at cutoff $\varepsilon$, reveal the first $d$ last-pair
generations inside each active gap. At fixed even $M\to\infty$, this
finite geometry converges in total variation to
$C_\infty(\varepsilon)$ with independent gap trees generated by
\eqref{eq:split-kernel}. Conditional on a nonempty active configuration,
the probability that a finite-sample daughter is empty before the
specified depth tends to zero.
\end{lemma}
\begin{proof}
Differentiating \eqref{eq:density-delay}, almost everywhere, gives
\[
 \partial_t w_t(x)=w_t(t)\ind_{\{x>3t\}}
                 \int_t^{x-2t}w_t(a)w_t(x-a-t)\,da.
\]
The bound \eqref{eq:two-sided-f} gives $w_t(x)\to0$ as $t\downarrow0$,
for every fixed $x>0$. Integration in $t$ proves that
\eqref{eq:split-kernel} has mass one. This defines consistent finite
tree laws and hence an infinite recursive tree.

For the fixed-count assertion, condition by integration on a nonempty
active configuration with $k$ gaps. If every interval is split for $d$
generations, the terminal intervals number $k3^d$, an even number.
Let their lengths and inherited cutoffs be $(x_\ell,t_\ell)$. Their
lengths sum to one, and $x_\ell\ge t_\ell>0$. Repeated use of the last-pair
bijection gives the joint density
\begin{equation}\label{eq:tree-coefficient}
 M![\lambda^M]e^\lambda\prod_{\ell=1}^{k3^d}
                         w_{t_\ell}^{\lambda}(x_\ell).
\end{equation}
The reference integration has the same cyclic-root factor $1/k$ as
before and one $da\,ds$ for each revealed pair. Replacing one interval
factor by its three daughters accounts for its two exposed labels;
there is no further combinatorial factor.

At each fixed geometry expand every terminal factor using
\eqref{eq:finite-convolution}. All component lengths are at least
$\min_\ell t_\ell>0$ and together sum to one. Their total number is even.
The empty and full subsets in coefficient extraction therefore give
two, and every other subset tends to zero as even $M\to\infty$.
Thus \eqref{eq:tree-coefficient} tends pointwise to
$2\prod_\ell w_{t_\ell}(x_\ell)$.
This is exactly the active-configuration density times the product of
the split kernels: all parent denominators telescope.

Its total mass is $\PP(C_\infty(\varepsilon)\ne\varnothing)$ because
each kernel has mass one. The finite densities are nonnegative, and
their total mass is at most
$\PP(C_M(\varepsilon)\ne\varnothing)$, which converges to that same
number. Fatou's lemma forces convergence of their masses; the elementary
Scheffe argument then gives convergence in $L^1$. The missing mass,
corresponding to a daughter that becomes empty too early, tends to zero.
Add the empty active configuration as a separate atom. This proves
total-variation convergence of the finite geometry without conditioning
on a Poisson count or exchanging an infinite-depth limit with $M$.
\end{proof}

\begin{corollary}[An effective finite-depth error]\label{cor:tree-rate}
Let $\tau_M(\varepsilon,d)$ denote total variation between the two
finite-geometry laws in Lemma~\ref{lem:finite-tree-limit}, including
early-empty daughters as separate finite-sample states. For
$0<\alpha<1/3$, put $\eta=\varepsilon\alpha^d$ and
\begin{align*}
 r_\alpha&=\frac{b_+}{b_-}\alpha^{-1}
                         e^{-\kappa(1/\alpha-1)},&
 n_{\varepsilon,d}&=\frac{3^d-1}{2\varepsilon},\\
 C_{\varepsilon,d,\eta}
   &=\sum_{\substack{2\le k\le\lfloor1/\varepsilon\rfloor\\k\ {\rm even}}}
       \frac{(\eta^{-1}e^{1/\eta})^{k3^d}}
                         {k\,(k3^d-1)!}.
\end{align*}
Then
\begin{equation}\label{eq:tree-TV-rate}
 \tau_M(\varepsilon,d)\le
 \min\{1,\Delta_M(\varepsilon)+n_{\varepsilon,d}r_\alpha
            +C_{\varepsilon,d,\eta}(1-2\eta)^M\}.
\end{equation}
\end{corollary}
\begin{proof}
At every node of state $(t,x)$, \eqref{eq:two-sided-f} gives
\[
 \PP(s\le\alpha t)=\frac{w_{\alpha t}(x)}{w_t(x)}
 \le\frac{b_+}{b_-}\alpha^{-1}
          e^{-\kappa(x/t)(1/\alpha-1)}\le r_\alpha.
\]
Here $\alpha t<x/3$, since $t\le x$. There are at most
$n_{\varepsilon,d}$ revealed nodes. Outside a limiting event of
probability at most $n_{\varepsilon,d}r_\alpha$, each child cutoff is
at least $\alpha$ times its parent's, so every terminal cutoff is at
least $\eta$.

Restrict the finite and limiting densities to geometries with this last
property. The coefficient argument bounds their pointwise difference,
with $q=k3^d$ terminal intervals, by
$(\eta^{-1}e^{1/\eta})^q(1-2\eta)^M$.
Terminal gap lengths are unit-Jacobian simplex coordinates for the
revealed geometry. With the cyclic-root factor $1/k$, its volume is at
most $1/(k(q-1)!)$. Integration gives the last term of
\eqref{eq:tree-TV-rate}. The empty active atom differs by at most
$\Delta_M(\varepsilon)$. On the remaining geometries bound the limiting
mass by $n_{\varepsilon,d}r_\alpha$; the finite mass there is controlled
by normalization and the discrepancies on the other sectors.
The total-variation bound follows. Sending first $M$ to infinity and
then $\alpha$ to zero also recovers the finite-depth convergence.
\end{proof}

\begin{lemma}[A contracting backward mass query]\label{lem:query-contraction}
For the limiting interval tree, a query of the cumulative mass
$\mu([0,z])$ can be reduced after $d$ generations to queries on terminal
intervals whose total length is a random variable $Z_d(t,x,z)$.
For two completions of the unseen intervals, each having the prescribed
mass and spatial support and not charging the queried boundaries, the
difference between the queried masses is at most $Z_d$. The natural
finite-sample measures and the tree completions used below satisfy this
boundary condition almost surely. Uniformly in $0<t\le x$ and $0\le z\le x$,
\begin{equation}\label{eq:query-contraction}
 \EE Z_d(t,x,z)\le x\rho^d,\qquad
 \rho=1-\frac{f(6)}2\left(\frac{b_-}{b_+}\right)^2\in(0,1).
\end{equation}
\end{lemma}
\begin{proof}
At the split $(a,s,c)$, where $c=x-a-s$, let $F_i$ be the three daughter
cumulative masses, in their local coordinates. Put
$\alpha=\min(s,a/2)$ and $\beta=\min(s,c/2)$. Reading the two cells of
the final collapse map simultaneously gives
\begin{equation}\label{eq:cdf-query}
 \mu([0,z])=
 \begin{cases}
 F_1(\min(z,a-\alpha)),&z<a,\\
 F_1(a-\alpha)+s-F_2(s/2)+F_3(\beta),&a<z<a+s,\\
 a+s+F_3(\max(z-a-s,\beta)),&z>a+s.
 \end{cases}
\end{equation}
Endpoint queries are known from the total mass and need no descendants.
All other queries follow either one daughter or all three, with
coefficients of absolute value one. At depth zero give an unresolved
interval budget equal to its length. Iterating \eqref{eq:cdf-query}
proves the asserted completion bound: sum the lengths of the queried
terminal intervals. This budget is always at most $x$.
For finite-sample measures and for the uniform terminal completions used
below, atoms do not lie at the queried thresholds, except on null affine
coincidence sets. The limiting finite tree densities have the same
property. Values at such boundaries consequently do not change the
argument.

It remains to prove a uniform loss of budget in one generation.
Under \eqref{eq:split-kernel},
\[
 \PP(s\le x/6)\ge12f(6)>0.
\]
Indeed this probability is one when $t<x/6$; otherwise it equals
$w_{x/6}(x)/w_t(x)$, and $w_t(x)\le1/(2x)$.
Given $s$, the conditional density of $a$ is proportional to
$f(a/s)f((x-s-a)/s)$. The common exponential factor cancels, and
\eqref{eq:two-sided-f} bounds this density below by
\[
 \frac{(b_-/b_+)^2}{x-3s}\quad(s<a<x-2s).
\]
For $s\le x/6$, restrict $a$ to
$[x/4,\,3x/4-s]\setminus[z-s,z]$.
This set has length at least $x/6$. On it, the query is outside the
middle daughter and hence follows just one side daughter, whose length
is at most $3x/4$. The conditional probability of this event is at
least $(b_-/b_+)^2/6$. Everywhere else the retained total length is at
most $x$. Consequently its expectation is at most
\[
 x\left(1-\frac{12f(6)}{24}
                     \left(\frac{b_-}{b_+}\right)^2\right)=\rho x.
\]
This bound holds for every state and threshold. Conditional independence
of daughter trees and induction now give \eqref{eq:query-contraction}.
\end{proof}

\begin{theorem}[Marked convergence and a limiting noise probability]
\label{thm:noise-convergence}
For every fixed $0<\varepsilon<1/2$, the active configuration
$C_M(\varepsilon)$ with the full-process basin weights $\theta_v(B)$
attached to its points has a joint weak limit as even $M\to\infty$.
In particular $A_{M+1}^{\rm late}(\varepsilon)$ converges. Consequently
\begin{equation}\label{eq:noise-limit-exists}
 A_N\longrightarrow A_\infty,\qquad
 p_N\longrightarrow p_\infty=1-A_\infty,\qquad
 p_N^{\rm pos}\longrightarrow p_\infty
\end{equation}
as odd $N\to\infty$. The proved bounds become
$c_*\le p_\infty\le u_*$, with the exact constants from
\eqref{eq:cstar} and \eqref{eq:ustar}.

There is also a characterization by finite tree integrals. Given
$C_\infty(\varepsilon)$, reveal $d$ generations of every gap tree,
put Lebesgue measure on its terminal intervals, apply the revealed
collapse maps, then apply the finite future maps of the active
configuration. Let $a_{\varepsilon,d}$ be the expected sum of squared
masses on the active labels so obtained, set to zero for an empty
active configuration. Then
\begin{equation}\label{eq:limit-approximation}
 |A_\infty-a_{\varepsilon,d}|
 \le8e^\gamma(1+e^\gamma)\varepsilon
      +K e^{-(1/\varepsilon-1)/4}
      +\frac6\varepsilon\rho^d.
\end{equation}
This proves existence and gives a rigorous approximation scheme.
Writing
$B(\varepsilon)=8e^\gamma(1+e^\gamma)\varepsilon+
K e^{-(1/\varepsilon-1)/4}$, a finite-sample bound is
\begin{equation}\label{eq:finite-noise-error}
 |p_{M+1}-p_\infty|
 \le B(\varepsilon)+\frac{12}{\varepsilon}\rho^d
       +\Delta_M(\varepsilon)+\tau_M(\varepsilon,d)
       +\EE\Theta_{M+1}^2.
\end{equation}
Use \eqref{eq:tree-TV-rate} for any $0<\alpha<1/3$ to bound $\tau_M$.
These constants are deliberately conservative. The numerical candidate
$p_\infty\approx0.658$ and a useful sharp rate in $N$ remain unproved.
\end{theorem}
\begin{proof}
Fix $\varepsilon$. If $k$ points are active, there are $k/2$ future
removal maps. Their composition $G$ is the identity or a constant on
each of finitely many intervals. Its breakpoints are among the three
cell boundaries of each map, so there are at most $3k/2$. To see this,
on a constant piece of the current composition, postcomposition stays
constant; on an identity piece only the new map's three boundaries can
appear. Isolated boundary values can be disregarded here.
Every constant value belongs to the original active configuration.

Just before these future maps, the insertion measure is a sum of measures
supported inside the $k$ active gaps, each with mass its gap length.
For a late label $v$, its final basin mass is the measure of the constant
intervals on which $G=v$. The possible singleton $\{v\}$ in an identity
region has zero mass: none of the earlier internal maps targets an
active endpoint. Thus the late weights can be recovered from cumulative
mass queries at the breakpoints of $G$ inside the gaps.

Replace each unknown gap measure by its depth-$d$ tree approximation.
If the cumulative mass discrepancy at a breakpoint is at most $Z_d$,
the total $\ell^1$ discrepancy of the late-weight vector is at most
$2\sum Z_d$. Indeed consecutive interval masses are differences of
cumulative masses, and changing interval labels has $\ell^1$ cost at
most two. Since both weight vectors have coordinates in $[0,1]$,
their squared norms differ by at most $4\sum Z_d$.
Condition on the active geometry in its limiting law. Apply
Lemma~\ref{lem:query-contraction} separately to each breakpoint query;
each gap length is at most one and $k\le1/\varepsilon$. The resulting
expected errors are at most
\[
 \frac3\varepsilon\rho^d
 \quad\hbox{for the weight vector, and}\quad
 \frac6\varepsilon\rho^d
 \quad\hbox{for its squared norm}.
\]
Multiple queries in one gap need not be independent; linearity of
expectation is sufficient.

These estimates first compare two depths on a single infinite limiting
tree. The vectors of approximate weights form an $L^1$ Cauchy sequence
on that space, and so define a limiting marked configuration. Next
compare the true finite-sample weights with their depth-$d$ approximations.
The same deterministic query budget applies. By
Lemma~\ref{lem:finite-tree-limit}, the revealed geometry converges in
total variation, so the approximations and the bounded query budgets
converge in distribution and expectation. The probability of an empty
daughter before depth $d$ tends to zero; on that exceptional event use
the error bound one for agreement, or two for the weight vector.
Hence, writing
$a_\varepsilon=\lim_{M\to\infty}A_{M+1}^{\rm late}(\varepsilon)$,
the limit exists and
\[
 |a_\varepsilon-a_{\varepsilon,d}|
        \le(6/\varepsilon)\rho^d.
\]
More explicitly, existence follows because the upper and lower limits
are at distance at most $(12/\varepsilon)\rho^d$ for every $d$.
The vector estimate and the same approximation argument for bounded
Lipschitz test functions give the asserted joint weak convergence,
including the positions and the random count. This does not assert
total-variation convergence of the marks.

The early-label bound \eqref{eq:early-label-bound}, first with $M$
sent to infinity and then with $\varepsilon\downarrow0$, now proves
convergence of $A_N$. It also gives
\[
 0\le A_\infty-a_\varepsilon
       \le8e^\gamma(1+e^\gamma)\varepsilon
                          +K e^{-(1/\varepsilon-1)/4}.
\]
Combine this with the depth error to obtain
\eqref{eq:limit-approximation}. Finally
$\EE\Theta_N^2\to0$ and \eqref{eq:noise-identity} give the two noise
limits and the previously proved lower and upper bounds.
All limits used here are ordered: first $M$ at fixed $\varepsilon,d$,
then $d$, and finally $\varepsilon\downarrow0$.

For the finite error, maximally couple the finite and limiting revealed
geometries. They disagree with probability at most
$\tau_M(\varepsilon,d)$. When they agree, the deterministic completion
budget compares the true finite weights to the limiting depth-$d$
approximation; its expected squared-norm error is at most
$(6/\varepsilon)\rho^d$. A mismatch costs at most one.
Therefore
\[
 |A_{M+1}^{\rm late}(\varepsilon)-a_\varepsilon|
 \le(12/\varepsilon)\rho^d+\tau_M(\varepsilon,d).
\]
The two discarded early-label contributions lie respectively in
$[0,B(\varepsilon)+\Delta_M(\varepsilon)]$ and $[0,B(\varepsilon)]$,
so their difference is at most $B(\varepsilon)+\Delta_M(\varepsilon)$
in absolute value. Add the exact available-mass second moment to obtain
\eqref{eq:finite-noise-error}.
\end{proof}

\subsection{The active-location measure process and its two time scales}
\label{sec:measure-process}

\status{Working proofs added 10 October 2026}
For a sample of size $n$, let
\[
 \Xi_n(r)=\sum_{v\text{ active after cutoff }r}\delta_v,
 \qquad K_n(r)=\Xi_n(r)(\T).
\]
Here $r$ is separation: physical time is $r/2$. These measures count
\emph{currently active locations}; their weights are one, not insertion
basin masses. For odd $n$, also put $\widehat\Xi_n(r)=\Xi_n(r)/K_n(r)$.
Normalization by $n$ describes the early population, whereas normalization
by the current count preserves a nonzero late probability measure.

\begin{theorem}[Odd entrance law and full future paths]
\label{thm:odd-entrance}
For each $0<r<1/2$, as odd $n\to\infty$, $\Xi_n(r)$ converges in total
variation to a random finite counting measure $\Xi^{\rm o}(r)$.
With $h_r(x)=r^{-1}f(x/r)\ind_{\{x\ge r\}}$, its density on an odd
$k$-point sector is
\begin{equation}\label{eq:odd-sector}
 2\prod_{i=1}^k h_r(x_i)
 \quad\hbox{against }\frac1k\,da\,dx_1\cdots dx_{k-1},
 \qquad \sum_i x_i=1.
\end{equation}
For $k=1$, the root $a$ is the only point and its single cyclic gap has
length one. In particular,
\begin{equation}\label{eq:odd-count}
 \PP(K^{\rm o}(r)=k)=\frac2k h_r^{*k}(1),\qquad k\ge1\text{ odd}.
\end{equation}
For every $\varepsilon>0$, the entire path after $\varepsilon$ converges
in total variation. The limits fit together into a process on $r>0$
with finitely many jumps on every interval bounded away from zero.
Forward evolution is greedy deletion: at a collision separation $s$,
\begin{equation}\label{eq:counting-evolution}
 \Xi^{\rm o}(s)=\Xi^{\rm o}(s-)-\delta_u-\delta_v,
 \qquad (a,s,c)\longmapsto a+s+c
\end{equation}
for the three gaps adjacent to the removed pair. Once a single point
remains it stays forever. The normalized paths converge as well.
\end{theorem}
\begin{proof}
For $k\ge3$ the shielding decomposition in
Theorem~\ref{thm:active-limit} applies without alteration except for
parity. For $k=1$, cut the circle at its sole active point and use two
protected copies of that point as endpoints of an interval of length one.
Lemma~\ref{lem:shield} identifies the shielding event, including its
cutoff, with survival of this point. Thus the exact finite density is
\[
 G_{n,k}(\mathbf x)=n![\lambda^n]e^\lambda
                         \prod_i h_r^\lambda(x_i)
\]
also for every odd $k$, including one. In its convolution expansion the
total number of component lengths and $n$ are both odd. Consequently
the empty and full subsets in coefficient extraction again contribute
two; all other subsets have modulus at most $(1-2r)^n$.
The previous proof gives the valid bound
\begin{equation}\label{eq:odd-TV}
 \|\mathcal L(\Xi_n(r))-\mathcal L(\Xi^{\rm o}(r))\|_{\rm TV}
 \le\min\{1,D_r^{\rm o}(1-2r)^n\},\quad
 D_r^{\rm o}=\sum_{\substack{1\le k\le\lfloor1/r\rfloor\\k\ {\rm odd}}}
                  \frac{(r^{-1}e^{1/r})^k}{k!}.
\end{equation}
There is no empty sector. To check normalization directly, take the
Laplace transform in total circumference $x$. If $H_r(s)=\widehat h_r(s)$,
then
\[
 \sum_{k\ge1\text{ odd}}\frac2k H_r(s)^k
 =2\operatorname{arctanh}H_r(s)=E_1(rs).
\]
Inversion gives $\sum_{k\text{ odd}}(2/k)h_r^{*k}(x)=1/x$ for $x>r$;
evaluate at $x=1$. All sums at a fixed $x$ are finite.

The future path is a measurable deterministic function of the active
configuration. Total variation contracts under that map. Consistency at
the cutoffs $2^{-j}$ constructs one process on $r>0$; on each such
interval the usual finite greedy path supplies its right-continuous
version. It has at most $\lfloor1/r\rfloor$ points at cutoff $r$ and hence
only finitely many later collisions. Normalization by the positive count
is another measurable map. This proves the path assertions.
\end{proof}

\begin{corollary}[Last collision and final position for odd samples]
\label{cor:odd-last}
Let $R_n$ be the last collision separation for odd $n\ge3$. Then
\begin{equation}\label{eq:odd-last}
 R_n\ \Longrightarrow\ R_*\overset d=\frac1{1+Y_1+Y_2},\qquad
 \PP(R_*\le r)=\frac2r f(1/r),\quad 0<r<1/3,
\end{equation}
where $Y_1,Y_2$ are independent with density $f$.
The last physical collision time is $R_*/2$ in the limit.
The final location $U$ is uniform on the circle and independent of $R_*$.
For $r\ge1/3$, $\Xi^{\rm o}(r)=\delta_U$.
\end{corollary}
\begin{proof}
The event $R_n\le r$ is exactly $K_n(r)=1$. Formula
\eqref{eq:odd-count} gives its limiting probability $2h_r(1)$.
The delay identity \eqref{eq:density-delay}, at $y=1/r$, identifies this
with $\PP(Y_1+Y_2>1/r-1)$. This is the continuous distribution in
\eqref{eq:odd-last}. It has no mass at zero, which completes the
convergence assertion. Three active points have a gap at most $1/3$,
so every odd finite process is already a singleton by that cutoff.
Rotation invariance makes the final location uniform even conditional
on the rotation-invariant last separation; hence the independence.
\end{proof}

\begin{proposition}[A stochastic rule in reverse time]
\label{prop:reverse-process}
Start with $\delta_U$ at separation $1/3$ and run the cutoff $r$ downward.
Conditionally on the current configuration, different gaps evolve
independently until the next split. A gap $(u,u+x)$ has rate density
\begin{equation}\label{eq:reverse-rate}
 \frac{h_r(a)h_r(r)h_r(x-a-r)}{h_r(x)}\,da\,|dr|,
 \qquad r<a<x-2r,
\end{equation}
for insertion of the two points $u+a,u+a+r$ (modulo one).
The rate is zero when $x\le3r$. This constructs exactly
$\Xi^{\rm o}$ in reverse time, including the singleton's gap of length one.
\end{proposition}
\begin{proof}
The fixed-count tree proof of Lemma~\ref{lem:finite-tree-limit} applies
to the odd sector too: there are now $k3^d$ terminal intervals, an odd
number, so the endpoint subsets still contribute two. The singleton is
handled by the same cut-circle identification as above. Thus, conditional
on the current geometry at cutoff $t$, its gap trees are independent with
kernel \eqref{eq:split-kernel}.
For one gap of length $x$, its probability of no split while lowering
the cutoff from $t$ to $r$ is $h_r(x)/h_t(x)$. Moreover,
\[
 \partial_r h_r(x)=h_r(r)\int_r^{x-2r}
                         h_r(a)h_r(x-a-r)\,da\quad(x>3r).
\]
Dividing the split density by the probability of no earlier reverse
split gives \eqref{eq:reverse-rate}. Its integral is
$\partial_r\log h_r(x)$; conditional daughter trees have the same rule.
There is no explosion above a positive cutoff because every gap then
has length at least that cutoff. At $r=1/3$ the odd process is already
the uniform singleton, so its reverse construction gives all the stated
laws. Forward in time the randomness is in the entrance configuration;
given that configuration, the future has no additional random clocks.
\end{proof}

\begin{theorem}[Deterministic empirical limit on the microscopic clock]
\label{thm:early-fluid}
For fixed $s\ge0$ and every continuous real $\phi$ on $\T$, as $n\to\infty$
through either parity,
\begin{equation}\label{eq:early-fluid}
 \frac1n\langle\Xi_n(s/n),\phi\rangle
 \longrightarrow q(s)\int_{\T}\phi(x)\,dx\quad\hbox{in }L^2,
 \qquad
 q(s)=\exp\left(-\int_0^{2s}\frac{1-e^{-u}}u\,du\right).
\end{equation}
For each finite $S$, convergence holds uniformly over $0\le s\le S$ in
probability. The mass profile satisfies
\begin{equation}\label{eq:fluid-ode}
 q(0)=1,\qquad q'(s)=-\frac{1-e^{-2s}}s q(s),\qquad
 q(s)\sim\frac1{2e^\gamma s}\quad(s\to\infty).
\end{equation}
\end{theorem}
\begin{proof}
First take a Poisson sample of intensity $n$ on the circle and cutoff
$s/n$, with $s>0$. Summing the exact sector densities over the initial
sample count cancels coefficient extraction and gives the density
$\prod_i h_{s/n}^{n}(x_i)$, including $k=1$; the empty sector contributes
to neither of the following moments.
Set $h_s^{(1)}=h_s^1$ and
$\mathcal U_s=\sum_{j\ge1}(h_s^{(1)})^{*j}$.
The scalar transform at $\lambda=1$ shows that $h_s^{(1)}$ is a probability
density, with mean
\begin{equation}\label{eq:finite-intensity-mean}
 m_s=2e^\gamma s\exp(E_1(2s)).
\end{equation}
Indeed its transform is
$\tanh((E_1(sz)-E_1(s(z+2)))/2)=1-m_sz+o(z)$ as $z\downarrow0$.

Here the renewal-density theorem applies with all its hypotheses
verified. The density is continuous on $[s,\infty)$ and bounded on
compact intervals, from its Poisson series or finite convolution formula.
It has an exponential envelope: in a long shielded interval, any interior
Poisson point with no other point within distance $s$ survives the cutoff,
which prevents clearance. In disjoint blocks of length $2s+1$, the event
of exactly one point in the middle unit interval and no other point in
the block has probability $e^{-(2s+1)}>0$; such a point is isolated by
more than $s$. All blocks lying inside the interval must avoid this event.
Since $h_s^{(1)}(x)$ is the corresponding Poisson shielding probability
at intensity one, independence bounds it by $C_s e^{-c_s x}$.
Consequently it is directly Riemann integrable and nonarithmetic, with
finite mean. Applying the key renewal theorem to
$h_s^{(1)}*(\delta_0+\sum_{j\ge1}(h_s^{(1)})^{*j})$ gives
\[
 \mathcal U_s(y)\longrightarrow m_s^{-1},\qquad
 \sup_{y\ge0}\mathcal U_s(y)<\infty.
\]
The bound also follows by writing the renewal sum as an expectation
along points spaced at least $s$ and summing the exponential envelope.
We use the usual key renewal theorem: for nonarithmetic increments of
mean $m$ and directly Riemann integrable forcing $b$, the convolution of
$b$ with the renewal measure tends to $m^{-1}\int b$; see
\href{https://ocw.mit.edu/courses/6-262-discrete-stochastic-processes-spring-2011/931ffa0940899c27f34b71ad64fd2bb0_MIT6_262S11_chap04.pdf}{Gallager, \emph{Discrete Stochastic Processes}, Theorem 4.7.3}.

Scaling gives $h_{s/n}^n(x)=n h_s^{(1)}(nx)$.
Choosing one distinguished active point cancels the cyclic-root factor.
Choosing an ordered pair splits the cycle into two nonempty strings of
gaps. The first intensity and ordered second factorial intensity at
displacement $x\in(0,1)$ are therefore, respectively,
\[
 n\mathcal U_s(n),\qquad
 n^2\mathcal U_s(nx)\mathcal U_s(n(1-x)).
\]
The renewal limit and its uniform bound prove convergence of the first
moment and second factorial moment of the normalized test integral to
$m_s^{-1}\int\phi$ and $m_s^{-2}(\int\phi)^2$.
The diagonal second-moment term is $O(1/n)$, so the variance vanishes.
The identity $m_s^{-1}=q(s)$ follows from
$\gamma+\log(2s)+E_1(2s)=\int_0^{2s}(1-e^{-u})\,du/u$.

For de-Poissonization, couple the first $n$ points and the first $P_n$
points of one iid uniform sequence, where $P_n$ is independent Poisson
of mean $n$. Repeated application of Lemma~\ref{lem:discrepancy} bounds
their active-set symmetric difference at every cutoff by $|P_n-n|$.
The difference of their normalized test integrals has squared expectation
at most $\|\phi\|_\infty^2/n$. This proves the fixed-count assertion
without a monotonicity assumption in sample size.
At $s=0$ the result is the ordinary iid empirical second-moment calculation.
For nonnegative $\phi$, the empirical integral decreases with $s$;
a finite grid and continuity of $q$ give uniform convergence in probability
on compact time intervals. Subtracting a constant handles signed $\phi$.
Differentiation and $E_1(2s)\to0$ prove \eqref{eq:fluid-ode}.
\end{proof}

\begin{theorem}[The small-cutoff end of the random entrance law]
\label{thm:entrance-fluid}
Let $m=\EE Y=2e^\gamma$ and $K^{\rm o}(r)=\Xi^{\rm o}(r)(\T)$.
For every continuous real $\phi$,
\begin{equation}\label{eq:entrance-fluid}
 r\langle\Xi^{\rm o}(r),\phi\rangle
 \longrightarrow m^{-1}\int_{\T}\phi(x)\,dx\quad\hbox{in }L^2
 \quad(r\downarrow0).
\end{equation}
In particular $rK^{\rm o}(r)\to m^{-1}$ and
$\widehat\Xi^{\rm o}(r)\to dx$ weakly in probability.
The same unnormalized assertion holds for the even entrance law, with
the counting measure zero on its empty configuration.
\end{theorem}
\begin{proof}
Define the locally finite renewal densities
\[
 R_{\rm o}(y)=\sum_{j\ge0}f^{*(2j+1)}(y),\qquad
 R_{\rm e}(y)=\sum_{j\ge1}f^{*(2j)}(y).
\]
The proof of Lemma~\ref{lem:gap-squares} gives $R_{\rm o}\le1$;
also $R_{\rm e}=f*R_{\rm o}\le1$.
Use the key renewal theorem with increment density $f*f$, mean $2m$,
and forcing $f$ or $f*f$. These forcings are directly Riemann integrable
by the exponential bounds of Lemma~\ref{lem:positive-renewal}; the
increment law has a density and hence is nonarithmetic. Thus
\begin{equation}\label{eq:parity-renewal}
 R_{\rm o}(y),R_{\rm e}(y)\longrightarrow\frac1{2m}.
\end{equation}
Put $R_{{\rm o},r}(x)=r^{-1}R_{\rm o}(x/r)$ and similarly for the even
density. From \eqref{eq:odd-sector}, the first intensity is
$2R_{{\rm o},r}(1)$. The ordered second factorial intensity is
\begin{equation}\label{eq:odd-two-point}
 2\big[R_{{\rm o},r}(x)R_{{\rm e},r}(1-x)
       +R_{{\rm e},r}(x)R_{{\rm o},r}(1-x)\big],\quad 0<x<1.
\end{equation}
Indeed the two strings of gaps between distinguished points must have
opposite parities. Multiplying the intensities by $r$ and $r^2$,
respectively, and using \eqref{eq:parity-renewal} gives limits $m^{-1}$
and $m^{-2}$ almost everywhere, bounded by two and four.
Dominated convergence and the diagonal term $O(r)$ prove the $L^2$ claim.
Division by the count gives the normalized assertion.
For even parity the first intensity is $2R_{{\rm e},r}(1)$, and the
second is $2[R_{{\rm o},r}(x)R_{{\rm o},r}(1-x)+
R_{{\rm e},r}(x)R_{{\rm e},r}(1-x)]$. The same bounds and limits apply.
\end{proof}

\begin{lemma}[A factorial renewal remainder]\label{lem:factorial-renewal}
Put $U=R_{\rm o}+R_{\rm e}$ and $c=1/m=1/(2e^\gamma)$, extending the
densities by zero below their support. For $y\ge1$,
\begin{equation}\label{eq:renewal-delay}
 yU(y)=\frac12+\int_0^{y-1}U(v)\,dv,\qquad
 \frac14\le U(y)\le\frac12.
\end{equation}
For every integer $j\ge1$ and every $y\ge j$,
\begin{align}
 |U(y)-c|&\le\frac1{4j!},\label{eq:renewal-factorial}\\
 a_0(y)&\le\frac1{j!},\qquad
 \left|R_{\rm o}(y)-\frac c2\right|,
 \left|R_{\rm e}(y)-\frac c2\right|
 \le\frac3{8j!}.\label{eq:parity-factorial}
\end{align}
In particular $R_{\rm o},R_{\rm e}\le1/2$ everywhere. Also
\begin{equation}\label{eq:renewal-excess}
 \int_0^\infty(U(y)-c)\,dy=c-\frac12.
\end{equation}
The integrand in \eqref{eq:renewal-excess} has finite absolute moments
of every order.
\end{lemma}
\begin{proof}
The scalar transform gives
$\widehat U(s)=(e^{E_1(s)}-1)/2$. Differentiating and inverting proves
the first identity in \eqref{eq:renewal-delay}; on $[1,2]$ it gives
$U(y)=1/(2y)$. For integers $j\ge2$ and $j\le y\le j+1$,
\[
 U(y)=\frac jyU(j)+\frac1y\int_j^yU(v-1)\,dv.
\]
Thus the range on $[j,j+1]$ is contained in the range on $[j-1,j]$.
If $D_j$ is the oscillation on $[j,j+1]$, then
$U'(y)=(U(y-1)-U(y))/y$ gives
$D_j\le D_{j-1}\log(1+1/j)\le D_{j-1}/j$.
Since $D_1=1/4$, $D_j\le1/(4j!)$. The already established limit $c$
belongs to each of these nested ranges. This proves the bounds for $U$.

Write $\varrho=e^\gamma\varphi$, with $\varphi$ from the Dickman
construction in Lemma~\ref{lem:positive-renewal}. Its delay equation gives
$v\varrho(v)=\int_{v-1}^v\varrho(u)\,du$ for $v\ge1$.
As $\varrho$ is nonincreasing and equals one on $[0,1]$,
$\varrho(v)\le\varrho(v-1)/v$. Iteration yields
$a_0(y)=\varrho(y-1)/y\le1/j!$ for $y\ge j$.
The positive representation of $f$ gives
\[
 R_{\rm o}-R_{\rm e}=\frac12a_0,
 \qquad R_{\rm o}=\frac12U+\frac14a_0,
 \qquad R_{\rm e}=\frac12U-\frac14a_0.
\]
These prove \eqref{eq:parity-factorial}; the global upper bound also
uses nonnegativity and $R_{\rm o},R_{\rm e}\le U$.
Finally the factorial remainder makes $U-c$ integrable with every
polynomial weight. As $s\downarrow0$,
\[
 \widehat U(s)=\frac c s+c-\frac12+O(s),
\]
using $E_1(s)=-\gamma-\log s+s+O(s^2)$. Subtract $c/s$ and use
dominated convergence to obtain \eqref{eq:renewal-excess}.
\end{proof}

\begin{lemma}[Uniform renewal comparison at high intensity]
\label{lem:uniform-renewal}
Let $h^{(\alpha)}=h_1^\alpha$ for intensity $\alpha>0$, and define
\[
 U_\alpha=\sum_{j\ge1}(h^{(\alpha)})^{*j},\qquad
 U_\infty=\sum_{j\ge1}f^{*j}=R_{\rm o}+R_{\rm e}.
\]
Then $U_\infty\le2$, $U_\infty(y)\to1/m$ as $y\to\infty$, and
\begin{equation}\label{eq:uniform-renewal}
 \sup_{y\ge0}|U_\alpha(y)-U_\infty(y)|\le2e^{-2\alpha},
 \qquad \alpha\ge1.
\end{equation}
\end{lemma}
\begin{proof}
The first two assertions follow from \eqref{eq:parity-renewal} and its
proof. The scalar transform gives, for $z>0$,
\[
 \widehat U_\alpha(z)=\frac12
   \left(\exp(E_1(z)-E_1(z+2\alpha))-1\right),\qquad
 \widehat U_\infty(z)=\frac12(e^{E_1(z)}-1).
\]
Put $v_\alpha(y)=e^{-2\alpha y}\ind_{\{y\ge1\}}/y$,
$\beta_\alpha=\int v_\alpha=E_1(2\alpha)$, and let
\[
 \nu_\alpha=\sum_{j\ge1}\frac{(-1)^j}{j!}v_\alpha^{*j}.
\]
This is a signed integrable density with
$\|\nu_\alpha\|_1\le e^{\beta_\alpha}-1$ and
$\|\nu_\alpha\|_\infty\le e^{-2\alpha}e^{\beta_\alpha}$.
Taking transforms and then inverting proves
\[
 U_\alpha-U_\infty=\tfrac12\nu_\alpha+U_\infty*\nu_\alpha.
\]
Consequently its uniform norm is at most
$e^{\beta_\alpha}(\tfrac12e^{-2\alpha}+2\beta_\alpha)$.
For $\alpha\ge1$, use
$\beta_\alpha\le e^{-2\alpha}/(2\alpha)\le1/8$ and
$e^{\beta_\alpha}\le(1-\beta_\alpha)^{-1}\le8/7<4/3$.
This gives \eqref{eq:uniform-renewal}.
\end{proof}

\begin{lemma}[Fixed-count first and second intensities]
\label{lem:intensity-coefficients}
Let $I_n(r)$ be the constant first intensity of $\Xi_n(r)$ and
$J_n(r;x)$ its ordered second factorial intensity at displacement
$x\in(0,1)$. Write $I_\infty^{[n]}(r),J_\infty^{[n]}(r;x)$ for those
of the entrance law having the parity of $n$. For $0<r\le e^{-1}$,
\begin{align}
 |I_n(r)-I_\infty^{[n]}(r)|
   &\le\tfrac12 r^{-3}e^{-2nr},\label{eq:first-intensity-error}\\
 \sup_{0<x<1}|J_n(r;x)-J_\infty^{[n]}(r;x)|
   &\le\tfrac14 r^{-6}e^{-2nr}.\label{eq:second-intensity-error}
\end{align}
Density versions are chosen by the formulas in the proof; boundary values
on null threshold sets do not affect any application.
\end{lemma}
\begin{proof}
For general intensity $\lambda$, put
$U_r^\lambda=\sum_{j\ge1}(h_r^\lambda)^{*j}$ and
$g_r^\lambda(y)=\ind_{\{y\ge r\}}(1-e^{-2\lambda y})/y$.
The transform identity $\tanh z/(1-\tanh z)=(e^{2z}-1)/2$ gives the
locally finite positive convolution expansion
\begin{equation}\label{eq:renewal-expansion}
 U_r^\lambda(y)=\frac12\sum_{j\ge1}
                               \frac{(g_r^\lambda)^{*j}(y)}{j!}.
\end{equation}
The Poisson sector formulas yield exactly
\[
 I_n(r)=n![\lambda^n]e^\lambda U_r^\lambda(1),\qquad
 J_n(r;x)=n![\lambda^n]e^\lambda
                         U_r^\lambda(x)U_r^\lambda(1-x).
\]
Expand each product $\prod_i(1-e^{-2\lambda y_i})$, whose component
lengths sum to one. Empty and full subsets contribute
$1+(-1)^{n+j}$, where $j$ is the total component count. These terms
give exactly the limiting intensities of parity $n$. Every other subset
has $r\le\sum_{i\in S}y_i\le1-r$, so its extracted coefficient has
modulus at most $(1-2r)^n\le e^{-2nr}$.
For clarity, the signed full-subset convolution series is
$U_-=-(R_{{\rm o},r}-R_{{\rm e},r})$, whereas the empty-subset series is
$U_+=R_{{\rm o},r}+R_{{\rm e},r}$. Thus the first retained density is
$U_++(-1)^nU_-$ and the second is
$U_+(x)U_+(1-x)+(-1)^nU_-(x)U_-(1-x)$, exactly the parity formulas above.

For the remaining absolute sums put $g_r(y)=\ind_{\{y\ge r\}}/y$ and
$V_r(y)=\sum_{j\ge1}2^j g_r^{*j}(y)/j!$.
For $y\le1$ all factors can be restricted to $[r,1]$, where their
uniform norm is $r^{-1}$ and their integral is $\ell=\log(1/r)\ge1$.
Hence
\[
 V_r(y)\le\frac{e^{2\ell}-1}{r\ell}\le r^{-3}.
\]
The first remainder is at most $e^{-2nr}V_r(1)/2$; the second is at
most $e^{-2nr}V_r(x)V_r(1-x)/4$. These are the claimed bounds.
\end{proof}

\begin{theorem}[The full intermediate-time regime]
\label{thm:full-mesoscopic}
For either parity, if $r_n\to0$ and
$nr_n\to\infty$, then for every continuous real $\phi$ on $\T$,
\begin{equation}\label{eq:full-mesoscopic}
 r_n\langle\Xi_n(r_n),\phi\rangle
       \longrightarrow m^{-1}\int_{\T}\phi(x)\,dx\quad\hbox{in }L^2.
\end{equation}
For every $0<a<b<\infty$, uniformly in $a\le u\le b$ in probability,
\[
 r_n\langle\Xi_n(r_nu),\phi\rangle
       \longrightarrow \frac1{mu}\int_{\T}\phi(x)\,dx.
\]
In particular the normalized active-location measure tends to uniform
measure throughout this regime. An arbitrary convention can be used when
the even process is empty, since that event has probability tending to zero.
\end{theorem}
\begin{proof}
First use a Poisson sample of intensity $n$. Set $\alpha_n=nr_n$.
The first and second factorial intensities, multiplied by $r_n$ and
$r_n^2$, respectively, are
\[
 U_{\alpha_n}(1/r_n),\qquad
 U_{\alpha_n}(x/r_n)U_{\alpha_n}((1-x)/r_n).
\]
Lemma~\ref{lem:uniform-renewal} compares these uniformly with $U_\infty$.
They tend to $1/m$ and $1/m^2$ almost everywhere and are uniformly
bounded. The diagonal second-moment term is $O(r_n)$, so the same moment
argument as before proves the Poisson assertion for every such sequence.

There are two overlapping ways to return to fixed count. If
$r_n\le n^{-2/3}$, the add-one coupling bounds the squared difference
of scaled test integrals by
$nr_n^2\|\phi\|_\infty^2\le n^{-1/3}\|\phi\|_\infty^2$.
If $r_n>n^{-2/3}$, use Lemma~\ref{lem:intensity-coefficients} directly.
After multiplication by $r_n$ and $r_n^2$, its errors are bounded by
$\tfrac12n^{4/3}e^{-2n^{1/3}}$ and
$\tfrac14n^{8/3}e^{-2n^{1/3}}$, respectively. They vanish. The entrance-law
intensities have the required small-$r$ limits for both parities by
Theorem~\ref{thm:entrance-fluid}. Thus the fixed-count first and second
moments converge on this subsequence too. These two cases cover every
$n$, proving \eqref{eq:full-mesoscopic}. The dividing power $2/3$ is only
a device in the proof, not a process transition scale.

For nonnegative $\phi$, the test integral decreases with $u$.
A finite grid and continuity of $1/u$ on $[a,b]$ give uniform convergence
in probability. Subtract a constant for signed tests, and divide by the
scaled count for normalized measures.
\end{proof}

\begin{corollary}[Every vanishing cutoff has a uniform empirical profile]
\label{cor:all-vanishing-cutoffs}
For any deterministic $r_n\ge0$ tending to zero,
\[
 \frac{\Xi_n(r_n)}{nq(nr_n)}\longrightarrow dx
 \quad\hbox{weakly in probability},\qquad nq(nr_n)\longrightarrow\infty.
\]
Consequently normalization by the current count gives uniform measure
in probability, under the same empty-state convention for even samples.
\end{corollary}
\begin{proof}
On subsequences where $nr_n$ is bounded, use the uniform microscopic
law and the positive minimum of $q$ on compact intervals. On subsequences
where $nr_n\to\infty$, use \eqref{eq:full-mesoscopic} and
$q(s)\sim1/(ms)$. Every subsequence has a further subsequence of one
of these types, which proves the assertion. The same two cases show
that the deterministic denominator tends to infinity.
\end{proof}

These results specify the scales. On the clock $s=nr=2n\,t_{\rm physical}$,
the mass $n^{-1}\Xi_n$ follows $q(s)dx$ for bounded $s$.
The matching relation $q(s)\sim(2e^\gamma s)^{-1}$ agrees with the
small-$r$ entrance law. At fixed positive physical time, the natural
nonzero limit is the random counting measure, or its normalization by
the current count; normalization by $n$ instead gives zero. No additional
$n$-dependent magnification near the last collision is needed on the
unit circle. The full intermediate regime is covered by
Theorem~\ref{thm:full-mesoscopic}; every vanishing cutoff has a uniform
normalized profile. The next result gives the leading covariance and a
Gaussian count law at the small-cutoff end of either random entrance process.

\begin{theorem}[Entrance-process fluctuations]\label{thm:entrance-fluctuations}
For either entrance parity let $\Xi(r)$ be its counting measure and
$K(r)=\Xi(r)(\T)$, with the even measure zero on its empty state.
Set $c=1/(2e^\gamma)$ and $v=2c^2=1/(2e^{2\gamma})$.
For continuous real $\phi,\psi$ on $\T$,
\begin{equation}\label{eq:entrance-covariance}
 r\operatorname{Cov}(\langle\Xi(r),\phi\rangle,
                     \langle\Xi(r),\psi\rangle)
       \longrightarrow v\int_{\T}\phi(x)\psi(x)\,dx.
\end{equation}
Moreover,
\begin{equation}\label{eq:entrance-count-clt}
 \sqrt r\big(K(r)-c/r\big)\ \Longrightarrow\ \mathcal N(0,v).
\end{equation}
In particular $\operatorname{Var}K(r)\sim v/r$, and its ratio to
$\EE K(r)$ tends to $e^{-\gamma}<1$.
More generally, the fields
\[
 G_r(\phi)=\sqrt r\left(\langle\Xi(r),\phi\rangle
                            -\frac c r\int_{\T}\phi\right)
\]
converge jointly on every finite family of continuous tests to a centered
Gaussian field $G$ with covariance $\EE G(\phi)G(\psi)=v\int\phi\psi$.
For the current-count-normalized measure $\widehat\Xi(r)$,
\begin{equation}\label{eq:normalized-gaussian}
 r^{-1/2}\left(\langle\widehat\Xi(r),\phi\rangle-\int_{\T}\phi\right)
 \ \Longrightarrow\ B(\phi):=\frac{G(\phi)-G(1)\int\phi}{c},
\end{equation}
jointly, where
\[
 \EE B(\phi)B(\psi)=2\left(\int_{\T}\phi\psi
                              -\int_{\T}\phi\int_{\T}\psi\right).
\]
Both field convergences hold also in the negative Sobolev space
$H^{-a}(\T)$ for every $a>1/2$. Use any probability measure on the even
empty event. These are small-cutoff results for the entrance laws.
Theorem~\ref{thm:full-finite-spatial-CLT} separately gives simultaneous
finite-count one-time limits throughout the mesoscopic regime, with
finite-intensity centering. Section~\ref{sec:temporal-spatial-CLT}
separately proves temporal Gaussian and cumulative path limits.
\end{theorem}
\begin{proof}
Put $R=1/r$, $D=a_0/2$, $U=R_{\rm o}+R_{\rm e}$, and take
$\sigma=1$ for even parity, $\sigma=-1$ for odd parity. The first and
ordered second factorial intensities are
\[
 r^{-1}\big(U(R)-\sigma D(R)\big),\qquad
 r^{-2}\big[U(x/r)U((1-x)/r)+\sigma D(x/r)D((1-x)/r)\big].
\]
These are the parity formulas already proved. Write $e(y)=U(y)-c$ and
$C_{\phi,\psi}(x)=\int_{\T}\phi(a)\psi(a+x)\,da$.
By Lemma~\ref{lem:factorial-renewal}, $e$ is bounded and integrable,
has a factorial tail, and $\int e=c-1/2$; $D$ also has a factorial tail.
After changing variables $x=ry$, the off-diagonal part of $r$ times
the covariance is, up to a term tending to zero,
\[
 \int_0^R\big[c e(y)+c e(R-y)+e(y)e(R-y)
                      +\sigma D(y)D(R-y)\big]C_{\phi,\psi}(ry)\,dy.
\]
The replacement of the squared first intensity by $c^2$ costs $o(1)$
because its error, multiplied by $R$, tends to zero. Dominated convergence
in the first two integrals gives $2c(c-1/2)\int\phi\psi$.
The other two integrals tend to zero: split at $R/2$, use the factorial
tail on the long argument and integrability on the short one.
The diagonal contribution tends to $c\int\phi\psi$. Their sum is
$2c^2\int\phi\psi$, proving \eqref{eq:entrance-covariance}.

We prove the Gaussian assertion by comparing the circle with ordinary
renewals. Rescale the circle to circumference $R$. Tilt its law by
$K/\EE K$, choose one of its points uniformly, and put that point at zero.
The rooted cyclic gaps then have density
\begin{equation}\label{eq:palm-renewal-bridge}
 \frac{\prod_{i=1}^k f(y_i)}{R_p(R)},\qquad
 \sum_i y_i=R,\quad k\equiv p\pmod2,
\end{equation}
where $R_1=R_{\rm o}$ and $R_0=R_{\rm e}$.
Indeed $\EE K=2R R_p(R)$ and the unrooted count probability is
$2R f^{*k}(R)/k$, so the tilt cancels the cyclic-root factor.
An independent uniform rotation recovers the tilted unrooted law.
The total variation between the original and tilted laws is at most
$\EE|K/\EE K-1|/2=O(R^{-1/2})$ by the covariance result just proved.

Let $S_0=0$, $S_j=Y_1+\cdots+Y_j$ for iid increments of density $f$.
Put $L=R^{1/4}$ and $\tau=\min\{j:S_j\ge R-L\}$.
On prefixes with $S_\tau<R$, the density of the stopped prefix under
\eqref{eq:palm-renewal-bridge}, relative to its iid law, is
\[
 \frac{R_{p-\tau\bmod2}(R-S_\tau)}{R_p(R)}.
\]
On $S_\tau\le R-L/2$ this ratio tends to one uniformly, by
\eqref{eq:parity-factorial}. Under the iid law this event has probability
tending to one: $\tau\le\lceil R\rceil+1$, and the probability any of
those increments exceeds $L/2$ is at most
$7(\lceil R\rceil+1)e^{-L/8}$.
Integrating the density on this event shows that its bridge probability
also tends to one. Normalization then proves total-variation convergence
of the stopped prefixes. This argument includes in the exceptional mass
the possible bridge prefix ending exactly at $R$.
The remaining arc has length $L$ and contains at most $L+1$ points,
since all gaps are at least one. Its contribution divided by $\sqrt R$
vanishes. Thus the Gaussian problem reduces to an ordinary renewal
process on $[0,R]$.

For completeness, its required central limit argument is elementary.
Write $N(t)=\max\{j:S_j\le t\}$. The iid increment variance is
$s_Y^2=\EE Y^2-m^2=4e^\gamma=2m$.
For each fixed positive $t$, Chebyshev's inequality applied to partial
sums gives $N(Rt)-cRt=O_{\PP}(\sqrt R)$.
Kolmogorov's maximal inequality shows that the centered sum of any
$O(\sqrt R)$ consecutive increments near index $\lfloor cRt\rfloor$,
divided by $\sqrt R$, tends to zero in probability. The overshoot is
also $o_{\PP}(\sqrt R)$, from the exponential tail and a union bound.
Inverting $S_j$ therefore gives, jointly for finitely many $t$,
\[
 \frac{N(Rt)-cRt}{\sqrt R}
 =-\frac{S_{\lfloor cRt\rfloor}-m\lfloor cRt\rfloor}
           {m\sqrt R}+o_{\PP}(1).
\]
The iid central limit theorem on disjoint blocks gives Gaussian limits
with covariance $v\min(t,u)$, since $s_Y^2/m^3=v$.
This proves the test-function assertion for step functions.

To pass to continuous tests, the ordinary renewal process has first
intensity $U(x)$ and ordered second intensity $U(x)U(y-x)$ for $0<x<y$.
Using $U=c+e$, $U\le1/2$, and the finite absolute first moment of $e$
gives, for bounded $b$ and $R\ge1$,
\[
 \EE\left[\left(\sum_{S_j<R}b(S_j/R)-cR\int_0^1b\right)^2\right]
       \le C R\|b\|_\infty^2.
\]
Indeed the off-diagonal covariance is bounded by a constant times
$\|b\|_\infty^2\int_{0<x<y<R}(|e(y-x)|+|e(y)|)\,dx\,dy$,
which is $O(R)\|b\|_\infty^2$; the diagonal is also $O(R)$, and the
bias including the root at zero is bounded. Uniform step approximation
now proves the assertion for continuous tests.
For each fixed rotation the covariance is the same $v\int\phi\psi$.
Integrating characteristic functions over the independent uniform
rotation, and then removing the tilt by total variation, proves the
assertion for the entrance law, including \eqref{eq:entrance-count-clt}.

The covariance calculation also gives a bound for
$\EE|G_r(\phi)|^2$ by $C\|\phi\|_\infty^2$, uniformly for small $r$.
Apply it to the trigonometric basis. For $1/2<b<a$, the expected squared
$H^{-b}$ norm is uniformly bounded, and the embedding
$H^{-b}\subset H^{-a}$ is compact (truncate the Fourier series).
This gives tightness in $H^{-a}$, and the finite-dimensional limits
identify the Gaussian field there. Finally $rK(r)\to c$ and
\[
 r^{-1/2}(\widehat\Xi(r)-dx)
       =\frac{G_r-G_r(1)\,dx}{rK(r)}
\]
on the nonempty event. The continuous mapping theorem gives
\eqref{eq:normalized-gaussian} and its Sobolev version. The even empty
probability tends to zero and has no effect on these distributional limits.
\end{proof}

There is a relevant existing coalescence literature. The transform
$\tanh(E_1(s)/2)$ also occurs in the pure triple-merge case of
\href{https://arxiv.org/abs/1112.4630}{Faggionato, Roberto and Toninelli,
\emph{Universality for one-dimensional hierarchical coalescence processes
with double and triple merges}, Theorem 2.12 and Remark 2.14}.
That comparison concerns the scalar gap law; no equivalence of their
epoch dynamics with the finite-circle process is assumed here, and no
novelty claim is made for that transform.

\subsection{Sharper scalar and query budgets for the noise approximation}
\label{sec:sharper-budgets}

The measure-process calculation also improves a quantitative input to
the original noise problem. Write $R_{{\rm o},\varepsilon}$ as above and set
\begin{equation}\label{eq:sharp-cutoff-budget}
 \mathcal B(\varepsilon)=\PP(L_\infty\le\varepsilon)
  +2\int_0^1x^2h_\varepsilon(x)
                       R_{{\rm o},\varepsilon}(1-x)\,dx.
\end{equation}
This is the exact expectation of the bounded functional $H$ used in
Theorem~\ref{thm:early-labels}, including its value one on the empty
even configuration. Thus $\mathcal B\le B$, and all noise approximation
bounds using $B(\varepsilon)$ remain valid with $\mathcal B(\varepsilon)$.
Moreover,
\begin{equation}\label{eq:sharp-cutoff-asymptotic}
 \mathcal B(\varepsilon)\sim
 \frac{\EE Y^2}{\EE Y}\varepsilon
 =2(1+e^\gamma)\varepsilon.
\end{equation}
To prove this, the integral in \eqref{eq:sharp-cutoff-budget} equals
$2\varepsilon\int_0^{1/\varepsilon}y^2 f(y)
R_{\rm o}(1/\varepsilon-y)\,dy$.
Use \eqref{eq:parity-renewal}, $R_{\rm o}\le1$, and the finite second
moment. The empty probability is exponentially small. This replaces
the old cutoff bound's leading coefficient $8e^\gamma(1+e^\gamma)$ by
the asymptotic coefficient $2(1+e^\gamma)$ for the exact scalar budget.
The asymptotic itself is not a certified error bound at a specified cutoff.

\begin{proposition}[An explicit finite-cutoff error]
\label{prop:certified-cutoff}
Write $R=1/\varepsilon$, $\mu_2=\EE Y^2=4e^\gamma(1+e^\gamma)$,
and let $K=\EE e^{Y/4}<7$ as in \eqref{eq:last-TV}.
For any integer $j\ge1$ with $v:=R-j\ge8$,
\begin{equation}\label{eq:finite-cutoff-budget}
 \mathcal B(\varepsilon)\le
 \varepsilon\left(2(1+e^\gamma)+\frac{3\mu_2}{4j!}\right)
 +K e^{-(R-1)/4}+\varepsilon K v^2e^{-v/4}.
\end{equation}
In particular the following outward bounds are certified:
\[
 \mathcal B(1/100)<0.055641,\qquad
 \mathcal B(1/200)<0.027821,\qquad
 \mathcal B(1/1000)<0.005565.
\]
These bound the early-label error in agreement; they are not enclosures
of $p_\infty$ itself.
\end{proposition}
\begin{proof}
In the scaled integral
$2\varepsilon\int_0^R y^2f(y)R_{\rm o}(R-y)\,dy$, use
\eqref{eq:parity-factorial} on $y\le R-j$ and
$R_{\rm o}\le1/2$ on the remaining part. The result is at most
\[
 \varepsilon\left(\frac{\mu_2}{m}+\frac{3\mu_2}{4j!}\right)
       +\varepsilon\EE[Y^2\ind_{\{Y>v\}}].
\]
For $v\ge8$, $y^2e^{-y/4}$ decreases on $[v,\infty)$, so the last
expectation is at most $Kv^2e^{-v/4}$. The empty-state probability
is at most $Ke^{-(R-1)/4}$, proving the inequality.

The saved certificate \texttt{certify\_cutoff\_budget.py} evaluates
the right side using rational arithmetic, $j=10$, $K<7$, and
$e^\gamma<891/500$. It certifies the last inequality from
$\gamma<H_{1024}-10\log2$, a rational lower series for $\log2$,
and a positive Taylor upper bound for the exponential. Negative
exponentials are bounded above by reciprocals of positive Taylor lower
sums. Thus no floating-point quadrature or statistical error estimate
enters these three bounds.
\end{proof}

There is also an exact state-dependent recursion for the expected query
budget. Let $b_0(t,x,z)=x\ind_{\{0<z<x\}}$.
For a split $(a,s,c)$ with $c=x-a-s$, put
$\alpha=\min(s,a/2)$ and $\beta=\min(s,c/2)$ and define
\[
 \mathcal S b=
 \begin{cases}
 b(s,a,\min(z,a-\alpha)),&z<a,\\
 b(s,a,a-\alpha)+b(s,s,s/2)+b(s,c,\beta),&a<z<a+s,\\
 b(s,c,\max(z-a-s,\beta)),&z>a+s.
 \end{cases}
\]
Set the endpoint values to zero and iterate the positive operator
\begin{equation}\label{eq:expected-query-recursion}
 b_{d+1}(t,x,z)=\int_0^{\min(t,x/3)}\int_s^{x-2s}
              \mathcal K_{t,x}(a,s)\,\mathcal S b_d\,da\,ds.
\end{equation}
Then $b_d=\EE Z_d$, $b_{d+1}\le b_d\le x\rho^d$, and
$b_d(t,x,z)=x b_d(t/x,1,z/x)$. These assertions follow from the query
identity, conditional independence of daughter trees, positivity, and
the original contraction lemma. In particular the first step is explicit:
\begin{equation}\label{eq:first-query-budget}
 b_1(t,x,z)=\int\mathcal K_{t,x}(a,s)
 \big[a\ind_{\{z<a\}}+x\ind_{\{a<z<a+s\}}
                  +c\ind_{\{z>a+s\}}\big]\,da\,ds,
 \quad 0<z<x.
\end{equation}

For the even limiting active configuration at $\varepsilon$, take the
breakpoints of its future collapse-map composition, as in the proof of
Theorem~\ref{thm:noise-convergence}. For each distinct breakpoint query
$j$, denote its initial gap length and local coordinate by $x_j,z_j$.
Define the quantity to be zero when that configuration is empty, and put
\begin{equation}\label{eq:adaptive-noise-budget}
 \mathcal Q_{\varepsilon,d}
 =4\EE\sum_j b_d(\varepsilon,x_j,z_j).
 \qquad
 |A_\infty-a_{\varepsilon,d}|
       \le\mathcal B(\varepsilon)+\mathcal Q_{\varepsilon,d}.
\end{equation}
The deterministic squared-mass error is at most four times the sum of
the query budgets. Taking its conditional expectation gives the bound;
multiple queries in one gap need not be independent. Also
$\mathcal Q_{\varepsilon,d}\le6\rho^d/\varepsilon$, so this retains the
previous certificate while allowing the actual states to improve it.
Computing these integrals with certified quadrature and a useful total
error is still required to enclose $p_\infty$ near its numerical candidate.

The same mass-query identity also gives a modest improvement of the
noise lower bound without evaluating the full approximation integral.

\begin{proposition}[A half-interval noise lower bound]\label{prop:half-noise}
Let $Z$ be the limiting mass in $[0,1/2]$ of a unit interval conditioned
on full shielding, and put $H=\EE[Z(1-Z)]$. Then
\begin{equation}\label{eq:half-disagreement}
 2H\ge\eta:=\frac{\log(8/5)}{12-\log(8/5)}.
\end{equation}
With $F(s)=\tanh(E_1(s)/2)$ and $c_*$ from \eqref{eq:cstar},
\begin{equation}\label{eq:cdagger}
 p_\infty\ge c_\dagger
 :=c_*+\eta\left[
     \int_0^\infty se^{-s}F(s)\,ds
        +\frac12\log\frac32-\frac1{12}\right]
 \approx0.1954170357.
\end{equation}
A stronger bound is obtained by putting
\[
 C_0=\frac1{24}\log\frac85,\qquad
 C_1=\frac1{48}+\frac1{24}\log2-\frac1{12}\log\frac53,
 \qquad
 \eta_+=\frac{\log(3/2)/6}{1-2C_0-2C_1/5}.
\]
Then $2H\ge\eta_+>\eta$, and
\begin{equation}\label{eq:cddagger}
 p_\infty\ge c_\ddagger
 :=c_*+\frac{\eta_+}{\eta}(c_\dagger-c_*)
 \approx0.2016982230.
\end{equation}
The exact expressions define the bounds; certified decimal enclosures
are given below. No new finite-$N$ rate is asserted here.
\end{proposition}
\begin{proof}
Define $Z$ as the $L^1$ limit of the midpoint query with uniform terminal
completions; Lemma~\ref{lem:query-contraction} makes the completion error
tend to zero. Reflection gives $\EE Z=1/2$. By shielding and scaling,
the normalized midpoint mass has this same law for every interval of
length $x$ with cutoff $x/3\le t\le x$.

For a unit interval, write its last split as $(a,s,c)$ with $a+s+c=1$.
Restrict to the event
\[
 E=\{s\le a,c\le2s\}.
\]
On $E$, both $a,c\le1/2$, so the midpoint lies between the last internal
pair. Moreover, $\alpha=a/2$, $\beta=c/2$, and all three daughter
intervals are fully shielded at cutoff $s$. Conditional on the split,
\eqref{eq:cdf-query} therefore reads
\[
 Z=aZ_1+s(1-Z_2)+cZ_3,
\]
where the $Z_i$ are independent copies of $Z$. Consequently
\[
 \EE[Z(1-Z)\mid a,s,c]
    =H(a^2+s^2+c^2)+\frac{as+ac+sc}{2}\qquad\hbox{on }E.
\]
Since $f(y)=1/(2y)$ for $1\le y\le3$, the split density on $E$ is
$1/(4asc)$. Set
\[
 C=\int_E\frac{as+ac+sc}{8asc}\,da\,ds,
 \qquad D=\int_E\frac{a^2+s^2+c^2}{4asc}\,da\,ds.
\]
The first integral can be evaluated directly:
\begin{align*}
 C={}&\frac18\int_{1/5}^{1/4}
       \left[2\log\frac{2s}{1-3s}+5-\frac1s\right]ds\\
    &+\frac18\int_{1/4}^{1/3}
       \left[2\log\frac{1-2s}{s}+\frac1s-3\right]ds
     =\frac1{24}\log\frac85.
\end{align*}
Here the respective $a$ ranges are $[1-3s,2s]$ and $[s,1-2s]$.
The inequality $a^2+s^2+c^2\ge as+ac+sc$ gives $D\ge2C$.
Discarding the nonnegative contribution from $E^c$ yields
$H\ge C+DH\ge C+2CH$, which proves \eqref{eq:half-disagreement}.

For the improvement retain also
$E_+\setminus E$, where $E_+=\{s\ge1/6,\ a,c\le1/2\}$
within the valid split region. Here the midpoint is still between the
last pair, and $a,c\le3s$, so the split density remains $1/(4asc)$.
Write the three independent nonnegative contributions to $Z$ as
$X_1,X_2,X_3$, bounded respectively by $a,s,c$, with
$X_2=s(1-Z_2)$. Its mean is $s/2$ and
$\EE[X_2(s-X_2)]=s^2H$.
Reflection gives $\EE X_1\ge a/2$ and $\EE X_3\le c/2$, since
their thresholds are respectively at least and at most the midpoint.
Writing $u=\EE X_1/a\ge1/2$ and $v=\EE X_3/c\le1/2$, their
cross term has expectation
$ac(u+v-2uv)=ac[1/2+2(u-1/2)(1/2-v)]\ge ac/2$.
The two cross terms involving $X_2$ have expectations $as/2$ and
$cs/2$, independently of the other two means. Keeping these terms
and the middle diagonal gives
\[
 \EE[Z(1-Z)\mid a,s,c]\ge s^2H+(as+ac+sc)/2.
\]
The contribution of the two middle cross terms over $E_+\setminus E$ is
\begin{align*}
 C_1
 &=\frac14\left[
   \int_{1/6}^{1/4}\log\frac1{1-2s}\,ds
   -\int_{1/5}^{1/4}\log\frac{2s}{1-3s}\,ds\right]\\
 &=\frac1{48}+\frac1{24}\log2-\frac1{12}\log\frac53>0.
\end{align*}
The respective symmetric $a$ ranges are $[1/2-s,1/2]$ and
$[1-3s,2s]$; above $s=1/4$ the two events agree.
The additional $ac/2$ term has integral
\[
 C_2=\frac18\left[\int_{1/6}^{1/4}1\,ds
                 -\int_{1/5}^{1/4}(5-1/s)\,ds\right]
     =\frac18\log\frac54-\frac1{48}.
\]
In particular $C_0+C_1+C_2=\log(3/2)/12$.
Pointwise the ratio of the coefficient of $H$ to the two middle
cross terms is $2s/(1-s)\ge2/5$. Combining this with
$C=C_0$ and $D\ge2C_0$ gives
\[
 H\ge\frac1{12}\log\frac32+(2C_0+2C_1/5)H,
\]
which proves $2H\ge\eta_+$.

To apply this interval law to the original fixed-count model, the last
pair and its arc trees must converge jointly. For $L_M=r\ge\delta>0$,
revealing finitely many generations in the two arcs gives the same
coefficient as \eqref{eq:tree-coefficient}, with the product over their
terminal intervals and reference measure $dr$ times the split coordinates.
The proof of Lemma~\ref{lem:finite-tree-limit} applies unchanged: the
number of terminal intervals is even, the empty and full subsets give
the limiting factor two, and the kernel denominators telescope.
The limiting marginal density is $2w_r(r)w_r(1-r)$; conditional on $r$,
the two arc trees are independent. Their total mass on $r\ge\delta$
equals the limiting last-pair probability, so Fatou and Scheffe also
remove the finite-sample empty-daughter cases. The query contraction
then passes to the midpoint masses. Finally let $\delta\downarrow0$,
using the scalar last-pair tail bound. This supplies the marked
convergence needed below, rather than inferring it from the scalar law.

The opposite-arc insertions already contribute at least
$c_*=\EE[L_\infty(1-L_\infty)]$ by Lemma~\ref{lem:last-comparison}.
Two insertions in the short arc always yield last-endpoint labels.
Conditional on its tree, their disagreement probability is $2Z(1-Z)$,
so their additional limiting contribution is $2H\EE L_\infty^2$.
When $L_\infty>1/3$, the long arc also has length less than
$2L_\infty$, and all its insertions yield the last endpoints.
Its normalized midpoint mass has the same fully shielded law, giving
the further contribution
$2H\EE[(1-L_\infty)^2\ind_{\{L_\infty>1/3\}}]$.
These insertion events are disjoint. The density
$1/(2r(1-r))$ on $[1/3,1/2]$ gives
\[
 \EE[(1-L_\infty)^2\ind_{\{L_\infty>1/3\}}]
   =\frac12\log\frac32-\frac1{12}.
\]
Use \eqref{eq:half-disagreement} and
$\EE L_\infty^2=\int_0^\infty se^{-s}F(s)\,ds$ to conclude.
Using $2H\ge\eta_+$ in the same inequality gives
\eqref{eq:cddagger}.
\end{proof}

\paragraph{Certified decimal versions of the exact noise bounds.}
The integer-interval calculation in
\texttt{certify\_noise\_constants.py} and
\texttt{noise-constant-certificate.json} certifies
\begin{equation}\label{eq:certified-noise-interval}
 0.2016982228\le p_\infty\le0.8112015673.
\end{equation}
It also gives the outward enclosures
\[
 \begin{array}{c@{\ \le\ }c@{\ \le\ }c}
 0.186818771115&c_*&0.186818771276\\
 0.195417035654&c_\dagger&0.195417035816\\
 0.201698222874&c_\ddagger&0.201698223035\\
 0.811201567283&u_*&0.811201567291.
 \end{array}
\]
These enclose the stated exact expressions; they do not certify the
candidate near $0.658$.

Here are the analytic error bounds used in the certificate.
Put $G(s)=e^{-2s}F(s)^2$ and $J(s)=se^{-s}F(s)$.
Then $c_*=1/4-\frac12\int_0^\infty G(s)\,ds$.
The proved bound $\EE e^{Y/4}<7$, together with
$e^{1/2}<5/3$ and $e^{1/4}<4/3$, gives
\[
 \sup_{s\ge0}|G^{(4)}(s)|<503808,
 \qquad \sup_{s\ge0}|J^{(4)}(s)|<17920.
\]
Indeed $G(s)=\EE e^{-sT}$ for $T=2+Y_1+Y_2$, with independent
copies $Y_1,Y_2$ of $Y$, and
$\EE T^4\le4!4^4\EE e^{T/4}<82\cdot4!4^4$.
For $R=1+Y$, differentiating $se^{-sR}$ gives
$|J^{(4)}(s)|\le5\EE R^3<17920$, using $u e^{-u}\le1$.
Composite Simpson integration on $[0,16]$ with mesh $2^{-12}$
therefore has errors at most $16\cdot503808/(180\cdot2^{48})$
and $16\cdot17920/(180\cdot2^{48})$.
The tails are at most $e^{-64}/4$ and $33e^{-32}/4$, since $Y\ge1$.

Every endpoint value is enclosed using integers with denominator
$10^{50}$ and outward rounding at every operation. The exponential
integral uses
$E_1(s)=-\gamma-\log s-\sum_{j\ge1}(-s)^j/(j\,j!)$,
with its alternating remainder after $128$ terms;
$F(0)=1$ is assigned exactly. Logarithms use the positive
atanh series after binary range reduction, and positive exponential
series use geometric remainder bounds. Finally, put
$H_n=\sum_{j=1}^n1/j$. For $n=2^{18}$,
\[
 H_n-\log n-\frac1{2n}<\gamma
       <H_n-\log n-\frac1{2(n+1)}.
\]
For these elementary bounds, write
$H_n-\log n-\gamma=\sum_{k=n}^\infty
\int_k^{k+1}(k+1-x)/(x(k+1))\,dx$ and bound each integrand using
$k\le x\le k+1$; the lower sum is at least its decreasing integral.
No floating-point quadrature is used in the certificate.
The separate $431$ comparisons with $100$-digit standard-library
decimal logarithms and exponentials validate the implementation,
not the analytic bounds.

\subsection{Several-point perturbations and the full noise profile}
\label{sec:perturbations}

\status{Deterministic coupling and finite profile identities proved}
The user's broader question concerns the entire configuration, its
matching history, and perturbations of several points. These are
different observables, even under the same coupling.

For finite location sets $A,B\subset\T$, let $C_r(A)$ be the active set
after separation cutoff $r$, and let $M_r(A)$ be the set of matching
edges accepted by then. Scan the pair distances of $A\cup B$ in a
common increasing order, accepting an edge in each process exactly
when both endpoints are still active there. Distances are distinct
in the random model; a common deterministic ordering of ties suffices
for the statements below. Put
\[
 d=|A\mathbin\triangle B|,\qquad
 D_r=|C_r(A)\mathbin\triangle C_r(B)|.
\]
A point active in only one process will be called a discrepancy.

\begin{theorem}[Contraction under arbitrary finite edits]
\label{thm:edit-contraction}
For $0\le r\le s$,
\begin{equation}\label{eq:edit-contraction}
 D_s\le D_r\le d,\qquad D_r\equiv d\pmod2.
\end{equation}
Every effective change of $D_r$ is a decrease by two. If the active
sets ever agree, their entire future paths agree. In particular,
resampling $k$ starting locations gives $D_r\le2k$ simultaneously
at all cutoffs, for every realization.
\end{theorem}
\begin{proof}
An edge accepted in both processes removes two common points and
leaves the discrepancy set unchanged. If an edge is accepted in
only one process, at least one endpoint is a discrepancy. With
exactly one discrepant endpoint, the other endpoint is common:
the discrepancy transfers to it and changes which process contains
it. With two discrepant endpoints, both belong to the same process,
and both disappear. These exhaust the possibilities. No new
discrepancy is born, proving the assertions.
\end{proof}

This proof uses only a shared ordered list of pair events. In
particular it holds for greedy matching in any finite metric space.
It strengthens the bound obtained by chaining add-one couplings:
the discrepancy count itself decreases with time.

\begin{corollary}[Closeness of active measures]\label{cor:edit-measures}
Write $\Xi_A(r)=\sum_{v\in C_r(A)}\delta_v$ and define $\Xi_B(r)$ similarly.
For every bounded real $\phi$,
\[
 |\langle\Xi_A(r)-\Xi_B(r),\phi\rangle|
       \le D_r\|\phi\|_\infty.
\]
If the two active counts are $m,n>0$, and their common active count
is $c$, their normalized measures satisfy
\begin{equation}\label{eq:edit-TV}
 d_{\rm TV}(\widehat\Xi_A(r),\widehat\Xi_B(r))
 =1-\frac{c}{\max(m,n)}
 =\frac{\max(m-c,n-c)}{\max(m,n)}
 \le\frac{D_r}{\max(m,n)}.
\end{equation}
Here probability total variation means $\sup_E|\mu(E)-\nu(E)|$.
For intrinsic circular distance, their Wasserstein distance $W_1$
is at most one half of the left side of \eqref{eq:edit-TV}.
\end{corollary}
\begin{proof}
The signed counting measure has total variation norm $D_r$.
Each common point carries mass $\min(1/m,1/n)=1/\max(m,n)$ in
the common part of the two probability measures. Couple that
common part identically, and transport the remaining mass by at
most the circle's diameter $1/2$.
\end{proof}

\begin{corollary}[Robustness under a growing number of edits]
\label{cor:edit-robustness}
Let $A_N$ be an iid uniform sample, $r_N\to0$, and let $B_N$ be any
possibly data-dependent modification with
\[
 |A_N\mathbin\triangle B_N|=o_{\PP}\bigl(Nq(Nr_N)\bigr),
\]
where $q$ is the function in \eqref{eq:early-fluid}. Then
\[
 \frac{\Xi_{B_N}(r_N)}{Nq(Nr_N)}\longrightarrow dx,\qquad
 \widehat\Xi_{B_N}(r_N)\longrightarrow dx
 \quad\hbox{weakly in probability}.
\]
The edited count is positive with probability tending to one.
Thus $o(N)$ edits are harmless on the microscopic clock, and
$o(1/r_N)$ edits are harmless throughout the intermediate regime
$r_N\to0$, $Nr_N\to\infty$.
\end{corollary}
\begin{proof}
Corollary~\ref{cor:all-vanishing-cutoffs} gives
$K_{A_N}(r_N)/(Nq(Nr_N))\to1$. The signed test-integral and count
differences are bounded by the edit budget. Divide by $Nq(Nr_N)$
and then by the edited count. The two stated scales follow from
positivity of $q$ on compact intervals and
$q(s)\sim(2e^\gamma s)^{-1}$.
\end{proof}

More generally, for any deterministic normalization $a_N$,
$a_N|A_N\triangle B_N|\to0$ in probability makes their centered
counting fields asymptotically equal against bounded tests, with
the same centering. In particular, edits $o_{\PP}(r_N^{-1/2})$ are
invisible on the $\sqrt{r_N}$ fluctuation scale whenever a baseline
fluctuation theorem is available. This transfer principle does not
by itself establish the still-open simultaneous finite-count CLT.

\begin{proposition}[The changed matching consists of paths]
\label{prop:edit-forest}
The graph formed by $M_r(A)\triangle M_r(B)$, together with isolated
vertices from $A\triangle B$, is a forest of alternating paths.
Each path either has one initial discrepancy and one currently
active discrepancy, or has two initial discrepancies and no active
discrepancy. An untouched initial discrepancy is the isolated
one-vertex version of the first case. Hence there are at most $d$
paths. If $J_r$ counts paths of the second kind, then
\[
 D_r=d-2J_r.
\]
Along each growing end, successive edge distances increase.
There is no bound on the number of changed edges in terms of $d$
alone.
\end{proposition}
\begin{proof}
Initially there are $d$ isolated discrepancy vertices. A common
accepted edge appears in both matchings and contributes nothing
to their difference. A transfer extends one growing path by an
edge to a previously common active vertex, which has no earlier
matching edge in either process. The path's active end changes
side, so edge colors alternate. An event between two discrepancies
joins two growing paths and leaves a path with two initial ends
and no active end. It cannot join the ends of a single growing
path, since that path has only one active end. Thus no cycle or
branch occurs, and the stated component count is preserved.

For the absence of a history bound, fix $m\ge1$ and put
\[
 x_i=\frac{i(i+1)}{8m(2m+1)},\qquad 0\le i\le2m.
\]
All points lie in $[0,1/4]$ and their consecutive gaps strictly
increase. On $A=\{x_0,\ldots,x_{2m}\}$ the matching pairs are
$(x_0,x_1),(x_2,x_3),\ldots$, leaving $x_{2m}=1/4$.
Replace $x_0$ by $3/4$ to obtain $B$. Its pairs are
$(x_1,x_2),(x_3,x_4),\ldots$, leaving $3/4$.
Every internal consecutive gap is less than $1/4$, whereas
$3/4$ is at distance at least $1/4$ from every other point of $B$.
The claimed matchings therefore follow in increasing-gap order.
Their $2m$ edges are all different, despite a single resampling
and $D_r\le2$ at every time. Their final positions are antipodal.
\end{proof}

For two odd samples of equal size, let
$T=\inf\{r:C_r(A)=C_r(B)\}$, allowing $T=\infty$. Equality is
absorbing, and both processes have reached their singleton states
by $r=1/3$. Consequently
\begin{equation}\label{eq:coalescence-noise}
 \PP(\hbox{different final positions})
    =\PP(T=\infty)=\tfrac12\EE D_{1/3}.
\end{equation}
For one resampled point, $D_r\in\{0,2\}$, so
$\EE D_r=2\PP(T>r)$. For several resampled points this last
identity is replaced by the decreasing discrepancy profile
$\EE D_r$, which measures more than final-label disagreement.

\paragraph{The resampling profile.}
For odd $N$, let $p_N(k)$ be the survivor-label disagreement
probability when a uniformly chosen set of $k$ coordinates is
resampled independently, keeping their labels. Exchangeability
makes this equal to the probability for any fixed set of $k$ labels.
Thus $p_N(1)=p_N$ and $p_N(0)=0$. Let $p_N^{\rm pos}(k)$ be its
position counterpart. Almost surely, distinct input labels have
distinct locations across the two samples except for the shared
coordinates. Therefore
\begin{equation}\label{eq:k-label-position}
 0\le p_N^{\rm pos}(k)-p_N(k)\le k/N.
\end{equation}
Indeed the difference event requires the same resampled label to
win both runs, and the first winner belongs to the resampled set
with probability $k/N$.

\begin{proposition}[Coupled limits share their Gaussian field]
\label{prop:coupled-entrance}
Fix $k$. The unique joint entrance limit established in
Theorem~\ref{thm:joint-perturbed-limit} has at most $2k$ discrepant
locations at every positive cutoff.

In every such coupling, put $c=1/(2e^\gamma)$ and
$G_r^i=\sqrt r(\Xi^i(r)-c\,dx/r)$ for $i=A,B$. As $r\downarrow0$,
\begin{equation}\label{eq:coupled-gaussian}
 (G_r^A,G_r^B)\Longrightarrow(G,G)
 \quad\hbox{in }H^{-a}(\T)\times H^{-a}(\T),\qquad a>1/2,
\end{equation}
where $G$ is the white noise in Theorem~\ref{thm:entrance-fluctuations}.
The two normalized empirical-measure fluctuation fields likewise
converge jointly to the same Gaussian bridge.
\end{proposition}
\begin{proof}
The deterministic discrepancy bound passes to the marked joint limit.
On that limit it holds at each dyadic cutoff and, by contraction,
simultaneously at all later cutoffs.

For $a>1/2$, the norm $C_a=\|\delta_x\|_{H^{-a}}$ is finite and
independent of $x$, as follows directly from the summable Fourier
weights. Thus
\[
 \|G_r^A-G_r^B\|_{H^{-a}}\le2kC_a\sqrt r\longrightarrow0.
\]
The marginal Gaussian theorem gives \eqref{eq:coupled-gaussian}.
For normalized measures, \eqref{eq:edit-TV} gives
\[
 r^{-1/2}\|\widehat\Xi^A(r)-\widehat\Xi^B(r)\|_{H^{-a}}
 \le\frac{4kC_a\sqrt r}{r\max(K^A(r),K^B(r))}
 \longrightarrow0
\]
in probability, because $rK^i(r)\to c$. Apply the normalized
marginal theorem. These are statements about the entrance coupling
followed by $r\downarrow0$, not a simultaneous
finite-$N$ fluctuation theorem.
\end{proof}

\begin{theorem}[A finite spectral formula for all resampling sizes]
\label{thm:noise-profile}
There are nonnegative weights $w_{N,j}$, $0\le j\le N$, with
\[
 \sum_{j=0}^N w_{N,j}=1,\qquad w_{N,0}=1/N,\qquad w_{N,1}=0,
\]
such that
\begin{equation}\label{eq:k-noise-spectrum}
 1-p_N(k)=\sum_{j=0}^N
       w_{N,j}\frac{\binom{N-k}{j}}{\binom Nj},
 \qquad 0\le k\le N.
\end{equation}
A numerator with $j>N-k$ is zero. The profile $p_N(k)$ is
increasing and discretely concave, and
$p_N(N-1)=p_N(N)=1-1/N$.

For $\nu_N=\sum_jw_{N,j}\delta_{j/N}$ and $1\le k\le N$,
\begin{align}
 \int x\,\nu_N(dx)&=p_N(1),\label{eq:spectral-mean}\\
 0\le p_N(k)-\int[1-(1-x)^k]\,\nu_N(dx)
       &\le\frac{k(k-1)}{2N}.\label{eq:profile-replacement}
\end{align}
In particular
\[
 p_N(1)\le p_N(k)
   \le1-(1-p_N(1))^k+\frac{k(k-1)}{2N}.
\]
\end{theorem}
\begin{proof}
Represent the survivor label by the unit vector
$F_N(X)=e_{\mathrm{winner}(X)}\in\mathbb R^N$. On the product space
of the independent coordinates, define its orthogonal components
by
\[
 F_{N,S}=\sum_{T\subseteq S}(-1)^{|S|-|T|}
                        \EE[F_N\mid X_T].
\]
Each component depends only on $X_S$ and has conditional mean
zero on integrating out any coordinate of $S$. These properties
give $F_N=\sum_S F_{N,S}$ and orthogonality of distinct components:
integrate a coordinate belonging to one index set but not the other.
Set $w_{N,j}=\sum_{|S|=j}\EE\|F_{N,S}\|^2$. Their sum is
$\EE\|F_N\|^2=1$. Uniformity of the survivor label gives $w_{N,0}=1/N$.
Conditional on the position of any one labelled point, rotation
invariance makes the distribution of every survivor label independent
of that position. Thus $\EE[F_N\mid X_i]=\EE F_N$ and $w_{N,1}=0$.

If coordinates in $R$ are resampled, conditional independence gives
agreement probability
$\EE\|\EE[F_N\mid X_{R^c}]\|^2
 =\sum_{S:S\cap R=\varnothing}\EE\|F_{N,S}\|^2$.
Averaging over $|R|=k$ gives \eqref{eq:k-noise-spectrum}.
For $a_j(k)=\binom{N-k}{j}/\binom Nj$, Pascal's identity gives
\[
 (-1)^\ell\Delta^\ell a_j(k)
    =\frac{\binom{N-k-\ell}{j-\ell}}{\binom Nj}\ge0,
 \quad k+\ell\le N,
\]
with zero on the right if $j<\ell$.
This proves monotonicity and concavity, and the endpoint identity
uses $w_{N,1}=0$.

Formula \eqref{eq:spectral-mean} is the case $k=1$.
For a set of $j$ marked coordinates, sampling $k$ coordinates
without replacement hits it at least as often as sampling with
replacement. Couple the two procedures until the first repeated
draw; the probability of a repeat is at most $\binom{k}{2}/N$.
This proves \eqref{eq:profile-replacement}. Finally apply Jensen's
inequality to the concave function $1-(1-x)^k$.
\end{proof}

\noindent
For example, the proved $p_3(1)=1/2$ forces
$(w_{3,0},w_{3,1},w_{3,2},w_{3,3})=(1/3,0,1/2,1/6)$.
Thus $p_3(2)=p_3(3)=2/3$.

\begin{corollary}[Subsequential profiles and rare independent noise]
\label{cor:spectral-limits}
Every subsequence of odd $N$ has a further subsequence on which
$\nu_N$ converges weakly to a probability measure $\nu$ on $[0,1]$.
On that subsequence, simultaneously for every fixed integer $k\ge1$,
\begin{equation}\label{eq:limiting-noise-profile}
 p_N(k)\longrightarrow p(k)=1-\int(1-x)^k\,\nu(dx),
 \qquad \int x\,\nu(dx)=p_\infty.
\end{equation}
The same limits hold for position disagreement. Existence of the
limits $p_N(k)$ for every fixed $k$ is equivalent to weak convergence
of the full sequence $\nu_N$.

Independently resample each coordinate with probability $\lambda/N$,
where $\lambda\ge0$ is fixed and $N\ge\lambda$, and denote the
disagreement probability by $g_N(\lambda)$. Then on the same subsequence
\begin{equation}\label{eq:rare-noise-profile}
 g_N(\lambda)\longrightarrow
       g(\lambda)=1-\int e^{-\lambda x}\,\nu(dx).
\end{equation}
Every such limit obeys
\[
 p_\infty(1-e^{-\lambda})
       \le g(\lambda)\le1-e^{-\lambda p_\infty},
 \qquad
 p_\infty\le p(k)\le1-(1-p_\infty)^k.
\]
\end{corollary}
\begin{proof}
Compactness of $[0,1]$, \eqref{eq:profile-replacement}, and convergence
of $p_N(1)$ give \eqref{eq:limiting-noise-profile}.
Equation \eqref{eq:k-label-position} identifies the position limits.
Conversely, limits for all $k$ specify integrals of every polynomial
in $1-x$; uniform polynomial approximation determines a unique
probability measure on $[0,1]$. Compactness then gives full convergence.

Independent-coordinate resampling leaves a component of order $j$
unchanged with probability $(1-\lambda/N)^j$. This converges uniformly
over $0\le j\le N$ to $e^{-\lambda j/N}$, giving
\eqref{eq:rare-noise-profile}. Concavity and the chord bound for
$1-e^{-\lambda x}$ on $[0,1]$ give the displayed inequalities.
\end{proof}

Thus the onset scale for independent-coordinate noise is already
identified: if the resampling probability $\epsilon_N$ satisfies
$N\epsilon_N\to0$, disagreement tends to zero because the probability
of any resampling does. At $\epsilon_N=\lambda/N$ with fixed
$\lambda>0$, the preceding bounds keep every subsequential limit
strictly between zero and one. Section~\ref{sec:coupled-limit}
proves uniqueness of that profile. Theorem~\ref{thm:no-zero-atom}
proves that disagreement tends to one above this scale.

\paragraph{Concrete questions opened by these results.}
The following remain conjectures or open tasks, not consequences of
the one-coordinate limit.
\begin{enumerate}
\item \textbf{Coupled evolution rule.} Section~\ref{sec:history-markov}
gives a reversible insertion rule on complete entrance histories.
Find a smaller sufficient state and useful quantitative estimates
for its discrepancy motion and coalescence.
\item \textbf{The limiting noise spectrum.} Characterize $\nu$ and
quantify its mass near zero. Theorem~\ref{thm:no-zero-atom} now proves
$\nu(\{0\})=0$ and $p_N(k_N)\to1$ whenever $k_N\to\infty$.
Its proof uses reversible insertion and endpoint regularity;
$w_{N,0}\to0$ alone would not exclude spectral mass accumulating
from orders $j=o(N)$. Theorem~\ref{thm:logarithmic-spectrum} below
gives a logarithmic decay rate for $1-p(k)$.
\item \textbf{Changed matching edges.}
Section~\ref{sec:growing-perturbations} proves expected growth of
order $k(1+\log(N/k))$, uniformly for $1\le k\le N$.
Section~\ref{sec:sublinear-history} identifies the coefficient for
every $k=o(N)$. Sections~\ref{sec:dense-history}--\ref{sec:history-concentration}
prove a dense-edit profile and relative concentration for growing
counts. Section~\ref{sec:fixed-history-LLN} includes fixed nonzero
counts. Evaluate the profile more explicitly and determine centered
fluctuations.
\item \textbf{Spatial response.} Determine the joint law of the
discrepancy locations, their signed action on smooth test functions,
and the loss of overlap as a function of cutoff and perturbation size.
The deterministic edit bounds give sufficient stability scales;
they do not establish that those scales are sharp for random edits.
\end{enumerate}

\subsection{Unique coupled limits from separated insertion gaps}
\label{sec:coupled-limit}

\status{Working proof: coupled entrance laws and the full noise spectrum}
The missing input in the preceding subsection is the joint law, including
which active points are shared. It can be recovered by inserting all
perturbed coordinates into one background and first stopping at a small
positive cutoff. Only insertions in the same or neighbouring surviving
gaps can interact before that cutoff.

For a finite background $B_M$ and a nonempty active configuration
$C_M(\varepsilon)$, let $I_i$ be its cyclic gaps, of lengths $x_i$.
Compose the background collapse maps up to cutoff $\varepsilon$.
Their pushforward of Lebesgue measure on $I_i$, in local coordinates,
is denoted by $\mu_{M,i}$; it has mass $x_i$. These are cutoff measures,
not the full-process basin weights.

\begin{lemma}[Convergence of the complete cutoff gap measures]
\label{lem:gap-measure-limit}
Along either fixed parity of $M$, the active configuration together with
its finite collection of measures $\mu_{M,i}$ converges jointly weakly.
The empty configuration is a separate state. Conditional on a nonempty
limiting configuration, the gap measures are independent functionals of
the limiting gap trees.

Let $\mu_i^{(d)}$ be obtained by completing depth-$d$ terminal intervals
with Lebesgue measure. For equal-mass measures, use $W_1$ to mean total
transport cost, without normalizing their masses. Then the limiting
measure $\mu_i$ satisfies
\begin{equation}\label{eq:whole-gap-approximation}
 \EE[W_1(\mu_i,\mu_i^{(d)})\mid C(\varepsilon)]
       \le x_i^2\rho^d.
\end{equation}
Almost surely it gives no mass to either readout boundary
$\min(\varepsilon,x_i/2)$ or $\max(x_i-\varepsilon,x_i/2)$.
\end{lemma}
\begin{proof}
For two completions below depth $d$, the cumulative-mass difference at
$z$ is bounded by $Z_d(\varepsilon,x_i,z)$. On a line, the transport
cost of equal-mass measures is the integral of their absolute cumulative
difference. Integrating the contraction bound over $0<z<x_i$ gives
$x_i^2\rho^d$, uniformly in the deeper completion. The same bound between
successive depths is summable. Completeness of the space of measures of
mass $x_i$ on $[0,x_i]$ therefore constructs $\mu_i$, almost surely and
in expected transport cost, and proves \eqref{eq:whole-gap-approximation}.

At each fixed depth, the finite background geometry converges in total
variation by Lemma~\ref{lem:finite-tree-limit} and its odd-parity version
in Theorem~\ref{thm:odd-entrance}. Couple the geometries identically outside
that error event. Integrating their deterministic cumulative budgets
then gives
\[
 \limsup_{M\to\infty}\EE\sum_i
 W_1(\mu_{M,i},\mu_{M,i}^{(d)})\le\rho^d,
\]
since $\sum_i x_i^2\le1$; early-empty daughter states have vanishing
probability. The converging-together argument proves the joint weak
limit. Independence follows from that of the limiting gap trees.

For either displayed boundary $h=h(x_i)$, every finite-depth
approximation gives $\{h\}$ zero mass almost surely: its atoms are
internal tree vertices with continuous split coordinates, and its
remaining part is diffuse. Apply the cumulative budget at $h-\delta$
and $h+\delta$. For deeper completions their expected mass in this
window is at most the depth-$d$ window mass plus $2x_i\rho^d$.
Weak convergence first bounds the limiting mass of the open window
by this quantity. Send $\delta\downarrow0$, then $d\to\infty$.
This proves the boundary assertion, conditionally on the active geometry.
\end{proof}

For a nonempty finite active set $C$ whose pair separations exceed
$\varepsilon$, define the signed one-point response
\begin{equation}\label{eq:cutoff-response}
 \eta_\varepsilon(C,y)=
 \begin{cases}
 -\delta_{v(y)},&d(y,C)\le\varepsilon,\\
 \delta_y,&d(y,C)>\varepsilon,
 \end{cases}
\end{equation}
where $v(y)$ is the nearest point of $C$. Ties are null in the laws
under consideration. In a gap of length $x$, the two decision
boundaries are exactly those in Lemma~\ref{lem:gap-measure-limit}.

\begin{lemma}[Separated insertions and their error probability]
\label{lem:separated-insertions}
Insert $q$ independent uniform probes into a common iid background,
with $q$ fixed. Call their gaps separated at cutoff $\varepsilon$ if
the background is nonempty and no two probes lie in the same or
neighbouring cyclic gaps. Let $Y_j$ be the image of probe $j$ under
the background maps before the cutoff. On this event, for every
subset $S$ of probes,
\begin{equation}\label{eq:isolated-readout}
 \Xi_{B_M+\{z_j:j\in S\}}(\varepsilon)
   =\Xi_{B_M}(\varepsilon)
      +\sum_{j\in S}\eta_\varepsilon(C_M(\varepsilon),Y_j).
\end{equation}
Conditional on the background, a probe chooses gap $i$ with probability
$x_i$, and its image has distribution $\mu_{M,i}/x_i$ in that gap.
The failure probability is at most
\begin{equation}\label{eq:gap-isolation-error}
 \beta_{M,q}(\varepsilon):=
 \PP(C_M(\varepsilon)=\varnothing)
 +3\binom q2\EE\!\left[
   \ind_{\{C_M(\varepsilon)\ne\varnothing\}}\sum_i x_i^2\right].
\end{equation}
For either fixed background parity,
$\beta_{M,q}(\varepsilon)\to\beta_q(\varepsilon)=O_q(\varepsilon)$
as first $M\to\infty$, then $\varepsilon\downarrow0$.
\end{lemma}
\begin{proof}
The closest-pair map has a stronger cutoff interpretation than its
survivor readout: after the background pair at distance $s$ is removed,
the augmented process agrees, at every cutoff at least $s$, with the
reduced background and the mapped insertion. If the insertion meets
one endpoint first, the other endpoint becomes its replacement.
Outside the two cells, removing the common pair leaves the same
discrepancy; a previous collision with another background point
remains frozen until that point's background removal. This is the
event-by-event argument of Proposition~\ref{prop:map}, and also applies
when the remaining background is a singleton.

Iterate through the background events before $\varepsilon$. No pair
in $C_M(\varepsilon)$ has distance at most $\varepsilon$, so the
remaining insertion either survives as $Y_j$ or removes its nearest
active endpoint. This gives \eqref{eq:cutoff-response} for one probe.
Its discrepancy remains in the closure of its original active gap:
it cannot cross a surviving endpoint, and after removing such an
endpoint the missing-point discrepancy is frozen through the cutoff.

For probes in separated gaps, these closed regions are disjoint.
Two positive discrepancies in different regions cannot meet first
across an intervening active endpoint. Negative discrepancies are
frozen, and cannot cause an interaction through the cutoff. Induction
over the common event list consequently superposes the single-probe
histories, proving \eqref{eq:isolated-readout}, simultaneously for all
subsets. The conditional sampling description is the definition of
pushforward.

For two uniform probes, the conditional probability of equal or
neighbouring gaps is at most
$\sum_i x_i(x_{i-1}+x_i+x_{i+1})\le3\sum_i x_i^2$.
This remains an upper bound with one or two gaps, where neighbours
can be counted twice. A union bound proves
\eqref{eq:gap-isolation-error}. The bounded gap functional converges
by the known active-configuration TV theorem. For even parity its
limiting expectation is \eqref{eq:gap-square-identity}. For odd parity it is
\[
 2h_\varepsilon(1)+2\int_0^1x^2h_\varepsilon(x)
                   R_{{\rm e},\varepsilon}(1-x)\,dx.
\]
The first term is the singleton sector. The exponential bound on $f$,
$R_{\rm o},R_{\rm e}\le1/2$, and $\EE Y^2<\infty$ make both gap
expectations $O(\varepsilon)$. The even empty probability is
exponentially small; the odd one is zero.
\end{proof}

\begin{theorem}[A unique joint limit with shared-point identities]
\label{thm:joint-perturbed-limit}
For every fixed $k$, the original and $k$-resampled odd processes
converge jointly after every positive cutoff, retaining the identities
of shared background points. The limit is unique. This includes
configurations at finitely many positive cutoffs and their full finite
future collision histories. This theorem asserts weak convergence.
The separate quantitative coupling in Lemma~\ref{lem:finite-family-rate}
below strengthens the positive-cutoff spatial insertion-family
comparison to TV.
Distinct limiting background identities have distinct locations almost
surely. In particular the final agreement probability converges:
\begin{equation}\label{eq:all-fixed-noise-limits}
 p_N(k)\longrightarrow p(k),\qquad
 p_N^{\rm pos}(k)\longrightarrow p(k).
\end{equation}
Joint convergence with shared identities also holds for any fixed finite family obtained by adding
prescribed subsets of finitely many independent uniform probes to one
common background of fixed parity.
\end{theorem}
\begin{proof}
For two copies take a background of size $M=N-k$ and $q=2k$
independent probes, partitioned into the two insertion sets. At cutoff
$\varepsilon$, mark every active location by the subset of copies in
which that point identity is present. Distinct probe locations are
separate point identities even when their coordinate labels coincide.
For a finite family use
nonempty subsets of its index set as marks. For now the compact state
space allows atoms with disjoint marks to coincide in location.
Each individual configuration has uniformly bounded count and
separation at least $\varepsilon$, giving tightness.

Define an auxiliary marked law $Q_{M,\varepsilon}$ by the sampling
description and response in Lemma~\ref{lem:separated-insertions},
putting a cemetery state on the failure event. Each positive response
has the mark of the copies using that probe; each negative response
removes those copies from the endpoint's mark. By the deterministic
lemma this state equals the true state on the separated event. Hence
their probability TV distance is at most $\beta_{M,q}(\varepsilon)$.
Lemma~\ref{lem:gap-measure-limit}, its zero-boundary-mass assertion,
and conditional sampling give a unique weak limit
$Q_\varepsilon$ for $Q_{M,\varepsilon}$.

Consider any subsequential limit $P_\varepsilon$ of the true marked
state. Lower semicontinuity of probability TV under joint weak
convergence gives
\begin{equation}\label{eq:coupled-limit-approximation}
 d_{\rm TV}(P_\varepsilon,Q_\varepsilon)
       \le\beta_q(\varepsilon).
\end{equation}
One can see this directly by taking the supremum of differences of
integrals of continuous functions with values in $[0,1]$; on this
compact metric space that supremum equals probability TV.

Fix target cutoffs bounded below by $r_0>\varepsilon$.
Forward evolution, retaining marks, is continuous
$P_\varepsilon$-almost surely. Each marginal is the known entrance
law, so ties between its collision distances and collisions at a
specified cutoff have probability zero. Coincidences between different
copies do not affect their separate greedy evolutions. Thus every
target subsequential law is a pushforward of some $P_\varepsilon$.
By \eqref{eq:coupled-limit-approximation} it lies within
$\beta_q(\varepsilon)$ in TV of the corresponding pushforward of
$Q_\varepsilon$, whether or not that auxiliary law has exceptional
ties. Any two target limits therefore have TV distance at most
$2\beta_q(\varepsilon)$. Send $\varepsilon\downarrow0$ to prove
uniqueness and hence convergence of the full sequence. The argument
also applies to the finite future event records of each copy.

On the separated event, the union of all distinct surviving identities
at cutoff $\varepsilon$ has pair distances greater than $\varepsilon$:
background points were already separated, a positive response is
farther than $\varepsilon$ from them, and different occupied gaps are
nonadjacent. Future evolution only deletes locations. Consequently,
in any limiting state after $r_0$, the probability that two different
identities have distance less than $\varepsilon/2$ is at most
$\beta_q(\varepsilon)$. Letting $\varepsilon\downarrow0$ rules out
accidental coincidences.

At cutoff $1/3$ each odd copy is a singleton. Agreement of their
background identity is the event that the common-mark count is one,
a continuous discrete-count readout on the marked space. A resampled
coordinate can cause label agreement without position agreement with
probability at most $k/N$, by \eqref{eq:k-label-position}. This proves
\eqref{eq:all-fixed-noise-limits}. If identities of resampled coordinates
are also requested at an earlier positive cutoff, the probability any
of them is still active is at most
$2k\max(1,\lfloor1/\varepsilon\rfloor)/N$ by exchangeability, so it
vanishes. The finite-family argument is identical, using its fixed
total number $q$ of probes and the appropriate subset marks.
\end{proof}

\begin{corollary}[A limiting spectrum and a stationary noise evolution]
\label{cor:noise-evolution}
The full sequence $\nu_N$ converges weakly to a probability measure
$\nu$ on $[0,1]$. Its fixed-$k$ and rare-noise profiles are
\[
 p(k)=1-\int(1-x)^k\,\nu(dx),\qquad
 g(\lambda)=1-\int e^{-\lambda x}\,\nu(dx),\qquad
 \int x\,\nu(dx)=p_\infty.
\]
The measure $\nu$ is unique. Theorem~\ref{thm:no-zero-atom} proves
$\nu(\{0\})=0$; its further identification remains open.

There is also a stationary limiting process under successive random
resamplings. Give each of the $N$ coordinates a refresh clock of rate
$1/N$, keeping the balloon evolution as a function of the current
input. On every bounded refresh-time interval, the coupled active
configurations at any fixed positive balloon cutoff, and their future
histories, have a unique joint limit. The final-location process
$U(\lambda)$ has uniform marginals and at most a rate-one Poisson
number of candidate jumps. It satisfies
\begin{equation}\label{eq:noise-time-survivor}
 \PP(U(0)\ne U(\lambda))=g(\lambda),\qquad
 \EE[\hbox{actual jumps in }[0,T]]=Tp_\infty.
\end{equation}
Section~\ref{sec:history-markov} identifies a reversible Markov rule
on complete entrance histories. No Markov property for the observed
fixed-cutoff configuration or survivor alone is asserted.
\end{corollary}
\begin{proof}
The fixed-$k$ limits and Corollary~\ref{cor:spectral-limits} imply full
spectral convergence and both formulas. Equivalently, their polynomial
moments determine $\nu$ uniquely.

The superposed refresh clock has rate one for every $N$. On $[0,T]$
let its count be $J\sim\operatorname{Poisson}(T)$, independent of the
locations. Conditional on $J=j$, a repeated coordinate has probability
at most $j(j-1)/(2N)$. If none repeats, the $j+1$ successive inputs
are a fixed family formed from one background and $2j$ independent
probes. Apply Theorem~\ref{thm:joint-perturbed-limit}, then truncate
the Poisson tail. The candidate times have the same continuous law
for every $N$ and are independent of the input family. This proves
convergence of the corresponding finite jump records and constructs
the limiting refresh-time paths. Consistency gives a process on all
bounded intervals, and stationarity follows from finite-$N$ stationarity.

The end-to-end disagreement is the Poisson mixture of $p(j)$:
$\sum_j e^{-\lambda}\lambda^j p(j)/j!
=1-\int e^{-\lambda x}\nu(dx)$. For each finite $N$, the mean number
of actual final-position changes on $[0,T]$ is
$T p_N^{\rm pos}(1)$. The marked finite-family limit determines which
candidate jumps are actual changes, including equal consecutive
states; different limiting identities do not accidentally coincide.
The Poisson candidate count supplies uniform integrability. Passing
to the limit gives the second assertion in
\eqref{eq:noise-time-survivor}.
\end{proof}

\subsection{Logarithmic growth of changed matching histories}
\label{sec:history-growth}

\status{Working proof: logarithmic order for every fixed resampling size}
Let $\mathcal L_n$ be the sum of all matching separations for $n$ iid
uniform points. For odd $N$, let $H_N(k)$ count edges in the symmetric
difference of the completed matchings before and after $k$ coordinate
resamplings, treating fresh locations as new vertices. Write
$H_N^{\rm lab}(k)$ for the corresponding count when edges are identified
only by their coordinate-label pairs.

\begin{lemma}[Total matching length and a deterministic response bound]
\label{lem:history-lower}
For either parity,
\begin{equation}\label{eq:total-matching-length}
 \EE\mathcal L_n\sim\frac{\log n}{4e^\gamma}.
\end{equation}
For odd $N\ge3$ and every $1\le k\le N$,
\begin{equation}\label{eq:history-lower}
 \EE H_N(k)\ge\EE H_N^{\rm lab}(k)
 \ge\max\left\{2\EE\mathcal L_{N-1}-2,\
                       k(1-1/N)^2\right\}.
\end{equation}
\end{lemma}
\begin{proof}
The count path gives the exact identity
\[
 2\mathcal L_n=\int_0^{1/2}
                 (K_n(r)-\ind_{\{n\ {\rm odd}\}})\,dr.
\]
Put $a_n=\sqrt{\log n}$ and $b_n=1/\log n$.
The contribution below $a_n/n$ is at most $a_n$, and that above $b_n$
is at most $\log(1/(2b_n))$, using $K_n(r)\le1/r$ for $r\le1/2$.
The full intermediate theorem implies
\[
 \sup_{a_n/n\le r\le b_n}
   |r\,\EE K_n(r)-(2e^\gamma)^{-1}|\longrightarrow0.
\]
Otherwise a deterministic sequence of violating cutoffs would
contradict its $L^2$ conclusion. Integrating over the middle range
proves \eqref{eq:total-matching-length}.

Fix an even background $B$ with distinct pair distances. For each
background edge $e$, let $s_e$ be its separation and let $m_e$ be the
Lebesgue mass of insertions for which that edge is absent from the
augmented matching. The cutoff interpretation of the collapse map
shows that $m_e$ is exactly the mass moved by its two cells.
If at least four background points remain, let
$L_e=a_e+s_e+c_e$ be the new gap length after the event.
For the last pair put $L_e=1$. The middle child interval contributes
all its mass $s_e$, and the cells lie in the parent interval. Thus
\[
 s_e\le m_e\le L_e.
\]
For a three-gap merger, the increase of the gap-square sum is
$2a_es_e+2a_ec_e+2s_ec_e\ge s_eL_e$.
For the final pair, formally join its two gaps into a gap of length
one; the increase is $2s_e(1-s_e)\ge s_e=s_eL_e$.
Telescoping therefore gives $\sum_e s_eL_e\le1$.

Two independent insertions disagree about the presence of the
background edge $e$ with probability $2m_e(1-m_e)$. These edges alone
give
\[
 \EE[H_N^{\rm lab}(1)\mid B]\ge
 2\sum_e m_e(1-m_e)
 \ge2\sum_e s_e(1-L_e)\ge2\mathcal L(B)-2.
\]
For general $k$, apply the product-space decomposition from
Theorem~\ref{thm:noise-profile} to the vector of indicators of all
coordinate-label matching edges. Its squared norm is $(N-1)/2$,
and the expected squared difference under resampling is
$\EE H_N^{\rm lab}(k)$. Its nonnegative spectral formula makes this
expectation increasing and concave in $k$.
Each possible label edge has probability $1/N$ by exchangeability.
Independent inputs consequently give
$\EE H_N^{\rm lab}(N)=(N-1)^2/N$.
Monotonicity and the chord bound prove \eqref{eq:history-lower}.
Finally, identifying edges by label pairs can only decrease their
symmetric difference. In fact
$0\le H_N(k)-H_N^{\rm lab}(k)\le2k$ almost surely: the extra terms
are common label edges incident to refreshed coordinates.
\end{proof}

The upper bound requires uniformity in finite Poisson intensity; the
limiting gap-density envelope alone does not supply it. The following
positive representation gives that uniformity.

\begin{lemma}[A uniform finite-intensity exponential envelope]
\label{lem:finite-intensity-envelope}
Write $f_\alpha(y)=w_1^\alpha(y)\ind_{\{y\ge1\}}$ and
$\beta_\alpha=\min(\alpha,1)$, where $w$ is the finite-intensity
clearance field of Theorem~\ref{thm:scalar-transform}.
There are constants $0<b_{\rm h}^-\le b_{\rm h}^+<\infty$,
independent of $\alpha>0$, and numbers $\kappa_\alpha>0$, such that
\begin{equation}\label{eq:finite-intensity-envelope}
 b_{\rm h}^-\beta_\alpha e^{-\kappa_\alpha y}
 \le f_\alpha(y)\le
 b_{\rm h}^+\beta_\alpha e^{-\kappa_\alpha y},
 \qquad y\ge1.
\end{equation}
The exponents satisfy $\kappa_\alpha<\alpha$ and, for $\alpha\ge1$,
$\kappa_\alpha\le\log2$.
\end{lemma}
\begin{proof}
Let
\[
 \operatorname{Ein}(v)=\int_0^v\frac{1-e^{-u}}u\,du,\qquad
 c_\alpha=\frac{e^{\operatorname{Ein}(2\alpha)}}{2\alpha}.
\]
Take a Poisson sum $V_\alpha$ with intensity
$(1-e^{-2\alpha u})\,du/u$ on $(0,1)$, and independently
$E_\alpha\sim\operatorname{Exp}(2\alpha)$.
The density $g_\alpha$ of $Z_\alpha=V_\alpha+E_\alpha$ has
$g_\alpha(0)=1/c_\alpha$. Its Laplace exponent has Levy density
$k_\alpha(u)/u$, where
\[
 k_\alpha(u)=
 \begin{cases}1,&0<u<1,\\ e^{-2\alpha u},&u>1.\end{cases}
\]
Differentiating the Laplace transform and inverting gives
$y g_\alpha(y)=(k_\alpha*g_\alpha)(y)$, and hence, for $y>1$,
\begin{equation}\label{eq:finite-Dickman-derivative}
 y g_\alpha'(y)=
 -(1-e^{-2\alpha})g_\alpha(y-1)
 -2\alpha\int_1^y e^{-2\alpha u}g_\alpha(y-u)\,du.
\end{equation}
Below one the derivative is zero. These equations also justify the
piecewise differentiation, starting from the constant density on
$[0,1]$. Thus $g_\alpha$ is nonincreasing and
$a_\alpha=-c_\alpha g_\alpha'$ is a probability density supported on
$[1,\infty)$.

For $s>0$, integration by parts and the Poisson transform give
\[
 \widehat a_\alpha(s)
 =1-c_\alpha s\,\EE e^{-sZ_\alpha}
 =1-\exp\{-E_1(s)+E_1(s+2\alpha)\}.
\]
The scalar transform therefore yields the positive series
\begin{equation}\label{eq:finite-positive-geometric}
 f_\alpha=\sum_{j\ge1}2^{-j}a_\alpha^{*j}.
\end{equation}
The moment generating function of $a_\alpha$ is finite below
$2\alpha$, and at $t=\alpha$ it is
$1+c_\alpha\alpha\,\EE e^{\alpha Z_\alpha}>2$.
There is consequently a unique $\kappa_\alpha\in(0,\alpha)$ with
$\int e^{\kappa_\alpha y}a_\alpha(y)\,dy=2$.
For $\alpha\ge1$, support in $[1,\infty)$ also gives
$\kappa_\alpha\le\log2$.

We give the uniform bounds needed for renewal, rather than invoking
a parameter-uniform renewal theorem. In this paragraph write
$\beta=\beta_\alpha$,
$\bar g_\alpha(z)=g_\alpha(z/\beta)/\beta$, and
$\bar a_\alpha(z)=a_\alpha(z/\beta)/\beta$.
The quantity $\beta c_\alpha$ lies in a fixed compact subinterval of
$(0,\infty)$. The random variables $\beta Z_\alpha$ have uniformly
bounded exponential moments of every order less than two: their
Poisson part has Levy intensity bounded by $du/u$ on $(0,1)$,
and their exponential part has rate at least two.
Monotonicity of $\bar g_\alpha$, its bounded value at zero, and an
exponential tail bound imply
$\bar g_\alpha(z)\le C_b e^{-bz}$ for each fixed $b<2$.

Their densities also have a uniform positive lower bound on each
fixed compact interval. For $\alpha\le1$, the event $V_\alpha=0$
gives
$\bar g_\alpha(z)\ge2e^{-\operatorname{Ein}(2)}e^{-2z}$.
For $\alpha\ge1$, the Levy description decomposes $Z_\alpha$ into
the Dickman variable $V$ of Lemma~\ref{lem:positive-renewal} and an
independent Poisson sum on $(1,\infty)$ of total intensity
$E_1(2\alpha)$. Hence
$g_\alpha(z)\ge e^{-E_1(2)}\varphi(z)$, where $\varphi$ is the
strictly positive Dickman density.

Scale \eqref{eq:finite-Dickman-derivative}. If $\alpha\le1$, it reads
\[
 \bar a_\alpha(z)=\frac{\beta c_\alpha}{z}
 \left[(1-e^{-2\alpha})\bar g_\alpha(z-\alpha)
       +2\int_\alpha^z e^{-2u}\bar g_\alpha(z-u)\,du\right],
 \qquad z\ge\alpha.
\]
For $\alpha\ge1$ use \eqref{eq:finite-Dickman-derivative} without
rescaling. The bounds just proved imply, uniformly in $\alpha$,
\begin{equation}\label{eq:uniform-renewal-input}
 \bar a_\alpha(z)\le C e^{-3z/2},\qquad
 \inf_{\substack{\alpha>0\\\beta_\alpha\le z\le L}}
                       \bar a_\alpha(z)>0\quad(L\ge1\text{ fixed}).
\end{equation}
For the lower bound when $\alpha\le1$, the two bracketed terms,
after division by $z$, bound respectively constant multiples of
$\alpha/z$ and $(z-\alpha)/z$. For the upper bound the integral
length cancels $1/z$ near zero; away from zero the exponential
integral is bounded uniformly. When $\alpha\ge1$, $z\ge1$ and
$2\alpha/(2\alpha-3/2)\le4$, which gives the same upper bound.

Put
\[
 q_\alpha(z)=\tfrac12 e^{(\kappa_\alpha/\beta_\alpha)z}
                                    \bar a_\alpha(z).
\]
These are probability densities. Since
$\kappa_\alpha/\beta_\alpha\le1$, they have a uniform exponential
upper envelope and the compact lower bounds in
\eqref{eq:uniform-renewal-input}. In particular their increments have
a common positive probability $p$ of lying in $[2,3]$.
For their renewal epochs $S_j$, including $S_0=0$, the expected
number in any interval of length one is at most $1+1/p$: start at
the first epoch in that interval and wait for an increment of length
at least two. Thus
\[
 u_\alpha(z):=\sum_{j\ge1}q_\alpha^{*j}(z)
       =\EE\sum_{j\ge0}q_\alpha(z-S_j)
\]
is bounded above uniformly, by summing the exponential envelope over
unit intervals.

For completeness this also gives a uniform lower bound at large $z$.
Let $T$ be the first renewal epoch after $t$. Its overshoot has tail
at most
\[
 \overline Q_\alpha(b)+
       \bigl(\sup_{\alpha,z}u_\alpha(z)\bigr)
                       \int_b^\infty\overline Q_\alpha(v)\,dv.
\]
Choose a fixed $b$ making this at most $1/2$, and set $L=b+2$.
For $z\ge L$, apply the bound at $t=z-L$.
With probability at least $1/2$, $z-T\in[2,L]$; the next renewal
has density at least $\inf_{\alpha,\,2\le v\le L}q_\alpha(v)>0$
at distance $z-T$. Regeneration gives a uniform positive lower
bound for $u_\alpha(z)$. On $\beta_\alpha\le z\le L$, simply use
$u_\alpha(z)\ge q_\alpha(z)$ and
\eqref{eq:uniform-renewal-input}.
Finally, tilting \eqref{eq:finite-positive-geometric} gives
$f_\alpha(y)/\beta_\alpha
 =e^{-\kappa_\alpha y}u_\alpha(\beta_\alpha y)$.
The two renewal bounds prove \eqref{eq:finite-intensity-envelope}.
\end{proof}

\begin{lemma}[A logarithmic bound for the insertion ancestry]
\label{lem:insertion-height}
For an even background $B_M$, form its ternary gap-merger trees and
join the final two gaps by one binary root. For an independent
uniform location $z$, let $D_M(z)$ be the number of ancestral mergers
of the initial gap containing $z$, including that final root.
There is $\rho_{\rm h}<1$, independent of $M$, such that
\begin{equation}\label{eq:height-tail}
 \PP(D_M(z)\ge d+2)
 \le\min\{1,eM^{3/2}\rho_{\rm h}^{\,d}\},
 \qquad M\ge2\text{ even},\ d\ge0.
\end{equation}
In particular $\EE D_M(z)\le C(1+\log M)$.
\end{lemma}
\begin{proof}
First use a Poisson input of intensity $\lambda$. Inside a shielded
interval of length $x$ and cutoff $t\le x$, the empty probability is
$\lambda e^{-\lambda x}/w_t^\lambda(x)$, and the split density is
\[
 \frac{w_s^\lambda(a)w_s^\lambda(s)
                         w_s^\lambda(x-a-s)}{w_t^\lambda(x)}\,da\,ds.
\]
Conditional on the split separation $s$, the density of $a$ is
proportional to $w_s^\lambda(a)w_s^\lambda(x-a-s)$.
Scaling Lemma~\ref{lem:finite-intensity-envelope} makes its exponential
factor constant in $a$. Its density is therefore bounded below by
$\theta/(x-3s)$, where $\theta=(b_{\rm h}^-/b_{\rm h}^+)^2>0$.

If $s\le x/6$, the event
$x/4\le a\le3x/4-s$ has conditional probability at least $\theta/3$,
and every daughter then has length at most $3x/4$.
If $s>x/6$, every daughter has length less than $2x/3$.
Consequently the expected length of the daughter containing an
independent uniform location is at most
$\rho_{\rm h}x$, with $\rho_{\rm h}=1-\theta/12$.
Empty intervals stop the ancestry. Conditional on the daughter chosen,
the location remains uniform and its subtree has the corresponding
Poisson shielding law. After $d$ generations the expected surviving
interval length is at most $x\rho_{\rm h}^d$.

Conditioning on shielding increases the probability of an empty
Poisson interval, because the empty event implies shielding.
Thus its probability of being nonempty is at most
$1-e^{-\lambda x}\le\lambda x$.
The probability that the selected ancestry has at least $d+1$
internal mergers is therefore at most $\lambda x\rho_{\rm h}^d$.
For a nonempty even Poisson circle, condition on its final pair.
The two arcs have independent Poisson interiors conditioned on
shielding by that pair's separation. Choosing the arc containing a
uniform location and adding the final binary root gives a tail at
most $\lambda\rho_{\rm h}^d$ for depth $d+2$.

Now set $\lambda=M$ and condition on the exact Poisson count $M$.
The conditional input is exactly the original iid fixed-count sample.
Moreover
$\PP(\operatorname{Poisson}(M)=M)\ge e^{-1}M^{-1/2}$.
For this elementary bound, concavity of $\log$ gives
$\log M!\le(M+1/2)\log M-M+1$ by the trapezoidal integral bound.
Dividing the Poisson tail by this count probability proves
\eqref{eq:height-tail}. Summing its minimum with one gives the
logarithmic expectation. No limit is inferred by conditioning a
Poisson asymptotic statement: the bound before conditioning is uniform.
\end{proof}

\begin{theorem}[Changed histories have logarithmic order for fixed $k$]
\label{thm:logarithmic-history}
For a universal finite $C$, every odd $N\ge3$, and $1\le k\le N$,
\begin{equation}\label{eq:history-upper}
 \EE H_N(k)\le\min\{N-1,Ck(1+\log N)\}.
\end{equation}
Together with \eqref{eq:history-lower}, this proves
$\EE H_N(k)=\Theta_k(\log N)$ for every fixed $k\ge1$.
More precisely the lower coefficient is at least
$1/(2e^\gamma)$, uniformly over $k\ge1$.
There is also the tail bound
\begin{equation}\label{eq:history-tail}
 \PP(H_N(k)\ge4k(d+2))
 \le\min\{1,2ke(N-1)^{3/2}\rho_{\rm h}^{\,d}\},\qquad d\ge0.
\end{equation}
Section~\ref{sec:growing-perturbations} sharpens the expected order
uniformly over $k$. Section~\ref{sec:sublinear-history} supplies the
sharp coefficient for every $k=o(N)$.
Sections~\ref{sec:dense-history}--\ref{sec:history-concentration}
supply the dense-edit law and growing-count concentration.
Section~\ref{sec:late-history-variance} proves the centered CLT
throughout the growing sublinear range;
Section~\ref{sec:dense-history-CLT} gives the positive-density CLT.
\end{theorem}
\begin{proof}
In an add-one coupling, a background edge is absent precisely when
its collapse map moves the current insertion image. The two completed
matchings have equal numbers of edges, so their symmetric difference
has twice as many edges as there are such moves.
A move can occur only at an ancestor of the original insertion gap:
all preceding maps preserve the mass and support inside the current
active gaps. Therefore the add-one cost is at most $2D_M(z)$.
With the even background of size $M=N-1$, compare the two independent
insertions through that background. Their distance is at most
$2D_M(U)+2D_M(V)$. Lemma~\ref{lem:insertion-height} gives the mean
bound for one resampling and its tail in \eqref{eq:history-tail}.

Perform $k$ resamplings successively. Each intermediate input is still
iid uniform, although the costs need not be independent. The triangle
inequality, linearity of expectation, and a union bound give the
factors $k$ in \eqref{eq:history-upper}--\eqref{eq:history-tail}.
The deterministic bound $H_N(k)\le N-1$ completes the upper bound.
Lemma~\ref{lem:history-lower} supplies the logarithmic lower bound.
\end{proof}

\subsection{Uniform resampling sizes through exchangeable gaps}
\label{sec:growing-perturbations}

\status{Working proof: the expected history scale for all $k$}
The preceding logarithmic bound can be sharpened without a
parameter-uniform tree argument. The relevant finite symmetry concerns
the order of the current gaps, not their independence.
It is the circular bag argument of
\href{https://www.phys.ens.psl.eu/~derrida/PAPIERS/1994/bray-epl.pdf}
{Bray, Derrida and Godr\`eche, \emph{Non-Trivial Algebraic Decay in a
Soluble Model of Coarsening}, Europhys.\ Lett.\ 27 (1994), 175--180},
p.~177. Their deterministic shortest-domain rule has the same
three-gap merger. We state the conditional-order assertion needed
here; their persistence observable is different from our survivor
noise and does not identify its spectrum.

\begin{lemma}[Exchangeability at a fixed active count]
\label{lem:gap-order}
Start with $M$ iid uniform locations. At every active count $m\ge3$,
root the cyclic gaps at an independently uniform active point.
Conditional on their unordered lengths, their order is uniform.
The same assertion holds when each gap is decorated with its pushed
origin measure in coordinates measured from its left endpoint.
Let $s_{M,m}$ be their minimum and let $L_{M,m}$ be the parent gap
length created by the next merger. Then
\begin{align}
 \EE[L_{M,m}\mid\hbox{unordered gaps}]
   &=s_{M,m}+\frac{2(1-s_{M,m})}{m-1}
       \le\frac3m,\label{eq:parent-mean}\\
 \EE[s_{M,m}L_{M,m}]&\le\frac3{m^2}.
 \label{eq:weighted-parent}
\end{align}
For the last pair in an even sample use $L_{M,2}=1$; the two
inequalities remain valid. The expected insertion ancestry therefore
also has the elementary bound
\[
 \EE D_M(z)\le1+\frac32\sum_{j=2}^{M/2}\frac1j
 \qquad(M\ge2\hbox{ even}).
\]
\end{lemma}
\begin{proof}
Initially the rooted gap vector is exchangeable Dirichlet. Given
any unordered collection of distinct current lengths, rotate the
minimum to the first position. Its left and right neighbours are
a uniform ordered pair from the other $m-1$ lengths. Conditional
on this pair, the other lengths retain uniform order. Contracting
the three consecutive gaps leaves their sum and that uniform
ordering; rooting uniformly again makes the whole new gap vector
exchangeable. Mixing over the neighbour pair preserves this
property, even after conditioning on the new unordered collection.
This proves the assertion by induction. The argument does not
assert independence of the lengths.
For decorated gaps, the new measure is a deterministic function of
the three ordered daughter measures and lengths, and the other marks
are unchanged. The same induction therefore applies to the unordered
collection of decorated gaps.

The neighbours have conditional total mean $2(1-s_{M,m})/(m-1)$.
Since $s_{M,m}\le1/m$, this proves both bounds.
An independent uniform insertion belongs to a parent interval with
probability its length; summing these probabilities over the merger
tree proves the ancestry bound, including its final binary root.
\end{proof}

\begin{lemma}[Shielded edges under several insertions]
\label{lem:edge-insertions}
Fix a background with distinct pair distances and a matched edge
$e$ of separation $s_e$. Let $I_e$ be its parent interval, of length
$L_e$, with the full circle used for a final binary root.
If no inserted point belongs to $I_e$, then $e$ remains matched.
If an odd number of inserted points belongs to the shorter arc of
length $s_e$ between its endpoints, then $e$ is absent.

Consequently, if $q_e(k)$ is its conditional absence probability
under $k$ independent uniform insertions, then
\begin{equation}\label{eq:edge-insertion-probabilities}
 q_e(k)\le\min(1,kL_e),\qquad
 1-q_e(k)\ge(1-L_e)^k,\qquad
 q_e(k)\ge k s_e(1-L_e)^{k-1}.
\end{equation}
\end{lemma}
\begin{proof}
For a parent interval with distinct endpoints, both endpoints remain
active through the event $e$. Every original interior vertex is
matched within that interval, at distance strictly smaller than its
distance to either endpoint. Indeed either endpoint would otherwise
have provided an earlier partner. Every point outside the interval
is still farther away: a shortest circular arc to it crosses an
endpoint. Inducting through the original interior matching events
therefore preserves every such edge when all insertions are outside,
regardless of any changes to the exterior or to the boundary points.
The full-circle cases are immediate.

If $e$ is accepted, its two endpoints remained active until
separation $s_e$. No earlier matching edge could cross either of
them: the crossed endpoint would have been a closer available
partner. All points in the shorter arc must already have disappeared
internally, so their initial number was even. That number was even
in the background; an odd number of new points in the arc makes
acceptance impossible. The probability of exactly one insertion
there is $k s_e(1-s_e)^{k-1}$, which is at least the last bound in
\eqref{eq:edge-insertion-probabilities}. The other bounds follow
from the first assertion and a union bound.
\end{proof}

\begin{lemma}[Matching length above an active-count threshold]
\label{lem:truncated-matching-length}
Put $c=(2e^\gamma)^{-1}$. If integers $M,q$ satisfy $q\ge2$ and
$M/q\to\infty$, then, for either parity,
\begin{equation}\label{eq:truncated-matching-length}
 \sum_{\substack{q\le m\le M\\m\equiv M\ ({\rm mod}\ 2)}}
       \EE s_{M,m}
       \sim\frac c2\log(M/q).
\end{equation}
Terms with fewer than two active points are omitted.
\end{lemma}
\begin{proof}
Whenever $m\to\infty$ and $m/M\to0$, the intermediate-time theorem
at the two cutoffs $(c-\delta)/m$ and $(c+\delta)/m$ implies
$m s_{M,m}\to c$ in probability. To see this, $s_{M,m}>r$ is
equivalent, off collision times, to $K_M(r)\ge m$.
The two limiting count ratios straddle one. Since
$0\le m s_{M,m}\le1$, convergence also holds in mean.

Write $t=M/q$ and $a=\exp(\sqrt{\log t})$.
Uniformly for $aq\le m\le M/a$, we thus have
$m\,\EE s_{M,m}=c+o(1)$: a violating sequence would satisfy the
same two hypotheses. The parity-restricted harmonic sum on this
range is $\tfrac12\log t+o(\log t)$.
The omitted ranges have total expected contribution at most
$O(1+\log a)=o(\log t)$, using $s_{M,m}\le1/m$.
This proves the result.
\end{proof}

\begin{theorem}[The uniform expected history scale]
\label{thm:uniform-history}
There are absolute constants $0<c_{\rm H}\le C_{\rm H}<\infty$
such that, for every odd $N\ge3$ and $1\le k\le N$,
\begin{equation}\label{eq:uniform-history}
 c_{\rm H}\,k\bigl(1+\log(N/k)\bigr)
 \le\EE H_N^{\rm lab}(k)\le\EE H_N(k)
 \le\min\{N-1,C_{\rm H}\,k(1+\log(N/k))\}.
\end{equation}
More precisely, for every sequence $1\le k=k_N=o(N)$,
\begin{equation}\label{eq:uniform-history-coefficient}
 \liminf_{N\to\infty}
 \frac{\EE H_N^{\rm lab}(k)}
      {k\log(N/k)}\ge\frac1{e^\gamma}.
\end{equation}
These expectation bounds alone do not identify the leading constant.
Theorem~\ref{thm:sharp-small-edits} below identifies it for
$k=o(\log N)$. Fluctuations and lower bounds with high probability
remain open.
\end{theorem}
\begin{proof}
Use a common background of size $M=N-k$, with independent sets of
$k$ probes producing the two samples. It suffices for the upper
bound to take $k<N/2$, since otherwise $H_N(k)\le N-1\le2k$.
For either augmented matching let $R$ count missing background edges.
Its symmetric difference with the background has
$2R+\lfloor k/2\rfloor$ edges. The triangle inequality, then
Lemmas~\ref{lem:gap-order} and \ref{lem:edge-insertions}, give
\begin{align*}
 \EE H_N(k)
 &\le k+4\sum_{m\equiv M\ ({\rm mod}\ 2)}\EE\min(1,kL_{M,m})\\
 &\le k+4\sum_{\substack{2\le m\le M\\m\equiv M\ ({\rm mod}\ 2)}}
                                      \min(1,3k/m)\\
 &\le C_{\rm H}k(1+\log(N/k)).
\end{align*}
The first sum is over matching events. Concavity of $\min(1,\cdot)$
justifies its replacement by the conditional parent-length mean.

For the lower bound, conditional independence of the two probe
sets makes their disagreement about background edge $e$ have
probability $2q_e(k)(1-q_e(k))$. In fact the spatial history cost
is at least twice the number of these disagreements. To check this,
write $a,b$ for the missing-background-edge counts, $J$ for the
number missing in both copies, $I$ for their common nonbackground
edges, and $d=\lfloor k/2\rfloor$. A common nonbackground edge
has two background vertices, since the fresh locations in the two
copies are distinct. Each matched background vertex in it belongs
to an edge missing in both copies. There is at most one unmatched
background vertex, so $2I\le2J+1$ and therefore $I\le J$.
Consequently
\[
 H_N(k)=2(a+b)+2d-2J-2I
       \ge2(a+b-2J)+2d.
\]
Using $H_N^{\rm lab}(k)\ge H_N(k)-2k$ gives
\[
 \EE H_N^{\rm lab}(k)
 \ge4k\,\EE\sum_e s_e(1-L_e)^{2k-1}-2k.
\]
Keep only events with active count $m\ge q=\max(3,k)$.
Bernoulli's inequality and \eqref{eq:weighted-parent} give
\begin{equation}\label{eq:history-truncated-lower}
 \EE H_N^{\rm lab}(k)\ge
 4k\left(
 \sum_{\substack{q\le m\le M\\m\equiv M\ ({\rm mod}\ 2)}}\EE s_{M,m}
       -C\,k/q\right)-2k,
\end{equation}
since $\sum_{m\ge q}3/m^2\le C/q$.
When $k=o(N)$, Lemma~\ref{lem:truncated-matching-length} proves
\eqref{eq:uniform-history-coefficient}.

Finally this sequential asymptotic implies a uniform positive
constant in \eqref{eq:uniform-history}. Otherwise take a sequence
for which the lower ratio tends to zero. If $N/k$ has a bounded
subsequence, the chord bound
$\EE H_N^{\rm lab}(k)\ge k(1-1/N)^2$ contradicts that conclusion.
Otherwise pass to $N/k\to\infty$ and use
\eqref{eq:uniform-history-coefficient}. Bounded $N$ cannot cause
a failure, again by the strictly positive chord bound.
\end{proof}

\begin{corollary}[History cost for any mixture of random edits]
\label{cor:mixed-history}
From an odd iid sample of size $N\ge3$, delete $d$ uniformly chosen
coordinates and add $a$ fresh uniform points. Put $q=d+a$, and
assume $1\le q\le N$. If $H_N(d,a)$ counts the spatial symmetric
difference of the completed matchings, then, with absolute constants,
\[
 \EE H_N(d,a)\asymp q(1+\log(N/q))
\]
uniformly over all such $N,d,a$. If $q=o(N)$, its lower asymptotic
coefficient relative to $q\log(N/q)$ is at least
$c=(2e^\gamma)^{-1}$.
\end{corollary}
\begin{proof}
Use a common background of size $M=N-d$, with $d$ old and $a$
fresh independent probes. For $q<N/2$, the upper parent-cover
argument bounds the cost by
$q+4\sum_{m\le M}\min(1,3q/m)=O(q(1+\log(N/q)))$.
For $q\ge N/2$, the two matching sizes give a deterministic $O(q)$
bound.

For a background edge of separation $s$ and parent length $L$,
the two copies disagree if exactly one of the $q$ probes lies
in its short arc and all the others lie outside its parent.
This event has probability $qs(1-L)^{q-1}$.
The common-new-edge count is still bounded by the common-missing-edge
count, as in Theorem~\ref{thm:uniform-history}, so the full cost is
at least twice this background disagreement count.
Keeping active counts $m\ge\max(3,q)$ and using
\eqref{eq:weighted-parent} and
Lemma~\ref{lem:truncated-matching-length} proves the lower
coefficient for $q=o(N)$.

Every initial discrepant vertex except at most the two final
survivors is incident to a changed edge. Each such edge accounts
for at most two of the $q$ vertices, so $H_N(d,a)\ge(q-2)/2$.
For $q\ge3$ this is at least $q/6$.
For $q=1$, addition changes the matching size; deletion removes a
matched vertex with probability $1-1/N$. For $q=2$, pure addition
or deletion changes the matching size, and resampling has the
positive chord bound of Lemma~\ref{lem:history-lower}.
These cover the small cases. The same subsequence argument as
before gives a uniform positive constant for all $q$.
\end{proof}

\subsection{Insertion and deletion have the same limiting response law}
\label{sec:edit-universality}

\status{Working proof: finite edit types and deletion basin weights}
Write $\xi_M=M^{-1}\sum_{u\in B_M}\delta_u$ for the initial empirical
probability measure. All collapse maps below are those of the
unperturbed background.

\begin{lemma}[Replacing the initial mass by empirical mass]
\label{lem:empirical-seed}
At a fixed positive cutoff $\varepsilon$, push both $dx$ and $\xi_M$
through the background collapse maps. Restricted to its active gaps,
their measures differ by a quantity tending to zero in probability
in bounded-Lipschitz distance, summed over the gaps.
The total empirical mass at the active endpoints is $O(1/(M\varepsilon))$.
The assertion holds along either parity of $M$.

For even $M$, let $\theta_v(B_M)$ be its full insertion basin weights
and let $d_v(B_M)$ be the deletion-graph indegrees. For every fixed
$\varepsilon>0$,
\begin{equation}\label{eq:late-deletion-masses}
 \sum_{v\in C_M(\varepsilon)}
       |d_v(B_M)/M-\theta_v(B_M)|
       \longrightarrow0\quad\hbox{in }L^1.
\end{equation}
\end{lemma}
\begin{proof}
Set
\[
 \Delta_M=\sup_{I\ {\rm circular\ interval}}
                  |\xi_M(I)-|I||+1/M.
\]
Then $\EE\Delta_M\to0$. One elementary proof first controls empirical
counts on a fixed grid by their variances, brackets every interval
by grid intervals, and then refines the grid. Boundedness upgrades
convergence in probability to convergence in mean. This uniform
control applies to random intervals selected from the same sample;
no independence from the tree geometry is needed.

Reveal $d$ generations of each cutoff gap tree. There are at most
$C_{\varepsilon,d}$ terminal intervals and revealed vertices.
On a terminal interval $I$, complete the two measures by Lebesgue
mass $|I|$ and by a constant density of total mass $\xi_M(I)$,
respectively. Their completed measures differ in total variation
by at most $C_{\varepsilon,d}\Delta_M$, even after applying the
revealed maps. Initial empirical atoms on revealed vertices have total mass at most
$C_{\varepsilon,d}/M$ and can be discarded with that error.

For a cumulative query, the unresolved empirical budget is at most
the length budget $Z_d$ plus $C_{\varepsilon,d}\Delta_M$.
Indeed its terminal interval masses sum to their lengths with that
error. Comparing each true measure with its completion therefore gives
\[
 |F_{\rm emp}(z)-F_{\rm Leb}(z)|
       \le2Z_d+C_{\varepsilon,d}\Delta_M.
\]
The same estimate controls any mass terms in the cumulative recursion.
Finite-tree total-variation convergence and
Lemma~\ref{lem:query-contraction} imply that the limsup of the expected
integrated error, summed over gaps, is at most $2\rho^d$.
The total-mass errors also vanish. The line formula for transport,
or integration against Lipschitz tests, proves the first assertion
by sending $d$ to infinity. For fixed $\varepsilon$ the number of
gaps is bounded. Original active endpoints carry only their original
mass $1/M$, since earlier maps never target them.

For the second assertion, an insertion at the position of an
original point $u$ immediately annihilates with $u$, leaving
$B_M-u$. Equivalently, the full collapse-map composition sends $u$
to $D_{B_M}(u)$. This follows also by taking an insertion sufficiently
close to $u$, so that this is the first collision.
Thus pushing $\xi_M$ through all the maps gives masses $d_v/M$.
The late-label readout in the proof of
Theorem~\ref{thm:noise-convergence} uses finitely many cumulative
queries in the cutoff gaps. Apply the preceding budget comparison
to those queries: its limsup expected error is
$O_\varepsilon(\rho^d)$, then let $d\to\infty$.
The endpoint mass discarded above is negligible. This proves
\eqref{eq:late-deletion-masses}.
\end{proof}

\begin{lemma}[All limiting insertion mass lies on tree vertices]
\label{lem:atomic-gap-mass}
The limiting pushed measure of a shielded interval $(t,x)$ is
purely atomic, supported on its countable set of internal collision
vertices. This holds for every $0<t\le x$.

More quantitatively, reveal $d$ generations of its tree and let
$D_d(t,x)$ be the total length of terminal intervals reached from
the root using only side daughters, never a middle daughter.
There is a universal $\rho_{\rm a}\in(0,1)$ such that
\begin{equation}\label{eq:side-mass-decay}
 \EE D_d(t,x)\le
 x\min\left(1,(x/t)\rho_{\rm a}^{\,d}\right).
\end{equation}
\end{lemma}
\begin{proof}
Choose a uniform location in the root interval, independently of
the tree, and follow its unique daughter interval backwards.
Kill this ancestry when it chooses the middle daughter.
Its conditional probability of surviving $d$ generations is
$D_d(t,x)/x$. Conditional on a chosen side daughter, the location
remains uniform in it and its tree has the corresponding shielding law.

Write $R=x/t\ge1$ and let $S$ be the last split separation.
The scalar field and its exponential envelope give, for $v\ge1$,
\[
 \PP(t/S>v)=\frac{w_{t/v}(x)}{w_t(x)}
 \le K v\exp\{-\kappa R(v-1)\},
 \qquad K=b_+/b_-.
\]
Thus, for a universal $K_0$,
\begin{equation}\label{eq:ancestry-ratio-bound}
 \EE(t/S)\le1+\frac{K_0}{R}.
\end{equation}
Indeed integrate the displayed tail, using
$\int_1^\infty v e^{-\kappa R(v-1)}\,dv
 =(\kappa R)^{-1}+(\kappa R)^{-2}$.
Choose $u\in(0,1)$ so small that
$K u^{-1}e^{-\kappa(u^{-1}-1)}\le1/2$.
Uniformly over all states, $\PP(S\ge ut)\ge1/2$.
The probability of choosing the middle daughter is consequently
$\EE(S/x)\ge b/R$, where $b=u/2>0$.

Conditional on $S=s$, the balanced-daughter argument of
Lemma~\ref{lem:insertion-height}, now using $b_\pm$, bounds the
expected chosen daughter length by $\rho_*x$ with a fixed
$\rho_*<1$. Its new cutoff is $s$. Killing middle choices only
reduces this bound, so \eqref{eq:ancestry-ratio-bound} yields
\[
 \EE[R_{\rm next}\ind_{\{\rm survive\}}\mid R]
      \le\rho_*R+\rho_*K_0.
\]
Write $B=\rho_*K_0$, $\eta=(1-\rho_*)/2$, and choose
\[
 A=\max\{1,(B+1)^2/(4\eta b)\},\qquad
 \delta=\min\{\eta,1/(2A)\},\qquad \rho_{\rm a}=1-\delta.
\]
For the killed transition and $R\ge1$, the drift and killing bounds give
\begin{align*}
 \EE[(R_{\rm next}+A)\ind_{\{\rm survive\}}\mid R]
 &\le\rho_*R+B+A-Ab/R\\
 &\le(1-\delta)(R+A).
\end{align*}
Indeed $(1-\rho_*-\delta)R+Ab/R
\ge\eta R+Ab/R\ge2\sqrt{\eta Ab}\ge B+1>B+\delta A$.
Iterating this inequality and using $R_d\ge1$ on survival gives
\[
 \PP(\hbox{alive after }d)
 \le\frac{R_0+A}{1+A}\rho_{\rm a}^{\,d}
 \le R_0\rho_{\rm a}^{\,d},\qquad R_0=x/t.
\]
This proves \eqref{eq:side-mass-decay}. A first-moment bound on
$R_d$ alone would give only a polynomial estimate; the added constant
in this killed Lyapunov function uses the uniform chance of absorption
when $R$ is bounded.

Let $V^{(d)}$ be the finite set of vertices revealed by depth $d$.
Every origin interval that enters a revealed middle daughter sends
all its mass to that daughter's endpoint pair at the parent
collapse, regardless of the unresolved finer maps.
Subsequent revealed maps leave that mass on $V^{(d)}$.
Every deeper completion therefore places at least $x-D_d(t,x)$
mass on this finite set. The same is true of the limiting measure,
by the closed-set inequality for weak convergence.
Since $D_d\downarrow0$ almost surely by
\eqref{eq:side-mass-decay}, its total mass lies on
$\bigcup_d V^{(d)}$.
\end{proof}

The side-daughter recursion is the dry-mass recursion in the
coarsening model cited in Section~\ref{sec:growing-perturbations}.
The proof here uses depth in the limiting tree and supplies only
\eqref{eq:side-mass-decay}; it does not import a physical-time
persistence exponent from that model.

\begin{theorem}[Universality over finitely many insertion and deletion types]
\label{thm:edit-universality}
Start with an iid uniform background $B_M$ of a fixed parity.
Choose finitely many distinct original labels uniformly, and finitely
many independent fresh uniform points. Regard deleting each chosen
label or adding each fresh point as one edit. Form any fixed finite
family of copies by prescribing a subset of these edits in each copy.

At every positive cutoff, the family's joint limiting law, with
shared-point identities and future histories, depends on the edit
subsets and the background parity, but not on which edits were
deletions and which were insertions. It equals the law obtained
by making every edit an independent insertion into the background.

In particular, start with odd $M$ and make any mixture of $2k$
uniform deletions and independent uniform insertions, with $k$ fixed.
The probability that its final survivor differs from the original
one tends to $p(k)$ from Corollary~\ref{cor:noise-evolution}.
Thus adding $2k$ fresh points, deleting $2k$ uniform original points,
and resampling $k$ coordinates have the same limiting survivor noise.
\end{theorem}
\begin{proof}
Represent a deletion by an extra point at the selected original
location, whose distance-zero pair annihilates immediately. Distinct
such pairs commute; equivalently take sufficiently small disjoint
perturbations and then let their displacements tend to zero.
The collapse-map readout therefore applies to deletion seeds too.

At cutoff $\varepsilon$, a uniform deletion seed hits a background
active point with probability at most $1/(M\varepsilon)$.
Otherwise it lies in a cutoff gap. Sampling a fixed number $q$ of
original labels without replacement differs from independent
sampling from $\xi_M$ by probability at most $\binom q2/M$.
By Lemma~\ref{lem:empirical-seed}, the conditional mapped-seed laws
in all gaps converge jointly to the same measures as fresh uniform
seeds. Products for finitely many independent seeds converge too.
The readout boundaries have zero limiting mass by
Lemma~\ref{lem:gap-measure-limit}.

The probability that two seed gaps coincide or are adjacent has the
same limiting bound $O_q(\varepsilon)$ as in
Lemma~\ref{lem:separated-insertions}: empirical gap masses differ
from lengths by at most $\Delta_M$, and there are at most
$1/\varepsilon$ gaps. On the separated event all subset readouts
are exact for the mixed edits, just as for insertions alone.
Every true subsequential marked law is therefore within
$O_q(\varepsilon)$ in total variation of the same limiting
auxiliary law. Push to a fixed later cutoff, then send
$\varepsilon\downarrow0$, exactly as in
Theorem~\ref{thm:joint-perturbed-limit}.

Fresh points themselves are active at $\varepsilon$ with probability
$O_q(1/(M\varepsilon))$, by exchangeability in each copy.
Selected original points, in any copy that retains them, have the same
bound. Hence these identities cause no surviving exceptional marks.
The union-separation argument again excludes accidental coincidence
of different limiting identities. This proves the claimed joint
law and future-history convergence.

When $2k$ edits are made to an odd background both copies remain
odd. One permitted assignment is $k$ deletions and $k$ fresh
insertions paired as coordinate resamplings. Its agreement law is
the already established $p(k)$, so all assignments have that limit.
Position and common-label comparisons differ only when one of the
finitely many edited labels wins, an event of probability $O_k(1/M)$.
\end{proof}

\begin{corollary}[Deletion degrees approximate insertion basins]
\label{cor:deletion-basin}
For even $M$,
\begin{equation}\label{eq:deletion-basin-l2}
 \EE\sum_{v\in B_M}
      \bigl(d_v(B_M)/M-\theta_v(B_M)\bigr)^2\longrightarrow0.
\end{equation}
In fact convergence also holds in expected $\ell^1$ distance.
Their ranked sequences converge in law in $\ell^1$ to the same
random probability vector. Its expected squared norm is $A_\infty$.
There is no loss of total mass into vanishing individual basins.
\end{corollary}
\begin{proof}
The deletion-graph identity gives exactly
\[
 \EE\sum_v(d_v/M)^2
    =\frac{M-1}{M}A_{M-1}+\frac1M\longrightarrow A_\infty,
 \qquad
 \EE\sum_v\theta_v^2=A_{M+1}\longrightarrow A_\infty.
\]
By \eqref{eq:late-deletion-masses}, their expected inner product
over the late labels differs from the expected late insertion
squared norm by $o(1)$. The early-label squared-mass bound tends
to zero as $\varepsilon\downarrow0$, so positivity and
Cauchy--Schwarz imply
$\EE\sum_v(d_v/M)\theta_v\to A_\infty$.
Expanding the square proves \eqref{eq:deletion-basin-l2}.

To strengthen this to total mass, construct the limiting two arc
trees of the final even pair. Their pushed measures are purely
atomic by Lemma~\ref{lem:atomic-gap-mass}, with total masses $L$
and $1-L$. The final collapse preserves atomicity and total mass.
Each resulting atom is a vertex removed by the background at a
strictly positive separation. Lemma~\ref{lem:atomic-gap-mass}
places every child's mass on its internal vertices, so it gives
zero mass to the child's endpoints before the parent collapse. Thus the atom
weights at revealed vertices agree with the limits of the named
finite basin weights; no extra mass arrives at a retained label
through a shrinking identity region.

It follows that the limiting insertion weights on labels active
at $\varepsilon$ increase in total mass to one as
$\varepsilon\downarrow0$. Their finite vectors converge, and
$\EE\sum_v\theta_v=1-\EE\Theta_{M+1}\to1$.
Consequently
\[
 \lim_{\varepsilon\downarrow0}\limsup_{M\to\infty}
 \EE\sum_{v\notin C_M(\varepsilon)}\theta_v=0.
\]
Equation~\eqref{eq:late-deletion-masses} and $\sum_vd_v/M=1$
give the same tail bound for deletion mass. Splitting the
$\ell^1$ difference into late and early labels now proves
$\EE\sum_v|d_v/M-\theta_v|\to0$.

The finite late vectors give converging-together approximations
to the full ranked vectors in $\ell^1$, by these total-mass tail
bounds. Decreasing rearrangement contracts $\ell^1$ distance
(match the ordered coordinates, or integrate the count above each
threshold). It follows that both ranked limits exist and agree.
Total mass is continuous in $\ell^1$, and the squared norms are
bounded and continuous there, giving the two stated normalizations.
\end{proof}

\subsection{Spatial displacement of an input point}
\label{sec:spatial-noise}

\status{Working proof: a continuous spatial noise profile}
The position from which an insertion starts can also be retained.
This requires more than convergence of its pushed mass measure.

\begin{lemma}[Convergence of insertion maps with the same origin]
\label{lem:origin-map}
In a limiting shielded interval, let $\Phi^{(d)}$ be the map obtained
from its first $d$ tree generations, using the identity map on each
terminal interval. There is a jointly measurable limiting map
$\Phi$, defined for almost every original location, such that
\begin{equation}\label{eq:origin-map-error}
 \EE\Leb\{z:\Phi^{(d)}(z)\ne\Phi(z)\}
       \le C x(x/t)\sigma^d,
\end{equation}
where $C<\infty$ and $\sigma\in(0,1)$ are universal.
Its pushforward of Lebesgue measure is the earlier limiting gap
measure, and almost all its values are internal tree vertices.

For an even circle background there is consequently a random
survivor map $S_\infty:\T\to\T$, with countable range up to a null
set. The finite backgrounds and this map can be coupled so that
\begin{equation}\label{eq:survivor-map-coupling}
 \EE\Leb\{z:S_{B_M}(z)\ne S_\infty(z)\}\longrightarrow0.
\end{equation}
The equality in these estimates means the same physical vertex,
not merely nearby positions. They assert these couplings, not
total-variation convergence of the law of the entire random map.
\end{lemma}
\begin{proof}
The cumulative-query argument has a pointwise version. For two
completions below a revealed depth $d$, the symmetric difference
of their threshold preimages
\[
 \{u:\Phi_1(u)\le z\}\ \triangle\ \{u:\Phi_2(u)\le z\}
\]
is contained in the union of the terminal origin intervals followed
by that query. Its length is therefore at most $Z_d(t,x,z)$.
To verify this, apply the recursion \eqref{eq:cdf-query} to a
single origin: only its own daughter contributes, possibly with
the indicator complemented in the middle daughter. A pruned
branch has already determined that indicator. Iterate this
observation. The natural completions satisfy the same null-boundary
conditions as in Lemma~\ref{lem:query-contraction}.

Fix $1\le h\le d$ and compare depth $d$ with any deeper completion.
Origins that have not entered a middle daughter in the first $h$
generations have total length $D_h(t,x)$.
Every other origin has a first such daughter, say at level $j\le h$.
At its parent collapse its image becomes one of the two parent
vertices. Which vertex it becomes is determined by the midpoint
query in that middle daughter, with at least $d-j$ revealed
generations. Once the two completions make the same choice, all
remaining ancestor maps give identical outputs.

At any fixed level the middle-daughter lengths sum to at most $x$.
Conditioning on their states and using the threshold-preimage
bound thus bounds the expected disagreement by
\[
 \EE D_h(t,x)+x\sum_{j=1}^h\rho^{d-j}
 \le x(x/t)\rho_{\rm a}^{\,h}
                   +\frac{x\rho^{d-h}}{1-\rho}.
\]
Take $h=\lfloor d/2\rfloor$; the finitely many smaller depths are
absorbed into $C$, with
$\sigma=\max(\sqrt{\rho_{\rm a}},\sqrt\rho)<1$.
The consecutive-depth disagreement bounds are summable.
Fubini and Borel--Cantelli give eventual exact agreement for almost
every origin.
Define $\Phi$ by that eventual value. The bound for arbitrary
depths follows by comparison with a deeper completion and Fatou.
This construction is measurable. Its pushforward is the previous
weak limit; Lemma~\ref{lem:atomic-gap-mass} identifies its support.

For the circle use the two final-pair arc trees, then the final
collapse. Their initial ratios are $1$ and $(1-L)/L$.
Since $\EE(1/L)=1+\EE Y=1+2e^\gamma<\infty$, their combined expected
map error is $C\sigma^d$. At any fixed depth their finite-sample
geometries converge in total variation, by the last-pair tree
argument in Proposition~\ref{prop:half-noise}. Couple those
geometries identically outside their error event.
The finite-sample disagreement with a depth-$d$ map is bounded
by exactly the same finite origin budgets just used. These bounded
budgets converge at each fixed depth, first restricting $L$ away
from zero and then removing that restriction. Their limsup
expectation is at most $C\sigma^d$.
Sending $M$ to infinity and then $d$ to infinity, or choosing a
slowly increasing depth, proves \eqref{eq:survivor-map-coupling}.
\end{proof}

\begin{corollary}[A power bound for total early basin mass]
\label{cor:early-total-mass}
There are universal $C<\infty$ and $\beta>0$ such that, for
$0<\varepsilon<1/2$,
\begin{equation}\label{eq:early-total-mass}
 \limsup_{\substack{M\to\infty\\M\ {\rm even}}}
 \EE\sum_{v\notin C_M(\varepsilon)}\theta_v(B_M)
 \le C\varepsilon^\beta.
\end{equation}
The same bound holds with $\theta_v(B_M)$ replaced by $d_v(B_M)/M$.
These are iterated-limit bounds on total mass.
Corollary~\ref{cor:basin-rate} below gives a uniform finite-$M$ bound.
\end{corollary}
\begin{proof}
Put $J=1+K_0$, using \eqref{eq:ancestry-ratio-bound}.
At every node with inherited cutoff $t$, its revealed split separation
$S$ satisfies $\EE(S^{-1}\mid t,x)\le J/t$.
The two final-pair arc trees both start with cutoff $L$, and
$\EE L^{-1}<\infty$. Thus the expected sum of reciprocal separations
of all pairs revealed through depth $d$, including the final pair,
is at most $C(3J)^d$. A union bound gives probability at most
$C\varepsilon(3J)^d$ that any such pair was removed by
$\varepsilon$.

Outside this event, every nonidentity value of the depth-$d$ map is
a vertex removed after $\varepsilon$.
The limiting map has countable range, so its equality with the identity
map holds only on a null set of origins. By
\eqref{eq:origin-map-error}, its expected mass on early vertices is
therefore at most
\[
 C\sigma^d+C\varepsilon(3J)^d.
\]
Choose $d=\lfloor\log(1/\varepsilon)/\log(3J/\sigma)\rfloor$.
This gives the claim with
$\beta=\log(1/\sigma)/\log(3J/\sigma)>0$ for the limiting map.
Convergence of the finite late-label vector and of total insertion
mass transfers it to \eqref{eq:early-total-mass}.
The expected $\ell^1$ comparison in
Corollary~\ref{cor:deletion-basin} transfers it to deletion mass.
\end{proof}

\begin{theorem}[A spatial noise profile and stability under small moves]
\label{thm:spatial-noise}
For odd $N$, let $\zeta_N(h)$ be the survivor-label disagreement
probability when one coordinate $X_1$ is replaced by $X_1+h$ modulo
one, for a deterministic $h\in\T$. There is a continuous, even,
periodic function $\zeta:\T\to[0,1]$ such that
\begin{equation}\label{eq:spatial-profile}
 \sup_{h\in\T}|\zeta_N(h)-\zeta(h)|\longrightarrow0,\qquad
 \zeta(0)=0,\qquad \int_{\T}\zeta(h)\,dh=p_\infty.
\end{equation}
For a uniform $U$ independent of the limiting map,
$\zeta(h)=\PP(S_\infty(U)\ne S_\infty(U+h))$.
There are nonnegative weights $b_\ell=b_{-\ell}$, $\ell\in\mathbb Z$,
with
\begin{equation}\label{eq:spatial-spectrum}
 \sum_\ell b_\ell=1,\qquad b_0=A_\infty,\qquad
 \zeta(h)=\sum_\ell b_\ell(1-\cos(2\pi\ell h)).
\end{equation}
These spatial Fourier weights are different from the coordinate
resampling spectrum $\nu$.
For universal $C<\infty$ and $\alpha>0$,
\begin{equation}\label{eq:spatial-modulus}
 \zeta(h)\le C\|h\|_{\T}^{\alpha},\qquad
 0<\|h\|_{\T}\le1/2.
\end{equation}
In particular $\zeta$ is H\"older continuous with this exponent.

In particular, moving a fixed number of coordinates by shifts
tending to zero in probability, independently of the starting
sample, changes the survivor with probability tending to zero.
The shifts may depend on $N$ and need not be smaller than the
initial spacing $1/N$. Section~\ref{sec:full-tree-rate} supplies a
separate finite-sample rate and a sufficient stability condition for
a growing number of shifted coordinates.
\end{theorem}
\begin{proof}
Under \eqref{eq:survivor-map-coupling}, the finite and limiting
position-disagreement indicators can differ only if one of the
two map evaluations disagrees. Translation preserves Lebesgue
measure, so their probabilities differ by at most twice the
expected map mismatch, uniformly in $h$. Label and position
disagreement differ by at most $1/N$: the exceptional event
requires the inserted coordinate to win in the first copy.
This proves the uniform convergence assertion.

Enumerate the countable limiting vertex values and let $A_j$ be
their measurable origin basins, partitioning $\T$ up to a null set.
Define
\[
 b_\ell=\EE\sum_j\left|
                  \int_{A_j}e^{-2\pi i\ell x}\,dx\right|^2.
\]
Parseval and Tonelli give $\sum_\ell b_\ell=1$, while the basin
second moment gives $b_0=A_\infty$.
The autocorrelation identity for each indicator $\ind_{A_j}$
gives \eqref{eq:spatial-spectrum}; its series converges absolutely
and uniformly. This proves continuity, symmetry, and the value at
zero. Integrating gives $1-b_0=p_\infty$, also the consequence
of using an independent uniform shift in the finite model.

For the modulus, use the depth-$d$ full-circle map and replace
its identity regions by one auxiliary label. The resulting
piecewise constant map has at most $3^{d+1}$ boundaries: there
are at most $3^d$ revealed collapse maps including the final
one, and each adds at most three boundaries to a composition
of identity and constant pieces. Replacing identity regions
does not increase their boundary count. Its disagreement with
$S_\infty$ still has expected mass at most $C\sigma^d$, since
$S_\infty$ has countable range and hence equals the origin
only on a null set.
Translating this finite partition by $h$ changes at most
$3^{d+1}\|h\|_{\T}$ of its mass. Therefore
\[
 \zeta(h)\le C\sigma^d+3^{d+1}\|h\|_{\T}.
\]
Choose $d=\lfloor\log(1/\|h\|_{\T})/\log(3/\sigma)\rfloor$.
This proves \eqref{eq:spatial-modulus} with
$\alpha=\log(1/\sigma)/\log(3/\sigma)>0$.
The triangle inequality for disagreement of three map evaluations
gives $|\zeta(h)-\zeta(h')|\le\zeta(h-h')$, proving the global
H\"older assertion. No numerical value of this exponent is certified.

For a random shift independent of the input, condition on that
shift and use the uniform convergence and continuity at zero.
For finitely many shifts, perform them successively. Conditional
on the shift vector, each intermediate input remains iid uniform.
A union bound proves the label assertion. Position disagreement
can additionally occur only when one of the shifted labels wins,
whose probability is at most the number of shifted labels divided
by $N$.
\end{proof}

\subsection{A reversible Markov rule on complete entrance histories}
\label{sec:history-markov}

\status{Working proof: a reversible Markov rule and absence of a zero spectral atom}
The full history at positive balloon cutoffs retains the insertion maps
that are lost when only one active configuration is observed.
Write $\omega=(C_\omega(r):r\in\mathbb Q\cap(0,1/2])$ for this
history, with locations serving as point identities, and let
$\Pi_{\rm e},\Pi_{\rm o}$ be its even and odd entrance laws.
The state space is a measurable subset of a countable product of
spaces of finite counting measures. Collision times and the vertices
of each cutoff gap tree are measurable from $\omega$.

For each rational $r$, construct the origin map
$\Phi_{\omega,r}$ inside its active gaps as in
Lemma~\ref{lem:origin-map}; if $C_\omega(r)$ is empty, use the full
even survivor map instead. Define the insertion transformation
$\mathcal T_z$ by
\begin{equation}\label{eq:history-insertion-rule}
 C_{\mathcal T_z\omega}(r)=
 \begin{cases}
 C_\omega(r)+\eta_r(C_\omega(r),\Phi_{\omega,r}(z)),
      &C_\omega(r)\ne\varnothing,\\
 \delta_{\Phi_{\omega,r}(z)},&C_\omega(r)=\varnothing.
 \end{cases}
\end{equation}
Here configurations are identified with their counting measures,
and $\eta_r$ is \eqref{eq:cutoff-response}.
All readouts use the same origin $z$. On a null set of
$(\omega,z)$ the transformation can be defined arbitrarily.

\begin{theorem}[The full entrance history is a sufficient Markov state]
\label{thm:history-markov}
For $\Pi_{\rm e}$- or $\Pi_{\rm o}$-almost every $\omega$ and
Lebesgue-almost every $z$, \eqref{eq:history-insertion-rule} is a
consistent entrance history of the opposite parity. The transformation
is jointly measurable and depends only on the displayed state and $z$.
Define
\[
 (Kf)(\omega)=\int_\T f(\mathcal T_z\omega)\,dz,
 \qquad \Pi=\tfrac12(\Pi_{\rm e}+\Pi_{\rm o}).
\]
Then $\Pi_{\rm e}K=\Pi_{\rm o}$, $\Pi_{\rm o}K=\Pi_{\rm e}$,
and $K$ is self-adjoint on $L^2(\Pi)$.
For almost every $z$, $\mathcal T_z$ preserves $\Pi$ and exchanges
its parity sectors. Moreover
\begin{equation}\label{eq:insertion-involutions}
 \mathcal T_z^2\omega=\omega,\qquad
 \mathcal T_z\mathcal T_w\omega=\mathcal T_w\mathcal T_z\omega
\end{equation}
almost everywhere under $\Pi(d\omega)\,dz$ and
$\Pi(d\omega)\,dz\,dw$, respectively.

The limiting refresh-time evolution of the complete odd entrance
history is the stationary reversible Markov process with transition
semigroup
\begin{equation}\label{eq:history-refresh-semigroup}
 P_\lambda=\exp\{\lambda(Q-I)\},\qquad
 Q=K^2\big|_{\rm odd}.
\end{equation}
Thus each rate-one candidate event applies two independent uniform
insertions using \eqref{eq:history-insertion-rule} in succession.
The operator $Q$ is a positive self-adjoint contraction on
$L^2(\Pi_{\rm o})$.
This gives a sufficient state; it does not assert a Markov property
for the survivor or for one fixed-cutoff configuration.
\end{theorem}
\begin{proof}
At a fixed nonempty cutoff, the active configuration and any fixed
depth of all its gap trees converge in total variation along either
parity. The same-origin proof of Lemma~\ref{lem:origin-map} applies
to each gap. Its total expected depth-$d$ map error in the limiting
configuration is at most
\[
 C\sigma^d r^{-1}\sum_i x_i^2\le C\sigma^d/r.
\]
The limsup finite-sample error obeys the same bound by the finite-tree
argument there. Empty even configurations use the two final-pair
trees and their $C\sigma^d$ bound. These approximations also retain
the independent uniform origin $z$. The zero readout-boundary mass
and the exact finite cutoff rule therefore identify the joint limit
of $(\omega_M,z,\omega_{M+z})$ as
$(\omega,z,\mathcal T_z\omega)$ at any finite collection of cutoffs.
The known joint entrance limit supplies consistency. A countable
intersection extends it to all rational cutoffs, and the pointwise
tree construction supplies measurability. It uses the collision
vertices recorded in $\omega$, with no independent tree resampling.

An iid sample with one independent point added is again iid, so the
two identities for $\Pi_{\rm e}K$ and $\Pi_{\rm o}K$ follow.
To prove reversibility, reverse the pair consisting of an even iid
sample and its one-point augmentation. Conditional on the augmented
odd sample, the removed point is a uniformly selected original
label. By Theorem~\ref{thm:edit-universality}, its joint entrance law
with this deletion equals that of an odd sample and one fresh
insertion. Hence
\[
 \Pi_{\rm e}(d\omega)K(\omega,d\omega')
   =\Pi_{\rm o}(d\omega')K(\omega',d\omega).
\]
This identity at all finite collections of rational cutoffs proves
self-adjointness on the complete state space.

We can also retain the location of the reversed deletion. In an iid
sample of size $M+1$, a uniformly tagged point has conditional location
law $\xi_{M+1}$. For bounded continuous $h$,
$\xi_{M+1}(h)\to\int h$ in $L^1$, even when multiplied by any bounded
function of the sample history. Consequently the tag is asymptotically
uniform and independent of the complete limiting history. The
forward graph identification then gives
$(\mathcal T_z)_\#\Pi_{\rm e}=\Pi_{\rm o}$ and the reverse parity
identity for almost every $z$.

For the deletion graph identification with that same origin, combine
the origin-map budget with Lemma~\ref{lem:empirical-seed}.
At a fixed depth, every unresolved origin set is a finite union of
terminal intervals from the side-ancestry and threshold queries.
Replacing its Lebesgue mass by empirical mass changes the bound by
at most $C_{r,d}\Delta_M$; the same is true after intersecting with
an interval used to test the retained origin. Finite-depth map
pieces are constant or the identity on finitely many intervals.
Their empirical graph measures therefore have the same weak limit
as their Lebesgue graph measures. First let $M\to\infty$, then
$d\to\infty$ in the origin-map bound. The tagged deletion, exactly
a coincident insertion, is thus identified by
$\mathcal T_z$ on the reversed history too. The same joint limit is
supported on $\omega'=\mathcal T_z\omega$ and
$\omega=\mathcal T_z\omega'$, proving the involution identity.

For successive independent insertions, retain the fresh origins.
The joint finite-family limit identifies each consecutive pair by
\eqref{eq:history-insertion-rule}, and the next origin remains uniform
and independent of the entire preceding family. Thus the limiting
chain is obtained by successive applications of $K$; no continuity
assumption on $\mathcal T_z$ as a function of $\omega$ is needed.
Using the joint family consisting of neither, either, and both of two
fresh insertions identifies both orders of composition with the same
last history. This proves the commutation identity in
\eqref{eq:insertion-involutions}.

Conditional on finitely many distinct refreshed coordinates, the
successive resampled inputs are a fixed family of cumulative deletions
and insertions. Edit-type universality replaces each such event by two
fresh insertions, preserving the whole joint family. The probability
of a repeated refreshed coordinate tends to zero on bounded refresh
intervals. Mixing over the independent rate-one Poisson candidate
clock gives \eqref{eq:history-refresh-semigroup}, including all rational
balloon cutoffs by a countable-product argument.
Self-adjointness of $K$ gives that of $Q$, and, extending an odd-sector
function by zero on the even sector,
$\langle f,K^2f\rangle_\Pi=\|Kf\|_{L^2(\Pi)}^2\ge0$.
Stationarity and contractivity follow from the parity marginals.
\end{proof}

\begin{corollary}[The zero atom measures conditional atomicity]
\label{cor:zero-atom-invariant}
Let $\mathcal I$ be the invariant sigma-field of $Q$ under
$\Pi_{\rm o}$, and let $\Lambda_\omega$ be the conditional distribution
of the survivor location $U(\omega)$ given $\mathcal I$. Then
\begin{equation}\label{eq:zero-atom-invariant}
 \nu(\{0\})=
 \EE\sum_{x\in\T}\Lambda_\omega(\{x\})^2.
\end{equation}
Only the countably many atoms contribute to the sum. Consequently,
$\nu(\{0\})=0$ is equivalent to almost-sure diffuseness of this
conditional distribution. Ergodicity of $Q$ is sufficient, since
then $\Lambda_\omega$ is the uniform distribution.
Theorem~\ref{thm:no-zero-atom} below proves this diffuseness directly.
Ergodicity of $Q$ remains open.
\end{corollary}
\begin{proof}
For the dyadic partition $\mathcal D_m$ of the circle into $2^m$
half-open arcs, put $f_D(\omega)=\ind_{\{U(\omega)\in D\}}$ and
\[
 a_m(\lambda)=\sum_{D\in\mathcal D_m}
                   \langle f_D,P_\lambda f_D\rangle_{\Pi_{\rm o}}.
\]
For each $\lambda$, this decreases as $m$ increases, with limit
$\PP(U(0)=U(\lambda))$. For each $m$, it also decreases in $\lambda$:
the spectral theorem for the self-adjoint contraction $Q$ represents
each summand as an integral of $e^{-\lambda(1-s)}$ over $0\le s\le1$.
The same theorem gives strong $L^2$ convergence of $P_\lambda$ to the
projection onto the fixed functions of $Q$. These are precisely the
$\mathcal I$-measurable functions: equality in the Markov contraction
forces a fixed function to agree at the two ends of almost every
stationary transition. Therefore
\[
 \lim_{\lambda\to\infty}a_m(\lambda)
       =\EE\sum_{D\in\mathcal D_m}\Lambda_\omega(D)^2.
\]
Both limits are decreasing infima, so they commute. On the left this gives
\[
 \lim_{\lambda\to\infty}\int e^{-\lambda x}\nu(dx)=\nu(\{0\}).
\]
On the right, continuity from above for
$\Lambda_\omega\otimes\Lambda_\omega$ on the decreasing same-cell
events gives its mass on the diagonal, the sum in
\eqref{eq:zero-atom-invariant}. Bounded convergence completes the proof.
\end{proof}

\begin{lemma}[Essential continuity at a protected endpoint]
\label{lem:origin-endpoint}
Almost surely, the limiting origin map $\Phi$ of a shielded interval
$[0,x]$ satisfies
\begin{equation}\label{eq:origin-endpoint}
 \mathop{\rm ess\,lim}_{z\downarrow0}\Phi(z)=0,
 \qquad
 \mathop{\rm ess\,lim}_{z\uparrow x}\Phi(z)=x.
\end{equation}
Here an essential limit means that for every $b>0$ some neighbourhood
of the endpoint has the asserted error at most $b$, outside a
Lebesgue null set. This holds simultaneously in all rational-cutoff
gaps of either entrance history.

There are universal $C<\infty$ and $\beta_{\rm e}>0$ such that,
in a state $(t,x)$, for $0<h\le t$ and $0<b\le x$,
\begin{equation}\label{eq:endpoint-power}
 \PP\left\{\mathop{\rm ess\,sup}_{0<z<h}\Phi(z)>b\right\}
 \le\min\{1,C(x/b)(h/t)^{\beta_{\rm e}}\}.
\end{equation}
The right endpoint has the reflected bound. At an entrance cutoff
$r$, with $0<h\le r/3$, the probability that some active vertex
fails to be removed by almost every insertion in its $h$-neighbourhood
is at most
\begin{equation}\label{eq:uniform-endpoint-removal}
 C r^{-1}(h/r)^{\beta_{\rm e}}.
\end{equation}

Put $V(\omega)=\bigcup_{r\in\mathbb Q\cap(0,1/2]}C_\omega(r)$.
For almost every entrance history,
$V(\mathcal T_z\omega)\subseteq V(\omega)$ for almost every $z$.
Moreover, for every point $v\in\T$ and fixed positive rational $r$,
an insertion whose origin is uniform in $(v-h,v+h)$ satisfies
\begin{equation}\label{eq:targeted-removal}
 \PP_z\{v\in C_{\mathcal T_z\omega}(r)\}\longrightarrow0
       \quad(h\downarrow0).
\end{equation}
The statement permits $v$ to be selected measurably from the history.
\end{lemma}
\begin{proof}
Consider a threshold query $\Phi(z)\le b$ for origins close to the
left endpoint, and follow only the leftmost daughter. At a node of
length $x$ with split $(a,s,c)$, let
$\ell=a-\min(s,a/2)\ge a/2$. For an origin in its left daughter,
the daughter output $y$ lies in $[0,a]$, and the parent map sends
$y\le\ell$ to itself and $y>\ell$ to $a+s$.
Thus if $b\ge a+s$ this query is settled affirmatively on the
entire left daughter, independently of all deeper completions.
Otherwise it recurses there with threshold
$b'=\min(b,\ell)$.

The lengths on this leftmost branch tend to zero almost surely.
Indeed the conditional balanced-split bound used in
Lemma~\ref{lem:insertion-height} gives a uniform positive probability
that every daughter has length at most $3x/4$.
Until the threshold-to-length ratio reaches $1/2$, any unresolved
threshold remains the original $b$; otherwise
$b'\ge a/2$ already. Hence either the query is settled or this ratio
reaches $1/2$ in finitely many steps. Once reached, it stays at least
$1/2$ along every further unresolved left step.

There is a uniform positive probability of $a+s\le x/2$ at any
node. To verify this, put $\theta=(b_-/b_+)^2$.
The event $S\le x/6$ has probability one if $t\le x/6$;
otherwise its probability is
\[
 \frac{w_{x/6}(x)}{w_t(x)}
 \ge (b_-/b_+)e^{-5\kappa}>0.
\]
Conditional on $S=s\le x/6$, the density of $a$ is at least
$\theta/(x-3s)$ on $(s,x-2s)$.
The subinterval $(s,x/2-s)$ has length at least $x/6$, so has
conditional probability at least $\theta/6$.
The resulting uniform positive success probability applies after
every unresolved step. Once the threshold ratio is at least $1/2$,
the query is therefore settled in finitely many further steps almost
surely. Its finite successful left daughter is an origin neighbourhood
on which every deeper completion answers affirmatively.
Apply this to all positive rational $b$ and use reflection for the
right endpoint. This proves \eqref{eq:origin-endpoint}.

For the quantitative assertion, the balanced-split bound gives
$\EE x_j\le x\rho_{\rm l}^j$ along the leftmost branch, for a
fixed $\rho_{\rm l}<1$. If the threshold ratio has not reached
$1/2$ after $j$ unresolved steps, the unchanged threshold $b$ is
less than $x_j/2$. The probability of this event is at most
$x\rho_{\rm l}^j/(2b)$.
After reaching that ratio, each further step succeeds with a uniform
probability at least $\pi>0$, as just proved.
Splitting depth $d$ in half therefore bounds the probability of an
unsettled query by $C(x/b)\rho_{\rm e}^d$, for some
$\rho_{\rm e}<1$.

Put $J=1+K_0>1$ from \eqref{eq:ancestry-ratio-bound}.
For the split separations $S_j$ along this branch,
$\EE S_j^{-1}\le J^j/t$.
The probability any of its first $d$ separations is below $h$ is
at most $C(h/t)J^d$.
If neither event occurs, the successful daughter has length at least
its split separation, hence at least $h$, and certifies the threshold
on all origins in $(0,h)$.
Choose $d=\lfloor\log(t/h)/\log(J/\rho_{\rm e})\rfloor$;
the bounded small depths are absorbed in $C$.
This proves \eqref{eq:endpoint-power} with
$\beta_{\rm e}=\log(1/\rho_{\rm e})/\log(J/\rho_{\rm e})>0$.
Finally apply it conditionally in all cutoff gaps with $t=r$ and
$b=r/3$, for both endpoints. Their lengths sum to one, so a union
bound gives \eqref{eq:uniform-endpoint-removal}. An image within
$r/3$ of an active endpoint has that endpoint as its unique nearest
active point and removes it at cutoff $r$.

There are countably many cutoff gaps. Their limiting maps have atomic
images on internal collision vertices, all in $V(\omega)$; the full
even map has the same property. Taking a countable intersection in
\eqref{eq:history-insertion-rule} proves the subset assertion.
If $v\notin V(\omega)$, it proves \eqref{eq:targeted-removal}
immediately. Otherwise choose a rational $r_0\le r$ strictly before
the removal of $v$. It is a protected endpoint of both adjacent gaps
at $r_0$. By \eqref{eq:origin-endpoint}, every sufficiently close
origin, except a null set, maps sufficiently close to $v$ that its
cutoff response is exactly $-\delta_v$.
Thus $v$ is absent at $r_0$ and remains absent at $r$, since forward
balloon evolution only removes locations. This proves the last claim
simultaneously for all $v\in V(\omega)$; that set is countable.
\end{proof}

\begin{theorem}[No zero spectral atom and growing-resampling disagreement]
\label{thm:no-zero-atom}
The invariant conditional survivor law $\Lambda_\omega$ of
Corollary~\ref{cor:zero-atom-invariant} is almost surely diffuse.
Consequently
\begin{equation}\label{eq:no-zero-atom}
 \nu(\{0\})=0,\qquad p(k)\longrightarrow1,\qquad
 g(\lambda)\longrightarrow1.
\end{equation}
For every sequence $k_N\to\infty$ with $k_N\le N$ and odd $N\to\infty$,
\begin{equation}\label{eq:growing-noise-decorrelation}
 p_N(k_N)\longrightarrow1.
\end{equation}
The same holds for position disagreement and for independent-coordinate
noise whenever its expected number of refreshed coordinates tends to
infinity. This argument gives no decay rate; a logarithmic one follows
in Theorem~\ref{thm:logarithmic-spectrum}. Ergodicity of the full
history kernel remains open.
\end{theorem}
\begin{proof}
We show that every $\mathcal I$-measurable point $X\in\T$ satisfies
$\PP(U=X)=0$. First extend $X$ from odd to even histories so that it
is invariant under one insertion as well as two.
On the odd sector put $F=e^{2\pi iX}$ and on the even sector put
$F=K(e^{2\pi iX})$. Since $Qe^{2\pi iX}=e^{2\pi iX}$,
self-adjointness between parity sectors gives
\[
 \|F\|_{L^2(\Pi_{\rm e})}^2
   =\langle e^{2\pi iX},Qe^{2\pi iX}\rangle_{\Pi_{\rm o}}=1.
\]
But $|F|\le1$, so $|F|=1$ almost surely on the even sector too.
Also $KF=F$. Stationarity and equality in the Markov contraction
give $F(\mathcal T_z\omega)=F(\omega)$ for almost every
$(\omega,z)$. Taking its circle-valued argument supplies the desired
extension of $X$.

For $0<h<1/2$, write
$a_h(z)=(2h)^{-1}\ind_{\{\|z\|_\T<h\}}$.
Start with $\omega\sim\Pi_{\rm o}$ and, conditional on $\omega$,
choose $Z$ with density $a_h(z-X(\omega))$.
Set $\omega'=\mathcal T_Z\omega$.
The measure-preserving involution property and invariance of $X$ imply,
for every bounded measurable $f$,
\begin{align*}
 \EE f(\omega')
 &=\int_\T\!\int a_h(z-X(\omega))f(\mathcal T_z\omega)
                                  \,\Pi_{\rm o}(d\omega)\,dz\\
 &=\int_\T\!\int a_h(z-X(\omega'))f(\omega')
                                  \,\Pi_{\rm e}(d\omega')\,dz
 =\Pi_{\rm e}(f).
\end{align*}
All transformations are valid under this adaptive sampling because
its density relative to $\Pi_{\rm o}(d\omega)\,dz$ is bounded for
each fixed $h$. Also $X(\omega')=X(\omega)$ almost surely.

Now choose an independent uniform $W$ and put
$\omega''=\mathcal T_W\omega'$. Then $\omega''\sim\Pi_{\rm o}$
and $X(\omega'')=X(\omega')$.
Let $B_r(\omega')$ be the total insertion basin mass carried by
vertices outside $C_{\omega'}(r)$ in the even state.
Conditionally on $\omega'$, the last insertion has exactly its
basin distribution. Hence
\begin{equation}\label{eq:adaptive-atom-removal}
 \PP(U=X)=\PP(U(\omega'')=X(\omega''))
 \le \PP\{X(\omega)\in C_{\omega'}(r)\}
            +\EE_{\Pi_{\rm e}} B_r.
\end{equation}
For fixed positive rational $r$, the first term tends to zero as
$h\downarrow0$ by Lemma~\ref{lem:origin-endpoint} and bounded
convergence, including history-dependent choices of $X$.
The second term tends to zero as $r\downarrow0$ by atomic total-mass
localisation; Corollary~\ref{cor:early-total-mass} gives the stronger
bound $C r^\beta$. Thus $\PP(U=X)=0$.

If $\Lambda_\omega$ had atoms with positive probability, select a
largest atom, breaking ties by its representative in $[0,1)$, and
use any fixed point on the event of no atoms. This is a measurable
function of a probability measure, for example by nested dyadic
intervals. It gives an $\mathcal I$-measurable $X$ with
$\PP(U=X)=\EE\Lambda_\omega(\{X\})>0$, a contradiction.
Equation~\eqref{eq:no-zero-atom} now follows from
Corollary~\ref{cor:zero-atom-invariant} and the spectral formulas.

For fixed $k$, eventually $k_N\ge k$, so monotonicity gives
$\liminf_Np_N(k_N)\ge p(k)$. Let $k\to\infty$.
Position disagreement is at least label disagreement.
For independent-coordinate noise the number $J_N$ of refreshed
coordinates tends to infinity in probability when
$\EE J_N\to\infty$ (its variance is at most its mean).
For each fixed $k$, its disagreement probability is at least
$p_N(k)\PP(J_N\ge k)$. The same two limits complete the proof.
\end{proof}

\begin{corollary}[Disjoint active configurations after growing resampling]
\label{cor:late-disjointness}
Fix a positive balloon cutoff $r$. For odd $N\to\infty$ and
$k_N\to\infty$, $k_N\le N$, the original and $k_N$-resampled active
configurations have no common location at cutoff $r$ with probability
tending to one. Their expected number of common active coordinate
labels also tends to zero. Thus their future histories after $r$ have
no shared vertex with probability tending to one.

For the stationary limiting history process,
\[
 \EE|C_{\omega(0)}(r)\cap C_{\omega(\lambda)}(r)|
       \longrightarrow0\qquad(\lambda\to\infty),
\]
and the same holds at integer steps of $Q$.
The cutoff is fixed: no simultaneous assertion for $r=r_N\downarrow0$
or independence of the two configurations is proved.
\end{corollary}
\begin{proof}
For the even history, let $G_r=(\Phi_{\omega,r})_\#\Leb$ be the
origin-image probability measure at cutoff $r$, with the full-map
convention in an empty state. It is supported on $V(\omega)$.
Thus, for each fixed $r$,
\[
 \EE_{\Pi_{\rm e}}G_r(V(\omega)\setminus C_\omega(\delta))
       \longrightarrow0\qquad(\delta\downarrow0).
\]
Repeat the adaptive insertion proof with any invariant point $X$.
If $X$ is absent at cutoff $\delta<r$ in the intermediate even state,
its occurrence at cutoff $r$ after the last insertion requires a
positive response with image exactly $X$.
Consequently
\[
 \PP_{\Pi_{\rm o}}\{X\in C_\omega(r)\}
 \le \PP\{X\in C_{\omega'}(\delta)\}
       +\EE_{\Pi_{\rm e}}G_r(V\setminus C(\delta)).
\]
Send the adaptive neighbourhood radius to zero, then
$\delta\downarrow0$. The left side is zero. Measurable selection of
an atom now shows that the conditional expected counting measure
$\Gamma_r=\EE[C_\omega(r)\mid\mathcal I]$ is almost surely diffuse.

Use $f_D(\omega)=C_\omega(r)(D)$ in the dyadic-partition argument of
Corollary~\ref{cor:zero-atom-invariant}.
Counts are bounded by $\max(1,\lfloor1/r\rfloor)$, so all squared
norms and limiting diagonal masses are integrable. The limiting
expected overlap is $\EE\sum_x\Gamma_r(\{x\})^2=0$.
The argument also works for $Q^k$: its spectrum is nonnegative, so
each dyadic same-cell inner product decreases in $k$.

For finite $N$ let $A_N(r,k)$ be the expected common active-label
count after $k$ resamplings. Apply the product-space decomposition to
the vector $(\ind_{\{i\ {\rm active\ at}\ r\}})_{i=1}^N$.
It expresses $A_N(r,k)$ as a nonnegative sum of the avoidance factors
in \eqref{eq:k-noise-spectrum}; hence it is nonincreasing in $k$.
For fixed $k$, common-label and common-location counts differ in
expectation by at most
$k\,\EE K_N(r)/N$, which tends to zero.
The coupled limit and bounded active counts give
\[
 A_N(r,k)\longrightarrow
 a(r,k):=\int |C_\omega(r)\cap C_{\omega'}(r)|
                 \,\Pi_{\rm o}(d\omega)Q^k(\omega,d\omega').
\]
Eventually $k_N\ge k$, so
$\limsup_N A_N(r,k_N)\le a(r,k)$; now let $k\to\infty$.
Location overlap is no larger than label overlap, and its count is
integer valued, proving the probability statement. Forward evolution
only deletes locations, which proves the future-history assertion.
\end{proof}

\begin{corollary}[All growing mixtures of random additions and deletions]
\label{cor:growing-mixed-noise}
Let $N$ and $N'=N-d+a$ be positive odd integers, with
$0\le d\le N$ and $a\ge0$. Delete $d$ uniformly selected original
points from an iid sample of size $N$, and add $a$ independent uniform
points. Give fresh points fresh labels. The survivor agreement
probability obeys the finite bound
\begin{equation}\label{eq:cross-size-agreement}
 \PP(U_N=U_{N'})
       \le\sqrt{(1-p_N(d))(1-p_{N'}(a))}.
\end{equation}
Here equality of locations is the same as equality of their common
original label, and $p_n(0)=0$.

As $d+a\to\infty$, this agreement tends to zero, uniformly over all
compatible odd sizes. At every fixed positive cutoff the probability
that the two active configurations have any common location also tends
to zero. In particular these assertions cover growing pure addition
and pure deletion counts, while keeping both sample sizes odd.
\end{corollary}
\begin{proof}
Condition on the common background $B$ of size $N-d=N'-a$.
Represent each survivor by a unit coordinate vector, $F$ and $F'$,
in the union of both label sets. Given $B$, the two extra point sets
are independent. Hence
\[
 \PP(U_N=U_{N'})
 =\EE\langle\EE[F\mid B],\EE[F'\mid B]\rangle
 \le\sqrt{\EE\|\EE[F\mid B]\|^2\,
             \EE\|\EE[F'\mid B]\|^2}.
\]
Two conditionally independent copies of $F$ retain exactly $B$ and
resample its $d$ other coordinates, using the same coordinate labels.
The first squared norm therefore has expectation $1-p_N(d)$;
the second has expectation $1-p_{N'}(a)$. This proves
\eqref{eq:cross-size-agreement}.
If $d+a\to\infty$, every subsequence has a further subsequence
with either $d\to\infty$ or $a\to\infty$.
The corresponding sample size also tends to infinity, and
Theorem~\ref{thm:no-zero-atom} makes that factor tend to zero.
This sequential argument also proves the stated uniformity.

For active configurations use instead the vectors of active-label
indicators. If $A_n(r,j)$ denotes the finite expected common active-label
count under $j$ resamplings, the same conditional Cauchy--Schwarz
argument bounds the mixed expected overlap by
\[
 \sqrt{A_N(r,d)A_{N'}(r,a)}.
\]
Each factor is bounded by $\max(1,\lfloor1/r\rfloor)$, and the
factor corresponding to a growing edit count tends to zero by
Corollary~\ref{cor:late-disjointness}. Markov's inequality gives
the asserted disjointness probability.
\end{proof}

\subsection{A finite-sample rate from the complete last-pair trees}
\label{sec:full-tree-rate}

\status{Working proof: a polynomial rate, with uncalibrated constants}
The complete origin map can be approximated using the two final-pair
arc trees. This avoids introducing an early-label cutoff before
comparing the finite and limiting survivor maps.

\begin{lemma}[A quantitative full-root geometry coupling]
\label{lem:full-root-rate}
For even $M\ge2$, let $\tau_M^{\rm last}(d)$ be the total-variation
distance between the finite and limiting last pair together with the
first $d$ generations in both arc trees. Include premature empty
daughters as separate finite-sample states. Put
\[
 J=1+K_0,\qquad C_0=2\EE L^{-1}=2(1+2e^\gamma),
 \qquad q=2\cdot3^d,
\]
where $K_0$ is from \eqref{eq:ancestry-ratio-bound}.
For $d\ge0$ and $0<\eta<1/2$,
\begin{equation}\label{eq:full-root-TV}
 \tau_M^{\rm last}(d)\le\min\left\{1,
 C_0\eta(3J)^d
 +\frac{(\eta^{-1}e^{1/\eta})^q}{(q-1)!}(1-2\eta)^M\right\}.
\end{equation}
\end{lemma}
\begin{proof}
Orient the last pair by its shorter arc and include its uniform
rotation. With both arc trees fully revealed to depth $d$, the
remaining intervals number $q$ and their lengths sum to one.
The same last-pair decomposition used in
Proposition~\ref{prop:half-noise} gives finite density
\[
 M![\lambda^M]e^\lambda
                  \prod_{\ell=1}^{q}w_{t_\ell}^\lambda(x_\ell)
\]
and limiting density $2\prod_\ell w_{t_\ell}(x_\ell)$, against
$dr$ and all revealed $da\,ds$. There is no cyclic-root factor
here: the final short-arc endpoint determines the root. At depth zero
this is the normalized last-pair density; revealing a split replaces
one interval factor by its three daughter factors, with unit Jacobian.

For every limiting node, \eqref{eq:ancestry-ratio-bound} gives
$\EE(S^{-1}\mid t,x)\le J/t$. All three daughter cutoffs equal
$S$. Hence, conditional on $L$, the expected sum of reciprocal
terminal cutoffs in the two depth-$d$ trees is at most
$2(3J)^d/L$. Markov's inequality gives
\[
 \PP\{\min_\ell t_\ell<\eta\}\le C_0\eta(3J)^d.
\]
This controls very small cutoffs jointly, without imposing a fixed
lower bound on every split ratio. On the remaining geometry,
every terminal cutoff is at least $\eta$.
The pointwise coefficient estimate from
Corollary~\ref{cor:tree-rate} bounds the density difference there by
$(\eta^{-1}e^{1/\eta})^q(1-2\eta)^M$.
Terminal lengths are unit-Jacobian simplex coordinates for the
$q-1$ independent split variables; the constrained region has volume
at most $1/(q-1)!$.

Both full geometry laws have total mass one. Their total variation is
at most the limiting bad-event mass plus the integrated density
difference on the good geometry. This also controls all finite
premature-empty states, without a separate conditioning argument.
The common independent rotation can be coupled identically, giving
\eqref{eq:full-root-TV} for physical locations as well.
\end{proof}

\begin{theorem}[A polynomial rate for the full map and noise profiles]
\label{thm:polynomial-noise-rate}
Put
\[
 \chi=\frac{\log(1/\sigma)}{4\log(3J/\sigma)}\in(0,1/4),
\]
where $\sigma$ is from Lemma~\ref{lem:origin-map} and $J=1+K_0$.
There are universal $C<\infty$ and $N_0$ such that finite even backgrounds
of size $M=N-1$ can be coupled to the limiting survivor map with
\begin{equation}\label{eq:full-map-rate}
 \EE\Leb\{z:S_{B_M}(z)\ne S_\infty(z)\}
 \le C N^{-\chi},
 \qquad N\ge N_0\ \hbox{odd}.
\end{equation}
Consequently, after changing $C$ if necessary,
\begin{equation}\label{eq:noise-polynomial-rate}
 \max\left\{|p_N-p_\infty|,
          \sup_{h\in\T}|\zeta_N(h)-\zeta(h)|\right\}
 \le C N^{-\chi}.
\end{equation}
The prefactor, exponent and threshold are not numerically calibrated. This is
an asymptotic rate, not a usable numerical error bar at a specified
sample size and not the much faster scalar last-pair TV rate.
\end{theorem}
\begin{proof}
The proof of Lemma~\ref{lem:origin-map} supplies a deterministic
depth-$d$ origin budget, capped at one, bounding the disagreement
between any two deeper completions of the same full-root geometry.
It consists of the side-only mass at depth $\lfloor d/2\rfloor$
and the unresolved midpoint-query origin intervals for earlier
middle choices. Its limiting expectation is at most $C\sigma^d$:
the two root ratios are $1$ and $(1-L)/L$, with
$\EE(1/L)<\infty$.

Maximally couple the two depth-$d$ geometries. On their agreement
event, the deterministic budget compares the true finite map directly
with the limiting map, irrespective of their conditional deeper laws.
On a geometry mismatch the disagreement mass is at most one.
Thus the expected full-map error is at most
\begin{equation}\label{eq:full-map-parameter-error}
 C\sigma^d+\tau_M^{\rm last}(d).
\end{equation}
In particular this argument uses a bounded function of the revealed
geometry, not a nonuniform finite-sample contraction assumption.

For large $M$, set
\[
 d=\left\lfloor\frac{\log M}{4\log(3J/\sigma)}\right\rfloor,
 \qquad \eta=M^{-1/4}.
\]
Both the origin budget $\sigma^d$ and the limiting bad-geometry
mass $\eta(3J)^d$ are at most $C M^{-\chi}$.
Moreover $q=2\cdot3^d\le2M^{1/4}$.
Ignoring the factorial only enlarges the coefficient error in
\eqref{eq:full-root-TV}. Its logarithm is at most
\[
 2M^{1/4}\bigl(\tfrac14\log M+M^{1/4}\bigr)-2M^{3/4}
       \le -M^{2/3}
\]
for all sufficiently large $M$.
Substitution in \eqref{eq:full-map-parameter-error} proves
\eqref{eq:full-map-rate}.

Under this coupling, either two independent uniform probes or the
two probes $U,U+h$ have a disagreement indicator mismatch only if
one of their map evaluations differs. This costs at most twice the
expected map mismatch, uniformly in $h$.
The finite label/position correction is at most $1/N$ in both cases.
Absorb that term into \eqref{eq:full-map-rate} to obtain
\eqref{eq:noise-polynomial-rate}.
\end{proof}

\begin{corollary}[A sufficient growing-shift stability scale]
\label{cor:growing-spatial-stability}
Shift $q_N$ prescribed coordinates by a random vector
$(H_{N,1},\ldots,H_{N,q_N})$ independent of the iid starting sample.
The shifts may be dependent on each other. With $\alpha>0$ from
Theorem~\ref{thm:spatial-noise} and $\chi>0$ from
Theorem~\ref{thm:polynomial-noise-rate}, the survivor-label
disagreement probability is at most
\[
 C\sum_{j=1}^{q_N}\EE\|H_{N,j}\|_\T^\alpha
       +Cq_N N^{-\chi}.
\]
In particular it tends to zero when both terms do.
The position-disagreement probability has the additional upper error
$q_N/N$, which also vanishes under these conditions.
\end{corollary}
\begin{proof}
Perform the shifts successively. Conditional on the shift vector,
each intermediate sample is still iid uniform. Apply the uniform
spatial profile rate, its H\"older modulus, and a union bound.
The only extra position changes with unchanged label require a
shifted label to win, of probability at most $q_N/N$.
Finally $q_N N^{-\chi}\to0$ implies $q_N/N\to0$, since $\chi<1$.
\end{proof}

\begin{corollary}[Quantitative basin and deletion-mass convergence]
\label{cor:basin-rate}
With $\chi$ from Theorem~\ref{thm:polynomial-noise-rate} and
$\beta=4\chi$, for every sufficiently large even $M$,
\begin{equation}\label{eq:deletion-basin-rate}
 \EE\sum_{v\in B_M}|d_v(B_M)/M-\theta_v(B_M)|\le C M^{-\chi}.
\end{equation}
The ranked insertion and normalized deletion vectors can each be
coupled to the limiting mass partition with expected $\ell^1$
distance at most $C M^{-\chi}$. Moreover, uniformly for
$0<\varepsilon<1/2$,
\begin{equation}\label{eq:finite-early-mass}
 \EE\sum_{v\notin C_M(\varepsilon)}
       \left(\theta_v(B_M)+d_v(B_M)/M\right)
       \le C\left(\varepsilon^\beta+M^{-\chi}\right).
\end{equation}
In particular the total mass on labels removed before any vanishing
cutoff tends to zero without an order of limits.
\end{corollary}
\begin{proof}
Use the empirical interval discrepancy $\Delta_M$ from
Lemma~\ref{lem:empirical-seed}. It satisfies
$\EE\Delta_M\le C M^{-1/2}$: symmetrization bounds the expected
supremum of the centered empirical distribution function by twice
the expected maximum of a length-$M$ Rademacher walk, divided by
$M$. The $L^2$ maximal inequality bounds that maximum by
$2\sqrt M$. A circular interval is a difference of two distribution
function values; the additional $1/M$ handles endpoints.

On a shared depth-$d$ full-root geometry, the deterministic unresolved
origin set is a union of at most $C3^d$ intervals. Outside it, every
deeper completion has the same constant output on at most $C3^d$
pieces: the origin has entered a middle daughter and its endpoint
query has settled. Assign the unresolved set to one revealed vertex.
The resulting finite-valued map approximates the true map with error
at most its unresolved origin mass. For empirical input mass that
error is at most the Lebesgue error plus $C3^d\Delta_M$;
the same discrepancy bounds the total variation between the two
approximate pushforwards. Revealed endpoints cost at most
$C3^d/M$, already included in this estimate. No independence of
the empirical measure and the revealed geometry is used.

Maximal geometry coupling and the deterministic budget therefore give
\[
 \EE\sum_{v\in B_M}|d_v/M-\theta_v|
 \le C\left(\sigma^d+\tau_M^{\rm last}(d)
                         +3^d\EE\Delta_M\right).
\]
Here the empirical full pushforward gives exactly $d_v/M$; the
Lebesgue pushforward gives the insertion masses and a possible
identity part, whose total mass is also covered by the unresolved
budget. Take the depth and $\eta=M^{-1/4}$ from
Theorem~\ref{thm:polynomial-noise-rate}. Then $3^d\le M^{1/4}$,
so the extra empirical error is at most $CM^{-1/4}$.
This proves \eqref{eq:deletion-basin-rate}. The full-map coupling
bounds total variation of its pushforward measures by twice its
origin mismatch. Decreasing rearrangement contracts $\ell^1$
distance, proving the two ranked-vector assertions.

For the early-mass assertion, use instead
\[
 d=\left\lfloor
 \frac{\min\{\log(1/\varepsilon),\tfrac14\log M\}}
      {\log(3J/\sigma)}\right\rfloor,
 \qquad \eta=M^{-1/4}.
\]
The finite and limiting geometry error is still at most
$CM^{-\chi}$, and
$\sigma^d\le C(\varepsilon^\beta+M^{-\chi})$.
The probability a revealed vertex has removal separation at most
$\varepsilon$ is at most $C\varepsilon(3J)^d\le C\varepsilon^\beta$.
Outside this event and the unresolved origin set, the true map
therefore takes values only at vertices in $C_M(\varepsilon)$.
This bounds the insertion tail. Add
\eqref{eq:deletion-basin-rate} for the deletion tail to obtain
\eqref{eq:finite-early-mass}.
\end{proof}

\subsection{Logarithmic spectral tails and shrinking-cutoff overlap}
\label{sec:spectral-rate}

\status{Working proof: quantitative decay for the limiting refresh process}
The endpoint estimate can be combined with spectral localization.
The argument uses finite-time survivor probabilities to approximate
a low-energy projection; it does not assume that projection is positive.

\begin{lemma}[Backward retention of tagged late points]
\label{lem:backward-retention}
Let a stationary insertion chain alternate between the two entrance
parities. If its terminal state is odd and $X$ is its survivor, then
after reversing $j$ insertion steps,
\begin{equation}\label{eq:backward-retention}
 \PP\{X\notin C_{\rm initial}(r)\}\le Cj\sqrt r.
\end{equation}
More generally, if the terminal tags are all its active points at a
cutoff $r_0\ge r$, their expected number missing from the initial
$r$-configuration is at most $Cj\sqrt r/r_0$.
\end{lemma}
\begin{proof}
Put $G_r=\max_i x_i$ for the largest active gap, and put $G_r=0$
on the empty even state. The two parity gap-square bounds in the
proof of Lemma~\ref{lem:separated-insertions} give
$\EE G_r\le(\EE\sum_i x_i^2)^{1/2}\le C\sqrt r$.
A uniform insertion can remove a specified active point only when
its origin lies in one of that point's adjacent gaps. Conditional on
the entire history, this has probability at most $2G_r$.

Each fixed-origin transformation preserves the parity laws and is an
involution. Thus the reversed stationary insertion chain again uses
independent uniform origins. Starting with a tagged terminal point,
bound the probability of its first removal at each reversed step by
$2G_r$ and sum expectations. Every intermediate state has the
stationary law of its parity, proving \eqref{eq:backward-retention}.
For the second assertion there are at most $1/r_0$ terminal tags,
all active initially in the reversed chain at cutoff $r$. Sum the
same conditional first-removal bounds over these tags and the $j$
steps. No independence of the tags and the history is needed.
\end{proof}

\begin{theorem}[Logarithmic spectrum and overlap bounds]
\label{thm:logarithmic-spectrum}
There are universal $C<\infty$ and $a>0$ such that
\begin{equation}\label{eq:log-spectral-tail}
 \nu([0,x])\le \frac{C}{[\log(e/x)]^a},\qquad 0<x<1,
\end{equation}
and consequently
\begin{equation}\label{eq:log-noise-decay}
 1-p(k)\le\frac{C}{[\log(k+2)]^a},\qquad
 1-g(\lambda)\le\frac{C}{[\log(\lambda+2)]^a}.
\end{equation}
One admissible exponent is
\[
 a=\frac{\log(10/9)}{\log(10+10/\beta_{\rm e})}>0,
\]
where $\beta_{\rm e}$ is from Lemma~\ref{lem:origin-endpoint}.
It is not numerically calibrated.

For the stationary odd chain with kernel $Q$, write
$\mathcal A_r(k)=\EE|C_{\omega_0}(r)\cap C_{\omega_k}(r)|$.
Uniformly for $0<r\le1/3$,
\begin{equation}\label{eq:shrinking-cutoff-overlap}
 \mathcal A_r(k)\le\frac C r
 \min\left\{1,
       \left(\frac{\log(e/r)}{\log(k+2)}\right)^a\right\}.
\end{equation}
The same bound holds for the continuous refresh process with $k$
replaced by $\lambda$. In particular, if $r_k\downarrow0$ and
\[
 r_k^{-1}\left(\frac{\log(e/r_k)}{\log(k+2)}\right)^a\longrightarrow0,
\]
the limiting configurations at $r_k$ are disjoint with probability
tending to one. For example $r_k=(\log(k+2))^{-b}$ works for every
$0<b<a$ once $r_k\le1/3$.
These assertions concern the limiting entrance process.
Section~\ref{sec:finite-family-rate} separately supplies a uniform
finite-$N$ transfer for growing $k$.
\end{theorem}
\begin{proof}
First consider the survivor. For a finite dyadic partition
$\mathcal D_m$, put $f_D=\ind_{\{U\in D\}}$ on the odd sector and
define the probability measure
\[
 \nu_m(B)=\sum_{D\in\mathcal D_m}
   \langle f_D,\ind_B(I-Q)f_D\rangle_{\Pi_{\rm o}}.
\]
The notation $\ind_B(I-Q)$ denotes a spectral projection.
The Laplace transform of $\nu_m$ is the same-cell probability
$a_m(\lambda)$ used in Corollary~\ref{cor:zero-atom-invariant}.
Hence $\nu_m$ converges weakly to $\nu$.

Fix $0<x<b<1$ and an integer $n\ge1$. In the direct sum over cells,
write
\begin{gather*}
 h_D=\ind_{[0,x]}(I-Q)f_D,\qquad g_D=Kh_D,
 \qquad u_D=Q^n f_D,\qquad v_D=Ku_D,\\
 \mu=\sum_D\|h_D\|_2^2=\nu_m([0,x]).
\end{gather*}
The vectors $h,u$ live on odd histories and $g,v$ on even histories.
The last two are nonnegative and each has coordinates summing to one.
Put $E=\sum_D\|h_D-u_D\|_2^2$. The spectral theorem gives
\begin{equation}\label{eq:spectral-window-error}
 E\le (nx)^2\mu+\nu_m([0,b])-\mu+e^{-2nb}.
\end{equation}
Indeed $1-(1-y)^n\le ny$ for $y\le x$, and on $y>x$ split the
integral at $b$. Also $\|g-v\|_2\le\sqrt E$ and
$\|u\|_2\le\sqrt\mu+\sqrt E$.

Choose an origin radius $h_0<r/6$, and make the cells shorter than
$h_0$. For the center $c_D$ of $D$, let
\[
 a_D(z)=\frac1{2h_0}\ind_{\{\|z-c_D\|_\T<h_0\}},
 \qquad
 (L_D w)(\omega)=\int a_D(z)w(\mathcal T_z\omega)\,dz.
\]
Each $L_D$ is a contraction between the parity $L^2$ spaces, since
each $\mathcal T_z$ preserves the respective laws. Let
$S(w,w')=\sum_D\langle w_D,L_Dw'_D\rangle_{\Pi_{\rm o}}$.
For the low-energy vectors,
\[
 \sum_D\int\EE_{\Pi_{\rm o}}
       |h_D(\omega)-g_D(\mathcal T_z\omega)|^2\,dz
   =\mu-\|g\|_2^2\le x\mu.
\]
Because $a_D\le1/(2h_0)$, expanding the corresponding locally
weighted squares yields
\begin{equation}\label{eq:local-low-energy}
 S(h,g)\ge\mu\left(1-\frac x2-\frac{x}{4h_0}\right).
\end{equation}
On the other hand, blockwise contraction and the preceding norm bounds
give
\begin{equation}\label{eq:local-projection-error}
 |S(h,g)-S(u,v)|\le2\sqrt{\mu E}+E.
\end{equation}
In particular, the projection error has no factor equal to the
number of cells or to $1/h_0$.

Write $A_D(\omega)=\ind_{\{C_\omega(r)\cap D\ne\varnothing\}}$.
Lemma~\ref{lem:backward-retention}, applied to $2n$ and $2n+1$
insertion steps, respectively, gives
\[
 \sum_D\EE_{\Pi_{\rm o}}u_D^2(1-A_D)
  +\sum_D\EE_{\Pi_{\rm e}}v_D^2(1-A_D)
       \le C(n+1)\sqrt r.
\]
Here $u_D^2\le u_D$ and $v_D^2\le v_D$.
By \eqref{eq:uniform-endpoint-removal}, outside an odd-state event
of probability at most $Cr^{-1}(h_0/r)^{\beta_{\rm e}}$, every
active vertex is removed by almost every insertion within distance
$2h_0$. If $A_D(\omega)=1$ and $z$ belongs to the support of $a_D$,
this removes the unique active point in $D$; the cutoff response is
exactly its deletion, so $A_D(\mathcal T_z\omega)=0$.

Thus in the integral defining $S(u,v)$ at least one of the two
active-cell indicators is zero, unless that exceptional event occurs.
Cauchy--Schwarz bounds the two missing-cell contributions by
$C\sqrt{(n+1)\sqrt r}$. For the exceptional event, use
$\sum_Du_D^2\le1$ and
$\sum_D\int a_D(z)\EE v_D(\mathcal T_z\omega)^2\,dz=\|v\|_2^2\le1$.
Consequently
\begin{equation}\label{eq:local-positive-small}
 0\le S(u,v)\le
 C\sqrt{(n+1)\sqrt r}
   +C\sqrt{r^{-1}(h_0/r)^{\beta_{\rm e}}}.
\end{equation}
All constants are independent of the partition mesh.

Put $B=8+10/\beta_{\rm e}$ and $A=B+2$. For sufficiently small
$b$, take
\[
 n=\lceil b^{-2}\rceil,\qquad r=b^8,
 \qquad h_0=b^B,\qquad x=b^A.
\]
The right side of \eqref{eq:local-positive-small} is at most $Cb$,
while $x/h_0=b^2$ and $(nx)^2\to0$.
Combining \eqref{eq:local-low-energy}--\eqref{eq:local-positive-small}
and using $2\sqrt{\mu E}\le\mu/2+2E$ gives, after reducing the
universal upper bound on $b$ if necessary,
$\mu\le7E+Cb$. Equation~\eqref{eq:spectral-window-error} then implies
\begin{equation}\label{eq:spectral-scale-contraction}
 \nu_m([0,b^A])\le\frac9{10}\nu_m([0,b])+Cb.
\end{equation}
For example it suffices in this last step to have
$(nx)^2\le1/36$; then $7/(8-7(nx)^2)<9/10$.
First send the mesh to zero at continuity points of $\nu$.
Approximation from above extends \eqref{eq:spectral-scale-contraction}
to every sufficiently small $b$, with $\nu$ in place of $\nu_m$.

Iterate at $b_j=b_0^{A^j}$. The error terms $Cb_j$ are summable
even after multiplication by $(10/9)^j$, so
$\nu([0,b_j])\le C(9/10)^j$.
Monotonicity between these scales proves
\eqref{eq:log-spectral-tail} with $a=\log(10/9)/\log A$.
Splitting the integrals for $1-p(k)$ and $1-g(\lambda)$ at
$k^{-1/2}$ and $\lambda^{-1/2}$ proves
\eqref{eq:log-noise-decay}; bounded small parameters are absorbed in $C$.

For the configuration assertion fix $r_0\le1/3$, put $R=1/r_0$,
and replace $f_D$ by $R^{-1/2}C_\omega(r_0)(D)$, using cells shorter
than $r_0$. Each unnormalized coordinate is at most one and their
sum is at most $R$. Accordingly the vectors $u,v$ still satisfy
$\sum_Du_D^2,\sum_Dv_D^2\le1$, and the second assertion of
Lemma~\ref{lem:backward-retention} gives the same normalized
missing-cell bound whenever $r\le r_0$.
All estimates above therefore hold unchanged, provided $b^8\le r_0$.
The spectral measures now have total mass at most one. Their weak
limit satisfies $\int(1-x)^k\,\nu_{r_0}(dx)=\mathcal A_{r_0}(k)/R$:
refining the cells leaves exactly the common spatial points.
Compactness and uniqueness of these polynomial moments on $[0,1]$
justify this weak limit.

Start the iteration at
$b_0=\min\{b_*,r_0^{1/8}\}$, with fixed sufficiently small $b_*>0$.
The constants in the iteration remain uniform in $r_0$ and
$\log(1/b_0)\le C\log(e/r_0)$. Thus its spectral tail is at most
\[
 C\min\left\{1,
     \left(\frac{\log(e/r_0)}{\log(e/x)}\right)^a\right\}.
\]
Split its discrete or continuous time transform at the inverse
square root of the time parameter and multiply by $R$.
This proves \eqref{eq:shrinking-cutoff-overlap}. Markov's inequality
gives the disjointness statements. No finite-$N$ uniformity is used.
\end{proof}

\subsection{Uniform finite-sample noise and overlap rates}
\label{sec:finite-family-rate}

\status{Working proof: quantitative transfer from entrance histories}
The origin budgets also apply to all gaps at a cutoff that decreases
with sample size. Together with separated insertions, this gives a
finite-family coupling whose error is explicit in the number of probes.

\begin{lemma}[A polynomial coupling for finite insertion families]
\label{lem:finite-family-rate}
Put
\[
 \chi_{\rm f}=\frac{\chi}{1+8\chi}
   =\frac{\log(1/\sigma)}
          {4\{\log(3J)+3\log(1/\sigma)\}}>0.
\]
For every sufficiently large $M$, an iid background of size $M$ of
either parity and $q$ independent
uniform probes, the family obtained from every subset of the probes
can be coupled to its canonical entrance insertion family so that
all spatial configurations and future histories agree after cutoff
$\varepsilon_M=M^{-\chi_{\rm f}}$, except on an event of probability
at most
\begin{equation}\label{eq:finite-family-rate}
 \min\{1,C(q+1)^2 M^{-\chi_{\rm f}}\}.
\end{equation}
Agreement retains the coincidence relation between points in different
copies. This is not a comparison of the original integer coordinate
labels or of the entire histories down to cutoff zero.
\end{lemma}
\begin{proof}
At cutoff $\varepsilon$, reveal $d$ generations of every nonempty
background gap tree. With $m$ root gaps there are $m3^d$ terminal
intervals. The fixed-count density is \eqref{eq:tree-coefficient},
with cyclic-root factor $1/m$; the same formula holds in the odd
sector by Proposition~\ref{prop:reverse-process}. The full and empty
subsets in coefficient extraction give the limiting factor two,
because $m3^d$ has the same parity as $M$.

All root cutoffs equal $\varepsilon$, and $m\le1/\varepsilon$.
Iterating $\EE(S^{-1}\mid t,x)\le J/t$ therefore gives
\[
 \EE\sum_{\text{terminal }i}t_i^{-1}
       \le \varepsilon^{-2}(3J)^d.
\]
The limiting mass of geometries with some $t_i<\eta$ is at most
$\eta\varepsilon^{-2}(3J)^d$. On the complementary geometry,
integrate the coefficient bound from Corollary~\ref{cor:tree-rate}.
The sector with $m$ roots has reference volume at most
$1/[m(m3^d-1)!]$. Bounding each such volume by one and summing at
most $1/\varepsilon$ sectors gives the error
\[
 \varepsilon^{-1}\exp\left\{
   \varepsilon^{-1}3^d\bigl(\log(1/\eta)+1/\eta\bigr)-2\eta M
                         \right\}.
\]
Discard the limiting empty even state at cost at most
$7e^{-(1/\varepsilon-1)/4}$; odd states are nonempty.
Normalization then controls all finite premature-empty states too.
Thus the total-variation geometry error is bounded by the sum of
these three quantities.

Write $\ell=\log(1/\sigma)$ and $L_0=\log(3J)$, and take
\[
 \varepsilon=M^{-\chi_{\rm f}},\qquad \eta=M^{-1/4},
 \qquad d=\left\lfloor\frac{\log M}{4(L_0+3\ell)}\right\rfloor.
\]
Then
\[
 \eta\varepsilon^{-2}(3J)^d\le M^{-\chi_{\rm f}},
 \qquad \sigma^d\le C M^{-\chi_{\rm f}},
 \qquad \varepsilon^{-1}3^d\le M^{1/4}.
\]
The coefficient error is consequently at most $e^{-M^{2/3}}$ for
all sufficiently large $M$. The geometry coupling fails with
probability at most $CM^{-\chi_{\rm f}}$.

On shared geometry the deterministic origin budgets compare any
two deeper completions. Their limiting expected total over the
cutoff gaps is at most
\[
 C\sigma^d\varepsilon^{-1}\EE\sum_i x_i^2\le C\sigma^d,
\]
using the gap-square bounds for either parity. The probability that
any of the $q$ common probes has different finite and limiting
cutoff images is therefore at most $Cq\sigma^d$, in addition to
the geometry error. This uses a deterministic budget on the shared
geometry and no uniform finite-count contraction assumption.

The limiting probability that the probes fail to occupy distinct,
nonadjacent gaps, or that the background is empty, is at most
$C(q+1)^2\varepsilon$, by Lemma~\ref{lem:separated-insertions}.
On the common good event the finite gap configuration is the same,
all probe images coincide, and the separated-insertion readout
holds simultaneously for every subset. Thus every copy has exactly
the same spatial configuration in the two families at cutoff
$\varepsilon$. Its deterministic future history then agrees too.
Summing the three errors proves \eqref{eq:finite-family-rate}.
\end{proof}

\begin{theorem}[Uniform logarithmic bounds in finite samples]
\label{thm:finite-logarithmic-noise}
For universal $C<\infty$, the exponent $a>0$ from
Theorem~\ref{thm:logarithmic-spectrum}, every odd $N\ge3$, and
$1\le k\le N$,
\begin{equation}\label{eq:finite-log-noise}
 1-p_N(k)\le\frac{C}{[\log(k+2)]^a}.
\end{equation}
Let $\mathcal A^{\rm lab}_{N,r}(k)$ be the expected common
active-label count for the original and $k$-resampled samples at
separation $r$. Uniformly for $0<r\le1/3$,
\begin{equation}\label{eq:finite-log-overlap}
 \EE|C_r(X)\cap C_r(X^{[k]})|
 \le\mathcal A^{\rm lab}_{N,r}(k)
 \le\frac Cr\min\left\{1,
       \left(\frac{\log(e/r)}{\log(k+2)}\right)^a\right\}.
\end{equation}
Consequently, for every $k_N\to\infty$ with $k_N\le N$ and every
$0<b<a$, the two finite configurations are disjoint with probability
tending to one at cutoff
$r_N=(\log(k_N+2))^{-b}$, once this is at most $1/3$.
Their later histories then share no vertex.
For independent-coordinate resampling with probability $s$, the same
two bounds hold with $k$ replaced by $Ns$ and with a changed constant.
\end{theorem}
\begin{proof}
For $1\le j\le N/2$, use a common background of size $M=N-j$
and two independent sets of $j$ probes. The entrance pair is exactly
the pair defining $p(j)$; this identity holds separately for every
integer $j$ by the fixed-family theorem, so there is no new limiting
exchange when $j$ is subsequently chosen as a function of $N$.
Lemma~\ref{lem:finite-family-rate} and the label/position correction
at most $j/N$ give
\begin{equation}\label{eq:finite-profile-rate}
 |p_N(j)-p(j)|\le C(j+1)^2 N^{-\chi_{\rm f}}.
\end{equation}
For $r\ge2N^{-\chi_{\rm f}}$, the coupling is already valid at
cutoff $r$, since $M^{-\chi_{\rm f}}\le2N^{-\chi_{\rm f}}$.
Each common-point count is at most $1/r$. The additional discrepancy
between common labels and common positions is at most the number
of refreshed labels active in the original sample, with expectation
$j\EE K_N(r)/N\le j/(Nr)$. Therefore
\begin{equation}\label{eq:finite-overlap-rate}
 |\mathcal A^{\rm lab}_{N,r}(j)-\mathcal A_r(j)|
       \le \frac Cr(j+1)^2N^{-\chi_{\rm f}}.
\end{equation}

Choose $j=\min\{k,\lfloor N^{\chi_{\rm f}/4}\rfloor\}$.
For large $N$, this lies between one and $N/2$,
$\log(j+2)\ge c\log(k+2)$, and both errors above, before the
factor $1/r$, are at most $CN^{-\chi_{\rm f}/2}$.
The finite spectral formulas make $1-p_N(k)$ and
$\mathcal A^{\rm lab}_{N,r}(k)$ nonincreasing in $k$.
Compare them at $j$, apply Theorem~\ref{thm:logarithmic-spectrum},
and absorb the polynomial error into the logarithmic term. For the
overlap bound this absorption is uniform in $r$, since
\[
 \min\left\{1,
    \left(\frac{\log(e/r)}{\log(k+2)}\right)^a\right\}
       \ge [\log(N+2)]^{-a},\qquad r\le1/3.
\]
This proves the assertions when $r\ge2N^{-\chi_{\rm f}}$.
For smaller $r$, the logarithmic ratio in
\eqref{eq:finite-log-overlap} is bounded below by a positive universal
constant, so the trivial bound $\mathcal A^{\rm lab}_{N,r}(k)\le1/r$
suffices. Bounded $N$ are covered by increasing $C$; indeed the
right sides' logarithmic factors have a positive lower bound there.

Common physical points always have a common unrefreshed label,
outside null accidental coincidences. This proves the first inequality
in \eqref{eq:finite-log-overlap}. Substitution of $r_N$ and Markov's
inequality give disjointness, since the remaining bound is
$O((\log k_N)^{b-a}(\log\log k_N)^a)\to0$.
Forward evolution only deletes locations, proving the history assertion.
For independent-coordinate noise with $Ns\ge2$, condition on the binomial number
$K$ of refreshed coordinates, of mean $Ns$. Its variance is at most
$Ns$, so $\PP(K<Ns/2)\le4/(Ns)$. On the complement apply the
just-proved bounds at $Ns/2$; the exceptional probability, with an
extra factor $1/r$ for overlap, is absorbed into the logarithmic bound.
Bounded $Ns$ are covered by the trivial agreement and count bounds.
\end{proof}

\begin{corollary}[Quantitative bounds for all random edit mixtures]
\label{cor:mixed-logarithmic-noise}
Suppose $N$ and $N'=N-d+a'$ are positive odd integers, with
$0\le d\le N$ and $a'\ge0$. Delete $d$ uniform labels and add $a'$
fresh uniform points, and put $q=d+a'\ge1$.
The survivor agreement probability is at most
\[
 C[\log(q+2)]^{-a/2}.
\]
The expected number of common active points at cutoff $0<r\le1/3$
is at most
\[
 \frac Cr\min\left\{1,
      \left(\frac{\log(e/r)}{\log(q+2)}\right)^{a/2}\right\}.
\]
Thus $q\to\infty$ gives disjoint configurations at
$r_q=(\log(q+2))^{-b}$ for every $0<b<a/2$, uniformly over the
compatible initial and final sample sizes.
\end{corollary}
\begin{proof}
Use the two conditional Cauchy--Schwarz inequalities in
Corollary~\ref{cor:growing-mixed-noise}. At least one of $d,a'$ is
at least $q/2$. Apply Theorem~\ref{thm:finite-logarithmic-noise}
to that factor in its own sample size and bound the other agreement
factor by one, or the other active-count factor by $1/r$.
The cases with a sample size one satisfy these trivial bounds as well.
Taking square roots gives both estimates. The cutoff conclusion
follows as in Theorem~\ref{thm:finite-logarithmic-noise}.
\end{proof}

\subsection{A second-moment reduction for the history coefficient}
\label{sec:history-coefficient}

\status{Working proof: a capture-mass reduction for the sharp constant}

\begin{lemma}[Uniform parent-interval second moments]
\label{lem:parent-second-moment}
At active count $m$ in a size-$M$ iid background, let
$T_{M,m}$ be the sum of squared current gap lengths. There is a
universal $C$ such that
\[
 \EE T_{M,m}\le C/m,\qquad \EE L_{M,m}^2\le C/m^2.
\]
The second assertion includes the last even pair, with $L_{M,2}=1$.
Consequently $\EE\sum_e L_e^2\le C$ uniformly in the initial size.
\end{lemma}
\begin{proof}
Condition on the unordered gaps at count $m\ge3$. Put $s$ for their
minimum and $T$ for their square sum. The two neighbours $a,c$ of
the minimum are a uniform ordered pair of the other gaps, so
\[
 \EE(a+c)=\frac{2(1-s)}{m-1},\quad
 \EE(a^2+c^2)=\frac{2(T-s^2)}{m-1},\quad
 \EE(ac)=\frac{(1-s)^2-(T-s^2)}{(m-1)(m-2)}.
\]
The merger increases the square sum by $2(as+ac+sc)$. Its conditional
mean is at most
\[
 \frac4{m(m-1)}+\frac2{(m-1)(m-2)}\le\frac{15}{m^2},
\]
using $s\le1/m$. The initial Dirichlet square sum has expectation
$2/(M+1)$. Summing the preceding bound over counts of the same parity
above $m$ gives $\EE T_{M,m}\le10/m$.
Expand $(a+s+c)^2$ and use the three conditional identities above.
They give $\EE L_{M,m}^2\le50/m^2$; the last pair satisfies this
bound directly. Summation proves the last assertion.
\end{proof}

\begin{proposition}[One-resampling history cost through captured origin mass]
\label{prop:history-capture}
For each background edge $e$, let $W_e$ be the Lebesgue measure of
insertion origins for which that edge disappears under one insertion.
For an even background of size $M=N-1$,
\begin{equation}\label{eq:history-capture-linearization}
 \EE H_N(1)=4\EE\sum_e W_e+O(1),\qquad
 \EE H_N^{\rm lab}(1)=4\EE\sum_e W_e+O(1),
\end{equation}
with absolute bounded errors. If $W_{M,m}$ denotes the captured mass
at count $m\ge4$, then
\begin{equation}\label{eq:capture-mean-bounds}
 \EE s_{M,m}\le\EE W_{M,m}
 \le\EE s_{M,m}+\frac{1-\EE s_{M,m}}{m-1}.
\end{equation}
In particular
\begin{equation}\label{eq:history-one-coefficient-bounds}
 e^{-\gamma}\le
 \liminf_{N\to\infty}\frac{\EE H_N(1)}{\log N}
 \le\limsup_{N\to\infty}\frac{\EE H_N(1)}{\log N}
 \le2+e^{-\gamma},
\end{equation}
and the same bounds hold for $H_N^{\rm lab}(1)$.
For every fixed $k$, the triangle inequality also gives
$\limsup\EE H_N(k)/(k\log N)\le2+e^{-\gamma}$.
Neither the existence nor the value of the sharp coefficient is
asserted by these bounds.

More precisely, immediately before the next collision at separation
$s=s_{M,m}$, define the total corner mass
\[
 A_{M,m}=\sum_i\left[
   \mu_i([0,\min(s,x_i/2)])+
   \mu_i([x_i-\min(s,x_i/2),x_i])\right]\in[0,1].
\]
There are no midpoint atoms almost surely. Then
\begin{equation}\label{eq:capture-corner-identity}
 \EE W_{M,m}
   =\EE s_{M,m}+\frac{\EE A_{M,m}-\EE s_{M,m}}{m-1}.
\end{equation}
\end{proposition}
\begin{proof}
Use two independent uniform insertions in the common even background.
In the notation of Theorem~\ref{thm:uniform-history}, the exact edge
identity is $H_N(1)=2(a+b)-2J-2I$, with $0\le I\le J$.
Conditional on the background,
$\EE(a+b)=2\sum_eW_e$ and $\EE J=\sum_eW_e^2$.
Since $W_e\le L_e$, Lemma~\ref{lem:parent-second-moment} bounds the
expected correction by an absolute constant. The label cost differs
by at most two, proving \eqref{eq:history-capture-linearization}.

Just before an edge of separation $s$ is removed, let its neighbouring
gap lengths be $a,c$. The central gap contributes all its mass $s$.
The additional captured masses are the pushed masses within distance
$\min(s,a/2)$ and $\min(s,c/2)$ of the two endpoints in the neighbouring
gaps. Conditional on the unordered decorated gaps, each neighbour is
uniform among the other $m-1$ marks by Lemma~\ref{lem:gap-order}.
Their two mean contributions add to $(A_{M,m}-s)/(m-1)$: the minimum
gap has length $s$, so its two corner masses sum to its whole mass $s$.
This proves \eqref{eq:capture-corner-identity}. Since
$s\le A_{M,m}\le1$, it also proves \eqref{eq:capture-mean-bounds}.
The final pair contributes at most one and can be omitted in an
asymptotic calculation.

Lemma~\ref{lem:truncated-matching-length} gives
$\sum_m\EE s_{M,m}\sim\log M/(4e^\gamma)$, while
$\sum_{4\le m\le M,\ m\ {\rm even}}1/(m-1)\sim\tfrac12\log M$.
Substitute the capture bounds in
\eqref{eq:history-capture-linearization} to obtain
\eqref{eq:history-one-coefficient-bounds}.
For fixed $k$, perform its resamplings successively and sum the
one-step history distances, whose marginal distributions are the same.
\end{proof}

Equation~\eqref{eq:capture-corner-identity} isolates a sufficient
condition for the sharp coefficient:
\[
 \EE A_{M,m}\longrightarrow\frac12
 \quad\hbox{whenever }m\to\infty,\quad m/M\to0.
\]
If this condition holds, the truncated harmonic-sum argument gives
$\EE H_N(1)\sim(1+e^{-\gamma})\log N$.
Lemma~\ref{lem:corner-mass-law} below proves it on a range sufficient
for the logarithmic asymptotic. At a deterministic cutoff $r$, the
analogous corner mass $A_M(r)$ satisfies the exact identity
\[
 \EE A_M(r)=\frac{1-\{\EE K_{M+1}(r)-\EE K_M(r)\}}2,
\]
by the add-one count difference. Evaluating it at the random next
collision separation requires a further argument; an unconditional
fixed-cutoff identity cannot simply be evaluated at that stopping time.

The conditional mean gap measure cannot be replaced by Lebesgue
measure using shielding and reflection alone. For example, take two
uniform interior points in $[0,1]$, conditioned to match each other
before either protected endpoint. Write their gaps as $(a,s,c)$.
Their conditional density is $6$ on $a,c\ge s$, $a+s+c=1$.
If $\mu$ is Lebesgue mass pushed through that one collapse, then
\[
 \EE\mu([0,z])=z-\frac56z^3,\qquad 0\le z\le1/8.
\]
To check it, the relevant region has $a\le2z$; rescale $a,s$ by $z$.
On the triangle $0\le s\le a\le2$, put $\alpha=\min(s,a/2)$.
The two incoming atomic contributions integrate to
\[
 \int(\alpha+s/2)\ind_{\{a+s\le1\}}\,da\,ds=1/18,
 \qquad \int(3s/2)\ind_{\{a\le1\}}\,da\,ds=1/4,
\]
while the removed original mass integrates to
$\int[\min(1,a+2s)-(a-\alpha)]_+\,da\,ds=4/9$.
Their difference, multiplied by the density $6$ and by $z^3$, is
$-5z^3/6$. This is a finite counterexample to that mean-measure shortcut,
not an identification of the limiting marked gap law.

\subsection{Corner-mass laws and the sharp one-point history constant}
\label{sec:sharp-one-history}

\status{Working proof: the leading mean coefficient for one edit or one resampling}

\begin{lemma}[Corner mass at deterministic cutoffs and collision times]
\label{lem:corner-mass-law}
For either entrance parity, the total corner mass satisfies
$A(r)\to1/2$ in $L^2$ as $r\downarrow0$.
For the finite size-$M$ process, the same holds if
\begin{equation}\label{eq:marked-cutoff-regime}
 r_M\to0,\qquad \frac{Mr_M}{\log(1/r_M)}\longrightarrow\infty.
\end{equation}
At a fixed active count it follows that
\begin{equation}\label{eq:corner-at-count}
 \EE|A_{M,m}-1/2|\longrightarrow0
 \quad\hbox{if }m\to\infty,\qquad \frac{m\log m}{M}\to0,
 \qquad m\equiv M\pmod2.
\end{equation}
This coefficient argument has an extra logarithmic condition.
Lemma~\ref{lem:full-corner-law} removes it by a separate Poisson
comparison, proving the full regime $Mr_M\to\infty$.
\end{lemma}
\begin{proof}
For a limiting gap of scaled length $Y\sim f$, let $B$ be its total
corner mass at inherited cutoff one. Then $0\le B\le Y$.
Conditional on the cyclic gap lengths, the gap marks are independent
and have these shielding laws. In the count-biased rooted bridge
\eqref{eq:palm-renewal-bridge}, attach the conditional marks to the
stopped prefix used in the entrance-fluctuation proof. This is a
Markov kernel and cannot increase the prefix TV error. The comparison
is therefore with iid pairs $(Y_i,B_i)$. The unexposed arc has length
at most $R^{1/4}$ on the same good event, where $R=1/r$, and its
total reward is bounded by that length. The renewal-reward law of
large numbers gives
\[
 A(r)=R^{-1}\sum_i B_i\longrightarrow\frac{\EE B}{\EE Y}
       \quad\hbox{in probability}.
\]
For completeness, this law follows by applying the iid strong law
to $Y_i$ and $B_i$, then inverting the partial sums of $Y_i$.
The count tilt has vanishing TV error as already proved.
Boundedness of $A(r)$ upgrades the convergence to $L^2$.

The add-one identity and the parity renewal formulas give exactly
\[
 \EE A^{\rm e}(r)=\frac12-\frac{a_0(1/r)}{2r},\qquad
 \EE A^{\rm o}(r)=\frac12+\frac{a_0(1/r)}{2r}.
\]
The factorial tail of $a_0$ identifies the limiting reward ratio as
$1/2$. In particular, if
$\psi(y)=\EE\mu_{1,y}([0,\min(1,y/2)])$, reflection gives the sum rule
\begin{equation}\label{eq:corner-reward-sum-rule}
 \int_1^\infty\psi(y)f(y)\,dy=\frac{\EE Y}{4}
                                           =\frac{e^\gamma}{2}.
\end{equation}

We next justify finite-count transfer in
\eqref{eq:marked-cutoff-regime}. Fix a length truncation $D>1$,
a depth $d$, and a lower cutoff ratio $0<\delta<1$.
Call a gap tree good if its root length is at most $Dr$, it is
fully revealed through depth $d$, and every terminal cutoff is at
least $\delta r$. For any nonnegative function of this geometry
bounded by its root length, its additive sum over good gaps has
finite-count first and second moments differing from the same
entrance-parity moments by
\begin{equation}\label{eq:local-marked-coefficient-error}
 C_{D,d,\delta}r^{-6}e^{-2M\delta r}+O_D(r).
\end{equation}
The $O_D(r)$ is needed only in the second moment.

Here is the coefficient argument for this claim. For one distinguished
gap with root length $x<1$, the unconditioned Poisson density, after
revealing its tree, is
\[
 \left(\prod_{i=1}^{3^d}w_{t_i}^{\lambda}(x_i)\right)
                         U_r^{\lambda}(1-x)
\]
against its root location, root length and split coordinates.
The outside gaps have been summed, so no unbounded product of
terminal factors remains there. For two nonadjacent distinguished
gaps, multiply their local factors and the two renewal densities
of the intervening arcs. Extract the coefficient
$M![\lambda^M]e^\lambda$ in each expression.

Expand the local factors by \eqref{eq:finite-convolution} and the
outside renewal densities by \eqref{eq:renewal-expansion}. Each local
tree has at most $D/\delta$ convolution components in total, because
they all have length at least $\delta r$ and sum to at most $Dr$.
After including the absolute subset sums, each terminal factor is bounded
by $C_{D,d,\delta}/r$. The local tree reference variables number
$3^d$ when the root length is included; rescaling them by $r$
therefore bounds the integrated absolute local factors by a constant.
Each outside absolute renewal series is bounded by $r^{-3}$, exactly
as in Lemma~\ref{lem:intensity-coefficients}.

All component lengths together sum to one. The empty and full subsets
give the corresponding entrance-parity moment density. Every other
subset has total length between $\delta r$ and $1-\delta r$, so its
extracted coefficient has modulus at most $e^{-2M\delta r}$.
This proves the displayed error for first moments and for the
nonadjacent part of second moments. Diagonal and adjacent-gap terms
are bounded directly by $CDr$: each retained reward is at most its
gap length, retained lengths are at most $Dr$, and all gap lengths
sum to at most one. Singleton sectors cannot be retained once
$Dr<1$. These observations prove
\eqref{eq:local-marked-coefficient-error}. Under
\eqref{eq:marked-cutoff-regime} its right side tends to zero.

Use the corner mass of the depth-$d$ uniform completion as this
reward, putting it zero on bad gaps. In the entrance law its additive
sum converges to a constant by the same marked renewal-reward argument.
The first and second moment comparison gives the same conclusion for
the finite sample. The approximation error is controlled separately.
The expected two-query budget summed over limiting gaps is at most
$2\rho^d$. The good-geometry first-moment comparison transfers this
bound to finite samples, up to a vanishing error. For the bad-gap
mass, apply that comparison to the reward equal to the root length
on good gaps and subtract from one. Its limiting upper bound is
\[
 c\,\EE[Y\ind_{\{Y>D\}}]+C\delta(3J)^d,
 \qquad c=(\EE Y)^{-1}.
\]
The second term follows by summing inverse terminal cutoffs, and
the entrance even empty probability tends to zero. This normalization
also controls finite premature-empty gaps. First choose $d$ large,
then $D$ large and $\delta$ small. The total approximation error tends
to zero, proving $A_M(r_M)\to1/2$ in $L^1$, hence in $L^2$ by boundedness.

Finally let $s=s_{M,m}$ under the hypotheses of
\eqref{eq:corner-at-count}, and fix $0<\epsilon<c/2$.
The count law places $s$ between $r_-=(c-\epsilon)/m$ and
$r_+=(c+\epsilon)/m$ with probability tending to one, and puts all
counts in that time window between
$(1-C\epsilon)m$ and $(1+C\epsilon)m$, with probability tending to one.
Both deterministic cutoffs satisfy \eqref{eq:marked-cutoff-regime}.

Between background collisions the collapse map is constant and its
corner regions expand, so $A_M(r)$ is nondecreasing. A collision
changes this mass only inside its parent interval, hence its downward
jump is at most that interval's length. Therefore throughout the
bracket,
\[
 A_M(r_-)-\sum_{e\text{ in bracket}}L_e
 \le A_M(s-)\le
 A_M(r_+)+\sum_{e\text{ in bracket}}L_e.
\]
On the high-probability count event, the sum is bounded by the sum
over a deterministic count window of relative width $C\epsilon$.
Its expectation is at most $C\epsilon$, using
$\EE L_{M,j}\le3/j$. The complementary event costs at most one
when bounding the deviation of $A_{M,m}$.
The endpoint corner masses tend to $1/2$ in $L^1$; send first $M$
to infinity and then $\epsilon$ to zero. This proves
\eqref{eq:corner-at-count}.
\end{proof}

\begin{theorem}[Sharp one-point mean history constants]
\label{thm:sharp-one-history}
For one coordinate resampling in an odd size-$N$ sample,
\begin{equation}\label{eq:sharp-one-history}
 \EE H_N(1)\sim\EE H_N^{\rm lab}(1)
                   \sim(1+e^{-\gamma})\log N.
\end{equation}
The expected spatial history cost of one fresh uniform addition or
one uniform deletion in a size-$M$ sample, of either parity, is
\[
 \frac{1+e^{-\gamma}}2\log M+o(\log M).
\]
These are mean asymptotics. Section~\ref{sec:small-edit-history}
first extends them to counts $o(\log N)$;
Section~\ref{sec:sublinear-history} reaches every sublinear count.
Section~\ref{sec:fixed-history-LLN} also proves relative concentration
for fixed edit counts. Section~\ref{sec:finite-history-CLT} proves centered Gaussian laws for every fixed edit family;
Section~\ref{sec:late-history-variance} treats every growing sublinear count.
\end{theorem}
\begin{proof}
Uniformly for parity-compatible counts
$\log M\le m\le M/(\log M)^2$, Lemma~\ref{lem:corner-mass-law}
gives $\EE A_{M,m}=1/2+o(1)$: any violating sequence satisfies
\eqref{eq:corner-at-count}. The omitted ranges have harmonic sum
$O(\log\log M)$ and $0\le A_{M,m}\le1$. Consequently
\[
 \sum_{\substack{m\ge3\\m\equiv M\ ({\rm mod}\ 2)}}
       \frac{\EE A_{M,m}}{m-1}\sim\frac14\log M.
\]
Also $\sum_m\EE s_{M,m}\sim\log M/(4e^\gamma)$, while
$\sum_{m\ge3}\EE s_{M,m}/(m-1)=O(1)$.
The decorated-gap identity therefore gives
\[
 \EE\sum_eW_e
     =\frac{1+e^{-\gamma}}4\log M+o(\log M).
\]
The final even pair and the bounded terminal odd stages do not affect
the coefficient. Proposition~\ref{prop:history-capture} proves
\eqref{eq:sharp-one-history} by taking $M=N-1$.
For one addition the exact spatial cost is twice the missing-background
edge count, plus the matching-size difference (zero or one). Its mean
is $2\EE\sum_eW_e+O(1)$. Reversing an iid add-one pair is uniform
deletion, proving the last assertion.
\end{proof}

\subsection{Bounded pair interactions and sharp small-edit history means}
\label{sec:small-edit-history}

\status{Working proof: the sharp coefficient for every fixed edit count}

\begin{lemma}[A bounded mixed second difference of the matching]
\label{lem:matching-rectangle}
Let $F(B)$ be the indicator vector of the completed spatial matching
of $B$, in the space of all unordered pairs of the locations under
consideration. If $B$ consists of $M$ iid uniform points and $X,Y$
are independent uniform probes, then, uniformly in $M\ge0$,
\begin{equation}\label{eq:matching-rectangle}
 \EE\|F(B+X+Y)-F(B+X)-F(B+Y)+F(B)\|_1\le C.
\end{equation}
\end{lemma}
\begin{proof}
Use closed parent intervals, which does not affect their lengths.
An accepted edge is protected under arbitrary insertions and deletions
outside its parent interval. Indeed all original interior pairs are
closer than either boundary and hence than every exterior point;
the shielding induction in Lemma~\ref{lem:edge-insertions} applies
even when exterior points are deleted.

We first prove a deterministic witness statement. For an edge $e$
not incident to $x$ or $y$, if its indicator has a nonzero mixed
difference on the four input sets, then some corner where $e$ is
accepted has a parent interval containing both $x$ and $y$.
If exactly one corner is positive, or the two positive corners are
diagonally opposite, both adjacent flips kill $e$ at a positive
corner. Exterior protection gives the witness immediately.

It remains to consider three positive corners. Root the square at
the positive corner opposite the unique zero, and call its parent
interval $I$. The two flips at this root may insert or delete their
respective locations $x,y$. Write $I_x,I_y$ for the parents in the
other positive corners. A full-circle parent is already a witness,
so suppose all three are proper intervals. If both source locations
were outside $I$, exterior protection would preserve $e$ in the
zero corner. Thus, after exchanging them if necessary, $x\in I$.
Since the second flip kills $e$, $y\in I_x$ and $x\in I_y$.

If $x\in I_x$ or $y\in I_y$ we are done. Otherwise compare the
active configurations immediately before the distance of $e$.
Each one-point flip changes that configuration by exactly one point.
For a common edge, adding one active point can only shorten its
parent interval, while removing one can only enlarge it; other
changes leave the parent unchanged. Hence $x\in I\setminus I_x$
forces $I_x\subset I$, so $y\in I$ too. Similarly $I_y\subset I$.
Both positive neighbouring configurations must therefore be
$C+\{u\}$ and $C+\{v\}$, where $C$ is the rooted corner's active set.
If $u=v$, their parent intervals coincide, contradicting
$x\notin I_x$ and $x\in I_y$. If $u\ne v$, the active set in the
zero corner, being at symmetric-difference distance one from each,
is either $C$ or $C+\{u,v\}$. In either case both endpoints of $e$
are still active, so the greedy scan accepts $e$ at its distance.
This contradiction proves the witness statement.

For a size-$n$ iid sample let $\xi_n$ be its empirical point measure.
Symmetrization and the $L^2$ maximal inequality for Rademacher walks
give $\EE\Delta_n^2\le C/n$ for the interval discrepancy used in
Lemma~\ref{lem:empirical-seed}. Consequently
\[
 \EE\sum_e\xi_n(I_e)^2
 \le 2\EE\sum_e L_e^2+n\EE\Delta_n^2\le C,
 \qquad
 \EE\sum_e L_e\xi_n(I_e)\le C,
\]
by Lemma~\ref{lem:parent-second-moment}. These estimates include
interval endpoints through the definition of $\Delta_n$.

In the corner with neither probe present, the conditional witness
probability for edge $e$ is $L_e^2$. With one probe present it is
$L_e\xi_n(I_e)$: the present probe is a uniform tag of the iid
sample, and the other source remains independent uniform.
With both present it is at most $2\xi_n(I_e)^2$, by sampling two
distinct tags without replacement. Thus the expected total number
of parent witnesses in the four corners is bounded. Each nonzero
mixed indicator has magnitude at most two. Edges incident to $x$
or $y$ contribute at most eight in total, since each point has degree
at most one in each matching. This proves
\eqref{eq:matching-rectangle}; samples with fewer than two points
are covered directly.
\end{proof}

\begin{theorem}[Sharp mean coefficients for small edit counts]
\label{thm:sharp-small-edits}
Let $h_M$ be the expected spatial matching distance between an iid
size-$M$ background and that background with one fresh uniform point.
For two independent sets of $d$ and $b$ probes added to a common
size-$M$ iid background, put $q=d+b$. Their expected spatial history
distance satisfies the uniform finite bound
\begin{equation}\label{eq:history-linearization}
 \left|\EE H(B+X_1+\cdots+X_d,
                   B+Y_1+\cdots+Y_b)-q h_M\right|\le Cq^2.
\end{equation}
Consequently, for every $1\le k=k_N=o(\log N)$ and odd $N\to\infty$,
\begin{equation}\label{eq:sharp-small-resampling}
 \EE H_N(k)\sim\EE H_N^{\rm lab}(k)
       \sim(1+e^{-\gamma})k\log(N/k).
\end{equation}
For any mixture of $d$ uniform deletions and $b$ fresh additions with
$q=d+b\ge1$ and $q=o(\log N)$, its expected spatial history cost is
\begin{equation}\label{eq:sharp-small-mixed}
 \frac{1+e^{-\gamma}}2\,q\log(N/q)\,(1+o(1)).
\end{equation}
These include every fixed resampling count and every fixed nonzero
mixture of random edits. The uniform $O(q^2)$ error only gives the stated range.
Section~\ref{sec:sublinear-history} uses a common stopping cutoff
to reach all sublinear edit counts.
Section~\ref{sec:history-concentration} gives relative concentration
for growing counts, and Section~\ref{sec:fixed-history-LLN} treats
fixed counts. Section~\ref{sec:finite-history-CLT} proves centered Gaussian laws for every fixed edit family;
Section~\ref{sec:late-history-variance} treats every growing sublinear count.
\end{theorem}
\begin{proof}
Write $\Delta_zF(B)=F(B+z)-F(B)$. Discrete telescoping gives
\[
 F(B+Z_1+\cdots+Z_j)-F(B)
       =\sum_{i=1}^j\Delta_{Z_i}F(B)+R_j,
 \qquad \EE\|R_j\|_1\le Cj^2.
\]
Indeed expand each increment along the preceding probes; each
correction is a mixed difference from
Lemma~\ref{lem:matching-rectangle}, on an iid enlarged background
independent of its two remaining probes.

For two single-probe increments on the same background, their supports
overlap only on missing background edges and common nonbackground
edges. If $J$ is the number missing in both copies and $I$ the number
of common new edges, then $I\le J$: each endpoint of a common new
edge belongs to a background edge missing in both, except for at
most one unmatched background point. Hence $2I\le2J+1$.
Conditional independence gives $\EE[J\mid B]=\sum_eW_e^2$, so
\[
 \EE|\operatorname{supp}(\Delta_XF(B))
          \cap\operatorname{supp}(\Delta_YF(B))|
       \le2\EE\sum_eW_e^2\le C.
\]
For vectors with coordinates in $\{-1,0,1\}$, the loss between
the sum of their $\ell^1$ norms and the norm of their signed sum is
at most twice the sum of their pairwise support overlaps.
Apply this to the $q$ independent single-probe increments, with
positive signs for the first family and negative signs for the second.
Their expected summed norms equal $q h_M$, and the expected loss is
$O(q^2)$. The two telescoping remainders give
\eqref{eq:history-linearization}.

For resampling use $M=N-k$ and $d=b=k$.
Theorem~\ref{thm:sharp-one-history} gives
$h_M\sim(1+e^{-\gamma})\log M/2$.
Since $k=o(\log N)$, the error $O(k^2)$ is negligible relative to
$k\log(N/k)$, and $\log M\sim\log(N/k)$.
The label correction is at most $2k$, proving
\eqref{eq:sharp-small-resampling}.
For a deletion/addition mixture the common retained background has
size $M=N-d$, and its deleted original points form an independent
uniform probe set of size $d$. Apply
\eqref{eq:history-linearization} and the same asymptotic with $q$
in place of $2k$ to obtain \eqref{eq:sharp-small-mixed}.
\end{proof}

\subsection{The sharp history coefficient throughout the sublinear regime}
\label{sec:sublinear-history}

\status{Working proof: all sublinear random edit counts}

\begin{lemma}[Corner mass throughout the mesoscopic regime]
\label{lem:full-corner-law}
The finite corner-mass law strengthens to
\[
 A_M(r_M)\longrightarrow\tfrac12\quad\hbox{in }L^2
 \qquad\hbox{if }r_M\to0,\quad Mr_M\to\infty.
\]
Consequently $\EE|A_{M,m}-1/2|\to0$ whenever $m\to\infty$,
$m/M\to0$, and $m\equiv M\pmod2$.
\end{lemma}
\begin{proof}
We first need estimates uniform in Poisson intensity. Write
$\alpha=\lambda r$, $R=1/r$, and let $P_\lambda$ be the initial
Poisson count. For $\alpha\ge1$ and $r\le1/4$,
\begin{equation}\label{eq:poisson-isolation-uniform}
 \PP(C_{P_\lambda}(r)=\varnothing)
 +\EE\!\left[\ind_{\{C_{P_\lambda}(r)\ne\varnothing\}}
                      \sum_i x_i^2\right]\le Cr.
\end{equation}
Here and below a corner mass in the empty state is zero.
To verify this bound, the lower envelope in
Lemma~\ref{lem:finite-intensity-envelope}, integrated against $dy$,
implies $\kappa_\alpha\ge c\min(\alpha,1)$: $f_\alpha$ has mass
one and $\kappa_\alpha\le\min(\alpha,1)$.
The upper envelope and Lemma~\ref{lem:uniform-renewal} therefore give
\[
 h_r^\lambda(x)\le Cr^{-1}e^{-cx/r},\qquad
 U_r^\lambda(x)\le C/r\qquad(\lambda r\ge1).
\]
The gap-square expectation is exactly
\[
 h_r^\lambda(1)+\int_0^1x^2h_r^\lambda(x)
                                  U_r^\lambda(1-x)\,dx\le Cr.
\]
An empty state either started empty or had its last pair removed at
some $s\le r$. The Poisson last-pair density is
$w_s^\lambda(s)w_s^\lambda(1-s)$. For $s\le1/\lambda$ the envelope
bounds this by $C\lambda^2e^{-c\lambda/2}$; for
$1/\lambda\le s\le r$ it is at most $Cs^{-2}e^{-c/(2s)}$.
Their integrals, together with $e^{-\lambda}$, are at most
$Ce^{-c'/r}\le Cr$. This proves \eqref{eq:poisson-isolation-uniform}.

The Poisson add-one derivative gives a useful exact mean identity:
\[
 1-2\EE A_{P_\lambda}(r)
   =\partial_\lambda U_r^\lambda(1)
   =e^{-2\lambda}
       +2\int_r^1U_r^\lambda(1-y)e^{-2\lambda y}\,dy.
\]
Differentiate the positive convolution exponential
\eqref{eq:renewal-expansion} to obtain the last equality.
The series are locally finite in space, so differentiation is
justified term by term. Thus
\begin{equation}\label{eq:poisson-corner-mean}
 0\le1-2\EE A_{P_\lambda}(r)
       \le e^{-2\lambda}+C e^{-2\alpha}/\alpha.
\end{equation}

We also need concentration, without assuming convergence of the
finite-intensity marks. Conditional on the Poisson cyclic gaps, their
marks are independent. In scaled coordinates let $B_\alpha(y)$ be
the corner reward in a shielded gap of length $y$, so
$0\le B_\alpha(y)\le y$. Put
\[
 b_\alpha(y)=f_\alpha(y)\EE B_\alpha(y),\qquad
 b_\alpha^* =\int b_\alpha(y)\,dy,\qquad c=(2e^\gamma)^{-1}.
\]
Uniformly for $\alpha\ge1$, $b_\alpha(y)\le Cy e^{-cy}$,
with a possibly different positive constant in the exponent.
The one-gap moment formula gives
\[
 \EE A_{P_\lambda}(r)
   =\int b_\alpha(y)U_\alpha(R-y)\,dy+o(1)
   =c b_\alpha^*+o(1)
\]
as $r\to0$, $\alpha\to\infty$. The first error is the exponentially
small singleton term. The second uses the uniform comparison of
$U_\alpha$ to $U_\infty$ and its renewal limit.

For the second moment, diagonal and adjacent-gap rewards contribute
$O(r)$, by \eqref{eq:poisson-isolation-uniform} and
$2x_ix_{i+1}\le x_i^2+x_{i+1}^2$. The remaining ordered pairs give
\[
 r\iint b_\alpha(y)b_\alpha(z)
      \int_0^{R-y-z}U_\alpha(t)U_\alpha(R-y-z-t)\,dt\,dy\,dz.
\]
Take the integrand as zero when $y+z>R$. The inner integral times
$r$ tends to $c^2$, uniformly for $y,z$ in compact sets, and is
uniformly bounded. Indeed discard fixed-length end intervals and
use the renewal limit on the remaining interval, then let that fixed
length increase. The exponential envelope controls the tails even
though $b_\alpha$ varies. Hence
$\EE A_{P_\lambda}(r)^2=c^2(b_\alpha^*)^2+o(1)$.
Equation~\eqref{eq:poisson-corner-mean} identifies
$cb_\alpha^*\to1/2$, proving the Poisson $L^2$ law.

Here is a de-Poissonization estimate for the marks. Couple every
sample size using prefixes of one iid sequence. The add-one identity
expresses $A(B,r)$ as
$[1-\EE_Y\{K(B+Y,r)-K(B,r)\}]/2$.
For two separated probes the mixed difference of $K$ vanishes by
Lemma~\ref{lem:separated-insertions}; in every case its magnitude is
at most two. Consequently
\[
 \EE|A_{j+1}(r)-A_j(r)|\le
 \PP(C_j(r)=\varnothing)+3\EE\!\left[
      \ind_{\{C_j(r)\ne\varnothing\}}\sum_i x_i^2\right]=:v_j(r).
\]
For a Poisson variable $P$ of integer mean $M$, put $p_j=\PP(P=j)$.
The probability that the path of counts from $M$ to $P$ crosses
$j\leftrightarrow j+1$ is at most $C\sqrt M\,p_j$.
For $j\ge M$, divide the upper tail by $p_j$ and sum
$M^\ell/[(j+1)\cdots(j+\ell)]$; this ratio decreases in $j$ and
is at most $1/p_M\le C\sqrt M$. The lower-tail ratio for $j<M$
is $\sum_{\ell=0}^j(j)_\ell/M^\ell$, increasing in $j$, and has
the same bound. Telescoping and independence of $P$ give
\begin{equation}\label{eq:corner-depoissonization}
 \EE|A_P(r)-A_M(r)|
 \le C\sqrt M\sum_jp_jv_j(r)\le C\sqrt M\,r,
 \qquad Mr\ge1.
\end{equation}
Thus the fixed-count law follows when $r\le M^{-2/3}$.
When $r>M^{-2/3}$, the hypotheses of
Lemma~\ref{lem:corner-mass-law} hold automatically, and give the
same conclusion. Boundedness upgrades $L^1$ transfer to $L^2$.
The collision-time bracket in that lemma now applies whenever
$m\to\infty$ and $m/M\to0$, proving the last assertion.
\end{proof}

\begin{lemma}[Empirical parent masses before a prescribed active count]
\label{lem:tagged-parent-tail}
For a size-$n$ iid sample, let $I_e$ be the closed parent interval of
edge $e$, $L_e=|I_e|$, and $\xi_n$ its empirical point measure.
Let $\mathcal E_{n,\ell}$ be the edges accepted with at least $\ell$
points active immediately before acceptance. Uniformly for
$2\le\ell\le n$,
\begin{equation}\label{eq:tagged-parent-tail}
 \EE\sum_{e\in\mathcal E_{n,\ell}}
       \bigl[L_e^2+L_e\xi_n(I_e)+\xi_n(I_e)^2\bigr]\le C/\ell.
\end{equation}
\end{lemma}
\begin{proof}
Lemma~\ref{lem:parent-second-moment} gives the assertion for $L_e^2$
by summing $C/m^2$ over $m\ge\ell$.
A proper parent interval spans $j$ consecutive original gaps for
some $1\le j<n$, and contains $j+1$ original points, including its
endpoints. Distinct mergers have distinct proper parent intervals:
after a merger its interior is empty, so the same interval cannot
become a parent again. There is at most one full-circle parent.

If its length is at least $j/(2n)$, its empirical mass is at most
$2j/n\le4L_e$. For the other intervals bound the selected family by
all $n$ circular intervals spanning $j$ original gaps. Each has a
$\operatorname{Beta}(j,n-j)$ length. The order-statistic identity
and the exponential bound for a binomial tail give
\[
 \PP\{\operatorname{Beta}(j,n-j)<j/(2n)\}
 =\PP\{\operatorname{Bin}(n-1,j/(2n))\ge j\}
 \le (e^{1/2}/2)^j.
\]
For example apply Markov's inequality to $2^{\operatorname{Bin}}$.
The expected sum of squared empirical masses of all exceptional
intervals is therefore at most
\[
 \sum_{j=1}^{n-1}n\frac{4j^2}{n^2}(e^{1/2}/2)^j\le C/n.
\]
Thus the selected empirical square sum is bounded by sixteen times
the parent-length square sum plus $C/n$. The full-circle parent
already has empirical mass equal to its length. Finally use
$2L_e\xi_n(I_e)\le L_e^2+\xi_n(I_e)^2$.
\end{proof}

\begin{proposition}[A uniform reduction to early single-probe captures]
\label{prop:early-history-linearization}
Let $B$ have $M$ iid uniform points, let two independent probe sets
have sizes $d,b$, and put $q=d+b\ge1$. Choose a parity-compatible
integer $m$ with $2q+4\le m<M$. Then
\begin{equation}\label{eq:early-history-linearization}
 \left|\EE H(B+X_d,B+Y_b)
          -2q\sum_{\substack{j>m\\j\equiv M\ ({\rm mod}\ 2)}}
                         \EE W_{M,j}\right|
 \le C\left(m+q+\frac{q^2}{m-q}\right).
\end{equation}
The two probe sets are denoted by $X_d,Y_b$; $W_{M,j}$ is the
single-probe capture mass in Proposition~\ref{prop:history-capture}.
\end{proposition}
\begin{proof}
Stop the background immediately after its count reaches $m$, at
separation $\tau$. In every augmented input using at most $q$ probes,
the active count at $\tau$ lies between $m-q$ and $m+q$.
Restrict all matching vectors to edges whose distance is at most
this same $\tau$. This common restriction is essential: separate
stopping counts in the different inputs would not preserve mixed
differences edge by edge.

The parent-witness statement in Lemma~\ref{lem:matching-rectangle}
also covers edges incident to probes. If an edge is incident to
just one probe, a nonzero mixed difference means flipping the other
probe kills that edge in a positive corner. Its closed parent
contains the incident probe and, by exterior protection, the other
probe. An edge joining the two probes is a witness itself.
Thus there is no additive constant for incident edges in the
localized bound.

Every positive-corner witness for an edge of distance at most $\tau$
occurs at active count at least $m-q$. Discarding the random cutoff
and retaining only this necessary count condition, the uniform-tag
argument in Lemma~\ref{lem:matching-rectangle}, now with
Lemma~\ref{lem:tagged-parent-tail}, bounds each expected restricted
mixed difference by $C/(m-q)$. After discarding the cutoff the
condition is a function of the unlabelled corner sample, so a present
probe really is a uniform tag; no independence from $\tau$ is used.
All corner sizes are between $M$ and $M+q$.

Discrete telescoping of the restricted vectors now has expected
remainder $Cq^2/(m-q)$. The same bound controls overlaps of their
single-probe increment supports. Indeed for any Boolean coordinate
$f$ on a two-probe square, if both increments from the base corner
are nonzero, then
$|f_{11}-f_{10}-f_{01}+f_{00}|\ge1$.
The support-overlap inequality from
Theorem~\ref{thm:sharp-small-edits} consequently gives
\[
 \left|\EE H_{\le\tau}(B+X_d,B+Y_b)-q h_M^\tau\right|
          \le Cq^2/(m-q),
\]
where $h_M^\tau$ is the expected single-addition history distance
restricted to this background cutoff.

For a single addition, the two active counts at $\tau$ differ by
exactly one. Counting edges gives the restricted history distance as
twice the number of missing background edges before $\tau$, plus
zero or one. Thus
\[
 h_M^\tau=2\sum_{j>m,\ j\equiv M\ ({\rm mod}\ 2)}\EE W_{M,j}+O(1).
\]
After $\tau$ the two augmented copies accept at most $m+q/2$ edges
in total. This bounds their remaining symmetric difference and
proves \eqref{eq:early-history-linearization}.
\end{proof}

\begin{theorem}[Sharp history means for every sublinear edit count]
\label{thm:sharp-sublinear-history}
For every $1\le k=k_N=o(N)$ and odd $N\to\infty$,
\begin{equation}\label{eq:sharp-sublinear-resampling}
 \EE H_N(k)\sim\EE H_N^{\rm lab}(k)
        \sim(1+e^{-\gamma})\,k\log(N/k).
\end{equation}
For any mixture of $d$ uniform deletions and $b$ fresh additions,
with $q=d+b\ge1$ and $q=o(N)$, the expected spatial history cost is
\begin{equation}\label{eq:sharp-sublinear-mixed}
 \frac{1+e^{-\gamma}}2\,q\log(N/q)\,(1+o(1)).
\end{equation}
The mixed-edit assertion permits either initial parity.
Sections~\ref{sec:dense-history}--\ref{sec:history-concentration}
treat dense resampling and relative concentration;
Section~\ref{sec:late-history-variance} proves the centered CLT
for every growing sublinear edit count.
\end{theorem}
\begin{proof}
The full corner law, the decorated-gap capture identity, and
Lemma~\ref{lem:truncated-matching-length} give
\begin{equation}\label{eq:truncated-capture-asymptotic}
 \sum_{j>m,\ j\equiv M\ ({\rm mod}\ 2)}\EE W_{M,j}
       \sim\frac{1+e^{-\gamma}}4\log(M/m)
       \qquad(m=o(M),\ m\ge2).
\end{equation}
Here is the required uniform harmonic-sum argument. On any central
range where $j\to\infty$ and $j/M\to0$, the mean corner mass is
$1/2+o(1)$ uniformly, by the sequential formulation of
Lemma~\ref{lem:full-corner-law}. Choose its lower and upper endpoints
so that the omitted harmonic sums are $o(\log(M/m))$; this is
possible because $M/m\to\infty$. For bounded $m$, let the lower
endpoint diverge sufficiently slowly. The corner contribution is
therefore $\tfrac14\log(M/m)(1+o(1))$, the separation contribution
is $e^{-\gamma}\log(M/m)/4$, and
$\sum_{j>m}\EE s_{M,j}/(j-1)=O(1/m)$.
This proves \eqref{eq:truncated-capture-asymptotic}.

For two probe sets on a common background choose $m$ of the proper
parity with $4q+4\le m\le4q+5$. For $q=o(M)$ this is admissible,
and Proposition~\ref{prop:early-history-linearization} has error
$O(q)$. Since $\log(M/m)=\log(M/q)+O(1)$, the expected cost is
$((1+e^{-\gamma})/2)q\log(M/q)(1+o(1))$.
For resampling, take $M=N-k$, $d=b=k$; for a mixed edit take
$M=N-d$. In either case $M\sim N$. The resampling label correction
is at most $2k$ and is negligible, proving both assertions.
\end{proof}

\subsection{A deterministic dense-edit profile and its matching spectrum}
\label{sec:dense-history}

\status{Working proof: a law of large numbers at positive edit density}

\begin{lemma}[Uniformly bounded label-correction moments]
\label{lem:history-label-correction}
For every odd $N$ and $1\le k\le N$,
\[
 0\le\EE[H_N(k)-H_N^{\rm lab}(k)]\le4.
\]
Moreover $\EE[(H_N(k)-H_N^{\rm lab}(k))^2]\le32$.
The pathwise bound remains $2k$.
\end{lemma}
\begin{proof}
Condition on the unchanged labelled points and on the two unlabelled
sets of refreshed locations. The assignments of the $k$ refreshed
labels are independent uniform permutations in the two copies.
An unchanged point paired to a refreshed point in both matchings
has the same partner label with probability $1/k$.
There are at most $k$ such unchanged points, so these common label
edges have conditional expected count at most one.
If the two matchings have $a,b$ edges with both endpoints refreshed,
the conditional expected common-label count of that type is
$ab/\binom{k}{2}\le1$ for $k\ge2$, since $a,b\le\lfloor k/2\rfloor$.
For $k=1$ this type is absent. Every extra common label edge reduces
the symmetric difference by two, proving the mean bound.

For the second moment, let $J_1,J_2$ be the respective counts of
these two types. Given the two spatial matchings, fix the first
label assignment and randomize the second. Any $r$ distinct
unchanged-to-refreshed coincidences prescribe $r$ different
images of this permutation, so
$\EE(J_1(J_1-1)\mid\text{geometry})\le1$ and $\EE J_1^2\le2$.
For $k\ge4$ the other type has conditional factorial second moment
\[
 \EE(J_2(J_2-1)\mid\text{geometry})
    =\frac{4a(a-1)b(b-1)}{k(k-1)(k-2)(k-3)}
    \le\frac23.
\]
Here the $r=2$ calculation chooses two disjoint edges in each
matching and their two possible orientations. For $k<4$,
$J_2\le1$. Thus $\EE J_2^2\le2$ in all cases, and
$H_N(k)-H_N^{\rm lab}(k)=2(J_1+J_2)$ gives the bound $32$.
\end{proof}

\begin{theorem}[The macroscopic history profile]
\label{thm:dense-history-profile}
There is a continuous increasing concave function
$\mathfrak h:[0,1]\to[0,1]$, with $\mathfrak h(0)=0$ and
$\mathfrak h(1)=1$, such that whenever $k_N/N\to\theta$,
\begin{equation}\label{eq:dense-history-LLN}
 \frac{H_N(k_N)}N\longrightarrow\mathfrak h(\theta),\qquad
 \frac{H_N^{\rm lab}(k_N)}N\longrightarrow\mathfrak h(\theta)
                         \quad\hbox{in }L^2.
\end{equation}
In fact the corresponding squared errors with
$\mathfrak h(k/N)$ tend to zero uniformly over $0\le k\le N$.
At small edit density,
\begin{equation}\label{eq:dense-history-small-theta}
 \mathfrak h(\theta)\sim(1+e^{-\gamma})\theta\log(1/\theta)
                      \qquad(\theta\downarrow0).
\end{equation}

One intrinsic description uses three independent Poisson processes
on $\mathbb R$, of intensities $1-\theta,\theta,\theta$.
The first is common to both inputs, and each other process belongs
to one input. For $\theta<1$, let $\varrho(\theta)$ be the Palm
probability that a common point has the same matching partner in
both completed infinite-line processes. Then
\begin{equation}\label{eq:dense-history-palm}
 \mathfrak h(\theta)=1-(1-\theta)\varrho(\theta).
\end{equation}
The function is specified by this law, not by an explicit evaluation
of $\varrho$. Section~\ref{sec:dense-history-CLT} separately proves
the positive-density centered CLT and a positive variance formula.
\end{theorem}
\begin{proof}
Fix a finite separation $s>0$ on the line. A gap longer than $s$
in the union of the three inputs separates all edges considered
before $s$. Its components are finite almost surely. Their point
counts have geometric tails, uniformly for $0\le\theta\le1$,
because the union rate lies in $[1,2]$. Component lengths also have
exponential tails. Run the greedy scan independently in each
component, in each of the two input colours. These constructions
are consistent as $s$ increases, and define the infinite histories.

The number of changed edges before $s$ is an additive reward over
these renewal components, bounded by their union point count.
The iid component lengths and rewards have finite second moments
(indeed exponential moments). Centering each reward by its mean
per unit length, summing the independent centered increments, and
bounding the first and last residual components proves a renewal
reward law in $L^2$ on an interval of length $L\to\infty$.
Denote the limiting changed-edge density by $\mathfrak h_s(\theta)$.
Closing the interval into a circle alters only the components meeting
the seam, of bounded expected size, so it gives the same limit.
The event of no separating gap has exponentially small probability
and its contribution also vanishes after division by $L$.

For fixed $s$, $\mathfrak h_s$ is continuous in $\theta$.
Couple two intensity triples by their common Poisson parts. The
intensity of differing points is at most $3|\theta-\theta'|$.
Their effects are confined to components of the union of both
triples, whose rate is at most three. The expected point count
of a component seen from one added point is bounded by a constant
depending only on $s$. Campbell's formula therefore bounds the
change of reward density by $C_s|\theta-\theta'|$.
The component moment bounds are uniform, so the same $L^2$ law
holds for any sequence $\theta_L\to\theta$.

Each marginal input has rate one. Its active point density at
separation $s$ is $q(s)$ from Theorem~\ref{thm:early-fluid}; this
also follows by taking the line limit in the finite-Poisson renewal
density. Thus every marginal point is eventually matched almost
surely, since $q(s)\to0$. At most half the active points can give
further edges in each marginal. It follows that
\[
 0\le\mathfrak h(\theta)-\mathfrak h_s(\theta)\le q(s),
 \qquad \mathfrak h=\lim_{s\to\infty}\mathfrak h_s.
\]
This convergence is uniform in $\theta$, so $\mathfrak h$ is
continuous. The same tail bound in expectation holds for a large
Poisson circle, with $q(s)+o(1)$ on the right. It follows from
the microscopic count law applied to each marginal. Sending first
the circle length to infinity and then $s$ to infinity proves the
Poisson full-history law in $L^1$.

Return to exactly $N-k,k,k$ points of the three types on the unit
circle. Couple each type to a Poisson count with that mean using
prefixes of one iid uniform sequence. The expected sum of count
differences is $O(\sqrt N)$. Every intermediate marginal sample
is iid conditional on its size, and the single-addition history
mean is at most $C(1+\log(n+2))$. The triangle inequality therefore
bounds the expected change in the history distance by
$C\sqrt N\log(N+2)=o(N)$; Cauchy--Schwarz handles the random sizes,
using their Poisson second moments. After rescaling the circle to
length $N$, the Poisson input is precisely the one above with
$\theta=k/N$. This proves the fixed-count $L^1$ law. Since
$0\le H_N(k)/N\le1$, it is an $L^2$ law too.
Lemma~\ref{lem:history-label-correction} transfers it to label edges.
Compactness of $[0,1]$ and continuity of $\mathfrak h$ give the
asserted uniformity by taking a convergent subsequence of any
putative violating edit fractions.

The finite label-history means are increasing and concave by
Lemma~\ref{lem:history-lower}; their limit has the same properties.
The endpoint values follow from identical inputs and disjoint
inputs, respectively. Each infinite marginal matching has edge
intensity $1/2$. Common matching edges have intensity
$(1-\theta)\varrho(\theta)/2$ by the Palm counting identity,
which proves \eqref{eq:dense-history-palm}.

Finally, if \eqref{eq:dense-history-small-theta} failed along
$\theta_j\downarrow0$, choose odd $N_j$ large enough that
$k_j=\lfloor\theta_jN_j\rfloor\ge1$, the rounding error is negligible,
and the mean convergence in \eqref{eq:dense-history-LLN} has error
$o(\theta_j\log(1/\theta_j))$. Then $k_j=o(N_j)$ and
Theorem~\ref{thm:sharp-sublinear-history} contradicts that failure.
\end{proof}

\begin{corollary}[An integer-valued matching spectrum]
\label{cor:matching-spectrum}
There is a unique probability vector $(\pi_j)_{j\ge2}$ such that
\begin{equation}\label{eq:matching-spectrum}
 1-\mathfrak h(\theta)=\sum_{j\ge2}\pi_j(1-\theta)^j,
                         \qquad 0\le\theta\le1.
\end{equation}
The normalized finite matching spectral weights converge to this
vector in $\ell^1$. These are unscaled integer coordinate orders,
not the survivor spectrum $\nu$ on scaled orders $j/N$.
In particular $\mathfrak h$ is real analytic on $(0,1)$,
$\mathfrak h(\theta)\ge2\theta-\theta^2$, and
\begin{equation}\label{eq:matching-spectrum-truncated-mean}
 \sum_{2\le j\le t}j\pi_j\sim(1+e^{-\gamma})\log t.
\end{equation}
This truncated-moment law does not assert an asymptotic for each
individual $\pi_j$ or for its tail probabilities.
\end{corollary}
\begin{proof}
Let $F$ be the vector of label-edge indicators and $F_S$ its
orthogonal product-space components. Put
\[
 a_{N,j}=\frac2N\sum_{|S|=j}\EE\|F_S\|^2,\qquad j\ge1.
\]
The finite spectral formula is
\[
 \frac{\EE H_N^{\rm lab}(k)}N
  =\sum_{j\ge1}a_{N,j}
        \left(1-\frac{\binom{N-k}{j}}{\binom Nj}\right),
 \qquad \sum_{j\ge1}a_{N,j}=(1-1/N)^2.
\]
Conditioning any label-edge indicator on one coordinate leaves its
mean constant, by common rotation invariance. Hence $a_{N,1}=0$.
Take a subsequence on which each fixed $a_{N,j}$ converges to $\pi_j$.
For $k/N\to\theta>0$, the avoidance factor is bounded by
$(1-k/N)^j$ and tends to $(1-\theta)^j$ for fixed $j$.
The tail of the agreement sum is therefore uniformly geometric.
The limiting identity is
$1-\mathfrak h(\theta)=\sum_{j\ge2}\pi_j(1-\theta)^j$.
Continuity at zero and monotone convergence give $\sum_j\pi_j=1$,
excluding escaping spectral mass. Uniqueness of power-series
coefficients gives convergence of the whole sequence, and convergence
of nonnegative total masses upgrades it to $\ell^1$ convergence.
The analytic and lower-bound assertions follow directly.

Write $c_{\rm H}=1+e^{-\gamma}$. Concavity and the secant slopes
between $\theta$ and $(1\pm\epsilon)\theta$, using
\eqref{eq:dense-history-small-theta}, give
$\mathfrak h'(\theta)\sim c_{\rm H}\log(1/\theta)$.
Put $L(v)=\sum_j j\pi_je^{-vj}$; the substitution
$v=-\log(1-\theta)$ gives $L(v)\sim c_{\rm H}\log(1/v)$.
For $0<\delta<1$, take $v=t^{-1-\delta}$ to bound the truncated
sum above by $e^{vt}L(v)$, with limiting upper coefficient
$c_{\rm H}(1+\delta)$.
Take $v=t^{-1+\delta}$ for the lower bound. The omitted part is
at most $\sup_{j>t}j e^{-vj}\sum_j\pi_j=o(1)$, so the limiting
lower coefficient is $c_{\rm H}(1-\delta)$.
Send $\delta\downarrow0$ to obtain
\eqref{eq:matching-spectrum-truncated-mean}.
\end{proof}

\begin{proposition}[The quadratic endpoint coefficient]
\label{prop:matching-quadratic-endpoint}
With $q$ from Theorem~\ref{thm:early-fluid},
\begin{equation}\label{eq:matching-pi-two}
 \pi_2=\frac12\int_0^\infty q'(x)^2\,dx
      =\int_0^\infty e^{-2\operatorname{Ein}(u)}
                       \left(\frac{1-e^{-u}}u\right)^2du
                       \in(1/3,1/2).
\end{equation}
Consequently $1-\mathfrak h(\theta)
 =\pi_2(1-\theta)^2+O((1-\theta)^3)$ as $\theta\uparrow1$.
The certified decimal enclosure below evaluates this exact integral.
\end{proposition}
\begin{proof}
In a rate-one Poisson input on the line, insert points at $0,x$ and
let $p(x)$ be the probability that they form an accepted edge.
Campbell's formula identifies $p(x)\,dx$ as the accepted-edge
intensity with separation in $dx$. The active intensity therefore
satisfies $q(s)=1-2\int_0^s p(x)\,dx$, so
\[
 p(x)=-q'(x)/2
     =q(x)\frac{1-e^{-2x}}{2x}
\]
almost everywhere. In particular $\int p=1/2$ and $0<p(x)<1$
for positive $x$.

Put $\varepsilon=1-\theta$, and insert $0,x$ as common points in
the coupled inputs. Let $p_\theta(x)$ be the probability that both
matchings accept that edge. Then
\[
 1-\mathfrak h(\theta)
       =2\varepsilon^2\int_0^\infty p_\theta(x)\,dx.
\]
The decision for this edge uses only separation at most $x$.
The finite-component construction consequently shows, as
$\theta\uparrow1$, that $p_\theta(x)\to p(x)^2$: the background
common intensity vanishes, leaving independent rate-one backgrounds
with the same two forced endpoints. Also $p_\theta(x)\le p(x)$,
since each marginal has rate one. Dominated convergence and
\eqref{eq:matching-spectrum} identify $\pi_2=2\int p(x)^2dx$.
The substitution $u=2x$ gives \eqref{eq:matching-pi-two}.
For the lower bound, $\operatorname{Ein}(u)<u$ and
$(1-e^{-u})/u>e^{-u/2}$ for $u>0$, so the final integrand exceeds
$e^{-3u}$. For the upper bound put
$f(u)=e^{-2\operatorname{Ein}(u)}$ and $g(u)=(1-e^{-u})/u$.
Then $f'=-2fg$, $g'<0$, and integration of $(fg)'$ gives
$\pi_2=(1+\int_0^\infty f g')/2<1/2$.
The remaining spectral mass contributes at most $\varepsilon^3$.
\end{proof}

The script \path{certify_matching_spectrum.py} gives the enclosure
below. Its exact integer endpoints are recorded in
\path{matching-spectrum-certificate.json}.
\begin{equation}\label{eq:matching-pi-two-certificate}
 0.388474956436\le\pi_2\le0.388474956440.
\end{equation}
It uses the same outward arithmetic with denominator $10^{50}$ as
the survivor-bound certificate, but evaluates a different constant.
On $[0,32]$, the alternating entire series for $\operatorname{Ein}$
is truncated at degree $128$ with its first omitted-term bound.
Composite Simpson integration uses $65536$ intervals.

Here is a derivative bound independent of the implementation.
Let $V_1,V_2$ be independent Dickman variables with Laplace transform
$e^{-\operatorname{Ein}(u)}$, as in
Lemma~\ref{lem:positive-renewal}, and let $U_1,U_2$ be independent
uniform variables on $[0,1]$. The integrand is the Laplace transform
of $T=V_1+V_2+U_1+U_2$. Its first four cumulants are
$3,7/6,2/3,29/60$, obtained by expanding the elementary Laplace
transforms. Hence
\[
 \sup_{u\ge0}|(\EE e^{-uT})^{(4)}|
     \le\EE T^4=4697/30<157.
\]
The Simpson error is at most $32\cdot157/(180\cdot2048^4)$.
For the tail, $\operatorname{Ein}(u)=\gamma+\log u+E_1(u)$ and
$0<E_1(u)\le e^{-u}/u$ give, with $R=32$,
\[
 \frac{e^{-2\gamma}}{3R^3}
      e^{-2e^{-R}/R}(1-e^{-R})^2
 \le\int_R^\infty e^{-2\operatorname{Ein}(u)}
                   \left(\frac{1-e^{-u}}u\right)^2du
 \le\frac{e^{-2\gamma}}{3R^3}.
\]
The $129$ separate $100$-digit decimal comparisons validate the
implementation; the displayed bounds justify the certificate.
This enclosure says nothing new about the survivor-noise constant
$p_\infty$.

\subsection{Concentration at every growing sublinear edit count}
\label{sec:history-concentration}

\status{Working proof: relative concentration; sparse centered limits follow below}

\begin{lemma}[Concentration of the truncated total capture mass]
\label{lem:capture-concentration}
Let $\mathcal W_{M,m}$ be the total single-probe capture mass over
events with active count greater than $m$. For parity-compatible
$m\ge2$ with $m=o(M)$,
\begin{equation}\label{eq:capture-concentration}
 \frac{\mathcal W_{M,m}}{\log(M/m)}
           \longrightarrow\frac{1+e^{-\gamma}}4
                                      \quad\hbox{in }L^1.
\end{equation}
\end{lemma}
\begin{proof}
Use deferred revelation of the cyclic order of decorated gaps.
Initially reveal only their unordered lengths. At each step reveal
the ordered neighbours of the minimum, merge the three decorated
gaps, and retain their full internal trees in the new mark.
Conditional on the entire revelation so far, the remaining cyclic
order is uniform: conditioning on the two neighbours selects one
ordered three-gap block, and contracting that block is a bijection
onto cyclic orders of the remaining marks. The revealed internal
trees impose no further order among the remaining blocks.
Absolute rotation is irrelevant to all quantities here.

In this filtration, just before count $j$ is reduced, $s_{M,j}$
and $A_{M,j}$ are known, and the decorated-gap identity holds
conditionally:
\[
 \EE[W_{M,j}\mid\mathcal F_j]
     =s_{M,j}+\frac{A_{M,j}-s_{M,j}}{j-1}.
\]
The centered differences are martingale differences in forward
merger time. Their sum above $m$ has second moment at most
\[
 \sum_{j>m}\EE W_{M,j}^2
          \le\sum_{j>m}\EE L_{M,j}^2\le C/m.
\]
Thus it remains to analyze the predictable sum.
Lemma~\ref{lem:full-corner-law} gives
$\EE|A_{M,j}-1/2|\to0$ throughout $j\to\infty$, $j/M\to0$.
The count inversion gives $j s_{M,j}\to c=(2e^\gamma)^{-1}$ in
probability in the same regime, hence in $L^1$ because
$0\le j s_{M,j}\le1$.
Use a central range as in the proof of
\eqref{eq:truncated-capture-asymptotic}. The triangle inequality
for the weighted $L^1$ errors gives predictable sum
$((c+1/2)/2)\log(M/m)+o_{L^1}(\log(M/m))$.
The omitted harmonic ranges are negligible by the deterministic
bounds on $A_{M,j}$ and $j s_{M,j}$; the extra
$\sum s_{M,j}/(j-1)$ is $O(1/m)$.
Adding the martingale sum proves \eqref{eq:capture-concentration}.
\end{proof}

\begin{theorem}[A law of large numbers for growing sparse edits]
\label{thm:sparse-history-LLN}
If $k_N\to\infty$ and $k_N=o(N)$, then
\begin{equation}\label{eq:sparse-history-LLN}
 \frac{H_N(k_N)}{k_N\log(N/k_N)}
       \longrightarrow1+e^{-\gamma}\quad\hbox{in }L^1,
\end{equation}
and the same holds for label edges. For a random mixture of
$d$ deletions and $b$ additions with $q=d+b\to\infty$,
$q=o(N)$, and either initial parity,
\[
 \frac{H_N(d,b)}{q\log(N/q)}
          \longrightarrow\frac{1+e^{-\gamma}}2
                                      \quad\hbox{in }L^1.
\]
In particular the matching-history lower bound holds with
probability tending to one at every growing sublinear edit count.
Section~\ref{sec:fixed-history-LLN} separately treats fixed edit
counts. No centered fluctuation law is asserted here.
\end{theorem}
\begin{proof}
We first record a uniform second-moment bound for a single addition:
\begin{equation}\label{eq:one-history-second-moment}
 \EE H(B_M,B_M+Z)^2\le C(1+\log(M+2))^2
\end{equation}
for either background parity. For even $M$, sum the ancestry tail
in Lemma~\ref{lem:insertion-height}, using $H\le2D_M$.
For odd $M$, the final singleton has one full-circle gap tree.
In the Poisson description, condition on that singleton at cutoff
$1/3$; its interior is a Poisson interval conditioned on shielding,
with root length one. The same uniform daughter contraction in
that lemma gives $\PP(D\ge d+1)\le\lambda\rho_{\rm h}^d$ before
conditioning on the initial count. Conditioning at $\lambda=M$
costs at most $e\sqrt M$, exactly as there. A single addition has
cost at most $2D+1$, the extra edge matching the final discrepancy
to the original singleton. This proves
\eqref{eq:one-history-second-moment}; bounded sizes are immediate.

Use a common size-$M=N-d$ background, with $q=d+b$ independent
probes, and choose $m\asymp4q$ with the background parity.
Let $\tau$ be its cutoff on reaching $m$, and let
$G_B(z)=H_{\le\tau}(B,B+z)$. The proof of
Proposition~\ref{prop:early-history-linearization} gives the stronger
random-variable comparison
\begin{equation}\label{eq:history-random-linearization}
 \EE\left|H(B+X_d,B+Y_b)-\sum_{i=1}^qG_B(Z_i)\right|\le Cq.
\end{equation}
Indeed the telescoping remainder and the support-overlap loss were
bounded in expected absolute value, and the late history cost is
bounded deterministically by $m+q/2$. Conditional on $B$, the
$G_B(Z_i)$ are iid, with mean
\[
 \mu_B=2\mathcal W_{M,m}+\delta_B,\qquad 0\le\delta_B\le1.
\]
Also $G_B(Z)\le H(B,B+Z)$, so
\eqref{eq:one-history-second-moment} and conditional variance give
\[
 \EE\left|\sum_{i=1}^qG_B(Z_i)-q\mu_B\right|
       \le C\sqrt q\,(1+\log(M+2)).
\]
Relative to $q\log(M/q)$ this tends to zero whenever
$q\to\infty$ and $q=o(M)$. To check all ranges, if $q\le\sqrt M$
then $\log(M/q)\ge\tfrac12\log M$; if $q>\sqrt M$, the ratio
is at most $C\log M/M^{1/4}$ for large $M$.
Lemma~\ref{lem:capture-concentration} gives
$\mu_B/\log(M/q)\to(1+e^{-\gamma})/2$ in $L^1$.
The error in \eqref{eq:history-random-linearization} is negligible
since $\log(M/q)\to\infty$. This proves the mixed-edit law.
Take $d=b=k$ and use the bounded mean label correction to obtain
\eqref{eq:sparse-history-LLN}.
\end{proof}

\begin{corollary}[Relative concentration for all growing resampling counts]
\label{cor:all-growing-history-concentration}
For every $k_N\to\infty$ with $k_N\le N$,
\[
 \frac{H_N(k_N)}{\EE H_N(k_N)}\longrightarrow1
                                      \quad\hbox{in }L^1,
\]
and likewise for label edges. The same assertion holds under
independent-coordinate resampling when its expected count $Ns_N$
tends to infinity.
\end{corollary}
\begin{proof}
Every subsequence has a further subsequence on which $k_N/N$ has
a limit. At a positive limit use
Theorem~\ref{thm:dense-history-profile}, whose limiting mean is
positive. At zero use Theorem~\ref{thm:sparse-history-LLN} and the
sharp mean. This gives the first assertion in $L^1$.

For independent resampling put $K\sim\operatorname{Bin}(N,s_N)$,
$\mu=Ns_N$, and $f_N(k)=\EE H_N^{\rm lab}(k)$, extended linearly
between integers. This function is increasing and concave, with
$f_N(0)=0$, so for $0<a\le1\le b$,
$a f_N(\mu)\le f_N(a\mu)\le f_N(\mu)$ and
$f_N(\mu)\le f_N(b\mu)\le b f_N(\mu)$ whenever the arguments are
in range. Binomial concentration gives $f_N(K)/f_N(\mu)\to1$
in probability. Its upper bound by $\max(1,K/\mu)$ is uniformly
integrable, so this convergence is in $L^1$ and
$\EE f_N(K)\sim f_N(\mu)$.

The already-proved conditional relative errors are uniformly small
for all sufficiently large deterministic counts: otherwise choose
a violating sequence of counts tending to infinity. Averaging over
$K$ preserves this bound, while any fixed bounded range costs at
most its relative mean weight and has probability tending to zero.
Thus the random-count history divided by $\EE f_N(K)$ tends to one
in $L^1$. The uniform mean label correction transfers the assertion
to spatial edges, since the chord bound gives
$\EE f_N(K)\ge\mu(1-1/N)^2\to\infty$.
\end{proof}

\subsection{Dilation ergodicity and the fixed-edit history law}
\label{sec:fixed-history-LLN}

\status{Working proof: relative concentration also for fixed nonzero edits}

\begin{lemma}[A stationary line entrance law and logarithmic capture bands]
\label{lem:line-entrance-dilation}
There is a stationary full entrance history on $\mathbb R$ whose
configuration at separation $r$ is the stationary renewal process
with gap density $h_r(x)=r^{-1}f(x/r)$ and intensity $c/r$,
where $c=(2e^\gamma)^{-1}$. Conditional on its gaps, their full past
trees are independent with the shielded kernels already constructed.
This history law is invariant under simultaneous dilation of space
and separation. Dilation by $e$ has no nonconstant $L^2$
eigenfunctions and, in particular, is ergodic.

Insert one source at the origin using the same-origin maps.
Let $J(a,b)$ count background edges of separation in $(a,b]$ missing
from the inserted history. Then, for $0<a<b<\infty$,
\begin{equation}\label{eq:line-capture-mean}
 \EE J(a,b)=a_{\rm H}\log(b/a),\qquad
 a_{\rm H}=\frac{1+e^{-\gamma}}4,
\end{equation}
and $J(e^{-n},1)/n\to a_{\rm H}$ almost surely and in $L^1$.
This is ergodicity of line dilations, not ergodicity of the circle
refresh kernel $Q$.
\end{lemma}
\begin{proof}
At cutoff one, take a stationary renewal process with iid gaps
having density $f$, and attach conditionally independent full gap
trees. The gap covering the origin has the size-biased gap law,
with a uniform origin position inside it. This gives the past at
every positive cutoff. For later finite cutoffs, a gap exceeding
that cutoff separates the greedy scan; such gaps occur in both
directions almost surely, so every relevant component is finite.
This defines the future consistently.

To check the stated laws and consistency of the past, one may use
the circle entrance laws on circumferences tending to infinity.
For any finite consecutive string of gaps, the outside renewal
factor tends to $c$, after scaling to cutoff one, by
Lemma~\ref{lem:factorial-renewal}. The resulting rooted law is
exactly the stationary renewal law just specified. Conditional
finite gap trees have the same kernels, and deeper trees are
obtained by their consistent product construction. At any smaller
positive cutoff only finitely many descendants occur in a bounded
root interval. Finite components justify passage to larger cutoffs.
Thus all finite collections of positive-cutoff local histories have
the prescribed limits. Scaling the root gap law and split kernels
now proves invariance under simultaneous space-time dilation.

The complete marked renewal process at cutoff one is translation
ergodic. Indeed a translation-invariant event does not depend on
the root's position inside its gap; at successive renewal boundaries
it is an invariant event of the iid sequence of marked gaps, hence
has probability zero or one. The whole entrance history is a
measurable function of this marked process, so it is translation
ergodic too. Translation acts strongly continuously on $L^2$:
bounded continuous functions of finitely many local configurations
are dense, and translation is continuous for such functions.

Let $D$ be dilation by $e$, acting unitarily on complex $L^2$, and
$T_x$ translation by $x$. If $Dg=\lambda g$ with $|\lambda|=1$,
the phase cancels in the affine conjugacy, giving
\[
 \|T_xg-g\|_2=\|T_{e^{-n}x}g-g\|_2\longrightarrow0.
\]
Thus every dilation eigenfunction is translation invariant and
therefore constant, with eigenvalue one. This proves the assertion.

The same-origin response is defined inside the gap containing zero
at each cutoff. The deterministic query budgets and the compatible
cutoff readouts give a consistent one-point perturbation. On a
fixed band $(a,b]$ it affects only finitely many points: at cutoff
$a$, take the gap containing the source and the neighbouring strings
up to gaps longer than $b$. Their lengths and point counts have
finite moments, and no event before $b$ crosses those barriers.

At separation $r$, the collision intensity per unit spatial length
and per unit separation is
\[
 (c/r)h_r(r)=c/(2r^2).
\]
This is the gap-density flux through the cutoff; $f(1)=1/2$.
Equivalently multiply the reverse-splitting rate
\eqref{eq:reverse-rate} by the stationary parent-gap intensity;
the parent density cancels and leaves
$(c/r)h_r(a)h_r(r)h_r(c')\,da\,dc'\,dr$.
Its two neighbouring gaps have independent density $h_r$, with
their conditional shielded marks. The central gap contributes
captured origin mass $r$. By the corner sum rule
\eqref{eq:corner-reward-sum-rule}, the two neighbouring contributions
have total mean $2r\int\psi(y)f(y)\,dy=r e^\gamma$.
The mean captured mass per collision is therefore $r(1+e^\gamma)$.
Integrating origin locations and applying stationarity gives
\eqref{eq:line-capture-mean} by mass transport. All terms are
integrable, since capture mass is bounded by the three-gap length.

The costs in successive unit logarithmic bands are translates
under $D$ of the integrable observable $J(1,e)$. Dilation ergodicity
and the ergodic theorem prove the final assertion, also in $L^1$
by truncation of that observable.
\end{proof}

\begin{lemma}[Local insertion histories on mesoscopic ratio bands]
\label{lem:local-history-band}
Let $B_M$ be an iid circle background and $Z$ an independent uniform
source. Write $J_M(a,b)$ for its number of missing background edges
with separation in $(a,b]$ after inserting $Z$.
If $r_M\to0$ and $Mr_M/\log(1/r_M)\to\infty$, then, for every
fixed $0<t_0<\cdots<t_l<\infty$,
\begin{equation}\label{eq:local-history-band}
 \bigl(J_M(r_Mt_{i-1},r_Mt_i)\bigr)_{i=1}^l
       \ \Longrightarrow\ \bigl(J(t_{i-1},t_i)\bigr)_{i=1}^l.
\end{equation}
This compares local spatial histories and their edge counts, not
raw integer coordinate labels.
\end{lemma}
\begin{proof}
Rotate the independent source to zero, and start at cutoff
$t_0r_M$. In the line entrance law reveal the gap containing zero
and neighbouring gaps out to barriers longer than $t_lr_M$ on
both sides. For fixed $t_l/t_0$, the number of intervening renewal
gaps is tight, since a gap exceeds the barrier threshold with
positive probability. Their total length divided by $r_M$ is tight
too. Choose a bound $D$ for both the number and scaled total length.
Inside this finite string reveal depth $d$ of the gap trees, and
require terminal cutoffs at least $\delta r_M$.

We spell out the finite-count comparison, since unmarked local
convergence would not determine the inserted history. For a fixed
consecutive string and its revealed trees, the Poisson density is
the product of its terminal factors $w_{t_i}^\lambda(x_i)$ times
$U_{t_0r_M}^\lambda(1-\sum x_i)$ for the outside string.
Integrate the left endpoint of the source gap over the interval
that makes this gap cover zero. This roots the string uniquely;
there is no additional cyclic-root factor. Extract
$M![\lambda^M]e^\lambda$.

On the compact geometry just chosen, expand the local factors and
the outside renewal factor as in
\eqref{eq:local-marked-coefficient-error}. The local components
number at most $D/\delta$, because every length is at least
$\delta r_M$ and their sum is at most $Dr_M$.
Each terminal factor costs at most $C_{D,d,\delta}/r_M$ in the
absolute expansion, canceled by rescaling its local reference
variables; the root-position integration only improves the bound.
The outside absolute renewal factor costs at most $C r_M^{-3}$.
Every proper subset of all component lengths has sum between
$\delta r_M$ and $1-\delta r_M$. Thus an upper bound for the
integrated remainder is
\[
 C_{D,d,\delta}r_M^{-6}e^{-2M\delta r_M}=o(1).
\]
The empty and full subsets give exactly the corresponding
circle-entrance geometry of the required parity. Its outside
renewal factor tends, after scaling, to the stationary line factor
$c$, so the local geometry converges in total variation.

The limiting probability of a bad terminal cutoff tends to zero
as $\delta\downarrow0$ for fixed $D,d$, by the inverse-cutoff
moment bound. Normalization of the local rooted law transfers
this bad-event bound to the finite sample, including premature
empty trees and small-count exceptional sectors. The source gap
has a size-biased root law. Its finite moments and
Lemma~\ref{lem:origin-map} bound the expected same-origin
completion error by $C_D\sigma^d$. The deterministic budget on
shared geometry compares arbitrary deeper finite and limiting
completions; no uniform finite-sample contraction is being assumed.

On the remaining event the base configuration and its one-point
response agree exactly in the revealed string, after rescaling.
Their evolution through $t_lr_M$ is contained between the two
barriers and gives the same edge counts. Threshold ties are null
under the limiting continuous geometry laws. Let first $M\to\infty$,
then $D,d\to\infty$ and $\delta\downarrow0$, choosing each
truncation in that order as needed. This proves
\eqref{eq:local-history-band}.
\end{proof}

\begin{theorem}[Fixed-edit laws and uniform relative history concentration]
\label{thm:all-history-concentration}
For one uniform addition or deletion in a size-$M$ background of
either parity,
\begin{equation}\label{eq:one-history-LLN}
 \frac{H}{\log M}\longrightarrow\frac{1+e^{-\gamma}}2
                                      \quad\hbox{in }L^1.
\end{equation}
Consequently, for every $1\le k_N=o(N)$, including fixed counts,
\[
 \frac{H_N(k_N)}{k_N\log(N/k_N)}\longrightarrow1+e^{-\gamma}
                                      \quad\hbox{in }L^1,
\]
also for label edges. Every nonzero random mixture of
$q_N=o(N)$ additions and deletions has the analogous $L^1$ law
with coefficient $(1+e^{-\gamma})/2$ per edit.
Moreover
\begin{equation}\label{eq:uniform-relative-history-concentration}
 \sup_{1\le k\le N}
    \EE\left|\frac{H_N(k)}{\EE H_N(k)}-1\right|
       \longrightarrow0,
\end{equation}
and likewise for label edges. Section~\ref{sec:finite-history-CLT}
separately proves the one-addition/deletion central limit theorem.
\end{theorem}
\begin{proof}
Let $V_M$ be the total missing-background edge count for one
uniform insertion. The sharp mean gives
$\EE V_M/\log M\to a_{\rm H}$ for either parity.
Fix $0<\epsilon<1/2$ and partition separations from
$r_-=M^{-1+\epsilon}$ up to $r_+=M^{-\epsilon}$ into unit
logarithmic bands. Their number $n_M$ satisfies
$n_M/\log M\to1-2\epsilon$.
For a fixed cap $K$, sum the truncated costs
$J_M(r_-e^i,r_-e^{i+1})\wedge K$ over the complete bands,
and call the sum $S_M^K$. Put
$a_K=\EE[J(1,e)\wedge K]$.

Group these bands into blocks of a fixed length $L$.
Lemma~\ref{lem:local-history-band} gives the limiting law of each
block, uniformly over its starting scale in this central range:
any violating sequence still has $r\to0$ and
$Mr/\log(1/r)\to\infty$. Since the capped block costs are bounded,
their expected absolute deviation from $La_K$ converges too.
The triangle inequality over blocks consequently gives
\[
 \limsup_{M\to\infty}\EE|S_M^K/n_M-a_K|
 \le\EE\left|\frac1L\sum_{i=0}^{L-1}
                    (J(e^i,e^{i+1})\wedge K)-a_K\right|.
\]
The incomplete block costs at most $LK/n_M$ and vanishes.
No independence between the finite blocks is required.
Dilation ergodicity sends the right side to zero as $L\to\infty$.
Thus $S_M^K/\log M\to(1-2\epsilon)a_K$ in $L^1$.

The crucial comparison is positive: $0\le S_M^K\le V_M$.
Combining this with the already-proved full mean yields
\[
 \limsup_{M\to\infty}\EE|V_M/\log M-a_{\rm H}|
       \le2\{a_{\rm H}-(1-2\epsilon)a_K\}.
\]
By \eqref{eq:line-capture-mean}, $a_K\uparrow a_{\rm H}$.
Send $K\to\infty$ and then $\epsilon\downarrow0$.
This proves $V_M/\log M\to a_{\rm H}$ in $L^1$ without a
uniform-integrability assumption on the individual finite band costs.
The single-addition history is $2V_M$ plus zero or one;
reversing the iid add-one coupling gives uniform deletion.
This proves \eqref{eq:one-history-LLN}.

For any fixed nonzero edit count, the proof of
Theorem~\ref{thm:sharp-small-edits} bounds in expected absolute value
the difference between the multi-edit history and the sum of its
single-probe costs by $Cq^2$. Apply \eqref{eq:one-history-LLN} to
each cost. Growing sublinear counts were already covered by
Theorem~\ref{thm:sparse-history-LLN}; a subsequence argument combines
the bounded and growing cases. The mean label correction is
negligible on every one of these scales.

Finally, any sequence of edit fractions has a convergent subsequence.
At a positive limit use the dense-edit $L^2$ law and its positive
mean. At zero use the sublinear $L^1$ law just proved and the sharp
mean, whether its edit counts are bounded or grow. A violating
sequence for the supremum in
\eqref{eq:uniform-relative-history-concentration} is therefore
impossible. The same argument applies to label edges.
\end{proof}

\begin{corollary}[History limits under rare coordinate refreshes]
\label{cor:poisson-history-limit}
If coordinates are independently refreshed with probability $s_N$
and $Ns_N\to\lambda<\infty$, then
\[
 \frac{H_N}{\log N}\Longrightarrow
             (1+e^{-\gamma})\operatorname{Poisson}(\lambda).
\]
More precisely, if $K_N$ is the actual refreshed count, then
$H_N/\log N-(1+e^{-\gamma})K_N\to0$ in $L^1$.
This holds for label edges too.

If instead every coordinate refreshes at rate $1/N$, let
$\mathcal H_N(t)$ be the matching distance from the initial sample
at refresh time $t$. Using the same rate-one candidate Poisson clock
$P(t)$ for every $N$, for each fixed $T<\infty$,
\[
 \sup_{0\le t\le T}
 \left|\frac{\mathcal H_N(t)}{\log N}
                    -(1+e^{-\gamma})P(t)\right|
                 \longrightarrow0\quad\hbox{in probability}.
\]
This history limit counts candidate refreshes; it is distinct from
the survivor-location process and its actual-jump rate $p_\infty$.
\end{corollary}
\begin{proof}
Conditional on $K_N=k$, use the fixed-$k$ $L^1$ law. The bound
$\EE H_N(k)\le Ck(1+\log N)$ and uniform integrability of binomial
counts with bounded means remove the count truncation. This proves
the stronger random-count statement and then the Poisson limit.

For the clock version, truncate the number of candidate events at
a fixed $L$. Their coordinate choices repeat with probability at
most $L(L-1)/(2N)$. On the complement, after event $j\le L$ the
sample is a $j$-coordinate resampling of the initial one, so the
fixed-$j$ laws and a finite union bound control all event times.
The distance is constant between events. Finally let $L\to\infty$
using the Poisson tail of $P(T)$.
\end{proof}

\begin{proposition}[Summable log-band covariances in the line entrance law]
\label{prop:line-history-variance}
In the line entrance process of
Lemma~\ref{lem:line-entrance-dilation}, put
$F_n=J(e^{-n},e^{1-n})-a_{\rm H}$ for $n\in\mathbb Z$.
There are $C<\infty$ and $\eta>0$ such that
\begin{equation}\label{eq:line-history-covariance}
 |\operatorname{Cov}(F_0,F_n)|\le Ce^{-\eta|n|}.
\end{equation}
Consequently the absolutely convergent quantity
\[
 v_{\rm H}=\EE F_0^2+2\sum_{n\ge1}\EE(F_0F_n)>0
\]
satisfies
\begin{equation}\label{eq:line-history-variance}
 \operatorname{Var}J(e^{-T},1)\le C(T+1),\qquad
 \frac{\operatorname{Var}J(e^{-T},1)}T\longrightarrow v_{\rm H}.
\end{equation}
In particular the logarithmic-band average has $L^2$ error
$O(T^{-1/2})$ for $T\ge1$. The variance constant is not numerically
calibrated. A simultaneous finite-$M$ variance asymptotic is not
established here.
\end{proposition}
\begin{proof}
For a source at $x$ write $F(x)=J_x(1,e)-a_{\rm H}$.
This stationary spatial field has moments of every order.
At cutoff one, its affected component is contained in the source
gap and the neighbouring strings up to gaps longer than $e$.
The numbers of neighbouring gaps have geometric tails, and every
gap length has an exponential envelope. A single inserted image
can join at most the two components adjacent to its source gap.
The number of captures is bounded by their total point count.

We first prove a translation modulus. Approximate the same-origin
map inside every cutoff-one gap by its depth-$d$ finite-valued
map, assigning the unresolved set to one revealed vertex.
Its expected error per unit spatial length is at most
$C\sigma^d$, by Lemma~\ref{lem:origin-map} and the finite second
moment of the root gap length. Its piece-boundary intensity is at
most $C3^d$, since the stationary gap intensity is finite.
Equal mapped images give equal signed cutoff responses and hence
equal future band costs. Therefore
\[
 \PP\{F(0)\ne F(h)\}\le C(\sigma^d+|h|3^d).
\]
For $|h|\le1$, balance these terms to get a bound $C|h|^{\alpha_0}$,
where $\alpha_0=\log(1/\sigma)/\log(3/\sigma)>0$.
The uniform fourth moment then gives
\begin{equation}\label{eq:line-band-translation-modulus}
 \|F(h)-F(0)\|_2\le C|h|^\kappa,
                         \qquad\kappa=\alpha_0/4>0.
\end{equation}

We next establish the spatial-average bound
\begin{equation}\label{eq:line-band-spatial-average}
 \EE\left|\frac1L\int_0^L F(x)\,dx\right|^2\le C/L,
                                      \qquad L\ge1.
\end{equation}
Divide the cutoff-one marked renewal sequence at gaps longer than
$2e$. Take each block to consist of a short-gap string and the
following long gap. Under the boundary Palm law the marked blocks
$B_i$ are iid, with lengths and point counts having exponential
tails. For a source in the short-gap string, the band response
stays in that block. For a source in its long gap, the image cannot
be within $e$ of both endpoints, so it affects only the block on
one side or the other. Thus the centered block reward
$R_i=\int_{B_i}F(x)\,dx$ depends only on $B_i,B_{i+1}$.
It has finite second moment, as do the boundary rewards under the
stationary size-biased block law, and $\EE R_i=0$ by the stationary
mean of $F$.

For detail on the random number of blocks, put
$g(B_i)=\EE[R_i\mid B_i]$. The variables $g(B_i)$ are iid centered,
while $R_i-g(B_i)$ are martingale differences after revealing
$B_{i+1}$, relative to the preceding revealed blocks.
Including a block until the accumulated length first crosses $L$
is predictable for both sums. Each block has length greater than
$2e$, so the number included is at most $L/(2e)+2$.
Their two second-moment bounds are consequently $O(L)$.
The first and last incomplete blocks have bounded second moments
under the stationary block law. This proves
\eqref{eq:line-band-spatial-average}.

Now average a fine-scale band observable over sources
$x\in[-h,h]$, where $e^{-n}\le h\le1$. By dilation invariance,
the variance of this spatial average is at most $C/(he^n)$.
Stationarity and \eqref{eq:line-band-translation-modulus} show that
replacing $\EE(F_0F_n)$ by its spatially averaged pairing costs
at most $Ch^\kappa$. Cauchy--Schwarz bounds the averaged pairing,
giving
\[
 |\EE(F_0F_n)|\le C\{h^\kappa+(he^n)^{-1/2}\}.
\]
Choose $h=e^{-n/(1+2\kappa)}$ to obtain
\eqref{eq:line-history-covariance} with
$\eta=\kappa/(1+2\kappa)$. Negative lags follow by stationarity
under dilation.

For integer $T$, add covariances of the $T$ consecutive bands.
Absolute summability proves both statements in
\eqref{eq:line-history-variance}, and the limiting variance is
nonnegative. For a noninteger endpoint the extra partial band has
uniformly bounded second moment, so it does not change the limit.
The exact mean is $a_{\rm H}T$ by
\eqref{eq:line-capture-mean}, proving the $L^2$ estimate.

To show positivity, exponential covariance decay gives $F_0$ an
analytic nonnegative spectral density $w(t)$ for the unitary $D$,
with $2\pi w(0)=v_{\rm H}$. If $v_{\rm H}=0$, then $w(t)=O(t^2)$
near zero. Hence $|1-e^{it}|^{-2}$ is integrable against its spectral
measure, and spectral calculus supplies a real $G\in L^2$ with
$F_0=G-DG$. One may obtain $G$ as the $L^2$ limit of the real
resolvents $\sum_{j\ge0}u^jD^jF_0$ as $u\uparrow1$.
But $F_0=J(1,e)-a_{\rm H}$ and $J(1,e)$ is integer valued, so
\[
 D(e^{2\pi iG})=e^{2\pi i a_{\rm H}}e^{2\pi iG}.
\]
This contradicts Lemma~\ref{lem:line-entrance-dilation}, because
$a_{\rm H}\in(1/4,1/2)$. Therefore $v_{\rm H}>0$.
\end{proof}

\begin{theorem}[A functional central limit theorem for line capture histories]
\label{thm:line-history-CLT}
In the stationary line entrance law, for every fixed $t_0<\infty$,
\begin{equation}\label{eq:line-history-CLT}
 \left(\frac{J(e^{-Tt},1)-a_{\rm H}Tt}{\sqrt T}\right)_{0\le t\le t_0}
       \Longrightarrow (\sqrt{v_{\rm H}}\,B_t)_{0\le t\le t_0}
                                      \qquad(T\to\infty),
\end{equation}
where $B$ is standard Brownian motion and $v_{\rm H}>0$ is from
Proposition~\ref{prop:line-history-variance}. Use the right-continuous
version at collision separations; convergence is in the Skorokhod
space, with a continuous limit. The corresponding one-insertion
changed-edge count on these bands has drift $2a_{\rm H}$ and
variance coefficient $4v_{\rm H}$.

This is an entrance-law theorem. The local finite-count comparison
used for the preceding law of large numbers does not by itself
transfer a $\sqrt{\log M}$ fluctuation statement to finite samples.
\end{theorem}
\begin{proof}
We strengthen the covariance argument to products of finitely many
unit-log-band observables. For each fixed $p$, uniformly in ordered
indices $n_1\le\cdots\le n_p$ and each split $1\le j<p$,
\begin{equation}\label{eq:line-history-product-mixing}
 \left|\operatorname{Cov}\left(
       \prod_{i\le j}F_{n_i},\prod_{i>j}F_{n_i}\right)\right|
       \le C_p e^{-\eta_p(n_{j+1}-n_j)},\qquad \eta_p>0.
\end{equation}

Here are the uniformity details. After dilating so that the coarsest
band in a product has index zero, all its bands have upper cutoff
at most $e$. Every band below cutoff one is determined inside the
source's cutoff-one gap: a discrepancy cannot cross an endpoint
protected through cutoff one. The coarsest band's future uses only
the neighbouring components already described. Thus any such product,
centered by its mean, has a block reward depending only on $B_i,B_{i+1}$
in the long-gap construction of
\eqref{eq:line-band-spatial-average}, regardless of the finer indices.

The second moments of these block rewards are bounded uniformly in
the finer indices. To see this, apply Cauchy--Schwarz to the block
integral and use the stationary block Palm identity. It reduces
the required bound to a constant times
\[
 \EE\left[L(0)^2
              \left(1+\left|\prod_{i=1}^pF_{n_i}(0)\right|^2\right)\right],
\]
where $L(0)$ is the length of the block containing a uniform spatial
origin. It has moments of every order. H\"older's inequality and
the dilation-invariant moment bounds for each $F_{n_i}$ bound this
display by $C_p$. The martingale decomposition of block rewards
therefore proves the spatial-average variance bound $C_p/L$ for
every centered product, uniformly in its finer indices.

For the coarse product ending at index $m=n_j$, the same-origin
map estimate and H\"older's inequality give
\[
 \left\|\prod_{i\le j}F_{n_i}(h)
              -\prod_{i\le j}F_{n_i}(0)\right\|_2
       \le C_p(|h|e^m)^{\kappa_p},
 \quad |h|e^m\le1,\qquad \kappa_p=\alpha_0/(4p)>0.
\]
Indeed telescope the product and use its factors' $L^{2p}$ norms;
the $L^{2p}$ translation difference is bounded using a $4p$ moment
and the image-disagreement probability.
The fine product, whose coarsest index is $n=n_{j+1}$, has spatial
average variance at most $C_p/(he^n)$.
The averaging proof of \eqref{eq:line-history-covariance} now gives
\eqref{eq:line-history-product-mixing} by balancing
$(he^m)^{\kappa_p}$ and $(he^n)^{-1/2}$.

It follows that, for each $p\ge2$, the joint cumulant of $p$ ordered
band variables is bounded by
$C_p e^{-\eta_p G}$, where $G$ is the largest gap between successive
indices. For completeness, expand the joint cumulant as its finite
polynomial in subset moments. Across a split attaining $G$, replace
each spanning subset moment by the product of its two parts, using
\eqref{eq:line-history-product-mixing}. All remaining moment factors
are uniformly bounded. The fully factorized polynomial is zero,
since it is the cumulant of two independent groups. This proves
the claimed bound without an independence assumption on the actual
bands.

Sum over all $p$-tuples in a block of $n$ consecutive indices.
After choosing the first index, the sum over nonnegative successive
gaps of $e^{-\eta_p\max g_i}$ is finite. Thus every cumulant of
order $p\ge2$ of the centered block sum is $O_p(n)$.
Its variance divided by $n$ tends to $v_{\rm H}>0$.
All normalized cumulants of order at least three vanish, so the
moments converge to the Gaussian moments; their determinacy gives
the scalar central limit theorem. The same cumulant bounds for
bounded deterministic step weights, together with covariance
summability, give the Gaussian finite-dimensional limits in
\eqref{eq:line-history-CLT} by the Cram\'er--Wold argument.

The second- and fourth-cumulant bounds also give, for every block
of $n\ge1$ consecutive bands,
\[
 \EE\left|\sum_{i=1}^nF_i\right|^4\le Cn^2.
\]
Linear interpolation of the normalized partial sums consequently
has a uniform fourth-moment increment bound $C|t-s|^2$, including
intervals shorter than one mesh step. This proves tightness.
The maximum of $O(T)$ individual band costs, divided by $\sqrt T$,
tends to zero in probability by their common fourth moment and
a union bound. It controls both the interpolation error and the
incomplete terminal logarithmic bands. This proves the functional
limit with the stated right-continuous convention.

Finally, the inserted and background active sets differ by one
signed point at either endpoint of a band. The difference of their
accepted-edge counts over the band is therefore in $\{-1,0,1\}$.
Their changed-edge count is twice the missing-background count
plus that bounded correction, proving the last assertion.
\end{proof}

\subsection{A finite-sample transfer across growing logarithmic bands}
\label{sec:finite-history-CLT}

\status{Working proof: joint functional CLTs for every fixed edit family}

\begin{lemma}[A positive high-intensity comparison]
\label{lem:relative-gap-density}
Let $f_\alpha=w_1^\alpha$ on $[1,\infty)$, extended by zero,
and put $k_\alpha(y)=e^{-2\alpha y}f(y)$. Then
\begin{equation}\label{eq:positive-gap-comparison}
 f_\alpha=(f-k_\alpha)*
          \sum_{j\ge0}(f*k_\alpha)^{*j},
 \qquad \|\mathcal L(f_\alpha)-\mathcal L(f)\|_{\rm TV}
                  \le e^{-2\alpha}.
\end{equation}
Here the zeroth convolution power is a unit atom at zero.
There is a constant $C$ such that, for $\alpha\ge1$ and $y\ge1$,
\begin{equation}\label{eq:relative-gap-density}
 1-e^{-2\alpha}\le\frac{f_\alpha(y)}{f(y)}
                \le\exp(Cy e^{-2\alpha}).
\end{equation}
In a shielded root of length $x$ and inherited cutoff $t$, reveal
its complete configuration at $a\le t$, and hence all its internal
collisions between $a$ and $t$. If $\lambda a\ge1$, the finite-Poisson
and entrance conditional densities of this revealed geometry have
ratio between $e^{-b}$ and $e^b$, where
\begin{equation}\label{eq:full-skeleton-ratio}
 b=C(x/a)e^{-2\lambda a}.
\end{equation}
There is no fixed bound on the number of revealed leaves.
\end{lemma}
\begin{proof}
Write $F(s)=\tanh(E_1(s)/2)$ for the transform of $f$.
The finite-intensity transform and the subtraction formula for
$\tanh$ give
\[
 \widehat f_\alpha(s)
   =\frac{F(s)-F(s+2\alpha)}{1-F(s)F(s+2\alpha)}.
\]
Expanding the denominator proves the positive identity. If
$r_\alpha=\int k_\alpha\le e^{-2\alpha}$, the common subdensity
$f-k_\alpha$ has mass $1-r_\alpha$ and the remaining convolution
terms have total mass $r_\alpha$. This proves the TV bound and
the lower relative bound.

Use the limiting envelope
$b^-e^{-\kappa y}\le f(y)\le b^+e^{-\kappa y}$.
For $j\ge1$,
\[
 \frac{(f^{*(j+1)}*k_\alpha^{*j})(y)}{f(y)}
 \le\frac{b^+}{b^-}\frac1{j!}
       \left(\frac{(b^+)^2y e^{-2\alpha}}{2\alpha}\right)^j.
\]
Indeed $f^{*(j+1)}(v)\le(b^+)^{j+1}v^je^{-\kappa v}/j!$
and $\int e^{\kappa y}k_\alpha(y)\,dy
\le b^+e^{-2\alpha}/(2\alpha)$. Summing gives the upper bound
in \eqref{eq:relative-gap-density}, after increasing $C$.

For the last assertion let $x_1,\ldots,x_l$ be the terminal
gaps at cutoff $a$. Their sum is $x$, each is at least $a$,
and $l\le x/a$. The exact conditional density is
\[
 \frac{\prod_{i=1}^l w_a^\lambda(x_i)}{w_t^\lambda(x)}
\]
on the set whose deterministic internal scan clears by $t$
without removing either protected endpoint. The entrance density
has the same expression with the limiting $w$ fields. These
formulas include the no-split sector. Scale
\eqref{eq:relative-gap-density} for every terminal and the root.
The sum of the upper logarithmic errors is at most
$C(x/a)e^{-2\lambda a}$; the lower errors are at most the same
quantity since $-\log(1-e^{-2\lambda a})\le Ce^{-2\lambda a}$.
This proves \eqref{eq:full-skeleton-ratio}.
\end{proof}

\begin{lemma}[De-Poissonizing a complete small root interval]
\label{lem:complete-root-depoisson}
Put $w=(\log M)^{1/3}$, and suppose
\begin{equation}\label{eq:root-comparison-window}
 e^w/M\le t\le e^{-w}.
\end{equation}
Choose an independent uniform source, rotate it to zero, and record
the cutoff-$t$ gap containing it, its two endpoints, and every
original input point in its interior. Give the empty and singleton
sectors a separate exceptional value. The laws of this complete
local record for an exact size-$M$ input and a Poisson input of
intensity $M$ have TV distance at most $Ce^{-cw}$, uniformly in
\eqref{eq:root-comparison-window}, for some $c>0$ and all sufficiently
large $M$. The assertion holds for either parity of $M$.
\end{lemma}
\begin{proof}
Write $x$ for the root length, $z\in(0,x)$ for the distance of
the source from its left endpoint, and $K$ for one plus the number
of original interior points. Thus $K$ is odd and the left endpoint
is counted locally; the right endpoint belongs to the exterior.
For any specified local point configuration that is shielded and
clears by $t$, its Poisson density is
\[
 e^{-Mx}M^K U_t^M(1-x)
\]
against $dz$, $dx$ and the ordered interior-point coordinates.
There is no additional cyclic-root factor. The exterior has been
summed without revealing any of its points.

Put $L=1-x$, $n=M-K$, and
$A_n(L)=n![\lambda^n]e^{\lambda L}U_t^\lambda(L)$.
Write $(M)_K=M!/(M-K)!$. The exact-count density divided by the
Poisson density is
\begin{equation}\label{eq:local-count-density-ratio}
 \frac{(M)_K}{M^K}e^{Mx}\frac{A_n(L)}{U_t^M(L)},
 \qquad A_n(L)=L^{n-1}I_n(t/L).
\end{equation}
The identity on the right follows by spatial scaling. Uniformly
for $x\le wt$ and $K\le M/2$, the first-intensity coefficient
bound, the factorial parity remainder, and the high-intensity
renewal comparison give
\[
 A_n(L)=\frac c t L^n(1+o(e^{-w})),\qquad
 U_t^M(L)=\frac c t(1+o(e^{-w})).
\]
For example the relative coefficient error is bounded by
$Ct^{-2}e^{-Mt}$ and the renewal remainder by
$C/\lfloor L/t\rfloor!+Ce^{-2Mt}$. Under
\eqref{eq:root-comparison-window} these are $o(e^{-w})$:
$e^w/\log M\to\infty$ and $L/t\ge e^w/2$.

Thus \eqref{eq:local-count-density-ratio} is, up to that error,
\[
 R_{M,x}(K)=\frac{(M)_K}{M^K}(1-x)^{M-K}e^{Mx}.
\]
Uniform Stirling bounds, or the integral remainder for
$\sum_{j<K}\log(1-j/M)$, yield
\begin{equation}\label{eq:local-count-Stirling}
 |\log R_{M,x}(K)|
 \le C\left(x+\frac{(K-Mx)^2}{M}+\frac1M\right)
                 \quad (x,K/M\le1/2).
\end{equation}
For clarity, with $y=K/M$ its leading term is
\[
 M\left[-(1-y)\log\frac{1-y}{1-x}-y+x\right]
                  -\tfrac12\log(1-y)+O(M^{-1}).
\]
The bracket and its first derivative in $y$ vanish at $y=x$;
its second derivative is $-1/(1-y)$. Also
$y\le x+\tfrac12M(y-x)^2+\tfrac1{2M}$, which bounds the logarithmic
prefactor by the same right side.

It remains to show that this ratio is near one on most Poisson
local records. The source-root length has density
$xw_t^M(x)U_t^M(1-x)$ on $0<x<1$. The uniform finite-intensity
envelope, $U_t^M\le C/t$, and the empty/singleton estimate in
Lemma~\ref{lem:full-corner-law} give
\[
 \PP(x>wt\text{ or exceptional})\le C(w+1)e^{-cw}+Ct.
\]
Conditional on $x$, the complete local count has the exact mgf
\[
 \EE(e^{uK}\mid x)
 =e^{Mx(e^u-1)}\frac{w_t^{Me^u}(x)}{w_t^M(x)}.
\]
For $x\le wt$ and $|u|\le1/2$, Lemma~\ref{lem:relative-gap-density}
bounds the field ratio by two for large $M$. Chernoff's inequality
therefore gives, with $\delta=\sqrt M e^{-w/4}$,
\[
 \PP(|K-Mx|>\delta\mid x)
 \le4\exp\left[-c\min\left\{\frac{\delta^2}{Mx},\delta\right\}\right]
                       =o(e^{-w})
\]
uniformly on those roots. Indeed $Mx\le Mw e^{-w}$.
On the complementary good event $K\le M/2$ and
\eqref{eq:local-count-Stirling} is $O(we^{-w}+e^{-w/2})$.
The two densities consequently agree there to that relative error.
Their total masses are one, so the Poisson bad-event bound also
bounds the exact-count bad mass, up to this error. This proves
the TV assertion, including the exceptional sectors.
\end{proof}

\begin{lemma}[A growing-band coupling and its omitted strips]
\label{lem:growing-band-coupling}
Put
\[
 T=\log M,\quad w=T^{1/3},\quad r_-=e^w/M,\quad r_+=e^{-w},
 \quad a=e^{-w/4}r_-.
\]
For an iid background of either parity and one independent uniform
insertion, its missing-background-edge process on $(r_-,r_+]$
can be coupled to the stationary line entrance process so that
\begin{equation}\label{eq:growing-band-coupling}
 J_M(r_-,s)=J(r_-,s)\quad\hbox{for every }r_-\le s\le r_+
\end{equation}
except on probability $Ce^{-cw}$, for some $c>0$.
For each fixed $1\le p<\infty$, the omitted early and late counts satisfy
\begin{equation}\label{eq:history-strip-moments}
 \|J_M(0,r_-)\|_p+\|J_M(r_+,1)\|_p\le C_p w.
\end{equation}
Consequently this coupling can be chosen with
\begin{equation}\label{eq:full-capture-Lp-coupling}
 \|J_M(0,1)-J(r_-,r_+)\|_p\le C_p w.
\end{equation}
All statements concern spatial histories and preserve the physical
identities needed to decide which background edges are missing.
\end{lemma}
\begin{proof}
First apply Lemma~\ref{lem:complete-root-depoisson} at $t=r_+$.
The Poisson root covering the source has joint length/position
density $w_t^M(x)U_t^M(1-x)\,dx\,dz$, $0<z<x$.
On $x\le wt$, Lemma~\ref{lem:relative-gap-density} and the renewal
bounds compare this to
$(c/t)w_t(x)\,dx\,dz$, the line root law, with error $o(e^{-w})$.
The discarded root mass is $Ce^{-cw}$ in both laws. Conditional
on matching root geometry, reveal the entire root skeleton down
to $a$. By \eqref{eq:full-skeleton-ratio} its density ratio has
logarithm bounded by
\[
 C(wt/a)e^{-2Ma}
       =CwM e^{-7w/4}\exp(-2e^{3w/4})=o(e^{-w}).
\]
Thus the finite and line background skeletons inside the source's
high-cutoff root agree with probability $1-Ce^{-cw}$.

We must also identify the insertion image at $r_-$.
In the line's cutoff-$r_-$ gap covering the source, reveal
$d=\lfloor w/(8\log(3J))\rfloor$ tree generations, with
$J=1+K_0$ as in Lemma~\ref{lem:origin-map}.
The inverse-cutoff bound gives
\[
 \PP(\text{some revealed terminal cutoff}<a)
       \le C(a/r_-)(3J)^d\le Ce^{-w/8}.
\]
The size-biased root gap has uniformly bounded scaled moments.
The expected deterministic origin-disagreement budget for two
arbitrary deeper completions of this shared geometry is therefore
at most $C\sigma^d$. Hence the two insertion images at $r_-$ agree
except on probability $Ce^{-cw}$. This uses the limiting budget
on shared geometry; no finite-intensity contraction estimate for
the origin map is being assumed.

The two high-root endpoints are protected through $r_+$.
If an insertion removes one of them earlier, the resulting negative
discrepancy stays there until the endpoint's background collision,
which occurs later. Therefore a discrepancy cannot leave this
root before $r_+$, and its entire central history is determined
by the revealed skeleton and the image at $r_-$.
This proves \eqref{eq:growing-band-coupling}.

We record the moment bounds needed to use the exceptional events.
Summing the geometric ancestry tail in
Lemma~\ref{lem:insertion-height} gives, for every fixed $p$,
\begin{equation}\label{eq:all-history-moments}
 \EE J_M(0,1)^p\le C_p(1+\log M)^p.
\end{equation}
For odd $M$ use the singleton full-circle root at cutoff $1/3$,
as in the proof of \eqref{eq:one-history-second-moment};
the same finite-Poisson daughter contraction and exact-count
conditioning give the same geometric tail. This also explains
why all fixed moments, not only the second, are available.

For the early strip, apply the complete-root comparison at $t=r_-$.
All captures below this cutoff lie on the origin's ancestry inside
that shielded root. In its Poisson law, conditional on root length
$x$ and averaging the uniform origin position, the ancestry tail
is bounded by
\[
 \min(1,Mx\rho_{\rm h}^{\,d}).
\]
Its $p$th moment is at most
$C_p[1+\log(1+Mx)]^p$. The Poisson source-root length divided by
$r_-$ has a uniformly exponential, size-biased envelope, and
$Mr_-=e^w$. Thus the early count has $L^p$ norm at most $C_pw$.
Set its comparison value to zero on improper roots. The root TV
coupling, \eqref{eq:all-history-moments}, and H\"older's inequality
transfer this bound to the exact sample: the bad-event contribution
to the $L^p$ norm is at most
$C_p(1+\log M)e^{-cw/(2p)}=o(1)$, using the $2p$ moment.

For the late strip, use Lemma~\ref{lem:finite-family-rate} to couple
the finite background and its insertion to their circle entrance
family above $M^{-\chi_{\rm f}}$. For large $M$, this is below $r_+$;
the error is polynomially small in $M$. In either circle entrance
parity, follow the ancestry of a uniform origin from the final
root, using a binary top root in the even case and the singleton
full-circle root in the odd case. The limiting two-sided envelope
gives the same expected daughter-length contraction
$\rho_{\rm h}<1$ as in Lemma~\ref{lem:insertion-height}.
After $d$ generations the expected surviving interval length is
at most $C\rho_{\rm h}^d$. If at least $d+O(1)$ ancestral mergers
have separation at least $r_+$, that interval has length at least
$r_+$. Hence their count has tail at most
\[
 \min(1,C r_+^{-1}\rho_{\rm h}^d).
\]
Every captured background edge is one such ancestral merger.
The entrance late count therefore has $L^p$ norm
$C_p(1+\log(1/r_+))\le C_pw$.
H\"older's inequality and \eqref{eq:all-history-moments} absorb
the polynomial coupling error. This proves
\eqref{eq:history-strip-moments}.

Finally the line count on $(r_-,r_+]$ has every fixed $L^p$ norm
bounded by $C_p(1+\log M)$, by Minkowski's inequality for its unit
logarithmic bands. On the exceptional event in
\eqref{eq:growing-band-coupling}, use these $2p$ moments and
\eqref{eq:all-history-moments}. The central-count error is $o(1)$
in $L^p$. Adding the two strips proves
\eqref{eq:full-capture-Lp-coupling}. The same argument bounds the
supremum of the central cumulative-count error.
\end{proof}

\begin{theorem}[Finite one-edit history fluctuations]
\label{thm:finite-one-history-CLT}
Let $V_M=J_M(0,1)$ be the total number of background edges destroyed
by one independent uniform insertion into $M$ iid circle points,
of either parity. With $a_{\rm H}=(1+e^{-\gamma})/4$ and the strictly
positive line variance $v_{\rm H}$,
\begin{align}
 \EE V_M&=a_{\rm H}\log M+O((\log M)^{1/3}),
                                      \label{eq:finite-capture-mean-rate}\\
 \operatorname{Var}V_M&\sim v_{\rm H}\log M,\qquad
 \frac{V_M-a_{\rm H}\log M}{\sqrt{\log M}}
       \Longrightarrow\mathcal N(0,v_{\rm H}).
                                      \label{eq:finite-capture-CLT}
\end{align}
In fact, with the right-continuous collision convention,
\begin{equation}\label{eq:finite-capture-FCLT}
 \left(\frac{J_M(0,M^{u-1})-a_{\rm H}u\log M}{\sqrt{\log M}}
                  \right)_{0\le u\le1}
       \Longrightarrow (\sqrt{v_{\rm H}}B_u)_{0\le u\le1}.
\end{equation}
Convergence is in the Skorokhod space, and the scalar centered
moments of every fixed order converge to those of the displayed
Gaussian law.

For the changed-edge cost $H$ of one addition or one uniform
deletion, the mean coefficient is $2a_{\rm H}$ and the variance
coefficient is $4v_{\rm H}$. Thus its centered scalar limit is
$\mathcal N(0,4v_{\rm H})$, with normalization $\sqrt{\log M}$.
No value of $v_{\rm H}$ is certified here. The separate joint
comparison below handles resampling and other fixed edit families.
\end{theorem}
\begin{proof}
Use the preceding lemma and set
$L=T-2w=\log(r_+/r_-)$, where $T=\log M$.
By dilation invariance the coupled line total has the law of
$J(e^{-L},1)$, whose mean is exactly $a_{\rm H}L$ and variance
is asymptotic to $v_{\rm H}L$. The $L^2$ comparison error is
$O(w)=o(\sqrt T)$. It proves the mean bound, variance asymptotic,
and scalar central limit theorem by Theorem~\ref{thm:line-history-CLT}.
For every fixed moment order the $L^p$ comparison gives the same
conclusion; the line's moment convergence follows from the
cumulant bounds in that theorem, also using the next higher even
moment for uniform integrability.

For the path assertion, let
$s(u)=\max(0,\min(L,uT-w))$. Uniformly in $u\in[0,1]$, the finite
cumulative count differs in $L^p$ supremum norm by $O(w)$ from
the line cumulative count up to central log-time $s(u)$.
This follows from the last sentence of the coupling lemma and
the two nonnegative strip bounds. After scaling the top cutoff
to one, the latter path is
\[
 J(e^{-L},e^{-L+s})=J(e^{-L},1)-J(e^{-L+s},1),\qquad 0\le s\le L.
\]
The line functional theorem, followed by time reversal, gives
the limit $\sqrt{v_{\rm H}}(B_1-B_{1-u})$, which has the law of
Brownian motion. Since $s(u)/L\to u$ uniformly and
$|s(u)-uT|\le2w$, neither the deterministic time change nor
the change of centering affects the limit. This proves
\eqref{eq:finite-capture-FCLT}.

At any cutoff the one-addition changed-edge count is $2J_M$
plus a term in $\{0,1\}$, by the one-active-discrepancy identity.
In particular this holds for the complete matching. Reversing
the iid add-one construction gives a uniform deletion, with
$M-1$ in place of $M$, which changes no displayed asymptotic.
The bounded correction proves all stated one-edit conclusions.
\end{proof}

\begin{lemma}[Independent central histories at finitely many sources]
\label{lem:independent-local-histories}
Fix $q$. In a common iid background of size $M$, insert independent
uniform probes $Z_1,\ldots,Z_q$. Use $w,r_-,r_+,a$ from
Lemma~\ref{lem:growing-band-coupling}. The $q$ individual insertion
histories on $(r_-,r_+]$ can be coupled to $q$ independent stationary
line insertion histories, agreeing simultaneously throughout the
band except on probability $C_qe^{-c_qw}$. On the same good event,
the probes occupy pairwise nonadjacent cutoff-$r_+$ gaps, and
every subset of probes superposes its individual histories through
that cutoff.
\end{lemma}
\begin{proof}
Put $t=r_+$ and $\ell=e^{-w/2}$. Except on probability $C_q\ell$,
the circular distances between all probe locations are at least
$\ell$. Condition on such locations in cyclic order. Reveal their
cutoff-$t$ roots, with lengths $x_i\le wt$ and internal source
positions $z_i\in(0,x_i)$. Their closed intervals are disjoint and
nonadjacent, since $2wt<\ell$ for large $M$. Write $L_i$ for the
exterior arc between root $i$ and root $i+1$. Then
$L_i\ge\ell/2$, and $\sum_iL_i=1-x$, where $x=\sum_i x_i$.

For a Poisson input of intensity $M$, the joint local tree density
is the product of its local fields times
\begin{equation}\label{eq:several-root-outside-factor}
 \prod_{i=1}^q U_t^M(L_i).
\end{equation}
If every original interior point is specified and $K_i$ counts
those points plus the left endpoint, its local field is instead
$e^{-Mx_i}M^{K_i}$ on the appropriate shielding event.
These formulas follow by summing the independent outside gap
strings. The source positions identify the roots uniquely, so
neither a cyclic-root factor nor a permutation factor is present.
The probability of a discarded long or improper root is at most
$C_qe^{-cw}$, by the one-root marginal bound and a union bound.

We justify exact-count conditioning before passing to the entrance
trees. Put $K=\sum_iK_i$ and $n=M-K$. The outside coefficient is
\[
 A_{n,q}=n![\lambda^n]e^{\lambda(1-x)}
                              \prod_iU_t^\lambda(L_i).
\]
Uniformly for these root geometries and $n\ge M/2$,
\begin{equation}\label{eq:several-root-coefficient}
 A_{n,q}=(c/t)^q(1-x)^n(1+o(e^{-w})),\qquad
 \prod_iU_t^M(L_i)=(c/t)^q(1+o(e^{-w})).
\end{equation}
For detail, expand each renewal factor as in
\eqref{eq:renewal-expansion}. Every component length is at least
$t$, and their total is $1-x$. The empty subset contributes
$(1-x)^n\prod_i U_t(L_i)$, while the full subset contributes
$(-1)^n(1-x)^n\prod_i U_{-,t}(L_i)$, where
$U_{-,t}=-(R_{{\rm o},t}-R_{{\rm e},t})$.
The factorial remainder makes $U_t(L_i)=c/t+o(t^{-1}e^{-w})$
and $U_{-,t}(L_i)=o(t^{-1}e^{-w})$, since
$L_i/t\ge e^{w/2}/2$. Every proper nonempty subset has extracted
coefficient at most $(1-x)^n e^{-2nt/(1-x)}$ in modulus.
The product of the $q$ absolute renewal series is at most
$C_qt^{-3q}$, so its relative error is at most
$C_qt^{-2q}e^{-Mt}=o(e^{-w})$. This proves the first formula;
the high-intensity renewal comparison proves the second.

The finite/Poisson density ratio is consequently
$R_{M,x}(K)(1+o(e^{-w}))$, exactly as in
\eqref{eq:local-count-density-ratio}. Given the root lengths,
the local Poisson records are independent shielded intervals, so
\[
 \EE(e^{uK}\mid x_1,\ldots,x_q)
 =e^{Mx(e^u-1)}
                 \prod_i\frac{w_t^{Me^u}(x_i)}{w_t^M(x_i)}.
\]
The product ratio is at most $2^q$ for $|u|\le1/2$ and large $M$.
Since $x\le qwt$, the same Chernoff window
$|K-Mx|\le\sqrt M e^{-w/4}$ and the same Stirling bound prove
TV agreement of the complete finite and Poisson local records
with error $C_qe^{-c_qw}$. Normalization again handles the bad
exact-count records; no independence after conditioning is asserted.

Next \eqref{eq:several-root-outside-factor}, the scalar relative
bound, and \eqref{eq:several-root-coefficient} compare the joint
Poisson root geometry with the product of $q$ line root laws.
Conditional on the roots, their past trees are independent.
Apply \eqref{eq:full-skeleton-ratio} down to $a$ in every root,
and then the origin-map comparison at $r_-$ from
Lemma~\ref{lem:growing-band-coupling}. A union bound over the
fixed number of roots proves simultaneous agreement of all
individual central histories with independent line histories.

Finally the proof of Lemma~\ref{lem:separated-insertions} gives
more than its displayed cutoff readout: discrepancies stay in
their pairwise disjoint closed roots throughout the event scan
up to $t$. Negative endpoint discrepancies are frozen and
positive ones cannot cross an intervening protected endpoint.
Thus every subset superposes its individual histories throughout
this interval, with disjoint changed-edge supports. This proves
the last assertion.
\end{proof}

\begin{theorem}[Joint Gaussian laws for fixed edit families]
\label{thm:fixed-edit-history-CLT}
Fix $q$ and let $B_M$ be a common iid background, with $q$
independent uniform probes. For $S,T\subseteq\{1,\ldots,q\}$,
let $H_M^{S,T}(u)$ be the symmetric-difference size of their
accepted matchings through separation $M^{u-1}$, $0\le u\le1$.
Jointly over all pairs $S,T$,
\begin{equation}\label{eq:fixed-family-history-FCLT}
 \left(\frac{H_M^{S,T}(u)-2a_{\rm H}|S\triangle T|u\log M}
                  {\sqrt{\log M}}\right)_{0\le u\le1}
 \Longrightarrow
 \left(2\sqrt{v_{\rm H}}\sum_{i\in S\triangle T}B_i(u)
                                      \right)_{0\le u\le1},
\end{equation}
where $B_1,\ldots,B_q$ are independent standard Brownian motions.
The terminal joint moments of every fixed order converge as well.
In particular, if $r=|S\triangle T|>0$, then
\begin{align}
 \EE H_M^{S,T}(1)&=2a_{\rm H}r\log M+O_q((\log M)^{1/3}),
                                      \label{eq:fixed-edit-mean-rate}\\
 \operatorname{Var}H_M^{S,T}(1)&\sim4r v_{\rm H}\log M.
                                      \label{eq:fixed-edit-variance}
\end{align}
For each fixed resampling count $k\ge1$, both spatial and label
history costs consequently satisfy
\begin{equation}\label{eq:fixed-resampling-history-CLT}
 \frac{H_N(k)-(1+e^{-\gamma})k\log N}{\sqrt{\log N}}
           \Longrightarrow\mathcal N(0,8k v_{\rm H}),
 \qquad \operatorname{Var}H_N(k)\sim8k v_{\rm H}\log N.
\end{equation}
More generally a fixed mixture of $r$ uniform deletions and fresh
additions has Gaussian variance coefficient $4r v_{\rm H}$ and
mean coefficient $2a_{\rm H}r$, regardless of parity.
\end{theorem}
\begin{proof}
Set $t=r_+$ and use the joint central coupling just proved.
On its good event, for every $s\le t$, the changed-edge supports
of the individual probes are disjoint and
\[
 H_{\le s}(B_M+S,B_M+T)
       =\sum_{i\in S\triangle T}H_{\le s}(B_M,B_M+Z_i).
\]
The common probes have identical local histories and cancel.
The individual early strips have every fixed $L^p$ norm $O_p(w)$.
For the late part of any pairwise history, telescope through at
most $q$ one-edit steps. Each intermediate background is iid of
size $M+j$, $0\le j\le q$, and each newly added probe is independent
of it; reverse a step when removing a probe. The symmetric-difference
metric for edges of separation greater than $t$ obeys the same
triangle inequality. Its one-edit cost is at most twice the
corresponding late captured-edge count plus one. The late-strip
bound thus gives $L^p$ norm $O_{p,q}(w)$, uniformly over the finite
family. Using the slightly smaller cutoff $r_+(M+j)$ only enlarges
this late count and changes no bound.

Full family history costs have $L^p$ norms $O_{p,q}(\log M)$ by
the same telescoping and \eqref{eq:all-history-moments}. These
bounds absorb the exceptional event $C_qe^{-c_qw}$ by H\"older's
inequality. Consequently, simultaneously for all pairs, the
entire cutoff process differs in $L^p$ supremum norm by
$O_{p,q}(w)$ from twice the sum, over $S\triangle T$, of its
independent line central capture processes. The one-edit bounded
edge-count corrections do not affect this estimate.

Apply the independent line functional limits and the same
clamped log-time change as in
Theorem~\ref{thm:finite-one-history-CLT}. Since
$w=o(\sqrt{\log M})$, this proves
\eqref{eq:fixed-family-history-FCLT}. The $L^p$ estimates and
the line cumulant bounds give terminal joint moment convergence,
including \eqref{eq:fixed-edit-mean-rate} and
\eqref{eq:fixed-edit-variance}. In particular the limiting
covariance of two centered terminal distances, divided by
$\log M$, is
$4v_{\rm H}|(S\triangle T)\cap(S'\triangle T')|$.

For $k$ resamplings use $M=N-k$, $k$ original probes and $k$
independent replacement probes, with disjoint index sets $S,T$.
Then $r=2k$ and $\log(N-k)=\log N+o(1)$.
The deterministic label correction is at most $2k$, so it
changes neither the centered limit nor its variance asymptotic.
A fixed deletion/addition mixture has the same common-background
description with its total edit count $r$, proving the final claim.
\end{proof}

\begin{corollary}[Centered rare-refresh history limits]
\label{cor:rare-history-CLT}
Put $b_{\rm H}=1+e^{-\gamma}$. If each coordinate is independently
resampled with probability $s_N$ and $Ns_N\to\lambda<\infty$,
let $K_N$ be the actual number resampled. Then, for either spatial
or label history distance,
\begin{equation}\label{eq:rare-history-CLT}
 \left(K_N,\frac{H_N-b_{\rm H}K_N\log N}{\sqrt{\log N}}\right)
       \Longrightarrow
       \left(K,\sum_{j=1}^K G_j\right),
\end{equation}
where $K\sim\operatorname{Poisson}(\lambda)$ and the $G_j$ are iid
$\mathcal N(0,8v_{\rm H})$, independent of $K$.

If coordinates refresh at rate $1/N$, write $\mathcal H_N(t)$ for
the matching distance from the initial input and $P(t)$ for the
common rate-one candidate clock. Jointly with this clock, on each
bounded refresh-time interval,
\begin{equation}\label{eq:history-clock-CLT}
 \left(\frac{\mathcal H_N(t)-b_{\rm H}P(t)\log N}{\sqrt{\log N}}
                                                    \right)_{t\ge0}
       \Longrightarrow\left(\sum_{j=1}^{P(t)}G_j\right)_{t\ge0}
\end{equation}
in the Skorokhod space. The limit is a compound Poisson process
with Gaussian jump sizes, not a Brownian process in refresh time.
The centering uses the actual candidate count.
\end{corollary}
\begin{proof}
Conditional on $K_N=k$, apply the fixed-$k$ resampling CLT.
Binomial convergence and tightness of $K_N$ permit truncation at
any fixed number of edits, proving the joint limit
\eqref{eq:rare-history-CLT}; then send that truncation to infinity.
No uniform growing-$k$ centered limit is needed.

For the clock statement, condition on exactly $l$ candidate events
on the time interval. Repeated coordinates have probability at
most $l(l-1)/(2N)$. On the complementary event, expose the $l$
original values of those coordinates and their $l$ fresh values
as $2l$ independent probes over a common size-$(N-l)$ background.
At the $j$th event, the symmetric difference of the initial and
current probe subsets consists of the first $j$ original/fresh
pairs. The joint family theorem therefore represents the centered
terminal history distances by successive sums of independent
Gaussian pair contributions, each of variance $8v_{\rm H}$.
Event times are independent of the input and probe locations.
Convergence of this finite jump vector proves the Skorokhod
statement conditional on the event count; Poisson truncation
removes the conditioning. The fixed-count label correction is
bounded and hence disappears under the same truncation.
\end{proof}

\begin{remark}[A smaller error in the fixed-edit mean]
\label{rem:history-mean-loglog}
The error terms in \eqref{eq:finite-capture-mean-rate} and
\eqref{eq:fixed-edit-mean-rate} can be improved to $O_q(\log\log M)$.
Here are the uniformity checks for this refinement. The root and
central-skeleton comparisons remain valid, with error
$C_qe^{-c_qw}$, uniformly for
\[
 2\log\log M\le w\le\tfrac14\log M.
\]
Indeed $Mt\ge e^w\ge(\log M)^2$ controls the exterior coefficient,
$Ma=e^{3w/4}\ge(\log M)^{3/2}$ controls the complete skeleton,
and all root tails and origin budgets already have the stated
exponential dependence on $w$. The local-count Chernoff bound
is uniform since $\delta^2/(Mx)\ge e^{w/2}/(q w)$.
The comparison strips in the Poisson and entrance laws have $L^p$
cost $C_{p,q}w$. Their finite-sample transfer and the central
coupling add a total exceptional-event contribution bounded by
\[
 C_{p,q}(\log M)e^{-c_qw/(2p)}.
\]
For any fixed $p,q$, choose $w=A_{p,q}\log\log M$ with
$A_{p,q}>\max(2,2p/c_q)$. The comparison then has $L^p$ error
$O_{p,q}(\log\log M)$. Taking $p=1$ proves the asserted sharper
mean error; taking larger fixed $p$ also retains every central
limit and moment statement above. The constants are uncalibrated.
\end{remark}

\subsection{Gaussian histories for slowly growing edit counts}
\label{sec:slow-history-CLT}

\status{Working proof: centered limits when $q(\log\log M)^2=o(\log M)$}

\begin{lemma}[A uniform comparison for logarithmically many probes]
\label{lem:logarithmic-probe-comparison}
Write $T=\log M$. For every $B>0$ there is a finite $A_B$ such
that, with $w=A_B\log T$, the central-history coupling of
Lemma~\ref{lem:independent-local-histories} holds uniformly
for $1\le q\le T$, with failure probability at most
$C_BT^{-B}$. The $q$ comparison line histories are independent.

For any fixed $p\ge1$, choose $B>2p$ and increase $A_B$ if
necessary. If $d+b=q\le T$ and $H_M^{d,b}(u)$ compares $d$ original
probes with $b$ fresh probes over the common background, there
is a coupling such that
\begin{equation}\label{eq:logarithmic-probe-Lp}
 \left\|\sup_{0\le u\le1}
 \left|H_M^{d,b}(u)-2\sum_{i=1}^q J_i(e^{-L},e^{-L+s(u)})\right|
                                      \right\|_p
        \le C_pq\log T.
\end{equation}
Here $L=T-2w$, $s(u)=\max(0,\min(L,uT-w))$, the $J_i$ are
independent line capture processes, and the finite cutoff is
$M^{u-1}$. Constants are independent of $d,b,q,M$ in the stated
range; their numerical values are not calibrated.
\end{lemma}
\begin{proof}
We track the dependence on $q$ in the preceding root argument.
Put $t=e^{-w}$ and $\ell=e^{-w/2}$. Close source pairs have
probability at most $Cq^2e^{-w/2}$; a long or improper root costs
at most $Cq(w+1)e^{-c_0w}$ for an absolute $c_0>0$.
On the remaining event the total local length is
$x\le qwt=o(1)$ uniformly for $q\le T$, provided $A_B>2$.
All exterior arcs have length at least $\ell/2$.

The relative proper-subset coefficient remainder in
\eqref{eq:several-root-coefficient} is at most
\[
 (Ct^{-2})^q e^{-Mt}
       \le\exp\{O_{A_B}(T\log T)-M/T^{A_B}\},
\]
which is smaller than every power of $T^{-1}$. The empty/full
subset remainder is bounded by a constant times
$q/\lfloor\ell/(2t)\rfloor!$, also smaller than every such power.
The Poisson exterior comparison has the same bound plus
$Cqe^{-2Mt}$.

Conditional on the root geometry, the field ratio in the mgf
of the total local count is bounded more sharply than $2^q$:
\[
 \prod_i\frac{w_t^{Me^u}(x_i)}{w_t^M(x_i)}
 \le\exp\{Cq(w+1)e^{-2Me^{-1/2}t}\}\le2,
                                      \qquad |u|\le1/2,
\]
for large $M$, by \eqref{eq:relative-gap-density}.
With $\delta=\sqrt M e^{-w/4}$, the Chernoff exponent is at least
\[
 c\min\left\{\frac{e^{w/2}}{q w},\delta\right\}.
\]
For $A_B>4$ this exceeds every fixed multiple of $\log T$.
On that count window the Stirling error is
$O(qw e^{-w}+e^{-w/2})$.

The product of the full-root skeleton density ratios down to
$a=e^{3w/4}/M$ has logarithmic error at most
\[
 Cqwt\,a^{-1}e^{-2Ma}
  =CqwM e^{-7w/4}\exp(-2e^{3w/4}),
\]
again negligible when $A_B>4$. The sum of the $q$ inverse-cutoff
and origin-map errors is at most $Cqe^{-c_1w}$, for an absolute
$c_1>0$. Choose $A_B$ large enough that the source, root, count,
and map errors above sum to $C_BT^{-B}$, uniformly for $q\le T$.
All products have now been controlled explicitly; a constant
depending on an unbounded number of roots has not been suppressed.
Normalization and the same conditional couplings prove the
first assertion, including separated closed roots and full early
history superposition.

For the second assertion, each omitted single-probe strip has
$L^p$ norm $O_p(w)$ after increasing $A_B$, by the estimates of
Remark~\ref{rem:history-mean-loglog}. For the late multi-edit
strip, telescope through the $q$ single-edit backgrounds of sizes
between $M$ and $M+q$. The single-edit constants are uniform in
that range, and using their smaller high cutoffs only enlarges
the late counts. Thus the total strip contribution is $O_p(qw)$.
On the good event, disjoint closed roots give exact additivity
of the early histories. On its complement, the finite history
and the sum of line histories have $L^{2p}$ norms $O_p(qT)$.
Its contribution is at most $C_pqT^{1-B/(2p)}=O_p(q)$.
The bounded per-probe edge corrections cost $O(q)$ as well.
This proves \eqref{eq:logarithmic-probe-Lp}.
\end{proof}

\begin{theorem}[Slowly growing edit central limit theorem]
\label{thm:slow-edit-history-CLT}
Let $q=d+b\ge1$ vary with the common iid background size $M$ and
suppose
\begin{equation}\label{eq:slow-edit-range}
 q(\log\log M)^2=o(\log M).
\end{equation}
For the spatial matching history between $d$ original and $b$
fresh probes, of either parity,
\begin{equation}\label{eq:slow-edit-history-FCLT}
 \left(\frac{H_M^{d,b}(u)-2a_{\rm H}q u\log M}
                     {\sqrt{q\log M}}\right)_{0\le u\le1}
       \Longrightarrow (2\sqrt{v_{\rm H}}B_u)_{0\le u\le1}.
\end{equation}
The terminal centered moments of every fixed order converge.
Moreover, uniformly in this range,
\begin{align}
 \EE H_M^{d,b}(1)&=2a_{\rm H}q\log M+O(q\log\log M),
                                      \label{eq:slow-edit-mean}\\
 \operatorname{Var}H_M^{d,b}(1)&\sim4v_{\rm H}q\log M.
                                      \label{eq:slow-edit-variance}
\end{align}
The terminal centering may be replaced by
$2a_{\rm H}q\log(M/q)$, or by the exact mean.

Consequently, whenever $1\le k=k_N$ and
$k(\log\log N)^2=o(\log N)$, both spatial and label costs satisfy
\begin{equation}\label{eq:slow-resampling-CLT}
 \frac{H_N(k)-(1+e^{-\gamma})k\log(N/k)}
                      {\sqrt{k\log(N/k)}}
       \Longrightarrow\mathcal N(0,8v_{\rm H}),
 \qquad \operatorname{Var}H_N(k)\sim8v_{\rm H}k\log(N/k).
\end{equation}
The functional limit and the explicit leading-order centering here
have this restricted range. Section~\ref{sec:late-history-variance}
separately proves the terminal CLT centered by the exact mean for
every growing $k=o(N)$.
Section~\ref{sec:all-sparse-functional} removes the growth restriction
for exact-mean-centered cumulative histories on the logarithmic
clock ending at separation $1/q$, with the completed endpoint
also allowed. It does not replace this mean by the leading formula.
Section~\ref{sec:late-history-maximal} separately controls the
whole remaining late path.
Section~\ref{sec:dense-history-CLT} also gives the positive-density CLT.
\end{theorem}
\begin{proof}
By \eqref{eq:logarithmic-probe-Lp}, the normalized comparison
error is at most
\[
 C_p\frac{q\log\log M}{\sqrt{q\log M}}=o(1).
\]
The deterministic centering changes from $2a_{\rm H}q s(u)$ to
$2a_{\rm H}q uT$ by the same order, uniformly in $u$.

For completeness the line limit is uniform in the number of
independent copies. A centered sum over $n$ complete log bands
in one copy has every cumulant of order $j\ge2$ bounded by
$C_j n$, by Theorem~\ref{thm:line-history-CLT}.
Summing $q$ independent copies multiplies these cumulants by $q$.
The variance is $q(v_{\rm H}n+o(n))$, and all cumulants of order
$j\ge3$, after normalization by $(qn)^{j/2}$, tend to zero.
The same statement for deterministic step weights proves the
Gaussian finite-dimensional limits by the moment method.
The second- and fourth-cumulant bounds give the uniform increment
estimate $Cq^2n^2$ for fourth moments of complete-band sums.
Linear interpolation after scaling by $\sqrt{qT}$ is therefore
tight. The maximum of $O(T)$ centered single-band sums is
$o_{\PP}(\sqrt{qT})$ by their fourth moment $O(q^2)$ and a union
bound. Partial-band counts are bounded by complete positive band
counts; their means contribute at most $Cq=o(\sqrt{qT})$, since
\eqref{eq:slow-edit-range} implies $q=o(T)$.
This proves the summed line functional limit, including time
reversal and the clamped time change in the comparison lemma.
Its terminal moments converge by the same cumulant estimates.

The $L^p$ comparison now proves the finite path and moment
statements. Exact line means and its variance asymptotic give
\eqref{eq:slow-edit-mean}--\eqref{eq:slow-edit-variance}.
Also $q\log q\le q\log\log M=o(\sqrt{q\log M})$ in the stated
range, justifying both alternative centerings.

For resampling take $q=2k$, $M=N-k$. Then $\log M\sim\log(N/k)$,
and the difference between the displayed resampling centering
and $2a_{\rm H}q\log M$ is
$O(k\log(k+1))=o(\sqrt{k\log N})$.
The label correction is bounded by $2k$, which is also negligible
in every normalized fixed moment. This proves
\eqref{eq:slow-resampling-CLT} and its variance assertion.
\end{proof}

\begin{corollary}[Separating count fluctuations from history fluctuations]
\label{cor:slow-independent-history-CLT}
Independently refresh each coordinate with probability $s_N$.
Suppose $\mu_N=Ns_N\to\infty$ and
$\mu_N(\log\log N)^2=o(\log N)$, and let $K_N$ be the actual
refresh count. With $b_{\rm H}=1+e^{-\gamma}$,
\begin{equation}\label{eq:two-scale-history-CLT}
 \left(\frac{K_N-\mu_N}{\sqrt{\mu_N}},
 \frac{H_N-b_{\rm H}K_N\log(N/K_N)}{\sqrt{\mu_N\log N}}\right)
       \Longrightarrow (Z,G),
\end{equation}
where $Z\sim\mathcal N(0,1)$ and
$G\sim\mathcal N(0,8v_{\rm H})$ are independent. The product
$K_N\log(N/K_N)$ is taken as zero at $K_N=0$.
This holds for spatial and label costs. At the larger fluctuation
scale,
\begin{equation}\label{eq:count-dominated-history-CLT}
 \frac{H_N-b_{\rm H}\mu_N\log(N/\mu_N)}
                    {\sqrt{\mu_N}\log N}
          \Longrightarrow\mathcal N(0,b_{\rm H}^2).
\end{equation}
\end{corollary}
\begin{proof}
Binomial concentration gives $K_N/\mu_N\to1$ in probability,
and its centered characteristic function gives the usual Gaussian
limit for $(K_N-\mu_N)/\sqrt{\mu_N}$, since $s_N\to0$.
Conditional on $K_N=k$, Theorem~\ref{thm:slow-edit-history-CLT}
applies. Its characteristic-function convergence is uniform on
any window $|k/\mu_N-1|\le\varepsilon_N\to0$ having probability
tending to one: otherwise a violating deterministic sequence in
that window would still satisfy \eqref{eq:slow-edit-range} and
contradict the theorem. In that window the conditional
normalization $\sqrt{k\log(N/k)}$ is asymptotic to
$\sqrt{\mu_N\log N}$. Taking conditional characteristic functions
proves the product limit in \eqref{eq:two-scale-history-CLT}.

For $g(x)=x\log(N/x)$, Taylor's formula on the same binomial
window gives
\[
 g(K_N)-g(\mu_N)
   =(\log(N/\mu_N)-1)(K_N-\mu_N)+O_{\PP}(1),
\]
because $g''(x)=-1/x$. The centered history term in
\eqref{eq:two-scale-history-CLT} is smaller than
$\sqrt{\mu_N}\log N$ by a factor $1/\sqrt{\log N}$.
Also $(\log(N/\mu_N)-1)/\log N\to1$.
These observations prove \eqref{eq:count-dominated-history-CLT}.
\end{proof}

\subsection{A complex-intensity comparison for complete configurations}
\label{sec:complex-coefficient-rate}

\status{Working proof: a coefficient bound independent of the number of leaves}

\begin{lemma}[Complex clearance fields]
\label{lem:complex-gap-fields}
The locally finite identity \eqref{eq:positive-gap-comparison}
extends to every complex $\alpha$. For $y\ge1$,
\begin{align}
 |f_\alpha(y)|&\le2R_{\rm o}(y)\le1
                                      &&(\operatorname{Re}\alpha\ge0),
                                      \label{eq:complex-gap-bound}\\
 \left|\frac{f_\alpha(y)}{f(y)}-1\right|
   &\le\exp(Cy e^{-2\operatorname{Re}\alpha})-1
                                      &&(\operatorname{Re}\alpha\ge1).
                                      \label{eq:complex-relative-gap}
\end{align}
Also, for all complex $\lambda$ and $x\ge t>0$,
\begin{equation}\label{eq:complex-parity}
 w_t^{-\lambda}(x)=-e^{2\lambda x}w_t^\lambda(x).
\end{equation}
\end{lemma}
\begin{proof}
At each fixed $y$ the positive convolution identity has only
finitely many nonzero terms, since $f$ and $k_\alpha$ are supported
on $[1,\infty)$. Each such convolution is entire in $\alpha$,
as is the original finite-Poisson series. Equality on positive
real $\alpha$ therefore proves equality throughout the complex
plane. If $u=\operatorname{Re}\alpha\ge0$, then
$|k_\alpha|=k_u\le f$, so
\[
 |f_\alpha|\le(f+k_u)*\sum_{j\ge0}(f*k_u)^{*j}
              \le2\sum_{j\ge0}f^{*(2j+1)}=2R_{\rm o}\le1.
\]
The last bound is Lemma~\ref{lem:factorial-renewal}.
For $u\ge1$, subtract the $j=0$ term $f$ before applying absolute
values. The resulting bound is
$k_u+2\sum_{j\ge1}f^{*(j+1)}*k_u^{*j}$.
The envelope calculation in Lemma~\ref{lem:relative-gap-density}
then gives \eqref{eq:complex-relative-gap}, after increasing $C$.
Finally $w_t^\lambda(x)=\lambda e^{-\lambda x}
\sum_{j\ge0}\lambda^{2j}W_j(t,x)$ gives
\eqref{eq:complex-parity} directly.
\end{proof}

\begin{proposition}[A two-saddle coefficient estimate]
\label{prop:uniform-relative-coefficient}
Let $x_i\ge t_i\ge a>0$, $1\le i\le l$, with $\sum_i x_i=1$.
Write $w_t(x)=t^{-1}f(x/t)$ and
\[
 D_M=M![\lambda^M]e^\lambda\prod_{i=1}^l w_{t_i}^\lambda(x_i),
 \qquad D_\infty=2\prod_{i=1}^l w_{t_i}(x_i).
\]
There is a universal $C<\infty$ such that, whenever
$Ma\ge2$ and $a^{-1}e^{-Ma}\le1$, and $M\equiv l\pmod2$,
\begin{equation}\label{eq:uniform-relative-coefficient}
 \left|\frac{D_M}{D_\infty}-1\right|
 \le C\left\{a^{-1}e^{-Ma}
                 +\sqrt M\exp(-M/2+C/a)\right\}.
\end{equation}
For the opposite parity, $D_M=0$ exactly. The constants do not
depend on the leaf count $l$, their lengths, or their cutoffs.
\end{proposition}
\begin{proof}
Put $P(\lambda)=\prod_iw_{t_i}^\lambda(x_i)$ and
$P_\infty=\prod_iw_{t_i}(x_i)$. On the circle
$\lambda=Me^{i\theta}$, the arc $|\theta|\le\pi/3$ has
$\operatorname{Re}\lambda\ge M/2$. By
\eqref{eq:complex-relative-gap},
\begin{equation}\label{eq:coefficient-core-error}
 \left|\frac{P(Me^{i\theta})}{P_\infty}-1\right|
 \le\exp\left(C\sum_i\frac{x_i}{t_i}
                       e^{-2Mt_i\cos\theta}\right)-1
 \le C a^{-1}e^{-Ma}.
\end{equation}
The last step uses $\sum_i x_i/t_i\le1/a$ and
$a^{-1}e^{-Ma}\le1$; its constant may absorb $e^C$.

If $\operatorname{Re}\lambda\ge0$, the crude bound
\eqref{eq:complex-gap-bound} and the limiting lower envelope give
\[
 |P(\lambda)/P_\infty|
 \le(b^-)^{-l}\exp\left(\kappa\sum_i x_i/t_i\right)
 \le e^{C/a},
\]
since $l\le1/a$. For negative real part, reflect $\lambda$ to
$-\lambda$ and use \eqref{eq:complex-parity}; since the lengths
sum to one,
\begin{equation}\label{eq:coefficient-parity-symmetry}
 e^{-\lambda}P(-\lambda)=(-1)^l e^\lambda P(\lambda).
\end{equation}

Cauchy's coefficient formula on $|\lambda|=M$ is
\[
 D_M=\frac{M!M^{-M}}{2\pi}
       \int_{-\pi}^{\pi}e^{Me^{i\theta}}
               P(Me^{i\theta})e^{-iM\theta}\,d\theta.
\]
On the positive core arc, replacing $P$ by $P_\infty$ incurs
relative absolute error $Ca^{-1}e^{-Ma}$. Indeed
$M!M^{-M}\le C\sqrt M e^{-M}$ and
$\int_{-\pi/3}^{\pi/3}e^{M\cos\theta}\,d\theta
\le Ce^M/\sqrt M$.
The normalized integral of $e^{Me^{i\theta}}e^{-iM\theta}$
on that arc differs from one by at most
$C\sqrt M e^{-M/2}$, because its full-circle coefficient is
exactly one and $\cos\theta\le1/2$ off the arc.

The reflected core around $\theta=\pi$ has the same contribution
multiplied by $(-1)^{M+l}$, by
\eqref{eq:coefficient-parity-symmetry}. On the two remaining
arcs $|\cos\theta|\le1/2$. The crude bound and reflection give
their total relative contribution at most
$C\sqrt M\exp(-M/2+C/a)$.
For matching parity the two core contributions add to two,
proving \eqref{eq:uniform-relative-coefficient} after division
by $D_\infty=2P_\infty$. For opposite parity,
\eqref{eq:coefficient-parity-symmetry} makes $e^\lambda P(\lambda)$
an entire function of parity $l$, so the coefficient is zero.
\end{proof}

\begin{theorem}[Global configuration coupling near the microscopic scale]
\label{thm:near-microscopic-TV}
For an iid size-$M$ circle input of either parity and its
corresponding entrance law, the total-variation distance between
the entire spatial histories at and after cutoff $a$ is at most
\begin{equation}\label{eq:near-microscopic-TV}
 \min\left(1,C\left\{a^{-1}e^{-Ma}
             +\sqrt M\exp(-M/2+C/a)\right\}\right)
\end{equation}
whenever $Ma\ge2$ and $a^{-1}e^{-Ma}\le1$.
In particular, for each fixed $A\ge1$, at
$a=A\log M/M$ this distance is
$O_A(M^{1-A}/\log M)+O(e^{-M/3})$.
More generally it tends to zero whenever $a\to0$ and
$Ma-\log(1/a)\to\infty$.

The same relative estimate applies on any specified finite tree
geometry whose terminal lengths sum to one and whose terminal
cutoffs are at least $a$. It does not compare raw original labels
or trees continued all the way to cutoff zero.
\end{theorem}
\begin{proof}
On every nonempty active sector at cutoff $a$, apply
Proposition~\ref{prop:uniform-relative-coefficient} to its gap
lengths, all with cutoff $a$. Its finite density is $D_M$ and
its entrance density is $D_\infty$, with the same cyclic-root
reference measure. The pointwise relative error integrates to
at most the displayed bound, since the total entrance nonempty
mass is at most one. The difference of the empty masses equals
the negative of the difference of the nonempty masses, so the
same estimate controls total variation in even parity too.
The future history is a deterministic function of the cutoff
configuration, and total variation contracts under this map.

For $a=A\log M/M$, the first term is
$M^{1-A}/(A\log M)$, and $C/a=o(M)$ controls the second.
Under the more general criterion, $Ma\to\infty$, so again
$C/a=o(M)$ and both terms vanish. Finally every revealed split
replaces a parent factor by three daughter factors in the exact
coefficient density. The proposition therefore applies unchanged
on each such tree geometry. Exceptional premature-empty branches
are handled by normalization when that comparison is used.
\end{proof}

\begin{corollary}[Simultaneous finite-count spatial Gaussian limits]
\label{cor:finite-spatial-CLT}
Let $r_M\to0$ and $Mr_M-\log(1/r_M)\to\infty$.
With $c=(2e^\gamma)^{-1}$, for either input parity,
\[
 \sqrt{r_M}\left(\Xi_M(r_M)-\frac c{r_M}\,dx\right)
      \Longrightarrow G
                  \quad\hbox{in }H^{-s}(\T),\qquad s>1/2,
\]
where $G$ is the white-noise field with variance density
$2c^2=1/(2e^{2\gamma})$. The current-count-normalized measure
has the same bridge limit as \eqref{eq:normalized-gaussian}.
Any fixed family of copies obtained by $o_{\PP}(r_M^{-1/2})$
arbitrary edits shares this Gaussian field and bridge jointly.

The covariance convergence \eqref{eq:entrance-covariance} holds
for these finite samples under the same assumptions;
in particular $\operatorname{Var}K_M(r_M)\sim2c^2/r_M$.
The moment assertion uses the relative density bound and a separate
empty-state estimate, not just vanishing TV distance.
Theorem~\ref{thm:full-finite-spatial-CLT} extends this to the full regime
$Mr_M\to\infty$, with finite-intensity centering.
Section~\ref{sec:temporal-spatial-CLT} supplies temporal limits separately.
\end{corollary}
\begin{proof}
Theorem~\ref{thm:near-microscopic-TV} couples the finite configuration
at $r_M$ to the entrance configuration with probability tending
to one. Apply the corresponding, possibly $M$-dependent, scaling
maps and Theorem~\ref{thm:entrance-fluctuations}; TV contraction
transfers both field limits.
For the edited copies, discrepancy contraction bounds the
$H^{-s}$ field difference by $C_s\sqrt{r_M}$ times the number
of edited points. It vanishes in probability; the normalized
bridge transfer follows by the same count normalization.

For covariance transfer, first suppose $r_M\le M^{-1/3}$.
On every nonempty sector, the pointwise relative density error
is $o(1)$. The entrance second moment of each centered field test
is uniformly bounded, so integration bounds the corresponding
nonempty second-moment error by $o(1)$ as well. The empty sector
needs a separate estimate. The last-pair integration in the proof
of Lemma~\ref{lem:full-corner-law} bounds the Poisson empty
probability at intensity $M$ by $Ce^{-c/r_M}$ when $Mr_M\ge1$.
Exact-count conditioning costs at most $e\sqrt M$. Thus the
finite empty probability is at most $C\sqrt M e^{-c/r_M}$;
the entrance bound is $Ce^{-c/r_M}$. The squared centered test
at the empty configuration is at most $C/r_M$, so this part of
the second-moment error is negligible for $r_M\le M^{-1/3}$.

If instead $r_M>M^{-1/3}$, couple the entire future history at
$a=3\log M/M<r_M$. The failure probability is
$O(M^{-2}/\log M)$. Active counts are bounded by $1/r_M$, so
the squared centered test is bounded by $C/r_M$. The resulting
second-moment error is $O(M^{-2}/(r_M\log M))=o(1)$.
First moments and cross moments follow by the same estimates and
polarization. This proves the covariance and count-variance
conclusions. In particular $r_M=\log M/M$ suffices for all the
displayed weak and covariance limits.
\end{proof}

\begin{theorem}[Sharper full-map and finite-family rates]
\label{thm:sharper-map-family-rates}
Put $s_0=\log(1/\sigma)$, $\ell_0=\log(3J)$, and
\[
 \beta_{\rm map}=\frac{s_0}{\ell_0+s_0}=4\chi,\qquad
 \beta_{\rm fam}=\frac{s_0}{\ell_0+3s_0}=4\chi_{\rm f}.
\]
For an even background of size $M=N-1$, the full survivor maps
can be coupled so that
\begin{equation}\label{eq:sharper-full-map-rate}
 \EE\Leb\{z:S_{B_M}(z)\ne S_\infty(z)\}
        \le C\left(\frac{\log N}{N}\right)^{\beta_{\rm map}}.
\end{equation}
The same bound holds for $|p_N-p_\infty|$ and
$\sup_h|\zeta_N(h)-\zeta(h)|$.

For either background parity and any $q$ independent probes, put
$\varepsilon_M=(\log M/M)^{\beta_{\rm fam}}$.
Their complete insertion family can be coupled to its entrance
family with identical spatial configurations and future histories
after $\varepsilon_M$, except on probability at most
\begin{equation}\label{eq:sharper-family-rate}
 \min\{1,C(q+1)^2\varepsilon_M\}.
\end{equation}
In particular, uniformly for odd $N$ and $1\le j\le N/2$,
\begin{equation}\label{eq:sharper-fixed-profile-rate}
 |p_N(j)-p(j)|\le C(j+1)^2
                         \left(\frac{\log N}{N}\right)^{\beta_{\rm fam}}.
\end{equation}
These improve the earlier polynomial exponents by a factor four,
at the stated logarithmic cost. The constants and exponents remain
uncalibrated. The deletion-mass comparison has an additional
empirical-discrepancy term; Corollary~\ref{cor:sharper-basin-rate}
separately optimizes it.
\end{theorem}
\begin{proof}
In the full last-pair geometry coupling of
Lemma~\ref{lem:full-root-rate}, restrict to terminal cutoffs at
least $a=3\log M/M$. The limiting bad mass is
$Ca(3J)^d$. On the good geometry, replace its absolute coefficient
estimate by Proposition~\ref{prop:uniform-relative-coefficient}.
The relative error is $O(M^{-2}/\log M)+O(e^{-M/3})$, independently
of the number $2\cdot3^d$ of leaves. It integrates against a
probability subdensity, so normalization gives
\[
 \tau_M^{\rm last}(d)
       \le Ca(3J)^d+CM^{-2}/\log M.
\]
The same deterministic completion budget as before therefore
bounds the expected full-map disagreement by
$C\{a(3J)^d+\sigma^d+M^{-2}/\log M\}$.
Take
\[
 d=\left\lfloor
       \frac{\log(M/\log M)}{\ell_0+s_0}\right\rfloor.
\]
The first two terms are $O((\log M/M)^{\beta_{\rm map}})$,
and the last is smaller. This proves
\eqref{eq:sharper-full-map-rate}; the two-query argument from
Theorem~\ref{thm:polynomial-noise-rate} transfers it to both noise
profiles. The original-label correction is smaller still.

For the family assertion, start at cutoff
$\varepsilon=\varepsilon_M$ and reveal
\[
 d=\left\lfloor
       \frac{\log(M/\log M)}{\ell_0+3s_0}\right\rfloor
\]
generations in every nonempty cutoff gap. Their expected summed
inverse terminal cutoffs are at most
$C\varepsilon^{-2}(3J)^d$, as in
Lemma~\ref{lem:finite-family-rate}. Thus the limiting probability
of a terminal cutoff below $a$ is at most
\[
 Ca\varepsilon^{-2}(3J)^d\le C\varepsilon.
\]
The empty sector also has probability $O(\varepsilon)$.
On every retained count sector the same uniform relative
coefficient bound applies: all terminal lengths sum to one,
their parity is that of the input, and the cyclic reference
factor is identical in the two laws. Sum against the limiting
probability density and use normalization to control premature
empty finite trees. The entire marked geometry therefore has
coupling error $O(\varepsilon)$.

The expected deterministic origin-disagreement budget, summed
over cutoff gaps, is $C\sigma^d\le C\varepsilon$, using
$\EE\sum x_i^2=O(\varepsilon)$. A union bound for $q$ probes
costs $Cq\varepsilon$. Their failure to occupy distinct
nonadjacent gaps costs $Cq^2\varepsilon$. On the complement,
the separated-insertion argument works simultaneously for all
subsets, and the future histories are identical. This proves
\eqref{eq:sharper-family-rate}; when its right side exceeds one
the statement is vacuous.

For $j$ resamplings use the common background of size $N-j$
and its $2j$ probes. The coupled survivor agreement events give
\eqref{eq:sharper-fixed-profile-rate}; the label/position correction
$j/N$ is absorbed by the displayed bound. No independence of
the coupled configurations is inferred.
\end{proof}

\subsection{Finite spatial fluctuations throughout the mesoscopic regime}
\label{sec:full-finite-spatial-CLT}

\status{Working proof: finite-intensity centering and conditional-count comparison}

\begin{lemma}[Uniform finite-intensity renewal and Gaussian estimates]
\label{lem:finite-Poisson-spatial-CLT}
Let $\alpha\ge1$, $f_\alpha=w_1^\alpha$, and
\[
 U_\alpha=\sum_{j\ge1}f_\alpha^{*j},\quad
 D_\alpha=\sum_{j\ge1}(-1)^{j-1}f_\alpha^{*j},\quad
 c_\alpha=c\exp(-E_1(2\alpha)),\quad c=(2e^\gamma)^{-1}.
\]
There are $C,\theta>0$ such that, uniformly for $\alpha\ge1$,
\begin{equation}\label{eq:uniform-renewal-excess}
 |U_\alpha(y)-c_\alpha|+|D_\alpha(y)|\le Ce^{-\theta y}.
\end{equation}
Writing $e_\alpha=U_\alpha-c_\alpha$, the weighted $L^1$ and
$L^\infty$ norms of $e_\alpha-e_\infty$ and
$D_\alpha-D_\infty$, with weight $e^{\theta y}$, tend to zero
as $\alpha\to\infty$.

Consider a Poisson input of intensity $\lambda$ on the unit circle,
conditioned on either initial parity. If $r\to0$ and
$\alpha=\lambda r\to\infty$, its centered field
\begin{equation}\label{eq:finite-Poisson-Gaussian-field}
 \sqrt r\left(\Xi_{P_\lambda}(r)-\frac{c_\alpha}{r}\,dx\right)
\end{equation}
converges in $H^{-s}(\T)$, $s>1/2$, to white noise with variance
density $2c^2$. Its covariances converge on continuous tests,
and the normalized measure has the bridge limit
\eqref{eq:normalized-gaussian}.
\end{lemma}
\begin{proof}
Put $v_\alpha(y)=e^{-2\alpha y}\ind_{\{y\ge1\}}/y$ and
\[
 \nu_\alpha^- =\sum_{j\ge1}\frac{(-1)^j}{j!}v_\alpha^{*j},
 \qquad
 \nu_\alpha^+ =\sum_{j\ge1}\frac1{j!}v_\alpha^{*j}.
\]
For any fixed $0<\theta<1/2$, their exponentially weighted
$L^1$ and $L^\infty$ norms are bounded by $Ce^{-2\alpha}$.
Indeed the weighted norms of $v_\alpha$ are
$E_1(2\alpha-\theta)$ and at most $e^{-2\alpha+\theta}$;
sum the convolution exponential using one $L^\infty$ factor.

The transform calculation in Lemma~\ref{lem:uniform-renewal}
gives
$U_\alpha-U_\infty=\nu_\alpha^-/2+U_\infty*\nu_\alpha^-$.
Since $\int\nu_\alpha^-=e^{-E_1(2\alpha)}-1$, subtraction of
the constant $c_\alpha$ yields the exact identity
\begin{equation}\label{eq:finite-renewal-excess-identity}
 e_\alpha=e_\infty+\tfrac12\nu_\alpha^-
       +e_\infty*\nu_\alpha^-
       -c\int_y^\infty\nu_\alpha^-(z)\,dz.
\end{equation}
Similarly, from the transform $\widehat D_\alpha
=(1-e^{-E_1(s)+E_1(s+2\alpha)})/2$,
\begin{equation}\label{eq:finite-parity-excess-identity}
 D_\alpha=D_\infty-\tfrac12\nu_\alpha^+
                            +D_\infty*\nu_\alpha^+.
\end{equation}
The factorial tails of $e_\infty$ and $D_\infty=a_0/2$
make their exponentially weighted norms finite. Weighted
convolution and tail-integral bounds prove all asserted estimates,
after decreasing $\theta$ if needed.

Here are the Gaussian-transfer details. Rescale the circumference
to $R=1/r$. Conditional on parity $p$, nonempty sector densities
are $2\prod_i f_\alpha(y_i)$ against the usual cyclic-root measure,
divided by $Z_p=1+(-1)^p e^{-2\lambda}$.
The first and second intensities have exactly the parity renewal
formulas used in Theorem~\ref{thm:entrance-fluctuations}, with
$U,D,c$ replaced by $U_\alpha,D_\alpha,c_\alpha$ and the factor
$Z_p^{-1}$. The weighted estimates just proved make all boundary
and parity remainders uniformly negligible. The covariance
calculation there consequently gives a uniform second-moment
bound on bounded tests and, as $r\to0$, $\alpha\to\infty$,
the limit $2c^2\int\phi\psi$. The bias relative to $c_\alpha/r$
is $O(r^{-1}e^{-\theta/r})+O(r^{-1}e^{-2\lambda})$.

Tilt by the active count and root at a uniform active point.
The resulting parity bridge has density
$\prod_i f_\alpha(y_i)/R_{p,\alpha}(R)$, where
$R_{p,\alpha}=(U_\alpha+(-1)^{p+1}D_\alpha)/2$.
Its stopped-prefix ratio to iid gaps is the same outside-renewal
ratio as in \eqref{eq:palm-renewal-bridge}. Stop at distance
$R-R^{1/4}$. Equation~\eqref{eq:uniform-renewal-excess} makes
that ratio tend uniformly to one when the unexposed arc has
length at least $R^{1/4}/2$. The uniform exponential gap envelope
bounds the exceptional overshoot probability by
$CR e^{-cR^{1/4}}$. The count tilt has TV cost $O(R^{-1/2})$ by
the second-moment bound. The remaining arc contains at most
$R^{1/4}+1$ points and is negligible after division by $\sqrt R$.

It remains an iid triangular-array renewal CLT. The gap mean is
$m_\alpha=1/c_\alpha$, by \eqref{eq:finite-intensity-mean} and
scaling. Uniform exponential moments and
\eqref{eq:positive-gap-comparison} imply convergence of every
gap moment to its limiting value; in particular
$\operatorname{Var}Y_\alpha\to2/c$. The partial-sum inversion,
maximal inequality, and step-function approximation in the proof
of Theorem~\ref{thm:entrance-fluctuations} are uniform under these
moment bounds. They give the Gaussian field centered at
$c_\alpha/r$, with limiting variance
$(2/c)/(1/c)^3=2c^2$. The uniform covariance bounds give Sobolev
tightness by the same Fourier argument. Finally $rK\to c$ and
normalization gives the bridge limit. This also proves second
moment convergence, hence uniform integrability of the squared
centered field on each continuous test.
\end{proof}

\begin{lemma}[Conditional local limits for the original point count]
\label{lem:mesoscopic-depoisson-TV}
Let $r=r_M\to0$ and $\alpha=Mr\to\infty$. The active configuration
from exactly $M$ iid input points and that from a Poisson input
of intensity $M$, conditioned to have the parity of $M$, have
total-variation distance tending to zero. The same holds for
their entire future histories after $r$.

More precisely, let $\mathcal C$ be the active configuration of
the unconditioned Poisson input, and $P$ its initial count.
On configurations having the parity of $M$,
\begin{equation}\label{eq:conditional-count-local-CLT}
 \sqrt M\,\PP(P=M\mid\mathcal C)
          \longrightarrow\frac2{\sqrt{2\pi}}
                                      \quad\hbox{in probability}.
\end{equation}
There is an event $\mathcal G$ with
$\PP(\mathcal G^c)\le Ce^{-c/r}$ on which this conditional
probability is at most $C/\sqrt M$. Constants are uniform for
sufficiently large $\alpha$.
\end{lemma}
\begin{proof}
Conditional on a nonempty $\mathcal C$, its shielded gap interiors
are independent. If $x_i=ry_i$ is a gap length, count its left
endpoint and its original interior points as $K_i$, so that
$P=\sum_iK_i$. Each $K_i$ is odd and has conditional mgf
\begin{equation}\label{eq:conditional-gap-count-mgf}
 \EE(e^{uK_i}\mid\mathcal C)
  =e^{\alpha y_i(e^u-1)}
                       \frac{f_{\alpha e^u}(y_i)}{f_\alpha(y_i)}.
\end{equation}
Thus its $j$th cumulant is
$\alpha y_i+\mathcal D^j\log f_\alpha(y_i)$, with
$\mathcal D=\alpha\partial_\alpha$.

For $1\le j\le4$ and large $\alpha$,
\begin{equation}\label{eq:gap-count-cumulant-error}
 |\mathcal D^j\log f_\alpha(y)|
 \le C_j(1+y)^j e^{-c\alpha}
                         \exp(C_j y e^{-c\alpha}).
\end{equation}
To verify this bound without differentiating an asymptotic
formula, apply Cauchy's derivative estimate to
$f_z(y)-f(y)$ on $|z-\alpha|\le\alpha/4$, using
\eqref{eq:complex-relative-gap}. It bounds
$\alpha^j|\partial_\alpha^j f_\alpha(y)|/f(y)$ by a constant
times $\exp(Cy e^{-3\alpha/2})-1$. At real $\alpha$,
$f_\alpha\ge(1-e^{-2\alpha})f$. Express the logarithmic
derivatives as polynomials in these derivative ratios. This
gives \eqref{eq:gap-count-cumulant-error}, including a bound of
the same form for the square of its $j=2$ error.

The uniform exponential gap envelope absorbs every polynomial
and the small exponential factor in this bound. The one-gap
intensity formula, with outside factor $U_r^M\le C/r$, therefore
gives
\begin{equation}\label{eq:summed-count-cumulant-errors}
 \EE\left[\ind_{\{\mathcal C\ne\varnothing\}}
       \sum_i|\mathcal D^j\log f_\alpha(y_i)|\right]
                \le Cr^{-1}e^{-c\alpha},\qquad 1\le j\le4.
\end{equation}
The singleton contribution obeys the same bound by the gap
envelope; the squared $j=2$ error does as well.

Write $\mu_{\mathcal C}=\EE(P\mid\mathcal C)$ and
$s_{\mathcal C}^2=\operatorname{Var}(P\mid\mathcal C)$.
Since $\sum_i\alpha y_i=M$ on a nonempty configuration,
\eqref{eq:summed-count-cumulant-errors} gives
$s_{\mathcal C}^2/M\to1$ in probability outside the empty event.
The Poisson empty probability is at most $Ce^{-c/r}$.
The total-variance identity is particularly useful here:
\[
 \operatorname{Var}\mu_{\mathcal C}
   =M-\EE s_{\mathcal C}^2
   \le M\PP(\mathcal C=\varnothing)+Cr^{-1}e^{-c\alpha}=o(M).
\]
Together with $\EE\mu_{\mathcal C}=M$, this proves
$(\mu_{\mathcal C}-M)/\sqrt M\to0$ in $L^2$.
No differentiation of the empty-state probability is required.

The conditional fourth central moment of $K_i$ is its fourth
cumulant plus three times its variance squared. Using
\eqref{eq:gap-count-cumulant-error} and
$\EE\sum_i x_i^2\le Cr$ gives
\begin{equation}\label{eq:conditional-count-Lindeberg}
 \frac1{M^2}\EE\left[
   \ind_{\{\mathcal C\ne\varnothing\}}
   \sum_i\EE\{(K_i-\EE[K_i\mid\mathcal C])^4\mid\mathcal C\}
                      \right]\le C(r+M^{-1})\longrightarrow0.
\end{equation}
Conditional independence, a Taylor expansion of characteristic
functions, and Cauchy--Schwarz now give, for each fixed real $z$,
\begin{equation}\label{eq:conditional-count-CF}
 e^{-izM/\sqrt M}\EE(e^{izP/\sqrt M}\mid\mathcal C)
                         \longrightarrow e^{-z^2/2}
                                      \quad\hbox{in probability}.
\end{equation}
For example the summed third-moment Taylor remainder divided by
$M^{3/2}$ is bounded by the square root of
\eqref{eq:conditional-count-Lindeberg} times the normalized
conditional variance; the maximum individual variance is
controlled by the same fourth-moment bound.

We supply a uniform Fourier tail bound to obtain a local limit.
If $1\le y_i\le2$, shielding gives exactly
$K_i\sim\operatorname{Poisson}(\alpha y_i)$ conditioned to be odd.
For all sufficiently large $a$, its characteristic function obeys
\begin{equation}\label{eq:odd-Poisson-CF-bound}
 \left|\frac{\sinh(ae^{i\theta})}{\sinh a}\right|
      \le\exp\{-ca\operatorname{dist}(\theta,\pi\mathbb Z)^2\}.
\end{equation}
Here is an elementary verification. By parity it suffices to take
$0\le\theta\le\pi/2$. For $\theta\le c_1/\sqrt a$, use two
independent copies of the odd Poisson variable and
$1-\cos z\ge z^2/2-z^4/24$. Its variance is at least $a/2$ and
its fourth central moment is at most $Ca^2$, from the derivatives
of $\log\sinh(ae^u)$. This proves the bound near zero.
For the remaining angles with $\cos\theta\ge1/2$, use
\[
 |\sinh(ae^{i\theta})/\sinh a|
 \le e^{-a(1-\cos\theta)}
                 \frac{1+e^{-2a\cos\theta}}{1-e^{-2a}}.
\]
The extra logarithm is $O(e^{-a})$, smaller than a fixed fraction
of $a(1-\cos\theta)$ there. If $\cos\theta<1/2$, the same
display is at most $3e^{-a/2}$, which proves the remaining range.

There is a $\delta>0$ such that
\begin{equation}\label{eq:short-gap-length-event}
 \mathcal G=\left\{\sum_{i:1\le y_i\le2}x_i\ge\delta\right\},
 \qquad \PP(\mathcal G^c)\le Ce^{-c/r}.
\end{equation}
For detail, tilt $f_\alpha(y)$ by
$e^{-\tau y\ind_{[1,2]}(y)}$, with fixed $\tau>0$.
Its mass is uniformly below one for large $\alpha$, by
$f_\alpha\to f$ on $[1,2]$. The uniform exponential gap moment
then permits a fixed $\eta>0$ for which the tilted density,
multiplied by $e^{\eta y}$, still has mass below one and bounded
supremum. Its renewal sum is consequently bounded by
$Ce^{-\eta R}$ at $R=1/r$. The cyclic density includes a factor
$R/k\le R$, so exponential Markov gives
$CR e^{-(\eta-\tau\delta)R}$ for the event in question on
nonempty sectors. Choose $\delta<\eta/(2\tau)$ and add the
exponentially small empty probability. This proves
\eqref{eq:short-gap-length-event}.

On $\mathcal G$, independence of the $K_i$ and
\eqref{eq:odd-Poisson-CF-bound} give
\[
 |\EE(e^{i\theta P}\mid\mathcal C)|
       \le e^{-cM\operatorname{dist}(\theta,\pi\mathbb Z)^2}.
\]
The other gaps have characteristic-function modulus at most one.
For the matching parity, Fourier inversion reduces to twice the
integral over $[-\pi/2,\pi/2]$. Set $z=\theta\sqrt M$ and apply
\eqref{eq:conditional-count-CF}; the last Gaussian bound supplies
an integrable dominating function on $\mathcal G$. It yields
\eqref{eq:conditional-count-local-CLT} and the bound
$C/\sqrt M$ on that event. Conditioning on either parity costs
at most a bounded factor and preserves these conclusions.

Finally let $\nu_M$ be the exact-count configuration law and
$\nu_M^{\rm P}$ the Poisson law conditioned on the matching parity.
Their Radon--Nikodym derivative is
\begin{equation}\label{eq:mesoscopic-count-density-ratio}
 \frac{d\nu_M}{d\nu_M^{\rm P}}(\mathcal C)
  =\frac{\PP(P\equiv M\!\pmod2)\PP(P=M\mid\mathcal C)}
               {\PP(P=M)}\longrightarrow1
                                      \quad\hbox{in probability}.
\end{equation}
Use $\PP(P=M)\sim(2\pi M)^{-1/2}$ and parity probability tending
to $1/2$. These are nonnegative densities of total mass one;
$\EE|R-1|=2\EE(1-R)_+$ upgrades their convergence to $L^1$.
This is the claimed TV convergence. Applying the deterministic
future evolution proves its path version.
\end{proof}

\begin{theorem}[The full finite mesoscopic spatial central limit theorem]
\label{thm:full-finite-spatial-CLT}
Let $r_M\to0$ and $Mr_M\to\infty$, with either parity of $M$.
Put $\alpha_M=Mr_M$ and
\[
 \rho_M=\frac{c\exp(-E_1(2\alpha_M))}{r_M}
        =M q(\alpha_M),\qquad c=(2e^\gamma)^{-1},
\]
where $q$ is the early density profile in \eqref{eq:fluid-ode}.
Then, for each $s>1/2$,
\begin{equation}\label{eq:full-finite-spatial-CLT}
 \sqrt{r_M}\bigl(\Xi_M(r_M)-\rho_M\,dx\bigr)
            \Longrightarrow G\quad\hbox{in }H^{-s}(\T),
\end{equation}
where $G$ is white noise of variance density $2c^2$.
For continuous real tests $\phi,\psi$,
\begin{equation}\label{eq:full-finite-spatial-covariance}
 r_M\operatorname{Cov}(\langle\Xi_M(r_M),\phi\rangle,
                       \langle\Xi_M(r_M),\psi\rangle)
          \longrightarrow2c^2\int_{\T}\phi\psi.
\end{equation}
The centered means in \eqref{eq:full-finite-spatial-CLT} tend to
zero on each continuous test. The current-count-normalized field
has the bridge limit \eqref{eq:normalized-gaussian}, and any
fixed family of copies with $o_{\PP}(r_M^{-1/2})$ arbitrary edits
shares both Gaussian limits jointly.

The finite-intensity centering matters. It may be replaced by
$c/r_M$ only when
$r_M^{-1/2}E_1(2Mr_M)\to0$; that condition is automatic in
Corollary~\ref{cor:finite-spatial-CLT}, but not throughout
$Mr_M\to\infty$. Temporal limits are proved separately in
Section~\ref{sec:temporal-spatial-CLT}.
\end{theorem}
\begin{proof}
Lemma~\ref{lem:mesoscopic-depoisson-TV} compares the exact input
to the matching-parity Poisson input at the same intensity $M$.
Apply Lemma~\ref{lem:finite-Poisson-spatial-CLT} and TV contraction
to the centered fields. The identity for $\rho_M$ follows from
$\operatorname{Ein}(2\alpha)=\gamma+\log(2\alpha)+E_1(2\alpha)$.
Since $r_MK_M\to c$, normalization gives the bridge. Discrepancy
contraction gives the arbitrary-edit assertion as in
Corollary~\ref{cor:finite-spatial-CLT}.

To justify moments, let $R_M$ be the density in
\eqref{eq:mesoscopic-count-density-ratio}. On the short-gap event
$\mathcal G$, Fourier inversion in the preceding proof gives
$R_M\le C$, while everywhere $R_M\le C\sqrt M$.
For $r_M\le M^{-1/3}$, the Poisson squared centered test is
uniformly integrable by its Gaussian and second-moment limits.
Thus $R_M\to1$ in probability and the boundedness on $\mathcal G$
transfer its first and second moments there, by truncation.
On $\mathcal G^c$ the squared centered test is at most $C/r_M$;
its exact-count contribution is at most
$C\sqrt M r_M^{-1}e^{-c/r_M}=o(1)$, by
\eqref{eq:short-gap-length-event}. The Poisson contribution
vanishes as well.

For $r_M>M^{-1/3}$, use
Corollary~\ref{cor:finite-spatial-CLT}, since its cutoff condition
is then automatic. The change from $c/r_M$ to $\rho_M$ is
negligible even in the centered mean: here $Mr_M>M^{2/3}$ and
\[
 \sqrt{r_M}|\rho_M-c/r_M|
       \le C r_M^{-1/2}E_1(2Mr_M)\longrightarrow0.
\]
Polarization proves the cross-covariance claim. Finally
$1-e^{-E_1(2Mr_M)}\sim E_1(2Mr_M)$, so the stated condition for
replacing the centering is exact. For example, at
$Mr_M=\log\log M$, the scaled centering difference has order
$\sqrt M/[(\log M)^2(\log\log M)^{3/2}]$ and diverges.
\end{proof}

\subsection{Temporal spatial Gaussian limits}
\label{sec:temporal-spatial-CLT}

\status{Working proof: covariance kernel and cumulative spatial path limits}

\begin{lemma}[Independent blocks for a bounded range of cutoffs]
\label{lem:temporal-renewal-blocks}
Fix $0<a\le b<\infty$. For $\alpha\in[1/a,\infty]$, set
\[
 h_a^\alpha(x)=a^{-1}f_{\alpha a}(x/a),\qquad
 \mu_\alpha(t)=\frac{c_{\alpha t}}t,
 \qquad c_z=c\exp(-E_1(2z)),
\]
with $f_\infty=f$, $c_\infty=c=(2e^\gamma)^{-1}$.
Take iid gaps $Y_i$ of density $h_a^\alpha$, and stop a block at
$J=\min\{i:Y_i>b\}$. Its length is $L=\sum_{i=1}^JY_i$;
its point core is
$\{0,Y_1,Y_1+Y_2,\ldots,Y_1+\cdots+Y_{J-1}\}$.
Let $A(t)$ be the number of these points remaining after their
internal greedy matching through cutoff $t$, $a\le t\le b$.
Successive blocks and their reward processes $A$ are iid, and
\begin{equation}\label{eq:temporal-block-mean}
 \frac{\EE A(t)}{\EE L}=\mu_\alpha(t).
\end{equation}
For some $\eta>0$, depending only on $a,b$,
$\sup_{\alpha\ge1/a}\EE e^{\eta(L+J)}<\infty$.
As $\alpha\to\infty$, the joint law of $L$ and any finite vector
$(A(t_1),\ldots,A(t_\ell))$ converges, with every fixed moment,
to its law with gap density $h_a^\infty$.
\end{lemma}
\begin{proof}
All active gaps at cutoff $a$ are at least $a$. A gap longer
than $b$ cannot be crossed by a collision through cutoff $b$:
later gaps are sums of earlier gaps. Consequently each point
core evolves independently of the other cores. The long gap at
its end is included in the block length, and every gap's left
endpoint belongs to exactly one core. This convention prevents
any duplication of boundary points. Successive blocks use
disjoint strings of iid gaps, proving their independence.

The probability $p_\alpha=\PP(Y_i>b)$ is bounded below uniformly
for $\alpha\ge1/a$, by the uniform exponential gap envelope.
The number $J$ is geometric with parameter $p_\alpha$; its short
gaps are at most $b$, and its last gap has a uniform exponential
moment. Decreasing $\eta$ proves the stated exponential bound.
The weighted convergence of $f_\alpha$ to $f$ gives convergence
of each fixed finite gap string. Greedy matching on a fixed
finite string is locally constant away from ties and threshold
equalities, which form a finite union of hyperplanes in its
continuous gap variables. Truncate $J$ and $L$, apply this
continuity, and remove the truncation using the uniform
exponential bound. This proves the assertion about all moments.

For clarity, the mean identity uses consistency of the one-time
renewal laws with the dynamics. The line Poisson process of
intensity $\alpha$ has, at cutoff $t$, stationary renewal gaps
$h_t^\alpha$ and intensity $c_{\alpha t}/t$. One can verify this
directly from the already proved circle sector densities: let
the circumference tend to infinity at fixed intensity $\alpha$.
The outside renewal factor tends to its constant by
\eqref{eq:uniform-renewal-excess}; the rooted local gap law is
the product of $h_t^\alpha$. Gaps longer than $b$ occur on both
sides of each bounded interval, so deterministic evolution
through $b$ passes to this local limit. The same argument for
$\alpha=\infty$ is Lemma~\ref{lem:line-entrance-dilation}.
The renewal-reward law applied to the blocks therefore gives
\eqref{eq:temporal-block-mean}. The stationary version has only
a length-biased boundary block, which does not affect that ratio.
\end{proof}

For $s,t>0$, choose any $a\le\min(s,t)$ and
$b\ge\max(s,t)$, and use the preceding block with
$\alpha=\infty$. Define
\begin{equation}\label{eq:temporal-covariance-block}
 V(s,t)=\frac{\EE\big[(A(s)-(c/s)L)(A(t)-(c/t)L)\big]}{\EE L}.
\end{equation}
The next proof shows that this definition does not depend on
$a,b$. It is an explicit expectation over a finite geometric
string of renewal gaps, not an evaluation in closed form.

\begin{theorem}[Finite-dimensional temporal spatial fluctuations]
\label{thm:temporal-spatial-CLT}
Fix $0<t_1<\cdots<t_\ell<\infty$. Let $r_M\to0$ and
$Mr_M\to\infty$, with either parity of $M$, and put
\[
 G_M(t)=\sqrt{r_M}\big(\Xi_M(r_Mt)-Mq(Mr_Mt)\,dx\big).
\]
For every $u>1/2$, the vector
$(G_M(t_1),\ldots,G_M(t_\ell))$ converges in
$(H^{-u}(\T))^\ell$ to a centered Gaussian vector of spatial
white-noise fields with
\begin{equation}\label{eq:temporal-spatial-covariance}
 \EE G(s,\phi)G(t,\psi)=V(s,t)\int_\T\phi\psi.
\end{equation}
The centered means and cross-covariances converge for continuous
real tests. The same assertion holds for either entrance law,
with centering $c/(r t)$, and for the parity-conditioned Poisson
input of intensity $M$, with the displayed finite-intensity
centering.

The kernel $V$ is continuous, symmetric, and positive
semidefinite, with
\begin{equation}\label{eq:temporal-covariance-scaling}
 V(t,t)=\frac{2c^2}{t},\qquad
 V(\lambda s,\lambda t)=\lambda^{-1}V(s,t).
\end{equation}
For current-count-normalized measures, the joint limits of
$r_M^{-1/2}(\widehat\Xi_M(r_Mt)-dx)$ have covariance
\begin{equation}\label{eq:temporal-bridge-covariance}
 \frac{st}{c^2}V(s,t)
       \left(\int_\T\phi\psi-\int_\T\phi\int_\T\psi\right).
\end{equation}
Any fixed family obtained by $o_{\PP}(r_M^{-1/2})$ arbitrary
edits shares these fields and bridges jointly. This theorem
asserts convergence at finitely many cutoff ratios; it does
not by itself assert path tightness. Theorem~\ref{thm:temporal-spatial-functional-CLT}
supplies a separate functional limit.
\end{theorem}
\begin{proof}
We first prove the Poisson assertion, allowing $\alpha=Mr\to\infty$
or $\alpha=\infty$ for the entrance law. Rescale the circle to
circumference $R=1/r$, so that its initial intensity is
$\alpha$ and its observation cutoffs are the fixed $t_i$.
Use the preceding blocks with $a=t_1$, $b=t_\ell$.
Write $m_\alpha=\EE L$ and
$Z_i(t)=A_i(t)-\mu_\alpha(t)L_i$. These are centered iid vectors
with uniform exponential moments, and their covariance matrices
converge as $\alpha\to\infty$.

Here are the needed vector renewal-reward details. On the line
rooted at a block boundary, write
$S_n=\sum_{i\le n}L_i$ and $N(x)=\max\{n:S_n\le x\}$.
The active count in $[0,Rx)$ at any of the specified cutoffs,
minus $\mu_\alpha(t)Rx$, is
\begin{equation}\label{eq:temporal-reward-inversion}
 \sum_{i\le N(Rx)}Z_i(t)+O_{\PP}(\log R),
\end{equation}
uniformly over a fixed finite set of $x$'s. Indeed the boundary
error is bounded by a constant times the largest $L_i+J_i$
encountered, and the uniform exponential moment bounds that
maximum by $O_{\PP}(\log R)$. Chebyshev's inequality gives
$N(Rx)-Rx/m_\alpha=O_{\PP}(\sqrt R)$.
For any deterministic window of $O(\sqrt R)$ indices around
$\lfloor Rx/m_\alpha\rfloor$, Kolmogorov's maximal inequality
bounds the centered reward sum divided by $\sqrt R$ by
$o_{\PP}(1)$. Increase the window's constant after applying
Chebyshev. Thus the random index in
\eqref{eq:temporal-reward-inversion} can be replaced by
$\lfloor Rx/m_\alpha\rfloor$.
The triangular-array iid central limit theorem, justified by the
uniform exponential moments, gives vector Brownian increments
in the spatial coordinate, with covariance per unit length
$\EE[Z(s)Z(t)]/\EE L$ in the limit. This proves
\eqref{eq:temporal-spatial-covariance} for spatial step tests.

To transfer from the line to the circle, tilt by the active count
at $a$ and root there. Its gap bridge is the finite-intensity
version of \eqref{eq:palm-renewal-bridge}. Stop the iid gap
prefix at distance $R-R^{1/4}$. The uniform renewal remainder
in Lemma~\ref{lem:finite-Poisson-spatial-CLT} couples this entire
prefix to the bridge in TV, with error tending to zero; removing
the count tilt costs $O(R^{-1/2})$. There may be a block crossing
the root and one crossing the end of the prefix. Discard these
blocks and the final arc. The two crossing blocks have
$O_{\PP}(\log R)$ points, and the final arc has at most
$C R^{1/4}+1$ points. Long-gap barriers make every other block's
whole future through $b$ identical in the line and circle.
The discarded contribution is $o_{\PP}(\sqrt R)$ simultaneously
at the specified times. Restore the independent uniform rotation;
the limiting covariance is rotation invariant.

Uniform approximation of continuous tests by step functions is
justified by the one-time second-moment bounds in
Lemma~\ref{lem:finite-Poisson-spatial-CLT}, at each of the finitely
many cutoffs. The same bounds hold for bounded measurable tests:
their proof bounds the absolute covariance kernel by an integrable
function. Apply Cauchy--Schwarz to a finite linear combination
of times. This avoids multiplying a deterministic modulus of
continuity by $\sqrt R$. Marginal Sobolev tightness gives tightness
in the finite product space, proving the field assertion.
The one-time Gaussian and second-moment limits make the squared
centered tests uniformly integrable. This also gives uniform
integrability of products at two times, so the joint weak limit
implies convergence of all stated cross-covariances and means.

This construction uses the same line history for any allowed
$a,b$. Its spatial covariance density is therefore independent
of that choice, proving the assertion following
\eqref{eq:temporal-covariance-block}. Positive semidefiniteness
is immediate from the covariance formula on a common block.
The one-time theorem gives $V(t,t)=2c^2/t$.
Dominated convergence on a common block, using absence of a
collision at a prescribed time and the exponential moment,
gives continuity. Dilation invariance of the line entrance law
scales block lengths by $\lambda$ and leaves counts unchanged;
\eqref{eq:temporal-covariance-block} gives the scaling relation.

For exact input count $M$, Lemma~\ref{lem:mesoscopic-depoisson-TV}
couples the full future from $r t_1$ to the matching-parity
Poisson future, with error tending to zero. This transfers the
joint field law. Moments require the additional argument already
used for Theorem~\ref{thm:full-finite-spatial-CLT}.
If $r\le M^{-1/3}$, the exact/Poisson density ratio at $rt_1$
is bounded on its short-gap event and is everywhere at most
$C\sqrt M$. Each centered squared test at any fixed $t_i$ is
bounded by $C/r$, while the bad event has probability at most
$Ce^{-c/r}$. Its contribution is therefore
$C\sqrt M r^{-1}e^{-c/r}=o(1)$. On the good event, the bounded
density ratio, its convergence to one in probability, and the
Poisson squared-test uniform integrability transfer first and
second moments, including products at distinct times.
If $r>M^{-1/3}$, use the full entrance-future coupling at
$3\log M/M<rt_1$, whose failure is $O(M^{-2}/\log M)$.
The deterministic squared-field bound $C/r$ makes its moment
error negligible. Moreover $Mr t_i\to\infty$ faster than
$M^{2/3}$ up to a fixed factor, so finite-intensity and entrance
centerings differ by $o(r^{-1/2})$. This proves all exact-count
moment assertions.

Finally $r\Xi_M(rt)(\T)\to c/t$ jointly at the finitely many
times. Subtracting the total-count fluctuation and dividing by
$c/t$ gives \eqref{eq:temporal-bridge-covariance}.
Discrepancy contraction bounds the field difference uniformly
in the cutoff by $C_u\sqrt r$ times the number of edited points.
This gives the shared Gaussian limits for the edited copies.
\end{proof}

\begin{lemma}[Temporal equicontinuity of block sums]
\label{lem:temporal-block-tightness}
Fix $0<a<b<\infty$ and use the blocks of
Lemma~\ref{lem:temporal-renewal-blocks}. For
$Z_\alpha(t)=A(t)-\mu_\alpha(t)L$ there is $C=C(a,b)$ such that
\begin{equation}\label{eq:temporal-block-increment}
 \EE|Z_\alpha(t)-Z_\alpha(s)|^k
       \le k!C^k|t-s|,\qquad k\ge2,
\end{equation}
uniformly for $\alpha\ge1/a$ and $a\le s,t\le b$.
If $\alpha_n\to\infty$, the sequential iid block sum
\[
 B_n(x,t)=n^{-1/2}\sum_{i\le\lfloor nx\rfloor}Z_{\alpha_n,i}(t),
       \qquad (x,t)\in[0,1]\times[a,b],
\]
has a continuous Gaussian limit in the uniform topology after
interpolation of its vanishing jumps. Its covariance is
\[
 (x\wedge y)\EE[Z_\infty(s)Z_\infty(t)].
\]
The same holds for the constant array $\alpha_n=\infty$.
\end{lemma}
\begin{proof}
Conditional on $J=j$, the $j-1$ short gaps are iid with density
$h_a^\alpha\ind_{[a,b]}/\PP(Y\le b)$, uniformly bounded in
$\alpha\ge1/a$. Every possible collision separation inside the
core is a sum of a nonempty consecutive string of these gaps.
A convolution of probability densities, one bounded by $C$,
is bounded by $C$. A union bound over at most $j^2$ pairs gives
\[
 \PP\{A(s)\ne A(t)\mid J=j\}\le Cj^2|t-s|.
\]
Since $|A(s)-A(t)|\le j$, the geometric tail of $J$ implies
\[
 \EE|A(s)-A(t)|^k
   \le C|t-s|\EE J^{k+2}\le k!C_1^k|t-s|.
\]
Here the factor $(k+2)!/k!$ is absorbed into $C_1^k$.
The last gap is independent of the short gaps conditional on
$J$ and has a uniform exponential moment. Also
\[
 \mu_\alpha'(t)
       =\frac{c_{\alpha t}}{t^2}(e^{-2\alpha t}-1)
\]
is uniformly bounded on $[a,b]$. Thus the increment of
$\mu_\alpha(t)L$ satisfies the same factorial moment bound,
proving \eqref{eq:temporal-block-increment}.

We give the tightness argument to specify the temporal
quantifiers. Centered independent increments satisfying
\eqref{eq:temporal-block-increment} obey the maximal Bernstein
bound
\begin{equation}\label{eq:temporal-maximal-Bernstein}
 \PP\left\{\max_{m\le n}\left|
       \sum_{i\le m}(Z_i(t)-Z_i(s))\right|
       >C\big(\sqrt{n|t-s|v}+v\big)\right\}\le2e^{-v},
       \qquad v\ge1.
\end{equation}
Indeed their centered exponential series is bounded by
$\exp(C|t-s|u^2/(1-C|u|))$ for $|u|<1/C$;
apply the exponential supermartingale maximal inequality and
optimize $u$, then repeat for $-u$.

Rescale $[a,b]$ to $[0,1]$ for the following estimates. On the
dyadic grid of mesh $2^{-j}$ there are $2^j$ adjacent pairs.
Use \eqref{eq:temporal-maximal-Bernstein} with
$v=C_2(j+y+1)$ and a union bound. Except on probability
$Ce^{-y}$, simultaneously at every level $j\ge0$ the maximum
normalized increment over all such pairs and all partial-sum
indices is at most
\[
 C\left(\sqrt{(j+y+1)2^{-j}}
                         +\frac{j+y+1}{\sqrt n}\right).
\]
Sum from $j_0$ through
$K_n=\lceil3\log_2(n+1)\rceil$. The first terms have sum at
most $C\sqrt{(j_0+y+1)2^{-j_0}}$; the second terms have sum
$O((K_n^2+(y+1)K_n)/\sqrt n)=o(1)$ for fixed $y$.
Within a finest-grid interval $I$, monotonicity of each $A_i$
bounds every uncentered partial-sum increment by
$\sum_{i\le n}(A_i(\inf I)-A_i(\sup I))$.
The finest-grid Bernstein bound controls its centered version;
its drift is at most $C2^{-K_n}\sum_{i\le n}L_i$.
After division by $\sqrt n$, both terms tend to zero in
probability, uniformly over the finest intervals.
Hence
\begin{equation}\label{eq:temporal-sequential-modulus}
 \lim_{\delta\downarrow0}\limsup_{n\to\infty}
 \PP\left\{\sup_x\sup_{|t-s|\le\delta}
        |B_n(x,t)-B_n(x,s)|>\varepsilon\right\}=0.
\end{equation}
The initial value at $a$ has a uniform exponential moment.
The same chaining estimates also bound
$\EE\sup_{x,t}|B_n(x,t)|^p$ uniformly for each fixed $p$;
integrate their exponential tail bounds.

At any finite grid of times, the iid vector partial-sum
invariance principle follows from the uniform fourth moments,
the maximal inequality, and the iid triangular-array central
limit theorem. It gives tightness and continuous limits in the
spatial index $x$. Combine this with
\eqref{eq:temporal-sequential-modulus} to get tightness on the
whole rectangle: a finite time grid controls spatial moduli,
and the temporal modulus fills its gaps.
Finite-dimensional limits identify the displayed covariance.
Finally the largest block reward among the first $n$ blocks is
$O_{\PP}(\log n)$. All spatial-index and temporal jumps divided
by $\sqrt n$ therefore vanish, so interpolation has no effect
on the limit.
\end{proof}

\begin{theorem}[A temporal functional limit for cumulative spatial counts]
\label{thm:temporal-spatial-functional-CLT}
Fix $0<a<b<\infty$. For $r=r_M\to0$, $Mr\to\infty$, define
\begin{equation}\label{eq:temporal-cumulative-field}
 F_M(x,t)=\sqrt r\big(\Xi_M(rt)([0,x))
                              -Mq(Mrt)x\big),
       \quad 0\le x\le1,\quad a\le t\le b.
\end{equation}
After continuous interpolation of its jumps, $F_M$ converges
in $C([0,1]\times[a,b])$ to a centered continuous Gaussian
field $F$ with
\begin{equation}\label{eq:temporal-cumulative-covariance}
 \EE F(x,s)F(y,t)=(x\wedge y)V(s,t).
\end{equation}
The interpolation can be chosen at uniform distance at most
$3\sqrt r$ from $F_M$. Consequently the counting-measure fields
$G_M(t)$ in Theorem~\ref{thm:temporal-spatial-CLT} converge in
$D([a,b],H^{-u}(\T))$ for every $u\ge1$, to a process with
continuous paths and covariance \eqref{eq:temporal-spatial-covariance}.

The normalized cumulative fields
\[
 r^{-1/2}\left(\widehat\Xi_M(rt)([0,x))-x\right)
\]
converge uniformly after interpolation to
$(t/c)(F(x,t)-xF(1,t))$. Their covariance is
\[
 \frac{st}{c^2}V(s,t)(x\wedge y-xy).
\]
Both assertions hold for either entrance parity, centered at
$c/(rt)$, and for the parity-conditioned Poisson input with its
finite-intensity centering. Any fixed family with
$o_{\PP}(r^{-1/2})$ arbitrary edits shares all these path limits.
This cumulative-field argument proves path tightness for $u\ge1$.
Corollary~\ref{cor:full-sobolev-temporal-CLT} separately extends it
to the full range $u>1/2$ using uniform Fourier maximal bounds.
\end{theorem}
\begin{proof}
First consider the ordinary line renewal process of blocks,
with intensity parameter $\alpha\to\infty$ or $\alpha=\infty$,
on an interval of length $R=1/r$. Its lengths have uniform
exponential moments and are at least $a$. For
$m_\alpha=\EE L$ and $N(Rx)=\max\{j:S_j\le Rx\}$, the
partial-sum maximal inequality and inversion give
\[
 \sup_{0\le x\le1}|N(Rx)-Rx/m_\alpha|=O_{\PP}(\sqrt R).
\]
Apply Lemma~\ref{lem:temporal-block-tightness} with a deterministic
number of blocks of order $R$ large enough to cover the whole
interval. Its spatial equicontinuity permits replacement of
$N(Rx)$ by $\lfloor Rx/m_\alpha\rfloor$, uniformly in both
$x$ and $t$, in the centered reward sum divided by $\sqrt R$.
The error between this sum and the actual cumulative count is
at most $C\max_{i\le CR}(L_i+J_i)=O_{\PP}(\log R)$,
simultaneously for all $x,t$: only the block containing $Rx$
is partial. This proves the uniform line limit with covariance
\eqref{eq:temporal-cumulative-covariance}.

The bridge comparison in the proof of
Theorem~\ref{thm:temporal-spatial-CLT} uses the whole gap prefix
at cutoff $a$. Its two crossing blocks and final arc have
$O_{\PP}(\log R)+O(R^{1/4})=o_{\PP}(\sqrt R)$ points.
That bound is uniform in $x$ and $t\in[a,b]$, since points only
disappear, and all other blocks are protected by gaps longer
than $b$. The count tilt again has vanishing TV cost. Thus the
same uniform limit holds for the rooted circle.
For a uniform rotation $z$, extend $F$ by
$F(x+1,t)=F(x,t)+F(1,t)$ and use
$F(z+x,t)-F(z,t)$. This is a continuous map in the uniform
topology, including at the wrap point. The Gaussian covariance
\eqref{eq:temporal-cumulative-covariance} has stationary
spatial increments, so the rotated limit has the same law.
This proves the Poisson and entrance assertions.
The full-future TV comparison at $ra$ in
Lemma~\ref{lem:mesoscopic-depoisson-TV} transfers the uniform
limit to exactly $M$ input points. Its target is the Poisson
process with $\alpha=Mr$, so the centering remains
$Mq(Mrt)=\mu_\alpha(t)/r$ at every time.

Here is one explicit interpolation convention. Use the finite
spatial grid of points active at $ra$, and the finite temporal
grid of their collision separations divided by $r$, together
with the rectangle endpoints. Interpolate the cumulative counts
bilinearly on this grid and keep the continuous deterministic
centering unchanged. A spatial jump is at most one and a
temporal jump removes at most two points, so the error is at
most $3\sqrt r$. Ties have probability zero. This also covers
the entrance law, whose configuration at $ra$ is finite.

For a periodic smooth test $\phi$, the signed centered measure
associated with a cumulative field $H$ satisfies
\[
 \langle\nu_H,\phi\rangle
       =\phi(0)H(1)-\int_0^1\phi'(x)H(x)\,dx.
\]
The trace bound $|\phi(0)|\le C\|\phi\|_{H^1}$ and
Cauchy--Schwarz show
$\|\nu_H\|_{H^{-1}}\le C\|H\|_\infty$.
Thus the uniform cumulative convergence implies continuous
$H^{-1}$ path convergence for the interpolated fields; the
same bound controls their distance from the original fields.
Embedding into $H^{-u}$ proves the assertion for every $u\ge1$,
including convergence in the usual Skorokhod topology.

Finally, uniformly on $[a,b]$,
\[
 r\Xi_M(rt)(\T)=\mu_{Mr}(t)+\sqrt r F_M(1,t)
                         \longrightarrow c/t
                                      \quad\hbox{in probability}.
\]
Division by this strictly positive limit gives the normalized
cumulative process and its covariance. An arbitrary choice of
normalized measure on the even empty event has no effect,
since that event throughout $[a,b]$ has probability tending
to zero. Discrepancy contraction bounds the cumulative-field
difference, uniformly over $x,t$, by $\sqrt r$ times the input
edit count. It proves the joint path assertion for edited copies.
\end{proof}

\begin{proposition}[Earlier-point counts inside a later gap]
\label{prop:two-time-gap-transform}
In the line entrance history, fix $0<s\le t$. Under the Palm
law at cutoff $t$, let $Y_t$ be the following gap length and
let $K_{s,t}$ count the cutoff-$s$ points in its half-open
interval, including the left endpoint. It is odd. Put
\[
 h_s(x)=s^{-1}f(x/s),\quad
 F_s(p)=\int_s^\infty e^{-px}h_s(x)\,dx,
 \quad R_{{\rm o},s}(x)=s^{-1}R_{\rm o}(x/s).
\]
For $p>0$ and $0\le z\le1$,
\begin{align}
 \EE[z^{K_{s,t}}e^{-pY_t}]
  &=\tanh\left(\int_t^\infty e^{-pu}b_s(u,z)\,du\right),
       \label{eq:two-time-marked-gap-transform}\\
 b_s(u,z)&=\sum_{\substack{j\ge1\\j\ {\rm odd}}}
                   \frac{z^j}{j}h_s^{*j}(u).\nonumber
\end{align}
The sum is locally finite. In particular,
\begin{equation}\label{eq:two-time-gap-first-mark}
 \EE[K_{s,t}e^{-pY_t}]
   =(1-F_t(p)^2)\int_t^\infty e^{-pu}R_{{\rm o},s}(u)\,du.
\end{equation}
The marked later gaps are iid, and $K_{s,t}\le Y_t/s$.
\end{proposition}
\begin{proof}
The cutoff-$s$ descendants of a shielded cutoff-$t$ gap are
its finite terminal leaves in the backward triple-split tree.
Every split increases their number by two. Counting each
terminal gap's left endpoint gives exactly $K_{s,t}$; the
right endpoint belongs to the next later gap. Since terminal
gaps have length at least $s$, $K_{s,t}\le Y_t/s$.
The marked renewal construction makes different later gaps
independent, including these finite tree marks.

For $v\ge s$, write
$g_v(x,z)=h_v(x)\EE[z^{K_{s,v}}\mid Y_v=x]$.
At $v=s$, $g_s(x,z)=z h_s(x)$.
The nonnegative tree recursion, with one factor $z$ per terminal
leaf, gives for fixed $x>v$, almost everywhere in $v$,
\[
 \partial_v g_v(x,z)
   =g_v(v,z)\ind_{\{x\ge3v\}}
          \int_v^{x-2v}g_v(y,z)g_v(x-v-y,z)\,dy.
\]
This is the same root-split recursion as the scalar transform,
with the three terminal-count marks multiplied. There are
finitely many leaves on each bounded $x$ interval, so marking
and differentiation commute there. Also $0\le g_v\le h_v$ for
$0\le z\le1$. Its Laplace transform $G_v(p,z)$ consequently
satisfies
\[
 \partial_v G_v(p,z)
       =-e^{-pv}g_v(v,z)(1-G_v(p,z)^2),
 \qquad G_s(p,z)=zF_s(p).
\]
The moving lower endpoint supplies the negative term.
As $v\to\infty$, $G_v(p,z)\le e^{-pv}\to0$.
Integrating the arctanh equation therefore gives
\[
 \int_s^\infty e^{-pu}g_u(u,z)\,du
          =\operatorname{arctanh}(zF_s(p))
          =\sum_{j\ {\rm odd}}\frac{z^j}{j}F_s(p)^j.
\]
Uniqueness of Laplace transforms identifies
$g_u(u,z)=b_s(u,z)$ almost everywhere; the locally finite
convolution expression fixes a density version. Integration
from $t$ to infinity proves
\eqref{eq:two-time-marked-gap-transform}.
Finally differentiate monotonically as $z\uparrow1$.
The exponential gap moment and $K_{s,t}\le Y_t/s$ justify
passing the derivative through expectation, and
$\partial_z b_s(u,1)=R_{{\rm o},s}(u)$.
This proves \eqref{eq:two-time-gap-first-mark}.
\end{proof}

\begin{corollary}[An explicit temporal covariance and a non-Markov limit]
\label{cor:explicit-temporal-kernel}
For $z\ge1$ put
\[
 T(z)=\int_z^\infty\left(R_{\rm o}(y)-\frac c2\right)\,dy.
\]
The covariance in Theorem~\ref{thm:temporal-spatial-CLT} is
\begin{equation}\label{eq:explicit-temporal-kernel}
 V(s,t)=\frac{2c}{t}T(t/s),\qquad 0<s\le t.
\end{equation}
In particular, for $1\le t/s\le3$,
\begin{equation}\label{eq:short-ratio-temporal-kernel}
 V(s,t)=\frac c t\left[c\left(\frac ts+1\right)
                                      -\log\frac ts\right].
\end{equation}
For a universal $C<\infty$,
\begin{equation}\label{eq:factorial-temporal-decorrelation}
 |V(s,t)|\le\frac C{t\lfloor t/s\rfloor!},\qquad s\le t.
\end{equation}
The log-time normalized total-count limit
\[
 Z(\tau)=\frac{e^{\tau/2}}{\sqrt2c}G(e^\tau,1),
                    \qquad \tau\in\mathbb R,
\]
is a stationary continuous Gaussian process of variance one.
Its correlation is
\begin{equation}\label{eq:log-time-count-correlation}
 \mathcal R(h)=\frac{e^{-|h|/2}}cT(e^{|h|}).
\end{equation}
It tends to zero faster than every exponential in $|h|$ and
is not Markov in its own filtration. This statement concerns
Gaussian count fluctuations in separation time, not the
circle's coordinate-refresh kernel $Q$.
\end{corollary}
\begin{proof}
Use the iid marked later gaps from
Proposition~\ref{prop:two-time-gap-transform} as renewal blocks.
Their two rewards are $K=K_{s,t}$ and $1$, with rates $c/s$
and $c/t$. A partial boundary gap has exponential moments,
so the same renewal-reward covariance formula gives
\begin{equation}\label{eq:later-gap-temporal-covariance}
 V(s,t)=\frac{\EE[(K-(c/s)Y_t)(1-(c/t)Y_t)]}{\EE Y_t}.
\end{equation}
Let $m=\EE Y_t=t/c$, $m_2=\EE Y_t^2=t^2(c^{-2}+2c^{-1})$,
and $d=c/(2s)$. The factorial renewal remainder gives
\[
 \int_t^\infty e^{-pu}R_{{\rm o},s}(u)\,du
      =\frac d p+T(t/s)-dt+O(p).
\]
Also $1-F_t(p)^2=2mp-(m^2+m_2)p^2+O(p^3)$.
Equation~\eqref{eq:two-time-gap-first-mark} thus yields
\[
 \EE K=2md=t/s,\qquad
 \EE(KY_t)=d(m^2+m_2)+2mdt-2mT(t/s)
             =\frac c s m_2-2mT(t/s).
\]
Insert these identities into
\eqref{eq:later-gap-temporal-covariance}; its numerator is
$2T(t/s)$, proving \eqref{eq:explicit-temporal-kernel}.

Since $R_{\rm o}=U/2+a_0/4$, $\int a_0=1$, and
$\int_0^\infty(U-c)=c-1/2$, one has
$\int_0^\infty(R_{\rm o}-c/2)=c/2$.
The density vanishes below one, so $T(1)=c$.
For $1\le y\le3$, no convolution with three gaps contributes
and $R_{\rm o}(y)=f(y)=1/(2y)$. Therefore
\[
 T(z)=\frac c2(z+1)-\frac12\log z,\qquad 1\le z\le3,
\]
which gives \eqref{eq:short-ratio-temporal-kernel}.
Integrate \eqref{eq:parity-factorial} over successive unit
intervals to obtain $|T(z)|\le C/\lfloor z\rfloor!$.
This proves \eqref{eq:factorial-temporal-decorrelation}.

Scaling of $V$ makes $Z$ stationary, and its variance is one.
Its continuous version follows from the functional theorem on
every compact time interval. The covariance formula is
\eqref{eq:log-time-count-correlation}; the factorial bound
implies its stated decay. A nondegenerate stationary Gaussian
Markov process with continuous covariance must have
$\mathcal R(u+v)=\mathcal R(u)\mathcal R(v)$ for $u,v\ge0$:
Gaussian conditional expectation is
$\EE[Z(u+v)\mid Z(u)]=\mathcal R(v)Z(u)$, and the Markov
property allows conditioning first on the full past at $u$.
Continuity and $\mathcal R(0)=1$ force
$\mathcal R(h)=e^{-\lambda h}$ for some finite $\lambda\ge0$.
This contradicts the factorial decay. Equivalently, on
$0\le h\le\log3$ the exact formula is
$\mathcal R(h)=\cosh(h/2)-(h/(2c))e^{-h/2}$,
which is not an exponential.
\end{proof}

\subsection{Sharper empirical basin rates}
\label{sec:sharper-basin-rate}

\status{Working proof: the empirical discrepancy is optimized separately}

\begin{corollary}[Improved deletion and ranked-basin comparison]
\label{cor:sharper-basin-rate}
Use $s_0=\log(1/\sigma)$, $\ell_0=\log(3J)$, and
$\beta_{\rm map}=s_0/(\ell_0+s_0)$ as in
Theorem~\ref{thm:sharper-map-family-rates}. Set
\[
 \beta_{\rm emp}=\frac{s_0}{2(\log3+s_0)},\qquad
 \mathcal E_M=\left(\frac{\log M}{M}\right)^{\beta_{\rm map}}
                                        +M^{-\beta_{\rm emp}}.
\]
For sufficiently large even $M$,
\begin{equation}\label{eq:sharper-deletion-basin-rate}
 \EE\sum_{v\in B_M}|d_v(B_M)/M-\theta_v(B_M)|
                                      \le C\mathcal E_M.
\end{equation}
The ranked normalized deletion vector can be coupled to the
limiting basin partition with expected $\ell^1$ error at most
$C\mathcal E_M$. The ranked insertion vector has the sharper
bound $C(\log M/M)^{\beta_{\rm map}}$ already supplied by the
full-map coupling. Uniformly for $0<\varepsilon<1/2$,
\begin{equation}\label{eq:sharper-early-basin-rate}
 \EE\sum_{v\notin C_M(\varepsilon)}
       \bigl(\theta_v(B_M)+d_v(B_M)/M\bigr)
       \le C\bigl(\varepsilon^{\beta_{\rm map}}+\mathcal E_M\bigr).
\end{equation}
The constants and exponents are still uncalibrated.
\end{corollary}
\begin{proof}
Take the terminal cutoff $a=3\log M/M$ in the full last-pair
trees. The complex coefficient comparison and inverse-cutoff
bound from Theorem~\ref{thm:sharper-map-family-rates} give
\[
 \tau_M^{\rm last}(d)
       \le C\bigl(a(3J)^d+M^{-2}/\log M\bigr).
\]
The empirical argument of Corollary~\ref{cor:basin-rate} remains
valid on every shared geometry: its additional cost is
$C3^d\Delta_M$, with $\EE\Delta_M\le C M^{-1/2}$.
Thus the error to optimize is
\begin{equation}\label{eq:three-basin-errors}
 C\left(\sigma^d+a(3J)^d+M^{-1/2}3^d
                                      +M^{-2}/\log M\right).
\end{equation}
Choose
\[
 d=\left\lfloor\min\left\{
      \frac{\log(1/a)}{\ell_0+s_0},
      \frac{\log M}{2(\log3+s_0)}\right\}\right\rfloor.
\]
The first restriction bounds $a(3J)^d$ by
$a^{\beta_{\rm map}}$; the second bounds $M^{-1/2}3^d$ by
$M^{-\beta_{\rm emp}}$. Since the depth is the smaller of the
two, $\sigma^d$ is at most a constant times their sum.
The final contour error is smaller. This proves
\eqref{eq:sharper-deletion-basin-rate}.
The full-map coupling bounds the insertion pushforward error;
decreasing rearrangement contracts $\ell^1$ distance. Add
\eqref{eq:sharper-deletion-basin-rate} for the deletion vector.

For the early-mass assertion, further restrict the chosen depth
by $\log(1/\varepsilon)/(\ell_0+s_0)$.
The probability a revealed vertex dies before $\varepsilon$ is
at most $C\varepsilon(3J)^d\le C\varepsilon^{\beta_{\rm map}}$.
The unresolved origin mass is at most
$C(\varepsilon^{\beta_{\rm map}}+\mathcal E_M)$, and the other
terms in \eqref{eq:three-basin-errors} retain their bounds.
Outside the bad geometry and unresolved origins, the common
map only takes values active at $\varepsilon$. This proves the
insertion tail bound. The deletion tail follows by adding
\eqref{eq:sharper-deletion-basin-rate}.
\end{proof}

\subsection{Temporal paths in the full spatial Sobolev range}
\label{sec:full-sobolev-paths}

\status{Working proof: uniform Fourier maximal bounds}

\begin{lemma}[A maximal bound uniform in spatial frequency]
\label{lem:uniform-Fourier-path-bound}
Fix $0<a<b<\infty$. Start the iid line blocks of
Lemma~\ref{lem:temporal-renewal-blocks} at a block boundary zero,
with intensity parameter $\alpha\to\infty$ or $\alpha=\infty$.
Let $\Xi_\alpha(t)$ be their active counting measure. For $R\ge1$
and real $q$, set
\[
 W_{R,\alpha}(q,t)=R^{-1/2}\left(
     \sum_{x\in\Xi_\alpha(t)\cap[0,R)}e^{iqx}
            -\mu_\alpha(t)\int_0^R e^{iqx}\,dx\right).
\]
For a constant depending only on $a,b$,
\begin{equation}\label{eq:uniform-Fourier-path-bound}
 \sup_{R\ge1}\sup_{q\in\mathbb R}
      \EE\sup_{a\le t\le b}|W_{R,\alpha}(q,t)|^2\le C,
\end{equation}
uniformly for sufficiently large $\alpha$, including $\infty$.
\end{lemma}
\begin{proof}
Within one block let $\mathcal A_i(t)\subset[0,L_i)$ be its
active point core, and define
\[
 Z_i(t,q)=\sum_{x\in\mathcal A_i(t)}e^{iqx}
                -\mu_\alpha(t)\int_0^{L_i}e^{iqx}\,dx,
 \quad m_\alpha(t,q)=\EE Z_i(t,q),
 \quad D_i(t,q)=Z_i(t,q)-m_\alpha(t,q).
\]
The bound $|e^{iqx}|=1$ and the collision-interval estimate in
Lemma~\ref{lem:temporal-block-tightness} give, uniformly in $q$,
\[
 \EE|D_i(t,q)-D_i(s,q)|^k\le k!C^k|t-s|,
 \qquad \EE|D_i(a,q)|^k\le k!C^k.
\]
Write $S_j=\sum_{i\le j}L_i$. For fixed $q$,
\[
 M_j(t,q)=\sum_{i\le j}e^{iqS_{i-1}}D_i(t,q)
\]
is a martingale in the block index, for each $t$; its phase
is measurable before the $i$th block is exposed. Conditional
factorial moment bounds for the real and imaginary parts
therefore give the same maximal Bernstein inequality as
\eqref{eq:temporal-maximal-Bernstein}.
The dyadic chaining proof applies without change to these
conditional bounds. On its finest grid, the absolute increment
is bounded by the total number of removed points in all exposed
blocks plus $C|t-s|(\sum_iL_i+n)$, so the same unweighted
monotonicity estimate controls the remainder even though the
Fourier rewards themselves are not monotone. Integrating the
resulting tail estimates gives
\begin{equation}\label{eq:Fourier-block-martingale-max}
 \EE\max_{j\le n}\sup_{a\le t\le b}|M_j(t,q)|^2\le Cn,
\end{equation}
uniformly in $q,n,\alpha$.

We must also control the nonzero block mean. Put
$\varphi_\alpha(q)=\EE e^{iqL}$. Since
$m_\alpha(t,0)=0$, the inequality $|e^{iqx}-1|\le|q|x$
and the block exponential moments give
\[
 |m_\alpha(t,q)|\le C\min(|q|,1).
\]
There is a matching denominator bound
\begin{equation}\label{eq:block-characteristic-denominator}
 |1-\varphi_\alpha(q)|\ge c_0\min(|q|,1).
\end{equation}
For small $q$, expand its imaginary part using the uniform
third moment and $\EE L\ge a$. Away from zero, the limiting
block length has an absolutely continuous distribution, hence
$|\varphi_\infty(q)|<1$ for $q\ne0$ and its transform tends
to zero at infinity. Its law is the geometric mixture of
convolutions of conditioned short gaps and a conditioned long
gap. The gap densities converge in $L^1$; truncating the
uniform geometric tail proves convergence of these block-length
laws in TV. Their characteristic functions thus converge
uniformly on the entire real line. Compactness away from zero
and the vanishing limiting transform prove
\eqref{eq:block-characteristic-denominator} for all large
$\alpha$.
It follows that
$b_\alpha(t,q)=m_\alpha(t,q)/(1-\varphi_\alpha(q))$ is bounded
uniformly when $q\ne0$.

The predictable phase sum has the exact decomposition
\[
 (1-\varphi_\alpha(q))\sum_{i\le j}e^{iqS_{i-1}}
  =1-e^{iqS_j}+N_j(q),\qquad
 N_j(q)=\sum_{i\le j}e^{iqS_{i-1}}
                     (e^{iqL_i}-\varphi_\alpha(q)).
\]
The second term is a martingale and
$\EE\max_{j\le n}|N_j(q)|^2\le4n$.
Multiply by the bounded $b_\alpha(t,q)$ and combine with
\eqref{eq:Fourier-block-martingale-max}. This yields
\[
 \EE\max_{j\le n}\sup_t
       \left|\sum_{i\le j}e^{iqS_{i-1}}Z_i(t,q)\right|^2
                         \le C(n+1).
\]
At $q=0$ the mean is zero and the first martingale alone
suffices. There are at most $\lceil R/a\rceil+1$ blocks
meeting $[0,R]$. The final partial block contributes at most
$C\max_i(J_i+L_i)$, whose squared expectation is
$O(\log^2(2+R))\le CR$. Divide by $R$ to obtain
\eqref{eq:uniform-Fourier-path-bound}.
\end{proof}

\begin{corollary}[Temporal Gaussian paths in every $H^{-u}$, $u>1/2$]
\label{cor:full-sobolev-temporal-CLT}
All measure-valued path limits in
Theorem~\ref{thm:temporal-spatial-functional-CLT} hold in
$D([a,b],H^{-u}(\T))$ for every $u>1/2$, with continuous
limiting paths. This includes the full finite mesoscopic regime
$r_M\to0$, $Mr_M\to\infty$, with centering $Mq(Mr_Mt)$,
both parity entrance laws, the finite-Poisson laws, normalized
bridges, and the shared limits for any fixed family of
$o_{\PP}(r_M^{-1/2})$ arbitrary edits.
\end{corollary}
\begin{proof}
First work with the line block process on $[0,R)$, rescaled
to the unit circle. Its $k$th spatial Fourier coefficient is
$W_{R,\alpha}(2\pi k/R,t)$.
For the projection $P_K$ onto modes $|k|\le K$,
Lemma~\ref{lem:uniform-Fourier-path-bound} implies
\begin{equation}\label{eq:uniform-Sobolev-path-tail}
 \EE\sup_{t\in[a,b]}\|(I-P_K)G_R(t)\|_{H^{-u}}^2
       \le C\sum_{|k|>K}(1+4\pi^2k^2)^{-u}\longrightarrow0
                                      \quad(K\to\infty),
\end{equation}
uniformly in $R$ and large $\alpha$. The finite Fourier
projections converge as continuous paths by the cumulative
functional theorem and integration by parts. Equation
\eqref{eq:uniform-Sobolev-path-tail} therefore gives tightness
and identifies the Gaussian limit in the full space. It also
makes the limit's Fourier truncations Cauchy in
$C([a,b],H^{-u})$ in probability, so the limit has continuous
paths. Interpolating each collision in time changes the original
field by at most $2C_u/\sqrt R$, because point masses have
bounded $H^{-u}$ norm for $u>1/2$.

The circle-prefix comparison discards only
$O_{\PP}(\log R)+O(R^{1/4})$ points, uniformly in cutoff time.
Its error in the $H^{-u}$ path supremum is at most $C_u$ times
that count divided by $\sqrt R$, which tends to zero.
The remaining gap prefix agrees in TV with probability tending
to one; the count tilt also has vanishing TV cost. Rotation
acts continuously and isometrically on $H^{-u}$. These facts
transfer the line path limit to both circle entrance and
finite-Poisson laws. The exact-count whole-future TV comparison
then transfers it to the finite sample. No expectation is
transferred through a TV failure event here; the uniform Fourier
moment estimate is used on the line to prove weak path tightness.

The total-count normalization converges uniformly to $c/t$,
so the continuous mapping argument gives the bridge process
in the same space. Finally the $H^{-u}$ norm of a discrepancy
measure is at most $C_u$ times its number of points.
Discrepancy contraction proves the shared path limit for the
stated edited families.
\end{proof}

\subsection{Higher parent moments and a quadratic probe bound}
\label{sec:parent-higher-moments}

\status{Working proof: a common-background $L^2$ interaction estimate}

\begin{lemma}[All fixed parent-length moments]
\label{lem:parent-all-moments}
For a size-$M$ iid uniform input and an attainable active count
$m\ge3$, put $T_{M,m}^{(p)}=\sum_i x_i^p$ for the current gaps.
For every integer $p\ge1$ there is $C_p<\infty$ such that
\begin{equation}\label{eq:parent-all-moments}
 \EE T_{M,m}^{(p)}\le C_p m^{1-p},\qquad
 \EE L_{M,m}^p\le C_p m^{-p}.
\end{equation}
The second bound includes the last even pair by setting
$L_{M,2}=1$. Consequently, for every fixed $p\ge1$,
\begin{equation}\label{eq:parent-square-sum-Lp}
 \left\|\sum_{e:\,\text{pre-collision count}>m}L_e^2
                                      \right\|_{L^p}\le C_p/m.
\end{equation}
The constants do not depend on $M$ or its parity.
\end{lemma}
\begin{proof}
Write $T_p=\sum_i x_i^p$, $s=\min_i x_i\le1/m$, and let
$a,c$ be the uniformly ordered neighbours of the minimum.
For positive integers $i,j$,
\begin{equation}\label{eq:gap-power-product}
 T_iT_j\le T_{i+j-1}.
\end{equation}
Indeed, give $x_i$ probability mass $x_i$; two increasing powers
of this size-biased variable have nonnegative covariance.
Conditional on the gap bag,
\[
 \EE(a^i+c^i)\le\frac{2T_i}{m-1},\qquad
 \EE(a^ic^j)
   \le\frac{T_iT_j}{(m-1)(m-2)}
   \le\frac{T_{i+j-1}}{(m-1)(m-2)}.
\]
Induct on $p$, starting from $T_1=1$. The increment under a
merger is $(a+s+c)^p-a^p-s^p-c^p$. Expand it into its positive
multinomial terms. A term with one neighbour power $i$ and
$s^k$, $i+k=p$, $k\ge1$, has expected size at most
$C m^{-k-1}\EE T_i\le C_p m^{-p}$ by induction.
A term involving both neighbour powers $i,j\ge1$ and $s^k$,
$i+j+k=p$, has expectation at most
$C m^{-k-2}\EE T_{i+j-1}\le C_p m^{-p}$.
All needed moment indices are smaller than $p$.
The initial Dirichlet gap law gives
$\EE T_{M,M}^{(p)}=p!/[(M+1)\cdots(M+p-1)]$.
Summing the expected increments over counts of the same parity
proves the first inequality in \eqref{eq:parent-all-moments}.
Now $L^p\le3^{p-1}(a^p+s^p+c^p)$ and the conditional neighbour
formula prove the second. The finitely many last-count cases
are absorbed into $C_p$.
Finally, Minkowski's inequality and the $2p$th moment bound give
\[
 \left\|\sum_{j>m}L_{M,j}^2\right\|_p
       \le\sum_{j>m}(\EE L_{M,j}^{2p})^{1/p}
       \le C_p\sum_{j>m}j^{-2}\le C_p/m.
\]
For noninteger $p$, increase it to the next integer.
\end{proof}

\begin{lemma}[The fully shielded minimum-gap mark]
\label{lem:minimum-gap-parent-mark}
At an attainable count $m\ge3$, decorate each active gap $i$
with $B_i$, the sum of squared lengths of all earlier parent
intervals contained in it. Write $s$ for the minimum gap and
$B_{\min}$ for its mark. Then
\begin{equation}\label{eq:minimum-gap-parent-mark}
 \EE[B_{\min}\mid\text{current geometry and original gap counts}]
                                     \le Cs^2,
 \qquad \EE B_{\min}\le C/m^2.
\end{equation}
Moreover, for $B=\sum_iB_i$ and $D=\sum_i x_i B_i$,
\begin{equation}\label{eq:parent-decoration-moments}
 \EE B\le C/m,\qquad \EE D\le C/m^2.
\end{equation}
\end{lemma}
\begin{proof}
Use the configuration immediately before the next collision;
its separation is the current minimum $s$. Given its endpoint
geometry and the numbers of original points inside its gaps,
the gap interiors are independent uniform samples subject to
shielding through cutoff $s$. This conditional description is
valid at a fixed active count as well: a candidate $m$-point
configuration whose interiors are all cleared through the
minimum of its gaps is exactly the unique state immediately
before the transition from $m$ to $m-2$. Thus the usual interval
factorization applies with cutoff equal to that minimum, not
to the preceding random collision time.

For the minimum gap the cutoff equals its whole length $s$.
If it contains $2n$ original interior points, identify its two
protected endpoints to make a circle of circumference $s$.
Its fully shielded history is the history of these $2n$ uniform
points and the distinguished point zero, conditioned on zero
being the final survivor. As long as zero survives, the two
arcs incident to it give exactly the two protected interval
boundaries, so the histories and parent lengths coincide.
A rotation-invariant statistic of the unlabelled circle history
has the same law under this conditioning: start with $2n+1$
iid points, distinguish one uniform label, and rotate it to
zero. Given the unlabelled geometry, that label is the survivor
with probability $1/(2n+1)$. Apply this observation to the sum
of squared parent lengths. Lemma~\ref{lem:parent-all-moments},
including the last three-point merger, bounds its expectation
by $C$ on the unit circle, uniformly in $n$.
Scaling gives $\EE B_{\min}\le Cs^2$ conditionally; for $n=0$
the mark is zero. Since $s\le1/m$, this proves
\eqref{eq:minimum-gap-parent-mark}.

Every earlier parent belongs to exactly one current gap, so
$B=\sum_{j>m}L_{M,j}^2$; the first bound in
\eqref{eq:parent-decoration-moments} follows from
\eqref{eq:parent-square-sum-Lp}. The gap-order lemma holds for
these decorations too, since their merge rule is deterministic.
At a merger of gaps $a,s,c$, with marks $B_a,B_s,B_c$ and
$L=a+s+c$, the new mark is $B_a+B_s+B_c+L^2$.
Consequently
\[
 \Delta D=(s+c)B_a+(a+c)B_s+(a+s)B_c+L^3.
\]
Conditional neighbour sampling bounds its expectation by
\[
 C\left(\frac{B}{m^2}+\frac{B_{\min}}m
                                  +\EE[L^3\mid\text{decorated bag}]\right).
\]
For example
$\EE(cB_a+aB_c)\le2B/[(m-1)(m-2)]$ and
$\EE[(a+c)B_s]\le2B_{\min}/(m-1)$; the remaining terms use
$s\le1/m$. Taking expectations gives $\EE\Delta D\le C/m^3$.
Initially $D=0$. Sum over earlier counts to obtain
$\EE D\le C/m^2$.
\end{proof}

\begin{proposition}[Quadratic occupancy of common-background parents]
\label{prop:common-parent-probe-L2}
Fix an attainable $m\ge3$ and let $\mathcal P_m$ be the parent
intervals of collisions whose pre-collision active count exceeds
$m$. On $\T^2$ define
\[
 h_m(z,z')=\sum_{I\in\mathcal P_m}\ind_I(z)\ind_I(z'),
 \qquad h_{m,1}(z)=\int_\T h_m(z,z')\,dz'.
\]
Then, uniformly in the initial size,
\begin{equation}\label{eq:common-parent-kernel-L2}
 \EE\|h_m\|_{L^2(\T^2)}^2\le C/m,
 \qquad \EE\|h_{m,1}\|_{L^2(\T)}^2\le C/m^2.
\end{equation}
For $q$ independent uniform probes, let $K_I$ count those in
$I$ and put $R_{q,m}=\sum_{I\in\mathcal P_m}K_I(K_I-1)$.
Writing $\mathcal B$ for the full background,
\begin{align}
 \EE\operatorname{Var}(R_{q,m}\mid\mathcal B)
            &\le C\left(\frac{q^3}{m^2}+\frac{q^2}{m}\right),
                 \label{eq:parent-probe-conditional-variance}\\
 \EE R_{q,m}^2
            &\le C\left(\frac{q^4+q^3}{m^2}
                                      +\frac{q^2}{m}\right).
                 \label{eq:parent-probe-second-moment}
\end{align}
In particular $\|R_{q,m}\|_2\le Cq$ when $m\ge q\ge2$.
This is a common-background occupancy bound. It does not by
itself control the matching-history telescoping remainder,
whose rectangle witnesses live in different positive corners.
\end{proposition}
\begin{proof}
Parent intervals are laminar, up to their endpoints. When a new
parent $I$ of length $L$ is formed, the old parents it contains
are precisely those in its three daughter gaps. Put
$B_*=B_a+B_s+B_c$. Adding $I$ to the kernel therefore increases
$\|h\|_2^2$ by $L^2+2B_*$, and increases $\|h_1\|_2^2$ by
$L^3+2LB_*$.
At pre-collision count $j$, conditional neighbour exchangeability
gives
\[
 \EE[B_*\mid\text{decorated bag}]
       =B_{\min}+\frac{2(B-B_{\min})}{j-1}.
\]
Its expectation is $O(j^{-2})$ by
Lemma~\ref{lem:minimum-gap-parent-mark}.
To control the other increment, expand
$(a+s+c)(B_a+B_s+B_c)$. The diagonal neighbour terms have mean
$2(D-sB_{\min})/(j-1)$; the cross-neighbour terms are at most
$2B/[(j-1)(j-2)]$. Terms involving $s$ or $B_{\min}$ are at
most $C(B/j^2+B_{\min}/j)$. Hence
\[
 \EE(LB_*)\le C\left(\frac{\EE D}{j}
              +\frac{\EE B}{j^2}+\frac{\EE B_{\min}}j\right)
                         \le C/j^3.
\]
The parent moment lemma bounds $\EE L^2$ and $\EE L^3$ by
$C/j^2$ and $C/j^3$. Summing the two kernel increments over
$j>m$ proves \eqref{eq:common-parent-kernel-L2}.

Finally
$R_{q,m}=\sum_{i\ne j}h_m(Z_i,Z_j)$ and, conditionally on
$\mathcal B$,
$\EE[R_{q,m}\mid\mathcal B]=q(q-1)\sum_I|I|^2$.
In the conditional variance expansion, terms with four distinct
probe indices cancel. Those sharing one index are bounded by
$Cq^3\|h_{m,1}\|_2^2$, and those sharing both by
$Cq^2\|h_m\|_2^2$. This gives
\eqref{eq:parent-probe-conditional-variance}.
Add the squared conditional mean and use
\eqref{eq:parent-square-sum-Lp} with $p=2$ to prove
\eqref{eq:parent-probe-second-moment}.
\end{proof}

\subsection{Fluctuations of the weighted early history}
\label{sec:parent-square-fluctuations}

\status{Working proof: an additive tree mark with a spatial Gaussian limit}

\begin{lemma}[Uniform moments of the parent-square tree mark]
\label{lem:parent-square-tree-mark}
Inside a shielded interval of length $x$ and cutoff $t\le x$,
let $\mathcal B$ be the sum of squared parent lengths over its
entire internal history. For both finite-Poisson shielding and
the entrance tree, uniformly in their parameters,
\begin{equation}\label{eq:parent-square-tree-moments}
 \|\mathcal B\|_{L^p}\le Cp x^2,\qquad p\ge1.
\end{equation}
If $\mathcal B^{(d)}$ retains only the first $d$ backward tree
generations, then
\[
 \|\mathcal B-\mathcal B^{(d)}\|_{L^p}
             \le C_p x^2\rho_{\rm h}^{d/p},
\]
where $\rho_{\rm h}<1$ is the uniform contraction in
Lemma~\ref{lem:insertion-height}. In particular the entrance
sum is finite almost surely, and
$\EE\exp(c_1\mathcal B/x^2)\le C_1$ for universal positive
$c_1,C_1$.
\end{lemma}
\begin{proof}
Choose an independent uniform location $z$ in the interval.
Let $L_j(z)$ be the length of its ancestral interval at backward
depth $j$, set to zero after the ancestry stops. The uniform
size-biased daughter contraction in the proof of
Lemma~\ref{lem:insertion-height} gives
$\EE L_j\le x\rho_{\rm h}^j$. Since $L_j\le x$,
$\|L_j\|_p\le x\rho_{\rm h}^{j/p}$.
For every tree,
\[
 \mathcal B=\int_0^x\sum_{j\ge0}L_j(z)
                 \ind_{\{\text{the depth-$j$ interval splits}\}}\,dz.
\]
Indeed a parent of length $\ell$ contributes $\ell$ on an
origin interval of length $\ell$. Jensen and Minkowski give
\[
 \|\mathcal B\|_p
   \le x\left\|\sum_{j\ge0}L_j(z)\right\|_p
   \le\frac{x^2}{1-\rho_{\rm h}^{1/p}}\le Cp x^2.
\]
Starting this sum at depth $d$ proves the truncation estimate.
The exponential-moment assertion follows by expanding its
positive series and using $p^p/p!\le e^p$.
The entrance estimate follows directly from the same split
bound, or by finite-depth convergence and monotone convergence.
\end{proof}

\begin{theorem}[A marked renewal central limit theorem for parent squares]
\label{thm:entrance-parent-square-CLT}
In either circle entrance law define
\[
 B(r)=\sum_{e:\,\text{separation}(e)\le r}L_e^2,
 \qquad \kappa=\frac c2+4+\frac2c,
 \qquad c=(2e^\gamma)^{-1}.
\]
For an odd final three-point collision its parent is the whole
circle; for an even final pair use the binary parent of length
one. The infinite sum of earlier terms is finite almost surely.
Let $(Y,\mathcal B)$ be a Palm gap and its full parent-square
mark at cutoff one in the line entrance law. Put
\[
 Z_0=1-cY,\qquad Z_1=\mathcal B-\kappa Y.
\]
Then, as $r\downarrow0$,
\begin{equation}\label{eq:parent-square-joint-CLT}
 \left(\sqrt r(K(r)-c/r),\,
              \frac{B(r)-\kappa r}{r^{3/2}}\right)
       \Longrightarrow\mathcal N\left(0,\,
                     c\,\EE[(Z_iZ_j)_{i,j=0,1}]\right).
\end{equation}
The centered means and covariance matrix converge as well.
The covariance matrix is positive definite.
Section~\ref{sec:weighted-history-integral} reduces its entries
to a one-dimensional integral and certifies the corresponding
active-count variance.
Moreover, for $0<\delta\le1/2$ and all sufficiently small $r$,
\begin{equation}\label{eq:parent-square-short-band-variance}
 \operatorname{Var}\{B(r)-B((1-\delta)r)\}\le C\delta r^3,
\end{equation}
and its mean is $\kappa\delta r+O(\delta r^{-C}e^{-c_2/r})$,
for suitable absolute constants $C,c_2>0$.
\end{theorem}
\begin{proof}
At cutoff one the line gaps are iid with their independent
conditional tree marks. Lemma~\ref{lem:parent-square-tree-mark}
and the exponential tail of $Y$ give every fixed moment of
$(Y,\mathcal B)$.
The mean mark rate can be evaluated by collision flux. At
separation $u$, the point intensity is $c/u$ and the gap density
at its lower endpoint is $1/(2u)$. Thus collisions have rate
$c/(2u^2)$ per unit spatial length and unit separation.
Their two neighbouring gaps are independent with law $uY$.
The parent length is $u(1+Y_1+Y_2)$, so Tonelli gives
\[
 \frac{\EE\mathcal B}{\EE Y}
   =\int_0^1\frac{c}{2u^2}
                       u^2\EE(1+Y_1+Y_2)^2\,du
   =\frac c2+4+\frac2c=\kappa,
\]
using $\EE Y=1/c$ and $\EE Y^2=c^{-2}+2c^{-1}$.
This proves that both $Z_0,Z_1$ are centered.

The covariance matrix is positive definite. The entrance
split kernel has total mass one, so the root itself is an
accepted parent and $\mathcal B\ge Y^2$ almost surely.
Also $Y$ has unbounded support, by \eqref{eq:two-sided-f}.
If $aZ_0+bZ_1=0$ almost surely and $b\ne0$, then
$\mathcal B=(\kappa+ac/b)Y-a/b$, contradicting that quadratic
lower bound for sufficiently large $Y$. If $b=0$, the
nonconstant gap law forces $a=0$.

Rescale the circle to length $R=1/r$. Its parent-square mark
sum is $r^{-2}B(r)$. The count-biased rooted renewal bridge
has independent marks conditional on its gaps. Attach those
same marks when coupling its stopped prefix to iid renewal
increments, as in \eqref{eq:palm-renewal-bridge}.
The iid vector renewal-reward CLT then gives covariance per
unit length $\EE(Z_iZ_j)/\EE Y=c\EE(Z_iZ_j)$.
For detail on the discarded marks, stop at $R-R^{1/4}$.
The longest gap is $O_{\PP}(\log R)$, and the sum of squared
gap lengths in the discarded arc is at most its length times
the longest gap. Conditional on the geometry, the $L^p$ norm
of its total mark is at most $C_p\sum y_i^2$ by the preceding
lemma. It is therefore $o_{\PP}(\sqrt R)$. A root-crossing
gap satisfies the same bound. The count tilt has vanishing TV
cost, which proves the joint weak limit in
\eqref{eq:parent-square-joint-CLT}.

We record the renewal calculation needed for moments and for
a shrinking band. Let $W$ be any real gap mark whose conditional
moments, multiplied by $f(y)$, have exponentially decreasing
upper envelopes. Define
$g(y)=f(y)\EE[W\mid Y=y]$ and
$h(y)=f(y)\EE[W^2\mid Y=y]$.
For parity $p\in\{0,1\}$, put
\[
 A_p=\sum_{k\equiv p\ (2),\ k\ge1}f^{*(k-1)},\qquad
 C_p=\sum_{k\equiv p\ (2),\ k\ge2}(k-1)f^{*(k-2)},
\]
where $f^{*0}=\delta_0$.
Differentiating the cyclic sector density
$2R f^{*k}(R)/k$ after marking its factors gives, on nonempty
sectors,
\begin{equation}\label{eq:marked-cyclic-moments}
 \EE\sum_iW_i=2R(g*A_p)(R),\qquad
 \EE\left(\sum_iW_i\right)^2
       =2R\bigl((h*A_p)(R)+(g*g*C_p)(R)\bigr).
\end{equation}
For large arguments, the factorial renewal estimates imply
\[
 A_p(y)=c/2+O(e^{-c_3y}),\qquad
 C_p(y)=c^2y/2+d_0+O(e^{-c_3y}),
\]
with the same finite constant $d_0$ for both parities.
For example, the sum of the two $C_p$ kernels is
$(\delta_0+U)^{*2}$, and their difference is
$(\delta_0-D)^{*2}$; write $U=c+e$ and use the integrable
exponential remainder. Atomic terms at zero are kept in the
convolutions and have exponentially small contributions at $R$.

Apply \eqref{eq:marked-cyclic-moments} to a linear combination
of $Z_0,Z_1$. Here $\int g=0$. Both the total mass and first
moment of $g*g$ vanish, so the affine part of $C_p$ contributes
only an exponentially small tail error. The leading second
moment is $cR\EE W^2$, and the mean is exponentially small.
This proves the covariance and centered-mean convergence.
The even empty sector is negligible: its probability is
exponentially small in $R$, and the full-circle parent-square
sum has all fixed moments by applying the tree-mark bound to
the last-pair roots. H\"older's inequality handles that sector.

For the short band, attach to a cutoff-one gap the mark
$W_\delta$, the sum of squared parents whose separation lies
in $[1-\delta,1]$. Conditional on $Y=y$, at most $2y$ terminal
gaps can occur at cutoff $1-\delta$, so $W_\delta\le Cy^3$.
If this mark is nonzero, the root's last split separation lies
in that band. The split kernel and the exponential bounds on
$f$ show
\[
 \PP(W_\delta>0\mid Y=y)\le C\delta y.
\]
Indeed, for $1/2\le u\le1$, the numerator density
$w_u(u)\int w_u(a)w_u(y-u-a)\,da$, divided by $w_1(y)$,
is at most $Cy$; its exponential ratio is bounded because
$e^{-\kappa_f y/u}\le e^{-\kappa_f y}$, where $\kappa_f$
is the gap-envelope exponent. Thus, for every fixed $j$,
\[
 \EE(W_\delta^j\mid Y=y)\le C_j\delta y^{3j+1}.
\]
Its mean reward per unit length is $\kappa\delta$, by the
same flux integral. Center the mark as
$W_\delta-\kappa\delta Y$.
Its conditional first-moment envelope is $O(\delta)$ times a
polynomial, its conditional second-moment envelope is also
$O(\delta)$ times a polynomial, and its unconditional mean is
zero. Equation~\eqref{eq:marked-cyclic-moments} now gives
variance at most $C\delta R$ and mean error
$O(\delta R^C e^{-c_2R})$ on the nonempty sectors.
For an empty even state, a nonzero band mark requires the
last-pair separation to lie in that band. Its separation density
and conditional full-root moments give the same exponential
bound with a factor $\delta$; otherwise the band mark is zero.
Scaling the mark by $r^2$ proves
\eqref{eq:parent-square-short-band-variance} and its mean claim.
\end{proof}

\begin{corollary}[Parent-square fluctuations at a fixed active count]
\label{cor:parent-square-count-CLT}
Let $B_{\infty,m}$ be the sum of squared parent lengths through
the collision that first leaves the parity-compatible count
$m$ active in an entrance history. Let $(G_0,G_1)$ have the Gaussian law in
\eqref{eq:parent-square-joint-CLT}. Then
\begin{equation}\label{eq:parent-square-count-CLT}
 m^{3/2}\left(B_{\infty,m}-\frac{\kappa c}{m}\right)
       \Longrightarrow c^{3/2}
                         \left(G_1+\frac\kappa cG_0\right).
\end{equation}
For a size-$M$ iid sample of the same parity, write $B_{M,m}$
for the corresponding sum. For every fixed $p\ge1$, whenever
$m=o(M/\log M)$ there is a coupling with
\begin{equation}\label{eq:parent-square-count-coupling}
 \|B_{M,m}-B_{\infty,m}\|_p\le C_p\frac{\log M}{M}.
\end{equation}
Consequently \eqref{eq:parent-square-count-CLT} also holds for
$B_{M,m}$ if $m\to\infty$ and
$m^{3/2}\log M/M\to0$. The Gaussian variance may be zero;
variance asymptotics at the random collision time require the
additional moment argument in Corollary~\ref{cor:parent-count-moments}.
\end{corollary}
\begin{proof}
Put $r=c/m$ and let $\tau_m$ be the first separation at which
$K$ reaches $m$. Monotonicity, the count law, and the temporal
count functional limit give
\begin{equation}\label{eq:active-count-cutoff-inversion}
 \frac{\tau_m-r}{r^{3/2}}\Longrightarrow G_0/c,
 \qquad \tau_m/r\longrightarrow1.
\end{equation}
Indeed, $K(\tau_m)=m$ and the field is asymptotically
continuous on any compact ratio interval containing one;
expand $c/\tau_m-c/r$. The parity restriction ensures the
count is reached exactly.

For any fixed $A$, the short-band variance bound gives,
at each point of a finite grid with mesh
$\varepsilon r^{3/2}$ in
$[r-Ar^{3/2},r+Ar^{3/2}]$,
\[
 B(t)-B(r)-\kappa(t-r)=o_{\PP}(r^{3/2}).
\]
The variance is $O(r^{7/2})$ for such a band, whereas the
squared normalization is $r^3$; its mean error is exponentially
small. Monotonicity of $B$ fills each grid interval with an
additional error at most $\kappa\varepsilon r^{3/2}$ plus
the endpoint errors. Let $\varepsilon\downarrow0$ and then
$A\uparrow\infty$, using \eqref{eq:active-count-cutoff-inversion}.
This proves
\[
 B(\tau_m)=B(r)+\kappa(\tau_m-r)+o_{\PP}(r^{3/2}),
\]
and hence \eqref{eq:parent-square-count-CLT}.

We prove the finite coupling without reconstructing a past
from an active-configuration TV comparison. Fix $p$, choose a
large constant $A_p$, and set $a=A_p\log M/M$.
Theorem~\ref{thm:near-microscopic-TV} couples the finite and
entrance configurations at $a$, and their complete futures,
except on probability $O(M^{1-A_p}/\log M)$.
The entrance count at $a$ is at least $c/(2a)$ except on
probability $Ce^{-c_4/a}$. For this lower-tail bound use the
cyclic sector densities. For $k\le cR/2$, $R=1/a$, the gap
exponential moment and a small positive tilt $\theta$ give
\[
 f^{*k}(R)\le C\exp\{-\theta R
                       +k\log\EE e^{\theta Y}\}
                         \le Ce^{-c_4R}.
\]
Sum over those sectors; their polynomial factors are absorbed
in the exponential. The same lower-count event holds for the
finite sample on the coupling event.

Choose an attainable $\ell\asymp1/a$ smaller than $c/(2a)$.
On $K_M(a)\ge\ell$, the finite pre-$a$ parent-square sum is
at most $B_{M,\ell}$, whose $L^p$ norm is $O(a)$ by
Lemma~\ref{lem:parent-all-moments}. On the complement use its
uniform global $L^{2p}$ bound and H\"older's inequality.
Choosing $A_p>2p+2$ makes this exceptional contribution $O(a)$.
Thus the finite pre-$a$ sum has $L^p$ norm $O(a)$.
For the entrance law, conditional tree moments give the same
bound: for gap lengths $x_i$ summing to one,
\[
 \left(\sum_i x_i^2\right)^p\le\sum_i x_i^{p+1},
 \qquad \EE\sum_i x_i^{p+1}\le C_p a^p.
\]
The second estimate follows from the one-gap intensity and its
exponential envelope. The empty state is exponentially unlikely
and has bounded full-tree moments.

Since $m=o(1/a)$, the two coupled histories both reach $m$
after cutoff $a$ on the good event. Every parent length accepted
between $a$ and that common transition is then identical.
Their difference at count $m$ is bounded by their two pre-$a$
sums. On the exceptional event use the uniform global
$L^{2p}$ bounds again. This proves
\eqref{eq:parent-square-count-coupling}. Multiplication by
$m^{3/2}$ transfers the weak limit under the stated scale
condition; it does not alone prove moment convergence of the
entrance statistic at its random transition time.
\end{proof}

\begin{lemma}[Higher centered moments of a shrinking history band]
\label{lem:parent-band-higher-moments}
For every integer $p\ge1$, uniformly for $0<\delta\le1/2$ and
small $r$ in either entrance parity,
\begin{equation}\label{eq:parent-band-higher-moments}
 \EE\left|\frac{B(r)-B((1-\delta)r)-\kappa\delta r}
                         {r^{3/2}}\right|^{2p}
       \le C_p\bigl(\delta^p+\delta r^{p-1}\bigr).
\end{equation}
The analogous full-cutoff centered parent-square statistic and
count field have uniformly bounded moments of every fixed order.
\end{lemma}
\begin{proof}
Here is a higher-order version of
\eqref{eq:marked-cyclic-moments}. Suppose a centered gap mark
$W$ has conditional moment densities
$g_j(y)=f(y)\EE[W^j\mid Y=y]$ whose exponentially weighted
absolute norms are at most $C_j\delta$, with $\int g_1=0$.
Differentiate the marked cyclic sector density $2R F^k/k$
$n$ times. Each set partition of $\{1,\ldots,n\}$ into
$\ell$ blocks gives the convolution of the corresponding
$g_j$'s with
\[
 D_{\ell,p_0}
   =\sum_{\substack{k\ge\ell\\k\equiv p_0\ (2)}}
          (k-1)_{\ell-1}f^{*(k-\ell)},
\]
multiplied by $2R$; here $p_0$ denotes parity and the subscript
on the falling factorial is its order. The sum and difference
of its two parity versions are respectively
\[
 (\ell-1)!(\delta_0+U)^{*\ell},\qquad
 (-1)^\ell(\ell-1)!(\delta_0-D)^{*\ell}.
\]
Here the difference means even minus odd. Consequently, on positive arguments,
$D_{\ell,p_0}$ is a polynomial of degree at most $\ell-1$
plus an exponentially decreasing remainder and atomic terms
at zero. This follows by writing $U=c+e$ and convolving; all
moments of $e$ and $D$ are exponentially integrable.

If the partition has $a$ singleton blocks, the convolution
of its moment densities has vanishing moments through order
$a-1$, because each singleton contributes the zero-mass
factor $g_1$. Its convolution with that polynomial therefore
has degree at most $\ell-1-a$. Put $b=\ell-a$, the number of
nonsingleton blocks. For $n=2p$ one has $b\le p$.
After multiplication by $2R$, the contribution is bounded by
$C_p\delta^\ell R^b$; if $b=0$ the polynomial contribution
vanishes entirely. Since $\delta\le1$, the sum of these bounds
is at most
$C_p[(\delta R)^p+\delta R]$.
The remainder terms satisfy the same bound, since their
absolute convolution norms are $O(\delta^\ell)$ and their
tails are exponential up to polynomial factors.
Thus
\begin{equation}\label{eq:marked-cyclic-higher-moments}
 \EE\left|\sum_iW_i\right|^{2p}
                  \le C_p[(\delta R)^p+\delta R].
\end{equation}

Use the centered band mark
$W=W_\delta-\kappa\delta Y$ from the preceding theorem.
Its moment densities satisfy these bounds, uniformly in
$\delta$, by the root-split estimate. The empty even sector
also has this bound: the nonzero band contribution requires
the last pair in the band and costs $O(\delta R^C e^{-cR})$;
the deterministic centering on the remaining empty states is
smaller. With $R=1/r$, scaling
\eqref{eq:marked-cyclic-higher-moments} proves
\eqref{eq:parent-band-higher-moments}.
For the full parent-square mark and for $1-cY$, use the same
argument with $\delta=1$ and their polynomial conditional
moment bounds. This gives the stated uniform centered moments.
\end{proof}

\begin{corollary}[Moment convergence at an active count and probe variance]
\label{cor:parent-count-moments}
The limit \eqref{eq:parent-square-count-CLT} holds with
convergence of every fixed centered moment. In particular,
with
\[
 v_{\rm B}=c^3\operatorname{Var}
                    \left(G_1+\frac\kappa cG_0\right)>0,
\]
one has
\begin{equation}\label{eq:parent-count-variance}
 m^3\operatorname{Var}B_{\infty,m}\longrightarrow v_{\rm B},
 \qquad \EE B_{\infty,m}=\kappa c/m+o(m^{-3/2}).
\end{equation}
These assertions transfer to $B_{M,m}$ if
$m^{3/2}\log M/M\to0$. More generally, if $m=o(M/\log M)$,
\begin{equation}\label{eq:finite-parent-count-variance-bound}
 \operatorname{Var}B_{M,m}
       \le C\left(m^{-3}+(\log M/M)^2\right).
\end{equation}
For the common-background probe statistic in
Proposition~\ref{prop:common-parent-probe-L2}, in this same range,
\begin{equation}\label{eq:unconditional-parent-probe-variance}
 \operatorname{Var}R_{q,m}
  \le C\left(\frac{q^2}m+\frac{q^3}{m^2}
            +\frac{q^4}{m^3}+q^4(\log M/M)^2\right).
\end{equation}
Thus $\operatorname{Var}R_{q,m}=O(q)$ if $m\asymp q\to\infty$
and $q^{3/2}\log M/M=O(1)$. This remains a statement about
common-background occupancy, not the full history remainder.
\end{corollary}
\begin{proof}
We verify uniform integrability at the random collision time.
Put $r=c/m$ and $U_r=(\tau_m-r)/r^{3/2}$.
The higher centered count moments, applied at the deterministic
cutoffs $r(1\pm u\sqrt r)$, and monotonicity give, for every
integer $j$,
\begin{equation}\label{eq:count-inversion-polynomial-tail}
 \PP(|U_r|>u)\le C_j u^{-2j},
               \qquad 1\le u\le(2\sqrt r)^{-1}.
\end{equation}
The centering difference is comparable to $u/\sqrt r$ on
this range. In particular the probability that
$\tau_m/r\notin[1/2,3/2]$ is $O_j(r^j)$.
Also $\tau_m\le1/m=r/c$ deterministically: the collision
entering count $m$ occurs from a count at least $m+2$, whose
minimum gap is at most $1/(m+2)$.

On $|U_r|\le A$, with $1\le A\le(2\sqrt r)^{-1}$,
monotonicity bounds $|B(\tau_m)-B(r)|$ by the two band
increments between $r$ and $r(1\pm A\sqrt r)$.
Lemma~\ref{lem:parent-band-higher-moments}, including the
deterministic drift, bounds their normalized $L^{2p}$ norms
by $C_p(A+1)$, for every fixed $p$.
Split according to $2^{k-1}<|U_r|\le2^k$ up to that range.
H\"older and \eqref{eq:count-inversion-polynomial-tail} bound
the contribution to the $p$th moment by
$C_p2^{kp}2^{-kj}$; choose $j>p$ and sum.
On the remaining event use $B(\tau_m)\le B(r/c)$.
The full-cutoff moment bound gives
$\|B(r/c)\|_{2p}\le C_pr$, while its exceptional probability
is $O_j(r^j)$. Choosing $j>p$ again makes the normalized
contribution bounded, and in fact negligible.
Together with the centered full-cutoff moments this proves
uniform boundedness of every fixed moment of
$(B(\tau_m)-\kappa r)/r^{3/2}$.

The weak limit in Corollary~\ref{cor:parent-square-count-CLT},
with this bound at a larger moment order, gives convergence of
all centered moments. This proves
\eqref{eq:parent-count-variance}.
The $L^p$ coupling \eqref{eq:parent-square-count-coupling}
transfers the moments under the stated scale restriction.
For the more general variance bound, use
$\operatorname{Var}(X+Y)\le2\operatorname{Var}X+2\EE Y^2$
and that coupling with $p=2$.
Finally total variance and
$\EE[R_{q,m}\mid\mathcal B]=q(q-1)B_{M,m}$ give
\eqref{eq:unconditional-parent-probe-variance} from
\eqref{eq:parent-probe-conditional-variance} and
\eqref{eq:finite-parent-count-variance-bound}.
Finally the covariance formula in
Theorem~\ref{thm:entrance-parent-square-CLT} gives
\[
 v_{\rm B}=c^4\EE(\mathcal B-2\kappa Y+\kappa/c)^2
 \ge c^4\int_{2\kappa}^\infty
          (y^2-2\kappa y+\kappa/c)^2 f(y)\,dy>0.
\]
The last inequality uses $\mathcal B\ge Y^2$ on this positive
tail. This is a strict analytic lower bound, not a calibrated
numerical enclosure.
\end{proof}

\begin{proposition}[Common-parent occupancy does not dominate history interaction]
\label{prop:common-parent-obstruction}
There are generic circle configurations $B_n$, two probes
$z_n,w_n$, and a cutoff at which the background has eight active
points, such that no accepted background parent contains both
probes, but the restricted matching-vector mixed difference has
$\ell^1$ norm $2n+1$. Moreover,
\[
 H_{\le\tau}(B_n,B_n+z_n)
   =H_{\le\tau}(B_n,B_n+w_n)=1,
 \qquad H_{\le\tau}(B_n,B_n+z_n+w_n)=2n+1.
\]
Thus neither a constant multiple of common-parent occupancy nor
that bound plus a constant times the probe or cutoff active
count can deterministically dominate this interaction.
\end{proposition}
\begin{proof}
Use a circle of circumference $L=128(n+3)^2$ and divide all
coordinates and separations by $L$ at the end. Put
\[
 x_j=\frac{j(j+5)}2+2,\quad 1\le j\le2n,
 \qquad B_n=\{0,x_1,\ldots,x_{2n},L/8,\ldots,7L/8\},
 \quad z_n=2,\quad w_n=L-1.
\]
Stop at separation $\tau=2n+5/2$ in these unscaled coordinates.
The background removes $(x_{2j-1},x_{2j})$ at separation
$2j+2$, for $1\le j\le n$, and nothing else through $\tau$.
Every corresponding parent starts at zero and ends at
$x_{2j+1}$, or at $L/8$ for the last pair. It contains $z_n$
but not $w_n$; the common-parent pair occupancy is zero.

The single probe $z_n$ removes zero at separation two, after
which all background pairs remain unchanged through $\tau$.
The probe $w_n$ alone does the same at separation one.
With both probes present, $w_n$ first removes zero. The remaining
probe $z_n$ removes $x_1$ at separation three; then
$(x_2,x_3),(x_4,x_5),\ldots,(x_{2n-2},x_{2n-1})$
are removed at separations $5,7,\ldots,2n+1$.
Thus every one of the $n$ background edges is absent, while
$n+1$ new edges are present, including $(0,w_n)$.
The two single-probe distances are one, and the two-probe
distance is $2n+1$.

In the mixed vector, $(0,w_n)$ cancels. The $n$ background
edges have coefficient $-1$, the $n-1$ new nonincident cascade
edges have coefficient $+1$, $(0,z_n)$ has coefficient $-1$,
and $(z_n,x_1)$ has coefficient $+1$. Its norm is $2n+1$.
The eight protected background points are zero and the seven
far points; the two probes occupy adjacent cutoff gaps.
All decisions through the cutoff use strict inequalities.
An arbitrarily small generic perturbation preserves them and
removes irrelevant distance ties among the far points.
\end{proof}

\subsection{Second moments for tagged parent witnesses}
\label{sec:tagged-parent-second-moments}

\status{Working proof: individual rectangles, with their own corner backgrounds}

\begin{lemma}[Empirical versions of the parent-overlap kernel]
\label{lem:tagged-parent-kernel-L2}
For a size-$M$ iid sample, let $\xi_M$ be its empirical measure
and let $\xi_M^{[2]}$ be the uniform probability measure on
ordered pairs of distinct original points. For the closed-parent
kernel $h_m$ in Proposition~\ref{prop:common-parent-probe-L2},
with attainable $m\ge3$,
\begin{equation}\label{eq:tagged-parent-kernel-L2}
 \EE\int h_m^2\,d(\Leb\otimes\xi_M)
       +\EE\int h_m^2\,d\xi_M^{[2]}
       \le C/m.
\end{equation}
The distinct-pair measure is essential: singleton intersections
of closed parent intervals then have zero mass.
\end{lemma}
\begin{proof}
We first record a summed empirical estimate for the minimum
arcs. At pre-collision count $j$, let $J_j$ be the closed arc
between the two colliding points, of length $s_j\le1/j$.
These arcs are distinct proper original-point intervals.
The beta/binomial argument in Lemma~\ref{lem:tagged-parent-tail}
applies to any selected family of distinct original-point intervals:
those spanning $k$ original gaps with length at least $k/(2M)$
have empirical mass at most four times their length, and the
expected sum of squared empirical masses of all exceptional
intervals is at most $C/M$. Consequently
\begin{equation}\label{eq:empirical-minimum-arcs}
 \EE\sum_{j>m}\xi_M(J_j)^2
       \le16\sum_{j>m}j^{-2}+C/M\le C/m.
\end{equation}
Only counts of the input parity enter these sums. Distinctness
holds because the two endpoints of $J_j$ are both removed at
that collision; no later collision has the same arc.

Decorate a current gap with
$B_i^{\rm emp}=\sum_{J\text{ earlier in gap }i}\xi_M(J)^2$.
At the minimum gap of length $s$, condition on the geometry
and its original count as in
Lemma~\ref{lem:minimum-gap-parent-mark}.
If the closed gap contains $k$ original points, identifying
its endpoints gives the fully shielded circle with $k-1$
points. Its empirical parent-square sum has expectation at
most a universal constant, by Lemma~\ref{lem:tagged-parent-tail}
with lower active count two. The conditional survivor-label
argument preserves that statistic too. Every global empirical
parent mass is at most $2k/M$ times its local empirical mass;
the factor two covers the identified endpoints of the full
root parent. Therefore
\[
 \EE[B_{\min}^{\rm emp}\mid\text{geometry and counts}]
                                      \le C\xi_M(I_{\min})^2.
\]
Sum this conditional estimate over $j>m$ and apply
\eqref{eq:empirical-minimum-arcs}. This gives the stronger
summed minimum-mark estimate
\begin{equation}\label{eq:empirical-minimum-mark}
 \sum_{j>m}\EE B_{\min,j}^{\rm emp}\le C/m.
\end{equation}
Also $\EE\sum_iB_i^{\rm emp}\le C/j$, by
Lemma~\ref{lem:tagged-parent-tail}.
The uniform cyclic-order argument applies to these decorations.

For the first kernel in \eqref{eq:tagged-parent-kernel-L2},
a new parent $I$ adds
$|I|\xi_M(I)+2\sum_{J\text{ in its daughters}}|J|\xi_M(J)$
to its squared norm. Any old parent outside those daughters
has at most a singleton intersection, of zero Lebesgue mass.
The daughter marks are bounded by half the physical-square
marks plus half the empirical-square marks. For the physical
marks their expected sum is $O(j^{-2})$, as in
Proposition~\ref{prop:common-parent-probe-L2}.
For empirical marks, conditional neighbour sampling gives
\[
 \EE[B_a^{\rm emp}+B_s^{\rm emp}+B_c^{\rm emp}]
 \le\EE B_{\min,j}^{\rm emp}
             +\frac{2}{j-1}\EE\sum_iB_i^{\rm emp}.
\]
Its sum over $j>m$ is at most $C/m$, by
\eqref{eq:empirical-minimum-mark} and the $C/j$ bound above.
The sum of the new-parent terms is already at most $C/m$
in expectation by Lemma~\ref{lem:tagged-parent-tail}.
This proves the first estimate.

For the distinct-tag kernel use the interval weight
$\xi_M^{[2]}(I\times I)=k_I(k_I-1)/(M(M-1))\le2\xi_M(I)^2$.
Singleton intersections contribute zero. Laminarity therefore
gives exactly the same new-parent and daughter-mark recurrence,
with this weight in place of $|I|\xi_M(I)$.
The empirical-square estimates just proved bound both terms
and give the second assertion.
\end{proof}

\begin{proposition}[An $L^2$ bound for one restricted matching rectangle]
\label{prop:matching-rectangle-L2}
Let $B$ have $M$ iid uniform points and let $z,w$ be independent
uniform probes. Write $F_{ab}$ for the matching-edge indicator
vector in $B+az+bw$, restricted to separation at most a common
possibly random cutoff $\tau$. Suppose all four corner
configurations have at least $\ell\ge6$ active points at
$\tau$, almost surely. Then
\begin{equation}\label{eq:matching-rectangle-L2}
 \EE\|F_{11}-F_{10}-F_{01}+F_{00}\|_1^2
       \le C/\ell.
\end{equation}
The bound also applies to any fixed two-probe rectangle within
a larger iid probe family, with a fixed subset of its other
probes included in $B$. It controls one rectangle; summing
their correlated contributions remains a separate issue.
\end{proposition}
\begin{proof}
The parent-witness argument in
Lemma~\ref{lem:matching-rectangle}, including its incident-edge
extension in Proposition~\ref{prop:early-history-linearization},
assigns every nonzero coordinate of the mixed vector to a
positive corner whose closed accepted parent contains both
probe locations. Each coordinate has absolute value at most
two. Every such witness before $\tau$ occurs at pre-collision
count at least $\ell$. Discard the random cutoff and sum all
parents satisfying this necessary count bound. The resulting
pathwise upper bound is twice the sum of the four corner
parent kernels evaluated at the two probe locations.

In a corner missing both probes these locations are independent
uniform points, so \eqref{eq:common-parent-kernel-L2} applies.
In a corner containing one probe, it is a uniform original tag
conditional on the unlabelled corner sample, while the absent
probe is independent uniform; use the first term of
\eqref{eq:tagged-parent-kernel-L2}. In a corner containing
both, their locations are a uniformly ordered pair of distinct
tags; use its second term. After discarding $\tau$, all these
parent kernels depend only on the unlabelled corner sample,
so this tag conditioning is valid even if $\tau$ was adaptive.
The corner sizes are $M,M+1,M+2$.
Replace the count threshold by the adjacent attainable count
in each parity; it is at least $\ell-2\ge4$.
Finally $(x_1+\cdots+x_4)^2\le4\sum_i x_i^2$ proves
\eqref{eq:matching-rectangle-L2}. Including the other fixed
probes in the iid background leaves this argument unchanged.
\end{proof}

\begin{corollary}[Probe-conditional variance of the actual remainder]
\label{cor:history-remainder-conditional-variance}
Let $B$ be an iid size-$M$ background and let $Z_1,\ldots,Z_q$
be independent uniform probes. Partition their indices into
fixed disjoint sets $P,Q$. Stop all matching vectors at the
separation $\tau(B)$ where the background first reaches an
attainable count $m\ge q+6$. Write $F_S$ for the restricted
matching vector of $B+\{Z_i:i\in S\}$ and set
\[
 R=\|F_P-F_Q\|_1-
                 \sum_{i=1}^q\|F_{\{i\}}-F_\varnothing\|_1.
\]
Then
\begin{equation}\label{eq:history-remainder-conditional-variance}
 \EE\operatorname{Var}(R\mid B)
                \le \frac{Cq(q-1)^2}{m-q}.
\end{equation}
This concerns the actual changing-background remainder.
It does not bound $\operatorname{Var}(\EE[R\mid B])$;
in particular it does not yet give an unconditional centered
history comparison at the conjectured scale.
\end{corollary}
\begin{proof}
The following deterministic calculation holds for any family
of Boolean edge vectors. Remove index $i\in P$ and write
$U=P\setminus\{i\}$, $V=Q$, $F_0=F_\varnothing$,
$e=F_{U\cup\{i\}}-F_U$, and $e_0=F_{\{i\}}-F_0$.
For Boolean vectors $A,C,D$,
$\|A-D\|_1-\|C-D\|_1=\langle1-2D,A-C\rangle$.
Hence the change $D_iR$ on inserting $i$, with the corresponding
single-probe cost subtracted, satisfies
\[
 |D_iR|\le\|e-e_0\|_1
       +2\,|\operatorname{supp}(e_0)
                       \cap\operatorname{supp}(F_V-F_0)|.
\]
The first term telescopes over $U$ into the sum of the norms
of the $|U|$ rectangles with directions $i,j$ and successive
prefixes of $U$ as their backgrounds. For the second, if a
Boolean coordinate flips under both $\{i\}$ and $V$ from
$F_0$, then
\[
 |F_{V\cup\{i\}}-F_V-F_{\{i\}}+F_0|\ge1
\]
at that coordinate. Telescoping this mixed vector over $V$
bounds the support intersection by another $|V|$ rectangles.
Thus $|D_iR|$ is at most twice a sum of $q-1$ rectangle norms.
For $i\in Q$ interchange $P,Q$.

Every corner in these telescopes contains at most $q$ probes,
so discrepancy contraction leaves at least $m-q$ points at
$\tau(B)$. Proposition~\ref{prop:matching-rectangle-L2} and
Cauchy--Schwarz give
\[
 \EE(D_iR)^2\le C(q-1)^2/(m-q).
\]
Given $B$, replace $Z_i$ by an independent copy and denote
the resulting remainder by $R^{(i)}$. The two insertion
increments have the same smaller probe family, and
$(R-R^{(i)})^2\le2(D_iR)^2+2(D_iR^{(i)})^2$.
Conditional Efron--Stein, followed by expectation in $B$,
therefore proves \eqref{eq:history-remainder-conditional-variance}.
All prefixes use fixed index orders, so their background
samples are iid before the corner estimate is applied.
\end{proof}

\subsection{A scalar integral for the weighted-history variance}
\label{sec:weighted-history-integral}

\status{Working proof: exact moment reduction}

\begin{proposition}[A one-dimensional integral for the parent-square variance]
\label{prop:weighted-history-integral}
Let $\beta$ be the expected full parent-square mark in a fully
shielded entrance interval of length one, and put $\mu=1/c=2e^\gamma$.
For $p>0$ write
\[
 F(p)=\EE e^{-pY}=\tanh(E_1(p)/2),\qquad
 C(p)=(1-\partial_p)^2F(p)^2.
\]
Then
\begin{equation}\label{eq:parent-beta-integral}
 \beta=\int_0^\infty
        \frac{p e^{-p}}{1-F(p)^2}\,C(p)\,dp.
\end{equation}
The integrand is positive, has continuous value $\kappa$ at
zero, and decays exponentially at infinity. The active-count
variance in Corollary~\ref{cor:parent-count-moments} is
\begin{equation}\label{eq:parent-variance-polynomial}
 \begin{split}
 v_{\rm B}={}&\frac{14}{9}+\frac{136}{9}c
       +\frac{683}{18}c^2
       +\left(\frac{884}{27}-\frac89\beta\right)c^3\\
       &+\left(\frac{76}{9}-\frac13\beta\right)c^4+c^5.
 \end{split}
\end{equation}
Thus this variance and the full two-dimensional covariance in
\eqref{eq:parent-square-joint-CLT} reduce to the same scalar
integral. These constants concern weighted parent histories,
not the survivor-noise probability.
\end{proposition}
\begin{proof}
Let $b_t(x)$ be the conditional mean parent-square mark in a
cutoff-$t$ gap of length $x$, and set
\[
 W_t(p)=\int_t^\infty e^{-px}w_t(x)b_t(x)\,dx
             =t^2W(pt),\qquad
 W(p)=\EE[\mathcal B e^{-pY}].
\]
The minimum gap is fully shielded, so $b_t(t)=\beta t^2$.
Differentiate the last-pair integral for the marked gap.
A newly formed parent of length $x$ has mark equal to $x^2$
plus its three daughter marks. At fixed $p$, the moving lower
endpoint and the triple merger therefore give
\[
 \partial_t W_t(p)=\frac{e^{-pt}}{2t}
 \left\{\beta t^2(F_t(p)^2-1)+2F_t(p)W_t(p)
                           +(t-\partial_p)^2F_t(p)^2\right\},
\]
where $F_t(p)=F(pt)$. The conditional tree-moment bound and
exponential gap tails justify differentiation and integration;
alternatively the identity follows first for truncated positive
marks and then by monotone convergence. With $t=1$ this reads
\begin{equation}\label{eq:parent-mark-transform-ODE}
 2W+pW'=\frac{e^{-p}}2
       \{\beta(F^2-1)+2FW+C\}.
\end{equation}
Since
$F'=-e^{-p}(1-F^2)/(2p)$,
an integrating factor after division by $p$ is
$I(p)=p^2/(1-F^2)$. Hence
\[
 (IW)'=-\tfrac12\beta p e^{-p}
                  +\frac{p e^{-p}C(p)}{2(1-F(p)^2)}.
\]
At zero, $I(p)\sim p/(2\mu)$ and $W(0)<\infty$.
At infinity, $W(p)=O(e^{-p}/p)$ by the conditional quadratic
mark bound, while $I(p)=O(p^2)$. Thus both boundary terms
vanish. Integrating and using $\int_0^\infty p e^{-p}dp=1$
proves \eqref{eq:parent-beta-integral}. Positivity follows from
\[
 C(p)=\EE[(1+Y_1+Y_2)^2e^{-p(Y_1+Y_2)}].
\]
At zero this is $1+8\mu+4\mu^2$, giving integrand limit
$(1+8\mu+4\mu^2)/(2\mu)=\kappa$.

Here are explicit moment identities to verify the variance
reduction. Put $m_j=\EE Y^j$, $b_j=\EE[\mathcal B Y^j]$,
$A_j=\EE(1+Y_1+Y_2)^j$, and $s_0=\EE\mathcal B^2$.
The first four gap moments from the Taylor expansion of $F$ are
\[
 \begin{gathered}
 m_1=\mu,\qquad m_2=\mu^2+2\mu,\\
 m_3=\tfrac32\mu^3+6\mu^2+\tfrac92\mu,\qquad
 m_4=3\mu^4+18\mu^3+30\mu^2+\tfrac{34}{3}\mu.
 \end{gathered}
\]
In particular $b_0=A_2/2=\kappa\mu$.
The first two derivatives of
\eqref{eq:parent-mark-transform-ODE} at zero give
\begin{align*}
 2b_1&=\beta\mu+(1+\mu)b_0+A_3/2,\\
 3b_2&=\beta(A_2-1)/2+2(1+\mu)b_1
                            +(1+2\mu+m_2)b_0+A_4/2.
\end{align*}
For completeness the second tree-mark moment closes without
requiring the second moment of the fully shielded mark.
On a merger the new mark is
$\mathcal B_s+\mathcal B_a+\mathcal B_c+(s+a+c)^2$.
Square this identity and use independence of the three conditional
daughter trees. The minimum-gap second moment cancels its moving
endpoint term at $p=0$. Scaling the second-mark transform by
$t^4$ leaves
\[
 3s_0=b_0^2+4\beta b_0
       +2\{b_2+2(1+\mu)b_1+(1+2\mu+m_2)b_0\}+A_4/2.
\]
All fourth moments needed here are finite by
Lemma~\ref{lem:parent-square-tree-mark}.
Finally
\[
 v_{\rm B}=c^4\{s_0-4\kappa b_1
                           +4\kappa^2m_2-\kappa^2\mu^2\}.
\]
Substitution gives \eqref{eq:parent-variance-polynomial}.
For the joint covariance the remaining entries are
$c^3(m_2-\mu^2)$, $c^2(\kappa m_2-b_1)$, and
$c(s_0-2\kappa b_1+\kappa^2m_2)$, in that order.
\end{proof}

\begin{corollary}[Certified weighted-history constants]
\label{cor:certified-parent-variance}
The exact constants in
Proposition~\ref{prop:weighted-history-integral} satisfy
\begin{equation}\label{eq:certified-parent-variance}
 \begin{gathered}
 1.875825564\le\beta\le1.880616052,\\
 9.525737325\le v_{\rm B}\le9.525841453.
 \end{gathered}
\end{equation}
These bounds evaluate the working proof's weighted-history
constants; they do not evaluate $v_{\rm H}$ or $p_\infty$.
\end{corollary}
\begin{proof}
The reproducible integer-interval calculation is
\nolinkurl{certify_weighted_history.py}; its output is
\nolinkurl{weighted-history-certificate.json}.
Here is its analytic enclosure rule. Write
\[
 u(p)=e^\gamma p e^{-\operatorname{Ein}(p)},\qquad
 A(p)=e^{\operatorname{Ein}(p)-p},\qquad B(p)=(1+u(p))^2.
\]
The integrand in \eqref{eq:parent-beta-integral} is
$e^{-\gamma}A(p)B(p)C(p)/4$.
The functions $A,C$ are nonincreasing and $B$ is nondecreasing:
$\operatorname{Ein}'(p)=(1-e^{-p})/p\le1$,
$u(p)=e^{-E_1(p)}$ increases, and $C$ is a positive Laplace
transform. Thus on a whole cell $[a,b]$ its integral lies between
\[
 \frac{e^{-\gamma}}4(b-a)A(b)B(a)C(b)
 \quad\hbox{and}\quad
 \frac{e^{-\gamma}}4(b-a)A(a)B(b)C(a).
\]
There is no unestimated quadrature error.
The values at zero are $A=B=1$ and $C=1+8\mu+4\mu^2$.
For positive arguments put $d=-F'=e^{-p}(1-F^2)/(2p)$;
then $F''=d\{1+(1-e^{-p}F)/p\}$ and
$C=F^2+4Fd+2d^2+2FF''$. These positive formulas give stable
outward endpoint evaluations.

For $p\ge8$, use $F\le e^{-p}/(2p)$, $1-F^2\ge1/2$,
$d\le e^{-p}/(2p)$ and $F''\le d(1+1/p)$.
The remaining integral is at most
$11e^{-24}/48<8.652\cdot10^{-12}$.
The calculation uses $32768$ cells of width $2^{-12}$ on
$[0,8]$, integer intervals with denominator $10^{50}$,
the existing harmonic enclosure of $\gamma$, and explicit
Taylor-series remainder bounds for the exponential and
$\operatorname{Ein}$. Substitution in
\eqref{eq:parent-variance-polynomial} with outward rounding
proves \eqref{eq:certified-parent-variance}.
The script additionally passes $195$ independent Decimal
endpoint comparisons. Those comparisons check the implementation;
the interval bounds and monotonicities certify the enclosure.
\end{proof}

\subsection{Spatial blocks for the actual early-history remainder}
\label{sec:history-cluster-remainder}

\status{Working proof: Poisson line blocks; finite-circle transfer remains separate}

\begin{lemma}[Occupied-gap clusters retain the whole early history]
\label{lem:history-occupied-clusters}
Fix a generic finite background and a separation cutoff $r$
with at least two active points. Given finitely many probes,
call a cutoff gap occupied if its interior contains a probe.
Group consecutive occupied gaps into maximal circular runs.
If all gaps are occupied, use the single whole-circle run.
Let $Z_D$ be the probes in run $D$, and let $F_S$ denote the
matching vector through separation $r$ after inserting subset $S$.
Then
\begin{equation}\label{eq:history-cluster-superposition}
 F_S-F_\varnothing
       =\sum_D(F_{S\cap Z_D}-F_\varnothing).
\end{equation}
The increment supports of different runs are disjoint.
Consequently every pairwise history distance is the sum of
its run distances, and the additivity remainder is the sum of
its run remainders. A run with at most one probe has zero
remainder. These assertions hold simultaneously for all probe
subsets and for every common restriction to earlier separation
bands.
\end{lemma}
\begin{proof}
Every original matching edge inside an empty cutoff gap has
its closed parent inside that gap. Exterior protection keeps
all these edges, at their original separations, under every
probe subset. Thus its original interior points are matched
in exactly the same way in every copy. A changed edge cannot
use any of them. Nor can a changed edge through separation
$r$ cross the emptied gap: the two remaining boundary points
are more than $r$ apart along that gap. If the shorter circle
arc runs the other way, it crosses another empty gap whenever
it connects different runs, and is also too long.

All changed edges therefore lie inside the closed union of
one occupied run. These closed unions have disjoint vertex
sets, since neighbouring runs are separated by an empty gap
of positive length. The greedy scan within a run has no accepted
edge joining it to another run or to an empty-gap interior.
Its internal scan is consequently determined by its own initial
points, endpoints, and probes. This proves the vector identity
and disjoint supports. Restricting its coordinates to any
common earlier band preserves both properties.
\end{proof}

\begin{theorem}[Uniform moments of Poisson line-block remainders]
\label{thm:line-block-remainder-moments}
Take a stationary Poisson input on $\mathbb R$ of intensity
$\alpha\ge1$, and independent probe processes of two colours,
with intensities $\eta_P,\eta_Q$ and
$0\le\eta=\eta_P+\eta_Q\le\eta_*$, where $\eta_*<\infty$ is fixed.
At separation one, partition the background gap sequence into
blocks ending at its first gap containing no probe. A block
consists of its preceding occupied run and this last empty gap;
zero-length occupied runs are allowed. Write $L_D$ for its full
length, $K_D$ for its probe count, and $R_D$ for the actual
history distance between the two colour inputs through
separation one, minus the sum of the single-probe distances,
within its occupied run.

Under the renewal Palm law the pairs $(L_D,R_D)$ are iid.
For some $a>0$ and every fixed $p\ge1$,
\begin{equation}\label{eq:line-block-remainder-moments}
 \sup_{\alpha\ge1,\,\eta\le\eta_*}
      \EE e^{a(L_D+K_D)}<\infty,\qquad
 \sup_{\alpha\ge1,\,\eta\le\eta_*}\EE|R_D|^p<\infty.
\end{equation}
The same assertions hold for the line entrance law, denoted by
$\alpha=\infty$. For fixed probe intensities, the block laws
converge as $\alpha\to\infty$, with all fixed joint moments.
In particular the variance of the sum of these actual block remainders on a
long Poisson line has order at most its spatial length. No
fixed-original-count circle comparison is asserted here.
\end{theorem}
\begin{proof}
If $\eta=0$, every reward is zero; assume $\eta>0$ below.
At cutoff one the background has iid gaps with the
finite-intensity renewal density, independently marked by
their conditional full pasts. Conditional on these gaps, the
probe restrictions are independent Poisson processes. The
marked gaps are therefore iid. A gap of length $Y$ is empty
with probability $e^{-\eta Y}$. The uniform exponential gap
bounds for $\alpha\ge1$ give a uniformly positive empty
probability and uniform exponential moments of the gap length.
Stopping the iid sequence at its first empty gap gives iid
blocks. For sufficiently small $a>0$,
\[
 \EE[e^{a(Y+K)};K>0]<1
\]
uniformly, so a geometric-series calculation proves the first
bound in \eqref{eq:line-block-remainder-moments}. The block
intensity $\nu$ is bounded above and below by positive constants.
Under the Palm law of a probe, its containing gap is length
biased and its remaining probes are still Poisson; the runs
on either side stop at independent empty gaps. Thus $K_D$
also has uniform exponential moments under this probe Palm law.
The localization lemma makes each block reward a function of
only that block's independent marks.

We next record a uniform single-insertion bound. For a
stationary Poisson background of any positive intensity,
the number $J$ of changed edges in a band $(u,2u]$ has
\begin{equation}\label{eq:line-single-band-moments}
 \sup \EE J^p\le C_p
\end{equation}
for every fixed $p$. At cutoff $u$ the insertion changes one
active point in the gap containing its origin. Expose the
neighbouring gap strings until encountering gaps longer than
$4u$ on both sides. The extra point can split only its original
gap; all other long gaps remain barriers through $2u$.
The number of exposed endpoints dominates $J$ up to a constant.
After scaling by $u$, each outside gap exceeds four with a
probability uniformly bounded below, over all finite-intensity
parameters, by Lemma~\ref{lem:finite-intensity-envelope}.
The two string counts are geometrically bounded, proving
\eqref{eq:line-single-band-moments}. The same argument applies
to the full history below $b$ if the initial intensity times
$b$ is bounded: use the original exponential gaps as barriers.

Give every probe an independent uniform priority. Telescope
from the background to each colour input in priority order.
At a tagged probe of priority $t$, its predecessor background
is a stationary Poisson process of intensity
$\alpha+t\eta_P$ or $\alpha+t\eta_Q$. For a fixed band let
$J_z$ be the sum of its telescoping cost and its single-probe
cost relative to the original background. The triangle inequality
bounds the absolute block remainder by $\sum_{z\in D}J_z$.
The localization lemma justifies using the global stationary
prefix: probes in other blocks cannot change this block's
history through separation one.

For any integer $p\ge1$,
$(\sum J_z)^p\le K_D^{p-1}\sum J_z^p$.
Assign each summand to its probe, instead of its block start.
Stationarity and Fubini give the mass-transport identity
\[
 \nu\,\EE\sum_{z\in D}K_D^{p-1}J_z^p
       =\eta\,\EE^{\rm probe}[K_D^{p-1}J_0^p].
\]
H\"older, the probe Palm exponential bound, and
\eqref{eq:line-single-band-moments} bound the right side by
$C_p$. Thus every band remainder has uniform moments of all
fixed orders, even though the entire single-probe history
need not have a uniform bound.

There is also a rare-interaction gain at small scales. If
no two probes of a block occupy the same or adjacent background
gaps at cutoff $b$, its entire remainder through $b$ vanishes
by Lemma~\ref{lem:history-occupied-clusters}. When
$\alpha b\ge1/2$, the finer gaps have
$\EE Y_b^2/\EE Y_b\le Cb$ and $\EE Y_b\le Cb$.
The expected number per unit length of probe pairs in the same
or adjacent finer gaps is consequently at most
\[
 C\eta^2\frac{\EE Y_b^2+(\EE Y_b)^2}{\EE Y_b}\le Cb.
\]
Every bad block contains at least one such pair; dividing by
$\nu$ gives a bad-block probability at most $Cb$.
Combining this with the $2p$th moment bound yields
\begin{equation}\label{eq:line-block-band-smallness}
 \|R_D((b/2,b])\|_p\le C_p b^{1/(2p)},
                    \qquad \alpha b\ge1/2.
\end{equation}

Put $n_\alpha=\lceil\log_2\alpha\rceil$ and
$b_\alpha=2^{-n_\alpha}\in(1/(2\alpha),1/\alpha]$.
The whole history below $b_\alpha$ obeys the same moment
argument, since every prefix intensity times $b_\alpha$
is at most $1+\eta_*$. Its bad-block probability is
$O(b_\alpha)$, so its $L^p$ norm is at most
$C_pb_\alpha^{1/(2p)}$. Sum
\eqref{eq:line-block-band-smallness} over the dyadic bands
between $b_\alpha$ and one. The geometric sum is uniform
in $\alpha$, proving the second bound in
\eqref{eq:line-block-remainder-moments}. The same argument
bounds the remainder below any $b$ by $C_pb^{1/(2p)}$
whenever $\alpha b\ge1$.

For the entrance law take the same marked renewal construction.
Finite collections of gap trees and fixed positive-cutoff
probe responses converge by the finite-intensity comparison
and contracting-query arguments in
Sections~\ref{sec:finite-history-CLT} and \ref{sec:coupled-limit}.
Truncate the block length, probe count, and number of gaps
first; their exponential bounds remove the truncations.
At a fixed positive lower cutoff the block reward therefore
converges in law, and its uniform higher moments give moment
convergence. The dyadic tail bound, uniform once $\alpha b\ge1$,
then removes the lower cutoff in every fixed $L^p$.
In the entrance construction the bad-block probabilities at
dyadic scales are summable, so eventually the probes are all
separated and the earlier remainder is identically zero.
This also defines its finite, full-past block remainder
without subtracting two divergent infinite history counts.
\end{proof}

\begin{corollary}[A spatial Gaussian limit for the Poisson remainder]
\label{cor:line-remainder-spatial-CLT}
Fix the two probe intensities. Let $\mathcal R_{\alpha,T}$
be the sum of $R_D$ over the stationary blocks whose left
endpoints lie in $[0,T]$. Set
\[
 \rho_\alpha=\frac{\EE R_D}{\EE L_D},\qquad
 v_{R,\alpha}=\frac{\EE(R_D-\rho_\alpha L_D)^2}{\EE L_D}.
\]
Uniformly in $\alpha\ge1$,
$\operatorname{Var}\mathcal R_{\alpha,T}\le C(T+1)$.
If $\alpha_T\to\infty$ and $T\to\infty$, then
\[
 \frac{\mathcal R_{\alpha_T,T}-\rho_{\alpha_T}T}{\sqrt T}
       \Longrightarrow\mathcal N(0,v_{R,\infty}),
 \qquad
 \frac1T\operatorname{Var}\mathcal R_{\alpha_T,T}
                       \longrightarrow v_{R,\infty}.
\]
The limiting variance is strictly positive if $\eta_P+\eta_Q>0$,
and is zero if both intensities vanish. Its value is not evaluated here.
\end{corollary}
\begin{proof}
Apply the iid renewal-reward argument to
$R_D-\rho_\alpha L_D$. The preceding theorem gives uniform
moments of every fixed order and convergence of those moments;
$\EE L_D$ is bounded above and below. Since $L_D\ge1$,
the renewal stopping index is at most $T+1$ after the first
boundary block. The stopped centered sum has second moment
at most $C(T+1)$ by the martingale square identity.
The overshoot and the stationary first block have uniformly
bounded moments, using the exponential length bound and
H\"older for the length-biased reward. This proves the variance
bound. The triangular-array Lindeberg theorem and renewal
count law give the stated Gaussian limit. Applying the same
stopped-sum bound at a higher moment order gives uniform
integrability and the variance limit. Stationarity makes the
mean exactly $\rho_\alpha T$ for the stated block-start sum.

For positivity, the final empty-gap length $E$ of a block is
independent of its occupied run and reward; it has a continuous,
nondegenerate density. If $R_D-\rho_\infty L_D=0$ almost surely,
conditioning on the occupied run forces $\rho_\infty=0$, and
then $R_D=0$ almost surely. We give a positive-probability
entrance event with $R_D=-1$.

Fix $0<\varepsilon<1/4$. At cutoff $\varepsilon$, choose
adjacent points $a,b$ with gap length in $(2\varepsilon,1)$,
and with the outside gaps on each side longer than two.
This local gap pattern has positive probability. Endpoint
regularity supplies sets of positive origin measure in the
left outside gap whose insertion removes $a$ by $\varepsilon$,
and in the right outside gap whose insertion removes $b$.
The two source gaps are nonadjacent, so their full histories
through $\varepsilon$ superpose exactly. Take one probe in
each set, of the same colour of positive intensity, and no
other probe in this cutoff-one gap or its two neighbouring
gaps. This probe event has positive conditional probability.

Between $\varepsilon$ and one the background accepts only
$(a,b)$ in this isolated component. Each single-probe copy
and the two-probe copy have already removed at least one of
its endpoints and accept no edge in the component. The band
contributes one to the two-probe distance and one to each
single-probe distance. Its remainder is $-1$, while the earlier
remainder is zero. The cutoff-one occupied run is just this
merged gap, so its block has $R_D=-1$. This contradicts an
identically zero block reward and proves positivity.
\end{proof}

\subsection{The early-history remainder on a Poisson circle}
\label{sec:circle-Poisson-remainder}

\status{Working proof: the circle transfer, before conditioning exact counts}

\begin{lemma}[Uniform circle-band moments]
\label{lem:circle-remainder-polynomial-moments}
On a circle of circumference $T\ge8$, take a Poisson background
of intensity $\alpha\ge1$ and the two probe colours of
Theorem~\ref{thm:line-block-remainder-moments}. Let $R_T$ be
the actual history additivity remainder through separation one.
For each fixed $p\ge1$,
\begin{equation}\label{eq:circle-remainder-polynomial-moments}
 \|R_T\|_p\le C_p T(1+\log(T+2)),
\end{equation}
uniformly in $\alpha$. The same holds conditional on either
parity of the original count.
\end{lemma}
\begin{proof}
We first extend \eqref{eq:line-single-band-moments} to the
circle in every band used above its microscopic scale.
Use the initial strip $b_\alpha=2^{-\lfloor\log_2\alpha\rfloor}$,
so $1\le\alpha b_\alpha<2$. At the lower cutoff $u$ of any
remaining band, rescale to circumference $T/u\ge8$.
The intensity parameter is at least one. Its renewal density
away from the atom at zero is uniformly bounded, and its gap
density has a uniform exponential envelope. The short-gap
subdensity $f\ind_{[1,4]}$ has mass at most $\theta<1$.
In the rooted cyclic law at a fixed origin, exposing $j$ short
neighbouring gaps therefore costs at most $C\theta^j$ after
the outside renewal density is integrated out. For the atomic
completion, the convolution with the origin gap has density
bounded by $C\theta_1^j$, $\theta_1<1$: exponentially tilt
by a small positive parameter and use the uniformly bounded
weighted density $x f(x)$ of that gap. This also covers a
string that wraps around the circle. Empty and singleton
background configurations add at most a constant to the
single-insertion band cost. Thus the geometric-component
proof gives uniform moments of that cost. For the initial
strip use the original exponential gaps; their intensity times
the strip width is bounded, as in the line proof. Conditioning
on original parity costs at most a constant because
$\alpha T\ge8$.

Poisson Campbell and H\"older now give an $L^p$ bound $C_pT$
for each block-free circle-band remainder: the total probe
count has moments of order $T^p$, and each probe's priority
prefix has the uniform single-band moments just proved.
At upper cutoff $b$, $\alpha b\ge1$, the cyclic one- and
two-gap moment estimates bound the expected number of probe
pairs in the same or adjacent gaps by $CTb$. Empty or singleton
background configurations have probability at most
$Ce^{-cT/b}$ and may be included in that bound. If none of
these events occurs, the remainder through $b$ is zero by
Lemma~\ref{lem:history-occupied-clusters}. Consequently
\[
 \|R_T((b/2,b])\|_p
                  \le C_pT\min\{1,(Tb)^{1/(2p)}\}.
\]
The initial strip has the same estimate. There are only
$O(1+\log(T+2))$ dyadic bands above $b=1/T$; the remaining
bounds sum geometrically. This proves
\eqref{eq:circle-remainder-polynomial-moments}.
\end{proof}

\begin{theorem}[Poisson-circle remainder variance and central limit theorem]
\label{thm:Poisson-circle-remainder-CLT}
Fix $\eta_P,\eta_Q$ with $0<\eta_P+\eta_Q<\infty$.
Use the circle setup of the preceding lemma, optionally
conditioning on either parity of the original count.
Let $\rho_\alpha,v_{R,\alpha}$ be the line block constants
in Corollary~\ref{cor:line-remainder-spatial-CLT}.
Uniformly for $\alpha\ge1$ and $T\ge8$,
\begin{equation}\label{eq:Poisson-circle-remainder-moments}
 \EE R_T=\rho_\alpha T+O(e^{-cT}),\qquad
 \operatorname{Var}R_T=v_{R,\alpha}T+O(e^{-cT}),
\end{equation}
after decreasing $c>0$ if necessary. Also
$\EE|R_T-\rho_\alpha T|^{2p}\le C_p T^p$ for every fixed
integer $p\ge1$. If $\alpha_T\to\infty$
and $T\to\infty$, then
\[
 (R_T-\rho_{\alpha_T}T)/\sqrt T
                  \Longrightarrow\mathcal N(0,v_{R,\infty}),
\]
with convergence of every fixed centered moment. The limiting
variance is strictly positive. Original and probe counts are
still Poisson, apart from the stated optional parity condition.
\end{theorem}
\begin{proof}
Put $f=f_\alpha$, $e(y)=f(y)e^{-\eta y}$ and $h=f-e$.
The line block-length density and its root-gap-parity signed
version are
\[
 g=\sum_{j\ge0}h^{*j}*e,\qquad
 g_-=-\sum_{j\ge0}(-h)^{*j}*e.
\]
Write $U_\alpha=\sum_{j\ge1}f^{*j}$ and
$D_\alpha=\sum_{j\ge1}(-1)^{j-1}f^{*j}$.
The two block renewal sums have the exact identities
\begin{equation}\label{eq:block-renewal-identities}
 U_g:=\sum_{j\ge1}g^{*j}=e+e*U_\alpha,
 \qquad
 U_-:=\sum_{j\ge1}g_-^{*j}=-e+e*D_\alpha.
\end{equation}
For example their transforms are respectively
$\widehat e/(1-\widehat f)$ and
$-\widehat e/(1+\widehat f)$.
The known uniform renewal estimates therefore give
$U_g(y)=\nu_\alpha+O(e^{-cy})$ and
$U_-(y)=O(e^{-cy})$, where
$\nu_\alpha=\widehat e(0)/\EE Y=1/\EE L_D$.
These estimates include exponentially weighted integral bounds.

Every block reward moment density has an exponential envelope.
Indeed its final empty gap is independent of its occupied run
and reward, and has an exponentially bounded density. Write
$L_D=O_D+E_D$ for the occupied and empty lengths. Convolution
bounds the density of $\EE[|R_D|^p;L_D\in dy]$ by
$Ce^{-cy}\EE[|R_D|^p e^{cO_D}]$. The latter expectation is
uniformly finite by H\"older and
Theorem~\ref{thm:line-block-remainder-moments}, after choosing
$c$ small. The same bound holds with a root-gap-parity sign.

On sectors having at least one empty gap, group the cyclic gap
law into $j$ blocks. Rooting at an empty-gap boundary replaces
the gap divisor by the block divisor. Conditional on original
parity $p$, the marked density is
\[
 \frac{T}{jZ_p}
   \left\{\prod_{i=1}^j dG_i
                 +(-1)^p\prod_{i=1}^j dG_{-,i}\right\},
 \qquad Z_p=1+(-1)^p e^{-2\alpha T},
\]
with total block length constrained to $T$.
Here $G$ is the joint line block law and $G_-$ weights it by
$(-1)^J$, with $J$ the number of original cutoff gaps in the
block. The occupied-run localization proves that its reward
is exactly the circle's actual remainder; the intervening
empty gap has length at least one and prevents a shorter
outside arc from changing the run's scan.

Set $Z=R_D-\rho_\alpha L_D$, whose mean is zero. Apply the
cyclic marked-moment expansion from
\eqref{eq:marked-cyclic-moments}, now with the two renewal sums
in \eqref{eq:block-renewal-identities}. For the unsigned term,
the first-moment density convolved with $\delta_0+U_g$ has
exponentially small value at $T$, since its total mass is zero.
The square term is
\[
 T\left\{h_Z*(\delta_0+U_g)
       +g_Z*g_Z*(\delta_0+U_g)^{*2}\right\}(T).
\]
Here $g_Z,h_Z$ are the first and second mark moment densities.
The first term is $T\nu_\alpha\EE Z^2$ plus an exponential
error. The convolution square of $\delta_0+U_g$ is affine
on the positive half-line up to an exponential error; both
its polynomial terms vanish against $g_Z*g_Z$, whose mass
and first moment are zero. All signed terms are exponentially
small by the second identity in
\eqref{eq:block-renewal-identities}. The same set-partition
expansion as in Lemma~\ref{lem:parent-band-higher-moments}
gives uniform normalized centered moments of every fixed order.

The omitted sector is harmless but requires a separate bound.
The density $h$ has mass strictly below one and a uniform
exponential moment. A small exponential tilt therefore makes
its renewal sum exponentially decreasing. The cyclic density
then gives probability $Ce^{-cT}$ that every active gap is
occupied. An empty active configuration has the same bound.
By \eqref{eq:circle-remainder-polynomial-moments} and
H\"older, every fixed reward moment on this exceptional sector
is exponentially small as well, uniformly in $\alpha$.
Thus the preceding moment calculations prove
\eqref{eq:Poisson-circle-remainder-moments} for the full law.

For the weak limit, tilt by the block count and root at a
uniform block start. The outside renewal ratio is now formed
from $U_g$ and $U_-$; the latter is exponentially small.
Stop the iid block prefix at distance $T-T^{1/4}$.
The uniform exponential length bound makes its outside ratio
tend to one off an exponentially unlikely overshoot event.
The block-count tilt has vanishing total-variation cost by
its first two cyclic moments. The discarded arc contributes
$o_{\PP}(\sqrt T)$: it contains at most $T^{1/4}+1$ block
starts, while the maximum reward among $O(T)$ blocks is
$O_{\PP}(T^{1/p})$ for any fixed $p>4$; the two crossing
blocks have the corresponding length-biased moment bounds.
The iid triangular-array renewal-reward CLT from
Corollary~\ref{cor:line-remainder-spatial-CLT} applies.
The higher cyclic moment bounds give uniform integrability,
proving convergence of every fixed centered moment.
\end{proof}

\begin{corollary}[Conditioning the original count by an unused reservoir]
\label{cor:original-count-remainder-variance}
Fix the two probe intensities with positive total intensity.
There are $C_0,C_1,\alpha_0<\infty$ such that the following holds.
On a circle of circumference $T$, take exactly $M$ iid uniform
background points, and independent Poisson probe clouds of
these intensities. If
\[
 C_0\log(M+2)\le T\le M/\alpha_0,
\]
then the actual additivity remainder through separation one
satisfies
\begin{equation}\label{eq:original-count-remainder-variance}
 \operatorname{Var}R_T\le C_1T.
\end{equation}
The probe counts remain Poisson. This assertion gives no
exact-probe-count or late-history comparison.
\end{corollary}
\begin{proof}
Start instead with a Poisson background of intensity
$\alpha=M/T$. At cutoff one, let
\[
 X=\sum_iY_i\ind_{\{Y_i\le2,\ \text{gap }i\text{ contains no probe}\}}.
\]
There are $\delta,c>0$ such that
\begin{equation}\label{eq:empty-gap-reservoir-tail}
 \PP(X<\delta T)\le Ce^{-cT}
\end{equation}
uniformly in $\alpha\ge1$. To verify this, mark a gap by its
contribution to $X$. Its negative-tilted density is
\[
 f_\alpha(y)\left[1-e^{-\eta y}
        +e^{-\eta y}e^{-\theta y\ind_{\{y\le2\}}}\right].
\]
For fixed small $\theta>0$, this density has mass uniformly
below one, since $f_\alpha$ is uniformly positive on $[1,2]$.
It retains a uniform exponential moment. Its defective renewal
sum is therefore exponentially decreasing after a small positive
length tilt. The cyclic sector formula bounds $\EE e^{-\theta X}$
by $Ce^{-c_2T}$, including the exponentially small empty sector.
Markov's inequality, with $\theta\delta<c_2$, proves
\eqref{eq:empty-gap-reservoir-tail}.

Reveal the cutoff geometry, all probes, and the complete past
inside every occupied gap. On $X\ge\delta T$, the remainder $R_T$ is measurable
with respect to this information, by occupied-run localization.
The original counts inside the empty gaps remain independent
odd variables conditional on their lengths; their past trees
were not used. On $X\ge\delta T$, restrict to the empty gaps
of lengths in $[1,2]$. The characteristic-function estimate
in \eqref{eq:odd-Poisson-CF-bound}, applied just to
this subcollection, gives
\[
 \sup_n\PP(P=n\mid\text{revealed information})
                   \le C(\alpha T)^{-1/2}=C M^{-1/2},
\]
for $\alpha\ge\alpha_0$. Indeed the product bound is
$\exp\{-c\alpha X\,\operatorname{dist}(t,\pi\mathbb Z)^2\}$;
Fourier inversion integrates its two lattice peaks. Additional
unrevealed empty-gap counts can only decrease the maximal
point mass. The known occupied counts supply a fixed offset.

Let $\mu=\EE_{\rm Pois}R_T$. The Poisson point probability
at its integer mean obeys $\PP(P=M)\ge cM^{-1/2}$.
Consequently on the good reservoir event the conditional
second moment of $R_T-\mu$, given $P=M$, is at most a constant
times its unconditional second moment, hence at most $CT$
by Theorem~\ref{thm:Poisson-circle-remainder-CLT}.
On the complementary event, use the global conditioning bound,
H\"older, the fourth centered moment $O(T^2)$ from that theorem,
and \eqref{eq:empty-gap-reservoir-tail}. Its contribution is
at most
\[
 C\sqrt M\,T e^{-cT/2}\le CT
\]
when $T\ge C_0\log(M+2)$, with $C_0$ sufficiently large.
Finally $\operatorname{Var}(R_T\mid P=M)$ is no larger than
this second moment about $\mu$. Conditioning a Poisson input
on $P=M$ gives the required iid background and leaves the
independent probe clouds unchanged.
\end{proof}

\subsection{Conditioning the probe counts through local increments}
\label{sec:fixed-probe-remainder}

\status{Working proof: negligible centered early remainder for all sparse edit counts}

In this subsection, $R_T(P,Q)$ is the actual history distance
between the background augmented by $P$ and by $Q$, through
separation one on the circle of circumference $T$, minus the
sum of the $|P|+|Q|$ single-probe distances from the background.
One colour may be absent. Constants below are uniform when
the total probe intensity belongs to a fixed compact subinterval
of $(0,\infty)$; the individual colour intensities may vanish.

\begin{lemma}[A local increment of the remainder]
\label{lem:local-remainder-increment}
Take the Poisson circle of Theorem~\ref{thm:Poisson-circle-remainder-CLT},
and insert one further, independent uniform probe $Z$ of either colour, including an absent colour. Write
\[
 D_ZR_T=R_T(P\cup\{Z\},Q)-R_T(P,Q)
\]
if that colour is $P$, and analogously for $Q$.
For every fixed $p\ge1$,
\begin{equation}\label{eq:local-remainder-increment}
 \sup_{\alpha\ge1,\ T\ge8}\EE|D_ZR_T|^p<\infty.
\end{equation}
Except on an event of probability $Ce^{-cT}$, the increment is
determined by the cutoff geometry, the empty/occupied flags,
and the complete probe and background records in an arc $J$
containing $Z$. Its length has a uniform exponential tail.
The same moment assertion holds for the entrance circle.
\end{lemma}
\begin{proof}
Force the gap containing $Z$ to be occupied, and extend its
occupied run in both directions until the first empty gaps.
Include these two guard gaps in $J$. Occupied-run localization
shows that the insertion affects at most two old runs and one
new run. The increment is the new run reward minus those old
run rewards. Outside runs cancel exactly, and the guard-gap
interiors are not needed for this computation.

The arc exploration has uniform exponential length and probe-count
bounds. On the line this follows from the geometric run of
occupied gaps, including the length bias of the origin gap.
For the circle the remaining outside renewal density is
uniformly bounded. A wrapping exploration is exponentially
unlikely: the occupied-gap subdensity has mass less than one
after a small positive length tilt, uniformly over bounded
total probe intensities. The origin length bias and at most
two guard gaps preserve that bound. Empty and singleton
cutoff configurations have the same exponential bound.

There is a direct moment bound that is uniform even when the
colour of $Z$ has intensity zero. In a separation band $I$,
the triangle inequality gives
\[
 |D_ZR_T(I)|\le J_Z(B+P;I)+J_Z(B;I)
\]
for colour $P$, where $J_Z(A;I)$ is the one-insertion history
cost in background $A$. Both backgrounds on the right are
stationary Poisson inputs, since $B+P$ has intensity
$\alpha+\eta_P$. The circle-band proof in
Lemma~\ref{lem:circle-remainder-polynomial-moments} therefore
gives uniform moments of every fixed order for this sum.
Use its microscopic strip
$b_\alpha=2^{-\lfloor\log_2\alpha\rfloor}$; the initial
intensity times this strip width is uniformly bounded for
both backgrounds.

At upper cutoff $b\ge b_\alpha$, if the gap containing $Z$
and its two neighbours contain no old probe, the new probe
forms a singleton run. Thus $D_ZR_T([0,b])=0$ by
Lemma~\ref{lem:history-occupied-clusters}. The expected total
length of these gaps is at most $Cb$, by the cyclic one- and
two-gap moment bounds. The bad-event probability is at most
$Cb$, including the exponentially unlikely empty/singleton
background sector. H\"older with the $2p$th band moment gives
\[
 \|D_ZR_T((b/2,b])\|_p\le C_pb^{1/(2p)}.
\]
The whole microscopic strip obeys the same bound, using
original exponential gaps for its single-insertion costs.
Summing dyadically proves \eqref{eq:local-remainder-increment}.
Conditioning on original parity costs a bounded factor.
At fixed $T$, marked convergence above a positive lower cutoff
and this uniform small-band bound pass to the entrance law.
The entrance assertion holds in either parity.
\end{proof}

\begin{lemma}[Two unused count reservoirs]
\label{lem:two-count-reservoirs}
Fix $0<u<v<\infty$. For every fixed $p\ge1$ there are
$A_p,C_p,\alpha_0<\infty$ such that the following holds.
The background consists of exactly $M$ iid uniform points,
and the two probe counts $k_1,k_2\ge0$ are fixed, with
$uT\le k_1+k_2\le vT$. All probe positions are independent
uniform points. If $A_p\log(M+2)\le T\le M/\alpha_0$,
adding an independent uniform probe of either colour satisfies
\begin{equation}\label{eq:fixed-count-local-increment}
 \EE|D_ZR_T|^p\le C_p.
\end{equation}
For a Poisson background of intensity $\alpha\ge1$, or an
entrance background, the same assertion holds for every
$T\ge8$, optionally with fixed original parity, when only
the probe counts are conditioned.
\end{lemma}
\begin{proof}
First use independent Poisson inputs with means $M,k_1,k_2$,
so $\alpha=M/T$ and $u\le\eta_1+\eta_2\le v$, where
$\eta_j=k_j/T$. Relabel the colours so that $k_1\ge k_2$.
Reveal the cutoff geometry, all empty/occupied flags, $Z$,
and the full local records in $J$ from the preceding lemma.
Call this information $\mathcal S$. The exploration of $J$
uses only the geometry, flags, and $Z$, so this adaptive
revelation does not examine any outside marks.

With probability at least $1-Ce^{-cT}$, the short empty gaps
outside $J$ have total length at least $\delta T$, and there
are at least $\delta T$ short occupied gaps outside $J$.
Here short means length in $[1,2]$. To check this, apply the
negative gap-reward tilt used in
\eqref{eq:empty-gap-reservoir-tail} both to short empty length
and to short occupied count. Each tilt decreases the mass by
a uniform positive amount, since $f_\alpha$ is uniformly
positive on $[1,2]$, the total probe intensity is bounded above,
and it is bounded away from zero. Their cyclic densities give
exponential lower tails. The exponential tail of $|J|$ then
removes its contribution: it removes at most $|J|$ gaps, and
at most $2|J|$ short-gap length. Denote this good event by $G$.
It belongs to $\mathcal S$, and on it $D_ZR_T$ is measurable
with respect to $\mathcal S$.

The unexamined original counts in the short empty gaps are
independent odd Poisson variables of parameters $\alpha Y_i$.
By \eqref{eq:odd-Poisson-CF-bound}, their sum has maximal point
mass at most $C/\sqrt M$ on $G$.
In a short occupied gap, the vector of probe counts consists
of independent Poisson variables of means $\eta_jY_i$,
conditioned on their sum being positive. These vectors are
independent across gaps, conditional on $\mathcal S$.
Their characteristic functions satisfy uniformly
\[
 |\phi_i(t_1,t_2)|\le
 \exp\{-c[\operatorname{dist}(t_1,2\pi\mathbb Z)^2
       +\eta_2\operatorname{dist}(t_2,2\pi\mathbb Z)^2]\}.
\]
Indeed $\eta_1\ge u/2$ and both parameters are bounded above.
The atoms $(1,0),(2,0)$ have probabilities bounded below by
$c$, while $(1,1)$ has probability at least $c\eta_2$.
Apply the identity for $|\phi_i|^2$ using two independent
copies and retain the pairs $((1,0),(2,0))$ and
$((1,0),(1,1))$. Fourier inversion for at least $\delta T$
such gaps gives the maximal vector point bound
\[
 C T^{-1/2}(1+\eta_2T)^{-1/2}
             =C\{T(1+k_2)\}^{-1/2}.
\]
This includes $k_2=0$, when the second coordinate is constant.

These two reservoirs are independent conditional on
$\mathcal S$, and independent of all remaining unexamined
variables. The first uses original counts in empty gaps;
the second uses probe counts in occupied gaps. Convolution
with the remaining joint counts cannot increase a maximal
point mass. Writing $N$ for the original Poisson count and
$\boldsymbol K$ for the probe count vector, we obtain on $G$
\begin{equation}\label{eq:joint-reservoir-point-bound}
 \sup_{n,\boldsymbol l}
 \PP(N=n,\boldsymbol K=\boldsymbol l\mid\mathcal S)
                \le C\{MT(1+k_2)\}^{-1/2}.
\end{equation}
The unconditional probability of the target counts is at
least $c\{MT(1+k_2)\}^{-1/2}$, by the Poisson point probabilities
at their integer means. Thus the contribution of $G$ to the
conditional $p$th increment moment is bounded by a constant
times its unconditional moment, which is bounded by
Lemma~\ref{lem:local-remainder-increment}. On $G^c$, global
conditioning and H\"older instead give
\[
 C\sqrt{MT(1+k_2)}
   (\EE|D_ZR_T|^{2p})^{1/2}\PP(G^c)^{1/2}
       \le C_p\sqrt M\,T e^{-cT/2}.
\]
This is bounded when $T\ge A_p\log(M+2)$, after increasing
$A_p$ and using $T\le M$. This proves
\eqref{eq:fixed-count-local-increment}.

If only probe counts are fixed, omit the original-count
reservoir and the conditioning on $N$, for either the Poisson
or entrance background. The same proof costs
only $C\sqrt{T(1+k_2)}\le CT$ on the exponentially rare bad event, which is
uniformly harmless for all $T\ge8$.
\end{proof}

\begin{theorem}[Exact-count variance of the early interaction remainder]
\label{thm:fixed-probe-remainder-variance}
Let $P,Q$ be disjoint independent probe
families on the unit circle, with fixed counts $d,a$ and
$q=d+a$. For an independent iid uniform background of size $M$, let
$R_{M;d,a}$ be the actual additivity remainder through
separation $1/q$. There are absolute $A,C,\alpha_0<\infty$ such that
\begin{equation}\label{eq:fixed-probe-remainder-variance}
 A\log(M+2)\le q\le M/\alpha_0
       \quad\Longrightarrow\quad
             \operatorname{Var}R_{M;d,a}\le Cq.
\end{equation}
For a Poisson background of intensity $\alpha q$ on the unit
circle, $\alpha\ge1$, or an entrance background, the analogous
variance bound holds for every $q\ge8$, in either original parity.
In particular the finite conclusion covers resampling
($d=a$) and pure additions or deletions. It concerns only
the early additivity remainder, not the complete history CLT.
\end{theorem}
\begin{proof}
Scale lengths by $q$, so $T=q$ and the separation cutoff is
one. Write the colour counts as $k_1\ge k_2\ge0$ by
relabelling if necessary. Replace them by independent
$K_j\sim\operatorname{Poisson}(k_j)$. Use the same independent
uniform sequence for its first $K_j$ or first $k_j$ points;
the counts are independent of every sequence and the background.
Clip $K_1$ to $[\lceil k_1/2\rceil,2k_1]$ and $K_2$ to
$[0,2q]$, obtaining $K_1^*,K_2^*$. Since $k_1\ge q/2$,
all intermediate total intensities lie in $[1/4,4]$.
Telescope from the clipped vector to the fixed vector one
colour at a time. Conditional on the
counts of other colours, the indicator that $K_j^*$ crosses
an index $l$ is independent of that step's input positions.
Minkowski and \eqref{eq:fixed-count-local-increment} with $p=2$
therefore give
\[
 \|R_{\boldsymbol K^*}-R_{\boldsymbol k}\|_2
 \le C\sum_j\sum_l
    \PP(K_j^*\text{ crosses }l)^{1/2}
 \le C\sum_j\sqrt{k_j}\le C\sqrt q.
\]
The second inequality follows by summing the Poisson
Chernoff bound $C\exp\{-c s^2/(k_j+s)\}$ at deviations
$s\ge0$. For $k_j=0$ that colour contributes zero.

For a finite background, the deterministic estimate
$|R_T|\le C(M+1)(K_1+K_2+1)$ controls unclipping.
Poisson exponential tails and their weighted moments give
\[
 \|(R_{\boldsymbol K}-R_{\boldsymbol K^*})
                       \ind_{\{\boldsymbol K\ne\boldsymbol K^*\}}\|_2
       \le C(M+1)(q+1)e^{-cq}\le1
\]
when $q\ge A\log(M+2)$, increasing $A$ if needed.
Thus $\|R_{\boldsymbol K}-R_{\boldsymbol k}\|_2\le C\sqrt q$.
Corollary~\ref{cor:original-count-remainder-variance} gives
$\operatorname{Var}R_{\boldsymbol K}\le Cq$, with constants
uniform over this compact total-intensity range. Centering the
coupling difference proves \eqref{eq:fixed-probe-remainder-variance}.

When the original background remains Poisson or has its
entrance law, the same telescoping argument uses the
probe-only part of Lemma~\ref{lem:two-count-reservoirs}.
Unclipping still has exponentially small second moment.
Indeed a fixed count vector in the clipped ranges has every
fixed reward moment bounded by a polynomial in $q$: condition
the Poisson circle at those counts, use their point probability
lower bound $c/q$, and apply
Lemma~\ref{lem:circle-remainder-polynomial-moments}, or its
entrance limit. This bound is uniform over the clipped
vectors. The unclipped Poisson reward has the same polynomial
moment bound directly. H\"older and the exponentially small
clipping probability prove the claim. Finally
Theorem~\ref{thm:Poisson-circle-remainder-CLT} bounds the
unconditioned Poisson or entrance variance by $Cq$.
\end{proof}

\begin{proposition}[A small-family comparison for the actual remainder]
\label{prop:small-family-remainder-comparison}
Fix $A<\infty$ and put $\chi=\chi_{\rm f}>0$ from
Lemma~\ref{lem:finite-family-rate}. For all sufficiently large
$M$ and $1\le q=d+a\le A\log(M+2)$, the finite and same-parity
entrance families can be coupled so that their actual remainders
through separation $1/q$ satisfy
\begin{equation}\label{eq:small-family-remainder-comparison}
 \|R_{M;d,a}-R_{\infty;d,a}\|_2
       \le C_A[\log(M+2)]^3 M^{-\chi/4}=o(1).
\end{equation}
The bound is uniform over the split $d+a=q$.
\end{proposition}
\begin{proof}
Write $\varepsilon=M^{-\chi}$, so $q^2\varepsilon\to0$.
First control the finite remainder below $\varepsilon$.
For some fixed $\theta,c>0$,
\begin{equation}\label{eq:early-count-lower-tail}
 \PP\{K_M(\varepsilon)\le\theta/\varepsilon\}
                \le C\sqrt M\,e^{-c/\varepsilon}.
\end{equation}
Indeed Poissonize the original count at mean $M$ and scale
by $\varepsilon$. The circumference is $1/\varepsilon$ and
the intensity parameter $M\varepsilon\ge1$. The uniform
exponential gap envelope gives
$f_\alpha^{*j}(T)\le Ce^{-sT}B^j$ for fixed $s>0,B<\infty$,
uniformly in $\alpha\ge1$. In the cyclic density $T/j$,
summing over $j\le\theta T$ is exponentially small if
$\theta\log B<s/2$. The empty sector has the same bound.
Conditioning the original Poisson count to be $M$ costs
at most $C\sqrt M$, proving \eqref{eq:early-count-lower-tail}.

Choose an attainable $m\asymp\varepsilon^{-1}$ with
$m\le\theta/\varepsilon$ and $m\ge2q+6$. Stop all inputs at
$\sigma=\min(\varepsilon,\tau_m(B))$, where $\tau_m(B)$ is
the separation at which the background first reaches $m$.
Discrepancy contraction leaves at least $m-q$ active points
in every corner. Telescope the deterministic insertion
inequality in the proof of
Corollary~\ref{cor:history-remainder-conditional-variance}
from no probes to all $q$ probes. It bounds $|R(\sigma)|$
by at most $Cq^2$ rectangle norms. Proposition~\ref{prop:matching-rectangle-L2}
and Minkowski therefore give
\[
 \|R(\sigma)\|_2\le Cq^2/\sqrt{m-q}
                         \le Cq^2\sqrt\varepsilon.
\]
For every fixed $p$, the finite ancestry tail gives an
$L^p$ bound $C_p\log(M+q+2)$ for one uniform insertion into
an iid background of size at most $M+q$. Telescoping a
fixed index order in each colour consequently bounds the
full remainder, and every band restriction, by
\begin{equation}\label{eq:finite-remainder-crude-moments}
 \|R_M\|_p\le C_pq\log(M+q+2).
\end{equation}
The histories through $\sigma$ and $\varepsilon$ agree if
$K_M(\varepsilon)>m$. H\"older and
\eqref{eq:early-count-lower-tail} now yield
\begin{equation}\label{eq:finite-remainder-small-cutoff}
 \|R_M([0,\varepsilon])\|_2
 \le Cq^2\sqrt\varepsilon
       +Cq\log(M+2)M^{1/8}e^{-c/(4\varepsilon)}.
\end{equation}

For the entrance family use the small-band estimate in the
proof of Lemma~\ref{lem:circle-remainder-polynomial-moments}.
On a circle of circumference $T\ge8$ with bounded total
probe intensity, summing the dyadic bands below $b\le1/T$
gives
\[
 \|R_T([0,b])\|_p\le C_pT(Tb)^{1/(2p)}.
\]
Take $T=\max(q,8)$, use Poisson probe means $d,a$, and set
$b=T\varepsilon$. Their joint probability at these exact
counts is at least $c/T$, also when one colour is absent.
Condition the second moment on these counts. The result is
\begin{equation}\label{eq:entrance-remainder-small-cutoff}
 \|R_\infty([0,\varepsilon])\|_2
                   \le CT^2\varepsilon^{1/4}
                   \le C(q+1)^2\varepsilon^{1/4}.
\end{equation}
The same count conditioning and the full-circle moment
bound give
\begin{equation}\label{eq:entrance-remainder-crude-moments}
 \|R_\infty([0,1/q])\|_4
                 \le C(q+1)^{5/4}(1+\log(q+2)).
\end{equation}
For $q<8$ the estimate first gives the history through $1/8$.
Every input has at most eight active points at that cutoff,
so the remaining band contributes at most $C(q+1)$ to
the remainder. Thus \eqref{eq:entrance-remainder-crude-moments}
holds for these finitely many counts as well.

Use Lemma~\ref{lem:finite-family-rate} to couple every subset
above $\varepsilon$, with failure probability
$C(q+1)^2\varepsilon$. On success their remainder contributions
from $(\varepsilon,1/q]$ coincide exactly. On failure,
H\"older, \eqref{eq:finite-remainder-crude-moments}, and
\eqref{eq:entrance-remainder-crude-moments} bound the $L^2$
cost by
\[
 C\{q\log(M+2)+(q+1)^{5/4}(1+\log(q+2))\}
                        [(q+1)^2\varepsilon]^{1/4}.
\]
Add \eqref{eq:finite-remainder-small-cutoff} and
\eqref{eq:entrance-remainder-small-cutoff}, and use
$q\le A\log(M+2)$. This proves
\eqref{eq:small-family-remainder-comparison}.
\end{proof}

\begin{corollary}[The centered early remainder is negligible for sparse edits]
\label{cor:all-sparse-remainder-variance}
There are absolute $C,\alpha_0<\infty$ such that, for every
$M$ and every split $q=d+a$ with $1\le q\le M/\alpha_0$,
\begin{equation}\label{eq:all-sparse-remainder-variance}
 \operatorname{Var}R_{M;d,a}\le Cq.
\end{equation}
Consequently, for any $q=q_M=o(M)$, uniformly over its split,
\[
 \frac{R_{M;d,a}-\EE R_{M;d,a}}
           {\sqrt{q\log(M/q)}}\longrightarrow0
                              \qquad\hbox{in }L^2.
\]
This removes the actual early additivity remainder at the
conjectured history fluctuation scale. The centered capture
sum and the late history after $1/q$ require separate comparisons,
supplied in Sections~\ref{sec:integrated-capture-fluctuations}--\ref{sec:late-history-variance}.
\end{corollary}
\begin{proof}
For $q\ge A\log(M+2)$ use
Theorem~\ref{thm:fixed-probe-remainder-variance}. Below this
threshold, Proposition~\ref{prop:small-family-remainder-comparison}
compares the remainder in $L^2$ to its entrance counterpart.
That counterpart has variance at most $Cq$ for $q\ge8$ by
the same theorem. For $q<8$ its second moment is uniformly
bounded by \eqref{eq:entrance-remainder-crude-moments}, which
also gives $Cq$ after increasing $C$. The finitely many
remaining small $M$ are covered by the deterministic finite
history bound. Dividing \eqref{eq:all-sparse-remainder-variance}
by $q\log(M/q)$ proves the last assertion.
\end{proof}

\subsection{Fluctuations of the integrated single-probe capture}
\label{sec:integrated-capture-fluctuations}

\status{Working proof for Poisson backgrounds; exact original count remains separate}

\begin{theorem}[Uniform centered moments of integrated capture]
\label{thm:Poisson-integrated-capture}
On a circle of circumference $T\ge8$, let $B$ be a Poisson
background of intensity $\alpha\ge1$, optionally conditioned
on either original parity. For an origin $z$, let $G_B(z)$
be the single-insertion matching-history distance through
separation one, and put
\[
 A_{\alpha,T}=\int_0^T G_B(z)\,dz.
\]
For every fixed integer $p\ge1$ there is $C_p<\infty$,
uniform in $\alpha,T$, such that
\begin{equation}\label{eq:integrated-capture-moments}
 \|A_{\alpha,T}-\EE A_{\alpha,T}\|_{2p}\le C_p\sqrt T.
\end{equation}
More precisely, for a dyadic band $(b/2,b]\subset(0,1]$
with $\alpha b\ge2$, its integrated cost $A(b/2,b]$ satisfies
\begin{equation}\label{eq:integrated-capture-band-moments}
 \|A(b/2,b]-\EE A(b/2,b]\|_{2p}\le C_p\sqrt{Tb}.
\end{equation}
The microscopic strip $[0,b_\alpha]$, where
$b_\alpha=2^{-\lfloor\log_2\alpha\rfloor}$, obeys the same
bound with $b=b_\alpha$.
\end{theorem}
\begin{proof}
At cutoff $b$, every insertion's changed history is confined
to its source gap, by
Lemma~\ref{lem:history-occupied-clusters} with one probe.
For a cutoff gap $I$ of length $x$, write
\[
 A_I=\int_I J_z((b/2,b])\,dz.
\]
At separation $b/2$ each of the two inputs has at most
$2x/b+2$ active points in this closed interval. Only those
points can participate in later edges of this band, and no
changed edge can cross another cutoff gap. Consequently
\begin{equation}\label{eq:integrated-capture-local-bound}
 0\le A_I\le Cx(1+x/b).
\end{equation}
This deterministic bound is independent of the depth of the
past tree. The mark is a function of that gap's full record.

Scale by $b$, putting $L=T/b$ and $\lambda=\alpha b$.
The gap lengths have density $f_\lambda$, independently
marked by their conditional pasts. With $Y=x/b$ and
$W=A_I/b$, \eqref{eq:integrated-capture-local-bound} gives
$W\le CY^2$. Thus every moment density
$f_\lambda(y)\EE[W^j\mid Y=y]$ has a uniform exponential
envelope. Let $a=\EE W/\EE Y$, which is uniformly bounded,
and center each mark as $V=W-aY$.

The cyclic marked-moment argument in
\eqref{eq:marked-cyclic-moments} and
Lemma~\ref{lem:parent-band-higher-moments} applies with
$f_\lambda$ and the uniform finite-intensity renewal excess.
For clarity, a product containing $s$ singleton mark factors
has its first $s$ convolution moments zero. It therefore
cancels the corresponding highest powers of the polynomial
renewal kernel. In a $2p$th moment the surviving power,
including the cyclic factor $L$, is at most $p$.
The exponential moment envelopes bound every remaining term
by $C_pL^p$. The signed parity terms and atomic renewal
remainders are exponentially small. This proves
\[
 \EE\left|\sum_I W_I-aL\right|^{2p}\le C_pL^p
\]
on the nonempty sectors; its first-moment error is exponentially
small as well. On the empty sector the band cost integrated
over all origins is at most $CL^2$ in scaled units, because
the two active counts at the band's lower endpoint are at
most $2L+1$. Its probability is $Ce^{-cL}$, so it has the
same moment bound. Multiplication by $b$ and centering at
the actual mean prove
\eqref{eq:integrated-capture-band-moments}.

For the microscopic strip, $\lambda=\alpha b_\alpha\in[1,2)$.
Let $N_I$ count the original points assigned to one gap,
including its left endpoint. The full single-insertion cost
inside that gap is at most $C(N_I+1)$, so its integrated
mark in scaled units is at most $CY(N_I+1)$.
Formula~\eqref{eq:conditional-gap-count-mgf} gives
\[
 f_\lambda(y)\EE[e^{uN_I}\mid Y=y]
      =e^{\lambda y(e^u-1)}f_{\lambda e^u}(y).
\]
For sufficiently small fixed $u>0$, the uniform gap envelope
bounds the right side by $Ce^{-cy}$, uniformly for
$\lambda\in[1,2]$. All microscopic mark moment densities
therefore have uniform exponential envelopes too, and the
same cyclic calculation applies. The empty sector is handled
by H\"older: the total original count is Poisson with mean
at most $2L$, the full integrated cost is at most $CL(N+1)$,
and the sector has probability $Ce^{-cL}$.

Finally partition $[0,1]$ into this microscopic strip and
the successive dyadic bands above it. Minkowski for the
centered sums gives
\[
 \|A_{\alpha,T}-\EE A_{\alpha,T}\|_{2p}
       \le C_p\sqrt T\left(\sqrt{b_\alpha}
                      +\sum_{b\in\{2b_\alpha,\ldots,1\}}\sqrt b\right)
       \le C_p\sqrt T.
\]
No independence between the different bands is used.
\end{proof}

\begin{corollary}[The Poisson-background conditional mean is negligible]
\label{cor:Poisson-capture-background-mean}
Take exactly $q$ independent uniform probes on a circle of
circumference $q\ge8$, and the Poisson background of intensity
$\alpha\ge1$. Put $S=\sum_{i=1}^qG_B(Z_i)$ and
$\mu_B=q^{-1}\int_0^qG_B(z)\,dz$. Then
\[
 \operatorname{Var}(q\mu_B)\le Cq,\qquad
 \|q(\mu_B-\EE\mu_B)\|_{2p}\le C_p\sqrt q.
\]
In particular, as $\alpha\to\infty$, uniformly in $q\ge8$,
\[
 \frac{q(\mu_B-\EE\mu_B)}{\sqrt{q\log\alpha}}
                   \longrightarrow0\quad\hbox{in }L^2.
\]
Thus the background component in
$S-\EE S=(S-q\mu_B)+q(\mu_B-\EE\mu_B)$ is negligible
at the conjectured sparse-history scale in this Poisson model.
A conditional variance limit for $S-q\mu_B$, and conditioning
the original count to be exact, are not supplied by this result.
\end{corollary}
\begin{proof}
Apply Theorem~\ref{thm:Poisson-integrated-capture} with $T=q$.
The probes are iid given $B$, so $\EE(S\mid B)=q\mu_B$.
The displayed decomposition and estimates follow.
\end{proof}

\begin{lemma}[Centered single-probe moments at finite Poisson intensity]
\label{lem:Poisson-point-capture}
Let $G_\alpha(0)$ be the full single-probe cost through
separation one on a stationary Poisson line of intensity
$\alpha\ge1$. Put
\[
 \Lambda_\alpha=1+\log\alpha,\qquad
 m_\alpha=\EE G_\alpha(0),\qquad
 s_\alpha^2=\operatorname{Var}G_\alpha(0).
\]
For every fixed $p\ge2$,
\begin{equation}\label{eq:Poisson-point-capture-moments}
 m_\alpha\le C\Lambda_\alpha,\qquad
 \|G_\alpha(0)-m_\alpha\|_p\le C_p\sqrt{\Lambda_\alpha},
 \qquad
 \frac{s_\alpha^2}{\log\alpha}\longrightarrow4v_{\rm H}.
\end{equation}
The centered moment bound also holds for $G_B(Z)$ on the
Poisson circle of circumference $T\ge8$, uniformly in $T$,
with either original parity. Moreover
\begin{equation}\label{eq:Poisson-point-capture-mean-bias}
 |\EE G_B(Z)-m_\alpha|\le Ce^{-cT}.
\end{equation}
\end{lemma}
\begin{proof}
The bounded-intensity range follows from the microscopic
strip moment argument above and finitely many bands.
For large $\alpha$ put $w=A_p\log\Lambda_\alpha$,
$r_-=e^w/\alpha$, $r_+=e^{-w}$, and
$a=e^{-w/4}r_-$. The Poisson version of
Lemma~\ref{lem:growing-band-coupling} needs no original-count
comparison. Its finite-intensity skeleton error on a source
root of length at most $wr_+$ is bounded by
\[
 C(wr_+/a)e^{-2\alpha a}
   =Cw\alpha e^{-7w/4}\exp(-2e^{3w/4}).
\]
Choose $A_p$ sufficiently large. The root-length tail, inverse
terminal-cutoff estimate, and origin-completion budget from
that lemma then couple the entire central band to the line
entrance capture process except with probability $Ce^{-cw}$.
On a circle, the complete source-root density has its usual
outside renewal factor. Its difference from the line factor
is exponentially small in $T/r_+$ on these roots, uniformly
in $\alpha,T$. The same coupling therefore holds there.

The omitted strips have $L^p$ norm at most $C_pw$, by the
uniform single-band moments and the microscopic initial strip.
The ancestry tail, integrated against the exponentially
bounded source-root law, bounds the full cost's $L^{2p}$ norm
by $C_p\Lambda_\alpha$, so increasing $A_p$ makes the coupling
failure contribute at most $C_p$ in $L^p$.
In the central band the total single-probe cost is twice
the missing-background-edge count plus an error of absolute
value at most one. The stationary line cumulant bounds and
CLT therefore give centered $L^p$ norm at most
$C_p\sqrt{\log\alpha}$ and variance asymptotic
$4v_{\rm H}\log\alpha$. Since $w=o(\sqrt{\log\alpha})$,
the omitted strips give \eqref{eq:Poisson-point-capture-moments},
including its circle moment bound.

Finally the first-moment version of the cyclic calculation
in Theorem~\ref{thm:Poisson-integrated-capture} compares each
band's mean with its line reward density. In scaled
circumference $T/b$, its error is a polynomial in $T/b$
times $b e^{-cT/b}$, with uniform coefficients; the same
holds for the microscopic strip. The line reward densities
sum to $m_\alpha$. Summing the exponential errors over the
dyadic bands and dividing by $T$ proves
\eqref{eq:Poisson-point-capture-mean-bias}.
\end{proof}

\begin{proposition}[Conditional capture variance]
\label{prop:Poisson-conditional-capture-variance}
Let $\sigma_B^2=\operatorname{Var}(G_B(Z)\mid B)$ for a uniform
probe on the Poisson circle of circumference $T\ge8$.
There is an absolute $D<\infty$ such that, whenever
$T\ge D\log(\Lambda_\alpha+2)$,
\begin{equation}\label{eq:Poisson-conditional-capture-variance}
 \EE\left|\frac{\sigma_B^2-s_\alpha^2}{\Lambda_\alpha}\right|^2
                       \le C/T.
\end{equation}
All constants are uniform in $\alpha\ge1$ and original parity.
\end{proposition}
\begin{proof}
For a unit-cutoff line gap of length $Y$ attach the mark
\[
 W=\int_{\text{gap}}(G(z)-m_\alpha)^2\,dz,\qquad
 V=W-s_\alpha^2Y.
\]
Campbell's identity gives $\EE V=0$. For sufficiently small
$\theta>0$, Jensen, the origin Palm identity, and H\"older give
\[
 \EE[e^{\theta Y}W^2]
 \le \EE\left[e^{\theta Y}Y
                \int_{\text{gap}}|G(z)-m_\alpha|^4\,dz\right]
 \le C\bigl(\EE|G_\alpha(0)-m_\alpha|^8\bigr)^{1/2}
 \le C\Lambda_\alpha^2.
\]
The omitted H\"older factor is a fixed exponentially weighted
moment of the length-biased gap. The same bound holds for
$\EE[e^{\theta Y}V^2]$, and the weighted first absolute moment
is at most $C\Lambda_\alpha$.

We also need a pointwise density bound for the atomic renewal
terms. The shielded-root argument in
Lemma~\ref{lem:insertion-height} gives, for a uniform origin in
a gap of length $y$, an ancestry tail at most
$\min(1,\alpha y\rho_{\rm h}^{\,j})$. Consequently its $p$th
cost moment is at most $C_p[1+\log(\alpha y+2)]^p$.
Jensen and the gap envelope therefore bound the first and
second moment densities of $V$ respectively by
\[
 C\Lambda_\alpha^2(1+y)^C e^{-cy},
 \qquad
 C\Lambda_\alpha^4(1+y)^C e^{-cy}.
\]
Use the cyclic second-moment formula for $\sum V$.
The bounded ordinary renewal part costs $CT\Lambda_\alpha^2$
by the weighted integral bounds. The affine term cancels
against the two zero-mass first-moment factors. Signed terms
and the atomic terms at zero cost at most
$C\Lambda_\alpha^4e^{-cT}$, after absorbing polynomials in $T$.
For an empty background at cutoff one, condition on its
last-pair roots. The same uniform ancestry tail bounds its
integrated squared cost moment by a polynomial in $T$ and
$\Lambda_\alpha$; its probability is $Ce^{-cT}$. This yields
\begin{equation}\label{eq:integrated-squared-capture-bound}
 \EE\left|
    \int_0^T(G_B(z)-m_\alpha)^2\,dz-s_\alpha^2T
       \right|^2
 \le CT\Lambda_\alpha^2+C\Lambda_\alpha^4e^{-cT}.
\end{equation}
The zero-probe history on an initially empty input is included
in this exceptional-sector estimate.

Write $\mu_B=T^{-1}\int G_B(z)\,dz$. Since
\[
 \sigma_B^2=T^{-1}\int(G_B(z)-m_\alpha)^2\,dz
                               -(\mu_B-m_\alpha)^2,
\]
Theorem~\ref{thm:Poisson-integrated-capture} with $p=2$ and
\eqref{eq:Poisson-point-capture-mean-bias} bound the $L^2$ norm
of the last term by $C/T+Ce^{-cT}$. Divide
\eqref{eq:integrated-squared-capture-bound} by
$T^2\Lambda_\alpha^2$. The exponential error is at most
$C/T$ if $D$ is sufficiently large. This proves
\eqref{eq:Poisson-conditional-capture-variance}.
\end{proof}

\begin{theorem}[The early-history CLT with a Poisson background]
\label{thm:Poisson-early-history-CLT}
Let $\alpha\to\infty$ and $q\to\infty$, with no further
relation between them. On a circle of circumference $q$,
take a Poisson background of intensity $\alpha$, optionally
conditioned on either original parity, and two independent
uniform probe families of fixed counts $d,a$, $d+a=q$.
Let $H^{\rm early}$ be their matching-history distance through
separation one. Uniformly over the split $d+a=q$,
\begin{equation}\label{eq:Poisson-early-history-CLT}
 \frac{H^{\rm early}-\EE H^{\rm early}}
          {\sqrt{4qv_{\rm H}\log\alpha}}
       \Longrightarrow\mathcal N(0,1),\qquad
 \frac{\operatorname{Var}H^{\rm early}}
                {q\log\alpha}\longrightarrow4v_{\rm H}.
\end{equation}
Equivalently, on the unit circle the expected background count
is $M=\alpha q$ and the separation cutoff is $1/q$.
The original count is not fixed. The later history is not
included in this theorem.
\end{theorem}
\begin{proof}
Put $S=\sum_{i=1}^qG_B(Z_i)$ and
$\mu_B=\EE(G_B(Z)\mid B)$. First suppose
$q\ge D\log(\Lambda_\alpha+2)$. The preceding proposition
gives $\sigma_B^2/\log\alpha\to4v_{\rm H}$ in probability
and in $L^1$. Given $B$, the summands of $S-q\mu_B$ are iid.
Their conditional Lindeberg expression tends to zero in
probability: its expectation is at most
\[
 \frac{C}{q\Lambda_\alpha^2}
       \EE|G_B(Z)-\mu_B|^4\le C/q,
\]
by Lemma~\ref{lem:Poisson-point-capture} and Jensen.
The conditional characteristic-function proof of the
Lindeberg theorem gives a Gaussian limit for
$(S-q\mu_B)/\sqrt{q\log\alpha}$. Corollary~\ref{cor:Poisson-capture-background-mean}
removes $q(\mu_B-\EE\mu_B)$ in $L^2$. Conditional variance
and that corollary also give the stated variance asymptotic
for $S$.

For completeness consider $q<D\log(\Lambda_\alpha+2)$.
The Poisson version of the growing-family root comparison
in Lemma~\ref{lem:logarithmic-probe-comparison} omits its
original-count Stirling ratio and has the same or smaller
error bounds. Rescale to the unit circle with mean count
$M=\alpha q$, and take $w=A\log\log M$ as there. It compares
the sum of the $q$ full single-probe histories to independent
line captures, with $L^2$ error $Cq\log\log M$.
The cutoff $1/q$ discards at most $Cq(1+\log(q+2))$ in
$L^2$: use the uniform single-band moments down to at most
eight active points, followed by the deterministic terminal
bound. Both errors are $o(\sqrt{q\log\alpha})$ in this range,
and $\log M/\log\alpha\to1$. The independent line-capture
CLT and its moment bounds give the same centered limit and
variance asymptotic for $S$.

Finally $H^{\rm early}=S+R$ exactly. The Poisson-background,
fixed-probe part of
Theorem~\ref{thm:fixed-probe-remainder-variance} gives
$\operatorname{Var}R\le Cq$, uniformly over the colour split.
Its centered contribution is negligible in $L^2$ relative
to $\sqrt{q\log\alpha}$. Cauchy--Schwarz also makes its
covariance with $S$ negligible at variance scale $q\log\alpha$.
This proves both assertions.
\end{proof}

\subsection{De-Poissonizing complete relative histories}
\label{sec:relative-history-depoisson}

\status{Working proof: exact-count early-history CLT; the late comparison follows below}

\begin{lemma}[An original-count reservoir outside all probe source gaps]
\label{lem:relative-history-depoisson}
Let $q=q_M\ge1$, $q=o(M)$, and choose $r=r_M$ so that
$r\to0$, $Mr\to\infty$, and $qr\to0$.
On the unit circle, take $q$ independent uniform probes.
The following data have total-variation distance tending to
zero between an exact size-$M$ background and a Poisson
background of mean $M$ conditioned to have parity $M$:
the probe origins, the background configuration at $r$, and
every original point inside a cutoff gap containing a probe.
On the empty or singleton cutoff sector reveal the whole
original input as well.

Consequently the same total-variation comparison holds for
the complete spatial relative matching histories of every
pair of probe subsets, simultaneously at all separations.
This assertion concerns differences of matching histories,
not the individual background history or its total input count.

There is an event $\mathcal G$, measurable from the displayed
data, with Poisson probability
$\PP(\mathcal G^c)\le Ce^{-c/r}$, on which the density of the
exact-count data relative to the matching-parity Poisson data
is at most a uniform constant.
\end{lemma}
\begin{proof}
Work first under an unconditioned Poisson input of mean $M$,
and denote its original count by $N$. Write $\mathcal C$ for
the cutoff configuration, with gaps $x_i$, and $O$ for the
indices of gaps containing at least one probe.
Let $\mathcal S$ be the sigma-field of the stated data.
On the nonspecial sectors, the unexamined gap counts remain
independent conditional on $\mathcal S$, with the same
conditional laws as given $\mathcal C$ alone. Probe occupancy
depends only on geometry and independent origins, and revealing
other gaps' full records does not change these conditional laws.

The expected revealed length tends to zero. Indeed, conditional
on $\mathcal C$,
\[
 \EE\left[\sum_{i\in O}x_i\mid\mathcal C\right]
 =\sum_i x_i[1-(1-x_i)^q]\le q\sum_i x_i^2.
\]
The cyclic gap-square bound therefore gives expected length
at most $Cqr=o(1)$, outside the exponentially unlikely
empty/singleton sectors. Put
$\mu_{\mathcal S}=\EE(N\mid\mathcal S)$ and
$s_{\mathcal S}^2=\operatorname{Var}(N\mid\mathcal S)$.
The count cumulant estimates in
Lemma~\ref{lem:mesoscopic-depoisson-TV}, summed only over
unexamined gaps, show
\[
 s_{\mathcal S}^2/M\longrightarrow1
             \quad\hbox{in probability and in }L^1.
\]
For the $L^1$ assertion, the expected missing leading variance
is at most $CMqr+o(M)$; the summed cumulant errors are
$Cr^{-1}e^{-cMr}=o(M)$. Total variance now gives
\[
 \EE(\mu_{\mathcal S}-M)^2=M-\EE s_{\mathcal S}^2=o(M).
\]
The conditional fourth-moment Lindeberg sum can only decrease
when the examined gaps are omitted. Its expectation divided
by $M^2$ is still $O(r+M^{-1})+o(1)$ by the proof of
\eqref{eq:conditional-count-Lindeberg}. Thus the conditional
central limit estimate for $N$, centered at $M$, remains valid
given $\mathcal S$.

For Fourier tails, the short background gaps $x_i\in[r,2r]$
have total length at least some $\delta>0$ except on probability
$Ce^{-c/r}$. This follows from the negative gap-reward tilt
used in \eqref{eq:empty-gap-reservoir-tail}, now without
a probe-empty mark. At most $q$ distinct gaps are occupied;
the total occupied short-gap length is therefore at most
$2qr$ deterministically. For large $M$ at least $\delta/2$
of short-gap length remains unexamined on that same event.
Call it $\mathcal G$. It excludes the empty/singleton sectors.
The odd-Poisson characteristic-function bound applied to this
reservoir gives
\[
 \left|\EE(e^{itN}\mid\mathcal S)\right|
 \le \exp\{-cM\operatorname{dist}(t,\pi\mathbb Z)^2\}
                        \quad\hbox{on }\mathcal G.
\]
The already revealed count contributes only a phase.
The conditional CLT and Fourier inversion consequently give,
on the parity-$M$ data,
\[
 \sqrt M\,\PP(N=M\mid\mathcal S)
                   \longrightarrow 2/\sqrt{2\pi}
                              \quad\hbox{in probability},
\]
and the left side is uniformly bounded on $\mathcal G$.
Divide by the matching-parity Poisson point probability.
The resulting density converges to one in probability and
has mean one. Its $L^1$ distance from one therefore tends to
zero, proving the data comparison and the bounded-density claim.

It remains to check what the data determine. Every edge inside
a probe-empty cutoff gap is protected under every probe subset,
at its original separation, so its contribution to every
relative matching vector is zero below $r$.
Occupied-run localization determines all other early differences
from the revealed records and endpoints, including adjacent-gap
cascades. It also determines every augmented configuration at
$r$. Their subsequent histories follow deterministically.
On the two special sectors the whole input was revealed.
Thus the entire family of relative-history paths is a
measurable function of the data. Total variation contracts
under this map.
\end{proof}

\begin{theorem}[The finite early-history central limit theorem]
\label{thm:finite-early-history-CLT}
Let $B_M$ consist of exactly $M$ iid uniform circle points.
Let two disjoint independent probe families have counts
$d,a$ with $q=d+a\to\infty$ and $q=o(M)$.
Write $H^{\rm early}_{M;d,a}$ for their spatial matching-history
distance through separation $1/q$. Then, uniformly over the split,
\begin{equation}\label{eq:finite-early-history-CLT}
 \frac{H^{\rm early}_{M;d,a}-\EE H^{\rm early}_{M;d,a}}
      {\sqrt{4qv_{\rm H}\log(M/q)}}
           \Longrightarrow\mathcal N(0,1),\qquad
 \frac{\operatorname{Var}H^{\rm early}_{M;d,a}}
                     {q\log(M/q)}\longrightarrow4v_{\rm H}.
\end{equation}
The history after separation $1/q$ is handled separately in
Section~\ref{sec:late-history-variance}.
\end{theorem}
\begin{proof}
Put $\alpha=M/q$ and let $S$ be the sum of the $q$ single-probe
costs through $1/q$. First note the uniform Poisson bound
\begin{equation}\label{eq:Poisson-capture-sum-centered-moments}
 \EE_{\rm Pois}|S-\EE_{\rm Pois}S|^{2p}
            \le C_p\{q(1+\log\alpha)\}^p
                       \qquad(q\ge8,\ \alpha\ge1).
\end{equation}
After scaling to circumference $q$, the conditional summands
are iid. The conditional even-moment expansion gives a bound
by
\[
 C_p\{q^p(\sigma_B^2)^p
          +q\,\EE[|G_B(Z)-\mu_B|^{2p}\mid B]\}.
\]
Jensen and Lemma~\ref{lem:Poisson-point-capture} bound its
expectation by $C_pq^p(1+\log\alpha)^p$.
Theorem~\ref{thm:Poisson-integrated-capture} bounds the
remaining conditional-mean fluctuation in $L^{2p}$ by
$C_p\sqrt q$.
This proves \eqref{eq:Poisson-capture-sum-centered-moments}.

Apply Lemma~\ref{lem:relative-history-depoisson} with
$r=(Mq)^{-1/2}$. The statistic $S$ is measurable from its data,
so the centered limit for $S$ proved in
Theorem~\ref{thm:Poisson-early-history-CLT} transfers with the same
deterministic centering $\EE_{\rm Pois}S$.
We justify transferring means and variances separately.
On $\mathcal G$ the exact-count density is bounded, and
\eqref{eq:Poisson-capture-sum-centered-moments} applies.
On $\mathcal G^c$, conditioning on $N=M$ costs at most
$C\sqrt M$. The full Poisson ancestry moments bound the
$L^{4p}$ norm of $S-\EE_{\rm Pois}S$ by
$C_pq\log(M+2)$. H\"older therefore bounds the exceptional
$2p$th moment by
\[
 C_p\sqrt M\{q\log(M+2)\}^{2p}e^{-c\sqrt{Mq}/2}.
\]
This is negligible compared with
$\{q(1+\log(M/q))\}^p$, uniformly for $q\le M$.
Hence the normalized exact-count variables are uniformly
integrable at every fixed moment order. Their mean tends to
zero under the Poisson centering, and their second moment
has the Poisson limit. Replacing that centering by
$\EE S$ gives the centered CLT and variance asymptotic for $S$.

Finally $H^{\rm early}_{M;d,a}=S+R_{M;d,a}$.
Corollary~\ref{cor:all-sparse-remainder-variance} makes the
centered remainder negligible in $L^2$ on
$\sqrt{q\log(M/q)}$. The limit and variance asymptotic
therefore pass from $S$ to $H^{\rm early}_{M;d,a}$.
\end{proof}

\subsection{Spatial locality of late bands and the full sparse-history CLT}
\label{sec:late-history-variance}

\status{Working proof: the complete centered history CLT for every growing sublinear edit count}

The late history can be controlled without a quantitative decay
estimate for common vertices. The relevant locality is that of an
accepted edge's parent interval, separately in each marginal input.
Fixing the probe counts introduces a weak spatial dependence which
we retain explicitly.

\begin{lemma}[Correlation of disjoint fixed-count restrictions]
\label{lem:binomial-restriction-correlation}
Let $\mathcal X$ consist of any finite number of independent
uniform samples of fixed counts on a circle of circumference $T$,
and an independent Poisson process. The records retain their sample
colours but not the individual labels within a sample.
For disjoint deterministic arcs $A,B$ of lengths strictly between
zero and $T$, put $p=|A|/T$, $s=|B|/T$.
If $U,V$ are centered square-integrable functions of the restrictions
of $\mathcal X$ to $A,B$, respectively, then
\begin{equation}\label{eq:binomial-restriction-correlation}
 |\EE UV|\le
 \sqrt{\frac{ps}{(1-p)(1-s)}}\,\|U\|_2\|V\|_2.
\end{equation}
The right side is zero if every fixed count is zero.
\end{lemma}
\begin{proof}
For one size-$n$ sample, let $K_A,K_B$ be its two counts.
Given these counts, the two unlabelled location records are
independent, so conditional expectation reduces the problem
to the two counts. Define $Q_j(x;p)$ by
\[
 (1+(1-p)t)^x(1-pt)^{n-x}
       =\sum_{j=0}^n Q_j(x;p)t^j.
\]
Independence of the $n$ categorical trials gives
\[
 \EE Q_j(K_A;p)Q_\ell(K_A;p)
   =\ind_{\{j=\ell\}}\binom nj[p(1-p)]^j,
\]
and, because the two categories are disjoint,
\[
 \EE Q_j(K_A;p)Q_\ell(K_B;s)
   =\ind_{\{j=\ell\}}\binom nj(-ps)^j.
\]
These polynomials are complete orthogonal bases for the count
spaces. After normalization the nonconstant correlations are
$(-\sqrt{ps/((1-p)(1-s))})^j$, $j\ge1$.
Their absolute values are at most the first one since $p+s\le1$.
This proves the assertion for one sample. For independent
samples, take tensor products of the bases: every nonconstant
term contains a factor bounded by the same first correlation.
The restrictions of the Poisson component are independent and
contribute no additional nonconstant correlation. This also
proves the degenerate zero-count assertion.
\end{proof}

\begin{lemma}[A uniform parent-interval localization bound]
\label{lem:late-band-locality}
On a circle of circumference $T\ge8$, let $B$ be a Poisson
background of intensity $\alpha\ge1$. Let $P,Q$ be independent
uniform samples with fixed counts $d,a\le\eta_*T$, independent of $B$,
where $\eta_*<\infty$ is fixed. Constants may depend on $\eta_*$.
For $1\le u\le T/8$, partition the circle into
$n=\lceil T/u\rceil$ equal cells $I_j$, of length $h=T/n$.
Assign an edge to the clockwise initial endpoint of its shorter
arc. Let $X_j$ count the edges in the spatial symmetric
difference of the two histories in $(u,2u]$ assigned to $I_j$.
Then $0\le X_j\le4$.

Let $W_{j,\ell}$ be the union of $I_j$ and the $\ell$ cells
on either side, and let $\mathcal F_{j,\ell}$ reveal the
three coloured input restrictions there. For
$0\le\ell\le n/16$,
\begin{equation}\label{eq:late-band-locality}
 \bigl\|X_j-\EE(X_j\mid\mathcal F_{j,\ell})\bigr\|_2
                  \le Ce^{-c\ell},
\end{equation}
uniformly in all the parameters and the split of the probes.
\end{lemma}
\begin{proof}
Each marginal active configuration at $u$ has pairwise
separation greater than $u$. A cell of length $h\le u$
therefore contains at most two active points, each used by
at most one later edge. This proves the bound of four.

We first check a uniform marginal domination. The uncoloured
input $B+P$, conditional on its total count, is an iid uniform
sample. Its count has law $k+\operatorname{Pois}(\lambda)$,
where $\lambda=\alpha T\ge8$ and $k=d\le\eta_*T\le\eta_*\lambda$.
Relative to a Poisson count of mean $\lambda+k$, its density is
\[
 r(n)=\frac{\PP(\operatorname{Pois}(\lambda)=n-k)}
             {\PP(\operatorname{Pois}(\lambda+k)=n)}
       \quad(n\ge k).
\]
For $k>0$ the successive ratio is
\[
 \frac{r(n+1)}{r(n)}
       =\frac{\lambda(n+1)}{(\lambda+k)(n+1-k)}.
\]
Thus the maximum is attained within one of $\lambda+k$.
The elementary Stirling bounds at these two Poisson means give
$\sup_n r(n)\le C\sqrt{(\lambda+k)/\lambda}\le C_{\eta_*}$.
For $k=0$ the density is one. The entire uncoloured input
law is consequently dominated by a uniform constant times
the Poisson law of intensity $\alpha+k/T$. The same holds
for $B+Q$.

At separation $2u$, the gap covering a deterministic location
in this Poisson comparison has an exponential tail on scale
$u$, uniformly in the intensity. Indeed its density is the
length-biased gap factor times the bounded outside renewal
density from Lemma~\ref{lem:finite-Poisson-spatial-CLT};
the gap envelope bounds the tail by $Ce^{-cL}$ at length $Lu$.
The atomic singleton term and the empty sector have probability
at most $Ce^{-cT/u}$ and obey the same bound whenever
$L\le T/(8u)$, after adjusting constants.
Marginal domination transfers these bounds to $B+P,B+Q$.

Make a second coloured input by independently resampling the
exterior of $W_{j,\ell}$ conditional on its interior record.
For each fixed-count sample the exterior count is the fixed
total minus the revealed count, and its exterior points are
independent uniform points there. Both complete inputs have
the original joint law.

If an accepted edge assigned to $I_j$ changes under this
replacement, its closed parent in a positive input must leave
$W_{j,\ell}$. This is the exterior-protection implication
proved in Lemma~\ref{lem:matching-rectangle}, applied to an
arbitrary exterior replacement. Every parent of an edge
accepted by $2u$ lies inside the closed gap containing that
edge at $2u$, unless the configuration has become empty.
A gap meeting $I_j$ and leaving $W_{j,\ell}$ covers one of
the two endpoints of $I_j$ and has length at least $\ell h$.
The preceding exponential tail and $h\ge u/2$ therefore
bound the probability of any such edge in the four marginal
inputs by $Ce^{-c\ell}$. Empty and singleton sectors are
included. The two scores differ by at most four. Finally
\[
 \EE\operatorname{Var}(X_j\mid\mathcal F_{j,\ell})
       =\tfrac12\EE(X_j-X'_j)^2\le Ce^{-c\ell}.
\]
Decreasing $c$ proves \eqref{eq:late-band-locality}.
No common protected endpoint in the two copies is required.
\end{proof}

\begin{theorem}[Linear late-history variance]
\label{thm:late-history-variance}
In the Poisson-background, fixed-probe model of
Lemma~\ref{lem:late-band-locality}, let $H^{>b}$ be the full
spatial history distance contributed by separations greater
than $b$. Uniformly for $1\le b\le T/2$,
\begin{equation}\label{eq:late-history-variance}
 \operatorname{Var}H^{>b}\le CT/b.
\end{equation}
The background may also be conditioned on either original
parity. In particular, at $T=q=d+a\ge8$ the variance after
separation one is at most $Cq$.

For exactly $M$ iid background points on the unit circle
and exactly $q=d+a$ independent probes, with $q\to\infty$
and $q=o(M)$, this implies
\begin{equation}\label{eq:finite-late-history-variance}
 \operatorname{Var}H^{>1/q}_{M;d,a}\le Cq
\end{equation}
for all sufficiently large $M$, with a constant independent
of the sequence and the split.
\end{theorem}
\begin{proof}
We first prove the band estimate
\begin{equation}\label{eq:late-band-variance}
 \operatorname{Var}H^{(u,2u]}\le CT/u
                    \qquad(1\le u\le T/8).
\end{equation}
Use the cells and scores of the preceding lemma. If $n<64$,
the deterministic bound $H^{(u,2u]}\le4n$ suffices.
Otherwise put $L=\lfloor n/16\rfloor$ and write
\[
 X_j-\EE X_j=\sum_{\ell=0}^L D_{j,\ell}+E_j,
\]
where $D_{j,0}=\EE(X_j\mid\mathcal F_{j,0})-\EE X_j$,
the later $D_{j,\ell}$ are successive conditional-expectation
differences, and $E_j=X_j-\EE(X_j\mid\mathcal F_{j,L})$.
All are centered and the locality bound gives
\[
 \|D_{j,\ell}\|_2\le Ce^{-c\ell},\qquad
 \|E_j\|_2\le Ce^{-cn}.
\]
All windows used have length at most $T/4$.
For fixed $j,\ell,m$, at most $C(\ell+m+1)$ cells $k$
have $W_{j,\ell}\cap W_{k,m}\ne\varnothing$.
For these pairs use Cauchy--Schwarz. For disjoint windows,
Lemma~\ref{lem:binomial-restriction-correlation} gives
\[
 |\EE D_{j,\ell}D_{k,m}|
 \le \frac{C\sqrt{(\ell+1)(m+1)}}{n}
                        e^{-c(\ell+m)}.
\]
Summing both types of pairs and both levels gives
\[
 \EE\left|\sum_{j=1}^n\sum_{\ell=0}^L D_{j,\ell}\right|^2
 \le Cn\sum_{\ell,m\ge0}
       \{\ell+m+1+\sqrt{(\ell+1)(m+1)}\}e^{-c(\ell+m)}
 \le Cn.
\]
Also $\|\sum_jE_j\|_2\le Cn e^{-cn}\le C\sqrt n$.
This proves \eqref{eq:late-band-variance}.

Partition $(b,T/2]$ into dyadic bands until their lower
endpoint is comparable to $T$, followed by the remaining
terminal interval. The terminal active counts, and hence
its edge cost, are bounded by an absolute constant.
Minkowski for the centered band sums gives
\[
 \|H^{>b}-\EE H^{>b}\|_2
       \le C\sqrt{T/b}\sum_{j\ge0}2^{-j/2}+C
       \le C\sqrt{T/b}.
\]
For $b>T/8$ the deterministic active-count bound gives the
same conclusion directly. The two original parities each
have probability bounded away from zero since $\alpha T\ge8$.
Conditioning the squared deviation from the unconditional
mean therefore preserves the variance bound.

For the fixed background, use
Lemma~\ref{lem:relative-history-depoisson} with
$r=(Mq)^{-1/2}$. Write $h$ for the mean of the late distance
in the matching-parity Poisson model. On its event
$\mathcal G$, the density bound and
\eqref{eq:late-history-variance}, scaled to circumference $q$,
give
\[
 \EE_M[(H^{>1/q}-h)^2\ind_{\mathcal G}]\le Cq.
\]
At separation $1/q$ there are at most $q+1$ active points
in each input, so $0\le H^{>1/q}\le q+1$ deterministically.
Global original-count conditioning costs at most $C\sqrt M$.
The exceptional contribution is consequently at most
\[
 C\sqrt M(q+1)^2e^{-c\sqrt{Mq}}=o(q).
\]
Variance is bounded by the second moment about any constant,
which proves \eqref{eq:finite-late-history-variance}.
This is a separate moment argument; it does not infer a
variance bound merely from total-variation convergence.
\end{proof}

\begin{theorem}[The complete sublinear-edit history CLT]
\label{thm:full-sparse-history-CLT}
Let $B_M$ contain exactly $M$ iid uniform circle points.
Add two independent uniform probe families of counts $d,a$,
and put $q=d+a$. For every $q\to\infty$ with $q=o(M)$,
the completed spatial matching distance $H_{M;d,a}$ satisfies,
uniformly over all splits,
\begin{equation}\label{eq:full-sparse-history-CLT}
 \frac{H_{M;d,a}-\EE H_{M;d,a}}
          {\sqrt{4qv_{\rm H}\log(M/q)}}
       \Longrightarrow\mathcal N(0,1),\qquad
 \frac{\operatorname{Var}H_{M;d,a}}
                 {q\log(M/q)}\longrightarrow4v_{\rm H}.
\end{equation}
Either input parity is allowed.
In particular, for $k\to\infty$, $k=o(N)$ coordinate
resamplings of an odd sample,
\begin{equation}\label{eq:full-sublinear-resampling-CLT}
 \frac{H_N(k)-\EE H_N(k)}
          {\sqrt{8kv_{\rm H}\log(N/k)}}
        \Longrightarrow\mathcal N(0,1),\qquad
 \frac{\operatorname{Var}H_N(k)}
                 {k\log(N/k)}\longrightarrow8v_{\rm H}.
\end{equation}
The same two assertions hold for the label-based history
$H_N^{\rm lab}(k)$, centered by its own expectation.
The positive constant $v_{\rm H}$ is the line capture
variance from Theorem~\ref{thm:line-history-CLT}; it remains
uncalibrated.
\end{theorem}
\begin{proof}
Decompose $H=H^{\rm early}+H^{>1/q}$.
Theorem~\ref{thm:finite-early-history-CLT} gives the centered
early CLT and variance. Theorem~\ref{thm:late-history-variance}
makes the centered late term $o_{L^2}(\sqrt{q\log(M/q)})$.
Cauchy--Schwarz controls their covariance, proving
\eqref{eq:full-sparse-history-CLT}.
For resampling put $M=N-k$ and $q=2k$; then
$\log((N-k)/(2k))/\log(N/k)\to1$.

Finally put $D=H_N(k)-H_N^{\rm lab}(k)$.
Lemma~\ref{lem:history-label-correction} gives
$0\le D\le2k$ and $\EE D\le4$, hence
$\EE D^2\le8k$. Its centered contribution is negligible
in $L^2$ on the displayed scale, proving the label assertion.
The deterministic $2k$ bound alone would not suffice here.
\end{proof}

\subsection{Centered fluctuations at positive edit density}
\label{sec:dense-history-CLT}

\status{Working proof: a positive dense-history variance and a centered Gaussian law}

\begin{corollary}[Late variance with bounded edit density]
\label{cor:dense-late-variance}
Fix $K<\infty$. For an exact size-$M$ uniform background
and independent probe families of fixed counts $d,a\le KM$,
\begin{equation}\label{eq:dense-late-variance}
 \operatorname{Var}H^{>r}_{M;d,a}\le C_K/r
                    \qquad(M^{-1}\le r\le1/2).
\end{equation}
No assumption that $(d+a)/M$ tends to zero is needed.
\end{corollary}
\begin{proof}
Start on circumference $T=M$ with background intensity one
and the stated fixed probe counts. Theorem~\ref{thm:late-history-variance}
gives variance at most $C_KM/b$ after $b\ge1$.
We transfer this bound by an original-count reservoir at
separation one, without claiming total-variation convergence.

Reveal the background cutoff configuration, all probes, and
complete original records only in occupied cutoff gaps; on
the empty or singleton sector reveal everything.
These data determine the complete relative history, by the
measurability argument in Lemma~\ref{lem:relative-history-depoisson}.
Let $X$ be the total length of probe-empty background gaps
whose lengths lie in $[1,2]$.
For Poisson probes of total intensity at most $2K$,
\eqref{eq:empty-gap-reservoir-tail} gives
$\PP(X<\delta_K M)\le C_Ke^{-c_KM}$.
Conditioning both probe counts to their integer means costs
at most $C\sqrt{(d+1)(a+1)}\le C_K(M+1)$.
After decreasing $c_K$, the same exponential bound holds
for the exact probe counts.

On $X\ge\delta_KM$, the unrevealed original counts in these
short empty gaps are independent Poisson variables of means
in $[1,2]$, conditioned to be odd. Their atoms at one and
three have uniformly positive probabilities. Hence the
modulus of each characteristic function is at most
$\exp\{-c\operatorname{dist}(t,\pi\mathbb Z)^2\}$.
There are at least $\delta_KM/2$ such gaps, so their sum has
maximal point probability $C_K/\sqrt M$.
The original Poisson point probability at $M$ is at least
$c/\sqrt M$. The density of the conditioned revealed data
is therefore bounded on this event.

Let $h_b$ be the Poisson-model mean of $H^{>b}$.
For the good event $\mathcal G=\{X\ge\delta_KM\}$,
$\EE_M[(H^{>b}-h_b)^2\ind_{\mathcal G}]\le C_KM/b$.
Deterministically $H^{>b}\le M/b+1$ for $b\le M/2$.
Global original-count conditioning and the exceptional
probability bound the remaining contribution by
$C_K\sqrt M(M+1)^2e^{-c_KM}$, which is absorbed into $C_KM/b$.
Set $b=Mr$ and use variance minimization about the true mean.
Finitely many small $M$ are covered by the deterministic
total input-count bound, after increasing $C_K$.
\end{proof}

\begin{lemma}[A finite-cutoff Gaussian law with exact colour counts]
\label{lem:dense-cutoff-conditional-CLT}
On a circle of circumference $T\to\infty$, let the three
independent uniform samples $B,P,Q$ have exact counts
$n_i=\lambda_iT$, $i=0,1,2$, where the $\lambda_i$ stay in a
compact subset of $(0,\infty)$ and tend to positive limits.
Fix $s<\infty$. The spatial history distance $H^{\le s}$
between $B+P$ and $B+Q$, centered by its exact-count mean,
has a Gaussian limit after division by $\sqrt T$.
Its variance divided by $T$ tends to
\begin{equation}\label{eq:dense-cutoff-variance}
 w_s=v_s-c_s^{\mathsf T}D^{-1}c_s,\qquad
 D=\operatorname{diag}(\lambda_0,\lambda_1,\lambda_2).
\end{equation}
Here $v_s$ and $c_s$ are respectively the history variance
and its three count covariances per unit length in the
unconditioned three-colour Poisson model.
The convergence includes every fixed centered moment.
These limiting quantities are continuous in the positive
intensity vector.
\end{lemma}
\begin{proof}
First use independent Poisson inputs of intensities $\lambda_i$.
A union gap longer than $s$ separates the two histories
through $s$. Successive components, each followed by its
first such gap, form iid blocks. Their lengths $L$, colour
counts $C=(C_0,C_1,C_2)$, and history reward $R$ have
exponential moments, uniformly for the intensities under
consideration. Indeed the number of shorter exponential
gaps before the first gap exceeding $s$ is geometric,
and $R\le C_0+C_1+C_2$.
Write
\[
 Z=(R,C_0,C_1,C_2),\qquad m=\EE Z/\EE L.
\]
The renewal-reward covariance is
\[
 \Sigma_s=\frac{\EE[(Z-mL)(Z-mL)^{\mathsf T}]}{\EE L}.
\]
The iid block CLT gives the joint centered limit of
$(H^{\le s},N_0,N_1,N_2)/\sqrt T$, with this covariance.
The seam component has every fixed moment bounded independently
of $T$, so closing the interval into a circle has no effect.
The same centered block martingale has $2p$th moment at most
$C_{p,s}T^p$: apply the even-moment martingale inequality up to
the first block crossing $T$, whose count is bounded by the
union point count plus one. Boundary residuals have uniform
moments by their length-biased exponential tails.
Thus joint moments also converge. The count block of
$\Sigma_s$ is exactly $D$, since the Poisson colour counts
are independent. Exponential domination of the block records
also proves continuity of all these quantities in the intensities.

We justify conditioning on all three counts by Fourier inversion.
Place disjoint cells around the circle, each with an interior
of length one and two empty-buffer candidates of length $s+1$.
There are at least $c_sT$ such cells. Reveal every buffer
and all input outside the cell interiors.
Whenever both buffers of a cell are empty, its interior is
isolated through $s$. Its history contribution and colour
counts are independent of the other such interiors.
The interior events \emph{zero union points} and \emph{one point
of colour $i$} have probabilities bounded below uniformly in
the intensities. On all four events the history reward is zero.
Pairing the zero event with each single-point event in the
squared modulus of the joint characteristic function gives
\[
 \left|\EE e^{itR_{\rm int}+i z\cdot C_{\rm int}}\right|
       \le \exp\left\{-c_s
                    \sum_{i=0}^2\operatorname{dist}(z_i,2\pi\mathbb Z)^2\right\},
\]
uniformly in real $t$. The empty-buffer indicators are
independent with uniformly positive probability. Multiplying
over their interiors and then averaging the buffers yields
\begin{equation}\label{eq:dense-count-Fourier}
 \left|\EE e^{itH^{\le s}+i z\cdot N}\right|
       \le C_s\exp\left\{-c_sT
                 \sum_{i=0}^2\operatorname{dist}(z_i,2\pi\mathbb Z)^2\right\}.
\end{equation}
The remaining bounded arc, if any, is included with the exterior.

In the Fourier formula for $N=(n_0,n_1,n_2)$, set
$z=h/\sqrt T$. The joint CLT supplies pointwise convergence
of the centered characteristic function, and
\eqref{eq:dense-count-Fourier} supplies the integrable
bound $C_se^{-c_s|h|^2}$. Gaussian integration, divided by
the corresponding Poisson count point probability, gives
\[
 \frac{H^{\le s}-\EE_{\rm Pois}H^{\le s}}{\sqrt T}
       \ \Big|\ \{N=n\}
       \ \Longrightarrow\ \mathcal N(0,w_s),
\]
with the Schur complement in \eqref{eq:dense-cutoff-variance}.
There is no unexamined lattice-period correction: the four
interior count atoms are $0,e_0,e_1,e_2$.

For moment transfer, retain only empty-buffer interiors whose
union count is at most one, and reveal this flag rather than
their count or position. Reveal all remaining records.
The truncated history is measurable from these data.
Except on probability $C_se^{-c_sT}$ there are at least
$\delta_sT$ retained interiors. Conditional on the data,
their count vectors are independent, supported on
$0,e_0,e_1,e_2$, with each probability bounded below.
Their total has maximal point probability $C_sT^{-3/2}$
by the same characteristic-function bound. The unconditional
target probability is at least $cT^{-3/2}$.
Thus the exact-count data density is bounded on this good
event. The Poisson centered $2p$th moment bound transfers
there. On its complement, global conditioning costs
$CT^{3/2}$, and the rough bound
$H^{\le s}\le N_0+N_1+N_2$ and H\"older make every normalized
fixed moment negligible. Uniform integrability now transfers
all moments, in particular the centering and variance, to
the exact-count mean.
\end{proof}

\begin{theorem}[The dense-edit history central limit theorem]
\label{thm:dense-history-CLT}
There is a continuous function $v_{\rm D}:(0,1)\to(0,\infty)$
such that, whenever $N$ is odd and $k/N\to\theta\in(0,1)$,
\begin{equation}\label{eq:dense-history-CLT}
 \frac{H_N(k)-\EE H_N(k)}{\sqrt N}
       \Longrightarrow\mathcal N(0,v_{\rm D}(\theta)),\qquad
 \frac{\operatorname{Var}H_N(k)}N\longrightarrow v_{\rm D}(\theta).
\end{equation}
The same statements hold for $H_N^{\rm lab}(k)$ centered by
its own mean. One intrinsic description of the variance is
\[
 v_{\rm D}(\theta)=\lim_{s\to\infty}
 \left\{v_s-c_s^{\mathsf T}
       \operatorname{diag}(1-\theta,\theta,\theta)^{-1}c_s\right\},
\]
with the Poisson block covariances of
Lemma~\ref{lem:dense-cutoff-conditional-CLT}.
In particular this is a conditioned variance, rather than
the unconditioned Poisson reward variance.
At small density,
\begin{equation}\label{eq:dense-variance-small-density}
 v_{\rm D}(\theta)\sim8v_{\rm H}\theta\log(1/\theta)
                                      \qquad(\theta\downarrow0).
\end{equation}
No explicit numerical evaluation of $v_{\rm D}$ is asserted.
The endpoint regime $k/N\to1$ is treated separately in
Sections~\ref{sec:endpoint-Poisson}--\ref{sec:endpoint-Gaussian}.
\end{theorem}
\begin{proof}
Rescale the original circle to circumference $T=N$.
The common background has count $N-k$ and the two probe
families each have count $k$. Their intensities tend to
$(1-\theta,\theta,\theta)$.
For each fixed $s$, the preceding lemma gives the exact-count
CLT for the history up to separation $s$ in these scaled units.

Use Corollary~\ref{cor:dense-late-variance} with
$M=N-k$, $d=a=k$, and original separation $r=s/N$.
On any compact subinterval of $0<\theta<1$, for all sufficiently
large fixed $s$ and $N$,
\[
 \operatorname{Var}(H_N(k)-H_N^{\le s/N}(k))
                       \le C_\theta N/s.
\]
The constants are uniform on such compact subintervals.
The square roots of the finite-cutoff limiting variances
therefore form a Cauchy family, with tail error $C_\theta/\sqrt s$.
Let their squared limit be $v_{\rm D}(\theta)$.
The same estimate proves convergence of the full variances
and, by the characteristic-function inequality
$|\EE e^{itX}-\EE e^{itY}|\le |t|\|X-Y\|_2$ applied to
centered normalized variables, the full Gaussian limit.
Uniformity on compact intensity sets and continuity of the
finite-cutoff covariances prove continuity of $v_{\rm D}$.

Here is a direct positive lower bound after conditioning the
counts. Partition circumference $N$ into cells with a unit
interior and unit buffers on each side. Call a cell suitable
if both buffers contain no union point and its interior
contains exactly two points of each colour $B,P,Q$.
The expected number of suitable cells is at least
$c_\theta N$ for large $N$: their fixed-count multinomial
probabilities converge to strictly positive Poisson probabilities.
Condition on all cell colour counts, every buffer record,
and every nonsuitable interior record. The remaining interior
positions are independent uniform samples within their cells.

In a suitable cell each input has four interior points. All
four are matched internally before separation one, while
the empty buffers prevent interaction with the exterior.
The entire future exterior history is consequently independent
of these interior positions. The interior history distance
has strictly positive conditional variance. For example,
open neighbourhoods of the configurations
\[
 \begin{array}{c|ccc}
     &B&P&Q\\ \hline
  A&(0.10,0.20)&(0.40,0.41)&(0.70,0.72)\\
  B&(0.20,0.80)&(0.21,0.78)&(0.23,0.76)
 \end{array}
\]
give respectively two and four changed spatial edges.
All three colour counts agree between the configurations.
These two open events have fixed positive conditional probabilities.
Conditional variance of the sum over suitable cells is
therefore at least a positive constant times their number.
Total variance yields $\operatorname{Var}H_N(k)\ge c_\theta N$,
so $v_{\rm D}(\theta)>0$.

Lemma~\ref{lem:history-label-correction} gives a uniform second
moment bound for $H_N(k)-H_N^{\rm lab}(k)$. Its centered
contribution divided by $\sqrt N$ tends to zero in $L^2$,
proving the label assertion.

Finally suppose \eqref{eq:dense-variance-small-density} failed
along some $\theta_j\downarrow0$. For each $j$ choose an odd
$N_j$ sufficiently large that $k_j=\lfloor\theta_jN_j\rfloor
\to\infty$ and the dense variance limit at that fixed density
approximates $v_{\rm D}(\theta_j)$ with error at most
$\theta_j\log(1/\theta_j)/j$. Since $k_j=o(N_j)$,
Theorem~\ref{thm:full-sparse-history-CLT} gives the conflicting
variance ratio $8v_{\rm H}$. This proves the small-density asymptotic.
\end{proof}

\subsection{Poisson shared edges near complete refresh}
\label{sec:endpoint-Poisson}

\status{Working proof: a rare-common-pair Poisson limit}

This is a Poisson limit in the refresh-density endpoint.
It is distinct from the physical last-collision law in
Section~\ref{sec:measure-process}.

\begin{theorem}[Poisson matching overlap at the square-root endpoint]
\label{thm:endpoint-Poisson}
Let $N$ be odd, let $k=N-M$, and suppose
\[
 M^2/N\longrightarrow\lambda\in[0,\infty).
\]
Let $J_N$ be the number of shared spatial edges in the completed
matchings after these $k$ coordinate resamplings. Then
\begin{equation}\label{eq:endpoint-Poisson}
 J_N\Longrightarrow\operatorname{Pois}(\kappa\lambda),
 \qquad
 H_N(k)=N-1-2J_N,
 \qquad
 \kappa=\pi_2/2=\int_0^\infty p(t)^2\,dt,
\end{equation}
where
$p(t)=-q'(t)/2=q(t)(1-e^{-2t})/(2t)$ as in
Proposition~\ref{prop:matching-quadratic-endpoint}.
All fixed moments of $J_N$ converge. The existing certificate
therefore gives
\[
 0.194237478218\le\kappa\le0.194237478220.
\]
More precisely, the shared edges, marked by their clockwise
initial endpoint on the unit circle and by $N$ times their
shorter-arc length, converge to a Poisson random measure
with intensity
\[
 \lambda\,dx\,p(t)^2\,dt,\qquad x\in\T,\quad t>0.
\]

For label-based histories the limit is different:
\begin{equation}\label{eq:endpoint-label-Poisson}
 \frac{N-1-H_N^{\rm lab}(k)}2
       \Longrightarrow\operatorname{Pois}(\kappa\lambda+1/2).
\end{equation}
The additional mean-$1/2$ variable is asymptotically independent
of the spatial shared-edge process. These assertions include
$M^2/N\to0$, when there are no spatial common edges with
probability tending to one.
\end{theorem}
\begin{proof}
Every shared spatial edge has two endpoints among the $M$
unchanged locations. Fix $s<\infty$ and retain only edges of
length at most $s/N$, writing their count as $J_{N,s}$.
For a factorial moment of order $\ell$, only $\ell$ vertex-disjoint
common pairs can contribute, since each matching has degree at
most one. There are $(M)_{2\ell}/2^\ell$ ordered choices of
unordered pairs, where $(m)_j=m(m-1)\cdots(m-j+1)$.
For each pair, integrating its shorter orientation and its
length $t/N\le s/N$ contributes $2\,dt/N$.

We verify the limit of the conditional acceptance probabilities.
Keep all pair anchors farther than $(2R+2s)/N$ from each other,
and require no other common point in their windows of radius
$(R+s)/N$. For fixed $\ell,s,R$, the excluded anchor volume
is $O_{\ell,s}(R/N)$, and the probability of an extra common
point there is $O_{\ell,s}(MR/N)$; both tend to zero.
On these separated windows the restrictions of the two size-$k$
probe samples converge jointly to independent rate-one Poisson
processes. An elementary multinomial-to-Poisson comparison
gives error $O_{\ell,s}(R^2/N+MR/N)$ for fixed $R$;
conditional positions within the windows are independent uniform.

The matching decision through $s$ is determined inside a
window whenever the probe union has a gap longer than $s$
on each side of the forced pair. Under the limiting rate-two
union input, failure has probability at most $C_se^{-c_sR}$:
test disjoint intervals of length $s+1$ on each side.
Thus, first sending $N\to\infty$ and then $R\to\infty$, the
$\ell$ pair decisions become independent. For a forced pair
at distance $t$, the two input decisions have probability
$p(t)^2$. It follows that
\begin{equation}\label{eq:endpoint-factorial-moments}
 \EE(J_{N,s})_\ell
       \longrightarrow
       \left(\lambda\int_0^s p(t)^2\,dt\right)^\ell.
\end{equation}
The same calculation with bounded test functions of the
anchors and lengths gives the corresponding factorial
moment measures. The simple domination
\[
 \EE(J_{N,s})_\ell
       \le (M)_{2\ell}(s/N)^\ell
       \le (M^2s/N)^\ell
\]
justifies the exponential generating-function expansion.
Hence the truncated marked process has the asserted Poisson
limit on $\T\times(0,s]$.

To remove the cutoff, condition on the unlabelled first
matching. Its $M$ common labels form a uniform subset of
the $N$ vertices. At separation $s/N$, the deterministic
packing bound leaves at most $N/s+1$ active vertices.
Consequently
\[
 \EE(J_N-J_{N,s})
 \le \frac{(M)_2}{(N)_2}\,\frac{N/s+1}{2}
 \le C\frac{M^2}{Ns}
 \qquad(1\le s\le N/2).
\]
As $M^2/N$ stays bounded, this makes the tail uniformly
small when $s\to\infty$. Also $\int_s^\infty p(t)^2dt\to0$,
since $0\le p\le1$ and $\int p=1/2$.
This proves both the full marked-process limit and the count limit.
The two matchings each have $(N-1)/2$ edges, proving the
identity for $H_N(k)$.

For moment convergence, dominate $J_N$ by $X_N$, the number
of first-matching edges with two common endpoints. Its exact
factorial moments are
\[
 \EE(X_N)_\ell
    =\left(\frac{N-1}{2}\right)_\ell
                      \frac{(M)_{2\ell}}{(N)_{2\ell}}.
\]
When $2\ell\le M=o(N)$ these are at most
$(C M^2/N)^\ell$, and they vanish if $2\ell>M$.
Thus $J_N$ has uniformly bounded exponential moments, giving
all fixed moment convergence. The value and certificate for
$\kappa$ are exactly half the already established quadratic
matching coefficient; no new numerical approximation is used.

Finally condition on the unchanged labelled points and both
unlabelled refreshed sets. Use the two extra common-label
edge counts $J_1,J_2$ from Lemma~\ref{lem:history-label-correction}.
There are at most $M$ possible unchanged-to-refreshed
coincidences, so $\EE(J_1\mid\text{geometry})\le M/k\to0$.
If $a,b$ count the refreshed-to-refreshed spatial edges in
the two copies, then deterministically $a/k,b/k\to1/2$,
since only $M=o(N)$ vertices have unchanged labels.
For every fixed $\ell$,
\[
 \EE((J_2)_\ell\mid\text{geometry})
       =\frac{2^\ell(a)_\ell(b)_\ell}{(k)_{2\ell}}
       \longrightarrow(1/2)^\ell
\]
uniformly over the possible geometries. Each factor in the
product on the left is at most one, since $a,b\le\lfloor k/2\rfloor$.
Its conditional generating functions therefore converge uniformly
to those of a mean-$1/2$ Poisson variable. This gives independence
from the spatial marked limit, which is measurable from the
geometry. Also the factorial moments of $J_1$ are at most
$(M/k)^\ell$, so it vanishes in every fixed moment.
The number of common label edges is $J_N+J_1+J_2$.
Adding the independent Poisson limits proves
\eqref{eq:endpoint-label-Poisson}, with fixed moment convergence
as well.
\end{proof}

\subsection{The Gaussian side of the complete-refresh endpoint}
\label{sec:endpoint-Gaussian}

\status{Working proof: the transition from rare shared edges to Gaussian overlap}

\begin{lemma}[A tagged late-overlap variance bound]
\label{lem:tagged-shared-tail}
Let $M\to\infty$, $M=o(N)$, and put $\Lambda=M^2/N$.
In the exact resampling model with common count $M$ and
two private counts $N-M$, let $J_N^{>b}$ count shared spatial
edges of length greater than $b/N$. There is a universal
constant $C$ such that, for all sufficiently large $N$,
\begin{equation}\label{eq:tagged-shared-tail}
 \operatorname{Var}J_N^{>b}\le C\Lambda/b
                                  \qquad(1\le b\le N/2).
\end{equation}
\end{lemma}
\begin{proof}
First keep the common count $M$ exact on circumference $N$,
and take independent Poisson private inputs of intensity
$1-\delta$, where $\delta=M/N$.
Either uncoloured marginal has count
$M+\operatorname{Pois}(N-M)$, and its density relative to
$\operatorname{Pois}(N)$ is uniformly bounded for $\delta\le1/2$,
by the successive-ratio calculation in
Lemma~\ref{lem:late-band-locality}.
Given the uncoloured input and its count $n$, the $M$ common
points are uniform tags. An edge has two common tags with
probability $(M)_2/(n)_2\le C\delta^2$ when $n\ge N/2$.
The exceptional event $n<N/2$ has probability $Ce^{-cN}$.

Use the cell decomposition for a band $(u,2u]$, now with
$X_j$ counting shared edges rather than changed edges.
There are at most two candidate edges per cell.
Every change of $X_j$ under conditional exterior resampling
requires a positive shared edge whose parent in at least
one of the two input histories leaves the local window.
Uniform tags, marginal domination, and the exponential
parent tail consequently give
\[
 \EE X_j^2\le C\delta^2,\qquad
 \|X_j-\EE(X_j\mid\mathcal F_{j,\ell})\|_2
                \le C\delta e^{-c\ell}.
\]
The term $Ce^{-cN}$ is absorbed here because $M\ge2$ and
$\ell\le N/(16u)$. This tagging is done separately in each
positive marginal input; it does not require independence
of the two matching decisions.
The only fixed-count input is the common sample, so
Lemma~\ref{lem:binomial-restriction-correlation} still applies.
The covariance summation in
Theorem~\ref{thm:late-history-variance}, with both local
$L^2$ norms carrying the factor $\delta$, gives
\[
 \operatorname{Var}J_N^{(u,2u]}\le C\delta^2N/u.
\]
For a terminal band with $N/u$ bounded, the same uniform-tag
argument bounds the second moment of its shared-edge count
by $C\delta^2$: there are only a bounded number of candidate
edges. Dyadic Minkowski thus gives
$\operatorname{Var}J_N^{>b}\le C\Lambda/b$
in this Poisson-private model.

To fix both private counts, regard the common points as
$M$ probes in each of the two independent private backgrounds,
whose mean count is $k=N-M$. Apply
Lemma~\ref{lem:relative-history-depoisson} at
$r=(kM)^{-1/2}$ in unit-circle coordinates to each background.
Its relative-history data determine all matching edges
incident to a common point: those edges are absent from
the corresponding private-only input. Thus the two data
sets determine the shared-edge process.
Conditional on the common locations, the private inputs
and their original counts are independent, so the two
conditional density bounds multiply. Passing to the prescribed
private parities costs only a bounded factor; equivalently the
density bound relative to unconditioned Poisson data gains that
factor. On the intersection of the good events the exact-count
density is bounded.
The probability of its complement before conditioning is
at most $Ce^{-c\sqrt{kM}}$, and conditioning both private
counts costs at most $Ck$.
Since $0\le J_N^{>b}\le M/2$, its exceptional centered
second moment about the Poisson-private mean is at most
$CkM^2e^{-c\sqrt{kM}}=o(\Lambda/b)$, uniformly for $b\le N/2$.
The good-event bound and variance minimization prove
\eqref{eq:tagged-shared-tail}.
\end{proof}

\begin{lemma}[Rare common pairs at a fixed cutoff]
\label{lem:rare-common-cutoff-CLT}
Suppose $\delta=M/N\to0$ and $\Lambda=N\delta^2\to\infty$.
For every fixed $s>0$, the exact-count number $J_{N,s}$ of
shared edges of length at most $s/N$ satisfies
\begin{equation}\label{eq:rare-common-cutoff-CLT}
 \frac{J_{N,s}-\EE J_{N,s}}{\sqrt\Lambda}
       \Longrightarrow\mathcal N(0,\kappa_s),\qquad
 \frac{\operatorname{Var}J_{N,s}}{\Lambda}\longrightarrow\kappa_s,
 \qquad \kappa_s=\int_0^s p(t)^2\,dt.
\end{equation}
All fixed centered moments converge.
\end{lemma}
\begin{proof}
Start with three Poisson colours on circumference $N$ of
intensities $\delta,1-\delta,1-\delta$.
At cutoff $s$, use the union components separated by gaps
longer than $s$ as in Lemma~\ref{lem:dense-cutoff-conditional-CLT}.
Write $R$ for a block's shared-edge reward and $L$ for its
length. The block length and total point count have uniform
exponential moments for $\delta\le1/2$.
Given the uncoloured block, each point is common with
probability $\delta/(2-\delta)$ independently.
A positive reward needs two common points; two different
shared edges need four. Thus, for every fixed $j\ge1$ and
some $\eta_s>0$,
\[
 \EE[e^{\eta_s L}R^j]\le C_{j,s}\delta^2,\qquad
 \EE[e^{\eta_s L}R(R-1)]\le C_s\delta^4.
\]
By the forced-pair formula and finite-window convergence
in Proposition~\ref{prop:matching-quadratic-endpoint},
the shared-edge density is
\[
 h_s(\delta)=\EE R/\EE L
       =\delta^2\int_0^s p_\delta(t)\,dt
       =\delta^2\{\kappa_s+o(1)\},
\]
where $p_\delta(t)$ is the probability that both inputs
accept two forced common points at distance $t$.
It tends to $p(t)^2$ and is bounded by one.
For the centered reward $Y=R-h_s(\delta)L$,
\[
 \frac{\EE Y^2}{\EE L}
       =h_s(\delta)+O_s(\delta^4)
       =\delta^2\{\kappa_s+o(1)\}.
\]
Its covariance per unit length with each colour count is
$O_s(\delta^2)$, using the same coloured block moment bounds.
The count covariance matrix is
$\operatorname{diag}(\delta,1-\delta,1-\delta)$.

The triangular block CLT therefore gives independent Gaussian
coordinates in the limit for the centered vector with
normalizations $\sqrt\Lambda,\sqrt M,\sqrt N,\sqrt N$.
Here are uniform Lindeberg and moment bounds.
The fourth moment of $Y$ is $O_s(\delta^2)$, that of the
centered common-count block is $O_s(\delta)$, and the private
ones have bounded fourth moments. The sum of their normalized
fourth-moment errors is
$O_s(\Lambda^{-1}+M^{-1}+N^{-1})=o(1)$.
The block-length law has uniformly bounded exponential moments,
so replacing the number of blocks by its renewal stopping index
preserves the limit; boundary residuals are negligible.
The even-moment martingale bound also gives
\[
 \EE_{\rm Pois}|J_{N,s}-\EE_{\rm Pois}J_{N,s}|^{2j}
       \le C_{j,s}(\Lambda^j+\Lambda)
       \le C'_{j,s}\Lambda^j
\]
for large $N$. In that inequality the predictable variance
is $O(\delta^2)$ per block and each higher increment moment
is $O_{j,s}(\delta^2)$.

We still have to condition all three original counts.
The empty-buffer argument of
Lemma~\ref{lem:dense-cutoff-conditional-CLT} is uniform here,
except that the single-common-point atom has probability
comparable to $\delta$. The other three atoms retain
probabilities bounded away from zero. It gives
\[
 \left|\EE e^{itJ_{N,s}+i z\cdot C}\right|
 \le C_s\exp\{-c_sN[
      \delta\operatorname{dist}(z_0,2\pi\mathbb Z)^2
       +\textstyle\sum_{i=1}^2\operatorname{dist}(z_i,2\pi\mathbb Z)^2]\},
\]
uniformly in the history Fourier parameter $t$.
In Fourier inversion scale $z_0$ by $M^{-1/2}$ and
$z_1,z_2$ by $N^{-1/2}$. The joint Gaussian limit and
the integrable Gaussian bound give the same history limit
after conditioning $C=(M,N-M,N-M)$, because the normalized
count covariances with the history tend to zero.

For moments, the retained zero/single-point interiors still
number at least $c_sN$ outside an event of probability
$C_se^{-c_sN}$. Their independent categorical count vectors
give maximal point mass $C_s/(N\sqrt M)$ by the displayed
weighted Fourier bound. This matches the three target Poisson
point probabilities, so the density bound on the retained
data is uniform. The shared reward is measurable from those
data. Transfer its centered moment estimate on the good event.
On the bad event, global conditioning costs $CN\sqrt M$;
the rough Poisson common-count moments are $O_j(M^j)$,
and the exponential event bound absorbs that cost.
Uniform integrability transfers the exact mean and all
centered moments, proving \eqref{eq:rare-common-cutoff-CLT}.
\end{proof}

\begin{theorem}[Gaussian overlap above the endpoint threshold]
\label{thm:endpoint-Gaussian}
If $M=N-k=o(N)$ and $\Lambda=M^2/N\to\infty$, then
\begin{equation}\label{eq:endpoint-Gaussian}
 \EE J_N\sim\kappa\Lambda,\qquad
 \operatorname{Var}J_N\sim\kappa\Lambda,\qquad
 \frac{J_N-\EE J_N}{\sqrt{\kappa\Lambda}}
                       \Longrightarrow\mathcal N(0,1).
\end{equation}
Consequently both spatial and label histories, centered
by their respective exact means, have Gaussian limits with
variance asymptotic
\[
 4\kappa\,\frac{(N-k)^2}{N}
       =2\pi_2\,\frac{(N-k)^2}{N}.
\]
The dense variance from Theorem~\ref{thm:dense-history-CLT}
therefore also satisfies
\begin{equation}\label{eq:dense-variance-near-one}
 v_{\rm D}(\theta)\sim2\pi_2(1-\theta)^2
                                      \qquad(\theta\uparrow1).
\end{equation}
\end{theorem}
\begin{proof}
For a fixed $s$, Lemma~\ref{lem:rare-common-cutoff-CLT}
gives the centered CLT with variance $\kappa_s\Lambda$.

Lemma~\ref{lem:tagged-shared-tail} bounds the omitted centered
part by $C\sqrt{\Lambda/s}$ in $L^2$.
Letting $s\to\infty$ proves the complete CLT and variance
asymptotic, since $\kappa_s\to\kappa$.
The single-pair calculation in
\eqref{eq:endpoint-factorial-moments}, divided by $\Lambda$,
still gives $\EE J_{N,s}/\Lambda\to\kappa_s$ when
$\Lambda\to\infty$: its errors depend on $M/N$, not on the
number of pair choices after this normalization.
The uniform-tag tail mean bound in that proof is at most
$C\Lambda/s$, proving the full mean asymptotic.

Use $H_N(k)=N-1-2J_N$ and the bounded second moment of
the label correction to obtain both history assertions.
Finally choose a diagonal sequence of sample sizes at any
densities $\theta_j\uparrow1$ so that
$(1-\theta_j)^2N_j\to\infty$ and the dense variances approximate
$v_{\rm D}(\theta_j)$ with error $o((1-\theta_j)^2)$.
The just-proved endpoint variance gives
\eqref{eq:dense-variance-near-one}.
As in the sparse theorem, replacing the exact centering by
its leading asymptotic is not asserted.
\end{proof}

\subsection{A uniform finite-size bound for the history mean}
\label{sec:profile-centering}

\status{Working proof: profile centering for the dense Gaussian law}

\begin{lemma}[The Poisson circle mean differs from the line by a constant]
\label{lem:Poisson-circle-profile-mean}
On circumference $T\ge8$, take independent Poisson colours
of intensities $1-\theta,\theta,\theta$, $0\le\theta\le1$,
and let $H_T^{\rm P}$ be the completed spatial history distance.
Uniformly in $T,\theta$,
\begin{equation}\label{eq:Poisson-circle-profile-mean}
 |\EE H_T^{\rm P}-T\mathfrak h(\theta)|\le C.
\end{equation}
\end{lemma}
\begin{proof}
Both marginal inputs have intensity one, uniformly in $\theta$.
For a band $(u,2u]$, $1\le u\le T/8$, partition circumference
$T$ into cells of length at most $u$ as in
Lemma~\ref{lem:late-band-locality}.
Couple the circle and line input records to agree on a
window of length $T/3$ around one cell, filling the respective
exteriors independently. The geometry inside this window is
the same in the two spaces.
If the assigned edge score differs, some positive accepted
edge has a parent leaving that window.
The marginal parent tails on the line and circle bound this
probability by $Ce^{-cT/u}$.
Each score is at most four, so summing over the cells gives
a mean difference at most $C(T/u)e^{-cT/u}$.

For the initial strip $[0,1]$, use unit cells and the
cutoff-one parent bound. A cell score is bounded by twice
the initial union point count in that cell, whose fixed
moments are uniformly bounded since its intensity is at
most two. H\"older therefore gives strip error $CTe^{-cT}$
after decreasing $c$.
The empty and singleton circle sectors are included in the
parent tail bounds.
Sum the initial strip and successive dyadic bands up to a
cutoff comparable with $T$; the sum of the displayed errors
is uniformly bounded. The remaining circle edge cost is
deterministically bounded by active packing. The expected
line edge cost assigned to an interval of length $T$ is
at most $Tq(cT)\le C$ beyond that cutoff, using the marginal
active density. Finally the full line reward density is
$\mathfrak h(\theta)$ by Theorem~\ref{thm:dense-history-profile}.
\end{proof}

\begin{theorem}[A uniform mean bias bound]
\label{thm:uniform-history-mean-bias}
For every odd $N$ and $0\le k\le N$,
\begin{equation}\label{eq:uniform-history-mean-bias}
 \left|\EE H_N(k)-N\mathfrak h(k/N)\right|
       +\left|\EE H_N^{\rm lab}(k)-N\mathfrak h(k/N)\right|
                       \le C\log^2(N+2).
\end{equation}
The constant is universal but uncalibrated.
\end{theorem}
\begin{proof}
Rescale to circumference $N$, and put $\theta=k/N$.
Write $Q$ for the Poisson colour law with means
$(N-k,k,k)$ and $P$ for its conditioning to those three
exact counts. Its global density obeys
\[
 dP/dQ\le A_N,\qquad A_N\le C(N+1)^{3/2},
\]
by the Poisson point probabilities at their integer means;
a colour of mean zero costs no factor.

We also need a much sharper comparison on a small arc $W$.
Put $p=|W|/N$. Conditional on its count, the unlabelled
restriction of a size-$n$ uniform sample is the same as that
of a Poisson input of mean $n$. The count laws are
$\operatorname{Bin}(n,p)$ and $\operatorname{Pois}(np)$.
Theorem 1, p.~474, of
\href{https://doi.org/10.1017/S0305004100061806}
{Barbour--Hall, \emph{On the rate of Poisson convergence}
(1984)}, gives
\[
 d_{\rm TV}(\operatorname{Bin}(n,p),\operatorname{Pois}(np))
        \le p(1-e^{-np})\le p.
\]
This imported inequality was checked in the
\href{https://user.math.uzh.ch/barbour/pub/Barbour/BHall_PoissonRate.pdf}
{authors' scan}; the zero-count case is immediate.
The three-colour restriction laws therefore have distance
at most $3p$, uniformly in $k$.

Here is the local-score comparison under the two exterior laws.
For a cell of length at most $u$ in band $(u,2u]$, let $X$
be its score under $P$. Keep its coloured input inside the
window of $\ell$ neighbouring cells on either side, and fill
the exterior with fresh independent Poisson inputs; call
the new score $Y$. Its full input law has distance at most
$3p$ from $Q$, since only the interior law differs.
Both scores lie in $[0,4]$.
Exterior protection implies that a difference needs a
positive parent leaving the window in one of the two
coupled configurations.
Under $P$ the probability of this event is at most
$A_NCe^{-c\ell}$; under the filled law it is at most
$Ce^{-c\ell}+3p$. Comparing the filled score once more
with $Q$ consequently gives
\begin{equation}\label{eq:local-mean-count-comparison}
 |\EE_P X-\EE_Q X|\le C(A_Ne^{-c\ell}+p),
 \qquad p\le C(\ell+1)u/N.
\end{equation}
This comparison does not assume that the conditioned
exterior is Poisson.

Choose $\ell=\lceil A\log(N+2)\rceil$ with $A$ large.
For dyadic $u\ge1$ satisfying $u\le N/[32(\ell+1)]$,
the window occupies at most a quarter of the circle and
the localization estimates apply. Summing
\eqref{eq:local-mean-count-comparison} over at most $CN/u$
cells costs $C\log(N+2)$ per band; the exponential error
is negligible by the choice of $A$.
There are at most $C\log(N+2)$ bands.
After the last such band the active-packing bound leaves
at most $C\log(N+2)$ further edges in either model.

For $[0,1]$ use unit cells and the same windows.
The score is bounded by twice the initial union count
inside the cell. In both the exact-count and Poisson laws
this count has a uniformly bounded exponential moment:
each exact restriction is binomial and its exponential
moment is bounded by that of a Poisson variable with
the same mean, while the total mean is at most two.
The filled law has the same interior count as the exact law.
Cap all scores at $B\log(N+2)$. Their discarded expectations
are at most $C(N+2)^{-4}$ for sufficiently large $B$.
The capped version of the preceding coupling estimate costs
at most
\[
 C\log(N+2)\{A_Ne^{-c\ell}+\ell/N\}
\]
per cell. Summing the $N$ cells gives $C\log^2(N+2)$.
The finitely many smaller $N$ for which these windows do
not fit are absorbed by increasing the constant.

Combining the strips and bands proves
$|\EE_P H-\EE_Q H|\le C\log^2(N+2)$.
Lemma~\ref{lem:Poisson-circle-profile-mean} replaces
$\EE_QH$ by $N\mathfrak h(k/N)$ at constant cost.
The bounded expected label correction in
Lemma~\ref{lem:history-label-correction} proves the second
term of \eqref{eq:uniform-history-mean-bias}.
\end{proof}

\begin{corollary}[Profile-centered dense histories]
\label{cor:profile-centered-history-CLT}
Whenever $k/N\to\theta\in(0,1)$,
\[
 \frac{H_N(k)-N\mathfrak h(k/N)}{\sqrt N}
        \Longrightarrow\mathcal N(0,v_{\rm D}(\theta)),
\]
and the same holds for the label-based history.
More generally the exact mean in any Gaussian regime above
may be replaced by $N\mathfrak h(k/N)$ whenever its standard
deviation divided by $\log^2(N+2)$ tends to infinity.
In particular this applies to the complete-refresh Gaussian
regime if $(N-k)/[\sqrt N\log^2(N+2)]\to\infty$.
The function $\mathfrak h$ is the established limiting profile;
this statement does not evaluate it or replace it by only
its leading endpoint expansion.
\end{corollary}
\begin{proof}
Use \eqref{eq:uniform-history-mean-bias} and Slutsky's theorem.
\end{proof}

\subsection{Quantitative spectrum convergence and endpoint centering}
\label{sec:endpoint-centering}

\status{Working proof: profile centering throughout the Gaussian endpoint}

\begin{lemma}[From the real profile to its low-order derivatives]
\label{lem:matching-spectrum-derivatives}
Use the finite weights $a_{N,j}$ from
Corollary~\ref{cor:matching-spectrum}, putting $a_{N,j}=0$ for $j>N$,
and define
\[
 G_N(z)=\sum_{j\ge2}a_{N,j}z^j,\qquad
 G(z)=\sum_{j\ge2}\pi_jz^j=1-\mathfrak h(1-z).
\]
For each fixed integer $r\ge0$, uniformly on $0\le z\le1/2$,
\begin{equation}\label{eq:matching-spectrum-derivatives}
 |G_N^{(r)}(z)-G^{(r)}(z)|
       \le C_r\frac{\log^{2r+3}(N+2)}N.
\end{equation}
In particular, for fixed $j\ge2$,
$|a_{N,j}-\pi_j|\le C_j\log^{2j+3}(N+2)/N$, and
\begin{equation}\label{eq:matching-spectrum-quadratic-error}
 |G_N(z)-G(z)|\le
          C z^2\frac{\log^7(N+2)}N,\qquad 0\le z\le1/2.
\end{equation}
\end{lemma}
\begin{proof}
We first compare sampling with and without replacement.
For $z=M/N\le1/2$ and $2\le j\le N$, put
$b_{N,j}(M)=(M)_j/(N)_j$, with value zero when $j>M$.
Then
\begin{equation}\label{eq:small-retention-avoidance}
 0\le z^j-b_{N,j}(M)\le 3z/N.
\end{equation}
For $j\le N/2$, write the falling-factorial ratio as
$\prod_{i<j}c_i$ with $c_i=\max(0,(M-i)/(N-i))$.
Each $c_i$ lies in $[0,z]$ and $z-c_i\le2i/N$.
Telescoping the product therefore bounds the difference by
$j(j-1)z^{j-1}/N\le3z/N$.
For $j>N/2$ use $z^j\le z2^{-(j-1)}\le3z/N$ instead.
The case $M=0$ is immediate.

The exact spectral formula gives
\[
 \frac{\EE H_N^{\rm lab}(N-M)}N
   = A_N-\sum_{j=2}^Na_{N,j}b_{N,j}(M),
 \qquad A_N=(1-1/N)^2.
\]
Since the weights are nonnegative and sum to $A_N\le1$,
\eqref{eq:small-retention-avoidance} and
Theorem~\ref{thm:uniform-history-mean-bias} imply
$|G_N(z)-G(z)|\le C\log^2(N+2)/N$ at every grid point
$z=M/N\le1/2$. Both derivatives $G_N',G'$ are at most one on
this interval: $j(1/2)^{j-1}\le1$ for $j\ge2$ and both
coefficient masses are at most one. Moving to the closest grid
point extends the same bound to the entire interval.

Here is an elementary derivative estimate with explicit logarithmic
loss. Set $L=\lceil4\log_2(N+2)\rceil$ and truncate $G_N-G$ at
degree $L$, obtaining a polynomial $P_L$. The discarded coefficient
mass in absolute value is at most two. Hence
\[
 \sup_{[0,1/2]}|P_L|
       \le C\log^2(N+2)/N+2^{-L}.
\]
Expand $P_L((x+1)/4)=\sum_{\ell=0}^L d_\ell T_\ell(x)$
in Chebyshev polynomials. The cosine integral formula for these
coefficients gives $|d_\ell|\le2\sup|P_L|$.
For every integer $r\ge0$,
$\sup_{[-1,1]}|T_\ell^{(r)}|\le\ell^{2r}$ for $\ell\ge1$.
Indeed $T_\ell'=\ell U_{\ell-1}$, and $U_{\ell-1}$ is a
nonnegative sum of Chebyshev polynomials with coefficient sum
$\ell$; induction on $r$ proves the bound.
The change of variable contributes $4^r$, so
\[
 \sup_{[0,1/2]}|P_L^{(r)}|
       \le C_r L^{2r+1}
                   \{\log^2(N+2)/N+2^{-L}\}.
\]
The differentiated tail is bounded by
$C_r L^r2^{-L}$ once $L\ge2r$, since its absolute coefficient
mass is at most two. Finitely many smaller $N$ are absorbed in
$C_r$. This proves \eqref{eq:matching-spectrum-derivatives}.
At zero its $j$th derivative is $j!(a_{N,j}-\pi_j)$.
Finally $G_N-G$ and its first derivative both vanish at zero.
Integrating the $r=2$ bound twice proves
\eqref{eq:matching-spectrum-quadratic-error}.
\end{proof}

\begin{theorem}[Endpoint profile centering without a logarithmic restriction]
\label{thm:endpoint-profile-centering}
Put $M=N-k$ and $\delta=M/N$. Uniformly in odd $N$ and
$0\le M\le N/2$, either history convention satisfies
\begin{equation}\label{eq:endpoint-profile-mean-bias}
 |\EE H_N(k)-N\mathfrak h(k/N)|
       \le C\{1+\delta^2\log^7(N+2)\}.
\end{equation}
Consequently, if $M=o(N)$ and $M^2/N\to\infty$,
\begin{equation}\label{eq:endpoint-profile-CLT}
 \frac{H_N(k)-N\mathfrak h(k/N)}
                  {\sqrt{2\pi_2}\,M/\sqrt N}
                  \Longrightarrow\mathcal N(0,1),
\end{equation}
also for labels.
\end{theorem}
\begin{proof}
For labels subtract the exact spectral formula in the preceding
proof from $N\mathfrak h(k/N)=N-NG(\delta)$.
The difference $N(A_N-1)=-2+1/N$ is bounded.
Replacing the falling-factorial sum by $G_N(\delta)$ costs
at most $3\delta$, by \eqref{eq:small-retention-avoidance}.
Replacing $G_N$ by $G$ costs at most
$C\delta^2\log^7(N+2)$, by
\eqref{eq:matching-spectrum-quadratic-error}.
This proves \eqref{eq:endpoint-profile-mean-bias} for labels.
The spatial convention differs in expectation by at most four.

Divide the bias by $\sqrt N\,\delta=M/\sqrt N$.
The result is at most
\[
 C\left\{\frac1{\sqrt N\,\delta}
                  +\frac{\delta\log^7(N+2)}{\sqrt N}\right\}=o(1)
\]
under the stated hypotheses. Apply
Theorem~\ref{thm:endpoint-Gaussian} and Slutsky's theorem.
\end{proof}

\begin{corollary}[Finite spectral expansions as Gaussian centers]
\label{cor:endpoint-polynomial-centering}
Fix an integer $J\ge2$. If
\[
 \sqrt N\ll M=N-k\ll N^{1-1/(2J)},
\]
then both history conventions satisfy
\begin{equation}\label{eq:endpoint-polynomial-CLT}
 \frac{H_N(k)-
       \left\{N-1-N\sum_{j=2}^J\pi_j(M/N)^j\right\}}
              {\sqrt{2\pi_2}\,M/\sqrt N}
          \Longrightarrow\mathcal N(0,1).
\end{equation}
In particular the explicit quadratic center
$N-1-\pi_2 M^2/N$ is valid throughout
$\sqrt N\ll M\ll N^{3/4}$, with $\pi_2$ given by
\eqref{eq:matching-pi-two} and its certified enclosure.
The coefficients of order three and higher remain unevaluated.
\end{corollary}
\begin{proof}
Nonnegativity and total mass one give
$0\le\sum_{j>J}\pi_j\delta^j\le\delta^{J+1}$.
The difference of the displayed center from the full profile
center, divided by $\sqrt N\,\delta$, is therefore bounded by
$1/(\sqrt N\,\delta)+\sqrt N\,\delta^J=o(1)$.
Apply \eqref{eq:endpoint-profile-CLT}.
\end{proof}

\begin{proposition}[A positive cubic coefficient and the sharp quadratic window]
\label{prop:matching-cubic-positive}
The third matching spectral coefficient satisfies
\begin{equation}\label{eq:matching-cubic-lower}
 \pi_3\ge\frac{25333744}{502211745}>\frac1{20}.
\end{equation}
Suppose $M=N-k=o(N)$ and $M^2/N\to\infty$.
If $M/N^{3/4}\to c\in[0,\infty)$, then for either convention,
\begin{equation}\label{eq:quadratic-center-critical-shift}
 \frac{H_N(k)-(N-1-\pi_2M^2/N)}
                   {\sqrt{2\pi_2}\,M/\sqrt N}
   \Longrightarrow
       \mathcal N\left(-\frac{\pi_3c^2}{\sqrt{2\pi_2}},1\right).
\end{equation}
If $M/N^{3/4}\to\infty$, this normalized expression tends to
$-\infty$ in probability. Thus a centered standard normal limit
with the quadratic center holds precisely when $M=o(N^{3/4})$
within the stated Gaussian endpoint regime.
The bound \eqref{eq:matching-cubic-lower} is not an evaluation of
$\pi_3$.
\end{proposition}
\begin{proof}
For $x>0$ let $f_x(\eta)$ be the indicator that the forced points
$0,x$ match in a rate-one Poisson input $\eta$ on the line.
Write $p(x)=\EE f_x(\eta)$ and
$p_x(z)=\EE f_x(\eta+\delta_z)$. These decisions are determined
by the finite union component at separation $x$.
Couple two inputs with common Poisson intensity $\delta$ and
private intensities $1-\delta$, and let $p_\delta(x)$ be the
probability that both accept the forced edge.
For their common input $C_\delta$, put
$g_\delta(C_\delta)=\EE[f_x(C_\delta+P_{1-\delta})\mid C_\delta]$.
Conditional independence gives
\[
 p_\delta(x)-p(x)^2=\operatorname{Var}g_\delta(C_\delta)\ge0.
\]

For a bounded interval $W$, take
$a(z)=p_x(z)-p(x)$ there and set
\[
 L=\int_W a(z)\{C_\delta(dz)-\delta\,dz\}.
\]
Then $\EE L=0$ and $\operatorname{Var}L=\delta\int_W a(z)^2dz$.
Reindexing the Poisson count in $W$ gives
\[
 \operatorname{Cov}(g_\delta,L)
   =\delta\int_W a(z)
         \{\EE g_\delta(C_\delta+\delta_z)-\EE g_\delta(C_\delta)\}\,dz
   =\delta\int_W a(z)^2dz.
\]
The reindexing follows by conditioning on the count and independent
uniform locations in $W$; the exterior is independent and can
remain conditioned throughout. The second equality uses
$C_\delta+P_{1-\delta}\stackrel{\rm law}=\eta$.
Cauchy--Schwarz, including the trivial zero-variance case, proves
\begin{equation}\label{eq:forced-edge-linear-projection}
 p_\delta(x)-p(x)^2
       \ge\delta\int_W(p_x(z)-p(x))^2\,dz.
\end{equation}

For $x>0$, the explicit formula
$p(x)=e^{-\operatorname{Ein}(2x)}(1-e^{-2x})/(2x)$ gives
$p(x)\ge e^{-3x}$:
$\operatorname{Ein}(2x)\le2x$, and the remaining factor is
$\int_0^1e^{-2xt}dt\ge e^{-x}$.
An added point at distance $a$ outside either endpoint has
the form $z=-a$ or $z=x+a$; take $0<a<x$.
Inside the edge, use $z=a$ or $z=x-a$ with $0<a<x/2$.
In each case, an empty Poisson interval of length $3a$
around the added point and its nearer endpoint makes those
two points match at separation $a$, regardless of the exterior.
The other forced endpoint lies farther away.
Thus $p_x(z)\le1-e^{-3a}$.
The two interior and two exterior regions are disjoint.
Put $u=3x$, $v=3a$ and
\[
 P(u,v)=1-u-v+\frac{u^2+v^2}{2}-\frac{u^3+v^3}{6}.
\]
The cubic Taylor lower bound for $e^{-y}$, $y\ge0$, gives
$p(x)-p_x(z)\ge e^{-u}+e^{-v}-1\ge P(u,v)$.
For $w=u+v\le4/3$,
\[
 P(u,v)=1-w+\frac{w^2}{2}-\frac{w^3}{6}
                              +uv(w/2-1)
       \ge 1-w+\frac{w^2}{4}-\frac{w^3}{24}\ge\frac1{81}.
\]
Here $uv\le w^2/4$, and the last polynomial decreases
on $[0,4/3]$. Thus this lower bound is positive throughout
the region used below.

By the Palm identity in
Proposition~\ref{prop:matching-quadratic-endpoint},
$G(\delta)=2\delta^2\int_0^\infty p_\delta(x)\,dx$ and
$\pi_2=2\int_0^\infty p(x)^2dx$.
For $c>0$, put
\[
 I(c)=\int_0^{4/9}\int_0^{\min(cx,4/9-x)}
                         P(3x,3a)^2\,da\,dx.
\]
Apply \eqref{eq:forced-edge-linear-projection} once with
$W=(-x,2x)$ and restrict its integral to the four disjoint
regions just described, using nonnegativity elsewhere.
The two exterior regions have $c=1$ and the two interior
regions have $c=1/2$, giving
\[
 \frac{G(\delta)-\pi_2\delta^2}{\delta^3}
                    \ge4\{I(1)+I(1/2)\}.
\]
Set $w=u+v$ and $t=v/(u+v)$.
The Jacobian from $(w,t)$ to $(x,a)$ is $w/9$.
Direct integration of the displayed polynomial gives
\[
 \begin{aligned}
 9I(1)&=\int_0^{4/3}\int_0^{1/2}
               wP(w(1-t),wt)^2\,dt\,dw=\frac{45988}{688905},\\
 9I(1/2)&=\int_0^{4/3}\int_0^{1/3}
               wP(w(1-t),wt)^2\,dt\,dw=\frac{2608408}{55801305}.
 \end{aligned}
\]
Consequently
$4\{I(1)+I(1/2)\}=25333744/502211745>1/20$.
The power-series identity for $G$ makes the left side tend to
$\pi_3$ as $\delta\downarrow0$, proving
\eqref{eq:matching-cubic-lower}.

The difference between the profile center and the quadratic
center is
$1-N\sum_{j\ge3}\pi_j\delta^j$, where $\delta=M/N$.
If $M/N^{3/4}\to c<\infty$, its normalized cubic term tends to
$-\pi_3c^2/\sqrt{2\pi_2}$.
The remaining terms are bounded in absolute value by
$C\sqrt N\,\delta^3=o(1)$; the constant term also vanishes.
Equation~\eqref{eq:endpoint-profile-CLT} proves
\eqref{eq:quadratic-center-critical-shift}.
If instead $M/N^{3/4}\to\infty$, nonnegativity of all higher
coefficients bounds the deterministic normalized shift above by
$-\pi_3(M/N^{3/4})^2/\sqrt{2\pi_2}+o(1)$, which tends to
$-\infty$, whereas the profile-centered variable is tight.
The necessity for a general sequence follows by taking a
subsequence on which $M/N^{3/4}$ has a positive finite or
infinite limit.
\end{proof}

\subsection{Maximal bounds for the early remainder and background mean}
\label{sec:history-process-remainders}

\status{Working proof: uniform path errors; the full functional CLT remains separate}

For a finite signed measure $\mu$ on a compact time interval,
write $\|\mu\|_{\rm TV}$ for its total variation and
$\|\mu\|_{\rm c}=\sup_t|\mu([0,t])|$ for its cumulative norm.
The supremum may be taken over a fixed countable dense set and
the endpoints, by right continuity.

\begin{lemma}[A maximal inequality for independent signed measures]
\label{lem:random-measure-maximal}
Let $\mu_1,\ldots,\mu_n$ be independent signed random measures
with $\sup_i\EE\|\mu_i\|_{\rm TV}^p\le K_p$ for some $p\ge2$.
Then
\begin{equation}\label{eq:random-measure-maximal}
 \left\|\left\|\sum_{i=1}^n(\mu_i-\EE\mu_i)\right\|_{\rm c}
                                      \right\|_p
                         \le C_p K_p^{1/p}\sqrt n.
\end{equation}
For iid marked renewal blocks $(L_i,\mu_i)$ with $L_i\ge1$,
uniform exponential length moments and all fixed joint moments
of $L_i,\|\mu_i\|_{\rm TV}$, put
$\rho=\EE\mu_i/\EE L_i$.
In the stationary line process, the sum $S_T$ over block starts
in an interval of length $T$ satisfies
\begin{equation}\label{eq:renewal-measure-maximal}
 \bigl\|\|S_T-T\rho\|_{\rm c}\bigr\|_p\le C_p\sqrt{T+1}.
\end{equation}
\end{lemma}
\begin{proof}
Symmetrize using independent copies and Rademacher signs.
It suffices to bound $\|\sum_i\varepsilon_i\mu_i\|_{\rm c}$.
Split each measure into its two Jordan parts. For a positive
part of mass $W_i$, sample a point $U_i$ independently with
conditional distribution $\mu_i/W_i$; a zero mass contributes zero.
Conditional Jensen bounds the $p$th moment of its Rademacher
cumulative sum by that of
$\sum_i\varepsilon_iW_i\ind_{\{U_i\le t\}}$.
Given the masses and sampled points, order the latter by $U_i$.
Their partial sums form a martingale in the independent signs.
The martingale maximal inequality and the Rademacher moment
bound give
\[
 \EE_\varepsilon\sup_t
       \left|\sum_i\varepsilon_iW_i\ind_{\{U_i\le t\}}\right|^p
       \le C_p\left(\sum_iW_i^2\right)^{p/2}.
\]
Ties only restrict the possible partial sums. Finally
$(\sum_iW_i^2)^{p/2}\le n^{p/2-1}\sum_iW_i^p$.
This proves \eqref{eq:random-measure-maximal}.
The Rademacher moment bound for even $p$ follows by expanding
the sum and retaining even multiplicities; general $p\ge2$
follows by the adjacent larger even moment and monotonicity
of conditional $L^p$ norms.

For the renewal assertion set $Z_i=\mu_i-L_i\rho$.
These are iid centered signed measures with bounded variation
moments. The cumulative norm of their successive partial sums
is a nonnegative submartingale. Applying the martingale maximal
inequality in the block index and then
\eqref{eq:random-measure-maximal} gives the same $\sqrt n$ bound
for the maximum over all indices up to $n$.
After removing the block containing the left endpoint, at most
$T+1$ blocks can start in the interval, since $L_i\ge1$.
Use that deterministic upper bound on the random index.
The difference between their total length and $T$ is bounded
by the two boundary-block lengths. Those blocks have the
stationary length-biased law, and their marks and lengths have
bounded fixed moments by the joint moment assumptions.
Their contribution, and the term involving $\rho$, is $O_p(1)$
in $L^p$, proving \eqref{eq:renewal-measure-maximal}.
\end{proof}

\begin{lemma}[The same maximal bound for cyclic renewal marks]
\label{lem:cyclic-measure-maximal}
Suppose a family of iid marked blocks $(L,\mu)$ has $L\ge1$,
uniform exponential length moments, and densities
\[
 \frac{\EE[\|\mu\|_{\rm TV}^p;L\in dy]}{dy}\le C_pe^{-c_py}
\]
for every fixed $p$. Suppose its ordinary renewal density is
uniformly bounded. Give cyclic blocks on circumference $T\ge8$
the density $T/j$ times the product block law, constrained to
total length $T$, with a normalization bounded above and below.
One may also add a parity-signed product law whose total variation
is dominated by the unsigned product law.
Allow an exceptional sector of probability $Ce^{-cT}$, provided
the total mark has every variation moment bounded by a polynomial
in $T$, including on that sector.
Then the total circle mark $S_T$ obeys
\begin{equation}\label{eq:cyclic-measure-maximal}
 \bigl\|\|S_T-\EE S_T\|_{\rm c}\bigr\|_p\le C_p\sqrt T.
\end{equation}
\end{lemma}
\begin{proof}
Center each block by $\rho L$ with
$\rho=\EE\mu/\EE L$. Split the circle's block starts between
two deterministic half circles. Reveal all records of the blocks
covering one half, including its two boundary blocks.
If their combined length $s$ is at most $3T/4$, their law has
uniformly bounded density relative to the corresponding
stationary line records.
Indeed if $v$ is the origin's position in its first block, the
line record density is
$\nu\,dv\prod dG_i$, $\nu=1/\EE L$.
Integrating the unexposed cyclic blocks replaces $\nu$ by
the outside renewal density at $T-s$, divided by the normalization.
The parity term has absolute value at most the same unsigned
renewal density. All these factors are uniformly bounded,
and $\nu$ is bounded away from zero.
The atomic completion at $T-s=0$ is absent on this event.

The sum of the centered marks with starts in that half circle
therefore satisfies the line bound
\eqref{eq:renewal-measure-maximal}, together with the
$O_p(1)$ boundary correction. Sum the bounds for the two halves.
If their covering records have length greater than $3T/4$,
one boundary block has length at least $T/8$.
The cyclic one-block formula and the exponential density
bounds give probability $Ce^{-cT}$ for this event, including
the atomic one-block completion.
On the regular sectors the total variation moments are
automatically polynomial: there are at most $T$ blocks, and
the one-block formula bounds the expected sum of their $p$th
variation powers by $C_pT$.
On the long-block event and the allowed exceptional sector, H\"older and
the polynomial total-variation moment bounds make every
required cumulative moment exponentially small.
This proves an $O(\sqrt T)$ bound centered by the measure $T\rho$.
Its expectation has cumulative norm at most the same bound;
subtracting the actual mean proves
\eqref{eq:cyclic-measure-maximal}.
\end{proof}

\begin{theorem}[The centered early remainder is uniformly small as a path]
\label{thm:early-remainder-maximal}
For an exact background of size $M$ and any split $q=d+a$,
let $R_{M;d,a}(r)$ be the actual history additivity remainder
through separation $r$, as in
Corollary~\ref{cor:all-sparse-remainder-variance}.
There are absolute $C,\alpha_0<\infty$ such that
\begin{equation}\label{eq:early-remainder-maximal}
 \EE\sup_{0\le r\le1/q}
       |R_{M;d,a}(r)-\EE R_{M;d,a}(r)|^2\le Cq,
                  \qquad 1\le q\le M/\alpha_0.
\end{equation}
For a Poisson or entrance background the same bound holds
with exactly $q\ge8$ probes; on circumference $q$ its cutoff
is one, and each fixed centered maximal $L^p$ norm is
at most $C_p\sqrt q$.
Consequently the centered remainder path divided by
$\sqrt{q\log(M/q)}$ vanishes in $L^2$ whenever $q=o(M)$,
uniformly in the two-colour split.
\end{theorem}
\begin{proof}
First take Poisson probe clouds on circumference $T\ge8$
and background intensity $\alpha\ge1$. Regard every block
remainder as a signed measure on collision separations in $[0,1]$.
The moment arguments in
Theorem~\ref{thm:line-block-remainder-moments} apply to its
total variation, not only its terminal value.
To check this strengthening, index each matching vector by
its physical edges and push its signed coordinates forward
by their fixed edge lengths. Total variation cannot increase
under this push-forward. In each band the telescoping
triangle inequality bounds the variation by the same sum
of single-insertion costs used there. If the probes occupy
separated fine gaps, the entire remainder measure vanishes
below that cutoff. Thus the rare-pair factor and the
geometric dyadic sum are unchanged.
The final empty gap is independent of the occupied-run
measure, so convolution with its density gives the exponential
variation-moment densities required by
Lemma~\ref{lem:cyclic-measure-maximal}.

The same strengthening applies to
Lemma~\ref{lem:circle-remainder-polynomial-moments}:
its band domination and its zero-remainder event concern
the whole edge vector. Hence the full circle measure satisfies
$\|\|R_T\|_{\rm TV}\|_p\le C_pT(1+\log(T+2))$, uniformly in
$\alpha$. The no-empty-gap and empty-configuration sectors
still have probability $Ce^{-cT}$.
The block renewal identities \eqref{eq:block-renewal-identities}
give a bounded ordinary renewal density, and the cyclic
marked law is exactly the one in
Theorem~\ref{thm:Poisson-circle-remainder-CLT}.
Lemma~\ref{lem:cyclic-measure-maximal} now proves
$\|\|R_T-\EE R_T\|_{\rm c}\|_p\le C_p\sqrt T$.
Finite-intensity marked convergence and the same dyadic
variation tail pass this bound to the entrance law.

Next condition the two probe counts. The local-increment
proof and its two-reservoir transfer also hold in variation.
At a single edge, the change of an absolute Boolean difference
is bounded by the changed edge indicator in the edited input.
Summing over edges in a band therefore gives
\[
 \|D_ZR_T|_I\|_{\rm TV}
       \le J_Z(B+P;I)+J_Z(B;I)
\]
for a probe of colour $P$.
The same fine-gap event makes the full increment measure
zero; thus the dyadic rare-interaction proof gives uniformly
bounded fixed variation moments of one increment.
The unused records in Lemma~\ref{lem:two-count-reservoirs}
determine the whole local measure, so its probe-only
conditioning argument applies without change.
The clipped Poisson-count coupling in
Theorem~\ref{thm:fixed-probe-remainder-variance}, now using
Minkowski in variation, costs $C_p\sqrt q$ in $L^p$.
The unclipping term is exponentially small by the polynomial
variation bounds and Poisson count tails.
Centering the coupling difference gives the claimed
Poisson-background and entrance maximal bounds.

Finally fix the original count, directly through the partial-record
reservoir. On the unit circle use a Poisson background of mean $M$
and set $r_*=1/\sqrt{Mq}$.
The data of Lemma~\ref{lem:relative-history-depoisson} determine
the entire remainder measure, including its single-probe terms.
Its bounded-density argument needs only
$Mr_*$ larger than a fixed constant and
$qr_*$ smaller than a fixed constant: the unoccupied short
gaps must retain a fixed positive total length.
Both conditions hold for $q\le M/\alpha_0$ with $\alpha_0$
sufficiently large, even though its stronger total-variation
convergence statement requires $q=o(M)$.
Thus on its good event the exact-count data density is bounded
by a constant, and the complement has Poisson probability
$Ce^{-c\sqrt{Mq}}$. Matching the prescribed original parity
costs only a constant.

For $q\ge8$ the fixed-probe Poisson remainder has every
uncentered variation moment bounded by a polynomial in $q$:
condition the Poisson probe counts at their means, with point
probability at least $c/(q+1)$, in the preceding circle bound.
For $q<8$, work first on circumference eight through cutoff one;
the remaining interval up to rescaled cutoff $8/q$ has bounded
variation by active packing. This gives the same uniform
polynomial moment bound for those finitely many counts.
Global conditioning of the original count costs at most
$C\sqrt M$. H\"older therefore bounds its bad-event contribution
to the maximal second moment by
$C\sqrt M(q+1)^C e^{-c\sqrt{Mq}}$, which is at most $Cq$.
On the good event the Poisson centered maximal second moment
transfers at constant cost. Subtracting the exact mean path
increases its $L^2$ bound by at most a factor two.
This proves \eqref{eq:early-remainder-maximal}; division by
$q\log(M/q)$ proves the final assertion.
\end{proof}

\begin{corollary}[A maximal bound for the integrated capture process]
\label{cor:integrated-capture-maximal}
On a circle of circumference $T\ge8$ with Poisson background
intensity $\alpha\ge1$, let
$A(r)=\int_0^T G_B(z;[0,r])\,dz$ for $0\le r\le1$.
Uniformly in $\alpha,T$, and optionally with prescribed
original parity, each fixed $p\ge2$ satisfies
\begin{equation}\label{eq:integrated-capture-maximal}
 \left\|\sup_{0\le r\le1}|A(r)-\EE A(r)|\right\|_p
                      \le C_p\sqrt T.
\end{equation}
With $T=q$, this makes the entire background conditional-mean
path negligible on the scale $\sqrt{q\log\alpha}$.
\end{corollary}
\begin{proof}
In one dyadic band $(b/2,b]$, scale by $b$ and attach to a
cutoff gap its positive measure of integrated capture times,
divided by $b$. Its total variation is at most $CY^2$ by
\eqref{eq:integrated-capture-local-bound}.
For the microscopic strip it is at most $CY(N_I+1)$, whose
moment densities have the uniform exponential envelope
proved in Theorem~\ref{thm:Poisson-integrated-capture}.
The cyclic gap law and its bounded renewal density therefore
satisfy Lemma~\ref{lem:cyclic-measure-maximal}.
On an empty sector, the variation is just the total positive
band cost, so the deterministic packing and microscopic count
bounds from that theorem also verify the exceptional-sector
hypothesis. Scaling back gives centered maximal $L^p$ norm
at most $C_p\sqrt{Tb}$ within that band.
Sum these bounds and the microscopic strip by Minkowski.
The geometric sum of $\sqrt b$ is bounded, which proves
\eqref{eq:integrated-capture-maximal}.
\end{proof}

\subsection{The unrestricted sparse functional limit on its natural clock}
\label{sec:all-sparse-functional}

\status{Working proof: cumulative early histories and the completed endpoint}

\begin{lemma}[Uniform centered logarithmic increments of one probe]
\label{lem:single-probe-log-increments}
On a stationary Poisson line of intensity $\alpha$, or on a
circle of circumference $T\ge8$, let $G(a,b]$ be the
single-insertion history cost in a separation interval $(a,b]$.
Suppose $\alpha a\ge1$ and $0<a\le b\le1$, and put
$\ell=\log(b/a)$. For each fixed $p\ge2$,
\begin{equation}\label{eq:single-probe-log-increments}
 \|G(a,b]-\EE G(a,b]\|_p\le C_p\sqrt{1+\ell}.
\end{equation}
Uniformly in these parameters, as $\ell\to\infty$,
\begin{equation}\label{eq:single-probe-increment-variance}
 \operatorname{Var}G(a,b]/\ell\longrightarrow4v_{\rm H}.
\end{equation}
The circle assertions allow either original parity.
The microscopic initial strip $[0,1/\alpha]$ has bounded
uncentered moments of every fixed order.
\end{lemma}
\begin{proof}
Scale $b$ to one, writing $\beta=\alpha b\ge e^\ell$.
The new circumference is at least eight.
Unit-log-band costs have uniform uncentered moments by
\eqref{eq:line-single-band-moments} and its circle version,
so the full interval has $L^p$ norm at most $C_p(1+\ell)$.
This settles bounded $\ell$.

For large $\ell$ set $w=A_p\log(\ell+2)$, taking $\ell$
large enough that $2w<\ell$.
Keep the central interval
$[r_-,r_+]=[e^{-\ell+w},e^{-w}]$ and truncate its source-root
skeleton at $a_*=e^{-w/4}r_-=e^{-\ell+3w/4}$.
Apply the Poisson skeleton and origin-query comparison used in
Lemma~\ref{lem:Poisson-point-capture}, with this smaller
logarithmic interval. Its finite-intensity skeleton error on a
source root of length at most $wr_+$ is bounded by
\[
 C(wr_+/a_*)e^{-2\beta a_*}
 \le Cw e^{\ell-7w/4}\exp\{-2e^{3w/4}\}.
\]
Here the only needed intensity inequality is $\beta\ge e^\ell$;
increasing the actual intensity improves this estimate.
Choose $A_p>4$, and then increase it for the fixed moment order.
The source-root tail, inverse terminal-cutoff estimate, and
origin-completion budget are exactly those of that comparison,
with $a_*/r_-=e^{-w/4}$; they contribute at most $Ce^{-cw}$.
On a circle the outside renewal factor differs from its line
factor by an exponentially small term in $T/r_+$.
Thus the central interval couples to a stationary line
entrance capture interval of logarithmic length $\ell-2w$
with failure probability $Ce^{-cw}$, uniformly in $\beta$.

Each omitted strip has uncentered $L^p$ norm at most $C_pw$.
The interval costs on both sides of the coupling have
$L^{2p}$ norm at most $C_p(1+\ell)$ by their unit-band
bounds. H\"older therefore makes the coupling failure cost
bounded in $L^p$ by increasing $A_p$.
In the coupled interval the changed-edge count is twice the
entrance capture count, up to an error of absolute value one.
The line cumulant bounds give centered $L^p$ norm $C_p\sqrt\ell$
and variance $4v_{\rm H}(\ell-2w)+o(\ell)$.
Since $w=o(\sqrt\ell)$, centering the omitted strips proves
both \eqref{eq:single-probe-log-increments} and
\eqref{eq:single-probe-increment-variance}.
The microscopic assertion is the original exponential-gap
component bound, at intensity times cutoff equal to one.
\end{proof}

\begin{theorem}[A Brownian history limit for every growing sparse edit count]
\label{thm:all-sparse-history-FCLT}
Let an exact iid background have size $M$, and let the two
independent uniform probe families have counts $d,a$ with
$q=d+a\to\infty$ and $q=o(M)$. Put
\[
 L=\log(M/q),\qquad
 r_M(t)=M^{-1}(M/q)^t,\qquad 0\le t\le1.
\]
Let $H_M(t)$ be their spatial matching-history distance
through separation $r_M(t)$. Uniformly over all splits,
and with either input parity,
\begin{equation}\label{eq:all-sparse-history-FCLT}
 \left(\frac{H_M(t)-\EE H_M(t)}
                      {\sqrt{4qv_{\rm H}L}}\right)_{0\le t\le1}
       \Longrightarrow (B_t)_{0\le t\le1}
\end{equation}
in the Skorokhod space, with continuous limit.
The value at $t=1$ may instead be the completed history cost,
centered by its completed expectation.

This clock runs from the microscopic separation $1/M$ to
$1/q$. The last assertion attaches the completed endpoint;
it does not parameterize the intervening separations after
$1/q$ as an additional path segment.
Section~\ref{sec:late-history-maximal} supplies that extension.
The centering is the exact mean at each time.
\end{theorem}
\begin{proof}
First use a Poisson background of mean $M$ on the unit circle,
with its prescribed parity and exactly $q$ uniform probes.
Scale by $q$, giving circumference $q$ and intensity
$\alpha=M/q$. The time cutoffs become
$b(t)=\alpha^{t-1}$, and $L=\log\alpha$.
Write $G_B(z,t)$ for one probe's cost through $b(t)$.
For independent uniform probes $Z_i$ define
\[
 X_q(t)=\sum_{i=1}^q
       \{G_B(Z_i,t)-\EE[G_B(Z_i,t)\mid B]\}.
\]
The history minus its exact mean is $X_q$ plus the centered
integrated-capture path and the centered additivity remainder.
Theorem~\ref{thm:early-remainder-maximal} and
Corollary~\ref{cor:integrated-capture-maximal} make those two
paths $O(\sqrt q)$ in each fixed maximal $L^p$ norm.
They are negligible after division by $\sqrt{qL}$.

We first prove finite-dimensional convergence of $X_q$.
For a fixed collection of times and real weights, put
$Y_B(z)=\sum_jc_jG_B(z,t_j)$, and let $s_{\alpha,c}^2$ be
the variance of the corresponding single-probe variable
on the Poisson line.
The one-time variance limit and
\eqref{eq:single-probe-increment-variance}, by polarization,
give
\begin{equation}\label{eq:single-probe-clock-covariance}
 \frac{s_{\alpha,c}^2}{L}\longrightarrow
       4v_{\rm H}\sum_{i,j}c_ic_j\min(t_i,t_j).
\end{equation}
The time-zero microscopic cost has bounded moments and makes
no contribution to this formula.

The conditional-variance proof of
Proposition~\ref{prop:Poisson-conditional-capture-variance}
applies to $Y_B$ too. Here are its required uniform bounds.
Each centered coordinate has $p$th moment at most $C_pL^{p/2}$,
by Lemma~\ref{lem:Poisson-point-capture} after scaling its
cutoff; fixed linear combinations obey the same bound.
Their squared-cost gap marks therefore have exponentially
weighted second moment at most $C_cL^2$.
Their pointwise mark densities are bounded by
$C_c(1+L)^4(1+y)^C e^{-cy}$, since each restricted cost is
at most the full positive capture cost and the ancestry
bound applies inside the source gap.
Finally the integrated linear-combination mean has centered
$L^4$ norm $C_c\sqrt q$ by the maximal integrated-capture bound;
its circle/line mean error is $C_ce^{-cq}$ by
\eqref{eq:Poisson-point-capture-mean-bias} at each cutoff.
The same cyclic second-moment cancellation consequently yields
\[
 \EE\left|
 \frac{\operatorname{Var}(Y_B(Z)\mid B)-s_{\alpha,c}^2}{L}
                         \right|^2\le C_c/q
       \quad\hbox{if }q\ge D_c\log(L+2).
\]
Conditional independence of the probes and their fourth
centered moments make the expected Lindeberg error at most
$C_c/q$. The conditional central limit argument and
\eqref{eq:single-probe-clock-covariance} prove the Gaussian
finite-dimensional limits.
If $q<D_c\log(L+2)$, use the Poisson version of the growing-root
comparison, exactly as in
Theorem~\ref{thm:Poisson-early-history-CLT}.
That comparison is simultaneous over the whole central
logarithmic interval and has path error $C_pq\log(L+2)$.
Its normalized error tends to zero since
$q\log^2(L+2)/L\to0$.
The independent line functional limits give the same
finite-dimensional laws in this remaining range.

For tightness, take $0\le s<t\le1$. The single-probe centered
increment bound and conditional Jensen imply
\[
 \EE\left|
 G_B(Z,t)-G_B(Z,s)
       -\EE[G_B(Z,t)-G_B(Z,s)\mid B]\right|^4
              \le C(1+L(t-s))^2.
\]
The fourth-moment expansion for the conditionally independent
probe sum, with Jensen also bounding the squared conditional
variance, gives
\begin{equation}\label{eq:capture-clock-increment-fourth}
 \EE|X_q(t)-X_q(s)|^4
             \le Cq^2(1+L(t-s))^2.
\end{equation}
Interpolate the normalized process $X_q/\sqrt{qL}$ on the
grid $t_j=j/\lceil L\rceil$.
At distinct grid points \eqref{eq:capture-clock-increment-fourth}
is bounded by $C|t_j-t_i|^2$ after normalization.
Linear interpolation gives that bound for arbitrary times
as well; within a single grid interval the fourth power of
the interpolation fraction compensates for its shorter length.
The fourth-moment tightness criterion therefore applies.
It also gives a uniform maximal fourth moment: sum the
$L^4$ norms of maximal dyadic increments, bounded by
$C2^{-h/4}$ at level $h$. The initial microscopic value has
bounded normalized fourth moment.

It remains to compare the actual path to this interpolation;
endpoint increment bounds alone would not suffice.
Within one grid interval, a probe's positive measure of
new edge times has total mass with uniformly bounded fourth
moment, since the separation ratio is at most $e$.
Conditional on $B$ these measures are iid over probes.
Lemma~\ref{lem:random-measure-maximal}, followed by averaging
over $B$, bounds the fourth moment of their centered cumulative
sum's maximum by $Cq^2$.
There are $O(L)$ grid intervals. Dividing by $(qL)^2$ and
summing shows that the normalized interpolation error tends
to zero in probability, with fourth moment at most $C/L$.
This proves tightness and a uniform maximal fourth moment
for the Poisson history centered by its mean.
Together with its finite-dimensional laws it proves
\eqref{eq:all-sparse-history-FCLT} in the Poisson model.

Now fix the original count. Lemma~\ref{lem:relative-history-depoisson}
with $r_*=1/\sqrt{Mq}$ compares the entire relative-history
path in total variation. To align the mean paths, use its
bounded data density on the good event and the Poisson
maximal fourth-moment bound just proved.
On the complement its probability is $Ce^{-c\sqrt{Mq}}$.
Global original-count conditioning costs $C\sqrt M$.
The whole uncentered history is at most $C(N_B+q)$, where
$N_B$ is the original Poisson count, and the Poisson mean
path is at most $C(M+q)$.
H\"older and Poisson moments therefore make the bad-event
fourth moment, even after conditioning and normalization,
negligible. Thus the exact-count path centered by the
Poisson mean has a uniformly bounded normalized maximal
fourth moment.

Couple the two relative-history paths to agree with probability
tending to one, using their total-variation comparison.
Their normalized difference in cumulative supremum norm has
expectation tending to zero by the maximal fourth-moment
bounds and H\"older. Consequently
\[
 \sup_{0\le t\le1}
 \frac{|\EE_{\rm exact}H_M(t)-\EE_{\rm Pois}H_M(t)|}
                         {\sqrt{qL}}\longrightarrow0.
\]
The full path law therefore transfers with its exact mean
centering. This argument uses a separate uniform-integrability
estimate and does not infer mean convergence from total
variation alone.

Finally Theorem~\ref{thm:late-history-variance} bounds the
centered history after separation $1/q$ by $C\sqrt q$ in
$L^2$. Changing only the value at $t=1$ to the completed
centered cost therefore changes the normalized path by
$o_{L^2}(1)$ in supremum norm. The same Brownian limit follows.
\end{proof}

\begin{corollary}[Resampling and label histories on the logarithmic clock]
\label{cor:resampling-history-FCLT}
For $k\to\infty$, $k=o(N)$, set $M=N-k$, $q=2k$,
$L=\log(M/q)$, and $r(t)=M^{-1}(M/q)^t$.
The cumulative spatial and label history distances, each
centered by its own mean and divided by
$\sqrt{8kv_{\rm H}L}$, converge to standard Brownian motion
on $0\le t\le1$. Either endpoint at $t=1$ may be replaced
by its completed centered cost.
One may replace $L$ in the normalization by $\log(N/k)$.
\end{corollary}
\begin{proof}
Apply Theorem~\ref{thm:all-sparse-history-FCLT}.
For labels, the difference from the spatial cumulative cost
is twice the number of accidental common label edges already
accepted in both copies. Each such edge enters at the larger
of its two acceptance separations, so this correction is
nonnegative, nondecreasing, and bounded by its completed value.
That value has second moment at most $32$ by
Lemma~\ref{lem:history-label-correction}.
The centered path correction is therefore negligible uniformly.
Finally $L/\log(N/k)\to1$.
\end{proof}

\subsection{Uniform control of the late path and the full history clock}
\label{sec:late-history-maximal}

\status{Working proof: the entire late path is negligible at the sparse scale}

\begin{lemma}[Bounded density on an unused fraction of the circle]
\label{lem:fixed-count-restriction-density}
Let $P_n$ be the law of exactly $n$ uniform points on a circle
and $Q_n$ the Poisson law of mean $n$.
For any deterministic measurable union of arcs $W$ occupying
a fraction $p\le1-\varepsilon$, the restricted laws satisfy
\begin{equation}\label{eq:fixed-count-restriction-density}
 \frac{d(P_n|_W)}{d(Q_n|_W)}\le C_\varepsilon,
\end{equation}
uniformly in $n$, including zero. The same bound, with a
constant for the fixed number of colours, holds for independent
colours and leaves any independent Poisson colour unchanged.
\end{lemma}
\begin{proof}
For a restriction containing $j$ points the density is
\[
 \frac{\PP\{\operatorname{Pois}(n(1-p))=n-j\}}
      {\PP\{\operatorname{Pois}(n)=n\}}.
\]
Conditional locations coincide in the two laws.
The numerator is at most $C/\sqrt{1+n(1-p)}$ and the
denominator is at least $c/\sqrt{1+n}$, by the elementary
Poisson point bounds. Their ratio is bounded when
$1-p\ge\varepsilon$. For $n=0$ the laws coincide.
Multiply the densities for independent colours.
\end{proof}

\begin{theorem}[A maximal late-history estimate]
\label{thm:late-history-maximal}
On circumference $T\ge8$, take a Poisson background of intensity
$\alpha\ge1$ and independent uniform probe families of exact
counts $d,a\le K\alpha T$, with fixed $K<\infty$.
Write $H^{(b,s]}$ for their spatial history distance restricted
to separations in $(b,s]$. For $1\le b\le T/2$ and every
fixed $p\ge2$,
\begin{equation}\label{eq:Poisson-late-history-maximal}
 \left\|\sup_{b\le s\le T/2}
   |H^{(b,s]}-\EE H^{(b,s]}|\right\|_p
                 \le C_{p,K}\sqrt{T/b}.
\end{equation}
Either prescribed original parity is allowed.
For an exact background of size $M$ on the unit circle, there
are $\alpha_0,C_p<\infty$ such that
\begin{equation}\label{eq:finite-late-history-maximal}
 \left\|\sup_{1/q\le s\le1/2}
   |H^{(1/q,s]}-\EE H^{(1/q,s]}|\right\|_p
                     \le C_p\sqrt q,\qquad
             2\le q=d+a\le M/\alpha_0.
\end{equation}
The supremum is zero if the interval is empty.
\end{theorem}
\begin{proof}
Fix a band $(u,2u]$, $1\le u\le T/8$, and partition the
circle into $n=\lceil T/u\rceil$ equal cells of length $h$.
Assign edges to their clockwise initial endpoint along their
shorter arc, as in Lemma~\ref{lem:late-band-locality}.
Let $X_j$ be the positive measure of assigned history-edge
separations in cell $j$. Deterministic active packing gives
$\|X_j\|_{\rm TV}\le4$.

Denote the mixed exact-probe law by $P$, and by $Q$ its
fully Poisson counterpart, with probe intensities $d/T,a/T$.
For a cell window $W_{j,\ell}$ containing the cell and
$\ell$ neighbours on either side, let
\[
 f_{j,\ell}=\EE_Q[X_j\mid
                   \hbox{all coloured input records in }W_{j,\ell}].
\]
This is a fixed measurable function of the interior records;
evaluate that same function under $P$.
For $0\le\ell\le\lfloor n/16\rfloor$ and $n\ge32$,
\begin{equation}\label{eq:late-measure-local-approximation}
 \left\|\|X_j-f_{j,\ell}\|_{\rm TV}\right\|_{L^p(P)}
       +\left\|\|X_j-f_{j,\ell}\|_{\rm TV}\right\|_{L^p(Q)}
                       \le C_{p,K}e^{-c_{p,K}\ell}.
\end{equation}
To prove it, retain the interior records and fill all exteriors
with fresh independent Poisson inputs of the $Q$ intensities.
The conditional mean of the new cell score, given the full
old input, is $f_{j,\ell}$.
If the two cell measures differ, a positive accepted edge in
one of the marginal inputs has its parent leaving the window.
This witnesses a difference at any band time simultaneously;
there is no additional union bound over times.

The old marginal input has count
$\operatorname{Pois}(\alpha T)+k$, $k=d$ or $a$.
The count-density calculation in
Lemma~\ref{lem:late-band-locality} bounds it relative to
$\operatorname{Pois}(\alpha T+k)$ by
$C\sqrt{1+k/(\alpha T)}\le C_K$.
Thus the same uniform marginal parent tail applies under $P$.
The filled law has bounded density relative to $Q$ by
Lemma~\ref{lem:fixed-count-restriction-density}, since each
window occupies at most one quarter of the circle.
Its parent tail has the same bound.
Both cell measures have mass at most four; their difference
has variation at most eight. Conditional Jensen proves
\eqref{eq:late-measure-local-approximation}.
The argument under $Q$ is the same without a density comparison.

Put $D_{j,0}=f_{j,0}-\EE_QX_j$ and
$D_{j,\ell}=f_{j,\ell}-f_{j,\ell-1}$ for $\ell\ge1$.
Under $Q$ these local signed measures have mean zero and
variation $L^p$ norm at most $C_{p,K}e^{-c_{p,K}\ell}$.
For each $\ell$, partition the cell indices into
$O(\ell+1)$ families whose windows are disjoint and whose
union occupies at most three quarters of the circle.
For detail, the window-overlap graph has maximum degree at
most $4\ell$, so colour it with at most $4\ell+1$ colours.
Split each colour family into two groups of nearly equal size.
Its windows were disjoint, and each has length at most $T/4$,
so each resulting union has length at most
$T/2+T/8<3T/4$.

Within a family of size $m$, the measures $D_{j,\ell}$ are
independent under $Q$. Lemma~\ref{lem:random-measure-maximal}
therefore gives cumulative maximal $L^p(Q)$ norm at most
$C_{p,K}\sqrt m\,e^{-c_{p,K}\ell}$ for their sum.
The whole sum depends only on that family's union of windows.
Its $P$-moment is at most a constant times its $Q$-moment
by \eqref{eq:fixed-count-restriction-density}.
Re-centering under $P$ costs at most a factor two.
Summing the families, using Cauchy--Schwarz on their sizes,
gives
\[
 \left\|\left\|\sum_j(D_{j,\ell}-\EE_PD_{j,\ell})
                                     \right\|_{\rm c}\right\|_{L^p(P)}
       \le C_{p,K}\sqrt{n(\ell+1)}\,e^{-c_{p,K}\ell}.
\]
Sum over $\ell$.
At the final radius $\lfloor n/16\rfloor$, the remaining
centered variation error is at most $C_{p,K}ne^{-c_{p,K}n}$
in $L^p$ by \eqref{eq:late-measure-local-approximation}.
For bounded $n$ the deterministic mass bound suffices.
We have proved the centered maximal band bound
$C_{p,K}\sqrt{T/u}$.

Decompose $(b,T/2]$ into dyadic bands up to separation
comparable with $T$. Minkowski sums their maximal bounds:
$\sum_{j\ge0}\sqrt{T/(2^jb)}\le C\sqrt{T/b}$.
The terminal sector has bounded deterministic variation by
active packing. If $b$ is already comparable with $T$,
packing alone gives the result.
This proves \eqref{eq:Poisson-late-history-maximal}.
Conditioning the original Poisson count on either parity
costs a constant, and re-centering gives the stated parity version.

For the exact-count assertion, first take $q\ge8$ and scale
the unit circle by $q$, using $\alpha=M/q$ and $b=1$.
The preceding estimate applies with $K=1$.
Use the partial-record reservoir at $r_*=1/\sqrt{Mq}$
as in the proof of Theorem~\ref{thm:early-remainder-maximal}.
For $q\le M/\alpha_0$, its good-data density is bounded and
its complement has Poisson probability $Ce^{-c\sqrt{Mq}}$.
The complete relative history, including the whole late path,
is measurable from those data.
After separation $1/q$ the total history cost is at most
$q+1$, deterministically. Global original-count conditioning
therefore bounds the bad-event $p$th moment by
$C\sqrt M(q+1)^p e^{-c\sqrt{Mq}}\le C_pq^{p/2}$.
The good-event maximal moment transfers at constant cost,
and re-centering at the exact mean completes the proof.
For $2\le q<8$ the deterministic packing bound suffices.
\end{proof}

\begin{corollary}[A functional limit covering the complete history]
\label{cor:full-history-two-clock}
Under the hypotheses of Theorem~\ref{thm:all-sparse-history-FCLT},
put $\alpha=M/q$, $L=\log\alpha$, and parameterize all
separations from $1/M$ to completion by
\[
 r_M(t)=
 \begin{cases}
  M^{-1}\alpha^t,&0\le t\le1,\\
  q^{-1}(q/2)^{t-1},&1\le t\le2.
 \end{cases}
\]
If $H_M(t)$ is the cumulative spatial history cost through
$r_M(t)$, then
\begin{equation}\label{eq:full-history-two-clock}
 \left(\frac{H_M(t)-\EE H_M(t)}
                     {\sqrt{4qv_{\rm H}L}}\right)_{0\le t\le2}
       \Longrightarrow (B_{t\wedge1})_{0\le t\le2}
\end{equation}
in the Skorokhod space, with continuous limit.
The same holds for resampling label histories.
Any deterministic continuous nondecreasing clock covering
$[1/q,1/2]$ may replace the second branch.
\end{corollary}
\begin{proof}
The first branch is Theorem~\ref{thm:all-sparse-history-FCLT}.
By \eqref{eq:finite-late-history-maximal}, the centered
increment over the entire second branch has maximal $L^2$
norm at most $C\sqrt q$. Dividing by $\sqrt{qL}$ makes it
vanish. Extend the first-branch process constantly from one
to two and apply the continuous mapping theorem and this
uniform comparison. At separation $1/2$ every matching is
complete. The monotone bounded label correction from
Corollary~\ref{cor:resampling-history-FCLT} handles labels.
The same maximal estimate is independent of the chosen
parameterization of the late separations.
\end{proof}

\subsection{Joint growing families and their overlap covariance}
\label{sec:joint-growing-history}

\status{Working proof: finitely many growing perturbation comparisons}

\begin{lemma}[Remainder paths with finitely many membership types]
\label{lem:finite-type-remainder}
Fix $m<\infty$. On a Poisson background of intensity $\alpha$
and circumference $q\ge8$, take $q$ independent uniform probes
with deterministic index subsets $S_j,T_j$, $1\le j\le m$.
Put $D_j=S_j\triangle T_j$, and define
\[
 R_j(r)=H(B+S_j,B+T_j;[0,r])
                  -\sum_{i\in D_j}G_B(Z_i;[0,r]),\qquad 0\le r\le1.
\]
Uniformly over the subsets and $\alpha\ge1$, each fixed
centered maximal $L^p$ norm of $R_j$ is at most
$C_{p,m}\sqrt q$. The same holds for the entrance background.
The original Poisson parity may be prescribed.
\end{lemma}
\begin{proof}
Give a probe its membership vector in the $2m$ subsets;
there are at most $R=2^{2m}$ types.
Poissonize these types independently at their prescribed
count means, first allowing their total intensity to lie
in any fixed compact subinterval of $(0,\infty)$.
Use occupied runs for the union of all types.
Lemma~\ref{lem:history-occupied-clusters} applies simultaneously
to every subset, including intersecting $S_j,T_j$.
A run with at most one probe has zero remainder: a probe
present in both inputs contributes neither a history difference
nor a single-probe term.

The total-variation band estimates of
Theorem~\ref{thm:early-remainder-maximal} apply to these remainders.
Telescope each of the two inputs from the background and
add the single-probe terms. This bounds a band's variation
by at most three sums of individual insertion costs.
Every Poisson priority-prefix background is still Poisson.
The fine-scale remainder measure vanishes unless two probes
of any types occupy the same or adjacent finer gaps.
Thus the same $O(b)$ rare-pair probability gives uniform
block variation moments. The independent final empty gap
smooths their moment densities, and the cyclic maximal
inequality gives a centered maximal bound $C_{p,m}\sqrt q$.
The no-empty sector has the same exponential probability
and polynomial variation moments.

For one new probe, exclusive membership in a comparison
costs at most its insertion into that input plus its
single-probe cost from $B$. Membership in both inputs costs
at most the two input insertion costs; there is no new
single-probe term. Membership in neither costs zero.
These are coordinatewise variation inequalities.
An isolated new fine-gap probe has zero remainder increment
in every case. Local increments therefore have the same
uniform variation moments and occupied-run localization as
Lemma~\ref{lem:local-remainder-increment}.

Here are the count-conditioning details for the extra types.
Let their Poisson intensities be $\eta_1,\ldots,\eta_R$,
with the largest first, so $\eta_1$ is bounded below by a
positive constant depending only on $R$ and the lower total
intensity. Leave the count vectors in the short occupied
gaps outside the local exploration unrevealed, as in
Lemma~\ref{lem:two-count-reservoirs}.
They are independent Poisson vectors conditioned to be nonzero.
The atoms $e_1,2e_1$ have probabilities bounded below,
and $e_1+e_j$ has probability at least $c_R\eta_j$.
The two-copy characteristic-function identity gives
\[
 |\phi_i(t)|\le
 \exp\left\{-c_R\left[
   \operatorname{dist}(t_1,2\pi\mathbb Z)^2+
   \sum_{j=2}^R\eta_j\operatorname{dist}(t_j,2\pi\mathbb Z)^2
                         \right]\right\}.
 \]
There are at least $c_Rq$ such gaps except with exponentially
small probability. Fourier inversion bounds the maximal
remaining count mass by
$C_Rq^{-1/2}\prod_{j=2}^R(1+k_j)^{-1/2}$, where
$k_j=q\eta_j$. This matches the target Poisson point
probability up to constants, and includes zero counts.
On the exceptional event the global conditioning factor
is at most a fixed power of $q$, absorbed by its exponential
tail and the increment's higher moments.
Thus exact-type local increments have bounded variation
moments uniformly, even for vanishing type proportions.

Clip the largest Poisson type count to half and twice its
mean, and all other counts to $[0,2q]$.
Every intermediate total intensity remains in a fixed compact
interval depending only on $R$.
The count-coupling argument now costs at most
$C_{p,R}\sum_j\sqrt{k_j}\le C_{p,R}\sqrt q$ in variation
$L^p$. The unclipping contribution is exponentially small,
using the polynomial variation moments.
Re-centering proves the assertion with all type counts fixed.
The entrance limit and either original parity follow by the
same dyadic variation tail and parity arguments as before.
\end{proof}

\begin{theorem}[Joint Brownian histories with overlap covariance]
\label{thm:joint-growing-history}
Let $B$ be an exact iid background of size $M$, with $q$ independent
uniform probes, where $q\to\infty$ and $q=o(M)$.
Fix $m<\infty$ pairs of deterministic subsets
$(S_j,T_j)$ of $\{1,\ldots,q\}$, and put $D_j=S_j\triangle T_j$.
The subsets may vary with $M$. Suppose
\begin{equation}\label{eq:joint-edit-overlap-matrix}
 \Gamma_{jk}=\lim_{M\to\infty}\frac{|D_j\cap D_k|}{q}
\end{equation}
exists for every $j,k$.
Use the complete two-part clock from
Corollary~\ref{cor:full-history-two-clock}, and let $H_j(t)$
be the cumulative spatial history distance between
$B+S_j$ and $B+T_j$ at that cutoff.
Then, jointly in $D([0,2],\mathbb R^m)$,
\begin{equation}\label{eq:joint-growing-history}
 \left(\frac{H_j(t)-\EE H_j(t)}
                  {\sqrt{4qv_{\rm H}\log(M/q)}}\right)_{j=1}^m
       \Longrightarrow \bigl(W_j(t\wedge1)\bigr)_{j=1}^m,
\end{equation}
where $W$ is a centered vector Brownian motion with
\[
 \operatorname{Cov}(W_j(s),W_k(t))=\Gamma_{jk}\min(s,t).
\]
Zero or singular limits are allowed. If each comparison is
a coordinate-resampling pair, the same joint conclusion holds
for its label histories.
\end{theorem}
\begin{proof}
The finite overlap matrices are Gram matrices of the vectors
$(\ind_{\{i\in D_j\}})_{i=1}^q/\sqrt q$, so their limit is
positive semidefinite.
Poissonize the background at mean $M$, keep all membership
counts fixed, and rescale to circumference $q$.
Write $L=\log(M/q)$ and use the same probe capture process
$G_B(Z_i,t)$ for every comparison.
By Lemma~\ref{lem:finite-type-remainder},
\[
 H_j(t)-\EE H_j(t)
   =\sum_{i\in D_j}
        \{G_B(Z_i,t)-\EE[G_B(Z_i,t)\mid B]\}
      +o_{L^2_{\rm max}}(\sqrt{qL})
\]
on the early branch. The background conditional-mean term
is $|D_j|/q$ times the centered integrated capture, hence
also has maximal $L^p$ norm at most $C_p\sqrt q$.
The displayed remainder estimate holds jointly over the
fixed number of comparisons.

For finite-dimensional limits, form an arbitrary fixed
linear combination of comparisons and times.
Each membership type contributes a fixed linear combination
of single-probe capture coordinates.
The weighted conditional-variance argument in
Theorem~\ref{thm:all-sparse-history-FCLT} applies to each of
these finitely many combinations.
Independence of the probes given $B$ makes their variances
add. Equation~\eqref{eq:single-probe-clock-covariance} identifies
their covariance divided by $qL$ as tending to
$4v_{\rm H}\Gamma_{jk}\min(s,t)$.
Its errors are uniform over the finite type list; individual
type counts need not diverge or have positive limiting proportions.
The expected Lindeberg error is still $C_m/q$.
For $q<D_m\log(L+2)$ use the simultaneous growing-root
comparison. Its independent line captures carry the same
membership indicators, giving the same overlap covariance
by their cumulant bounds.

Each coordinate has the fourth increment and intra-band
maximal estimates from the preceding functional theorem,
with $|D_j|\le q$ in place of the number of summands.
Together with the finite-type remainder bound this proves
joint tightness and normalized maximal fourth-moment bounds.
The partial-record coupling is already simultaneous over
all probe subsets. Its bounded-density and bad-event estimates
therefore transfer the vector path to exact original count.
The maximal fourth moments align every mean path uniformly,
as in Theorem~\ref{thm:all-sparse-history-FCLT}.

For the late branch of comparison $j$, treat
$B+(S_j\cap T_j)$ as its common background; its size is
$M+|S_j\cap T_j|$, and its differing probes have total
count $|D_j|\le q$.
The proof of \eqref{eq:finite-late-history-maximal} also
applies at cutoff $1/q$ when the actual probe count is
smaller than $q$: use circumference $q$ in its Poisson
estimate, and $r_*=1/\sqrt{(M+|S_j\cap T_j|)q}$ in the
partial-record reservoir. At most $q$ gaps can be occupied,
and the deterministic late cost is still at most $q+1$.
Thus its centered maximal $L^p$ norm is $C_{p,m}\sqrt q$.
A finite union of these bounds makes the whole late vector
negligible after normalization, proving
\eqref{eq:joint-growing-history}.

For coordinate-resampling comparisons, each cumulative
label correction is bounded by its completed correction,
whose second moment is at most $32$.
A finite sum of these bounds transfers the joint limit;
no independence between the label corrections is needed.
\end{proof}

\begin{corollary}[Brownian-sheet marginals in the edit fraction]
\label{cor:edit-fraction-sheet-marginals}
Take the nested additions
$P(u)=\{1,\ldots,\lfloor qu\rfloor\}$, $0\le u\le1$,
and compare $B$ with $B+P(u)$ on the complete clock above.
For every finite collection of edit fractions, the normalized
centered histories converge jointly as separation paths to
$W(u,t\wedge1)$, where
\[
 \operatorname{Cov}(W(u,s),W(v,t))
                   =\min(u,v)\min(s,t).
\]
Thus the edit-fraction marginals have the Brownian-sheet law.
This statement does not assert tightness jointly in the
continuous edit-fraction parameter.
Section~\ref{sec:early-history-sheet} proves that tightness
on the early branch.
\end{corollary}
\begin{proof}
The differing sets are nested, and
$|P(u)\cap P(v)|/q\to\min(u,v)$.
Apply Theorem~\ref{thm:joint-growing-history}.
\end{proof}

\subsection{A Brownian sheet for nested sparse additions}
\label{sec:early-history-sheet}

\status{Working proof: tightness in both edit fraction and early separation}

For a signed measure $\mu$ on $[0,1]^d$, write
$\|\mu\|_\Box=\sup_{x\in[0,1]^d}|\mu(\prod_j[0,x_j])|$.
All fields below are separable: their suprema equal suprema
over a fixed dense set together with one-sided limits.

\begin{lemma}[The rectangle version of the measure maximal inequality]
\label{lem:rectangle-measure-maximal}
For independent signed random measures $\mu_1,\ldots,\mu_n$
on $[0,1]^d$ and $p\ge2$,
\[
 \left\|\left\|\sum_i(\mu_i-\EE\mu_i)\right\|_\Box\right\|_p
 \le C_{p,d}\sqrt n\,
                  \max_i\|\|\mu_i\|_{\rm TV}\|_p.
\]
The renewal and cyclic conclusions of
Lemmas~\ref{lem:random-measure-maximal}
and~\ref{lem:cyclic-measure-maximal} hold with
$\|\cdot\|_\Box$ in place of the cumulative norm.
\end{lemma}
\begin{proof}
Symmetrize, split into Jordan parts, and sample one atom from
each positive part as in Lemma~\ref{lem:random-measure-maximal}.
Conditional on the masses $W_i$ and atoms $U_i$, the remaining
field is
$\sum_i\varepsilon_iW_i\ind_{\{U_i\le x\}}$,
where the inequality is coordinatewise.
Place these weights on the finite grid formed by their
coordinate values. Distinct grid cells contain disjoint
collections of independent signs.
In dimension two, the vector of column partial sums, as rows
are added, is a martingale. Its maximum absolute coordinate
is a nonnegative submartingale. The martingale maximal
inequality first removes the row maximum, at a factor
depending only on $p$. After all rows have been added, the
column sums have independent centered increments, so a second
application removes the column maximum. Iterating this argument
in $d$ coordinates gives
\[
 \EE_\varepsilon\sup_x
 \left|\sum_i\varepsilon_iW_i\ind_{\{U_i\le x\}}\right|^p
 \le C_{p,d}\left(\sum_iW_i^2\right)^{p/2}.
\]
The last expectation over masses is bounded as in the
one-dimensional lemma. The rectangle norm is convex, so
the additional block-index maximal inequality in its renewal
proof still applies. The cyclic proof uses only that line
bound, boundary variation moments, and the bounded density
of half-circle records; it is unchanged.
\end{proof}

\begin{lemma}[The remainder uniformly over all priority thresholds]
\label{lem:priority-remainder-maximal}
On a circle of circumference $q\ge8$, take a Poisson
background $B$ of intensity $\alpha\ge1$ and exactly $q$
independent uniform probes with independent uniform priorities.
Let $P(v)$ contain the probes of priority at most $v$, and put
\[
 R(v,r)=H(B,B+P(v);[0,r])
                 -\sum_{z\in P(v)}G_B(z;[0,r]).
\]
For each fixed $p\ge2$, uniformly in $\alpha$ and $q$,
\begin{equation}\label{eq:priority-remainder-maximal}
 \left\|\sup_{0\le v,r\le1}
                 |R(v,r)-\EE R(v,r)|\right\|_p\le C_p\sqrt q.
\end{equation}
The original Poisson parity may be fixed.
\end{lemma}
\begin{proof}
First use a marked Poisson probe process of intensity one.
Order its probes by priority. At a priority jump, take the
signed edge-time measure of the corresponding remainder
increment. Placing that measure at its priority defines a
signed measure on $[0,1]^2$, whose anchored cumulative
function is $R(v,r)$. Its variation is at most the sum of
the variation norms of the individual remainder increments.

Use occupied blocks for the full probe process, at cutoff
one. Their bivariate marks have uniformly bounded fixed
variation moments. Here is the required strengthening of
Theorem~\ref{thm:line-block-remainder-moments}.
In one separation band, the variation of the jump at a probe
$z$ is bounded by its insertion cost into its priority
predecessor input plus its insertion cost into $B$.
Under the marked probe Palm law, a probe of priority $v$
has a Poisson predecessor background of intensity $\alpha+v$.
The uniform single-band moments and the probe Palm exponential
moments of the full block size therefore give all fixed
moments of the sum of these costs, by exactly the
mass-transport and H\"older calculation in that theorem.
The same calculation allows a small exponential block-length
weight, by using higher moments.

If the full probes occupy distinct nonadjacent background
gaps at a finer cutoff $b$, every subset has zero remainder
through $b$. Thus the entire bivariate measure below $b$
vanishes on that event, including all its priority jumps.
The bad-block probability is $O(b)$. H\"older gives the
same band variation bound $C_pb^{1/(2p)}$ as before.
The microscopic initial strip has this bound too.
Dyadic summation gives uniform full variation moments.
The final empty gap is independent of the preceding occupied
run mark and smooths its exponentially weighted moment
densities. Lemma~\ref{lem:rectangle-measure-maximal} and the
block renewal identities consequently give
\eqref{eq:priority-remainder-maximal} before conditioning
the probe count.

For completeness the exceptional-sector estimate also
survives this strengthening. The global band variation is
bounded by the sum of the priority insertion and single-probe
costs, with $L^p$ norm at most $C_pq$. At cutoff $b$ its
nonzero event has probability at most $\min(1,Cqb)$.
Below the microscopic scale use the corresponding initial
component bound. Summing these band estimates gives
\begin{equation}\label{eq:priority-remainder-polynomial}
 \|\|R\|_{\rm TV}\|_p\le C_pq(1+\log(q+2)),
\end{equation}
uniformly in $\alpha$. The no-empty-block and empty/singleton
background sectors have probability $Ce^{-cq}$, so the
cyclic exceptional contribution is harmless.

We now condition the total probe count without revealing
it in the whole remainder field. In a deterministic quarter
circle, partition the complete background cutoff-one gaps
into consecutive disjoint triples. Call the middle gap
eligible when all three lengths are in $[1,2]$, both guard
gaps contain no probe, and the middle gap contains at most
one probe. Reveal the complete background, the eligibility
flags, and all probe records except those in eligible middle
gaps. Denote these data by $\mathcal S$.
In an eligible center, hide its zero-or-one count and, if
present, the probe's position and priority.

The field $R(v,r)$ is $\mathcal S$-measurable. Indeed an
eligible center is either empty or an isolated one-probe
run, whose remainder is identically zero for every priority
threshold and every separation through one.
The two empty guards separate all other runs under every
subset, by Lemma~\ref{lem:history-occupied-clusters}.
Changing the hidden count does not change their remainders.
This is why priorities, rather than their ranks, are retained
during count conditioning.

Except with probability $Ce^{-cq}$ there are at least $cq$
eligible centers. To verify this cyclic assertion, expose
the background gap records covering the quarter circle.
When their combined length is at most $q/2$, their density
relative to stationary line records is bounded by a constant:
the outside renewal factor and the origin length bias are
the same as in Lemma~\ref{lem:cyclic-measure-maximal}.
Failure requires a boundary gap of length at least $q/8$
and has exponentially small probability.
On the line, after the initial boundary gap, the first
$\lfloor c_0q\rfloor$ iid gaps all fit in the quarter circle,
except with exponentially small probability, for a sufficiently
small fixed $c_0>0$. This follows from their uniform exponential
length moment and the exponential tail of the initial
length-biased gap. Each disjoint triple is eligible with a
uniformly positive probability: the gap density is uniformly
positive on $[1,2]$, and the Poisson void and zero-or-one
probabilities are bounded below on those lengths.
The triple events are independent, so their elementary
Bernoulli exponential bound gives the claimed abundance.
The bounded record density transfers it to the circle.
The finitely many bounded values of $q$ may be absorbed
by enlarging constants.

Conditional on $\mathcal S$, the hidden counts are independent
Bernoulli variables with parameters $Y/(1+Y)\in[1/2,2/3]$.
For $m\ge cq$ such variables, the modulus of their product
characteristic function is at most
$\exp\{-cm\operatorname{dist}(t,2\pi\mathbb Z)^2\}$.
Fourier inversion bounds the maximal remaining count mass
by $C/\sqrt q$. Since $\PP(\operatorname{Pois}(q)=q)\ge c/\sqrt q$,
conditioning on the exact total $q$ changes the
$\mathcal S$-density by a bounded factor on the good event.
On its complement, global conditioning costs only $C\sqrt q$;
\eqref{eq:priority-remainder-polynomial}, H\"older, and its
exponential probability control every fixed moment there.
Transfer the centered Poisson maximal bound and then
subtract the new mean field, at a factor at most two.
This proves \eqref{eq:priority-remainder-maximal}.
Conditioning the independent original parity costs a bounded
factor; recentering once more gives the same conclusion.
\end{proof}

\begin{lemma}[The remainder uniformly over deterministic edit prefixes]
\label{lem:ranked-prefix-remainder}
For $q$ iid uniform probes in a fixed index order, let
\[
 R_k(r)=H(B,B+\{Z_1,\ldots,Z_k\};[0,r])
                  -\sum_{i=1}^kG_B(Z_i;[0,r]),\qquad 0\le k\le q.
\]
In the Poisson background of the preceding lemma,
\begin{equation}\label{eq:ranked-prefix-remainder}
 \left\|\max_{0\le k\le q}\sup_{0\le r\le1}
                   |R_k(r)-\EE R_k(r)|\right\|_p\le C_p\sqrt q.
\end{equation}
For an exact background of size $M$ on the unit circle,
the same bound holds through separation $1/q$, uniformly
for $q\le M/\alpha_0$, with a sufficiently large absolute
$\alpha_0$. Every fixed $p\ge2$ is allowed.
\end{lemma}
\begin{proof}
Work first with the exact-$q$ priority model and write
$m(v,r)=\EE R(v,r)$. Its mean is uniformly Lipschitz:
\begin{equation}\label{eq:priority-mean-lipschitz}
 \sup_r|m(v,r)-m(w,r)|\le Cq|v-w|.
\end{equation}
Indeed a single addition has expected remainder-increment
variation at most $C$, for every old probe count $0\le k<q$.
For $k=0$ the increment is zero. For $k\ge1$, rescale
circumference $q$ to $T'=\max(k,8)\le q$.
The old probe density belongs to $[1/8,1]$, the background
intensity is $\alpha q/T'\ge1$, and the target cutoff
$T'/q$ is at most one.
The probe-only part of Lemma~\ref{lem:two-count-reservoirs},
in its variation form proved in
Theorem~\ref{thm:early-remainder-maximal}, gives the required
bound. The mean number of priorities crossed between $w$
and $v$ is $q|v-w|$, proving
\eqref{eq:priority-mean-lipschitz}.

Let $U_{(k)}$ be the ordered priorities and put $U_{(0)}=0$.
The standard exponential-spacing representation can be checked
directly by the constant density on the spacing simplex:
$U_{(k)}=S_k/S_{q+1}$, where $S_k$ sums $k$ independent
mean-one exponential variables. The martingale maximal
inequality and the centered even-moment expansion give
$\|\max_{k\le q+1}|S_k-k|\|_p\le C_p\sqrt q$.
On $S_{q+1}\ge q/2$ this bounds the spacing ratios, and its
complement has exponential probability. Hence
\[
 \left\|\max_{0\le k\le q}|U_{(k)}-k/q|\right\|_p
                                     \le C_pq^{-1/2}.
\]
The probe positions in priority order are still iid uniform.
Evaluate the whole-field bound
\eqref{eq:priority-remainder-maximal} at $U_{(k)}$, and use
\eqref{eq:priority-mean-lipschitz} to replace
$m(U_{(k)},r)$ by $m(k/q,r)$, at cost $C_p\sqrt q$.
Subtracting the actual ranked-prefix mean costs at most
another factor two. This proves
\eqref{eq:ranked-prefix-remainder}.

For the exact original count, use the partial-record
reservoir at $r_*=1/\sqrt{Mq}$, as in
Theorem~\ref{thm:early-remainder-maximal}.
Its data determine all subset histories and all single-probe
terms, hence the entire ranked-prefix remainder field.
Its good-event density is bounded for $q\le M/\alpha_0$.
The uncentered variation moments of that field are polynomial
in $q$, uniformly in the Poisson background intensity:
apply global probe-count conditioning to
\eqref{eq:priority-remainder-polynomial}; ranking only pushes
the priority coordinate forward and cannot increase variation.
Thus the bad-event contribution, after original-count
conditioning, is at most
$C_p\sqrt M(q+1)^{C_p}e^{-c_p\sqrt{Mq}}$.
This is harmless for every fixed $p$.
The finitely many $q<8$ are treated on circumference eight,
with the remaining bounded interval controlled by packing,
as in that earlier theorem. Recenter to obtain the claim.
\end{proof}

\begin{theorem}[The early Brownian sheet for nested additions]
\label{thm:early-history-sheet}
Let $B$ consist of $M$ iid uniform circle points, independently
of $q$ iid uniform probes, where $q\to\infty$ and $q=o(M)$.
Put
\[
 L=\log(M/q),\qquad
 r_M(t)=M^{-1}(M/q)^t,\qquad
 H_M(u,t)=H(B,B+\{Z_1,\ldots,Z_{\lfloor qu\rfloor}\};
                                                   [0,r_M(t)]).
\]
Center each value by its own expectation and divide by
$\sqrt{4qv_{\rm H}L}$, obtaining $Z_M(u,t)$.
Let $\widetilde Z_M$ be its bilinear interpolation on
$u=i/q$, $t=j/\lceil L\rceil$. Then
\begin{equation}\label{eq:early-history-sheet}
 \widetilde Z_M\Longrightarrow W
       \quad\hbox{in }C([0,1]^2),\qquad
 \|Z_M-\widetilde Z_M\|_\infty\longrightarrow0
       \quad\hbox{in probability},
\end{equation}
where $W$ is the standard Brownian sheet,
$\EE W(u,s)W(v,t)=\min(u,v)\min(s,t)$.
Either background parity is allowed.
This theorem ends at separation $1/q$; a late-path bound
uniform over all edit fractions is a separate assertion.
Section~\ref{sec:complete-history-sheet} supplies that bound.
\end{theorem}
\begin{proof}
Poissonize the original background at mean $M$, optionally
condition its parity, and rescale to circumference $q$.
The intensity is $\alpha=M/q$ and the cutoffs are
$b(t)=\alpha^{t-1}$. Define the conditional empirical sums
\[
 X_k(t)=\sum_{i=1}^k
       \{G_B(Z_i;[0,b(t)])-\EE[G_B(Z_i;[0,b(t)])\mid B]\}.
\]
The difference between the centered history field and
$X_{\lfloor qu\rfloor}$ is the centered prefix remainder
plus $(\lfloor qu\rfloor/q)(A(b(t))-\EE A(b(t)))$.
Lemma~\ref{lem:ranked-prefix-remainder} and
Corollary~\ref{cor:integrated-capture-maximal} bound their
uniform $L^p$ norms by $C_p\sqrt q$.
They vanish after normalization by $\sqrt{qL}$.

Here are the two-parameter tightness details for $X$.
For any fixed integer $p>2$, conditional independence,
the centered moment expansion, and conditional Jensen give
\begin{align*}
 \EE|X_j(t)-X_i(t)|^{2p}
       &\le C_p(j-i)^p L^p,\qquad i<j,\\
 \EE|X_k(t)-X_k(s)|^{2p}
       &\le C_pk^p(1+L|t-s|)^p,\qquad s<t.
\end{align*}
The single-probe inputs are
Lemma~\ref{lem:single-probe-log-increments} and its bounded
microscopic strip. In the conditional expansion, every
nonzero term uses each probe index at least twice; H\"older
and Jensen bound all products by the corresponding
unconditional $2p$th moment. There are at most $C_pn^p$
terms when $n$ probes are summed. For large $M$, $L\ge1$.

On the two grids, division by $(qL)^p$ bounds coordinate
increments by $C_p|\Delta u|^p$ and $C_p|\Delta t|^p$.
Bilinear interpolation preserves the combined bound
$C_p(|\Delta u|+|\Delta t|)^p$: within a single mesh interval
the $2p$th power of the interpolation fraction absorbs its
shorter length, and crossing an interval boundary splits
into at most two such pieces and an intervening grid increment.
This gives tightness and uniformly bounded maximal $2p$th
moments for the interpolated fields.
Explicitly, on a dyadic grid of level $h$ there are $O(4^h)$
neighbour increments, so the $L^{2p}$ norm of their maximum
is at most $C_p2^{-h(1/2-1/p)}$.
Summing this geometric bound proves both equicontinuity and
the maximal moment assertion; the value at $(0,0)$ is zero.

We also control the discrepancy from interpolation.
For one probe, the same one-dimensional log-grid argument
and the single-band measure maximal bound give
\[
 \EE\sup_t
 |G_B(Z_i;[0,b(t)])-\EE[G_B(Z_i;[0,b(t)])\mid B]|^{2p}
                                                       \le C_pL^p.
\]
Consequently the maximum normalized jump in the edit
coordinate has $2p$th moment at most $C_pq^{1-p}$.
Within one time-grid interval, each probe contributes a
positive edge-time measure with uniformly bounded $2p$th
variation moment, since the separation ratio is at most $e$.
Conditionally on $B$ these measures are independent.
The cumulative norm of their centered partial sums is a
submartingale in the probe index.
The block-index maximal argument of
Lemma~\ref{lem:random-measure-maximal} therefore gives
\[
 \EE\max_{k\le q}\sup_{t\text{ in one interval}}
 |X_k(t)-X_k(t_{\rm left})|^{2p}\le C_pq^p.
\]
Summing over $O(L)$ intervals and normalizing gives
$C_pL^{1-p}\to0$. Both interpolation errors vanish.
The uniformly negligible remainder and integrated-capture
fields transfer this conclusion and the maximal moment
bounds to the centered Poisson history field.
Its finite-dimensional limits are those of
Corollary~\ref{cor:edit-fraction-sheet-marginals}.
This proves \eqref{eq:early-history-sheet} for that model.

Finally Lemma~\ref{lem:relative-history-depoisson} at
$r_*=1/\sqrt{Mq}$ compares all subset relative histories
simultaneously in total variation. The Poisson uniform
maximal moments transfer on its good event by the bounded
record density. On the complement, global original-count
conditioning costs $C\sqrt M$, while the whole raw field
is bounded by $C(N_B+q)$ and its mean by $C(M+q)$.
Poisson moments and the probability $Ce^{-c\sqrt{Mq}}$
make its normalized contribution negligible.
The coupling and these maximal moments imply
\[
 \sup_{u,t}
 \frac{|\EE_{\rm exact}H_M(u,t)-\EE_{\rm Pois}H_M(u,t)|}
                         {\sqrt{qL}}\longrightarrow0.
\]
Thus the exact-count fields have the same limits with their
own mean centers. Interpolation is a contraction in supremum
norm, so it transfers as well. This proves the theorem.
\end{proof}

\subsection{The early sheet for a nested coordinate-resampling path}
\label{sec:resampling-history-sheet}

\status{Working proof: fixed initial sample and nested replacements}

\begin{lemma}[Single-band insertion moments with finitely many forced points]
\label{lem:forced-point-band-moments}
Let $B$ be a stationary Poisson line of intensity $\alpha$,
or a Poisson circle of circumference $T\ge8$ and intensity
$\alpha\ge1$. Add at most $h$ deterministic points, then
insert one further deterministic point. For each fixed
$h$ and $p\ge1$, the resulting history distance in a band
$(b,2b]$, with $2b\le1$ on the circle, has $p$th moment
at most $C_{p,h}$, uniformly over the point locations and
the intensity. The whole initial strip $[0,b]$ has the
same bound when $\alpha b$ is bounded.
Null coincidences and distance ties may be excluded.
\end{lemma}
\begin{proof}
First suppose $\alpha b\ge1/2$. In the unaugmented
background configuration at cutoff $b$, a gap longer than
$4b$ that contains none of the forced points remains a
barrier through cutoff $2b$ in every relevant copy.
Its interior original edges are protected, and no accepted
edge can cross its empty interval before that cutoff,
even if a boundary endpoint is removed.
A forced point can occupy at most one such background gap.
Starting from the new point's gap, expose gaps in both
directions until $h+1$ gaps longer than $4b$ have appeared
on each side. At least one on each side contains none of
the other $h$ forced points. These two barriers confine
every new discrepancy relevant to the band.

The number of exposed background endpoints has bounded
fixed moments: each outside gap is long with a probability
uniformly bounded below, and the line gap sequence is iid.
The origin gap is length biased but does not add endpoints
to the outside strings. The circle version follows from
the same bounded outside-renewal comparison as the
single-band proof in
Theorem~\ref{thm:line-block-remainder-moments}.
A wrapping string with many endpoints has an exponential
tail; the finitely many cases with fewer than $2h+3$
endpoints have a deterministic bound. One may also see the
tail by using the occupied short-gap subdensity until
$h+1$ long gaps have occurred; a fixed number of intervening
long gaps adds only a polynomial factor to its exponential
tail. Finite perturbations change the active count inside
the protected interval by at most $h+1$.
The band contains at most a constant times this many
active endpoints from the two copies, proving the bound.

If $\alpha b<1/2$, use original Poisson gaps longer than
$8b$ instead, for the entire history through $2b$.
Their probabilities are bounded below, and the same
$h+1$-barrier argument applies to the original point counts.
For the stated initial-strip assertion the lower bound
on this long-gap probability depends only on the upper
bound for $\alpha b$. This proves all claims.
\end{proof}

\begin{lemma}[A uniform remainder for separate deletion and addition prefixes]
\label{lem:resampling-prefix-remainder}
On a circle of circumference $q=2k$, $k\ge4$, take a
Poisson background $B$ of intensity $\alpha\ge1$, and
independent uniform old and fresh probe families
$P=\{P_1,\ldots,P_k\}$ and $Q=\{Q_1,\ldots,Q_k\}$.
For $0\le i,j\le k$ let
\[
 A=B+P,\qquad
 A_{i,j}=B+\{P_{i+1},\ldots,P_k\}+\{Q_1,\ldots,Q_j\},
\]
and subtract from their relative history through $r$ the
single-probe costs of $P_1,\ldots,P_i,Q_1,\ldots,Q_j$.
Call the result $R_{i,j}(r)$. For every fixed $p\ge2$,
\begin{equation}\label{eq:resampling-prefix-remainder}
 \left\|\max_{i,j\le k}\sup_{0\le r\le1}
              |R_{i,j}(r)-\EE R_{i,j}(r)|\right\|_p
                                                   \le C_p\sqrt q.
\end{equation}
The original Poisson parity may be fixed.
The bound also transfers to an exact background of size
$M$ on the unit circle through cutoff $1/q$, for
$q\le M/\alpha_0$.
\end{lemma}
\begin{proof}
First give the old and fresh probes independent Poisson
intensities $1/2$, each with uniform priority marks.
Write $P_{\le v}$ and $Q_{\le w}$ for their priority prefixes.
Define
\[
 R(v,w,r)=H(B+P,B+P_{>v}+Q_{\le w};[0,r])
       -\sum_{z\in P_{\le v}\cup Q_{\le w}}G_B(z;[0,r]).
\]
Use occupied blocks for the full union $P\cup Q$.
Within one block, the field is the cumulative function of
a signed measure on $[0,1]^3$: put the one-threshold
increments of $R(v,0,r)$ on the face $w=0$, those of
$R(0,w,r)$ on the face $v=0$, and the mixed priority
increments in the interior. Here $R(0,0,r)=0$.

We verify its uniform block variation moments.
On either face, a priority jump is bounded in variation
by one insertion or deletion cost in the current input
plus the corresponding single-probe cost in $B$.
Under the one-probe Palm law the required current input,
with the tagged point removed, is Poisson.
For a mixed jump at an old and a fresh priority, fix a
physical edge and write $a$ for its initial matching
indicator and $f(v,w)$ for its current indicator.
The identity
\[
 |a-f|=a+(1-2a)f
\]
shows that the mixed history difference is
$(1-2a)\Delta_v\Delta_w f$; the single-probe terms have
zero mixed difference.
Its absolute value is at most the sum of the two old-point
insertion differences, with the fresh tagged point respectively
present and absent. Under the two-probe Palm law their
backgrounds are Poisson with at most one additional forced
point. Lemma~\ref{lem:forced-point-band-moments} bounds all
fixed moments uniformly, even when the two tagged positions
are fixed.

The mass-transport argument can now sum over pairs in a
block without losing a volume factor. Under the two-tag
Palm law with tag separation $z$, for sufficiently small
$a_0>0$,
\begin{equation}\label{eq:two-tag-block-weight}
 \EE^{0,z}\!\left[
  \ind_{\{\text{tags in the same block }D\}}
             e^{a_0(L_D+K_D)}\right]\le Ce^{-c|z|}.
\end{equation}
To check this, a common gap containing both tags has an
exponential length tail. Otherwise every intervening gap
must be occupied. Its tilted subdensity is
$f_\alpha(y)\EE[e^{a_0(y+K)};K>0\mid Y=y]$, which has mass
strictly less than one for small $a_0$, uniformly in $\alpha$.
The intervening convolution sum therefore has an exponential
tail. The two endpoint gap biases and the runs continuing
to the nearest empty gaps have uniform exponential moments
and preserve that estimate, after decreasing $a_0$.
This also proves the estimate with any fixed polynomial
in $L_D,K_D$ in place of the exponential.

For example the $p$th power of the mixed variation is
bounded by $K_D^{2p-2}$ times the sum of the $p$th powers
of its pair costs. Transport this sum to two tagged probes,
integrate their relative position $z$, and use H\"older,
\eqref{eq:two-tag-block-weight}, and the forced-point moment
bound. The integral is finite uniformly. The face terms
use the one-tag calculation already proved.
This gives every fixed band variation moment, also with
a small exponential block-length weight.

If all full probes occupy separated finer gaps, every
subset comparison has zero remainder through that cutoff,
including comparisons with old probes common to both inputs.
Hence the entire three-dimensional measure vanishes there.
The $O(b)$ bad-block probability gives the same dyadic
variation gain as in Lemma~\ref{lem:priority-remainder-maximal}.
The initial strip is covered by the forced-point initial
bound. The final empty gap smooths the resulting full
variation-moment densities. The rectangle and cyclic maximal
lemmas give centered field norm $C_p\sqrt q$.
For the exceptional sector, global pair summation bounds a
band variation in $L^p$ by $C_p(q+1)^2$.
The full-probe fine-pair event has probability
$\min(1,Cqb)$, so dyadic summation gives the raw bound
$C_p(q+1)^2(1+\log(q+2))$. This polynomial bound controls
the exponentially rare exceptional sectors.

Condition the two probe counts through the isolated-center
reservoir of Lemma~\ref{lem:priority-remainder-maximal}.
An eligible center contains at most one probe of either
color and has empty guards. Its remainder is zero for
every $v,w,r$, whether that probe is fresh, old and common,
or old and removed. Thus the entire field remains measurable
when the center count and all its marks are hidden.
The hidden count vector has atoms $0,e_P,e_Q$, with probabilities
proportional to $1,Y/2,Y/2$. All three are bounded below
for $Y\in[1,2]$.
There are $cq$ independent centers except with exponential
probability. Their characteristic functions give maximal
vector point mass $C/q$, matching
$\PP(\operatorname{Pois}(k)=k)^2$.
The good-event density is therefore bounded, and the
polynomial variation moments control its complement.
Recenter to obtain the uniform priority-field bound with
both counts exactly $k$.

It remains to rank the two priorities separately.
The mean field is $Ck$-Lipschitz in each threshold, uniformly
in the other threshold and separation.
Indeed changing one old point from common to removed changes
the remainder by the difference between two local increments:
insert that point in both inputs, or insert it only in
the initial input. Adding one fresh point is a single
exclusive increment.
The finite-type local-increment argument in the proof of
Lemma~\ref{lem:finite-type-remainder} bounds their expected
variation uniformly for every allocation of the other
$q-1$ probes to the four membership types.
Unused probes may be retained in that type count; the
total density is $(q-1)/q$, bounded away from zero.
Conditioning on how many priorities are crossed gives the
claimed Lipschitz bound.
Each color's ordered priorities differ uniformly from
$i/k$ by $O_{L^p}(k^{-1/2})$.
Evaluate the whole field at the two order statistics and
recenter, as in Lemma~\ref{lem:ranked-prefix-remainder}.
This proves \eqref{eq:resampling-prefix-remainder}.
Conditioning original parity costs a bounded factor.

Finally the partial-record reservoir determines every one
of these relative histories and their single-probe terms.
Its good-event density and exponentially small complement,
together with the polynomial raw variation bound just proved,
transfer \eqref{eq:resampling-prefix-remainder} to exact $M$
at $r_*=1/\sqrt{Mq}$. This is the same original-count
argument as in Lemma~\ref{lem:ranked-prefix-remainder}.
\end{proof}

\begin{theorem}[The early Brownian sheet for nested coordinate resampling]
\label{thm:resampling-history-sheet}
Start with $N$ iid uniform circle points. Select $k$ coordinates
independently of their positions, order them, and successively
replace them by independent fresh uniform points.
Suppose $k\to\infty$ and $k=o(N)$, and put
\[
 M=N-k,\qquad q=2k,\qquad L=\log(M/q),\qquad
 r(t)=M^{-1}(M/q)^t.
\]
Let $H_N(u,t)$ be the history distance from the fixed
initial sample after the first $\lfloor ku\rfloor$ replacements,
through separation $r(t)$.
Center by its exact expectation and divide by
$\sqrt{8kv_{\rm H}L}$.
Its bilinear interpolation on $u=i/k$, $t=j/\lceil L\rceil$
converges in $C([0,1]^2)$ to a standard Brownian sheet,
and the raw field differs uniformly from that interpolation
by a quantity tending to zero in probability.
The result holds for both spatial and coordinate-label
history distance, centered separately.
One may replace $L$ in the normalization by $\log(N/k)$.
The cutoff clock still ends at $1/(2k)$.
Corollary~\ref{cor:complete-history-sheet} extends it through completion.
\end{theorem}
\begin{proof}
By exchangeability the unselected $M$ points form an
independent background $B$, and the old and fresh selected
points form two independent ordered families $P,Q$ of size $k$.
The initial sample is $B+P$ and the sample after $j$ replacements
is $B+\{P_{j+1},\ldots,P_k\}+\{Q_1,\ldots,Q_j\}$.
In the Poisson-background model, the preceding lemma makes
the centered remainder uniformly $O_{L^p}(\sqrt q)$.
The integrated-capture mean term is $(2j/q)(A-\EE A)$
and has the same maximal bound.

The leading field is therefore the conditional empirical
sum of the first $j$ independent pairs of old and fresh
single-probe captures. Its index and time increment
moments are exactly those in the proof of
Theorem~\ref{thm:early-history-sheet}, with two independent
captures in each summand.
That proof gives two-dimensional tightness, maximal moments,
and the vanishing interpolation error.
The finite-dimensional covariance follows either by the
conditional variance calculation or by
Theorem~\ref{thm:joint-growing-history}: two edit prefixes
share $2\min(i,j)$ differing probe identities, so division
by $q=2k$ gives $\min(u,v)$.
The partial-record TV comparison and maximal moment argument
align the whole exact-$M$ mean surface and transfer the
spatial sheet limit.

For labels, write $D_j(r)$ for spatial minus label history
distance at prefix $j$. For each fixed $j$, this correction
is nonnegative and nondecreasing in $r$, and
$\EE D_j(1/2)^2\le32$ by
Lemma~\ref{lem:history-label-correction}.
Its permutation proof also applies to even $N$: it uses
only disjointness of matching edges, not the existence of
a single unmatched point.
A union bound over $0\le j\le k$, with no independence
assumption, gives
\[
 \PP\!\left(\max_{j\le k}\sup_r
       |D_j(r)-\EE D_j(r)|>\varepsilon\sqrt{kL}\right)
                                      \le \frac{C}{\varepsilon^2L}
                                      \longrightarrow0.
\]
This transfers the sheet and interpolation assertions to
the separately centered label field.
Finally $L/\log(N/k)\to1$, proving the normalization claim.
\end{proof}

\subsection{The complete history sheet}
\label{sec:complete-history-sheet}

\status{Working proof: the late path uniformly over edit prefixes}

\begin{lemma}[A variation-sensitive bound for bounded cumulative fields]
\label{lem:bounded-field-entropy}
Let $f_i$ be the anchored cumulative fields of independent
signed random measures $\mu_i$ on $[0,1]^d$.
Fix $d$ and $p\ge2$, and suppose
\[
 \|f_i\|_\infty\le A,\qquad
 \max_i\|\|\mu_i\|_{\rm TV}\|_p\le V.
\]
Then
\begin{equation}\label{eq:bounded-field-entropy}
 \left\|\sup_x\left|\sum_{i=1}^n(f_i(x)-\EE f_i(x))\right|
                                                   \right\|_p
 \le C_{p,d}A\sqrt n\,\sqrt{\log(2+V/A)}.
\end{equation}
The zero case $A=0$ is interpreted separately.
\end{lemma}
\begin{proof}
Symmetrize. Conditional on the fields, the Rademacher sum
has subgaussian increments in the metric
$d_f(x,y)^2=\sum_i|f_i(x)-f_i(y)|^2$.
Let $\nu_j$ be the $j$th coordinate projection of
$\sum_i|\mu_i|$, and put $V_*=\sum_i\|\mu_i\|_{\rm TV}$.
Since every field increment is bounded by $2A$,
\[
 d_f(x,y)^2\le
   2A\sum_{j=1}^d\nu_j((x_j\wedge y_j,x_j\vee y_j]),
 \qquad \operatorname{diam}(d_f)\le2A\sqrt n.
\]
Partition each control measure into intervals of mass
at most $\varepsilon^2/(2Ad)$. Retain both sides of atoms
larger than that mass; there are at most the same order
of such atoms. The resulting coordinate grid gives an
$\varepsilon$-net of size at most
$[C_d(1+AV_*/\varepsilon^2)]^d$.

Take nets at radii $2A\sqrt n\,2^{-h}$.
The conditional $L^p$ norm of the maximum of $m$
subgaussian increments of scale $\varepsilon$ is at most
$C\varepsilon\sqrt{p+\log(2m)}$.
This follows directly from the exponential sign moment
bound and integration of the union tail.
Chaining successive nets and summing $2^{-h}\sqrt{a+h}$
therefore gives
\[
 \left\|\sup_x\left|\sum_i\varepsilon_i f_i(x)\right|
                           \right\|_{L^p(\varepsilon)}
 \le C_{p,d}A\sqrt n
               \sqrt{\log(2+V_*/(An))}.
\]
One anchor value is bounded by the ordinary Rademacher
moment inequality. The chain converges in the finite
dimensional space of vectors $(f_i(x))_{i\le n}$;
including one-sided limits handles atoms.
Finally $\|V_*/(An)\|_p\le V/A$.
For $K=\max(1,V/A)$ use
$\log(2+z)\le\log(2+K)+z/K$ to take the remaining
$L^p$ norm. This proves \eqref{eq:bounded-field-entropy}.
\end{proof}

\begin{lemma}[Ranking transfers a centered maximal bound]
\label{lem:ranking-maximal-transfer}
Give finitely many independent iid probe families independent
uniform priorities. Let $F(\boldsymbol v,t)$ be a separable
field determined by their priority prefixes and other data
independent of the priorities. Let $Y_{\boldsymbol i}(t)$
be the corresponding field in the ranked probe orders.
If
\[
 \left\|\sup_{\boldsymbol v,t}
       |F(\boldsymbol v,t)-\EE F(\boldsymbol v,t)|\right\|_p
                                                        \le C,
\]
then
$\|\max_{\boldsymbol i}\sup_t
       |Y_{\boldsymbol i}(t)-\EE Y_{\boldsymbol i}(t)|\|_p\le2C$.
\end{lemma}
\begin{proof}
The ranked probe positions and all the other data are
independent of the priority order statistics $\boldsymbol U$.
Put $m(\boldsymbol v,t)=\EE F(\boldsymbol v,t)$.
Evaluating the whole-field bound at $\boldsymbol U_{\boldsymbol i}$
gives the bound $C$ about the random center
$m(\boldsymbol U_{\boldsymbol i},t)$.
Conditional Jensen over $\boldsymbol U$, keeping the ranked
data fixed, gives the same bound about the deterministic
center $c_{\boldsymbol i}(t)=
\EE_{\boldsymbol U}m(\boldsymbol U_{\boldsymbol i},t)$.
Its difference from the actual mean has supremum norm
at most $C$, by Jensen once more. Recenter to conclude.
\end{proof}

\begin{lemma}[A local shielding certificate]
\label{lem:local-shielding-certificate}
Let $I$ be a cell of length $a\in[u/2,u]$ on a circle of
circumference $T$, with $u\ge1$.
Surround it by a window $W$ extending $\ell$ cells in each
direction, where $\ell\ge16$ and
$(2\ell+1)a\le T/8$.
There is an event determined only by the input in $W$
on which every matching edge with an endpoint in $I$,
through separation $2u$, is unchanged by arbitrary changes
outside $W$.
For the restriction to $W$ of homogeneous Poisson input
of intensity $\lambda\ge1$, with at most $h$ additional
fixed points, its failure probability is at most
$C_he^{-c_h\ell}$, uniformly in the added locations.
\end{lemma}
\begin{proof}
The certificate is the existence, on each side of $I$,
of a fully shielded interval $[a_0,b_0]\subset W$ of length
greater than $2u$, at distance at least $2u$ from $I$.
Here fully shielded means that the greedy process on its
interior points and two artificial endpoint points clears
every interior point by $2u$ without touching either endpoint.
It suffices to search endpoint pairs among the local input
points, so this is an event of the local records.

Its shielding implication does not require the endpoints
to survive in an exterior completion.
At each internal collision in the shielded scan, both
distances to the two still-present artificial endpoints
exceed the collision distance. Every point outside the
interval is farther still. Induction on these internal
collisions shows that no exterior completion can change
them, even if it removes an actual endpoint earlier.
All interior points consequently clear in the prescribed
way. No edge of length at most $2u$ can cross the whole
interval. The two such intervals isolate the cell's history,
so its edge vector agrees with the finite-window computation.

To bound failure, fill the exterior by independent homogeneous
Poisson input and first omit the $h$ additional points.
At cutoff $2u$, the background gap law, scaled by $2u$,
has a uniform exponential envelope since $2\lambda u\ge2$.
In each buffer there are $h+1$ disjoint complete cutoff
gaps except with probability $C_he^{-c_h\ell}$.
Indeed failure leaves a bounded number of gaps covering a
buffer of length proportional to $\ell u$, and the
length-biased gap tails and bounded outside-renewal density
give this estimate. The empty/singleton sector has the
same exponential bound. Alternating candidate gaps makes
their closed intervals disjoint if needed.
Each cutoff gap has length greater than $2u$ and its
interior is fully shielded through that cutoff.
At most $h$ of these intervals contain an added point,
so one on each side remains a certificate for the augmented
local records. These intervals lie inside $W$, making
the exterior filling irrelevant to their local validity.
\end{proof}

\begin{theorem}[Late histories uniformly over all edit prefixes]
\label{thm:late-prefix-maximal}
In either of the following Poisson-background models on
circumference $q\ge8$, with background intensity $\alpha\ge1$,
the centered late-history field after separation one has
maximal $L^p$ norm at most $C_p\sqrt q$, for every fixed
$p\ge2$:
\begin{enumerate}
\item $q$ iid probes, comparing the background with every
      initial probe prefix;
\item $q=2k$ probes in two families of size $k$, comparing
      $B+P$ with $B+P_{>i}+Q_{\le j}$ for every $0\le i,j\le k$.
\end{enumerate}
The maximum includes every separation up to $q/2$.
Original parity may be prescribed.
For exact $M$ background points on the unit circle the same
bound holds after separation $1/q$, uniformly for
$q\le M/\alpha_0$.
\end{theorem}
\begin{proof}
First use marked Poisson probes: intensity one in case 1,
or two intensities $1/2$ in case 2.
Let $Q$ denote this full independent-input law and $P$
the law with the prescribed probe counts fixed.
Keep continuous priorities initially. There are $m=1$ or
$m=2$ priority coordinates, respectively.

Consider a band $(u,\min(2u,q/2)]$, $u\ge1$.
Partition the circle into $n=\lceil q/u\rceil$ equal cells.
Assign each changed edge to a fixed one of its endpoints.
Its cell score $X_j(\boldsymbol v,s)$ counts assigned
changed edges in $(u,s]$. For every priority state its
positive edge-time measure has mass at most four, by
active-point packing in each of the two inputs.
Let $f_{j,\ell}$ be the same field computed using only
the records in the window of $\ell$ cells on either side,
with the exterior empty. It also has height at most four.

There is a local-record event $E_{j,\ell}$ such that
$f_{j,\ell}=X_j$ for all priorities and band readouts,
and
\begin{equation}\label{eq:priority-local-certificate}
 Q(E_{j,\ell}^c)\le C(1+u\ell)^m e^{-c\ell}.
\end{equation}
For additions, check the shielding certificate for the
fixed background and for every local probe prefix.
The local input changes only at priorities of probes
inside the window. Under one-probe Palm conditioning,
its state at such a priority is the restriction of
homogeneous Poisson input plus at most one forced point.
Lemma~\ref{lem:local-shielding-certificate} and the marked
Poisson summation formula give
$C(1+|W|)e^{-c\ell}$.
For resampling, check the initial input and the grid of
local old/fresh priority corners. One- and two-probe Palm
conditioning leave homogeneous Poisson inputs with at most
two forced points, giving $C(1+|W|)^2e^{-c\ell}$.
These are all the possible local states.
When one is certified, arbitrary exterior priority changes
cannot affect the cell score. Thus local priorities suffice;
no union over the whole circle's probe count is taken.
This proves \eqref{eq:priority-local-certificate}.

The event is determined by $W$. Its bound transfers to
$P$ at a constant factor by
Lemma~\ref{lem:fixed-count-restriction-density}, since
the window occupies at most $1/8$ of the circle.
The assertion remains valid for marked restrictions:
given their counts and locations the priority laws agree.

The local fields are cumulative functions of signed
measures in $m+1$ coordinates. If $K_W$ is the full number
of probes in the window, their variation is at most
$C(1+K_W)^m$. In case 1 each priority jump is the difference
of two edge-time measures of mass at most four.
In case 2 use the axis jumps and mixed priority rectangles;
each rectangle uses four such measures and there are at
most $K_W^2$ rectangles.
Consequently under $Q$ their fixed variation moments are
bounded by $C_p(1+u\ell)^m$, uniformly in $\alpha$.

For $\ell\ge17$ put $D_{j,\ell}=f_{j,\ell}-f_{j,\ell-1}$.
On $E_{j,\ell}\cap E_{j,\ell-1}$ this entire field, and
hence its signed measure, is zero.
H\"older with the higher Poisson count moments yields
\begin{equation}\label{eq:local-field-difference-variation}
 \|\|D_{j,\ell}\|_{\rm TV}\|_{L^p(Q)}
                   \le C_p(1+u\ell)^4e^{-c_p\ell}.
\end{equation}
This is a bound on the whole priority and separation
measure, not only a fixed-priority value.

Choose $\ell_0=\lceil A_p\log(2+u)\rceil\ge16$ and
$\ell_*=\lfloor n/32\rfloor$.
If $\ell_0>\ell_*$, deterministic packing gives a centered
band supremum at most $8n$. Since
$n\le C_p\log(2+u)$ in this case, this is at most
$C_p\sqrt n\log(2+u)$.
We may therefore assume $\ell_0\le\ell_*$.

For a fixed window radius $\ell$, color the window-overlap
graph with $O(\ell)$ colors and split each family into two,
as in Theorem~\ref{thm:late-history-maximal}.
Each family has disjoint windows whose union occupies at
most $3/4$ of the circle.
Under $Q$ its local fields are independent.
For the base fields $f_{j,\ell_0}$,
Lemma~\ref{lem:bounded-field-entropy} gives the centered
family bound
$C_p\sqrt{n_{\rm fam}}\sqrt{\log(2+u\ell_0)}$.
The joint record density under $P$ is bounded by
Lemma~\ref{lem:fixed-count-restriction-density};
transfer this bound and recenter under $P$, at a factor
at most two. Summing the families gives
\[
 C_p\sqrt{n\ell_0}\sqrt{\log(2+u\ell_0)}
                                  \le C_p\sqrt n\log(2+u).
\]
For the successive differences use the ordinary rectangle
measure maximal bound and
\eqref{eq:local-field-difference-variation}.
The same grouping and density transfer give
$C_p\sqrt{n\ell}(1+u\ell)^4e^{-c_p\ell}$.
Their sum over $\ell>\ell_0$ is at most $C_p\sqrt n$
when $A_p$ is sufficiently large.

The residual $X_j-f_{j,\ell_*}$ has height at most eight
and vanishes on the local certificate.
By \eqref{eq:priority-local-certificate} under $P$ and
Minkowski, its sum has maximal $L^p$ norm at most
\[
 C_p n(1+u\ell_*)^{m/p}e^{-c\ell_*/p}.
\]
Since $\ell_*\ge A_p\log(2+u)$ and $\ell_*\asymp n$,
increasing $A_p$ bounds this by $C_p\sqrt n$.
Recenter and telescope the local fields.
We have proved the priority-band bound
\begin{equation}\label{eq:late-priority-band-maximal}
 \left\|\sup_{\boldsymbol v,s}
    |H^{(u,s]}(\boldsymbol v)-\EE H^{(u,s]}(\boldsymbol v)|
                                                    \right\|_p
                   \le C_p\sqrt{q/u}\log(2+u).
\end{equation}
All constants are uniform in the background intensity.

Dyadic summation for $u=1,2,4,\ldots$ uses
$\sum_h2^{-h/2}\log(2+2^h)<\infty$ and gives the full
late priority-field bound $C_p\sqrt q$.
Apply Lemma~\ref{lem:ranking-maximal-transfer} to one or
two colors. This proves the prescribed-prefix assertions,
centered by their own expectations.
Original parity can now be conditioned at bounded cost,
followed by recentering.

For exact original count, the partial-record data at
$r_*=1/\sqrt{Mq}$ determine every subset relative history.
Their good-event density is bounded for
$q\le M/\alpha_0$. In every input, the number of accepted
edges after $1/q$ is at most $Cq$ by deterministic packing,
so the entire raw field, including all prefixes, is bounded
by $Cq$. Global original-count conditioning and the
probability $Ce^{-c\sqrt{Mq}}$ make its bad-event
contribution harmless in every fixed moment.
Transfer the good-event maximal bound and recenter.
This proves the exact-count assertion.
\end{proof}

\begin{corollary}[Brownian sheets on the complete separation clock]
\label{cor:complete-history-sheet}
For nested additions as in Theorem~\ref{thm:early-history-sheet},
or nested coordinate resampling as in
Theorem~\ref{thm:resampling-history-sheet}, extend the clock by
\[
 r(t)=
 \begin{cases}
 M^{-1}(M/q)^t,&0\le t\le1,\\
 q^{-1}(q/2)^{t-1},&1\le t\le2.
 \end{cases}
\]
Keep each theorem's edit mesh, exact mean centering, and
normalization $\sqrt{4qv_{\rm H}\log(M/q)}$.
Bilinear interpolation on those edit grids and the time
grid $j/\lceil\log(M/q)\rceil$, $0\le j\le2\lceil\log(M/q)\rceil$,
converges in $C([0,1]\times[0,2])$ to
$W(u,t\wedge1)$.
The raw field is uniformly close to its interpolation in
probability. For resampling this holds for spatial and
label history distance, centered separately.
\end{corollary}
\begin{proof}
Theorem~\ref{thm:late-prefix-maximal} makes the full centered
late field uniformly $O_{L^p}(\sqrt q)$.
After division by $\sqrt{q\log(M/q)}$ it vanishes, so
the normalized field is uniformly close to its early field
extended constantly in separation time after one.
Apply the two early sheet theorems and the contraction
property of interpolation.
For labels the correction union bound in the resampling
theorem already includes every separation through completion.
\end{proof}

\subsection{Profile centering throughout the sparse Gaussian regime}
\label{sec:all-sparse-profile-centering}

\status{Working proof: the deterministic profile replaces the exact terminal mean}

\begin{lemma}[The Poisson profile mean with a prescribed common parity]
\label{lem:parity-Poisson-profile-mean}
The bound of Lemma~\ref{lem:Poisson-circle-profile-mean}
also holds after prescribing the parity of the common
Poisson color, uniformly for $0\le\theta\le1/2$.
\end{lemma}
\begin{proof}
On circumference $T$, let the common intensity be
$\lambda=1-\theta\ge1/2$.
For a window $W$ and $j$ common points in it, the
parity-conditioned restriction has density relative to
its unconditioned restriction
\[
 \frac{1+\varepsilon(-1)^j e^{-2\lambda(T-|W|)}}
      {1+\varepsilon e^{-2\lambda T}},
 \qquad \varepsilon=1\text{ for even parity},\
        \varepsilon=-1\text{ for odd parity}.
\]
If $|W|=T/3$, this differs uniformly from one by
$Ce^{-cT}$. The full parity-conditioned input law has
density at most a constant relative to the unconditioned
law, so its parent-localization failure bounds are
multiplied only by a constant.

Repeat the line/circle comparison in
Lemma~\ref{lem:Poisson-circle-profile-mean}.
The two window records can now disagree with probability
$Ce^{-cT}$. In a band $(u,2u]$ the additional mean error
is at most $C(T/u)e^{-cT}$, beside the previous
$C(T/u)e^{-cT/u}$ error.
The microscopic unit-cell scores have bounded higher
moments also after parity conditioning; H\"older bounds
their additional error by $CTe^{-cT}$ after decreasing $c$.
Both error sums are bounded uniformly in $T$.
Terminal packing and the expected line tail remain bounded,
proving the assertion.
\end{proof}

\begin{lemma}[The Poisson-background mean with exact probe counts]
\label{lem:fixed-probe-profile-mean}
Let $k\ge4$, $q=2k$, $M\ge q$, and $N=M+k$.
Take a Poisson background of mean $M$ on the unit circle,
optionally conditioned on either parity, and exactly $k$
independent uniform probes of each of two colors.
If $\mu_{M,k}^{\rm fix}$ is their expected completed spatial
history distance, then
\begin{equation}\label{eq:fixed-probe-profile-mean}
 \left|\mu_{M,k}^{\rm fix}-N\mathfrak h(k/N)\right|
                                                    \le C\sqrt k.
\end{equation}
The constant is independent of $M,k$ and the prescribed parity.
\end{lemma}
\begin{proof}
Scale to circumference $q$, so the background intensity
is $\alpha=M/q\ge1$. Compare the fixed probe counts
$(k,k)$ with independent Poisson counts of means $(k,k)$.
Use the same background and each color's iid location sequence.

Through separation one, write the history as the sum of
the single-probe captures plus the actual additivity
remainder. The expectations of the capture sums agree
exactly: both count models have total mean $q$, independent
of the background and the location sequences.
The uncentered count coupling in the proof of
Theorem~\ref{thm:fixed-probe-remainder-variance}, or its
variation version in
Theorem~\ref{thm:early-remainder-maximal}, gives
\[
 \|R_{\rm fixed}-R_{\rm Pois}\|_2\le C\sqrt q.
\]
That proof uses bounded fixed-count local increments,
clipping, and the polynomial Poisson circle remainder
moments, uniformly in $\alpha$ and either original parity.
Thus the early mean difference is at most $C\sqrt q$.
No factor $\log\alpha$ is incurred, since the capture
expectations cancelled before the remainder was compared.

For the late mean, consider a band $(u,\min(2u,q/2)]$
and $n=\lceil q/u\rceil$ cells.
Use the local finite-window score of the preceding section
at the final priority values, with window radius $\ell$.
At this one priority state its shielding failure probability
is $Ce^{-c\ell}$ under the Poisson probe law.
Prescribing background parity costs only a bounded factor.
The local fixed-probe restriction density is bounded, so
the same failure estimate holds with probe counts fixed.
Both cell scores have height at most four.

Let $p_W=|W|/q=(2\ell+1)/n$.
For each probe color, its restricted count is
$\operatorname{Bin}(k,p_W)$ or
$\operatorname{Pois}(kp_W)$, respectively.
The Barbour--Hall bound used in
Theorem~\ref{thm:uniform-history-mean-bias} gives distance
at most $p_W$. Given these counts, the point laws agree.
The two-color restriction distance is therefore at most
$2p_W$; the original background law is shared by the two
models, including its parity.
Consequently the whole band's mean difference, uniformly
over its cumulative readout, is at most
\[
 C\{\ell+n e^{-c\ell}\}.
\]
Take $\ell=\lceil A\log(n+2)\rceil$, with $A$ sufficiently
large. If this exceeds $n/32$, then $n$ is bounded by
an absolute constant depending only on $A$, and packing
gives the same estimate directly.
Otherwise the window and local-certificate arguments apply.
The band mean difference is thus at most $C\log(n+2)$.
Summing the dyadic bands gives
$C\log^2(q+2)$ for the complete late mean difference.

The full Poisson probe model, rescaled instead to
circumference $N$, has color intensities
$1-\theta,\theta,\theta$, $\theta=k/N\le1/3$.
Lemmas~\ref{lem:Poisson-circle-profile-mean}
and~\ref{lem:parity-Poisson-profile-mean} give its mean
$N\mathfrak h(k/N)+O(1)$.
Combining the comparisons gives error
$C\{\sqrt q+\log^2(q+2)+1\}\le C\sqrt q$,
which proves \eqref{eq:fixed-probe-profile-mean}.
\end{proof}

\begin{theorem}[Profile centering for every growing sparse resampling count]
\label{thm:all-sparse-profile-centering}
For odd $N\to\infty$, $k\to\infty$ and $k=o(N)$,
\begin{equation}\label{eq:all-sparse-profile-mean-bias}
 \frac{\EE H_N(k)-N\mathfrak h(k/N)}
                         {\sqrt{k\log(N/k)}}\longrightarrow0.
\end{equation}
The same holds for label histories. Consequently, for either
history convention,
\begin{equation}\label{eq:all-sparse-profile-CLT}
 \frac{H_N(k)-N\mathfrak h(k/N)}
                    {\sqrt{8kv_{\rm H}\log(N/k)}}
                            \Longrightarrow\mathcal N(0,1).
\end{equation}
This uses the full deterministic profile, not merely its
leading small-density equivalent.
\end{theorem}
\begin{proof}
Put $M=N-k$, $q=2k$, and $L=\log(M/q)$.
Compare the exact-$M$ model with a Poisson background of
mean $M$ and matching parity, keeping both probe counts exact.
The mean-path alignment in
Theorem~\ref{thm:all-sparse-history-FCLT} gives an early
mean difference $o(\sqrt{qL})$ through separation $1/q$.
It uses the all-subset relative-history TV coupling and
separate maximal moment uniform integrability.

The late mean difference is $O(\sqrt q)$ as well.
In the proof of
Theorem~\ref{thm:late-history-maximal}, center the
Poisson-background late path by its own mean before using
the original-count reservoir.
Its bounded good-event density transfers its maximal
$L^2$ norm $C\sqrt q$. On the complementary event the
uncentered late cost and its mean are bounded by $Cq$,
and global conditioning costs $C\sqrt M$ against the
probability $Ce^{-c\sqrt{Mq}}$.
Thus the exact-count late path has maximal $L^2$ distance
at most $C\sqrt q$ from the Poisson mean path, even before
recentering by its own mean.
Jensen gives the asserted late mean comparison.

Lemma~\ref{lem:fixed-probe-profile-mean} now compares that
Poisson-background, exact-probe mean with
$N\mathfrak h(k/N)$ at cost $C\sqrt q$.
Since $L\to\infty$ and $L/\log(N/k)\to1$, this proves
\eqref{eq:all-sparse-profile-mean-bias}.
The mean label correction is at most four, so the same
centering holds for labels.
Apply the complete sparse terminal CLT from
Section~\ref{sec:late-history-variance} to obtain
\eqref{eq:all-sparse-profile-CLT}.
\end{proof}

\begin{corollary}[One profile center throughout the growing Gaussian regimes]
\label{cor:all-growing-profile-CLT}
Let $N$ be odd, $k\to\infty$, and
$(N-k)^2/N\to\infty$. Then, for either history convention,
\[
 \frac{H_N(k)-N\mathfrak h(k/N)}
                    {\sqrt{N v_{\rm D}(k/N)}}
                          \Longrightarrow\mathcal N(0,1),
 \qquad
 \frac{\operatorname{Var}H_N(k)}
                         {N v_{\rm D}(k/N)}\longrightarrow1.
\]
No limiting edit fraction is required.
The functions $\mathfrak h$ and $v_{\rm D}$ are the
previously defined deterministic profile and variance
function; their interior values remain unevaluated.
\end{corollary}
\begin{proof}
Every subsequence has a further subsequence on which
$k/N$ tends to a point of $[0,1]$.
At an interior point use the dense CLT, its profile
centering, and continuity of $v_{\rm D}$.
At zero use Theorem~\ref{thm:all-sparse-profile-centering}
and $v_{\rm D}(\theta)\sim8v_{\rm H}\theta\log(1/\theta)$.
At one the assumption $(N-k)^2/N\to\infty$ gives the
Gaussian endpoint theorem and its profile centering, with
$v_{\rm D}(\theta)\sim2\pi_2(1-\theta)^2$.
The corresponding variance asymptotics hold in each regime.
All further-subsequence limits agree, proving the claim.
\end{proof}

\subsection{Deterministic centering of the whole resampling sheet}
\label{sec:profile-centered-history-sheet}

\status{Working proof: cutoff profiles, bounded edits, and uniform prefix means}

For $s\ge0$ and $0\le\theta\le1$, define
$\mathfrak h_s(\theta)$ to be the expected number per unit
length of spatial history-difference edges through separation
$s$ for the stationary line inputs with common/private
Poisson intensities $(1-\theta,\theta,\theta)$.
Assign an edge to its left endpoint to define this stationary
density. Finite original Poisson components define every
finite-cutoff history. Positivity and the marginal matching
intensities give
\begin{equation}\label{eq:cutoff-profile-basic}
 0\le\mathfrak h_s(\theta)\le1,\qquad
 \mathfrak h_s(0)=0,\qquad
 \mathfrak h_s(\theta)\uparrow\mathfrak h(\theta).
\end{equation}
The two marginal late edge intensities also give
$0\le\mathfrak h(\theta)-\mathfrak h_s(\theta)\le q(s)$,
where $q(s)$ is the unit-intensity active-point density.
This defines a deterministic cutoff profile; it is not
an explicit evaluation of that profile.

\begin{lemma}[Uniform cutoff-profile mean comparison]
\label{lem:cutoff-profile-mean}
For a resampling comparison write $H_N(k;r)$ for its
spatial history distance through separation $r$ and put
\[
 b_{N,k}=
 \sup_{1/(N-k)\le r\le1/2}
   \left|\EE H_N(k;r)-N\mathfrak h_{Nr}(k/N)\right|.
\]
For every growing $k=o(N)$,
\begin{equation}\label{eq:growing-cutoff-profile-mean}
 b_{N,k}=o(\sqrt{k\log(N/k)}).
\end{equation}
For each fixed $k\ge1$,
\begin{equation}\label{eq:fixed-cutoff-profile-mean}
 b_{N,k}=O_k(\log\log(N+3)).
\end{equation}
Both bounds hold for label histories as well.
The corresponding terminal bounds hold with
$\mathfrak h$ in place of $\mathfrak h_{N/2}$.
\end{lemma}
\begin{proof}
The line/circle proof of
Lemma~\ref{lem:Poisson-circle-profile-mean} is uniform
over its cumulative cutoff: terminate its dyadic sum at
any band and use the same local comparison for that
band's partial readout. Its absolute error sum is still
bounded. The same is true of the common-parity extension.
Thus, for the fully Poisson color model and $\theta\le1/2$,
\[
 \sup_{0\le r\le1/2}
  \left|\EE H_N^{\rm P}(r)-N\mathfrak h_{Nr}(\theta)\right|
                                                        \le C,
\]
also after prescribing common parity.
For cutoffs below the first unit band on circumference
$N$, use its microscopic cell-count moment bound.

The private-count comparison in
Lemma~\ref{lem:fixed-probe-profile-mean} is uniform over
cutoffs as well. Its early coupling holds in signed-measure
variation, and its late estimate explicitly covers
partial bands. Hence with background still Poisson of
matching parity, fixing both private counts $k$ changes
the cutoff mean by at most $C\sqrt k$.
For $k=1,2,3$, run the same early count coupling on
circumference eight, with probe intensities $k/8$ in each
color. These lie in a fixed compact interval, so the
probe-only increment bounds and polynomial remainder
moments give a bounded coupling cost.
After its unit cutoff deterministic packing bounds every
remaining history by a constant.
This proves the same mean conclusion for these finitely
many counts when $N$ is large.

For growing $k$, the original-count early mean-path
alignment and late mean comparison in
Theorem~\ref{thm:all-sparse-profile-centering} are already
uniform in cutoff. They prove
\eqref{eq:growing-cutoff-profile-mean}.

For fixed $k$ set $M=N-k$ and use the full-family coupling
in Theorem~\ref{thm:fixed-edit-history-CLT}, with
$w=A_k\log\log M$ as justified in
Remark~\ref{rem:history-mean-loglog}.
Its proof compares the entire cutoff process in
$L^1$ supremum norm with twice the sum of independent
line central capture processes, with error $O_k(w)$.
The good-event comparison is simultaneous throughout the
central interval. The two omitted strips have positive
cost and $L^1$ norm $O_k(w)$, which bounds every partial
strip readout, and the exceptional event is controlled
by the higher whole-history moments.
The same comparison holds for the matching-parity Poisson
background of mean $M$: retain the Poisson root/skeleton
part of the proof and omit the original-count coefficient
comparison. The outside parity remainder has the same
exponential bound. Both models therefore have mean paths
within $O_k(w)$ of the same clamped line-capture mean.
Their mutual mean-path difference is $O_k(\log\log M)$.
Combining this with the preceding private-count and
circle-profile bounds proves
\eqref{eq:fixed-cutoff-profile-mean}.

The mean spatial/label correction is at most four at
every cutoff, since its cumulative correction is bounded
by its completed correction. This proves the label claims.
Finally \eqref{eq:cutoff-profile-basic} gives
$N|\mathfrak h-\mathfrak h_{N/2}|\le Nq(N/2)\le C$,
so the terminal centers may be changed as asserted.
\end{proof}

\begin{theorem}[The complete resampling sheet with a deterministic center]
\label{thm:profile-centered-history-sheet}
Take the nested coordinate-resampling path of
Theorem~\ref{thm:resampling-history-sheet}, with odd $N$,
$k\to\infty$, $k=o(N)$, $M=N-k$, $q=2k$, and
$L=\log(M/q)$.
Use the complete two-part separation clock $r(t)$ from
Corollary~\ref{cor:complete-history-sheet}.
At edit fraction $u$, set $j=\lfloor ku\rfloor$.
Then that corollary's sheet limit and uniform interpolation
approximation hold with deterministic centering
\begin{equation}\label{eq:profile-centered-history-sheet}
 N\mathfrak h_{N r(t)}(j/N)
\end{equation}
in place of the exact mean, for both spatial and label histories.
In particular the completed history, centered by
$N\mathfrak h(j/N)$ and divided by $\sqrt{8kv_{\rm H}L}$,
converges as an edit-fraction path to standard Brownian motion.
\end{theorem}
\begin{proof}
We claim that
\begin{equation}\label{eq:uniform-prefix-profile-mean}
 \max_{0\le j\le k}
 \sup_{1/(N-k)\le r\le1/2}
 \frac{|\EE H_N(j;r)-N\mathfrak h_{Nr}(j/N)|}
                              {\sqrt{k\log(N/k)}}\longrightarrow0.
\end{equation}
The term $j=0$ is zero.
If this failed, choose a violating prefix sequence $j_N$.
Pass to a subsequence on which $j_N\to\infty$ or
$j_N$ is fixed. In the growing case
$j_N\le k=o(N)$, and
$j_N\log(N/j_N)\le k\log(N/k)$ for large $N$.
Moreover $1/(N-k)\ge1/(N-j_N)$.
Thus \eqref{eq:growing-cutoff-profile-mean} contradicts
the violation.
In the fixed case use \eqref{eq:fixed-cutoff-profile-mean}.
Here $k\log(N/k)/\log N\to\infty$:
for $k\le\sqrt N$ it is at least $k/2$, while for
$k>\sqrt N$ it diverges by $k=o(N)$.
This again contradicts a violation and proves
\eqref{eq:uniform-prefix-profile-mean}.
The same argument applies to labels.

Since $L/\log(N/k)\to1$, this uniform deterministic
center error is negligible at the sheet normalization.
Apply Corollary~\ref{cor:complete-history-sheet};
bilinear interpolation cannot increase the supremum error.
At completion, replacing the cutoff profile by
$\mathfrak h$ costs uniformly at most a constant, as in
the preceding lemma. The slice $W(u,1)$ is standard
Brownian motion in $u$, proving the final assertion.
\end{proof}

\begin{corollary}[The Gaussian law also permits bounded nonzero edit counts]
\label{cor:all-nonzero-profile-CLT}
For odd $N$, let $1\le k\le N$ satisfy
$(N-k)^2/N\to\infty$.
Then the normal limit and variance asymptotic of
Corollary~\ref{cor:all-growing-profile-CLT} hold, without
requiring $k\to\infty$.
\end{corollary}
\begin{proof}
On any subsequence with $k\to\infty$ use that corollary.
Otherwise pass to a further subsequence with fixed
$k\ge1$. Lemma~\ref{lem:cutoff-profile-mean} gives terminal
profile mean bias $O_k(\log\log N)=o(\sqrt{\log N})$.
The fixed-edit central limit and variance theorem gives
variance $8kv_{\rm H}\log N$, and
$Nv_{\rm D}(k/N)\sim8kv_{\rm H}\log N$.
The resulting limits agree with those on the growing
subsequences, proving the assertion.
\end{proof}

\subsection{The history process across all edit fractions}
\label{sec:dense-history-process}

\status{Working proof: a continuous Gaussian process under nested resampling}

Take independent iid uniform samples $(X_i)_{i\le N}$ and
$(Y_i)_{i\le N}$ on the unit circle, and put
$A_j=\{Y_1,\ldots,Y_j,X_{j+1},\ldots,X_N\}$.
Throughout this subsection $H_N(j;r)$ compares $A_0$ with
$A_j$ through separation $r$, and $H_N(j)$ denotes completion.
The order of replacement is fixed in advance.

\begin{lemma}[A uniform exponential label correction]
\label{lem:label-correction-exponential}
For every $t\ge0$, uniformly in $N,j$,
\begin{equation}\label{eq:label-correction-exponential}
 \EE\exp\{t(H_N(j)-H_N^{\rm lab}(j))\}
                       \le \exp(e^{4t}-1).
\end{equation}
Consequently, for every fixed $p<\infty$,
\[
 \left\|\max_{0\le j\le N}\sup_r
       |H_N(j;r)-H_N^{\rm lab}(j;r)|\right\|_p
                              \le C_p\log(N+2).
\]
The assertion holds for either sample parity.
\end{lemma}
\begin{proof}
Use the two extra common-label edge counts $J_1,J_2$
from Lemma~\ref{lem:history-label-correction}, with $k=j$.
For $k=0$ the correction is zero. Conditional on the
geometry and the first label assignment, any $r$ distinct
unchanged-to-refreshed coincidences prescribe $r$ different
images of the second permutation. There are at most $k$
eligible unchanged points, so
$\EE((J_1)_r\mid\text{geometry})\le1$.
For the refreshed-to-refreshed type, with $a,b\le\lfloor k/2\rfloor$
edges in the two copies,
\[
 \EE((J_2)_r\mid\text{geometry})
             =\frac{2^r(a)_r(b)_r}{(k)_{2r}}\le1.
\]
The expression is zero when $r>\min(a,b)$.
For the inequality, each remaining factor is at most one:
put $d=k-2h\ge2$ and use
$2(a-h)(b-h)\le2\lfloor d/2\rfloor^2\le d(d-1)$.
The factorial generating function therefore gives
$\EE e^{sJ_i}\le\exp(e^s-1)$ for $s\ge0$.
Since the correction is $2(J_1+J_2)$, Cauchy--Schwarz
proves \eqref{eq:label-correction-exponential}; independence
of $J_1,J_2$ is not needed.
For each $j$ the cumulative correction is nonnegative,
nondecreasing in $r$, and bounded by its terminal value.
A union bound in \eqref{eq:label-correction-exponential},
followed by tail integration, proves the maximal assertion.
Only disjointness of the edges in each matching was used,
so the proof does not require odd $N$.
\end{proof}

\begin{lemma}[The fixed-cutoff edit-fraction process]
\label{lem:dense-cutoff-process}
Fix $b<\infty$. Linearly interpolate the grid values
\[
 Z_{N,b}(j/N)=
   \frac{H_N(j;b/N)-\EE H_N(j;b/N)}{\sqrt N},
                    \qquad 0\le j\le N.
\]
Then $Z_{N,b}$ converges in $C([0,1])$ to a centered
continuous Gaussian process $G_b$. The convergence holds
jointly for any finite list of cutoffs, and includes
every fixed moment of the supremum norm.
For each integer $p\ge2$ there is $C_{p,b}$ such that
\begin{equation}\label{eq:dense-cutoff-index-moment}
 \EE|Z_{N,b}(j/N)-Z_{N,b}(i/N)|^{2p}
                   \le C_{p,b}|(j-i)/N|^p.
\end{equation}

Here is a block description of its covariance.
For $0<\theta_1<\cdots<\theta_m<1$, partition the priority
interval at these fractions. Use two Poisson colors for
each interval, old and fresh, each with intensity equal
to that interval's length. In the union process, blocks
end at the first gap longer than $b$.
Let $L$ be block length, $C$ its color-count vector, and
$R=(R_1,\ldots,R_m)$ its history rewards for these
resampling comparisons through $b$.
Writing $Z=(R,C)$ and $a=\EE Z/\EE L$, partition
\[
 \frac{\EE[(Z-aL)(Z-aL)^{\mathsf T}]}{\EE L}
       =\begin{pmatrix}V_b&B_b\\B_b^{\mathsf T}&D\end{pmatrix},
 \qquad D=\operatorname{diag}(\text{color intensities}).
\]
Then
\begin{equation}\label{eq:dense-process-kernel}
 \big(\Gamma_b(\theta_i,\theta_j)\big)_{i,j\le m}
                 =V_b-B_bD^{-1}B_b^{\mathsf T}.
\end{equation}
Endpoints may be included by omitting zero-length colors.
\end{lemma}
\begin{proof}
Rescale the circle to circumference $N$. For a fixed
fraction grid, condition the independent Poisson colors
to their exact counts
$\lfloor N\theta_\ell\rfloor-\lfloor N\theta_{\ell-1}\rfloor$
in each old and fresh family.
This is precisely the joint law of the required prefixes;
the unused fresh colors can also be retained.
The proof of Lemma~\ref{lem:dense-cutoff-conditional-CLT}
extends to this fixed finite color set and the vector
reward $R$. Union-gap blocks have exponential moments.
In an isolated interior, zero union points and one point
of any color all give the zero reward vector.
They provide the Fourier bound
$C_b\exp\{-c_bN\sum_\ell\operatorname{dist}(z_\ell,2\pi\mathbb Z)^2\}$
for the color counts, uniformly in all reward frequencies.
The zero-or-one interior reservoir has maximal count
probability $C_bN^{-d/2}$ for $d$ positive colors,
matching the exact target probability.
Thus Fourier inversion and the bounded-density moment
transfer give the joint Gaussian limit and the Schur
complement \eqref{eq:dense-process-kernel}.
All intensities are bounded away from zero for this fixed
grid after zero-length intervals are omitted.
This proves finite-dimensional convergence, including
moments. Several cutoffs use their largest cutoff for the
blocks and put all rewards in the same vector.

We give a uniform increment argument to establish
tightness; finite-dimensional convergence alone does not
give it. Fix $i<j$, put $h=(j-i)/N$, and split each
family into the three index ranges
$[1,i],(i,j],(j,N]$.
Let $P$ be the exact-count law of these at most six colors,
and $Q$ their independent Poisson laws with matching means.
The total union intensity under $Q$ is two.
Call the old and fresh colors in $(i,j]$ changed tags;
their combined intensity is $2h$.

Partition the circle into cells of length comparable
with $b+1$. Assign each short edge to its clockwise
initial endpoint, and let $X_\nu$ be the cell contribution
to $H_N(j;b/N)-H_N(i;b/N)$.
For a window of $\ell$ neighboring cells on each side,
let $f_{\nu,\ell}$ be the same increment computed with
empty exterior. There is a local certificate, a union-empty
interval longer than $b$ on each side of the cell, which
makes this computation agree with every exterior completion
and every prefix. Testing disjoint intervals of length
$b+1$ bounds failure by $C_be^{-c_b\ell}$ under $Q$.
The same bound holds with one forced point: omit the at
most one test interval containing it.
Every score is bounded by four times the union point
count in its cell, which has uniformly bounded fixed
moments, depending only on $b$.

If there is no changed tag in the window, then
$f_{\nu,\ell}=0$. Moreover
$f_{\nu,\ell}-f_{\nu,\ell-1}$ vanishes when both local
certificates hold. Summing over a changed tag in the larger
window, using its intensity $2h$, and then applying
H\"older to the one-forced-point certificate gives, for
each fixed $r\ge2$,
\[
 \EE_Q|f_{\nu,8}|^r\le C_{r,b}h,\qquad
 \EE_Q|f_{\nu,\ell}-f_{\nu,\ell-1}|^r
       \le C_{r,b}h(1+\ell)^{C_r}e^{-c_{r,b}\ell}.
\]
Increasing the initial fixed radius if needed changes only
the constants. This is a joint rare-tag and failed-shielding
bound, not a product of two asserted independent events.

For each radius, group disjoint windows into $O(\ell)$
families and split them so their unions occupy at most
$3/4$ of the circle. Under $Q$ the local increments in
one family are independent. Expansion of the centered
$2p$th moment has no singleton factors. The preceding
moment bounds therefore give a family norm at most
\[
 C_{p,b}(1+\ell)^{C_p}e^{-c_{p,b}\ell}
       \{\sqrt{n_{\rm fam}h}+(n_{\rm fam}h)^{1/(2p)}\}
\]
for successive-radius differences; omit the exponential
factor for the fixed initial radius.
Lemma~\ref{lem:fixed-count-restriction-density} transfers
each bound to $P$, uniformly even when some color counts
are zero or one. Recentering under $P$ costs at most two.
Summing families and then radii bounds the full local sum
by $C_{p,b}\{\sqrt{Nh}+(Nh)^{1/(2p)}\}$.

Stop at radius a fixed small fraction of the number of
cells. The remaining error has norm at most
$C_{p,b}N e^{-c_{p,b}N}$: use the local certificate,
the bounded restriction density, the cell count moments,
and Minkowski. Since $Nh=j-i\ge1$, this error and the
preceding bound are at most $C_{p,b}\sqrt{Nh}$.
Finitely many $N$ for which the windows do not fit are
covered by the deterministic total-count bound.
Dividing by $\sqrt N$ proves
\eqref{eq:dense-cutoff-index-moment}.

Linear interpolation extends this estimate to all
fractions with a constant factor: within one mesh interval
the increment is a proportional part of a grid increment;
between intervals split into the two end pieces and a
grid increment. The resulting bound is
$C_{p,b}|\theta-\eta|^p$ for the $2p$th moment.
Dyadic chaining with $p>1$ gives tightness in $C([0,1])$.
Using arbitrarily higher moments in the same chaining
also bounds every fixed supremum moment uniformly in $N$.
Since the process starts at zero, no separate anchor bound
is needed. Finite-dimensional convergence and uniform
integrability finish the proof.
\end{proof}

\begin{lemma}[Late histories uniformly over all resampling fractions]
\label{lem:dense-late-prefix}
For every fixed $p\ge2$, uniformly for $2\le b\le N/2$,
\begin{equation}\label{eq:dense-late-prefix}
 \left\|\max_{0\le j\le N}\sup_{b/N\le r\le1/2}
   \left|H_N(j;r)-H_N(j;b/N)
           -\EE\{H_N(j;r)-H_N(j;b/N)\}\right|\right\|_p
        \le C_p\sqrt{N/b}\log(2+b).
\end{equation}
There is no positive common-background assumption here.
\end{lemma}
\begin{proof}
Work on circumference $N$, giving the old and fresh
families independent uniform priorities. First let $Q$
be the law of two independent rate-one marked Poisson
inputs and $P$ its conditioning to $N$ points of each type.
For thresholds $(v,w)$ compare the full old input with
old priorities greater than $v$ and fresh priorities at
most $w$. Write its cumulative history field as $F(v,w;s)$.

Use the three squares
\[
 \mathcal R_a=[a,a+1/2]^2,\qquad a=0,1/4,1/2.
\]
Their union contains $\{|v-w|\le1/4\}$.
On each whole square the current input has Poisson
intensity $1-v+w\ge1/2$; the initial input has intensity one.
This finite rectangular cover avoids taking signed
variation on a nonrectangular strip.

Consider a band $(u,\min(2u,N/2)]$, $u\ge2$, and cells
of length in $[u/2,u]$.
The local shielding certificate of
Lemma~\ref{lem:local-shielding-certificate} applies at
intensity at least $1/2$: rescale space by $1/2$, so the
new intensity is at least one and the new band scale is
at least one. For a window of radius $\ell$, check the
initial input and every local old/fresh priority corner
within one square, including its boundary states.
The one- and two-tag Palm arguments in
Theorem~\ref{thm:late-prefix-maximal} give a local event
$E_{\nu,\ell,a}$ with
\[
 Q(E_{\nu,\ell,a}^c)
                  \le C(1+u\ell)^2e^{-c\ell}.
\]
The forced tags number at most two and all remaining
marginal intensities stay at least $1/2$.
On this event the local cell field agrees with the true
field on the entire square, for every band readout and
every exterior completion.
The bound transfers to $P$ through the two fixed-count
restriction densities.

Reanchor the field on the square's lower-left corner,
including boundary-axis values as atoms.
Its edge-time measures have mass at most four at every
priority state by active packing. If $K_W$ is the number
of old and fresh points in the window, the resulting
three-coordinate signed measure has variation at most
$C(1+K_W)^2$. This also bounds the measures on boundary
axes and the corner. A successive-radius difference is
identically zero on the square when both certificates
hold, so its reanchored measure is zero there too.
H\"older and Poisson count moments yield
\[
 \left\|\|f_{\nu,\ell}-f_{\nu,\ell-1}\|_{\rm TV,\mathcal R_a}
                                                 \right\|_{L^p(Q)}
                \le C_p(1+u\ell)^4e^{-c_p\ell}.
\]

All inputs to the localization sum in
Theorem~\ref{thm:late-prefix-maximal} now hold on this
rectangle. Explicitly, start at
$\ell_0=\lceil A_p\log(2+u)\rceil\ge16$.
The bounded-field entropy lemma and disjoint-window
grouping give
$C_p\sqrt{(N/u)\ell_0}\sqrt{\log(2+u\ell_0)}$
for the base fields. The successive differences sum to
$C_p\sqrt{N/u}$ by the displayed exponential variation
bound and the ordinary rectangle-measure maximal lemma.
The final local-certificate error is bounded by
$C_p n(1+u\ell_*)^{2/p}e^{-c\ell_*/p}$, where
$n=\lceil N/u\rceil$ and $\ell_*=\lfloor n/32\rfloor$.
It is at most $C_p\sqrt n$ when $\ell_0\le\ell_*$,
after increasing $A_p$. If $\ell_0>\ell_*$, packing
gives $8n\le C_p\sqrt n\log(2+u)$ directly.
The restricted densities are bounded on each disjoint
window family, whose union has length at most $3N/4$;
transfer and recenter just as in that theorem.
Thus on each square the centered band maximum under $P$
is at most $C_p\sqrt{N/u}\log(2+u)$.
Summing the three squares and the dyadic bands starting
at $b$ gives
$C_p\sqrt{N/b}\log(2+b)$.

It remains to pass from priorities to the actual equal
prefixes. Let $U_j,V_j$ be the two priority order
statistics, with $U_0=V_0=0$.
The binomial exponential-moment bound, applied at
$j/N\pm1/8$ and summed over $j$, gives
\[
 P\left(\max_j|U_j-j/N|>1/8
          \ \hbox{or}\ \max_j|V_j-j/N|>1/8\right)
                                      \le Ce^{-cN}.
\]
Consequently every $(U_j,V_j)$ lies in the three-square
cover outside this event.
For arbitrary thresholds the whole late field has
deterministic height at most $CN/b$, by packing.
The exceptional event therefore costs only
$C_p(N/b)e^{-c_pN}$ in maximal $L^p$ norm, and this is
absorbed by the claimed bound.
We have a maximal bound for the field evaluated at
$(U_j,V_j)$ about its priority mean evaluated there.
The ranked old and fresh locations are independent of
their order statistics. Conditional averaging over
those statistics and recentering, exactly as in
Lemma~\ref{lem:ranking-maximal-transfer}, proves
\eqref{eq:dense-late-prefix}.
\end{proof}

\begin{theorem}[The Gaussian edit-fraction process]
\label{thm:dense-history-process}
For odd $N$, linearly interpolate
\[
 Z_N(j/N)=\frac{H_N(j)-N\mathfrak h(j/N)}{\sqrt N},
                         \qquad 0\le j\le N.
\]
Then $Z_N$ converges in $C([0,1])$ to a centered continuous
Gaussian process $G$, with
\[
 G(0)=G(1)=0,\qquad
 \EE G(\theta)^2=v_{\rm D}(\theta)\quad(0<\theta<1).
\]
Its covariance is
\begin{equation}\label{eq:dense-process-full-kernel}
 \Gamma(\theta,\eta)=\lim_{b\to\infty}\Gamma_b(\theta,\eta),
\end{equation}
where \eqref{eq:dense-process-kernel} uses exact color-count
conditioning. The convergence in
\eqref{eq:dense-process-full-kernel} is uniform on $[0,1]^2$.
The same limit holds for label histories.
The raw field
$\{H_N(\lfloor N\theta\rfloor)-N\mathfrak h(\theta)\}/\sqrt N$
is uniformly close in probability to the interpolated one.
No Brownian-bridge or Markov identification of $G$ is asserted.
\end{theorem}
\begin{proof}
First center all grid values by their exact expectations.
Lemma~\ref{lem:dense-late-prefix} bounds, in maximal
$L^p$ norm after division by $\sqrt N$, the difference
between the full process and its cutoff-$b$ version by
$C_p b^{-1/2}\log(2+b)$.
Interpolation cannot increase this bound.
Lemma~\ref{lem:dense-cutoff-process} supplies a continuous
Gaussian limit for each fixed cutoff, jointly for any
finite collection of cutoffs.

Construct the countable joint Gaussian family
$(G_{2^a})_{a\ge1}$ using these consistent joint limits.
The same maximal bound transfers to each pair of
limiting cutoffs by lower semicontinuity. In particular,
\[
 \sum_{a\ge1}\EE\|G_{2^{a+1}}-G_{2^a}\|_\infty
                          \le C\sum_{a\ge1}(a+1)2^{-a/2}<\infty.
\]
They therefore converge uniformly almost surely to a
continuous Gaussian process $G$. The limiting maximal
$L^p$ tail bound and the prelimit bound show that the
full exact-centered process converges to $G$: apply the
triangle inequality first to a fixed cutoff and then
send the cutoff to infinity.
The fixed-cutoff supremum moments and the uniform late
bound also give uniformly bounded full supremum moments.
Cauchy--Schwarz consequently gives uniform convergence
of the covariance kernels in
\eqref{eq:dense-process-full-kernel}.
Consistency means that inserting more fraction-grid
points in the block description does not change an
already specified covariance.

The scalar dense CLT identifies the interior variances
as $v_{\rm D}(\theta)$. At zero the history is zero.
At complete refresh the two spatial matchings have no
common point almost surely, so $H_N(N)=N-1$ and its
exact-centered fluctuation is zero. Thus both endpoint
values of $G$ are zero.

The uniform mean bias in
Theorem~\ref{thm:uniform-history-mean-bias} is
$O(\log^2(N+2))=o(\sqrt N)$, so the deterministic profile
center gives the same limit. The exponential label
correction lemma makes the whole label path differ by
$O_{L^p}(\log(N+2)/\sqrt N)$.

Finally tightness of the interpolated process in
$C([0,1])$ makes its largest mesh jump tend to zero in
probability. Since $\mathfrak h$ is increasing and
concave, its oscillation on any interval of length $1/N$
is at most $\mathfrak h(1/N)=O(\log(N+2)/N)$.
The raw field and its interpolation thus differ by at
most that largest mesh jump plus
$O(\log(N+2)/\sqrt N)$. This proves the last assertion,
also for labels.
\end{proof}

\subsection{The small-edit boundary of the Gaussian history process}
\label{sec:small-edit-Gaussian-tangent}

\status{Working proof: a functional boundary limit from the two path theorems}

Let $G$ and $\Gamma$ be the centered Gaussian process and its
covariance from Theorem~\ref{thm:dense-history-process}.
The small-density variance asymptotic determines its scale
near zero but does not by itself determine its local path law.
The complete sparse path theorem supplies that additional information.

\begin{theorem}[The Brownian tangent at zero]
\label{thm:small-edit-Gaussian-tangent}
As $\delta\downarrow0$,
\begin{equation}\label{eq:small-edit-Gaussian-tangent}
 \left\{\frac{G(\delta u)}
                {\sqrt{8v_{\rm H}\delta\log(1/\delta)}}:
                   0\le u\le1\right\}
       \ \Longrightarrow\ \{B(u):0\le u\le1\}
                  \quad\hbox{in }C([0,1]),
\end{equation}
where $B$ is standard Brownian motion. Moreover,
\begin{equation}\label{eq:small-edit-kernel-tangent}
 \sup_{u,v\in[0,1]}
 \left|
 \frac{\Gamma(\delta u,\delta v)}
                   {8v_{\rm H}\delta\log(1/\delta)}
                  -\min(u,v)\right|\longrightarrow0.
\end{equation}
\end{theorem}
\begin{proof}
Use the linearly interpolated, profile-centered field $Z_N$
from Theorem~\ref{thm:dense-history-process}. For integers
$1\le k<N/3$, put
\[
 s_{N,k}=\sqrt{8kv_{\rm H}\log((N-k)/(2k))},\qquad
 X_{N,k}(u)=\frac{\sqrt N}{s_{N,k}}Z_N(ku/N).
\]
At $u=j/k$ this is exactly
\[
 \frac{H_N(j)-N\mathfrak h(j/N)}{s_{N,k}},
\]
and both sides have the same linear interpolation between
these grid points. Thus
Theorem~\ref{thm:profile-centered-history-sheet} gives
$X_{N,k}\Longrightarrow B$ for every sequence
$k\to\infty$, $k=o(N)$, with odd $N$.
No additional approximation of the mean is being made here.

For a fixed $0<\delta<1/3$, take $k=\lfloor N\delta\rfloor$.
The dense path theorem and continuity of restriction and
linear rescaling give
\begin{equation}\label{eq:fixed-delta-tangent-approximation}
 X_{N,\lfloor N\delta\rfloor}
       \Longrightarrow Y_\delta,\qquad
 Y_\delta(u)=
 \frac{G(\delta u)}
 {\sqrt{8v_{\rm H}\delta\log((1-\delta)/(2\delta))}} .
\end{equation}
Indeed $k/N\to\delta$ and the maps
$f\mapsto f((k/N)\,\cdot)$ converge uniformly on compact
subsets of $C([0,1])$ to $f\mapsto f(\delta\,\cdot)$.

Here is the diagonal argument that joins these two limits.
Let $\delta_j\downarrow0$ be arbitrary and let $d_{\rm BL}$
be a bounded-Lipschitz metric for weak convergence of laws
on $C([0,1])$ with its supremum norm.
For each $j$, choose an odd $N_j>N_{j-1}$ large enough that
$k_j=\lfloor N_j\delta_j\rfloor\ge j$ and
\[
 d_{\rm BL}\bigl(\mathcal L(X_{N_j,k_j}),
                         \mathcal L(Y_{\delta_j})\bigr)<1/j.
\]
Such a choice follows from
\eqref{eq:fixed-delta-tangent-approximation}, with $\delta_j$
held fixed during this choice. Now $k_j\to\infty$ and
$k_j/N_j\le\delta_j\to0$, so the sparse path theorem gives
$X_{N_j,k_j}\Longrightarrow B$. The displayed distance then
gives $Y_{\delta_j}\Longrightarrow B$.
Every sequence tending to zero has the same conclusion.
Finally,
\[
 \frac{\log((1-\delta)/(2\delta))}{\log(1/\delta)}
                      \longrightarrow1,
\]
which proves \eqref{eq:small-edit-Gaussian-tangent}.
This argument uses the sparse theorem for arbitrary
sublinear sequences; fixed-fraction convergence alone
would not justify the conclusion.

For completeness, path convergence of these centered
Gaussian processes also gives the uniform covariance claim
without a separate moment-transfer assertion.
Write $W_\delta$ for the process on the left of
\eqref{eq:small-edit-Gaussian-tangent}.
Tightness in $C([0,1])$ bounds their supremum norms in
probability. A centered normal variable has a fixed
positive probability of exceeding its standard deviation
in absolute value; hence
$\sup_{\delta,u}\EE W_\delta(u)^2<\infty$ for all sufficiently
small $\delta$.
Likewise tightness makes the modulus of continuity small
in probability, and the same normal-variable observation gives
\[
 \lim_{\eta\downarrow0}\limsup_{\delta\downarrow0}
 \sup_{|u-v|\le\eta}
       \EE|W_\delta(u)-W_\delta(v)|^2=0.
\]
Cauchy--Schwarz therefore makes their covariance kernels
uniformly equicontinuous as $\delta\downarrow0$.
Convergence of the centered Gaussian finite-dimensional laws
gives pointwise convergence of those kernels to $\min(u,v)$.
A finite grid and the preceding equicontinuity upgrade this
to uniform convergence, proving
\eqref{eq:small-edit-kernel-tangent}.
\end{proof}

\subsection{A Poisson representation of every matching spectral coefficient}
\label{sec:matching-Poisson-spectrum}

\status{Working proof: positive coefficients and sharp polynomial-centering windows}

Let $\eta$ be a unit-rate Poisson process on $\mathbb R$.
For $x>0$, let $f_x(\eta)$ indicate that the forced points
$0,x$ form an accepted edge after they are added to $\eta$.
Its decision is measurable from the finite component containing
the endpoints at separation $x$; define it arbitrarily on the
null exceptional configurations. Write $p(x)=\EE f_x(\eta)$,
as in Proposition~\ref{prop:matching-quadratic-endpoint}.
For $n\ge1$ and $\boldsymbol z=(z_1,\ldots,z_n)$, set
\[
 K_{n,x}(\boldsymbol z)
   =\EE\sum_{A\subseteq\{1,\ldots,n\}}(-1)^{n-|A|}
            f_x\left(\eta+\sum_{i\in A}\delta_{z_i}\right),
 \qquad K_{0,x}=p(x).
\]
Repeated positions and distance ties may be discarded in the
integrals below. These are mean iterated add-one differences,
not probabilities, and they may have either sign.

\begin{proposition}[The matching spectrum as squared Poisson kernels]
\label{prop:matching-Poisson-spectrum}
For every integer $n\ge0$,
\begin{equation}\label{eq:matching-Poisson-spectrum}
 \pi_{n+2}
   =\frac2{n!}\int_0^\infty
                  \int_{\mathbb R^n}K_{n,x}(\boldsymbol z)^2
                                      \,d\boldsymbol z\,dx.
\end{equation}
For $n=0$ the inner integral means $p(x)^2$.
All these integrals are finite and their sum is one.
In particular the previously bounded cubic coefficient has the
exact representation
\begin{equation}\label{eq:matching-cubic-integral}
 \pi_3=2\int_0^\infty\int_{\mathbb R}
          \bigl\{\EE f_x(\eta+\delta_z)-p(x)\bigr\}^2\,dz\,dx.
\end{equation}
This identifies its integrand but does not evaluate the integral.
\end{proposition}
\begin{proof}
We use the bilinear Fock-space identity of
\href{https://arxiv.org/pdf/0909.3205}{Last and Penrose,
\emph{Poisson process Fock space representation, chaos expansion
and covariance inequalities}}, Theorem~1.1, equation~(1.5).
For a Poisson process with any $\sigma$-finite intensity $\mu$,
that identity writes the inner product of two square-integrable
functionals as the sum of the inner products of their mean
iterated differences, with weights $1/n!$ and measures $\mu^n$.
Our functionals are bounded indicators, so the theorem applies.

Fix $0\le\delta\le1$ and use the marked space
$\mathbb R\times\{c,1,2\}$, with respective intensities
$\delta,1-\delta,1-\delta$.
The first indicator sees colors $c,1$ and the second sees
colors $c,2$. A difference in a color invisible to a functional
is zero. In the bilinear identity, therefore, every term
vanishes unless all its difference coordinates have color $c$.
On those coordinates both kernels equal $K_{n,x}$, since
each marginal input has rate one. If $p_\delta(x)$ is the
probability that both inputs accept the forced edge, the identity is
\begin{equation}\label{eq:forced-edge-Poisson-spectrum}
 p_\delta(x)=
   \sum_{n\ge0}\frac{\delta^n}{n!}
          \int_{\mathbb R^n}K_{n,x}(\boldsymbol z)^2
                                      \,d\boldsymbol z.
\end{equation}
Thus the power of $\delta$ counts common background points;
it does not count the two endpoints, which are already forced.

The pair Palm identity gives
$1-\mathfrak h(1-\delta)=2\delta^2\int_0^\infty p_\delta(x)\,dx$.
Tonelli applies to the nonnegative terms in
\eqref{eq:forced-edge-Poisson-spectrum}.
At $\delta=1$ its sum is
$\EE f_x(\eta)^2=p(x)$, and $2\int_0^\infty p(x)\,dx=1$.
Consequently the integrated coefficients are finite and sum
to one. Uniqueness of the power series in
Corollary~\ref{cor:matching-spectrum} identifies them with
$\pi_{n+2}$, proving the proposition.
\end{proof}

\begin{theorem}[Every matching spectral coefficient is positive]
\label{thm:matching-spectrum-positive}
For all $j\ge2$, $\pi_j>0$.
More explicitly, for $n\ge1$, put
$r_n=1/\{18(2^n-1)\}$. Then
\begin{equation}\label{eq:matching-spectrum-positive-bound}
 \pi_{n+2}\ge
       \frac{r_n^n\,4^{-n(n+1)/2}}{81}>0.
\end{equation}
These small explicit lower bounds are not numerical
approximations of the coefficients.
\end{theorem}
\begin{proof}
The case $j=2$ is already known.
Fix $n\ge1$, abbreviate $r=r_n$, and take
$1/9<x<1/3$. The formula for $p$ gives
$p(x)\ge e^{-3x}>1/3$.
Choose $z_i$ in the disjoint intervals
\[
 I_i=(r4^{\,i-n-1},\,2r4^{\,i-n-1}),
                         \qquad 1\le i\le n.
\]
All these points lie in $(0,r/2)$. In any nonempty subset,
its smallest point $z$ is less than half every other point
in that subset.

If $\eta((-r,2r))=0$, the points $0,z$ are each closer to
the other than to any other forced or Poisson point.
The point $x$ lies beyond $2r$.
Thus neither endpoint can be removed before separation $z$,
and they match then. In particular the forced edge $0,x$
is absent, whatever the exterior configuration.
For every nonempty $A\subseteq\{1,\ldots,n\}$ this proves
\[
 0\le\EE f_x\left(\eta+\sum_{i\in A}\delta_{z_i}\right)
                  \le1-e^{-3r}\le3r.
\]
Separate the empty-subset term in the definition of $K_{n,x}$.
The triangle inequality gives, throughout the displayed box,
\[
 |K_{n,x}(\boldsymbol z)|
       \ge p(x)-(2^n-1)3r\ge1/6.
\]
The box volume is
$\prod_i|I_i|=r^n4^{-n(n+1)/2}$.
Its $n!$ coordinate permutations are disjoint, and the kernel
is symmetric. Integrate its square over these boxes and over
$x\in(1/9,1/3)$ in \eqref{eq:matching-Poisson-spectrum}:
\[
 \pi_{n+2}\ge
 \frac2{n!}\,\frac29\,\frac1{36}\,
       n!r^n4^{-n(n+1)/2}.
\]
This is \eqref{eq:matching-spectrum-positive-bound}.
\end{proof}

\begin{corollary}[Every fixed polynomial-centering window is sharp]
\label{cor:all-polynomial-windows-sharp}
Fix $J\ge2$, let $M=N-k=o(N)$ with $M^2/N\to\infty$,
and define
\[
 C_{N,J}=N-1-N\sum_{j=2}^J\pi_j(M/N)^j.
\]
If $M/N^{1-1/(2J)}\to c\in[0,\infty)$, then, for
spatial or label history distance,
\begin{equation}\label{eq:all-polynomial-critical-shift}
 \frac{H_N(k)-C_{N,J}}{\sqrt{2\pi_2}\,M/\sqrt N}
       \Longrightarrow
       \mathcal N\left(-\frac{\pi_{J+1}c^J}{\sqrt{2\pi_2}},\,1\right).
\end{equation}
If that ratio tends to infinity, the normalized expression
tends to $-\infty$ in probability.
Consequently the centered standard normal limit with center
$C_{N,J}$ holds precisely when
$M=o(N^{1-1/(2J)})$ within this Gaussian endpoint regime.
\end{corollary}
\begin{proof}
Put $\delta=M/N$.
The difference of the full profile center from $C_{N,J}$,
divided by the displayed standard deviation, is
\[
 \frac1{\sqrt{2\pi_2}\sqrt N\,\delta}
 -\frac{\sqrt N}{\sqrt{2\pi_2}}
                    \sum_{j>J}\pi_j\delta^{j-1}.
\]
The first term vanishes. The leading term in the sum is
$\pi_{J+1}\sqrt N\,\delta^J$ and the remaining terms are
bounded by $\sqrt N\,\delta^{J+1}$, since the coefficients
are nonnegative and sum to one.
At a finite limiting ratio these terms give exactly the
shift in \eqref{eq:all-polynomial-critical-shift}.
Apply Theorem~\ref{thm:endpoint-profile-centering}.
At a diverging ratio, strict positivity of $\pi_{J+1}$
makes the deterministic shift tend to $-\infty$, while
the profile-centered fluctuation stays tight.
For a general sequence failing the stated little-oh
condition, take a subsequence with a positive finite or
infinite limiting ratio; its limit cannot be centered
standard normal. This proves necessity as well as sufficiency.
\end{proof}

\subsection{The covariance tangent near complete refresh}
\label{sec:upper-Gaussian-tangent}

\status{Working proof: uniform covariance, finite-dimensional and $L^2$ limits}

Write $\kappa_b=\int_0^b p(t)^2\,dt$ and
$\kappa=\pi_2/2=\lim_{b\to\infty}\kappa_b$.
The natural variables near complete refresh are shared edges.
Truncating the history distance itself retains marginal matching-count
fluctuations that need not be small on this boundary scale.

For a retention fraction $r\in[0,1]$, put
$M_N(r)=\lceil Nr\rceil$. Let $J_N(r)$ count the shared spatial
edges between the original input and its copy after
$N-M_N(r)$ replacements, and let $J_{N,b}(r)$ retain only edges
of length at most $b/N$. All sample-size limits below use odd $N$.
The complete identity $H_N(N-M_N(r))=N-1-2J_N(r)$ gives
\[
 K(r,s):=\lim_{N\to\infty}
       \frac{\operatorname{Cov}(J_N(r),J_N(s))}{N}
       =\frac14\Gamma(1-r,1-s).
\]
The covariance limit follows from the moment bounds in
Theorem~\ref{thm:dense-history-process}.

\begin{lemma}[Two retention levels at a fixed cutoff]
\label{lem:upper-tangent-cutoff}
For every fixed $b\ge1$, the limits
\[
 K_b(r,s)=\lim_{N\to\infty}
       \frac{\operatorname{Cov}(J_{N,b}(r),J_{N,b}(s))}{N}
\]
exist. There is a finite $C_b$ such that
\begin{equation}\label{eq:upper-tangent-cutoff}
 \sup_{0\le r,s\le\delta}
 \left|K_b(r,s)-\kappa_b\min(r,s)^2\right|
                   \le C_b\delta^3,\qquad 0<\delta\le1/2.
\end{equation}
\end{lemma}
\begin{proof}
Give independent rate-one Poisson inputs $O,F$ independent
uniform marks. At retention level $r$, the current input is
$A_r=O_{>1-r}\cup F_{\le1-r}$, and the original is $O$.
For a fixed list of levels, partition each mark interval at
their thresholds. Include all old and fresh colors, even
unused fresh points, so the union intensity is always two.
Its components separated by gaps longer than $b$, followed
by their separating gaps, form independent blocks.
Their lengths $L$ and point counts have exponential moments
depending on $b$, uniformly over the color partition.

Let $R_r$ be a block's shared-edge reward at level $r$,
$C_c$ its count of color $c$, and $\lambda_c$ that color's
intensity. Put $\mu_r=\EE R_r/\EE L$ and
\[
 V_{rs}=\frac{\EE[(R_r-\mu_rL)(R_s-\mu_sL)]}{\EE L},
 \qquad
 \mathcal B_{r,c}=
 \frac{\EE[(R_r-\mu_rL)(C_c-\lambda_cL)]}{\EE L}.
\]
The finite-color conditioning argument of
Lemma~\ref{lem:dense-cutoff-process} applies to these vector
rewards: isolated interiors with zero or one union point
give zero shared-edge reward at every level and retain the
count reservoir. Thus the exact-count covariance is
\begin{equation}\label{eq:shared-cutoff-Schur}
 K_b(r,s)=V_{rs}
       -\sum_{\lambda_c>0}
                 \frac{\mathcal B_{r,c}\mathcal B_{s,c}}{\lambda_c}.
\end{equation}
Each positive intensity is fixed when $N\to\infty$;
zero colors are omitted. No uniform conditional limit theorem
for vanishing color intensities is used.

Let $R_{r,s}$ count edges shared by all three matchings
$O,A_r,A_s$. Two distinct edges counted by $R_r,R_s$ belong
to the same original matching, so they have four different
old endpoints. Conditional on the unmarked block, the
original matching is independent of its marks. Consequently
\[
 \EE(R_rR_s-R_{r,s})\le C_br^2s^2,\qquad
 \EE[L^aR_r]\le C_{a,b}r^2\quad(a=0,1),
\]
and
$V_{rs}=\EE R_{r,s}/\EE L+O_b(r^2s^2)$.
Writing $w=\min(r,s)$, the pair Palm formula gives
\[
 \frac{\EE R_{r,s}}{\EE L}
      =w^2\int_0^b a_{r,s}(t)\,dt,
\]
where $a_{r,s}(t)$ is the probability that the forced
old points $0,t$, retained at both levels, match in all
three inputs. Positive displacement counts each unordered
pair once, so no further factor of two occurs here.

In the union component containing this forced pair, call
any other point changed if its old or fresh mark exceeds
$1-\max(r,s)$. The expected number of other points in
the component is at most $C_b$, uniformly for $t\le b$.
The probability of any changed point there is therefore
at most $C_b\max(r,s)$.
Without such a point, both current inputs agree locally
with the full fresh input plus the forced pair.
The union gaps shield the decisions, and the old and fresh
backgrounds are independent. Coupling to these two independent
decisions proves
$|a_{r,s}(t)-p(t)^2|\le C_b\max(r,s)$.
It follows that
\begin{equation}\label{eq:upper-unconditioned-covariance}
 V_{rs}=\kappa_b\min(r,s)^2+O_b(\delta^3)
                         \qquad(r,s\le\delta).
\end{equation}

We check that the Schur correction has the same error order,
including arbitrarily narrow color intervals.
The Poisson counting identity, applied first on a large
circle and then divided by its length, gives
\[
 \mathcal B_{r,c}
       =\lambda_c\,\EE D_0^cJ_{b,r}.
\]
Here $D_0^cJ_{b,r}$ is the total change in the cutoff
shared-edge count upon inserting one point of color $c$
at zero. Only its union component can change, and that
component has bounded point-count moments depending on $b$.
If the inserted point is old and retained at level $r$,
a nonzero change requires at least one further retained
old point. In every other case it requires at least two
retained old points. Averaging the other independent marks,
and bounding the absolute change by the component point count,
therefore gives
\[
 |\EE D_0^cJ_{b,r}|\le
 \begin{cases}
 C_br,&c\text{ is old and retained at }r,\\
 C_br^2,&\text{otherwise}.
 \end{cases}
\]
For levels at most $\delta$, the total intensity of old
colors retained anywhere in the list is at most $\delta$.
For those colors use
$|\mathcal B_{r,c}|\le C_b\lambda_c\delta$;
for all others use $C_b\lambda_c\delta^2$.
Since the total color intensity is two,
\[
 \left|\sum_{\lambda_c>0}
        \frac{\mathcal B_{r,c}\mathcal B_{s,c}}{\lambda_c}\right|
                    \le C_b(\delta^3+2\delta^4).
\]
This retains the intensity factor before dividing by
$\lambda_c$ in the Schur complement.
Together with \eqref{eq:upper-unconditioned-covariance},
it proves \eqref{eq:upper-tangent-cutoff}.
\end{proof}

\begin{theorem}[The upper-end covariance tangent]
\label{thm:upper-Gaussian-tangent}
As $\delta\downarrow0$,
\begin{equation}\label{eq:upper-Gaussian-kernel}
 \sup_{u,v\in[0,1]}
 \left|
 \frac{\Gamma(1-\delta u,1-\delta v)}{2\pi_2\delta^2}
                         -\min(u,v)^2\right|\longrightarrow0.
\end{equation}
Consequently
\[
 X_\delta(u)=\frac{G(1-\delta u)}{\sqrt{2\pi_2}\,\delta}
       \Longrightarrow B(u^2)
\]
in finite-dimensional distributions and in $L^2([0,1])$.
Convergence in $C([0,1])$ is not asserted here.
\end{theorem}
\begin{proof}
The tagged tail bound, Lemma~\ref{lem:tagged-shared-tail},
requires a vanishing retained fraction, whereas $K_b$ was
defined by first fixing that fraction.
The following diagonal argument respects this distinction.
Fix $b\ge1$, and take arbitrary $\delta_n\downarrow0$
and $u_n,v_n\in[0,1]$. Put
$r_n=\delta_nu_n$, $s_n=\delta_nv_n$.
If either is zero, its full and cutoff shared counts vanish.
Otherwise choose odd $N_n$ so large that the full and
cutoff covariances divided by $N_n$ approximate
$K(r_n,s_n)$ and $K_b(r_n,s_n)$ within $\delta_n^2/n$,
and $N_nr_n^2,N_ns_n^2\ge n$.
These choices use only fixed-parameter covariance convergence.

Along the resulting sequence both retained counts diverge,
both retained fractions tend to zero, and both
Gaussian endpoint parameters $M^2/N$ diverge.
For centered variables, the tagged tail bound gives
\[
 \left\|
 \frac{(J_N(r_n)-J_{N,b}(r_n))
             -\EE(J_N(r_n)-J_{N,b}(r_n))}{\sqrt N}
 \right\|_2
             \le C r_n/\sqrt b\,(1+o(1)).
\]
The scalar endpoint variance theorem gives
$\|(J_N(r_n)-\EE J_N(r_n))/\sqrt N\|_2
 =(\sqrt\kappa+o(1))r_n$, and similarly for $s_n$.
Bound the cutoff norm by the full norm plus the tail norm
and apply Cauchy--Schwarz to the covariance difference.
Since $r_ns_n\le\delta_n^2$, arbitrary choice of these
sequences proves
\begin{equation}\label{eq:upper-covariance-cutoff-removal}
 \limsup_{\delta\downarrow0}\sup_{u,v\in[0,1]}
 \frac{|K(\delta u,\delta v)-K_b(\delta u,\delta v)|}{\delta^2}
                         \le C/\sqrt b.
\end{equation}
The constant is independent of $b\ge1$.

Combine this with Lemma~\ref{lem:upper-tangent-cutoff},
then let $b\to\infty$.
It gives $K(\delta u,\delta v)/\delta^2\to
\kappa\min(u,v)^2$ uniformly.
Since $\Gamma(1-r,1-s)=4K(r,s)$ and $4\kappa=2\pi_2$,
this proves \eqref{eq:upper-Gaussian-kernel}.
Centered Gaussian finite-dimensional laws are determined
by their covariance matrices, giving the first law limit.

For the Hilbert-space assertion, let $P_m$ replace a
function by its average on each of $m$ equal subintervals.
Uniform covariance convergence gives convergence of the
finite-dimensional Gaussian vectors $P_mX_\delta$ and
\[
 \lim_{\delta\downarrow0}
       \EE\|X_\delta-P_mX_\delta\|_{L^2}^2
   =\EE\|B(u^2)-P_mB(u^2)\|_{L^2}^2
   =\frac1{6m}.
\]
The last equality follows by integrating
$\EE|B(u^2)-B(v^2)|^2=|u^2-v^2|$ within each cell.
Finite-rank approximation, first at fixed $m$ and then
as $m\to\infty$, proves convergence in $L^2([0,1])$.

For the supremum topology one still needs a uniform
modulus estimate, such as
$\EE|X_\delta(u)-X_\delta(v)|^2\le C|u-v|^\alpha$
for some $\alpha>0$, or a uniform shared-edge tail maximum.
The covariance estimate above has a cutoff error independent
of $|u-v|$. Its uniform convergence alone does not supply
that missing modulus.
\end{proof}

\subsection{The explicit count-conditioning correction}
\label{sec:dense-count-correction}

\status{Working proof: the Poisson count contribution in terms of the mean profile}

The conditioned block formula for $\Gamma$ subtracts a count
covariance. That entire correction can be expressed using
$\mathfrak h$ and its derivative, although the remaining
matching covariance is not explicitly evaluated.
For $0<\theta<1$, put
\begin{equation}\label{eq:dense-count-gradients}
 a(\theta)=\frac{\mathfrak h(\theta)
                    +(1-\theta)\mathfrak h'(\theta)}2,\qquad
 c(\theta)=\mathfrak h(\theta)-\theta\mathfrak h'(\theta).
\end{equation}

For the unconditioned model, take independent rate-one old and
fresh Poisson inputs on a circle of circumference $T$, with
independent uniform priorities. At fraction $\theta$, retain
the old points whose priorities exceed $\theta$ and insert
the fresh points whose priorities are at most $\theta$.
Let $H_T^{\rm P}(\theta)$ be the completed spatial matching
distance from the entire old input. Its two input counts
are random. Write $\Gamma_{\rm P}(s,t)$ for the limiting
covariance of these distances divided by $T$.

\begin{proposition}[The covariance removed by fixing the counts]
\label{prop:dense-count-correction}
The limits defining $\Gamma_{\rm P}$ exist. For $0<s\le t<1$,
\begin{align}
 \Gamma_{\rm P}(s,t)-\Gamma(s,t)
   &=2s\,a(s)a(t)+(t-s)c(s)a(t)
                        +(1-t)c(s)c(t).
       \label{eq:dense-count-correction}
\end{align}
In particular, with $v_{\rm P}(\theta)=\Gamma_{\rm P}(\theta,\theta)$,
\begin{equation}\label{eq:dense-Poisson-variance}
 v_{\rm P}(\theta)
   =v_{\rm D}(\theta)+(1-\theta)c(\theta)^2
                                      +2\theta a(\theta)^2.
\end{equation}
The correction is the covariance of
\begin{equation}\label{eq:dense-count-Gaussian}
 Z(\theta)=a(\theta)\{B_0(\theta)+B_1(\theta)\}
             +c(\theta)\{B_0(1)-B_0(\theta)\},
\end{equation}
where $B_0,B_1$ are independent standard Brownian motions.
Thus the unconditioned Gaussian finite-dimensional limits
can be represented as $G+Z$, with $Z$ independent of $G$.
The coefficients extend this process continuously to both
endpoints, with $Z(0)=0$ and
$Z(1)=\{B_0(1)+B_1(1)\}/2$.

Writing $c_{\rm H}=1+e^{-\gamma}$, the unconditioned variance
has the different small-edit asymptotic
\begin{equation}\label{eq:dense-Poisson-small-variance}
 v_{\rm P}(\theta)\sim
       \frac{c_{\rm H}^2}{2}\,
             \theta\log^2(1/\theta)
                         \qquad(\theta\downarrow0).
\end{equation}
This does not evaluate $v_{\rm H}$ or the conditioned
interior variance $v_{\rm D}$.
\end{proposition}
\begin{proof}
For three positive Poisson intensities
$\lambda=(\lambda_0,\lambda_1,\lambda_2)$, let $F_b(\lambda)$
be the mean history distance per unit length through
separation $b$, between the common input of intensity
$\lambda_0$ augmented by either private input.
The finite-component exponential moments justify
differentiating in each intensity. Poisson count
differentiation gives
\begin{equation}\label{eq:dense-mean-count-derivative}
 \lambda_i\partial_iF_b(\lambda)
   =\lim_{T\to\infty}
       \frac{\operatorname{Cov}(H_T^{\le b},N_i)}T,
\end{equation}
where $N_i$ is the total count of color $i$.
For example, at finite $T$ differentiation of the Poisson
count probabilities gives
$\partial_i\EE H_T^{\le b}
 =\operatorname{Cov}(H_T^{\le b},N_i)/\lambda_i$.
The fixed-$b$ component moments, including the seam
component, permit passage to the per-length limit,
uniformly on compact subsets of positive intensities.

We justify differentiation after completion.
The marginal-locality proof of
Theorem~\ref{thm:late-history-variance} applies also when
all three inputs are Poisson. In that case restrictions
to disjoint windows are independent. The bounded cell
scores and exponential localization therefore give
\[
 \operatorname{Var}H_T^{(b,b']}\le C_{\mathcal K}T/b
 \qquad (b'\ge b\ge1,\ \lambda\in\mathcal K),
\]
for every compact positive intensity set $\mathcal K$.
It follows by Cauchy--Schwarz and
$\operatorname{Var}N_i=\lambda_iT$ that
\[
 |\partial_iF_{b'}-\partial_iF_b|
                      \le C_{\mathcal K}/\sqrt b.
\]
Packing bounds the remaining mean density by $C/b$.
Hence $F_b$ and all its first derivatives converge uniformly
on compact sets. The completed density $F$ is $C^1$, and
\eqref{eq:dense-mean-count-derivative} passes to completion.
The same late-variance bound makes the unconditioned
covariances converge and removes the cutoff from their
joint Gaussian laws and their count Schur complements.

Spatial rescaling gives $F(t\lambda)=tF(\lambda)$, and
interchanging the private colors preserves $F$.
At $\lambda=(1-\theta,\theta,\theta)$, write
$a=\partial_1F=\partial_2F$ and $c=\partial_0F$.
Since $F(1-\theta,\theta,\theta)=\mathfrak h(\theta)$,
Euler's identity and differentiation along this curve give
\[
 (1-\theta)c+2\theta a=\mathfrak h,\qquad
 -c+2a=\mathfrak h'.
\]
Solving yields \eqref{eq:dense-count-gradients}.
Homogeneity is used only for the completed density:
the finite-cutoff identity is
$F_b(t\lambda)=tF_{tb}(\lambda)$.

For $s\le t$, split the old priorities into strips of
lengths $s,t-s,1-t$. Their mean-gradient pairs for the two
rewards are respectively
$(a(s),a(t))$, $(c(s),a(t))$, and $(c(s),c(t))$.
The fresh strip of length $s$ has pair $(a(s),a(t))$;
the later fresh strip contributes zero to the first reward.
Each strip count is an independent Poisson count.
Its Schur-complement term is its intensity times the
product of the two gradients, by
\eqref{eq:dense-mean-count-derivative}.
Summing gives \eqref{eq:dense-count-correction}.
Conditioning these finitely many strip counts gives
exactly the nested fixed-size comparison defining $\Gamma$.
The covariance of \eqref{eq:dense-count-Gaussian} is the
same strip sum, and Gaussian projection makes this
count component independent of the conditioned residual.
The diagonal gives \eqref{eq:dense-Poisson-variance}.

Finally, Corollary~\ref{cor:matching-spectrum} gives
\[
 \mathfrak h'(\theta)\sim c_{\rm H}\log(1/\theta),\qquad
 \mathfrak h(\theta)\sim c_{\rm H}\theta\log(1/\theta).
\]
Thus $a(\theta)\sim(c_{\rm H}/2)\log(1/\theta)$.
Concavity gives $0\le c(\theta)\le\mathfrak h(\theta)$,
so its squared contribution is
$O(\theta^2\log^2(1/\theta))$.
The conditioned variance is
$v_{\rm D}(\theta)\sim8v_{\rm H}\theta\log(1/\theta)$.
Equation~\eqref{eq:dense-Poisson-variance} now proves
\eqref{eq:dense-Poisson-small-variance}.
The same endpoint estimates, Brownian H\"older continuity,
and $\mathfrak h(1)=1$, $\mathfrak h'(1)=0$ give the
asserted continuous extension of $Z$.
\end{proof}


\section{Scope and open questions}\label{sec:frontier}
This publication selection retains the mathematical record through Section 7.61.
Its analytic proofs are working arguments awaiting independent human verification.
The finite Lean development is a separate companion, covering the stated finite
process results through \texttt{Stable.lean}. The asymptotic entrance, noise,
history and fluctuation results in this document are ordinary mathematical proofs.
The value near $0.658$ remains an uncertified estimate.
\section{Computational evidence}\label{sec:ledger}
The companion's three interval and rational certificate programs evaluate the
cutoff budget, the stated bounds on the noise constant, and the matching-spectrum
coefficient. Its README gives commands and scope. Earlier numerical observations
retain their reported status; this export makes no fresh verification claim for
unselected historical simulations. The finite Lean package includes its own
reproducible build and per-theorem axiom audit.
\end{document}
